\documentclass[11pt,a4paper]{article}
% Font-map declarations are used only by engines exposing this primitive.
\ifdefined\pdfmapfile
\pdfmapfile{+lm.map}
\pdfmapfile{+cm.map}
\pdfmapfile{+cmextra.map}
\pdfmapfile{+cyrillic.map}
\pdfmapfile{+euler.map}
\pdfmapfile{+latxfont.map}
\pdfmapfile{+symbols.map}
\fi
\usepackage[T1]{fontenc}
\usepackage{lmodern,amsmath,amssymb,amsthm,mathtools,mathrsfs}
\usepackage[margin=27mm]{geometry}
\usepackage{enumitem,booktabs,longtable,array,xcolor,microtype,fancyhdr}
\usepackage{tikz-cd}
\usepackage[colorlinks=true,linkcolor=blue!40!black,citecolor=blue!40!black,urlcolor=blue!40!black]{hyperref}
\setlength{\headheight}{14pt}
\pagestyle{fancy}\fancyhf{}\fancyhead[L]{Concurrent communication and composable security}\fancyfoot[C]{\thepage}
\setlist{nosep,leftmargin=*}
\providecommand{\vmode}{\leavevmode}
\newtheorem{theorem}{Theorem}[section]
\newtheorem{lemma}[theorem]{Lemma}
\newtheorem{proposition}[theorem]{Proposition}
\newtheorem{corollary}[theorem]{Corollary}
\theoremstyle{definition}
\newtheorem{definition}[theorem]{Definition}
\newtheorem{construction}[theorem]{Construction}
\newtheorem{example}[theorem]{Example}
\newtheorem{remark}[theorem]{Remark}
\newcommand{\E}{\mathcal E}
\newcommand{\Set}{\mathbf{Set}}
\newcommand{\Cat}{\mathbf{Cat}}
\newcommand{\op}{\mathrm{op}}
\newcommand{\id}{\mathrm{id}}
\newcommand{\Id}{\mathrm{Id}}
\newcommand{\Path}{\operatorname{Path}}
\newcommand{\Tr}{\mathbf{Tr}}
\newcommand{\Sc}{\operatorname{Sc}_{\square}}
\newcommand{\Un}{\operatorname{Un}_{\square}}
\newcommand{\Reach}{\operatorname{Reach}}
\newcommand{\Conf}{\operatorname{Conf}}
\newcommand{\En}{\operatorname{En}}
\newcommand{\pre}{\operatorname{pre}}
\newcommand{\post}{\operatorname{post}}
\newcommand{\dom}{\operatorname{dom}}
\newcommand{\lab}{\operatorname{lab}}
\newcommand{\Pow}{\mathcal P}
\newcommand{\Pf}{\mathcal P_{\!f}}
\newcommand{\eps}{\varepsilon}
\newcommand{\TCub}{\operatorname{TCub}}
\newcommand{\Cub}{\mathsf C_{\square}}
\newcommand{\cG}{\mathsf G}
\newcommand{\cH}{\mathsf H}
\newcommand{\supp}{\operatorname{supp}}
\newcommand{\obs}{\operatorname{obs}}
\newcommand{\colim}{\operatorname*{colim}}
\DeclareMathOperator{\im}{im}
\DeclareMathOperator{\Hom}{Hom}
\DeclareMathOperator{\Fun}{Fun}
\newcommand{\llbracket}{\mathopen{[\![}}
\newcommand{\rrbracket}{\mathclose{]\!]}}
\newcommand{\crCG}{\mathsf{CG}_{\mathrm{crypt}}}
\newcommand{\crEval}{\operatorname{ev}}
\newcommand{\crPK}{\mathsf{pk}}
\newcommand{\crVK}{\mathsf{vk}}
\newcommand{\crAEnc}{\mathsf{aenc}}
\newcommand{\crSEnc}{\mathsf{senc}}
\newcommand{\crADec}{\mathsf{adec}}
\newcommand{\crSDec}{\mathsf{sdec}}
\newcommand{\crSig}{\mathsf{sig}}
\newcommand{\crMac}{\mathsf{mac}}
\newcommand{\crOK}{\mathsf{ok}}
\newcommand{\crFail}{\mathsf{fail}}
\newcommand{\pcLang}{\pi_{\mathrm{occ}}}
\newcommand{\pcCore}{\mathsf{Core}}
\newcommand{\pcFwd}{\mathsf{Fwd}}
\newcommand{\pcCut}{\mathsf{cut}}
\newcommand{\Max}{\operatorname{Max}}
\newcommand{\CG}{\mathsf{CG}}
\newcommand{\Val}{\mathsf{Val}}
\newcommand{\Proc}{\mathsf{Proc}}
\newcommand{\Env}{\mathsf{Env}}
\newcommand{\Hist}{\mathsf{Hist}}
\newcommand{\ev}{\mathsf{ev}}
\newcommand{\ret}{\mathsf{return}}
\newcommand{\bind}{\mathsf{bind}}
\newcommand{\fail}{\crFail}
\newcommand{\Ty}{\mathsf{type}}
\newcommand{\Int}{\mathsf{Int}}
\newcommand{\Ev}{\mathsf{Ev}}
\newcommand{\Res}{\mathsf{Res}}
\newcommand{\Time}{\mathsf{Time}}
\newcommand{\Lift}{\mathsf{Lift}}

% Local administrative families and cryptographic values.
\newcommand{\PubRes}{\mathcal R}
\newcommand{\PolicyRes}{\mathcal S}
\newcommand{\PolicyCat}{\mathcal A}
\newcommand{\mathbfone}{\mathbf 1}
\newcommand{\Sub}{\operatorname{Sub}}
\newcommand{\Nat}{\operatorname{Nat}}
\newcommand{\Obs}{\mathsf{Obs}}
\newcommand{\Out}{\mathsf{Out}}
\newcommand{\DY}{\mathsf{DY}}

% Simulation and Scott-open information structure.
\newcommand{\Csc}{\mathcal C_{\mathrm{crypt}}}
\newcommand{\Frm}{\mathbf{Frm}}
\newcommand{\Img}{\operatorname{Im}}
\newcommand{\D}{\mathsf D}
\newcommand{\Op}{\mathcal O}
\newcommand{\Pub}{\mathsf{Pub}}
\newcommand{\Att}{\mathsf A}
\newcommand{\Beh}{\mathsf{St}}
\newcommand{\Views}{\mathsf V}
\newcommand{\Sec}{\mathsf{Sec}}
\newcommand{\Sim}{\mathsf{Sim}}
\newcommand{\eqd}{\simeq_{\mathrm d}}
\newcommand{\eqsc}{\approx_{\mathrm{Sc}}}
\newcommand{\Cfg}{\mathsf{Cfg}}
\newcommand{\Exec}{\mathsf{Exec}}
\newcommand{\Trace}{\operatorname{Trc}}
\makeatletter
\renewcommand*\l@subsection{\@dottedtocline{2}{1.5em}{3.1em}}
\renewcommand*\l@subsubsection{\@dottedtocline{3}{4.6em}{4.2em}}
\makeatother
\allowdisplaybreaks
\setlength{\emergencystretch}{3em}

\title{Concurrent Communication and Composable Security\\[4pt]\large Public Information, Protocol Evidence, and Online Simulation}

\author{
  Conor Sheridan\\[0.5ex]
  \href{mailto:siridconor@gmail.com}
       {\texttt{siridconor@gmail.com}}
}
\date{\today}

\begin{document}
\maketitle
\begin{abstract}
When does the information released by an ideal resource suffice to reproduce the public behaviour of a communication protocol? We study this question in concurrent games. A strategy records the occurrences of a protocol, their causal dependencies, and their finite consistency conditions. Payloads range over algebraic domains, including domains of Scott opens. Positive queries have finite witnesses, and singleton payloads recover the original discrete protocols. Secured interaction and hiding compose protocols. Characteristic tests show that two strategies agree on all contextual may observations exactly when they have the same finite displayed traces. Spans of serialized configurations witness this equality and compose through protocol contexts.

A coalition's observation map determines both its public assertions and the internal histories it cannot currently distinguish. One simulator must reproduce every tested interaction using the information available at the ideal ports when each reply is due. For each coalition in a correct conversion, security is equivalent to mutual factorization over the honest interface and to equality of attainable attacker-view images. The simulator witnesses compose sequentially and in parallel against joint attackers. For a continuously initialized deterministic service, future reply behaviour descends exactly when each of its observable predicates pulls back from the initial disclosure. This factorization constructs the simulator. Witness erasure and adaptive threshold disclosure illustrate the roles of public operations and the timing of release.

An occurrence-process language represents every ordinary strategy. The retained histories support distributed control, temporal execution, evidence, and policy. Complete finite display tables admit observational decision and bounded simulator synthesis. For unrestricted effective presentations, no effectively checkable system of finite security certificates is both sound and complete.
\end{abstract}

\clearpage
\begingroup
\small
\tableofcontents
\endgroup
\newpage

% The source is self-contained, including its bibliography and technical appendices.


\section{Introduction}\label{sec:introduction}
A protocol implementation records which input enabled a reply, which occurrence supplied a packet, and which choices could coexist. A participant sees only the events exposed at its ports. Cryptographic security asks whether a simulator interacting with an ideal resource can reproduce what a coalition sees, including the replies to requests the coalition chooses as the interaction proceeds.

We describe communication by receptive, courteous strategies on polarized event structures. Secured interaction matches occurrences at shared ports and checks their joint causal order. Hiding gives the composite strategy. Positive occurrence tags, waiting guards, and finite exclusions yield a process presentation of every ordinary strategy. They retain the distinctions needed, for example, to tell whether a reply can occur before an input or whether three pairwise compatible actions can occur together.

A packet may carry an observation. For an algebraic payload domain $D$, the frame $\Op(D)$ is itself an algebraic domain, and applying a received predicate to a received value has a Scott-open success condition. Section~\ref{subsec:scott-payloads} constructs an ordinary strategy for such readouts. Corollary~\ref{cor:exact-payload-observations} proves that its observable payload predicates are exactly the Scott opens. Section~\ref{subsec:open-streams} extends FIFO and processor observations to streams of these payloads. Exact atomic messages, including entire words, are recovered on discrete sorts. Packet occurrences and their suppliers still determine the order and multiplicity of deliveries. Corollary~\ref{cor:transmit-observation} applies the channel simulation to transmission of a Scott open and identifies the information its scheduling notifications release.

At an exposed interface a test specifies a strategy and a Scott-open success observation. Characteristic tests recognize individual finite displayed traces. Consequently trace equality characterizes contextual may equivalence for the arenas used here. Section~\ref{subsec:complete-trace-testing} gives the comparison with the may-testing theorem of Castellan, Clairambault, and Winskel~\cite{ccw-mays-musts}. A complete trace certificate retains a serialization on each side of every matched word. These certificates compose by pullback and pass through secured contexts and tensor.

A coalition's present information and its capacity to interact give two uses of the same boundary. Pulling Scott opens back along the observation map gives the frame of public assertions. Its kernel pair identifies histories with the same current view. Interactive security asks for one strategy $\sigma$ at the ideal ports, chosen before the honest inputs and public tests. Even when an ideal service eventually discloses a value, a simulator that must answer an earlier request can use only the information released by that time.

The main comparison is Theorem~\ref{thm:categorical-crypto-bridge}. A comparison of the real network with the network using a constructed simulator gives a factorization over the honest interface. Correctness gives the factorization in the opposite direction. Together they characterize equality of attainable attacker-view images and compose against joint attackers. The construction of the simulator is explicit for public reference services. In the initial-disclosure case, the behaviour of the private controller on every future request word descends through the disclosed value precisely when an online reply controller can be initialized from that value alone.

Broadbent and Karvonen~\cite[Theorems~4.6 and~4.9]{crypto-bk} give the categorical composition theorem for secure conversions and a reduction to initial attacks. Delaune, Kremer, and Pereira~\cite{crypto-dkp} give a symbolic account of simulation and protocol composition in the applied pi calculus using an observational preorder. Hayman~\cite{crypto-hayman} uses concurrent games on event structures for causal reasoning about signature exchange. Here the communication model determines the admissible tests, the causal information available to a simulator, and the execution evidence for replacing a real network by its simulated ideal counterpart.

Local views of the same histories determine distributed permissions. Finite return retains the transitions traversed by a computation, while historical evidence records the origin of a public action. The effective results determine when the comparison can be decided from a presentation of its resources.

\subsection{Results and reading order}
Section~\ref{sec:game-foundations} constructs the game category and its occurrence-process presentation. Section~\ref{sec:game-calculus} characterizes observation by finite traces and computes secured execution by configuration kernels. Sections~\ref{sec:protocol-games} and~\ref{sec:execution-games} develop communication histories, distributed control, witnessed return, and generated time.

The security argument in Section~\ref{sec:cryptographic-communication-games} can be read directly after the testing results. It proceeds from resource factorizations to public assertions and private indistinguishability, then to public operation algebras and online simulation. Theorem~\ref{thm:continuous-disclosure} and Corollary~\ref{cor:continuous-disclosure-controller} construct a continuous reply controller from the public disclosure. Theorem~\ref{thm:initial-disclosure-descent} gives the discrete instance. Theorem~\ref{thm:categorical-crypto-bridge} interprets service comparisons and network certificates as secure refinements. Witness erasure in Section~\ref{subsec:crypto-zk} and adaptive disclosure in Theorem~\ref{thm:adaptive-sharing} give examples with different disclosure times. Section~\ref{sec:protocol-evidence-policy} develops evidence and policy over histories. Section~\ref{sec:effective-foundations} proves the finite verification results and the obstruction for unrestricted effective presentations.

The appendices supply the instances and detailed comparisons used along this route. Appendix~\ref{app:communication-examples} constructs the communication media and their history maps. Appendix~\ref{app:cryptographic-examples} constructs channels, adaptive threshold disclosure, proof services, garbling and signed grants, with each conclusion stated at its own interface. Appendix~\ref{app:symbolic-implementation} realizes allocation, initialization and audit; Appendix~\ref{app:quantitative-comparison} explains the boundary of the qualitative symbolic interpretation. Appendix~\ref{app:technical-presentations} supplies the named elaboration, kernel and evidence-mode calculations. All definitions and proofs needed here are included in this paper; no supplemental automata account is required.

\subsection{Conventions and specified interfaces}
\begin{definition}[Ambient notation and categorical carriers]\label{def:formal-ambient-notation-and-categorical-carriers}
\leavevmode
\begin{enumerate}[label=(\roman*)]
\item $\Set$ denotes the category of sets and functions and $\Cat$ the category of small categories and functors at the stated size level. For a category $C$, $C^\op$ reverses arrows, $C(c,d)$ is its hom-set, $\id_c$ its identity, and $ba=b\circ a$ its composition. The category $\Fun(C,D)$ has functors as objects and natural transformations as arrows. The notation $\E$ in the administrative family constructions is instantiated by $\Set$; the game carriers are also sets. Categories and event sets are small in a fixed universe; their presheaf categories may be taken at a larger size level.
\item For a set $X$, $\Pow(X)$ and $\Pf(X)$ are its subsets and finite subsets. A \textbf{partial function} $f:X\rightharpoonup Y$ is a subset $\dom f\subseteq X$ with a function $\dom f\to Y$; $f[A]=\{f(a):a\in A\cap\dom f\}$. Their category is $\mathbf{Par}(\Set)$, with composition on the inverse-image intersection of domains. A \textbf{span} $X\rightsquigarrow Y$ is a pair $X\xleftarrow{s}R\xrightarrow{t}Y$; its chronological composite with $Y\xleftarrow{u}Q\xrightarrow{v}Z$ has carrier $R\times_YQ$ and endpoints $s\pi_R,v\pi_Q$.
\item A \textbf{directed graph} is $(G_0,G_1,s,t)$ with $s,t:G_1\to G_0$. Its \textbf{free path category} $\Path(G)$ has objects $G_0$ and arrows finite composable edge lists, including the empty list at each vertex; composition concatenates lists. For a set $X$, $X^*=\coprod_{n\geq0}X^n$, $X^+=\coprod_{n\geq1}X^n$, $\eps$ is the empty word, and chronological $uv$ executes $u$ before $v$. The \textbf{discrete category} $\operatorname{Disc}(X)$ has objects $X$ and only identity arrows.
\item A \textbf{presheaf} on $C$ is a functor $C^\op\to\Set$. Its \textbf{category of elements} has objects $(c,x)$ with $x\in X(c)$ and arrows $(c,x)\to(d,y)$ given by $f:c\to d$ with $X(f)y=x$. For a covariant $F:C\to\Set$, the covariant category of elements instead has $F(f)x=y$. The \textbf{family completion} $\operatorname{Fam}(C)$ has objects $(c_i)_{i\in I}$ and arrows $(u,(f_i)):(c_i)\to(d_j)$ with $u:I\to J$ and $f_i:c_i\to d_{u(i)}$. This is distinct from the slice notation $\mathsf{Fam}(X)=\Set/X$ used for evidence.
\end{enumerate}
The notation $\pcLang$ for occurrence syntax is introduced in Definition~\ref{def:pc-core}; its medium-indexed form is defined in Definition~\ref{def:pc-medium}. The abstract and introduction preview these subsequent declarations. Display isomorphisms are written $\cong$, operational reductions $\longrightarrow$, and observational equivalences are defined at their specified interfaces.
\leavevmode
\begin{enumerate}[label=(\roman*),start=5]
\item For $H:C\to D$ and $F:C\to B$, a \textbf{left Kan extension} is a functor $L:B\to D$ with $\lambda:H\Rightarrow LF$ for which composition with $\lambda$ induces a bijection $\operatorname{Nat}(L,K)\cong\operatorname{Nat}(H,KF)$, natural in $K:B\to D$. It is denoted $\operatorname{Lan}_F H$. The \textbf{Yoneda embedding} is $y:C\to\Set^{C^\op}$, $c\mapsto C(-,c)$. The extensions used below are constructed by their displayed colimits.
\item For $K:C^\op\times C\to\Set$, its \textbf{coend} $\int^c K(c,c)$ is the quotient of $\coprod_cK(c,c)$ identifying $K(f,1)z\in K(c,c)$ and $K(1,f)z\in K(d,d)$ for every $f:c\to d$ and $z\in K(d,c)$. A map from this quotient is exactly a family of functions out of $K(c,c)$ satisfying these equations. Iterated coends are the corresponding quotients by all the indicated relations.
\end{enumerate}
\end{definition}


\begin{center}\small
\begin{tabular}{@{}p{.21\linewidth}p{.70\linewidth}@{}}
\toprule Notation & Meaning\\\midrule
$\Delta M$, $\Op(D)$ & Discrete payload domain and frame of Scott-open payload predicates.\\
$\mathsf Q_D(I)$, $\mathsf S(D)$ & Positive query interface and domain of streams over $D$.\\
$T_A$, $[w]$ & Serialized trace domain on $A$ and the cylinder generated by $w$.\\

$A^\ast$ & Physical polarity reversal; compact tensor dual.\\
$\Pub_J(r)$ & Frame of assertions observable at coalition $J$.\\
$K_J(r)$ & Kernel pair of the coalition's observation map.\\
$GR$ & Controller-realizability closure on history languages.\\
$\Box_J$ & Greatest public subopen of a history assertion.\\
$\Trace(S)$ & Displayed finite trace language, without forgetting payloads.\\
$\operatorname{tr}^U$, $\mathsf{Ret}_U$ & Compact feedback and witnessed first return, respectively.\\
$\sigma\Vdash_H r\preccurlyeq s$ & Resource comparison witnessed by an online simulator.\\
$M$ & Comparison-certificate carrier, distinct from the simulator.\\
\bottomrule
\end{tabular}
\end{center}
Display isomorphism, finite causal comparison and contextual may equivalence remain distinct. Strategy 2-cells are event maps over a boundary; execution maps retain endpoints and faces; kernel maps are natural transformations; service maps satisfy a command-successor square. Their comparison theorems state the level at which each operates. In particular public current-state equivalence does not assert equality of future interactions. 


\section{Concurrent strategies and their process language}\label{sec:game-foundations}
A protocol implementation records which occurrences can coexist and which must precede others. The order and finite consistency relation of an event structure encode these two constraints. Secured matchings use both to construct interaction. Positive occurrence refinement then supplies a process presentation of the strategy display. We use the concurrent-game framework of~\cite{ccrw}.

\subsection{Events, configurations, and open interaction}
Events record individual occurrences, while configurations record compatible causal histories. The strategy conditions below specify how such histories may realize an interface: negative moves must remain available when the arena enables them, and additional causal waits have a prescribed polarity.

\begin{definition}[Event structures and configurations]\label{def:event-structure}
\leavevmode
\begin{enumerate}[label=(\roman*)]
\item An \textbf{event structure} is a tuple $E=(|E|,\leq_E,\mathsf{Con}_E)$, where $|E|$ is a set, $\leq_E\subseteq|E|^2$ is a partial order and $\mathsf{Con}_E\subseteq\Pf(|E|)$. Its \textbf{principal past} is $[e]_E=\{d\in|E|:d\leq_Ee\}$, required to be finite. The consistency conditions are
\[
\begin{gathered}
 \varnothing,\{e\}\in\mathsf{Con}_E,\qquad
 Y\subseteq X\in\mathsf{Con}_E\Rightarrow Y\in\mathsf{Con}_E,\\
 X\in\mathsf{Con}_E,\ d\leq_Ee\in X\Rightarrow X\cup\{d\}\in\mathsf{Con}_E.
\end{gathered}
\]
\item A \textbf{configuration} is an element of
\[
\begin{gathered}
 \Conf(E)=\{x\subseteq|E|:x=\mathord\downarrow_E x\text{ and }\Pf(x)\subseteq\mathsf{Con}_E\},\\
 \Conf_f(E)=\Conf(E)\cap\Pf(|E|).
\end{gathered}
\]
An event $e$ is \textbf{enabled} at $x$ when $e\notin x$ and $x\cup\{e\}\in\Conf(E)$. A \textbf{causal prefix} of $x$ is a configuration $y\subseteq x$. A \textbf{cover} $e\to_Ee'$ means $e<_Ee'$ with no $d$ strictly between them. A finite set $U\subseteq|E|\setminus x$ is \textbf{jointly enabled} at $x$ when $x\cup V\in\Conf(E)$ for every $V\subseteq U$.
\end{enumerate}
\end{definition}

\begin{definition}[Arenas and their boundaries]\label{def:arenas-boundaries}
An \textbf{arena}, or \textbf{game}, is an event structure $A$ with a polarity function $\mathrm{pol}_A:|A|\to\{-,+\}$. Write $x^\pm=x\cap\mathrm{pol}_A^{-1}(\pm)$. A \textbf{finite position} is $x\in\Conf_f(A)$; a \textbf{move} is a pair $(x,e)$ with $e$ enabled at $x$.
\leavevmode
\begin{enumerate}[label=(\roman*)]
\item The \textbf{parallel arena} $A\parallel B$ has carrier $|A|\amalg|B|$, order within the tagged summands, their inherited signs, and consistency tested in each summand. The \textbf{tensor unit} $I$ has empty carrier.
\item The \textbf{dual arena} $A^\ast$ has the same carrier, order and consistency and polarity $-\mathrm{pol}_A$. The \textbf{boundary from $A$ to $B$} is $A^\ast\parallel B$.
\end{enumerate}
\end{definition}

\begin{definition}[Event maps and ordinary strategies]\label{def:ordinary-strategy}
\leavevmode
\begin{enumerate}[label=(\roman*)]
\item A total \textbf{event-structure map} $f:S\to E$ is a function $|S|\to|E|$ such that $f[x]\in\Conf(E)$ and $f|_x$ is injective for every $x\in\Conf(S)$. A \textbf{prestrategy} on $A$ is such a map $\sigma:S\to A$ with $\mathrm{pol}_A\sigma=\mathrm{pol}_S$; this map is its \textbf{display}.
\item The display is \textbf{receptive} when, for every $x\in\Conf(S)$ and every negative $a\notin\sigma[x]$ with $\sigma[x]\cup\{a\}\in\Conf(A)$, there is a unique $s\notin x$ with $\sigma(s)=a$ and $x\cup\{s\}\in\Conf(S)$.
\item It is \textbf{courteous} when
\[
 s\to_St,\quad (\mathrm{pol}_S(s),\mathrm{pol}_S(t))\ne(-,+)
 \quad\Longrightarrow\quad \sigma(s)<_A\sigma(t).
\]
An \textbf{ordinary strategy} $A\to B$ is a receptive courteous display $S\to A^\ast\parallel B$. Distinct incompatible events of $S$ may have the same displayed label.
\end{enumerate}
\end{definition}
The constructions are set-valued. Section~\ref{sec:effective-foundations} equips their finite evidence with computational presentations.

\subsection{Configuration domains and finite evidence}\label{subsec:configuration-domains}
Later observations may concern infinite configurations, but successful Scott-open observations will have finite witnesses. The configuration-domain calculation establishes the compact basis needed for that passage from semantic observation to finite execution evidence.

\begin{definition}[Directed completeness and compact approximation]\label{def:formal-directed-completeness-and-compact-approximation}
\leavevmode
\begin{enumerate}[label=(\roman*)]
\item A subset $D$ of a poset is \textbf{directed} when it is nonempty and every finite subset has an upper bound in $D$. A \textbf{dcpo} is a poset in which every directed subset has a supremum. A map is \textbf{Scott-continuous} when it is monotone and preserves these suprema.
\item An element $k$ is \textbf{compact} when $k\leq\bigvee D$ for directed $D$ implies $k\leq d$ for some $d\in D$. A \textbf{compact basis} $K_X\subseteq X$ has, for each $x$, a directed lower approximation set $K_x=\{k\in K_X:k\leq x\}$ with supremum $x$. A dcpo with such a basis is \textbf{algebraic}. Empty and unpointed dcpos are permitted.
\end{enumerate}
\end{definition}

\begin{definition}[Configuration domains and ideal completion]\label{def:formal-configuration-domains-and-ideal-completion}
For the event structure $E$ of an arena or strategy constructed in Section~\ref{sec:game-foundations}, write
\[
 \mathsf D(E)=(\Conf(E),\subseteq),\qquad
 K_E=(\Conf_f(E),\subseteq).
\]
An ideal of a poset is a nonempty directed lower subset. Write $\operatorname{Idl}(K)$ for its ideal completion, ordered by inclusion. Directed families here are nonempty; the empty poset has empty ideal completion.
\end{definition}

\begin{proposition}[Configuration domain]\label{prop:effective-configuration-domain}
Directed joins in $\mathsf D(E)$ are unions. Its compact elements are exactly the finite configurations, and there is a canonical order isomorphism
\[
 \mathsf D(E)\cong\operatorname{Idl}(K_E),\qquad
 x\longmapsto\{k\in K_E:k\subseteq x\}.
\]
\end{proposition}

\begin{proof}\vmode
\leavevmode
\begin{enumerate}[label=(\roman*)]
\item \textbf{Directed joins.}
A directed union is down-closed.
Each finite subset of that union lies in one member and is therefore consistent.

\item \textbf{Compactness of finite configurations.}
If a finite configuration lies in a directed union, its finitely many events lie together in one member, proving compactness.

\item \textbf{Finite approximation.}
For any configuration $x$, the finite unions of the finite principal pasts of its events form a directed family of finite configurations with union $x$.

\item \textbf{Characterization of compact elements.}
A compact $x$ must belong to this family and hence is finite.

\item \textbf{Ideal-completion correspondence.}
Conversely the union of an ideal of finite configurations is a configuration.
Its finite subconfigurations already belong to the ideal, by directedness and lower closure.
These two operations are inverse.
\end{enumerate}
\end{proof}

\begin{example}[Causal waits and finite exclusions]\label{ex:causal-waits-and-exclusions}
Let the arena have two independent moves $n^-$ and $p^+$. One strategy offers $p$ immediately. A second has the same finite consistency relation and adds the dependency $n<p$. A test that supplies $n$ only after observing $p$, and reports success upon observing $p$, succeeds with the first strategy. With the second, a matching containing $p$ would require both $n<p$ and $p<n$. No secured matching can therefore report success.

Finite consistency carries independent information. On three positive occurrences one can permit each pair and exclude the triple. Any binary conflict on those occurrences that excludes the triple would also exclude one of its pairs. The process representation in Theorem~\ref{thm:pc-representation} recovers causal waits from covers and exclusions from minimal inconsistent downsets.
\end{example}
\subsection{Payload domains and observation ports}\label{subsec:scott-payloads}
A send has a payload and an occurrence. Its payload may be an ordinary message, information that can be refined, or a predicate that the recipient can apply to another value. The occurrence determines when the recipient acquires access. We describe payloads by algebraic dcpos $D$ and their observable predicates by $\Op(D)$, using the inclusion information order on configuration domains~\cite{effective-domains,pub-observable}. Write $K(D)$ for the compact elements of $D$. A Scott open is an upper set $U$ such that a directed supremum in $U$ has a member in $U$.

\begin{proposition}[Finite evidence and observation-valued payloads]\label{prop:open-payload-basis}
For $U\in\Op(D)$,
\[
 x\in U\quad\Longleftrightarrow\quad
 \exists k\in K(D)\quad k\le x,\ k\in U .
\]
The dcpo $\Op(D)$, ordered by inclusion, is algebraic. Its compact elements are the opens $U_F=\bigcup_{k\in F}\uparrow k$ for finite $F\subseteq K(D)$. The success set of applying a predicate to a value is Scott open in $\Op(D)\times D$.
\end{proposition}
\begin{proof}
The compacts below $x$ form a directed family with supremum $x$. Openness gives the forward implication, and upward closure the converse. Directed joins in $\Op(D)$ are unions. A directed cover of $U_F$ contains its finitely many generators together in one member, hence covers $U_F$ there. The finite unions of generators lying in $U$ form a directed family with union $U$. A compact $U$ therefore equals one of these finite unions. Finally the success set is
\[
 \{(V,x):x\in V\}
 =\bigcup_{k\in K(D)}
   \bigl(\uparrow_{\Op(D)}(\uparrow_D k)\bigr)\times(\uparrow_D k).
\]
Each factor is the principal open of a compact element. This also gives the finite evidence for success when both the predicate and the argument arrive progressively.
\end{proof}

\begin{construction}[Positive observation port]\label{con:positive-observation-port}
For an index set $I$ let $\mathsf Q_D(I)$ have a lane $(i,k)$ for each $i\in I$ and $k\in K(D)$. A lane has a negative request followed by a positive acknowledgement. Different lanes are independent. The state $P_x^I$ for $x\in D$ retains every request and just those acknowledgements whose $k$ satisfies $k\le x$. The inherited causal order makes the display courteous. Every unused request has its unique lift, so the display is receptive. Distinct indices in $I$ permit repeated requests for the same basic observation. For a handle issued during execution, activate its lanes after the issuance notification. Only lanes exported by the declared interface are available to a client.
\end{construction}

At an observation-only boundary the visible events name requests and acknowledgements. Giving every $x\in D$ a separate visible send label would allow the characteristic test of Section~\ref{subsec:complete-trace-testing} to recognize that exact label, even when $\{x\}$ is not open. A private value behind an issued handle therefore determines which replies can occur. Its identity is absent from the handle's public label. Two independently supplied handles for $x\in D$ and $y\in E$ give a handle for $(x,y)$: for a compact pair $(k,l)$ query the two component lanes and acknowledge the pair only after both replies. Both compact tests hold exactly when $(k,l)\le(x,y)$, and separate lanes preserve the ordering of repeated requests.

\begin{proposition}[Implementing a continuous observation]\label{prop:continuous-query-strategy}
A Scott-continuous map $f:D\to E$ has an ordinary strategy on observation ports which, supplied with the corresponding $P_x^{I\times K(E)}$, has the same finite displayed traces as $P_{f(x)}^I$. The strategy is uniform in $x$.
\end{proposition}
\begin{proof}
For output lane $a=(i,l)$ allocate private input lanes $(a,k)$ in $\mathsf Q_D(I\times K(E))$. For each compact $k$ with $l\le f(k)$, the output request enables a query for $k$. Its input acknowledgement enables a possible output acknowledgement of $a$. Make the alternative output acknowledgements of a single lane mutually inconsistent. All other consistency is componentwise. The resulting display into $\mathsf Q_D(I\times K(E))^*\parallel\mathsf Q_E(I)$ is locally injective on configurations. Its extra causal covers run from negative to positive. An external request and each enabled input acknowledgement have unique lifts, including after a different witness for the same output lane has succeeded. Thus it is an ordinary strategy.

After connection to $P_x$, an output acknowledgement of $l$ is possible exactly when $k\le x$ and $l\le f(k)$ for some compact $k$. This is equivalent to $l\le f(x)$ by continuity and compactness. Each finite serialization of $P_{f(x)}$ lifts: between a request and its successful answer insert the private query and acknowledgement for one such $k$. Finite induction suffices, and the input lanes of different output lanes are disjoint. Conversely, hiding a run erases its private queries and leaves a serialization of $P_{f(x)}$. In particular, applying this construction successively to represented inputs interprets composition of continuous maps on their observations.
\end{proof}

\begin{example}[A predicate received as data]\label{ex:receiving-an-open}
Take $D=\Pow(\mathbb N)$ with inclusion. A growing sensor record $x$ lists detected channels. A finite set $F$ witnesses the test $F\subseteq x$. A received predicate $V\in\Op(D)$ may progressively announce finite alternatives for $F$. The receiver can acknowledge as soon as one announced alternative is contained in its record. Proposition~\ref{prop:open-payload-basis} gives exactly these witnesses. Repeating the request in a fresh lane gives a separate acknowledgement occurrence.
\end{example}

\begin{proposition}[Recovery of atomic interfaces]\label{prop:discrete-observation-recovery}
Give a set $M$ the equality order, writing $\Delta M$ for the resulting domain. Every element is compact, $\Op(\Delta M)=\Pow(M)$, and the subposet of singleton opens is isomorphic to $\Delta M$. At an atomic port, relabelling each payload $m$ by the token $\{m\}$, with admissible values precisely these tokens, induces isomorphisms of the original arenas and strategy displays. This applies sort by sort to a signature of atomic data and public operations.
\end{proposition}
\begin{proof}
A nonempty directed subset of $\Delta M$ has only one element. The inverse on singleton tokens returns their unique member. Keep each event, polarity, causal edge, and consistent family, and use the payload bijection on its label. Relabel an operation by conjugation on its arguments and results, fixing control and failure tags. Matching labels in a connection, operation domains, and successor states correspond in both directions. Secured interaction, hiding, tensor, and the finite tests therefore transport along the isomorphisms. Public frames transport by inverse image. The negative requests of Construction~\ref{con:positive-observation-port} belong to a separately specified query interface. Atomic receipt preserves the original send event and its timing.
\end{proof}

\subsection{Finite secured interaction and hiding}\label{subsec:game-secured}
Connecting two displays requires a common history of their shared interface whose causal requirements can be executed together. We construct those histories as secured matchings, then hide the internal events while retaining their causal witnesses. The principal-history lemma first supplies the event structure realizing such a history family; it will also be used for participant and resource histories.

\begin{definition}[Principal-history realization]\label{def:principal-realization}
A \textbf{locally ordered configuration family} consists of $\varnothing\in\mathcal F\subseteq\Pf(L)$ and partial orders $\leq_u$ on $u\in\mathcal F$ whose ideals are exactly $\{v\in\mathcal F:v\subseteq u\}$. Its \textbf{principal-history carrier} is
\[
 |\operatorname{Pr}(\mathcal F)|=\{(a,[a]_u):u\in\mathcal F,\ a\in u\},\qquad
 [a]_u=\{b\in u:b\leq_u a\}.
\]
Equal pairs are one element. Order is inclusion of their principal histories, and a finite set is consistent when all its pairs occur in the principal-history family of one $u\in\mathcal F$. The label function is $(a,[a]_u)\mapsto a$. Lemma~\ref{lem:admin-family} verifies the event-structure and configuration properties of this construction.
\end{definition}

\begin{lemma}[Realizing a family of locally ordered configurations]\label{lem:admin-family}
Let $\mathcal F$ be a family of finite labelled sets containing the empty set. For each $u\in\mathcal F$, let a partial order on $u$ have as its ideals exactly the members of $\mathcal F$ contained in $u$. The principal ideal of a label $a$ is then unchanged when computed in any smaller member containing it. Construct events $(a,[a]_u)$, identify equal such pairs, order them by inclusion of their principal histories, and call a finite set consistent when all its events occur in one $u\in\mathcal F$. This is an event structure whose finite configurations correspond bijectively to $\mathcal F$.
\end{lemma}
\begin{proof}\vmode
\leavevmode
\begin{enumerate}[label=(\roman*)]
\item \textbf{Agreement of local orders.}
We first check that the local orders agree on overlaps, then recover the configurations from their principal histories.
Write $P_u(a)=[a]_u$.
If $v\subseteq u$ lies in $\mathcal F$, then $v$ is an ideal of $u$.
The ideals of the restricted order on $v$ are exactly the members of $\mathcal F$ contained in $v$, which are also the ideals of its specified order.
An order is determined by its ideals, so these orders agree.
In particular $P_v(a)=P_u(a)$ whenever $a\in v$.

\item \textbf{Principal predecessors.}
Put $\Theta(u)=\{(a,P_u(a)):a\in u\}$.
Each $P_u(a)$ belongs to $\mathcal F$.
If $(b,Q)$ is an event with $Q\subseteq P_u(a)$, then $Q$ is an ideal of $u$, and compatibility of the restricted orders gives $Q=P_u(b)$.
It follows that the predecessors of $(a,P_u(a))$ are exactly the events of $\Theta(P_u(a))$.
They form a finite set.

\item \textbf{Antisymmetry.}
Moreover, a principal history has a unique greatest label in its specified order.
Thus two events with the same principal-history set have the same label, and inclusion of histories defines a partial order on events.

\item \textbf{Consistency axioms.}
The empty set and all singletons are consistent, and consistency is hereditary by its common-bound definition.
If a finite set is witnessed by $\Theta(u)$, adjoining a predecessor of any of its events keeps it inside $\Theta(u)$.
This proves the remaining consistency axiom and also shows that $\Theta(u)$ is a configuration.

\item \textbf{Recovery of configurations.}
Conversely, let $C$ be a finite configuration.
Its consistency supplies $u\in\mathcal F$ with $C\subseteq\Theta(u)$.
The labels $v$ of $C$ form an ideal of $u$: whenever $(a,P_u(a))\in C$ and $b\leq_u a$, downward closure puts $(b,P_u(b))$ in $C$.
Hence $v\in\mathcal F$, and agreement of the restricted orders gives $C=\Theta(v)$.

\item \textbf{Inverse correspondence.}
Taking labels and applying $\Theta$ are inverse, proving the bijection and local injectivity.
This gives the principal-history realization of~\cite[Propositions~1--2]{profunctors} and verifies its event-structure properties.
\end{enumerate}
\end{proof}

This principal-history construction is applied only after its ideal description has been proved for the particular diagrams. The same label may occur in different principal histories; the resulting events are then distinct. Where compatible alternative causes invalidate a single prime-history description, Section~\ref{subsec:disjunctive-causal-interfaces} retains those alternatives and constructs their explicit comparison. No prime presentation of an arbitrary history family is presumed.

\begin{definition}[Secured matching pairs]\label{def:formal-secured-matching-pairs}
Let $\sigma:S\to A^\ast\parallel B$ and $\tau:T\to B^\ast\parallel C$ be receptive, courteous strategies. A \textbf{matching pair} $(s,t)$ consists of finite configurations with identical $B$ images. Local injectivity pairs their $B$ events uniquely. Form a finite directed graph by taking their two causal orders and identifying each matched pair of internal events. Its remaining events are visible at $A,C$. Call the matching \textbf{secured} when this graph is acyclic, and order it by transitive \textbf{closure}.
\end{definition}

\begin{lemma}[Secured matchings and their visible past]\label{lem:matching-past}
A matching pair is secured exactly when it can be built from the empty pair by enabled visible extensions and enabled synchronized internal extensions. Its down-closed subgraphs are precisely its matching subconfigurations. If it is secured, the down-closure of its visible events is a canonical matching subconfiguration, obtained by successively removing internal maximal events not in that down-closure.
\end{lemma}
\begin{proof}\vmode
\leavevmode
\begin{enumerate}[label=(\roman*)]
\item \textbf{Execution implies securedness.}
An execution order gives a linear order extending both causal orders, so implies acyclicity.

\item \textbf{Securedness gives an execution.}
Conversely a linear extension of the finite acyclic graph adds a visible event only after its predecessors on its own side, and an internal matched event only after the predecessors on both sides.
Thus every prefix projects to a configuration of each strategy, proving the execution assertion and the downset description.

\item \textbf{Visible causal past.}
The union of the principal downsets of visible events is a downset.

\item \textbf{Removal of unused internal events.}
Any event outside it is internal; if such events remain, choose a maximal one in that finite complement.
It is also maximal in the whole graph, since any successor in the visible down-closure would put it there.
Remove it from both sides.
Repetition leaves exactly the stated down-closure, independently of removal order.
\end{enumerate}
\end{proof}

\begin{definition}[Secured interaction and visible restriction]\label{def:formal-secured-interaction-and-visible-restriction}
The family of finite secured joint graphs is stable under taking ideals, with its inherited order. Applying the principal-history construction of Lemma~\ref{lem:admin-family} constructs their interaction event structure. Its events are the principal matching histories; consistency means joint occurrence in a secured matching. Hide the internal events by restricting to its visible events, retaining their inherited causal order and consistency. This constructs the ordinary composite $\tau\odot_{\mathrm{ord}}\sigma$. A finite configuration of this composite is determined by its visible events with their principal matching histories. Their union is exactly the visible-past matching of Lemma~\ref{lem:matching-past}; conversely such a matching gives that configuration. In particular internal causal witnesses remain in event identities after their labels are hidden.
\end{definition}

\begin{definition}[Matching-history posets and pruning]\label{def:matching-history-posets}
For $\sigma:S\to A^\ast\parallel B$ and $\tau:T\to B^\ast\parallel C$, let $\mathsf M_f(\sigma,\tau)$ be the set of secured matching graphs of Definition~\ref{def:formal-secured-matching-pairs}, ordered by node inclusion. Its \textbf{visible-past subposet} $\mathsf M_f^{\mathrm{vis}}$ consists of $m$ with $m=\mathord\downarrow_m\mathrm{Vis}(m)$. Define
\[
 \mathsf{prune}:\mathsf M_f\to\mathsf M_f^{\mathrm{vis}},\quad
 m\mapsto\mathord\downarrow_m\mathrm{Vis}(m),\qquad
 v:\Conf(T\bowtie_BS)\to\Conf(T\odot S),\quad x\mapsto x\cap\mathrm{Vis}.
\]
The same configuration map is denoted $\mathsf{vis}$ on inclusion-ordered configuration domains. Their well-definedness and comparison are proved in Proposition~\ref{prop:finite-secured-witnesses}.
\end{definition}

\begin{proposition}[Finite witnesses for interaction and hiding]\label{prop:finite-secured-witnesses}
Write $T\mathbin{\bowtie_B}S$ for the interaction event structure just constructed.
\leavevmode
\begin{enumerate}[label=(\roman*)]
\item Let $\mathsf M_f(\sigma,\tau)$ be the finite secured matching graphs, ordered by inclusion of their labelled nodes. Taking the union of principal histories gives an order isomorphism
$\Conf_f(T\mathbin{\bowtie_B}S)\cong\mathsf M_f(\sigma,\tau)$.
\item Let $\mathsf M_f^{\mathrm{vis}}(\sigma,\tau)$ consist of those matchings equal to the down-closure of their visible nodes. The same construction gives
$\Conf_f(T\odot S)\cong\mathsf M_f^{\mathrm{vis}}(\sigma,\tau)$.
\item Restriction to visible events corresponds to pruning to that down-closure:
\[
\begin{tikzcd}[column sep=large,row sep=large]
\Conf_f(T\mathbin{\bowtie_B}S)
 \arrow[r,"v"] \arrow[d,"\cong"']
 &\Conf_f(T\odot S) \arrow[d,"\cong"]\\
\mathsf M_f(\sigma,\tau)
 \arrow[r,"\mathsf{prune}"']
 &\mathsf M_f^{\mathrm{vis}}(\sigma,\tau).
\end{tikzcd}
\]
\end{enumerate}
\end{proposition}
\begin{proof}\vmode
\leavevmode
\begin{enumerate}[label=(\roman*)]
\item \textbf{Interaction configurations.}
A finite matching subconfiguration is precisely an ideal of the joint causal order, by Lemma~\ref{lem:matching-past}.
Lemma~\ref{lem:admin-family} therefore applies.
Its mutually inverse operations take a matching to its principal histories and a configuration of histories to their union.

\item \textbf{Visible configurations.}
For a configuration of visible events, take the union of their full principal matching histories.
It is secured and every node is below a visible one.
Conversely the visible principal histories of such a matching form a configuration of the hidden event structure.
Each operation recovers the other, and both preserve inclusion.

\item \textbf{Hiding and pruning.}
Removing internal events from the event configuration retains exactly the principal histories of its visible events.
Their union is its visible down-closure, proving the square.

\item \textbf{Directed extension.}
All of these operations concern finite histories; infinite configurations are subsequently obtained by their directed unions.
\end{enumerate}
\end{proof}

\subsection{Composition and compact duality}
Secured matching has supplied a candidate composite. We now verify that it remains an ordinary strategy and identify its identity and duality maps. The polarized Scott order records the configuration relation used by asynchronous copycat.

\begin{definition}[Scott configuration order]\label{def:scottorder}
For finite configurations $x,y$, put
\begin{equation}\label{eq:scottorder}
x\sqsubseteq_E y\quad\Longleftrightarrow\quad
x\supseteq^- x\cap y\subseteq^+y.
\end{equation}
The superscripts mean that the events removed from $x$ are negative and those added to obtain $y$ are positive.
\end{definition}

Interaction and hiding are those of Definitions~\ref{def:formal-secured-matching-pairs} and~\ref{def:formal-secured-interaction-and-visible-restriction}: the composite retains the principal causal witnesses of its visible events. \begin{definition}[Asynchronous copycat]\label{def:formal-asynchronous-copycat}
Copycat on $A$ consists of two copies of $A$, with their inherited orders and with each negative occurrence preceding its corresponding positive occurrence. A finite set is consistent when its causal \textbf{closure} projects to consistent sets in both copies. Its configurations are pairs $(x,y)$ with $y\sqsubseteq_A x$, where the Scott order is defined in \eqref{eq:scottorder}.
\end{definition}

\begin{theorem}[Constructed ordinary strategy composition]\label{thm:ordinary-category}
Interaction and hiding preserve receptive, courteous strategies. Their composition is associative up to the canonical witness isomorphism, and asynchronous copycat is its identity. With parallel product and polarity reversal, isomorphism classes of these strategies form a compact closed category.
\end{theorem}
\begin{proof}\vmode
\leavevmode
\begin{enumerate}[label=(\roman*)]
\item \textbf{Comparison with secured interaction.}
We verify the strategy conditions on the retained matching histories, then identify the structural comparisons.
The construction is the ordinary secured interaction of~\cite[Theorem~3.20, Proposition~4.1, and Section~5]{ccrw}.

\item \textbf{Composite display.}
A finite composite configuration determines its visible-past matching by Proposition~\ref{prop:finite-secured-witnesses}.
Its external image is a configuration: each component display preserves configurations, and the external arena is the parallel product of the two external factors.
Local injectivity follows from local injectivity on those factors.
Directed unions give the same statements for arbitrary configurations.

\item \textbf{Existence of negative lifts.}
Let $z$ be a composite configuration and let $m$ be the union of its retained principal matching histories.
Its component projections are configurations, also when $z$ is infinite.
If an external negative move $a$ is enabled at the image of $z$, it is enabled at the displayed configuration of the incident component, since the other boundary factor is independent.
Receptivity of that component gives a unique enabled occurrence above its projection of $m$.
Adjoin it to the matching.
It is a new maximal vertex, so no cycle is introduced.
Every previous vertex remains below an old visible vertex, and the new vertex is visible.
This gives a composite extension.

\item \textbf{Uniqueness of negative lifts.}
For uniqueness, consider any composite extension with label $a$.
Every immediate component predecessor of its negative root has, by courtesy, an arena-predecessor label in the same external factor.
Those labels already occur in the image of $z$.
Local injectivity in the extended matching identifies these predecessors with the corresponding old visible occurrences.
Their retained histories already lie in $m$, so the new root adds no other matching history.
Component receptivity therefore identifies it with the occurrence just constructed.

\item \textbf{Courtesy of visible covers.}
Let $u\to v$ be a visible cover.
In a finite matching history containing it, expand the dependency into a path of component covers.
If $u$ is positive, its first component cover reflects arena order by component courtesy.
Its successor is therefore in the same external arena factor and is visible.
Since there is no intermediate visible event, this successor is $v$, and the required arena order follows.
If $v$ is negative, the same argument applied to the last component cover gives that order.
Thus every cover other than a negative-to-positive cover is reflected, as required.

\item \textbf{Copycat event structure.}
Its order is generated by the arena orders and the negative-to-positive copy edges.
Projecting a path to arena labels gives a nondecreasing path; a cycle would therefore have to consist only of copy edges at one label, where there is only one direction.
Principal pasts are finite because their labels lie in the finite arena past of their root.
The consistency definition then gives the event-structure axioms.

\item \textbf{Copycat strategy conditions.}
At a configuration, an arena-enabled negative occurrence has all its copycat predecessors already present: a path into it must finish with an arena-order edge, and its preceding history is already in the configuration.
Adjoining that occurrence preserves consistency and gives the unique lift.
Every cover is an arena-order edge or a negative-to-positive copy edge, proving courtesy.

\item \textbf{Associativity.}
A composite configuration expands to the joint matching histories of its constituent strategies.
Both bracketings retain precisely the causal past of the ultimate visible events, so their comparison preserves constituent occurrences and principal histories.

\item \textbf{Removal of unused positive events.}
For the right copycat unit, let $b$ be the $B$-image of the source component of a finite matching and let $c$ be its external copycat configuration.
The copycat condition is $c\sqsubseteq_B b$.
Remove the source occurrences over $b\setminus c$, which are positive.
The remainder is a configuration: courtesy forces a cover from a removed positive to stay in the $B$ factor, and downclosure of $c$ forces its successor to be removed as well.

\item \textbf{Completion of pending negative events.}
Next add the missing negative moves in $c\setminus b$ in an arena order.
They are enabled when added and have unique lifts by receptivity.
Courtesy into a negative occurrence determines its immediate predecessors among the already matched arena predecessors.
The resulting source configuration has the same external boundary as the matching.

\item \textbf{Inverse unit correspondence.}
Conversely, from a source configuration insert the matching copycat occurrences at every $B$ event and take the visible past.
These operations are inverse on visible-past matchings.
Removed positive occurrences have no path to a visible event; added negative occurrences are the uniquely forced forwards of pending external negative moves.
Inserting and pruning therefore restores precisely the original retained principal histories.

\item \textbf{Left unit and arbitrary configurations.}
The dual calculation gives the left unit.
These configuration correspondences preserve principal histories and hence induce display isomorphisms; arbitrary configurations follow by directed union.

\item \textbf{Associativity and unit coherence.}
With four factors every associativity route retains the same joint visible past, proving the pentagon.
The unit comparisons insert and remove the same copies, proving the unit coherence.

\item \textbf{Tensor and compact duality.}
Parallel factors partition the joint histories, so these comparisons respect tensor and its symmetries.
Reading copycat as evaluation and coevaluation makes the same insertion-and-removal calculation the snake comparison.
Parallel product and polarity reversal therefore give the asserted compact closed category after passage to isomorphism classes.
\end{enumerate}
\end{proof}

\begin{definition}[Ordinary strategy category]\label{def:formal-ordinary-strategy-category}
We write $\mathsf{CG}$ for this category. For event-generated cubical carriers the tensor is also written $\boxtimes$; its relation to $\parallel$ is established in Section~\ref{sec:game-calculus}. The compact cups and caps are copycat with their boundary assignments changed. Thus compact \textbf{closure} is obtained from the same interaction construction used for communication.
\end{definition}

\subsubsection{Strategy maps and the two categorical levels}\label{subsec:strategy-levels}
Passing to isomorphism classes gives the ordinary strategy category. Some later comparisons also need maps between the underlying displays, so we retain those maps as 2-cells and make explicit how the bicategorical and quotient presentations are related.

\begin{definition}[Strategy hom-categories and adjunctions]\label{def:formal-strategy-hom-categories-and-adjunctions}
We retain the bicategory before taking isomorphism classes. For games $A,B$, let $\mathbb{CG}(A,B)$ have ordinary displays $\sigma:S\to A^\ast\parallel B$ as objects. A morphism $h:\sigma\Rightarrow\tau$ is a total event-structure map $h:S\to T$ with $\tau h=\sigma$:
\[
\begin{tikzcd}[column sep=large]
S \arrow[rr,"{h}"] \arrow[dr,"{\sigma}"'] && T \arrow[dl,"{\tau}"]\\
& A^\ast\parallel B .
\end{tikzcd}
\]
Thus $h$ preserves configurations and is injective on each configuration; it need not be globally injective. Its polarity preservation follows from the triangle. Vertical composition is function composition. The pullback universal property of secured interaction, followed by visible restriction, constructs horizontal composition on these maps. The associators and unitors are those of \cite[Section~4]{ccrw}. In particular composition is a functor between hom categories, and the interchange equation holds. We use the notation
\[
\mathsf{CG}(A,B)=\operatorname{Ob}\mathbb{CG}(A,B)/\cong.
\]
The compact closed category of Theorem~\ref{thm:ordinary-category} is this quotient. Strategy adjunctions later in the paper belong to $\mathbb{CG}$, whereas adjunctions of evidence functors belong to $\mathbf{Cat}$. A strategy adjunction $l:A\to B\dashv r:B\to A$ consists of maps $\eta:1_A\Rightarrow rl$ and $\epsilon:lr\Rightarrow1_B$ satisfying the triangle equations. Its two underlying composites are
\[
\begin{tikzcd}[column sep=large]
l \arrow[r,"{l\eta}"] \arrow[rr,bend right=25,"{1_l}"'] & lrl \arrow[r,"{\epsilon l}"] & l
\end{tikzcd}
\qquad
\begin{tikzcd}[column sep=large]
r \arrow[r,"{\eta r}"] \arrow[rr,bend right=25,"{1_r}"'] & rlr \arrow[r,"{r\epsilon}"] & r .
\end{tikzcd}
\]
Associators and unitors are inserted in these formulas; their coherence is the cited bicategory coherence. Keys are the specified 2-cells. No noninvertible 2-cell is replaced by an equality in $\mathsf{CG}$.
\end{definition}

\begin{lemma}[Horizontal maps on principal matching histories]\label{lem:horizontal-principal-map}
For display maps $h:S\to S'$ and $k:T\to T'$ over the same interfaces,
there are total event maps
\[
 k\bowtie h:T\bowtie_B S\longrightarrow T'\bowtie_B S',
 \qquad k\odot h:T\odot S\longrightarrow T'\odot S'.
\]
If $(r,[r]_m)$ is a principal joint history, map its component
occurrences by $h,k$, form the target matching graph $m'$ with its
target component orders, and put
\begin{equation}\label{eq:horizontal-principal-map}
 (k\bowtie h)(r,[r]_m)=(r',[r']_{m'}).
\end{equation}
The second map restricts this formula to visible roots. Both maps
commute with their boundary displays and preserve identity and
composition of display maps.
\end{lemma}
\begin{proof}\vmode
\leavevmode
\begin{enumerate}[label=(\roman*)]
\item \textbf{Order reflection.}
A total event map reflects order within a configuration: if
$h(s)\leq h(t)$ for $s,t\in x$, the configuration $h([t])$ contains
$h(s)$, and injectivity of $h|_x$ implies $s\leq t$.
The same holds for $k$.

\item \textbf{Securedness of mapped matchings.}
Matching labels are preserved, so the mapped
component configurations still match.
Every edge of their target
joint order lifts to a source relation; a target cycle would therefore
give a source cycle.
Thus $m'$ is secured and its root past exists.

\item \textbf{Configuration preservation.}
For a joint configuration $m$, its component images form a joint
configuration $m'$.
The images under
\eqref{eq:horizontal-principal-map} of its principal events are
exactly all the principal events of $m'$: every target node is the
image of a source node, and order reflection shows that its target
past is computed inside the image of that source node's past.
This proves configuration preservation.

\item \textbf{Local injectivity and boundary.}
Distinct roots in one joint
configuration have distinct images by component local injectivity and
the matching identification, proving local injectivity.
Boundary
labels and visibility are fixed, giving the display equations.

\item \textbf{Compatibility with composition.}
For composable $h,h'$ and $k,k'$, taking a principal past after the
first image and then after the second equals taking it after the
composite image.
Indeed every predecessor of the final root lifts,
by the preceding reflection argument, to the intermediate root past.
Hence
$(k'\bowtie h')(k\bowtie h)=(k'k)\bowtie(h'h)$, and likewise after
visible restriction; identities follow directly.
This proves the
horizontal functor and interchange in this presentation.
\end{enumerate}
\end{proof}

\begin{construction}[Flattening comparisons]\vmode\label{con:joint-flattening-comparisons}
\leavevmode
\begin{enumerate}[label=(\roman*)]
\item \textbf{Common joint-history carrier.}
For a finite composable list of displays, let $\mathsf J$ be its
secured joint graphs pruned to the causal past of the ultimate visible
nodes.

\item \textbf{Expansion and reconstruction.}
For every binary bracketing $b$, expansion of principal matching
histories gives a bijection
$F_b:\Conf_f(S_b)\to\mathsf J$.
Its inverse restricts a joint graph
to each subexpression, takes the principal pasts needed by that
subexpression's outgoing nodes, and regroups them.
The finite matching
and visible-past correspondence proves both inverse equations.

\item \textbf{Comparison on events.}
Thus the comparison from $b$ to $c$ is
$F_c^{-1}F_b$.
On events use this formula on a principal configuration
and take its root.
These event maps preserve consistency and display.

\item \textbf{Coherence equation.}
For any $b,c,d$,
\[
 (F_d^{-1}F_c)(F_c^{-1}F_b)=F_d^{-1}F_b.
\]
This equation proves the associativity coherence, including the
pentagon.

\item \textbf{Copycat unit comparisons.}
Copycat unit comparisons use the configuration adjustments
of ordinary composition, namely receptive negative extensions and
maximal positive removals; their inverse and triangle equations are
the ordinary unit theorem cited above.
The effective construction
below computes those same adjustments.
\end{enumerate}
\end{construction}

\subsection{Evaluation and compact transposition}\label{subsec:semantic-evaluation}
Compact duality lets us present an arrow as a state and express application by connection. These operations will allow the testing and security constructions to use the same strategy in either presentation, with its comparison witnesses retained.

\begin{definition}[Compact evaluation and currying]\label{def:formal-compact-evaluation-and-currying}
Put $[A,B]=A^\ast\otimes B$. Let $\eta_A:I\to A^\ast\otimes A$ be the copycat cup and let $\bar\epsilon_A:A^\ast\otimes A\to I$ be the symmetry-adjusted cap. With associators suppressed, define
\begin{align*}
\operatorname{ev}_{A,B}&=(\bar\epsilon_A\otimes1_B)(1_{A^\ast}\otimes\sigma_{B,A}),\\
\Lambda(f)&=(1_{A^\ast}\otimes f)(\sigma_{\Gamma,A^\ast}\otimes1_A)
                    (1_\Gamma\otimes\eta_A),\qquad f:\Gamma\otimes A\to B.
\end{align*}
\end{definition}
The cap and cup equations give the two commuting diagrams
\[
\begin{tikzcd}[column sep=large]
\Gamma\otimes A \arrow[r,"{\Lambda(f)\otimes1_A}"] \arrow[dr,"{f}"'] &{[A,B]}\otimes A \arrow[d,"{\operatorname{ev}_{A,B}}"]\\
&B
\end{tikzcd}
\qquad
\begin{tikzcd}[column sep=large]
\mathsf{CG}(\Gamma\otimes A,B) \arrow[r,shift left=2,"{\Lambda}"] &
\mathsf{CG}(\Gamma,[A,B]) \arrow[l,shift left=2,"{g\mapsto\operatorname{ev}(g\otimes1)}"].
\end{tikzcd}
\]
Indeed, expanding either composite pairs one cup with one cap; the corresponding snake equation cancels that pair, leaving $f$ or $g$. Naturality follows by composition with maps on $\Gamma$ and $B$. This is the adjunction $(-\otimes A)\dashv[A,-]$ induced by compact duality. \subsection{Guarded occurrence displays and refinements}\label{subsec:semantic-occurrence-displays}
\begin{definition}[Guard-and-exclusion display]\label{def:guard-exclusion-display}
For an arena $A$, $J\subseteq |A|^-\times |A|^+$ and $K\subseteq\Pf(|A|^+)\setminus\{\varnothing\}$, form the directed graph $D$ generated by strict arena order and $J$. The carrier $G_A(J,K)$ consists of those events whose entire $D$-past is finite, acyclic, arena-consistent and contains no member of $K$. Order these events by $D$-reachability. A finite subset is consistent when its $D$-closure is arena-consistent and contains no member of $K$. The display $g_{A,J,K}:G_A(J,K)\to A$ is inclusion, with inherited polarity.
\end{definition}
\begin{proposition}[Ordinary guarded displays]\label{prop:guard-exclusion-display}
The displayed carrier is an event structure and $g_{A,J,K}$ is receptive and courteous. Its finite configurations are precisely the finite $D$-downsets whose arena images are consistent and contain no forbidden cohort.
\end{proposition}
\begin{proof}\vmode
\leavevmode
\begin{enumerate}[label=(\roman*)]
\item \textbf{Surviving event structure.}
A surviving event's predecessors survive and have finite consistent pasts.
Reachability is a partial order on the surviving carrier.
The consistency test contains singletons, is hereditary, and is closed under adjoining predecessors.

\item \textbf{Inclusion display.}
Inclusion preserves configurations and is locally injective.

\item \textbf{Receptivity.}
If a negative arena event $n$ is enabled after the image of a configuration $C$, all generating predecessors of $n$ are arena predecessors and lie in $C$.
Their $D$-pasts also lie in $C$.
Adjoining $n$ creates neither a cycle nor an arena inconsistency nor a positive forbidden cohort.
Thus $n$ survives and has exactly one lift.
This is receptivity.

\item \textbf{Courtesy.}
A cover of the surviving order is either an arena dependency or a generating edge in $J$: a longer generating path would have a surviving intermediate event and would not be a cover.
The latter edges have sign $(-,+)$, giving courtesy.

\item \textbf{Configuration characterization.}
The stated characterization of configurations follows directly from downclosure and the consistency test.
\end{enumerate}
\end{proof}
\begin{definition}[Positive occurrence refinement]\label{def:formal-positive-occurrence-refinement}
A public positive move can have several incompatible realizations. The refinement records those realizations as distinct events. Let $P(A)$ be the positive causal subposet and let $T:P(A)^{\op}\to\Set$ be a presheaf of tags. For any arena event $a$, put
\[
 \widehat T(a)=\lim_{p\leq a,\ p\text{ positive}}T(p).
\]
The indexing order is opposite to the ancestor order. A tuple consists of mutually compatible positive-ancestor tags. The refined arena $A[T]$ has events $(a,t)$ with $t\in\widehat T(a)$; the predecessors of $(a,t)$ are $(b,t|_b)$ for $b\leq a$. A finite set is consistent when its causal \textbf{closure} projects injectively and consistently into $A$. Write $\rho(a,t)=a$.

Negative tags are determined by the positive ancestors already received; they are not independent choices. The tuple construction is the right Kan extension from the positive causal subposet. The following lemmas prove these assertions and construct the ordinary adapter $\pcFwd_\rho:A[T]\to A$ by relabelling the right boundary of copycat. The event map $\rho$ itself is not silently treated as a strategy arrow.
\end{definition}

\begin{definition}[Minimal inconsistent ideals]\label{def:formal-minimal-inconsistent-ideals}
Minimal inconsistent ideals mean inclusion-minimal inconsistent finite downsets, rather than minimal inconsistent arbitrary subsets.
\end{definition}

\begin{lemma}[Stability of a negative lift]\label{lem:pc-negative-stability}
Let $\sigma:S\to A$ be ordinary. If a negative $e$ extends a configuration $x\subseteq y$ and $\sigma(e)$ remains arena-enabled after $\sigma(y)$, then the same $e$ extends $y$.
\end{lemma}
\begin{proof}\vmode
\leavevmode
\begin{enumerate}[label=(\roman*)]
\item \textbf{Lift at the larger configuration.}
Receptivity supplies a lift $f$ after $y$ with $\sigma(f)=\sigma(e)$.

\item \textbf{Predecessors at the smaller configuration.}
Courtesy reflects the label of every immediate predecessor of $f$ strictly below this label.
All these labels already occur in $\sigma(x)$.
Local injectivity in $y$ identifies their occurrences with those in $x$, so the proper past of $f$ lies in $x$.

\item \textbf{Consistency and uniqueness.}
Consistency follows by restriction from $y\cup\{f\}$.
Thus $f$ also extends $x$, and receptive uniqueness there gives $f=e$.
\end{enumerate}
\end{proof}

\begin{lemma}[Intrinsic inconsistent ideals have positive maxima]\label{lem:pc-positive-conflict}
Let $I$ be a minimal inconsistent finite downset of an ordinary strategy. If the display is injective on $I$ and its image is an arena configuration, every maximum of $I$ is positive.
\end{lemma}
\begin{proof}\vmode
\leavevmode
\begin{enumerate}[label=(\roman*)]
\item \textbf{Maximal events and proper ideals.}
A principal past is consistent, so $I$ has at least two maxima.
Put $X=I\setminus\Max(I)$, where $\Max$ denotes the set of maximal elements.
Every proper ideal of $I$ is consistent.

\item \textbf{Exclusion of a negative maximum.}
A negative maximum $e$ would therefore extend $X$.
Its label remains enabled after $I\setminus\{e\}$ because the whole image is an injective arena configuration.
Lemma~\ref{lem:pc-negative-stability} would make $I$ consistent, a contradiction.
\end{enumerate}
\end{proof}

\begin{lemma}[The tag refinement]\label{lem:pc-tag-refinement}
The construction $A[T]$ is an arena with finite principal pasts. Its projection $\rho$ preserves configurations, is locally injective, preserves causal order, and has unique negative lifts. Moreover $\widehat T$ is the right Kan extension of $T$ along the inclusion of positive causal subposets, with opposite variance.
\end{lemma}
\begin{proof}\vmode
\leavevmode
\begin{enumerate}[label=(\roman*)]
\item \textbf{Principal pasts.}
The past of $(a,t)$ has exactly one occurrence $(b,t|_b)$ over each $b\leq a$.
It is finite and projects bijectively to $[a]$.

\item \textbf{Consistency.}
The stated consistency test is hereditary and closed under predecessors.

\item \textbf{Unique negative extension.}
If a negative $a$ is enabled after a projected configuration, all its positive ancestors already have selected tags.
These determine one compatible tuple, and hence its unique consistent extension in $A[T]$.

\item \textbf{Kan extension property.}
For the Kan property, a map from a whole-arena presheaf restricted to positives into $T$ sends an element over $a$ to the tuple of its restricted positive-ancestor images.
Naturality makes this tuple compatible and its components determine the extension uniquely.
\end{enumerate}
\end{proof}

\begin{lemma}[Canonical injective factorization]\label{lem:pc-injective}
For an ordinary $\sigma:S\to A$, let $T(p)$ consist of the positive source occurrences labelled $p$, with restriction maps selecting their positive arena ancestors. The map $i:S\to A[T]$ recording the label and positive-ancestor tuple is globally injective and ordinary, and $\rho i=\sigma$.
\end{lemma}
\begin{proof}\vmode
\leavevmode
\begin{enumerate}[label=(\roman*)]
\item \textbf{Restriction maps.}
The principal source history of an occurrence over $q$ contains a unique occurrence over every arena predecessor $p\leq q$: its image is an arena configuration and the display is locally injective.
Thus the restriction maps exist and compose.

\item \textbf{Separation of positive occurrences.}
Positives are distinguished by their own tags.

\item \textbf{Induction on negative occurrences.}
For negatives, induct on the finite height of the arena label $a$.
Courtesy puts the label of every immediate source predecessor strictly before $a$.
Occurrences over positive strict arena predecessors are fixed by the tuple; occurrences over negative strict arena predecessors are fixed inductively by the restricted tuple.

\item \textbf{Uniqueness of negative occurrences.}
The union of their principal source pasts is exactly the proper past of the negative occurrence: it lies in that past and contains all its immediate predecessors.
Two negative occurrences with the same tuple therefore have the same proper past.
Receptivity identifies them.
This proves global injectivity.

\item \textbf{Configuration preservation.}
Tuples agree on overlaps of a source configuration, so $i$ preserves configurations.

\item \textbf{Receptivity.}
A negative extension in $A[T]$ projects to an enabled negative in $A$; receptivity of $\sigma$ supplies its source lift with exactly the already determined tuple.

\item \textbf{Courtesy.}
Courtesy of $i$ follows because every non-$(-,+)$ cover is reflected in $A$, and the selected ancestry makes that reflection an order relation in $A[T]$.
\end{enumerate}
\end{proof}

\begin{construction}[Refined forwarding]\vmode\label{con:semantic-forwarding}
\leavevmode
\begin{enumerate}[label=(\roman*)]
\item \textbf{Forwarding display.}
For the projection $\rho:A[T]\to A$, take copycat on $A[T]$ and relabel only its right boundary by $\rho$.
Denote the resulting display by $\pcFwd_\rho:A[T]\to A$.

\item \textbf{Display properties.}
Configuration preservation and local injectivity follow from those of $\rho$.

\item \textbf{Receptivity.}
A right negative move lifts uniquely through $\rho$ and then through copycat; a left negative move has its unchanged copycat lift.

\item \textbf{Courtesy.}
Order preservation of $\rho$ and the $(-,+)$ copycat edges prove courtesy.

\item \textbf{Forwarding comparison.}
For the injective factorization $i:S\to A[T]$ of Lemma~\ref{lem:pc-injective}, matching takes place on $A[T]$.
External relabelling changes no secured matching history.
The copycat unit isomorphism consequently gives $\pcFwd_\rho\circ i\cong\rho i=\sigma$.
\end{enumerate}
\end{construction}
\subsection{Prefix arenas and generic operational carriers}\label{subsec:semantic-prefix-carriers}
To interpret processes with local state, we need carriers for prefixes, offers and committed extensions. The following constructions retain the selected occurrences and their dependencies so that operational execution can later be compared with the game display.

\begin{definition}[Event residual carrier]\label{def:event-residual-carrier}
For an event structure $U$ and $h\in\Conf_f(U)$, the \textbf{event residual} $U/h$ has inherited order and polarity and the following event and consistency carriers:
\[
\begin{aligned}
 |U/h|&=\{e\in|U|\setminus h:h\cup[e]_U\in\Conf_f(U)\},\\
 \mathsf{Con}_{U/h}&=\{X\in\Pf(|U/h|):h\cup\mathord\downarrow_UX\in\Conf_f(U)\}.
\end{aligned}
\]
\end{definition}
\begin{lemma}[Residual configurations]\label{lem:event-residual-configurations}
The residual is an event structure and
\[
 \Conf(U/h)=\{x\subseteq|U|\setminus h:h\cup x\in\Conf(U)\}.
\]
\end{lemma}
\begin{proof}\vmode
\leavevmode
\begin{enumerate}[label=(\roman*)]
\item \textbf{Residual event structure.}
A predecessor outside $h$ of a surviving event survives, by restriction of $h\cup[e]_U$.
Its residual past is therefore $[e]_U\setminus h$, which is finite.
The consistency formula contains the empty set and singletons, is hereditary, and is closed under adjoining residual predecessors because those predecessors are already in the displayed $U$-downclosure.

\item \textbf{Insertion of the prefix.}
If $x$ is a residual configuration, $h\cup x$ is downclosed in $U$: predecessors in $h$ are retained and all other predecessors lie in $x$.
Each finite subset is contained in $h\cup\downarrow_U X$ for a finite $X\subseteq x$ and hence is consistent.

\item \textbf{Removal of the prefix.}
Conversely, if $h\cup x$ is a configuration, every event of $x$ survives, its residual predecessors belong to $x$, and each finite $X\subseteq x$ passes the residual consistency test.

\item \textbf{Inverse correspondence.}
These constructions retain the same subset $x$.
\end{enumerate}
\end{proof}

\begin{definition}[Prefix-arena notation]\label{def:prefix-arena-carrier}
For a set $T$, an indexed family of arenas $B:T\to\operatorname{Ob}(\mathsf{CG})$, and $\epsilon\in\{-,+\}$, the \textbf{prefix arena} $\mathsf{Pref}^{\epsilon}_T(B)$ is the polarized event structure constructed in Construction~\ref{con:dependent-prefix-arenas}. Its branch-root occurrence is $r_v$ for $v\in T$; all unions below are tagged unions, so roots and continuation events are distinct.
\end{definition}

\begin{construction}[Dependent prefix arenas]\vmode\label{con:dependent-prefix-arenas}
\leavevmode
\begin{enumerate}[label=(\roman*)]
\item \textbf{Branch family.}
Let $T$ be a set and $(B_v)_{v\in T}$ a family of arenas.
For a sign $\epsilon\in\{-,+\}$ form the arena
$\mathsf{Pref}^{\epsilon}_T(B)$ as follows.

\item \textbf{Events and polarities.}
Its events are the tagged disjoint union
$\coprod_{v\in T}(\{r_v\}+|B_v|)$.
The sign of $r_v$ is $\epsilon$, and continuation events retain their signs.

\item \textbf{Root labels.}
The subject $a$ and payload $v$ are labels of $r_v$ when this arena is used for a named prefix.

\item \textbf{Causal order.}
Within branch $v$, place $r_v$ below every continuation event and retain the order of $B_v$.
There is no order between different branches.

\item \textbf{Consistency and branch exclusion.}
A finite set is consistent exactly when it is empty or contained in one branch and its continuation part belongs to $\mathsf{Con}_{B_v}$.
Thus different roots, and events in different branches, are inconsistent.
\end{enumerate}
\end{construction}

\begin{lemma}[Prefix carrier and configurations]\label{lem:prefix-carrier}
The prefix construction is an arena, with configurations precisely $\varnothing$ and $\{r_v\}\cup x$ for $v\in T$ and $x\in\Conf(B_v)$.
\end{lemma}
\begin{proof}\vmode
\leavevmode
\begin{enumerate}[label=(\roman*)]
\item \textbf{Finite principal pasts.}
The past of a continuation event is its finite past in $B_v$ with the one new root adjoined.

\item \textbf{Order and consistency.}
The order is a partial order, since branches are disjoint and each root is least in its branch.
Consistency contains singletons, is hereditary, and is closed under adjoining predecessors by the same properties in $B_v$.

\item \textbf{Configuration classification.}
A nonempty downclosed consistent set lies in one branch and contains its root; removing that root leaves a configuration of $B_v$.
Conversely adjoining the root to any such configuration gives a downclosed set all of whose finite subsets pass the stated consistency test.
\end{enumerate}
\end{proof}

\begin{definition}[Witnessed medium carriers]\label{def:calculus-medium-carrier}
For a participant set $P$, local action alphabets $(\mathsf{Act}_p)_{p\in P}$, and result-label set $L$, a \textbf{witnessed medium} consists of a store set $K_\Pi$, a family $W_\Pi:K_\Pi\to\Set$, and functions
\[
\begin{aligned}
 I&:\coprod_bW_\Pi(b)\longrightarrow\Pf(P),&
 \mathsf{lab}&:\coprod_bW_\Pi(b)\longrightarrow L,\\
 \mathsf{req}&:\coprod_{b,\omega\in W_\Pi(b)}I(\omega)
       \longrightarrow\coprod_{p\in P}\mathsf{Act}_p,&
 \mathsf{next}&:\coprod_bW_\Pi(b)\longrightarrow K_\Pi.
\end{aligned}
\]
where $\mathsf{req}(b,\omega,p)$ lies in the summand $\mathsf{Act}_p$. A \textbf{medium state} is $b\in K_\Pi$ and an \textbf{operation witness at $b$} is $\omega\in W_\Pi(b)$. Write $I=I(\omega)$, $a_i=\mathsf{req}(b,\omega,i)$, $\ell=\mathsf{lab}(b,\omega)$ and $b'=\mathsf{next}(b,\omega)$. The action alphabets and each concrete store and operation family are constructed at their respective protocol interfaces.
\end{definition}

\begin{definition}[Labelled cube faces and cubical arenas]\label{def:formal-labelled-cube-faces-and-cubical-arenas}
Fix a label set $\Lambda$ with $\mathrm{pol}:\Lambda\to\{-,+\}$. The \textbf{labelled face category} $\mathsf C_\square(\Lambda)$ has objects $(F,\ell)$ with $F$ finite and $\ell:F\to\Lambda$. A morphism $(F,\ell)\to(G,m)$ is $(u,\alpha)$ with $u:F\hookrightarrow G$, $mu=\ell$, and $\alpha:G\setminus u[F]\to\{0,1\}$; composition transports the assigned endpoints, and identity uses the empty assignment. A \textbf{polarized cubical carrier} is a presheaf $X:\mathsf C_\square(\Lambda)^\op\to\Set$. A \textbf{rooted cubical arena} is $(X,i_X)$ with $i_X\in X(\varnothing)$. Duality retains $X$ and reverses $\mathrm{pol}$. For two carriers, parallel tensor is
\[
 (A\boxtimes B)_F=\coprod_{F=F_A\amalg F_B}A_{F_A}\times B_{F_B},
\]
with the inherited labels and faces. The empty direction set gives the unit. Associativity and symmetry retain the same tagged partition of directions. This formula is the Day tensor on the face category: in its defining coend, an injection into a disjoint union partitions the retained directions between the two factors, and the endpoint assignments are absorbed by their face maps.

For a positive direction put $[1]_+=(0\to1)$, and for a negative direction put $[1]_-=(1\to0)$. In the following formulas $\int X$ is the category of cells and face maps of the presheaf $X$.
\end{definition}

\begin{definition}[Anchored instructions]\label{def:anchored-instruction}
For a residual-state set $S$, an \textbf{anchored instruction} records its operation, selected local offers, payloads, consumed or read resource occurrences, and outcome branch. At a state where its guards hold it gives the rule's partial update $a:D_a\hookrightarrow S\to S$. Its residual reuses these selected identifiers, rather than choosing a new occurrence at the same numerical position. Thus a queue removal tests the selected head, exchange tests the selected adjacent pair, and a capacity update tests its selected accept or overflow branch.
\end{definition}

\begin{definition}[Resource occurrence names]\label{def:formal-resource-occurrence-names}
Initial packet names are tagged boundary positions. New names are terms formed from a producer occurrence and its output position. A local producer is indexed by its owner and local step ordinal; a joint producer by its operation and finite map of participating owners to their ordinals. Copy and dispatch children retain their parent and child index. Disjoint constructors distinguish initial, produced, copied, and dispatch occurrences. This constructs exchange-invariant fresh names without selecting an unused atom. Erasure retains the chosen position and operation branch; every enabled unrefined rule has its unique current selection of resource identifiers and its prescribed fresh production names.
\end{definition}

\begin{definition}[Finite resource coefficients]\label{def:resource-coefficient}
At $q\in S$ take an $F$-indexed tuple of distinct enabled anchored instructions, with disjoint participant supports and with no occurrence consumed twice. Evaluate every permutation of this finite tuple. Retain it when all evaluations are defined with the same selected branches, and prefixes with the same executed subset $A\subseteq F$ yield the same named state $q_A$. Call the resulting set $\mathcal C_N(F)$. Equivalently, form the intersection of the permutation evaluation domains by successive pullbacks and then the equalizers of all equal-subset prefix maps. This is a finite-limit construction for each tuple, followed by a coproduct over tuples.
\end{definition}

\subsection{Calculus: interfaces, binding, and process representation}\label{subsec:calculus-foundations}
We now give syntax to the interface and occurrence constructions. The aim is a process representation whose interpretation preserves their existing connection operations. Binding and substitution organize the parameters; declarations then express waits, exclusions and refined occurrences.

\begin{construction}[Section signatures]\vmode\label{con:section-signatures}
\leavevmode
\begin{enumerate}[label=(\roman*)]
\item \textbf{Primitive signature.}
The primitive symbols name the constructed carriers and maps, with their displayed sources and targets.
The signature is extended in the concluding calculus subsection of each subsequent section.

\item \textbf{Witness-indexed symbols.}
Symbols are indexed by their specified values or witnesses.
For example, a particular grant witness indexes a resource operation; it is not replaced by a mere proposition saying that a grant exists.

\item \textbf{Finite syntax and binding.}
Type expressions and raw derivations are finite trees over these symbols and the rules below.
Binding uses finite sorted contexts and their extension by a distinguished variable.
The raw syntax is the union of its finite-depth stages.

\item \textbf{Congruence quotient.}
On each declared ordinary sort, take the intersection of all typed congruences containing the stated computation equations and closed under the constructors and capture-avoiding substitution.
The indiscrete relation at each sort makes this collection nonempty; its intersection is the least such congruence.

\item \textbf{Partial evaluation and soundness.}
Partial evaluations retain their stated success and failure branches, rather than becoming total value operations.
The interpretation proof checks each generating equation in the semantic category; the equations are not additional laws imposed on that category.

\item \textbf{Extension of the signature.}
A later section extends this signature by its new, already constructed operations.
It reuses the previous structural rules and prints only its additional rules and their verification.
This avoids repeating the same substitution proof in every section.
\end{enumerate}
\end{construction}
Value and process judgements retain their distinct contexts and substitutions. The simulator judgement introduced in Section~\ref{sec:public-information-secure-refinement} is an equation about already typed process terms. Its observation judgement ranges over Scott opens and is a semantic specification, with finite execution witnesses when an observation occurs. Concrete service results carry their actual transition records; Section~\ref{subsec:administrative-process-calculus} constructs parameter families and resource rules over protocol histories. No evidence judgement creates an executable operation without a supplied program. The generic substitution calculation is Theorem~\ref{thm:substitution-indexed-interpretation}; each later extension proves its constructor cases.

\begin{definition}[Value formation and constructors]\label{def:value-formation-constructors}
The sorted signature has its declared atomic sorts, channel sorts $\mathsf{Chan}(\tau)$, and the constructor symbols actually present in the relevant section. A value context is a finite list of distinct variables with these sorts. For a constructor $c:\tau_1\times\cdots\times\tau_n\to\tau$, the rules are
\[
 \frac{x:\tau\in\Gamma}{\Gamma\vdash x:\tau}\qquad
 \frac{\Gamma\vdash v_1:\tau_1\quad\cdots\quad\Gamma\vdash v_n:\tau_n}
 {\Gamma\vdash c(v_1,\ldots,v_n):\tau}.
\]
Value-sort and context formation are written explicitly as
\[
 \frac{\tau\text{ a declared carrier sort}}{\tau:\Ty_{\Val}},\quad
 \frac{\tau:\Ty_{\Val}}{\mathsf{Chan}(\tau):\Ty_{\Val}},\quad
 \frac{}{\varnothing\ \mathsf{vctx}},\quad
 \frac{\Gamma\ \mathsf{vctx}\quad\tau:\Ty_{\Val}\quad x\notin\Gamma}
 {\Gamma,x:\tau\ \mathsf{vctx}}.
\]
The channel carrier is the declared sorted name supply. Their interpretation is constructor evaluation. A partial destructor is introduced in Section~\ref{subsec:cryptographic-process-rules}; it is not silently made a total constructor. A \textbf{value substitution} $\theta:\Gamma\to\Val(\Delta)$ is a tuple of terms of the declared sorts, one for every variable of $\Gamma$.
\end{definition}

\begin{definition}[Interface formation]\label{def:interface-formation}
For each arena constructed in the semantic development there is an atomic interface symbol $\underline A$. Formal interface expressions are generated by
\[
 \frac{A\text{ a constructed arena}}{\underline A:\Int}\qquad
 \frac{}{I:\Int}\qquad
 \frac{A:\Int\quad B:\Int}{A\otimes B:\Int}\qquad
 \frac{A:\Int}{A^\ast:\Int}.
\]
Their meanings are respectively that arena, the empty arena, parallel arena, and polarity dual. Write $A\multimap B:=A^\ast\otimes B$. A \textbf{linear interface context} is generated by
\[
 \frac{}{\varnothing\ \mathsf{ctx}}\qquad
 \frac{\Delta\ \mathsf{ctx}\quad A:\Int\quad x\notin\operatorname{names}(\Delta)}
 {\Delta,x:A\ \mathsf{ctx}}.
\]
Coherence maps interpret changes of parenthesization, unit ports, exchange and the canonical dual identifications. No copy or discard of an arbitrary interface variable is a structural rule. Below the fixed value context $\Gamma$ is suppressed.
\end{definition}

\begin{definition}[Typed judgments and their relations]\label{def:calculus-judgment-sorts}
\leavevmode
\begin{enumerate}[label=(\roman*)]
\item A \textbf{value judgment} $\Gamma\vdash v:\tau$ has the sorted value context of Definition~\ref{def:value-formation-constructors}. An \textbf{interface judgment} $\Gamma;\Delta\vdash P:A$ has that value context and the separately owned interface context $\Delta$. At valuation $\rho$, its interpretation has type $\bigotimes\llbracket\Delta\rrbracket_\rho\to\llbracket A\rrbracket_\rho$ in $\mathsf{CG}$. When the value context is empty or fixed, the single context before $\vdash$ is the interface context; letters used for that single context do not turn it into a cartesian value context.
\item A \textbf{scoped-process judgment} $\Gamma;\Xi\vdash P\ \mathsf{proc}$ asserts membership in the finite binding syntax of Definition~\ref{def:formal-sorted-value-contexts-and-process-constructors}, with a separate finite set $\Xi$ of process variables. It does not itself assert an ordinary public-interface typing.
\item The \textbf{structural congruence} $P\equiv Q$ in a stated judgment sort is the least typed congruence generated by its displayed computation equations. A \textbf{witnessed transition} $N\xrightarrow{\ell,d}_\Pi N'$ is a derivation with source, target, label and retained operation witness as in Definition~\ref{def:pc-medium}. Equality of two paths under a constructed commuting cell does not equate either path to an empty execution. Observational equivalence is separately defined at the specified interfaces, in particular Definition~\ref{def:formal-contextual-boundary-equivalence} for contextual Scott observation.
\item A \textbf{Scott-observation judgement} $x\models U$ asserts membership of a configuration in a Scott open; the public instances and their finite evidence are constructed in Section~\ref{subsec:public-observations}. A \textbf{simulator judgement} retains a typed simulator process and its comparison equation. A \textbf{resource judgement} in the administrative application has actual parameter families, a separately allocated resource context, and a decision view, as constructed in Section~\ref{subsec:administrative-process-calculus}. These specification and evidence rules are interpreted in their displayed carriers.
\end{enumerate}
\end{definition}

\subsubsection{Polarized process declarations}\label{subsec:pc-core}
The internal language starts from the same interfaces as the game semantics. A declaration specifies the additional waiting and exclusion decisions made by a process; the arena specifies the moves the environment may offer. We write $\pcLang$ for this occurrence language. Its interface variables and occurrence addresses are distinct from the message names introduced later in this section.

\begin{definition}[Polarized core]\label{def:pc-core}
On an arena $A$, a declaration consists of
\[
 J\subseteq |A|^-\times |A|^+,
 \qquad K\subseteq\Pf(|A|^+)\setminus\{\varnothing\}.
\]
Write $\pcCore[A;J;K]:A$. A waiting edge $(n,p)\in J$ requires the receipt of $n$ before offering $p$; a clause $F\in K$ forbids completing its whole positive cohort. Starting from $C=\varnothing$, the one-shot operational rules are
\[
 \frac{n\in |A|^-\setminus C\quad C\cup\{n\}\in\Conf_f(A)}
 {C\xrightarrow{n}C\cup\{n\}},
 \qquad
 \frac{E_A(C,p)\quad J^{-1}(p)\subseteq C\quad
       \mathsf{safe}_K(C\cup\{p\})}
 {C\xrightarrow{p}C\cup\{p\}}.
\]
Here $E_A(C,p)$ means $p\in |A|^+\setminus C$ and $C\cup\{p\}\in\Conf_f(A)$, and $\mathsf{safe}_K(X)$ means $F\nsubseteq X$ for every $F\in K$. An anchored family $(C,U)$ is a concurrent cell when $U$ is finite, disjoint from $C$, and $C\cup V$ is reachable for every $V\subseteq U$.
\end{definition}

The negative rule is the receptive input skeleton of the declared interface. Thus a declaration offering no positive event can still accept its permitted inputs. It is not the inactive named process $0$. The only extra causal guards are negative-to-positive; same-polarity sequencing and dependencies into negative moves belong to the interface.

\begin{proposition}[Constructed typing of the core]\label{prop:pc-core-sound}
Every polarized declaration constructs an ordinary strategy on $A$. Its finite configurations and cells are exactly the reachable sets and anchored families of Definition~\ref{def:pc-core}.
\end{proposition}
\begin{proof}\vmode
\leavevmode
\begin{enumerate}[label=(\roman*)]
\item \textbf{Ordinary interpretation.}
Interpret $\pcCore[A;J;K]$ by $g_{A,J,K}$ of Definition~\ref{def:guard-exclusion-display}.
Proposition~\ref{prop:guard-exclusion-display} proves that this display is ordinary.

\item \textbf{Reachable configurations.}
Every rule application extends a finite consistent $D$-downset.
Conversely, a linearization of a finite configuration satisfies each displayed rule: the arena and $J$ predecessors are present and no forbidden cohort is completed.
Thus finite configurations are exactly the reachable sets.

\item \textbf{Concurrent cells.}
Applying this characterization to every face $C\cup V$ proves the cell assertion.
\end{enumerate}
\end{proof}


\subsubsection{Process typing, connection, and natural deduction}\label{subsec:pc-typing}
The declarations specify individual strategies. The typing rules now organize their connection and tensor operations, so that a typed process expression denotes the corresponding ordinary strategy and its structural conversions have explicit interpretations.

\begin{definition}[Interface contexts and process judgments]\label{def:formal-interface-contexts-and-process-judgements}
Use the formation rules of Definition~\ref{def:interface-formation} and the judgment sorts of Definition~\ref{def:calculus-judgment-sorts}. An interface context lists separately named occurrences; each connection uses its designated occurrence once. For a single incoming interface, $P:A\to B$ abbreviates that interface judgment. The complete introduction and elimination rules are Definition~\ref{def:interface-introduction-elimination}.

The \textbf{port-bending constructors} are the following two presentations of the compact transposition bijection:
\[
 \frac{\Delta,x:A\vdash P:B}{\Delta\vdash\mathsf{bend}_x(P):A^\ast\otimes B},\qquad
 \frac{\Delta\vdash R:A^\ast\otimes B}{\Delta,x:A\vdash\mathsf{unbend}_x(R):B}.
\]
They expose and reconnect the same port using copycat and introduce no scheduling event. The \textbf{connection abbreviation} $\pcCut_B(P,Q):A\to C$ is $\mathsf{connect}_B(P,Q)$ for $P:A\to B$ and $Q:B\to C$. Name restriction $(\mathsf{new}\ a)P$ instead binds a message name in the scoped process judgment.
\end{definition}

\begin{definition}[Interface introduction and elimination]\label{def:interface-introduction-elimination}
All concatenated contexts below have disjoint interface names.
\leavevmode
\begin{enumerate}[label=(\roman*)]
\item Identity, unit introduction, tensor introduction and connection are
\[
 \frac{}{x:A\vdash x:A}\quad
 \frac{}{\varnothing\vdash *:I}\quad
 \frac{\Delta\vdash P:A\quad\Lambda\vdash Q:B}
 {\Delta,\Lambda\vdash\langle P,Q\rangle:A\otimes B},
\]
\[
 \frac{\Delta\vdash P:A\quad\Lambda,x:A\vdash Q:B}
 {\Delta,\Lambda\vdash\mathsf{cut}_x(P,Q):B}.
\]
The first rule is forwarding; $*$ is the empty interface process. The tensor term is the independent parallel composition.
\item Tensor and unit elimination are
\[
 \frac{\Delta\vdash P:A\otimes B\quad\Lambda,x:A,y:B\vdash Q:C}
 {\Delta,\Lambda\vdash\mathsf{let}\ (x,y)=P\ \mathsf{in}\ Q:C},
\]
\[
 \frac{\Delta\vdash P:I\quad\Lambda\vdash Q:C}
 {\Delta,\Lambda\vdash\mathsf{let}\ *=P\ \mathsf{in}\ Q:C}.
\]
The first bundles the two continuation ports and connects $P$; the second connects the empty port. Neither operation duplicates a port.
\item Function-interface introduction and elimination are
\[
 \frac{\Delta,x:A\vdash P:B}{\Delta\vdash\lambda x.P:A\multimap B}
 \qquad
 \frac{\Delta\vdash F:A\multimap B\quad\Lambda\vdash Q:A}
 {\Delta,\Lambda\vdash\mathsf{app}(F,Q):B}.
\]
The term $\lambda x.P$ is port bending; it is not an input prefix on a message channel. The cup and cap have types $I\to A^\ast\otimes A$ and $A\otimes A^\ast\to I$, with the indicated coherence symmetries.
\item The semantic guarded/excluded occurrence display on $(A,J,K)$ gives the declaration constant
\[
 \frac{A:\Int\quad J\subseteq|\llbracket A\rrbracket|^-\times|\llbracket A\rrbracket|^+
 \quad K\subseteq\mathcal P_f(|\llbracket A\rrbracket|^+)\setminus\{\varnothing\}}
 {\varnothing\vdash\mathsf{Core}(A,J,K):A}.
\]
Its meaning is the display of Definition~\ref{def:guard-exclusion-display}. The finiteness, acyclicity and consistency tests there construct the surviving event carrier; they are not an extra premise on $J,K$ in this rule.
\end{enumerate}
\end{definition}

\begin{definition}[Interface constructor spellings]\label{def:calculus-interface-spellings}
The rules above use $x$ for forwarding $\mathsf{fwd}_x$, $*$ for the empty interface process $0:I$, $\langle P,Q\rangle$ for the typed independent composition $P\mid Q$, and $\mathsf{cut}_x(P,Q)$ for $\mathsf{connect}_x(P,Q)$. Also $\lambda x.P$ is $\mathsf{bend}_x(P)$, and $\mathsf{Core}(A,J,K)$ is $\pcCore[A;J;K]$. These are abbreviations in the interface judgment sort; scoped named syntax retains its own constructors.
\end{definition}

\begin{definition}[Interface computation equations]\label{def:interface-computation-equations}
Linear substitution means connection at the indicated occurrence. In addition to tensor coherence and the associative/unit laws of connection, impose the following proved equations:
\begin{align*}
 \mathsf{app}(\lambda x.P,Q)&\equiv\mathsf{cut}_x(Q,P),&
 \lambda x.\mathsf{app}(F,x)&\equiv F,\\
 \mathsf{let}\ (x,y)=\langle P,Q\rangle\ \mathsf{in}\ R
   &\equiv R[P/x,Q/y],&
 \mathsf{let}\ (x,y)=P\ \mathsf{in}\langle x,y\rangle&\equiv P,\\
 \mathsf{let}\ *=*\ \mathsf{in}\ Q&\equiv Q,&
 \mathsf{let}\ *=P\ \mathsf{in}\ *&\equiv P.
\end{align*}
The function eta equation uses a fresh $x$. Simultaneous substitution in the tensor equation uses separate input contexts for $P$ and $Q$.
\end{definition}

\begin{proposition}[Sound process typing and structural conversions]\label{prop:process-typing-structural-conversions}
The displayed formation, introduction, elimination and computation rules interpret in $\mathsf{CG}$. Connection is secured composition, independent processes use parallel tensor, and forwarding is copycat. Port bending and unbending are mutually inverse up to display isomorphism. Forwarding removal, tensor/unit conversions, reassociation of connection, and permutations of independent connections preserve the interpreted display. In particular,
\[
 \llbracket\pcCut_B(P,Q)\rrbracket
 \cong\llbracket Q\rrbracket\circ\llbracket P\rrbracket,
 \qquad
 \llbracket P\mid Q\rrbracket
 \cong\llbracket P\rrbracket\parallel\llbracket Q\rrbracket.
\]
\end{proposition}
\begin{proof}\vmode
\leavevmode
\begin{enumerate}[label=(\roman*)]
\item \textbf{Structural interpretations.}
A context denotes its ordered tensor.
Identity is copycat, unit introduction is $\id_I$, tensor introduction is the tensor of the two arrows, and connection is secured composition after tensoring by the unused context.
The ordinary-category theorem gives the well-typed composite and its coherent isomorphisms at the display level.

\item \textbf{Compact currying.}
Let $G=\llbracket\Delta\rrbracket$.
The cup and cap give the bijection
\[
 \CG(G\otimes A,B)\cong\CG(G,A^\ast\otimes B).
\]
Its forward direction bends the $A$ port; its inverse evaluates the displayed $A^\ast$ against $A$.

\item \textbf{Function equations.}
Expanding either round trip gives one cup followed by one cap, which cancels by the corresponding snake identity.
This proves both function computation equations.

\item \textbf{Tensor elimination.}
Tensor elimination interprets as $r\circ(p\otimes1)$ with the specified rearrangement of ports.
For $p=p_1\otimes p_2$, this is precisely the composite interpreting $R[P_1/x,P_2/y]$.

\item \textbf{Tensor and unit equations.}
For $R=\langle x,y\rangle$, it is the identity on the tensor output.
The unit calculation is the same calculation with the unitors.

\item \textbf{Independent substitution.}
Bifunctoriality proves permutations of independent substitutions.

\item \textbf{Soundness and quotient.}
The guarded/excluded event construction supplies the interpretation of each core constant.
Induction over the remaining derivation now gives soundness.
Proposition~\ref{prop:basic-interpretation-principle} gives the quotient interpretation.

\item \textbf{Application comparison.}
The beta calculation is the commuting triangle
\[
\begin{tikzcd}[column sep=large]
 G\otimes A\arrow[r,"\Lambda(f)\otimes1"]\arrow[dr,"f"']
 &(A^\ast\otimes B)\otimes A\arrow[d,"\mathsf{ev}_{A,B}"]\\
 &B.
\end{tikzcd}
\]
The triangle is the evaluation comparison of Definition~\ref{def:formal-compact-evaluation-and-currying}; Theorem~\ref{thm:ordinary-category} supplies the display-level structural comparisons.
\end{enumerate}
\end{proof}

\paragraph{Evaluation, application, and tensor elimination}\label{par:compact-evaluation-tensor-elimination}
\begin{definition}[Application and tensor-elimination interpretation]\label{def:formal-application-and-tensor-elimination-syntax}
For the constructors of Definition~\ref{def:interface-introduction-elimination}, write $[A,B]=A^\ast\otimes B$ and use the evaluation map of Definition~\ref{def:formal-compact-evaluation-and-currying}. Their recursive interpretations are
\[
 \llbracket\mathsf{app}(f,a)\rrbracket
 =\operatorname{ev}_{A,B}(\llbracket f\rrbracket\otimes\llbracket a\rrbracket),\qquad
 \llbracket\mathsf{let}\ (x,y)=p\ \mathsf{in}\ k\rrbracket
 =\llbracket k\rrbracket(\llbracket p\rrbracket\otimes1),
\]
with the port symmetries specified by the disjoint premise contexts.
\end{definition}
Contexts are disjoint in these rules. No diagonal or weakening of an owned interface is implicit. The equations just proved are structural equations; equality of arbitrary declaration constants is the displayed-event-structure equality decided on the finite fragment in Proposition~\ref{prop:pc-normalization}.

The detailed named-prefix, binding and scope elaboration is given in Appendix~\ref{subsec:pc-named}; it constructs the named syntax used in the protocol examples. The core declarations above already present the occurrence displays needed for the following readback.

\subsubsection{Positive refinement and process readback}\label{subsec:pc-refinement-proof}
The semantic factorization and forwarding construction are in Section~\ref{subsec:semantic-occurrence-displays}. The remaining readback identifies the process presentation.
\begin{proposition}[Readback and forwarding]\label{prop:pc-readback}
Identify $S$ with its image under $i$. Let $J$ consist of source covers absent from the refined arena and let $K$ consist of $\Max(I)$ for the intrinsic minimal inconsistent source ideals. Then $\pcCore[A[T];J;K]$ has exactly the configurations of $S$. Relabelled copycat gives an ordinary $\pcFwd_\rho$ and $\pcFwd_\rho\circ i\cong\sigma$.
\end{proposition}
\begin{proof}\vmode
\leavevmode
\begin{enumerate}[label=(\roman*)]
\item \textbf{Source configurations satisfy the core.}
Every source configuration respects its arena, $J$, and $K$.

\item \textbf{Lifting a core step.}
Conversely, induct on a core run already in $S$.
A negative step has its unique source lift.
A positive step is an occurrence of $S$, since its tag is such an occurrence.
Its source predecessors are present by the arena and $J$ tests.

\item \textbf{Exclusion of inconsistent extensions.}
If adjoining it were inconsistent, the resulting finite downset would contain a minimal inconsistent ideal $I$.
Its image is an arena configuration, so $I$ is intrinsic and its maximal cohort belongs to $K$, contradicting the rule.

\item \textbf{Recovery of the display.}
Thus the configurations coincide, as do the reconstructed order, consistency and cells.
Extra negative occurrences in the refined arena remain unreachable.

\item \textbf{Forwarding display.}
For forwarding, take copycat on $A[T]$ and relabel only its right boundary by $\rho$.
Configuration preservation and local injectivity follow from those of $\rho$.

\item \textbf{Forwarding strategy conditions.}
Negative lifts on the right use its receptivity, followed by copycat; lifts on the left are the unchanged copycat lifts.
Courtesy follows from order preservation and the $(-,+)$ copycat edges.

\item \textbf{Copycat unit comparison.}
In composition with $i$, the matching interface is still $A[T]$.
External relabelling changes no secured matching history.
The copycat unit isomorphism therefore identifies the source with $S$ and its external display with $\rho i=\sigma$.
\end{enumerate}
\end{proof}

\subsubsection{Occurrence refinement and representation}\label{subsec:pc-representation}
Readback has identified the guards and exclusions of an injective refined display. Combining it with positive refinement gives the representation theorem for an arbitrary ordinary display.


\begin{theorem}[Constructive process representation]\label{thm:pc-representation}
Every ordinary display $\sigma:S\to A$ has a constructed presentation
\begin{equation}\label{eq:pc-representation}
 \sigma\cong\pcFwd_\rho\circ\pcCore[A[T];J;K].
\end{equation}
The construction is finite when $S$ and $A$ are finite. Set-indexed declarations represent all ordinary strategies. Consequently the interpreted core language is surjective on the arrows of $\mathsf{CG}$ up to display isomorphism.
\end{theorem}
\begin{proof}\vmode
\leavevmode
\begin{enumerate}[label=(\roman*)]
\item \textbf{Positive refinement.}
Tag a positive arena event by its actual positive source occurrences.
Their unique positive arena ancestors give the restriction maps.
Lemma~\ref{lem:pc-injective} gives an injective ordinary display $i:S\to A[T]$.

\item \textbf{Waiting guards.}
Take $J$ to be its non-arena covers.
Courtesy gives their sign $(-,+)$.

\item \textbf{Forbidden cohorts.}
Take $K$ to be the maximal cohorts of its intrinsic minimal inconsistent downsets.
Lemma~\ref{lem:pc-positive-conflict} makes them positive.

\item \textbf{Readback and forwarding.}
The readback argument of Proposition~\ref{prop:pc-readback} identifies the core with $i$, and the forwarding unit calculation gives $\pcFwd_\rho i\cong\sigma$.

\item \textbf{Arrow representation.}
Apply this construction to $A^\ast\parallel B$ and transpose to obtain the arrow statement.
\end{enumerate}
\end{proof}

The chosen generators are relatively minimal at a fixed occurrence set and interface. Removing a non-arena cover permits its target without the required predecessor; removing an intrinsic minimal forbidden ideal permits that ideal. This minimality leaves open encodings that introduce private events or change the observation.

The size of the syntax is a separate issue. A finite syntax over countably many names cannot describe the unrestricted category: even one positive public move permits arbitrarily large sets of pairwise conflicting source realizations.

We use display isomorphism as equality of interpreted core terms. Section~\ref{subsec:pc-equations} gives the finite normalization calculation. Operational observations of named extensions are introduced with those extensions.

\paragraph{Session-process comparison.}\label{par:pc-session-comparison}

The session calculi of Caires--Pfenning~\cite{calculus-cp} and Wadler~\cite{calculus-wadler} organize intuitionistic and classical linear session types by introduction and cut-elimination rules. Here compact transposition supplies the interface function rules, and explicit prefix arenas supply the named prefix rules. The scoped operational language also retains the constructed choices, stores and unrestricted recursive syntax. Comparisons with session models therefore use the interfaces and observation maps stated in those comparisons.

For a strategy in the ordinary overlap of the session model of \cite{cy}, Theorem~\ref{thm:pc-representation} constructs its core declaration and readback returns that same display. The finite-definability comparison in \cite[Theorem~6.10]{cy} applies on its stated internal session interfaces; its testing theorem uses its own semantic observation. The general ordinary category here also permits finite-cohort exclusions and interface objects not generated by that session grammar.

\subsubsection{One process language through the constructions}\label{subsec:pc-stages}
Readback and typed connection now give a process presentation of ordinary strategies with their occurrence witnesses retained. The finite instance is stated below; its normalization and observational decision procedures are supplied in Section~\ref{sec:effective-foundations}. Subsequent extensions keep these interface and connection operations while adding the syntax required by their particular constructions.

\begin{corollary}[Finite process presentation of concurrent strategies]\label{cor:pc-finite-presentation}
On finite arenas, finite typed occurrence terms with equality computed by Proposition~\ref{prop:pc-normalization} present the category of finite ordinary strategies up to display isomorphism. The presentation preserves independent composition, connection, and compact duality.
\end{corollary}
\begin{proof}\vmode
\leavevmode
\begin{enumerate}[label=(\roman*)]
\item \textbf{Finite interpretation and representation.}
Finite declarations, refined forwarding, tensor and secured connection construct finite ordinary displays.
Theorem~\ref{thm:pc-representation} constructs a finite term for every finite display.

\item \textbf{Normalization and faithfulness.}
The normalization procedure compares the constructed finite order, consistency and boundary labelling by a finite bijection search; it identifies two terms exactly when their interpreted displays are isomorphic.
Hence interpretation is full and faithful on the resulting hom-sets and is the identity on arena objects.

\item \textbf{Structural preservation.}
Proposition~\ref{prop:process-typing-structural-conversions} and the copycat cup and cap identify the stated operations.

\item \textbf{Unrestricted representation.}
For arbitrary ordinary strategies the set-indexed readback theorem supplies representation; the finite decision procedure has its stated finite domain.
\end{enumerate}
\end{proof}

The subsequent extensions keep these interface and connection operations. Each adds particular syntax, operational rules, or judgement indices, and constructs their game and categorical interpretations at the same point:
\begin{center}\small
\begin{tabular}{@{}p{.12\linewidth}p{.38\linewidth}p{.41\linewidth}@{}}

\toprule Section & Process construction & Game and categorical construction\\\midrule
\ref{sec:game-calculus} & Execution evidence and tests & Finite traces, comparison certificates and secured kernels\\
\ref{sec:protocol-games} & Medium rules and permission offers & Protocol histories, control compilation and winning transfer\\
\ref{sec:execution-games} & Recursive services and routed execution & Finite return, relational feedback and time\\
\ref{sec:cryptographic-communication-games} & Calls, strict operations and simulator terms & Public frames, service descent and secure conversions\\
\ref{sec:protocol-evidence-policy} & Evidence and policy programs & Evidence adjunctions and retained histories\\
\ref{sec:effective-foundations} & Coded realizers and finite certificates & Effective lifts, decision procedures and synthesis bounds\\
\bottomrule
\end{tabular}
\end{center}
For a protocol with compatible alternative causes, the physical history and its causal refinement retain their specified comparison map. Named operational elaboration, history construction, and core readback are the successive constructions connecting that instance to the three presentations. Their witness comparisons are proved with the instances, and transported through execution in Section~\ref{sec:execution-games}.


We can therefore pass between an ordinary display and its occurrence-process presentation while retaining its boundary, causal waits and consistency witnesses. Connection and tensor have the same interpretation on both sides. The next question is observational: which of these distinctions can a test at the exposed interface detect?

\section{Interaction, observation, and trace certificates}\label{sec:game-calculus}
Which differences between two protocol implementations can a test detect at their common interface? Characteristic tests reduce contextual may comparison to inclusion of finite displayed trace languages. Certificates retain the configurations and linearizations witnessing a match. Configuration kernels calculate secured execution with its causal information.

\subsection{Contextual may observation}\label{subsec:contextual-may}
An observation asks whether some execution can reach a specified open set at the final exposed boundary. Quantifying this question over contexts gives the comparison that the subsequent characteristic tests will characterize.

\begin{definition}[Contextual Scott equivalence]\label{def:formal-contextual-boundary-equivalence}
For $t:I\to O$, write $t\models U$ when some configuration of its display has boundary in the Scott open $U\subseteq\D(O)$. Equally typed strategies satisfy $S\leq_{\mathrm{may}}T$ when every well-typed one-hole context $C[-]$ built from composition and tensor, and every Scott open $U$ of its final arena, satisfy
\[
 C[S]\models U\quad\Longrightarrow\quad C[T]\models U.
\]
Write $S\eqsc T$ for the two directions. A context uses only the declared exposed ports. This positive may observation compares interactive behaviour, not merely current public configurations; its domain-theoretic basis is~\cite{pub-observable}.
\end{definition}

\begin{theorem}[Contextual congruence and calculus correspondence]\label{thm:contextual-congruence-calculus}
The relation $\eqsc$ is a symmetric monoidal congruence. The quotient $\CG/{\eqsc}$ is compact closed, with the inherited cups and caps. Set-indexed occurrence-process tests give exactly this semantic equivalence. The same conclusions hold for the declared interface subcategories with their inherited ambient tests.
\end{theorem}
\begin{proof}\vmode
\leavevmode
\begin{enumerate}[label=(\roman*)]
\item \textbf{Contextual congruence.}
Substituting precomposition, postcomposition or tensor with a fixed strategy into a test context proves compatibility with those operations.

\item \textbf{Quotient structure.}
Hence the quotient operations are defined and inherit their equations, including the snake identities.
Dualization is expressed using cups, caps and tensor and also descends.

\item \textbf{Agreement of process and semantic tests.}
Readback presents every fixed strategy in a context, while interpreted process contexts are semantic contexts.
The finite execution evidence of Theorem~\ref{thm:evidence} identifies their success predicates.
Thus the two testing quantifications agree, without restricting arbitrary adversaries to finite programs.
\end{enumerate}
\end{proof}

\subsection{Complete finite-trace testing}\label{subsec:complete-trace-testing}
For a state $S:I\to A$, a test consists of a strategy $t:A\to O$ and a Scott open $U\subseteq\D(O)$. It succeeds when $tS\models U$. We construct such a test for each finite displayed trace and then show that trace inclusion suffices for success to pass from one strategy to another in every context.

For an arena $A$, let $\mathsf W(A)$ consist of words of distinct actual arena events whose underlying set is a finite configuration and whose order extends arena causality. Payloads and occurrence labels remain part of the letters. For a display $S\to A$, put
\[
 \mathsf E(S)=\coprod_{x\in\Conf_f(S)}\operatorname{Lin}(x),\qquad
 \lambda_S:\mathsf E(S)\to\mathsf W(A),\qquad
 \Trace(S)=\im\lambda_S.
\]
Here $\operatorname{Lin}(x)$ is the set of linearizations of the source order, and $\lambda_S$ applies the display letterwise. A trace forgets concurrency witnesses, but $\mathsf E(S)$ retains them. This separates the testing quotient from the witnessed execution presentation.

\begin{lemma}[Ordinary characteristic tests]\label{lem:characteristic-test}
For each $w\in\mathsf W(A)$ there is an ordinary test $\tau_w:A\to\Omega$, with one positive success event $\omega$ in $\Omega$, such that
\[
 \tau_wS\models\mathord\uparrow\{\omega\}
       \quad\Longleftrightarrow\quad w\in\Trace(S).
\]
\end{lemma}
\begin{proof}\vmode
\leavevmode
\begin{enumerate}[label=(\roman*)]
\item \textbf{Characteristic test.}
Write $P_w,N_w$ for the positive and negative letters in $w$.
On $A^\ast\parallel\Omega$ use the guard-and-exclusion display with
\[
\begin{aligned}
 J_w&=\{(p,n):p\in P_w,n\in N_w, p\text{ precedes }n\text{ in }w\}
       \cup\{(p,\omega):p\in P_w\},\\
 K_w&=\{\{n\}:n\in A^-\setminus N_w\}.
\end{aligned}
\]
All guards have negative-to-positive polarity at this test boundary, and all excluded singletons are positive there.

\item \textbf{Strategy conditions.}
Proposition~\ref{prop:guard-exclusion-display} makes this an ordinary strategy, also for infinite $A$.

\item \textbf{Success from a trace.}
A source configuration linearized by $w$ matches a successful test configuration: their joint order is extended by $w\omega$ and is secured.

\item \textbf{Restriction of a successful matching.}
Conversely, take a finite successful matching.
Its source configuration $x$ contains every letter in $P_w$ and no negative letter outside $N_w$.
Restrict it to letters of $w$.
This restriction is down-closed: at a first source cover from an omitted event to a retained one, the omitted event is positive; courtesy forces its label to be an arena predecessor of a retained label, contradicting downclosure of the underlying arena configuration of $w$.

\item \textbf{Completion of missing inputs.}
Add any missing negative letters in their $w$ order by receptivity.
Each is arena-enabled when added, and no newly added event can be a predecessor of an event already present.
The resulting configuration has exactly the letters of $w$.

\item \textbf{Recovery of the trace order.}
Every source cover except a negative-to-positive one agrees with the arena order and hence with $w$.
A negative-to-positive cover $n<p$ contrary to $w$ between originally retained events would meet the test edge $p<n$ and create a cycle.
A newly added negative cannot precede an already retained positive.
Thus every cover agrees with $w$, so $w$ linearizes the resulting configuration.
\end{enumerate}
\end{proof}

\begin{theorem}[Full abstraction for may observation]\label{thm:trace-full-abstraction}
For ordinary displays on the same arena,
\[
 S\leq_{\mathrm{may}}T\quad\Longleftrightarrow\quad
 \Trace(S)\subseteq\Trace(T),\qquad
 S\eqsc T\quad\Longleftrightarrow\quad\Trace(S)=\Trace(T).
\]
For arrows, apply this statement to their compact names on $A^\ast\parallel B$. The quotient is exactly the category of realizable trace languages, with composition induced by secured interaction.
\end{theorem}
\begin{proof}\vmode
\leavevmode
\begin{enumerate}[label=(\roman*)]
\item \textbf{Trace inclusion from testing.}
The forward implication follows by testing with every $\tau_w$.

\item \textbf{Finite contextual witness.}
For the reverse, take a successful contextual observation of $S$.
Algebraicity reduces it to a finite configuration, and secured composition represents it by a finite joint matching graph.

\item \textbf{Replacement by an equal trace.}
Linearize that graph.
Its restriction to the $S$ component is a word $w\in\Trace(S)$; choose a configuration of $T$ with that same word.
Replace the component by this configuration, identifying the same shared boundary letters.
The original global word, with the new occurrences in the same positions, extends every component order and therefore secures the replacement.

\item \textbf{Preservation of success.}
Hiding preserves its exposed configuration, which still belongs to the selected open.
The construction works for each finite context network, tensor and arbitrary Scott-open success.

\item \textbf{Equality and compact names.}
The reverse inclusion proves equality.
Compact naming and un-naming are contexts, giving the arrow statement.

\item \textbf{Composition of trace languages.}
More explicitly, a trace of $gf$ is the visible restriction of a finite joint word whose restrictions to the $A^\ast,B$ and $B^\ast,C$ boundaries are traces of $f,g$, identifying the shared $B$ events.
Such a word makes the matching secured.
Conversely linearize a finite secured representative.
If it includes hidden events outside the visible causal past, pruning them leaves source subconfigurations and the same visible word.

\item \textbf{Tensor and well-definedness.}
Tensor takes tagged shuffles.
These operations are well defined on trace equality by the established contextual congruence, and only languages realized by ordinary displays are included.
\end{enumerate}
\end{proof}

Theorem~10.3 of Castellan, Clairambault and Winskel~\cite{ccw-mays-musts} states the finite-trace characterization for stopping strategies on race-free games; their subsequent discussion explains why stopping information and neutral test events are unnecessary for the may result. The argument above establishes the characterization for the ordinary strategies and arenas defined here, using guard-and-exclusion tests and arbitrary Scott-open success observations. In particular, its converse supplies missing negative letters by receptivity rather than requiring the successful matching to contain the entire target word. Must testing uses additional stopping information. Winning conditions and temporal prefixes below retain their explicit execution carriers.

\subsubsection{Scott opens on serialized histories}\label{subsec:open-trace-testing}
Order the legal finite displayed words $\mathsf W(A)$ by prefix and form the trace domain $T_A=\operatorname{Idl}(\mathsf W(A))$. A compact point is $\downarrow w$ for a legal finite serialization $w$. Its cylinder is $[w]=\{I:w\in I\}$. For an ordinary display $S$ define its positive observation on $U\in\Op(T_A)$ by
\[
 \mathsf{May}_S(U)\quad\Longleftrightarrow\quad
 \exists v\in\Trace(S)\quad \downarrow v\in U .
\]

\begin{theorem}[Complete Scott-open trace observation]\label{thm:open-trace-testing}
For displays $S,T$ on $A$,
\[
 S\leq_{\mathrm{may}}T
 \quad\Longleftrightarrow\quad
 \forall U\in\Op(T_A)\;
       \bigl(\mathsf{May}_S(U)\Rightarrow\mathsf{May}_T(U)\bigr).
\]
Mutual implication is contextual Scott equivalence. The characteristic test $\tau_w$ implements the cylinder $[w]$. In particular open observations recover the original finite-word comparison.
\end{theorem}
\begin{proof}
The initial segment of a source serialization is a serialization of a subconfiguration. Hence $\Trace(S)$ is prefix closed and $\mathsf{May}_S([w])$ holds exactly when $w\in\Trace(S)$. Lemma~\ref{lem:characteristic-test} supplies the ordinary test. Every ideal is the directed union of its principal ideals, so each Scott open is a union of these cylinders. The existential definition of $\mathsf{May}_S$ preserves arbitrary unions. The criterion now follows in both directions from Theorem~\ref{thm:trace-full-abstraction}. The argument requires finite tests only for basis cylinders. Distinct successful alternatives can witness different opens, so their conjunction must be witnessed on a common execution.
\end{proof}

\begin{corollary}[The observations of a represented payload]\label{cor:exact-payload-observations}
Take one index in Construction~\ref{con:positive-observation-port} and write $Q=\mathsf Q_D(\{*\})$. The subsets of $D$ observable by a Scott-open trace test on $P_x$ are exactly $\Op(D)$:
\[
 \bigl\{\{x:\mathsf{May}_{P_x}(W)\}:W\in\Op(T_Q)\bigr\}
       =\Op(D).
\]
Moreover $x\le y$ exactly when $P_x\leq_{\mathrm{may}}P_y$.
\end{corollary}
\begin{proof}
For a legal word $w$ on $Q$, its requests are always accepted. It is a trace of $P_x$ exactly when $k\le x$ for each acknowledgement of $k$ in $w$. The payload predicate of $[w]$ is therefore a finite intersection of principal opens. For arbitrary $W$ take the union over its basis cylinders, giving a Scott open of $D$. Conversely write $r_k,a_k$ for the request and acknowledgement of $k$. For $U\in\Op(D)$ take
\[
                W_U=\bigcup_{k\in K(D)\cap U}[r_k\,a_k].
\]
A successful trace witnesses some $k\le x$ in $U$, and every such $k$ supplies the two-event trace. Proposition~\ref{prop:open-payload-basis} now gives precisely $x\in U$. By Theorem~\ref{thm:open-trace-testing}, $P_x\leq_{\mathrm{may}}P_y$ holds exactly when every compact below $x$ is below $y$. Taking their directed supremum gives $x\le y$.
\end{proof}

\begin{example}[Endpoint and interactive observations]\label{ex:trace-endpoint-separation}
Let the boundary have independent events $n^-,p^+,q^+$. In $T$, $p$ is initially available and $q$ waits for $n$. In $S$ take two conflicting occurrences of $p$. The eager one conflicts with $q$. The other waits for $n$ and can coexist with $q$, which also waits for $n$. Receptivity holds because $n$ is always available before its occurrence. All extra covers are negative-to-positive. Both strategies attain precisely
\[
 \varnothing,\quad\{n\},\quad\{p\},\quad\{n,p\},\quad\{n,q\},
 \quad\{n,p,q\}
\]
at the boundary. Thus all may tests of the endpoint configuration alone give the same results. The cylinder $[p\,n\,q]$ separates them: in $S$ an early $p$ selects the occurrence incompatible with $q$. This is why the frame of current public assertions and the open testing of interactive traces use different domains.
\end{example}

\subsection{Proof-relevant comparison and composable certificates}\label{subsec:trace-certificates}
Trace equality characterizes observation, but later constructions must also transport the evidence for that equality. A certificate therefore retains realizations of each shared word on both sides. Surjective legs make the comparison complete without choosing a matching realization for every witness.

\begin{definition}[Complete trace certificate]\label{def:trace-certificate}
A certificate from $S$ to $T$ is a span
\[
 \mathsf E(S)\xleftarrow{p}M\xrightarrow{q}\mathsf E(T),
 \qquad \lambda_Sp=\lambda_Tq,
\]
with both legs surjective. Its elements retain the two source configurations, their serializations, and any additional comparison witness. A map of certificates commutes with both legs.
\end{definition}

\begin{theorem}[Complete comparison calculus]\label{thm:complete-certificates}
The displays $S,T$ are contextually equivalent iff they admit a trace certificate. The canonical one is
\[
 \mathsf E(S)\times_{\mathsf W(A)}\mathsf E(T).
\]
Certificates compose by pullback over $\mathsf E(T)$ and have diagonal identities. They form a bicategory with the canonical rebracketing of witness tuples. Secured contexts and tensor transport certificates to certificates.
\end{theorem}
\begin{proof}\vmode
\leavevmode
\begin{enumerate}[label=(\roman*)]
\item \textbf{Completeness of certificates.}
Equal trace images make the canonical projections surjective; any surjective equal-word span proves equal images.
Apply Theorem~\ref{thm:trace-full-abstraction}.

\item \textbf{Pullback composition.}
Pullback of surjections is surjective in sets: for a given left witness choose a middle witness over its image, and conversely.
The pairing retains both witnesses, so no chosen matching function is required.

\item \textbf{Span coherence.}
Associativity and units identify the same tuples and give their coherence equations.

\item \textbf{Transport through a context.}
For context transport retain a certificate element, a finite context configuration, and a global linearization of their secured matching.
Replace the component serialization by its mate, as in the preceding proof, and hide the common interface.

\item \textbf{Surjectivity after hiding.}
Both projections cover all resulting serialized configurations.
Indeed any serialization of a visible configuration extends to a serialization of a finite matching representative: add its visible order to the joint partial order; a cycle would contradict the visible order being a linear extension, and then take a finite topological sort.
Surjectivity of the original certificate supplies the replacement.
The same construction in the reverse direction gives the other projection.

\item \textbf{Tensor transport.}
Tensor retains a pair of certificate elements and the shuffle of their displayed words.
These operations prove contextual transport without imposing extension-compatible choices of matches.
\end{enumerate}
\end{proof}

A certificate retains a comparison object as evidence of equality; the setoid construction of~\cite{benton-hofmann-nigam} provides a related but distinct use of proof relevance. There is also a testwise presentation. Let $W_{S,\tau}$ be the finite successful execution witnesses for a test $\tau$. Contextual equality is equivalent to the existence, for every $\tau$, of a span surjective onto $W_{S,\tau}$ and $W_{T,\tau}$. Their Cartesian product is the canonical carrier, and is empty when both sets are empty. Trace certificates assemble the comparison into one structural carrier independent of the success open.

\subsection{Causal certificates as a stronger proof method}\label{subsec:causal-certificates}
Some protocol comparisons preserve the order of each finite history, not just its possible words. The causal certificates below express that stronger property and provide a sufficient method for obtaining the complete observational certificates already constructed.

\begin{definition}[Finite displayed order carrier]\label{def:finite-displayed-orders}
For an arena $A$ let $\mathsf{Ord}_f(A)$ consist of pairs $(u,R)$, where $u\in\Conf_f(A)$ and $R$ is a partial order extending the arena order on $u$. A display $S\to A$ defines
\[
 b_S:\Conf_f(S)\to\mathsf{Ord}_f(A),\qquad
 b_S(x)=(\sigma[x],(\sigma|_x)_*\leq_x),\quad
 \mathcal B_A(S)=\im b_S.
\]
Local injectivity makes the transported order well defined. It extends arena order because the displayed principal past of each occurrence is an arena configuration containing its label.
\end{definition}

\begin{definition}[Finite causal comparison certificate]\label{def:crypto-causal-certificate}
A causal certificate is a span $\Conf_f(S)\xleftarrow p M\xrightarrow q\Conf_f(T)$ with $b_Sp=b_Tq$ and
\[
 \im(b_Sp)=\mathcal B_A(S),\qquad
 \im(b_Tq)=\mathcal B_A(T).
\]
Thus it covers displayed order images, not necessarily every source configuration. It exists exactly when these images agree; their equality pullback is a canonical certificate.
\end{definition}

\begin{lemma}[Finite causal behaviours suffice for contextual comparison]\label{lem:crypto-causal-congruence}
A causal certificate gives $S\eqsc T$. Equivalently it suffices to match every finite configuration, in each direction, by a display-preserving order isomorphism. No globally selected or extension-compatible family of matches is needed.
\end{lemma}
\begin{proof}\vmode
\leavevmode
\begin{enumerate}[label=(\roman*)]
\item \textbf{Equality of trace images.}
Matching orders have exactly the same displayed linearizations, so their trace images agree.
Theorem~\ref{thm:trace-full-abstraction} applies.

\item \textbf{Transport of causal witnesses.}
More concretely, transporting a secured matching along the order isomorphism preserves its acyclicity, visible causal past and displayed order; this also proves the stronger causal certificate transport used next.
\end{enumerate}
\end{proof}

\begin{construction}[Saturation and contextual composition]\vmode\label{con:certificate-interaction}
\leavevmode
\begin{enumerate}[label=(\roman*)]
\item \textbf{Saturated carrier.}
Write $X=\Conf_f(S)$, $Y=\Conf_f(T)$, and $B=\mathsf{Ord}_f(A)$.
Saturate an image-covering certificate by
\[
 \widehat M=X\times_B M\times_B Y,
\]
using its common map to $B$.

\item \textbf{Surjective endpoints.}
Both endpoint projections are now surjective, and $m\mapsto(p(m),m,q(m))$ retains every original witness.

\item \textbf{Composition and coherence.}
Saturated certificates compose by pullback over their actual middle configuration set.
Reassociation keeps the same tuple, giving the span coherence maps.

\item \textbf{Need for saturation.}
Pulling back the original unsaturated spans directly can lose coverage: their two middle witnesses may have the same displayed order but be different configurations.

\item \textbf{Transport through a context.}
For a context $K$, pair each element of $\widehat M$ with a finite context configuration giving a secured matching.
Transport the strategy nodes by the unique display-preserving order isomorphism and keep the context nodes.
The joint orders are isomorphic; take visible principal histories after pruning.

\item \textbf{Coverage and resaturation.}
Every finite composite configuration has such a finite representative, so the resulting span covers both displayed order images.
Saturate again if configuration-surjective legs are needed.

\item \textbf{Tensor transport.}
Tensor takes products and disjoint unions of orders.

\item \textbf{Trace-certificate conversion.}
Finally, adjoining a linearization of the common displayed order converts a saturated causal certificate into a trace certificate, with surjectivity on both serialization carriers.
\end{enumerate}
\end{construction}

\begin{example}[Causal equality is strictly finer than testing]\label{ex:causal-certificate-incomplete}
Let $A$ be flat with $n^-,p^+$. The eager display $E$ has independent $n,p$. Let $D$ have $n^-$, independent $p_0^+$, and $p_1^+$ with $n<p_1$, where $p_0$ conflicts with $p_1$ and both display as $p$. Both displays are receptive and courteous, and
\[
 \Trace(E)=\Trace(D)=\{\eps,n,p,np,pn\}.
\]
However $D$ has a displayed configuration ordered $n<p$, whereas $E$ does not. Thus trace certificates are complete for the declared observation, while causal certificates are a sufficient, stronger construction method. Neither comparison erases the witnesses from the original displays.
\end{example}

\subsection{Configuration and kernel calculation of strategies}\label{subsec:kernel-calculation}
The testing theorem identifies what an observer can distinguish. To calculate the execution witnesses themselves, we return to the finite configurations of a display $\sigma:S\to A^\ast\parallel B$. Its kernel has fibres $\Phi_\sigma(x,y)=\{z\in\Conf_f(S):\sigma[z]=x\sqcup y\}$. We will identify secured interaction with the reached kernel, retaining the witnesses in this formula. Appendix~\ref{subsec:technical-kernels} supplies the generators, raw coends, categories of elements and resource charts. The polarized Scott order used for this calculation is distinct from inclusion on $\D(S)$.

\subsection{Concurrent cells and Scott categories}\label{subsec:concurrent-cells-scott-categories}
A configuration kernel must retain independent extensions as well as positions. Configuration cubes record those extensions, and the polarized Scott category determines the directions in which kernel fibres restrict.


\begin{definition}[Event-generated configuration cubes]\label{def:formal-event-generated-configuration-cubes}
For a polarized event structure $E$, define $K(E)$ to have an $F$-cell $(x,U,\phi)$ consisting of a finite configuration $x$, a finite set $U$ of events disjoint from $x$, and a bijection $\phi:F\to U$, such that $x\cup V$ is a configuration for every $V\subseteq U$. Faces omit an event from $U$, placing it in $x$ for the upper face. Thus the events in $U$ are jointly enabled and pairwise concurrent.
\end{definition}

\begin{definition}[Scott cubes and their Kan extension]\label{def:scott-cubes-kan-extension}
For a labelled shape put
\[
 J_{\mathrm{Sc}}(F,\ell)=\prod_{f\in F}[1]_{\mathrm{pol}(\ell(f))}.
\]
Cofaces insert their endpoint constants. Define the ordinary Scott category by
\begin{equation}\label{eq:scottkan}
 \Sc(X)=\colim_{(c,x)\in\int X}J_{\mathrm{Sc}}(c)
 \quad\text{in }\Cat.
\end{equation}
Its nerve is the set-valued presheaf
$N_{\mathrm{Sc}}(C)_c=\Hom_{\Cat}(J_{\mathrm{Sc}}(c),C)$.
\end{definition}
\begin{proposition}[Universal property and tensor]\label{prop:scottmonoidal}
There is an ordinary adjunction $\Sc\dashv N_{\mathrm{Sc}}$ and coherent natural isomorphisms
\[
 \Sc(X\boxtimes Y)\cong\Sc(X)\times\Sc(Y),
 \qquad \Sc(X^\ast)\cong\Sc(X)^\op.
\]
\end{proposition}
\begin{proof}\vmode
\leavevmode
\begin{enumerate}[label=(\roman*)]
\item \textbf{Universal property.}
A functor out of the colimit is a family of functors from its categorical cubes, agreeing on every face and reparameterization.
Such families are exactly presheaf maps $X\to N_{\mathrm{Sc}}C$.

\item \textbf{Tensor comparison.}
On representables, disjoint union is sent to product of categorical cubes.
Ordinary $\Cat$ is cartesian closed: product with $C$ has the functor-category right adjoint $\Fun(C,-)$ and hence preserves colimits.
Substituting the Day coend into \eqref{eq:scottkan}, commuting the colimits, and applying this representable calculation gives the first isomorphism.

\item \textbf{Coherence.}
Every iterated comparison partitions the same directions into the same factors, proving coherence.

\item \textbf{Polarity reversal.}
Opposition is an involutive equivalence on ordinary categories, so preserves colimits; reversing each interval proves the second isomorphism.
\end{enumerate}
\end{proof}
\begin{definition}[Scott boundary notation]\label{def:formal-scott-boundary-notation}
For an arena $A$ write $C_A=\Sc(A)$ and let $i_A$ be the object given by its root. The morphism arena is $A^\ast\boxtimes B$, with Scott category $C_A^\op\times C_B$. The cubical carrier remains part of the game throughout. Its dimension counts simultaneous execution directions; replacing the associated categorical machinery by ordinary categories does not discard those cells.
\end{definition}

\begin{definition}[Scott adjunction unit and counit]\label{def:scott-adjunction-maps}
The \textbf{Scott unit} $\eta_X:X\to N_{\mathrm{Sc}}\Sc X$ sends a cell $x\in X(c)$ to its colimit-insertion functor $J_{\mathrm{Sc}}(c)\to\Sc X$. The \textbf{Scott counit} $\epsilon_C:\Sc N_{\mathrm{Sc}}C\to C$ is the functor induced by the evaluation cocone: on a cell $F:J_{\mathrm{Sc}}(c)\to C$, its restriction is $F$. The colimit property supplies both maps and the adjunction of Proposition~\ref{prop:scottmonoidal}.
\end{definition}

These maps obey
\[
\begin{tikzcd}[column sep=large]
\Sc X \arrow[r,"{\Sc\eta_X}"] \arrow[dr,"{1}"']
 &\Sc N_{\mathrm{Sc}}\Sc X \arrow[d,"{\epsilon_{\Sc X}}"]\\
 &\Sc X
\end{tikzcd}
\qquad
\begin{tikzcd}[column sep=large]
N_{\mathrm{Sc}}C \arrow[r,"{\eta_{N_{\mathrm{Sc}}C}}"] \arrow[dr,"{1}"']
 &N_{\mathrm{Sc}}\Sc N_{\mathrm{Sc}}C \arrow[d,"{N_{\mathrm{Sc}}\epsilon_C}"]\\
 &N_{\mathrm{Sc}}C .
\end{tikzcd}
\]
Each triangle evaluates the same cube functor; thus these equations also specify the adjunction's coherence maps.

\subsection{Secured interaction and reached kernels}\label{subsec:pc-secured-proof}
The interaction of $\sigma:S\to A^\ast\parallel B$ and $\tau:T\to B^\ast\parallel C$ is already determined by the matching histories of Section~\ref{subsec:game-secured}. The remaining task is to recover precisely those histories from the Scott kernels. Normalization and reachability identify the part of the kernel calculation that represents ordinary interaction and hiding.

\begin{lemma}[Coend matching normalization]\label{lem:pi0coend}
The raw set-valued coend $R=\Phi_\tau\star_B\Phi_\sigma$ is the quotient of matching pairs by insertion and deletion of matched internal events maximal on both sides. Securedness is invariant under this equivalence. Two secured pairs give the same coend class exactly when their visible-past matchings agree. Consequently the composite's configuration kernel injects into $R$, with image the classes of secured pairs.
\end{lemma}
\begin{proof}\vmode
\leavevmode
\begin{enumerate}[label=(\roman*)]
\item \textbf{Reduction to atomic Scott arrows.}
By Theorem~\ref{thm:eventscott}, every Scott arrow at the middle interface is a composite of single negative removals and positive additions.
It suffices to analyze the generating coend relation for one such arrow.

\item \textbf{Positive middle-interface move.}
For a positive $B$ addition, its lift on the $S$ side is removal of the corresponding positive strategy event in the reverse restriction direction, while its lift on the $T$ side is the unique addition of the corresponding negative event.
The removed positive event is maximal: courtesy reflects a causal successor of a positive event to the arena, where the added $B$ event is maximal and has no order relation to the external $A$ coordinates.
The newly added event on the other side is maximal in its new configuration.
Thus the coend relation adds or removes one matched event maximal on both sides.

\item \textbf{Negative move and converse.}
For a negative $B$ removal the roles are reversed.
Conversely any such matched maximal event gives this atomic relation, using its polarity and unique negative lift.

\item \textbf{Securedness and visible-past invariance.}
Adding a new maximal node to an acyclic joint graph, or deleting a maximal node, preserves acyclicity in both directions.
It also leaves the visible down-closure unchanged.
The same is therefore true along every coend zigzag.

\item \textbf{Normalization and completeness.}
By Lemma~\ref{lem:matching-past}, every secured pair is equivalent to its visible-past matching.
Hence equality of those matchings is sufficient, and invariance proves necessity.

\item \textbf{Composite configurations.}
The description of composite configurations in Proposition~\ref{prop:finite-secured-witnesses} gives the claimed injection and image, with all occurrence identities retained.
\end{enumerate}
\end{proof}

The following is the finite-configuration form of \cite[Theorem~2]{profunctors}. We retain the verification that its reachable matching classes are exactly the progress-word image used in this paper.
\begin{theorem}[Reached interaction is ordinary hiding]\label{thm:ordinarycomp}
There is a natural isomorphism of pointed set-valued kernels
\[
 \Reach(\Phi_\tau\star_B\Phi_\sigma)
 \cong\Phi_{\tau\odot_{\mathrm{ord}}\sigma}.
\]
Its cubical unstraightening is the configuration carrier of the ordinary composite. The comparison respects strategy composition, parallel product and compact duality, including the coherence maps of Theorem~\ref{thm:ordinary-category}.
\end{theorem}
\begin{proof}\vmode
\leavevmode
\begin{enumerate}[label=(\roman*)]
\item \textbf{Rooted secured classes.}
Let $U\subseteq R$ be the classes of secured pairs.
The empty pair belongs to $U$.

\item \textbf{Closure under visible restrictions.}
An atomic external restriction either adds an enabled negative visible event or removes a maximal positive visible event.
The former has its unique receptive lift, and the latter has its unique maximal-event removal by courtesy, as in Theorem~\ref{thm:scott-discrete-fibrations}.
It leaves the middle matching unchanged.
Adding or removing this maximal visible node preserves acyclicity, so $U$ is a subpresheaf.

\item \textbf{First reachability inclusion.}
The same argument proves closure under reverse positive restrictions: when a class restricts to a secured class, use its matching representative; its restriction is secured by Lemma~\ref{lem:pi0coend}, and restoring the one maximal visible node is a legal secured extension.
Proposition~\ref{prop:kernel-reachability} gives $\Reach(R)\subseteq U$.

\item \textbf{Construction of a progress history.}
Conversely build a secured pair in a linear extension of its joint order.
Each internal synchronized extension is an equality of coend classes by Lemma~\ref{lem:pi0coend}.
Each negative visible extension is a restriction, and each positive visible extension is a reverse positive restriction.
The resulting finite word is a progress history in $R$.
Thus $U\subseteq\Reach(R)$.

\item \textbf{Identification of configuration fibres.}
Matching normalization identifies $U$ with the composite's configuration fibres.
Restrictions act by the same visible extensions or removals and the same pruning of unused internal events, so the identification is natural and preserves the empty root.
Receptivity and courtesy of the composite, verified in Theorem~\ref{thm:ordinary-category}, allow Theorem~\ref{thm:scott-discrete-fibrations} to identify its cubical display.

\item \textbf{Structural coherence.}
For compatibility with composition, expand both bracketings to their joint matching histories.
The comparison retains the same visible principal histories under each associator and unitor.
Parallel histories split componentwise, and cups and caps are the same copycat histories with changed boundary assignments.
Thus the comparison also preserves parallel product, compact duality and their coherence maps.
\end{enumerate}
\end{proof}

\begin{definition}[Extension carriers and kernel positions]\label{def:semantic-extension-fibres}
For an ordinary display $\sigma:S\to A$, its \textbf{extension carrier} is $\mathsf{Ext}(S)=\{(x,s)\in\Conf_f(S)\times S:s\notin x,\ x\cup\{s\}\in\Conf_f(S)\}$. Its endpoints form the commuting square
\[
\begin{tikzcd}[column sep=large]
\mathsf{Ext}(S)
 \arrow[r,"{\lambda}"] \arrow[d,"{\partial}"']
 &\Conf_f(S)\times A \arrow[d,"{b}"]\\
\Conf_f(S)^2 \arrow[r,"{a}"']
 &\Conf_f(A)\times\mathcal P(A).
\end{tikzcd}
\]
Here $\lambda(x,s)=(x,\sigma(s))$, $\partial(x,s)=(x,x\cup\{s\})$, $b(x,v)=(\sigma[x],\{v\})$, and $a(x,y)=(\sigma[x],\sigma[y\setminus x])$. The occurrence $s$ is retained in the upper-left carrier even if another incompatible event has the same displayed label.

For a kernel $\Phi:A\to B$, a \textbf{kernel position} consists of an interface configuration and a fibre element. At a negative arena extension, presheaf restriction supplies the next fibre element. At a positive extension $c\to d$, a continuation of $z\in\Phi(c)$ is an element $w\in\Phi(d)$ satisfying $\Phi(c\to d)(w)=z$. Concurrent extensions are calculated by the terminal-corner formula of Lemma~\ref{lem:terminalcorner}. Theorem~\ref{thm:scott-discrete-fibrations} identifies these rules with the configuration rules for an ordinary strategy. These rules calculate the protocol examples directly from their selected occurrence commitments.
\end{definition}

\subsection{Calculus: configurations, execution witnesses, and equations}\label{subsec:calculus-configurations}
The kernel comparison now gives an execution interpretation for process terms. The following judgements expose finite configurations and witnessed extensions, and the equations identify exactly the rearrangements justified by the retained cells.

\subsubsection{Operational configurations of process terms}\label{subsec:native-game-rules}
\begin{definition}[Configuration and execution types]\label{def:calculus-configuration-types}
For a typed occurrence term with display $\sigma:S\to A$, the \textbf{configuration type} is $\mathsf{Cfg}(P)=\Conf_f(S):\mathsf{type}_{\Set}$. The \textbf{finite execution type} $\mathsf{Exec}_P(x,y)$ consists of finite composable lists of the occurrence extensions of Definition~\ref{def:formal-operational-extension-objects-and-kernel-positions}, with endpoints $x,y$. Its rules are
\[
 \frac{}{\varnothing:\mathsf{Cfg}(P)},\qquad
 \frac{x:\mathsf{Cfg}(P)\quad e\notin x\quad x\cup\{e\}\in\Conf_f(S)}
      {x\cup\{e\}:\mathsf{Cfg}(P)},
\]
\[
 \frac{}{\mathsf{nil}_x:\mathsf{Exec}_P(x,x)},\qquad
 \frac{p:\mathsf{Exec}_P(x,y)\quad (y,e)\in\mathsf{Ext}(S)}
      {\mathsf{snoc}(p,e):\mathsf{Exec}_P(x,y\cup\{e\})}.
\]
The \textbf{execution fold} into a family $Z(x,y)$ is specified by $z_x\in Z(x,x)$ and typed handlers $h_e:Z(x,y)\to Z(x,y\cup\{e\})$. Its computation is $\mathsf{fold}(\mathsf{nil}_x)=z_x$ and $\mathsf{fold}(\mathsf{snoc}(p,e))=h_e(\mathsf{fold}(p))$. The \textbf{cell quotient} identifies the two residual orders around each constructed square. A fold on that quotient is exactly a fold whose handlers satisfy those square equations; this is the quotient's universal property.
\end{definition}
\begin{proposition}[Configuration elimination and cell computation]\label{prop:calculus-configuration-elimination}
The configuration rules generate every finite configuration. Finite execution has the displayed unique fold. Two enumerations of one finite configuration differ by the constructed adjacent independent swaps.
\end{proposition}
\begin{proof}\vmode
\leavevmode
\begin{enumerate}[label=(\roman*)]
\item \textbf{Execution of a finite configuration.}
Remove a maximal event of a nonempty finite configuration.
The remainder is a configuration of smaller cardinality, so induction constructs an execution from the empty configuration.

\item \textbf{Soundness of the rules.}
Conversely each rule preserves downclosure and consistency by its premise.

\item \textbf{Execution fold.}
Recursion on the length of a list gives the fold, and induction on length proves its uniqueness and computation equations.

\item \textbf{Exchange of linearizations.}
Given two linearizations, move the first event of the second leftwards in the first.
Every crossed event is incomparable with it: an earlier predecessor would occur earlier in the second, and a successor could not occur earlier in the first.
All intermediate prefixes are ideals of the same consistent configuration, hence each crossing is its specified square.
Remove the common first event and repeat by finite induction.

\item \textbf{Quotient elimination.}
Functions constant under these squares factor uniquely through the quotient.
This proves the elimination principle with its precise independence equations.
\end{enumerate}
\end{proof}

\begin{definition}[Operational extension objects and kernel positions]\label{def:formal-operational-extension-objects-and-kernel-positions}
For an ordinary strategy $\sigma:S\to A$, its operational configurations are $x\in\Conf_f(S)$. Its witnessed moves and concurrent families are defined by
\[
 \frac{s\notin x\quad x\cup\{s\}\in\Conf_f(S)}
      {x\xrightarrow{\sigma(s)}x\cup\{s\}},
 \qquad
 \frac{U\cap x=\varnothing\quad
       (x\cup V\in\Conf_f(S))_{V\subseteq U}}
      {(x,U)\text{ is a concurrent cell}}.
\]
The witness of the first rule is $s$, not just its image label. Faces of the second rule execute some events and omit others. The face equations follow from union and restriction, as in Proposition~\ref{prop:family-cubes}. Thus these operational rules are read directly from a strategy's configurations.

The extension carrier and kernel-position families are those of Definition~\ref{def:semantic-extension-fibres}. Their occurrence and pullback witnesses interpret precisely the displayed rules.
\end{definition}

\subsubsection{Process equations and retained witnesses}\label{subsec:pc-equations}
The kernel constructions calculate the process operations of Section~\ref{subsec:pc-typing}. Their configurations retain reached occurrence histories, their cells retain residual families of commitments, and their fibre elements retain the suppliers and choices realizing a boundary configuration. Thus Theorem~\ref{thm:ordinarycomp} identifies the reached coend calculation with typed cut. These retained witnesses determine which process equations hold.

\begin{proposition}[Receipt and counter realization]\label{prop:pc-guard-store}
A finite core declaration is implemented exactly by persistent receipts and anonymous finite-cohort counters, with atomic output commitments. This implementation preserves its committed graph and finite cells.
\end{proposition}
\begin{proof}\vmode
\leavevmode
\begin{enumerate}[label=(\roman*)]
\item \textbf{Receipt and budget invariant.}
The receipt of $e$ is present exactly when $e\in C$.
Give each $F\in K$ initial budget $|F|-1$; after $C$ its residual budget is $|F|-1-|C\cap F|$.

\item \textbf{Individual enabling.}
An output checks its waiting receipts and atomically consumes from every incident counter.
These tests are exactly the core enabling tests.

\item \textbf{Concurrent enabling.}
A joint family must additionally be an arena cell and fit every residual counter; these are precisely the tests for every subset of that family.
No reservation transition or named choice of a budget slot is introduced.
\end{enumerate}
\end{proof}

\begin{example}[Two equations that discard witnesses]\label{ex:pc-equivalence-boundary}
The conventional processes $a!v.0$ and $a!v.0+a!v.0$ have the same labelled transition relation up to strong bisimulation. Their exact occurrence presentations have one and two conflicting positive realizations, respectively. The one-event fibres have cardinalities one and two, so their displays are not isomorphic. Similarly, the concurrent execution of independent $p$ and $q$ has a square; the sequential alternatives $p.q+q.p$ retain two schedules but do not supply that square. Neither idempotent choice nor this interleaving expansion is an equation of the occurrence language.
\end{example}

Display equality, operational reduction, and a chosen contextual observation consequently have distinct notation and uses. The named protocol instances below state their observations explicitly. An operational comparison is transported through a game interpretation only by its stated comparison theorem.


\subsection{Process evidence for observations}
Add an assertion judgement, with semantic assertions ranging over arbitrary Scott opens of the existing boundary domain. Its atoms are \([k]_J\) for finite boundary configurations \(k\), interpreted on \(R\) by \(B_k^{R,J}\). Finite conjunctions and arbitrary set-indexed disjunctions present the public frame after quotienting by their interpretation on \(R\). The latter quotient matters: resource-specific entailments are derived from \(R\), not just from the arena's compatibility equations.

For a finite configuration \(x:\Cfg(R)\), the elementary evidence rules are
\[
\frac{k\subseteq\partial_Jx}{x\models_R[k]_J},\qquad
\frac{x\models_R\varphi\quad x\models_R\psi}{x\models_R\varphi\wedge\psi},\qquad
\frac{i\in L\quad x\models_R\varphi_i}{x\models_R\bigvee_{i\in L}\varphi_i}.
\tag{Obs}
\]
The open assertions belong to the specification judgement. An arbitrary disjunction is not a new process operation or a program deciding membership in a non-effective open. The already constructed execution judgement supplies the operational evidence:
\[
\frac{p:\Exec_R(\varnothing,x)\qquad \partial_Jx\in U}
 {R\Downarrow_J U}.
\tag{May}
\]
Conjunction in \textup{(Obs)} uses the same configuration. Two different may-witnesses for \(U\) and \(V\) do not by themselves witness \(U\cap V\); incompatible branches give immediate counterexamples to such a rule.

\begin{theorem}[Finite operational evidence]\label{thm:evidence}
For any Scott open \(U\subseteq\D(A_J)\),
\[
\exists x\in\D(S_R)\;\partial_Jx\in U
\quad\Longleftrightarrow\quad
\exists x\in\Conf_f(S_R)\ \exists p:\Exec_R(\varnothing,x)\;\partial_Jx\in U.
\]
\end{theorem}
\begin{proof}\vmode
\leavevmode
\begin{enumerate}[label=(\roman*)]
\item \textbf{Finite observation witness.}
The left-hand condition is membership of a configuration in the Scott open \(\partial_J^{-1}U\).
Algebraicity supplies a finite subconfiguration already in that open.

\item \textbf{Operational realization.}
Proposition~\ref{prop:calculus-configuration-elimination} supplies a finite execution to it.

\item \textbf{Converse and exchange invariance.}
Conversely that finite configuration is a configuration of the interpreted strategy.
Independent reorderings yield the same boundary evidence by the existing cell equations.
This uses finite witnesses but selects no effective enumeration of arbitrary opens.
\end{enumerate}
\end{proof}


The observational criterion is now complete: finite trace equality characterizes contextual may equivalence, and serialized configurations provide composable certificates. Causal certificates remain available when a comparison must also preserve finite orders. The protocol constructions that follow use these criteria alongside the retained execution histories, choosing the comparison appropriate to each claim.

\section{Communication behaviours and distributed control}\label{sec:protocol-games}
The comparison theory can now be applied to communication networks. We first retain the histories and participant views of a network presentation, then ask which decisions can be implemented from the specified local observations. Permissions and their commitment compiler answer that question, with separate comparisons for histories, concurrency and winning conditions. The observational comparisons use the finite tests of Section~\ref{subsec:complete-trace-testing}. The medium implementations also retain the histories used by the control proofs.

\subsection{Protocol interfaces and network interpretation}\label{subsec:abstract-protocols}
A protocol declaration supplies a realizer, while a network expression records how several declarations are connected. We retain the expression's internal execution presentation because later control and audit arguments use occurrences that its external arrow hides.

\begin{definition}[Communication protocol and presentation]\label{def:abstract-communication-protocol}
An open communication protocol $P:A\to B$ is an ordinary strategy, represented by a receptive, courteous display $\sigma_P:S_P\to A^\ast\parallel B$. A protocol declaration retains this display and its typed occurrence presentation from Theorem~\ref{thm:pc-representation}. Let $\mathcal A_{\mathrm{comm}}$ be the declared communication arenas closed under the empty arena, parallel product and polarity reversal. Write $\mathsf G_{\mathrm{comm}}$ for the full subcategory of $\CG$ on these arenas. The declared protocols are distinguished arrows; arbitrary strategies at the same interfaces remain available as contexts.

A network presentation $N$ is a typed expression in these declarations and the constructors of Definition~\ref{def:formal-typed-network-category-and-interpretation}. Its external display is $S_N$. Its retained execution presentation $U_N$ is the secured interaction of its component displays with all internal matched occurrences retained, before the final hiding. Construct it from finite matching configurations: retain every component occurrence and matching witness, test acyclicity of the union of their causal orders, and take principal histories as in Section~\ref{subsec:game-secured}. Hiding the internal occurrences gives $S_N$. The execution presentation belongs to the specified network expression; it is not reconstructed from the isomorphism class of its external arrow.

For a declared operational implementation write $X_N$ for its physical execution carrier and retain its constructed comparison with the causal presentation. Such a comparison records its actual event, history, or cubical maps. A physical event identifier need not be an event of a prime causal presentation. The concrete comparisons are proved in Appendix~\ref{app:communication-examples}.
\end{definition}

\begin{definition}[Retained matching labels and boundary polarizations]\label{def:retained-network-labels}
For a fixed network expression, tag each component occurrence by its
component leaf. A joint node is either an unmatched external occurrence
$(i,s)$, or the matching record $((i,s),b,(j,t))$ on a declared internal
wire, with $\sigma_i(s)=b=\sigma_j(t)$ after the wire identification.
The joint graph is the quotient of the tagged component nodes by these
declared matches; its order is the transitive closure of component
orders. Its legal finite graphs are the secured ones with configuration
projections. Principal histories realize the underlying event structure
$U_N$.

The retained label of an internal node records both incident port
signs, $(\epsilon,-\epsilon)$; it is not a single inherited ordinary
arena sign. External nodes record their external port sign.
When a binary polarization of the retained event structure is used,
the typed expression designates each wire's source port; define
$\pi_N$ on an internal root to be that source port's sign and on
an external root to be its external sign. This gives a total sign
function while retaining both incident signs in the label record.
Write $\lambda_N:|U_N|\to L_N$ for this record-valued label map.
Hiding restricts to external-rooted principal histories, whose sign
is their external root sign, and gives the ordinary display $S_N$.
Full audit assigns its specified negative sign to every executed
direction. Control assigns its specified negative sign to physical
events and positive sign to permission commitments. These are distinct
polarizations of their declared carriers.
\end{definition}
\begin{lemma}[Labelled residual and hiding compatibility]\label{lem:retained-residual-labels}
The history, prefix-insertion and residual constructions on $U_N$
preserve $\lambda_N$ on surviving events. For every specified
polarization $\pi$ of its event carrier, residual formation inherits
$\pi$ by restriction. External hiding inherits precisely the external
port signs.
\end{lemma}
\begin{proof}\vmode
\leavevmode
\begin{enumerate}[label=(\roman*)]
\item \textbf{Restriction of labels and supports.}
All these operations restrict the surviving event set or insert a fixed
configuration; none changes a principal-history record.
Hence the
restricted label is $\lambda_N|_{|U_N/h|}$ and the restricted sign is
$\pi|_{|U_N/h|}$.

\item \textbf{Iterated residuals.}
Iterated residuals have the same surviving events, so
their label and sign equations are identities.

\item \textbf{Hiding and internal signs.}
Hiding retains the
root's external tag.
The two incident signs of an internal matching
are never identified as one ordinary sign.
\end{enumerate}
\end{proof}

\begin{theorem}[Compact interpretation of communication networks]\label{thm:abstract-communication-interpretation}
The category $\mathsf G_{\mathrm{comm}}$ is compact closed. Typed communication networks have a compact-preserving interpretation into it, extending their declared strategies. The interpretation commutes with value substitution. Independent networks have disjoint parallel event presentations; connection is secured interaction followed by hiding.
\end{theorem}
\begin{proof}\vmode
\leavevmode
\begin{enumerate}[label=(\roman*)]
\item \textbf{Interpretation of network constructors.}
Interpret a declaration by its display, forwarding by copycat, tensor by disjoint parallel, connection by secured interaction and hiding, and bending by the copycat cup and cap.
Theorem~\ref{thm:ordinary-category} proves that each resulting display is a strategy and that composition is coherent.

\item \textbf{Compact structure.}
Fullness and closure of the selected objects supply the same cups, caps and structural arrows as $\CG$; their snake and monoidal equations are inherited.

\item \textbf{Interpretation of expressions.}
Induction on a network expression therefore defines the interpretation, and these equations make it descend through the structural congruence.

\item \textbf{Substitution.}
Substitution on each declaration is the indexed interpretation of Section~\ref{par:typed-network-interpretation}; induction on tensor and connection gives the same comparison on networks.

\item \textbf{Retained execution witnesses.}
Reassociation of a finite matching retains its component witnesses and their union order, giving the corresponding external-display comparisons.
The retained internal presentation is used when a statement concerns internal execution or stopping, rather than only the external arrow.
\end{enumerate}
\end{proof}

\subsection{Streams of observable payloads}\label{subsec:open-streams}
A stream observation can specify which packet positions have appeared and what has become known about each payload. For an algebraic payload domain $D$, order the finite lists of compact elements by
\[
 P_D=K(D)^*,\qquad
 (k_1,\ldots,k_m)\preceq(l_1,\ldots,l_n)
 \quad\Longleftrightarrow\quad m\le n\ \text{and}\ 
 k_j\le l_j\text{ for }1\le j\le m .
\]
Put $\mathsf S(D)=\operatorname{Idl}(P_D)$. The empty list gives the least stream. The alphabet model takes $D=\Delta M$. Sending an entire finite string as a single atomic payload instead takes $D=\Delta(M^*)$.

\begin{proposition}[Stream observations and their compact basis]\label{prop:stream-open-basis}
The compact points of $\mathsf S(D)$ are the principal ideals $\downarrow u$. Their principal opens $[u]=\{I:u\in I\}$ form a basis. The observation that the first $m$ positions exist with payloads satisfying $U_1,\ldots,U_m\in\Op(D)$ is the open
\[
 \bigcup_{\substack{k_j\in K(D)\cap U_j\\1\le j\le m}}
       [(k_1,\ldots,k_m)] .
\]
For $D=\Delta M$ the order on $P_D$ is precisely the prefix order of $C_M$ in Appendix~\ref{app:communication-examples}.
\end{proposition}
\begin{proof}
Each ideal is the directed union of its principal subideals. A directed union containing $u$ has a member containing $u$, giving compactness and the basis claim. Whenever position $j$ exists in $I$, the $j$th coordinates of its lists of length at least $j$ form a directed set. Its supremum is the payload at that position. If each of the first $m$ payloads lies in its respective open, pick a witnessing list for each coordinate by Scott openness. Directedness places these finitely many witnesses below one list in $I$; its length-$m$ prefix gives a member of the displayed union. The reverse implication follows by upward closure. For $m=0$ the observation is the whole domain. In the discrete case coordinate approximation is equality, giving the ordinary prefix cylinders.
\end{proof}

\begin{proposition}[Continuous processors and kernel composition]\label{prop:open-stream-processor}
For a Scott-continuous $f:\mathsf S(D)\to\mathsf S(E)$ define a subterminal kernel on the compact positions by
\[
 H_f(u,v)=\{*\mid v\in f(\downarrow u)\}.
\]
The identity is supported precisely when $v\preceq u$, equivalently $[u]\subseteq[v]$. For composable processors $f,g$ there is a canonical coherent isomorphism $H_g\star H_f\cong H_{gf}$ of these poset kernels. On observations, $(gf)^{-1}=f^{-1}g^{-1}$.
\end{proposition}
\begin{proof}
Monotonicity in $u$ and lower closure in $v$ give the two functorial actions. Continuity and directed approximation give
\[
 w\in g(f(\downarrow u))
 \quad\Longleftrightarrow\quad
 \exists v\in f(\downarrow u)\quad w\in g(\downarrow v).
\]
The right side is the support of the coend. Any two intermediate witnesses $v,v'$ have a common upper bound in the ideal $f(\downarrow u)$. It still witnesses the output by monotonicity of $g$. The coend identifications through this bound identify both representatives. Each nonempty fibre is therefore a singleton. The unique maps on subterminal fibres give naturality and associativity, and the identity follows by testing the inclusion of cylinders on $\downarrow u$.
\end{proof}

An original rooted prefix processor $h:C_M\to C_N$ extends to streams by $\bar h(I)=\bigcup_{u\in I}\downarrow h(u)$. This union is an ideal, commutes with directed unions and agrees with $h$ on principal ideals. Consequently $H_{\bar h}(u,v)=C_N(v,h(u))$. The FIFO and prefix-processor strategies of Appendix~\ref{app:communication-examples} realize these special cases with their original timing. Proposition~\ref{prop:continuous-query-strategy} also realizes the observation of a continuous stream processor when the chosen interface exports its query lanes.

A basis witness need not specify a first supplying occurrence. In the configuration domain of two independent events $a,b$, the open $\uparrow\{a\}\cup\uparrow\{b\}$ has two compatible minimal witnesses at $\{a,b\}$. The finite ordinal of input prefixes used in Proposition~\ref{prop:native-prefix} does have a first supplier. When the supplier matters in a concurrent protocol, its occurrence and alternative enabling histories are retained as in Section~\ref{subsec:disjunctive-causal-interfaces}.

\subsection{Causal histories and residual communication}\label{subsec:abstract-protocol-histories}
To describe what a network can do next, we first retain the history that reaches its current configuration. Residual communication then records extensions of that history, including their supplying occurrences.

\begin{definition}[Protocol history and extension]\label{def:protocol-history-extension}
For a specified protocol execution presentation $U=U_N$, put $\mathsf H(U)=\Conf_f(U)$. Its history graph has an edge $(h,e):h\to h\cup\{e\}$ exactly when $e\notin h$ and $h\cup\{e\}$ is a configuration. Write $\En_U(h)$ for this set of enabled occurrences. A concurrent family at $h$ is a finite set $F$ of such occurrences with $h\cup F\in\Conf_f(U)$. Its faces are $h\cup J$, $J\subseteq F$. These are the configuration cubes $K(U)$ calculated in Section~\ref{sec:game-calculus}. The same definitions apply to the displayed source $S_P$ when only external strategy histories are being studied.
\end{definition}

\begin{theorem}[Histories and independent exchange]\label{thm:generic-history-exchange}
Paths from $h$ to $k$ in the history graph are the linear extensions of $k\setminus h$ with its inherited order. The path category modulo adjacent exchanges of independent extensions is canonically the inclusion poset $\mathcal H_U=(\Conf_f(U),\subseteq)$.
\end{theorem}
\begin{proof}\vmode
\leavevmode
\begin{enumerate}[label=(\roman*)]
\item \textbf{Paths and linearizations.}
An extension path adds each event of $k\setminus h$ once and only after its predecessors, so gives a linear extension.
Conversely each such ordering gives configuration prefixes: they are downsets of the consistent configuration $k$.

\item \textbf{Independent exchange.}
To compare two orderings, move the first event of one leftwards in the other.
Every crossed event is incomparable with it, since either strict order would contradict one of the two linear extensions.
Each exchange is the square of the two enabled events in $k$.
Remove the common first event and induct on the number remaining.

\item \textbf{Path classes and composition.}
Thus there is exactly one path class when $h\subseteq k$, and none otherwise.
Concatenation gives inclusion composition.
\end{enumerate}
\end{proof}

\begin{proposition}[Generic residual histories]\label{prop:generic-protocol-residual}
For $h\in\Conf_f(U)$, the residual of Definition~\ref{def:event-residual-carrier} has
\[
 |U/h|=\{e\notin h:h\cup[e]_U\in\Conf_f(U)\},\qquad
 X\in\mathsf{Con}_{U/h}\ \Longleftrightarrow\ h\cup\downarrow_U X\in\Conf_f(U).
\]
With inherited order and retained labels it is an event structure; every specified polarization is inherited as in Lemma~\ref{lem:retained-residual-labels}. Prefix insertion $j_h(x)=h\cup x$ gives
\[
 \Conf(U/h)\cong\{k\in\Conf(U):h\subseteq k\},\qquad
 K(U/h)\cong K(U)_{\geq h},\qquad
 (U/h)/x\cong U/(h\cup x).
\]
The cube carrier on the right is rooted at $h$; the last formula concerns finite residual $x$. These comparisons commute with further prefix insertion and residual formation.
\end{proposition}
\begin{proof}\vmode
\leavevmode
\begin{enumerate}[label=(\roman*)]
\item \textbf{Residual event structure.}
Principal pasts are finite, singletons pass the defining consistency test, and taking subsets or adjoining predecessors preserves it.

\item \textbf{Configuration correspondence.}
A residual downset $x$ is consistent exactly when every finite part of $h\cup x$ is consistent in $U$; its residual downward closure says precisely that $h\cup x$ is downward closed.
Removing $h$ gives the inverse configuration map.

\item \textbf{Cube comparison.}
On cubes the comparison is $(x,F)\mapsto(h\cup x,F)$, so it commutes with every face.

\item \textbf{Iteration coherence.}
Both iterated residuals have the same surviving events and the consistency test with $h\cup x$.
Their comparison is the identity on those events; all iteration coherences follow from union associativity.
No operational state evaluation is used.
\end{enumerate}
\end{proof}

\subsection{Participant views and operational evaluation}\label{subsec:native-resource-histories}
A distributed participant acts from a local view of a history. The following projection and evaluation constructions identify the information that can enter such a decision and connect it to the retained operational witnesses.

\begin{definition}[Supports and causal decision views]\label{def:protocol-causal-views}
A distributed protocol presentation records the finite participant support $\dom(e)$ of each occurrence. For any finite history define
\[
 v_p(h)=\downarrow_U\{d\in h:p\in\dom(d)\},\qquad
 D_{e,p}=v_p([e)),\qquad [e)=\{d:d<e\}.
\]
The downclosure is taken in $U$. These are finite configurations and are defined without a last-event convention. An observation is a specified function on these local views. In a sequential-worker presentation, each worker's events in a configuration form a chain; several workers may belong to one party.
\end{definition}
\begin{lemma}[Sequential-worker comparison]\label{lem:sequential-worker-decision-view}
In a sequential-worker presentation, if $e$ is enabled after $h$ and $p\in\dom(e)$, then $v_p(h)=D_{e,p}$. The worker-chain property is preserved by independent tensor, secured interaction of the declared worker interfaces, and restriction to visible occurrences.
\end{lemma}
\begin{proof}\vmode
\leavevmode
\begin{enumerate}[label=(\roman*)]
\item \textbf{Identification of the causal view.}
Every $p$-event $d\in h$ is comparable with $e$ in $h\cup\{e\}$.
The alternative $e<d$ contradicts downward closure of $h$, so $d<e$.
Conversely $[e)\subseteq h$.
Taking downclosures proves the equality.

\item \textbf{Sequential visits and tensor.}
Sequential worker unfolding orders its local visits.
Tensor uses distinct worker tags.

\item \textbf{Interaction and hiding.}
A secured matching configuration projects injectively to each worker configuration; its order contains the worker order, so all visits to that worker remain comparable.
Hiding retains the induced order on the surviving visits.
These observations prove the preservation assertions directly.
\end{enumerate}
\end{proof}

A concrete operational presentation may evaluate histories by a map $q:\mathsf H(U)\to Q$, together with witnessed maps on extensions. Its comparison theorem establishes whether that evaluation is well defined and which extensions it lifts. Abstract histories already evaluate into their own configuration graph by $h\mapsto h$.

The later theory therefore uses $\En_U(h)$ directly; an operational enabled set enters through its proved comparison. Section~\ref{subsec:local-resource-charts} constructs the physical resource charts and universal rooted displays. Appendix~\ref{app:communication-examples} treats compatible alternative physical causes with their causal comparisons.

\subsection{Causal interfaces for protocol histories}\label{subsec:protocol-history-causal-interfaces}
\begin{definition}[Audit and communication interfaces]\label{def:formal-audit-and-communication-interfaces}
A protocol has the strategy display of Definition~\ref{def:abstract-communication-protocol}; a concrete operational realization supplies its comparison with that display. The retained execution, public communication, and resource-sensitive boundaries expose different data. The full audit interface retains every operational event and the resource charts. A communication interface exposes accepted inputs and delivered outputs, with internal witnesses retained in the identities of their causal histories after hiding. A protocol interface additionally exposes the dependencies and conflicts governing capacity, credit, status, or name availability. These interfaces have comparison maps; their different kernels are not identified merely because they produce the same payload words.
\end{definition}

\subsubsection{Constructing an arena from causal histories}
\begin{definition}[Least-order causal interface]\label{def:causal-interface}
For a polarized event structure $S$, define $A_S$ on the same events, consistency and polarity, with order the reflexive transitive closure of
\[
 \{(s,t):s\to_St,\ (\mathrm{pol}_S(s),\mathrm{pol}_S(t))\ne(-,+)\}.
\]
The \textbf{causal-interface display} is $\sigma_S=\id_{|S|}:S\to A_S$. Lemma~\ref{lem:causal-interface} proves that this is an ordinary strategy; its kernel is calculated in Corollary~\ref{cor:causal-interface-kernel}.
\end{definition}

\begin{lemma}[Causal-interface construction]\label{lem:causal-interface}
For a polarized event structure $S$, $A_S$ is an arena and $\sigma_S:S\to A_S$ is receptive and courteous.
\end{lemma}
\begin{proof}\vmode
\leavevmode
\begin{enumerate}[label=(\roman*)]
\item \textbf{Arena and display.}
The arena order is contained in the source order, so its pasts are finite and consistency is closed under adjoining its predecessors.
Every source configuration is an arena configuration.

\item \textbf{Receptivity.}
At a source configuration an enabled arena-negative event has all its immediate source predecessors present, since all edges into negative events were retained.
Source downclosure supplies their predecessors; unchanged consistency gives the unique extension.

\item \textbf{Courtesy.}
Every extra immediate source edge has sign $(-,+)$, proving courtesy.
\end{enumerate}
\end{proof}


\subsection{Calculus: witnessed communication and protocol computation}\label{subsec:calculus-protocols}
The history constructions now supply the semantic carriers for communication steps. Their process rules retain the selected offers, values and medium updates, so that the operational interpretation records the same committed witnesses.

\paragraph{Typed committed transitions.}\label{par:typed-committed-transitions}
Instantiate Definition~\ref{def:formal-committed-witness-fibres} with the scoped typed offers and medium store. Its dependent product already checks all subjects, payload sorts and selected continuations, including the terminal empty-premise case. A medium operation is introduced in this typed witness fibre; Appendix~\ref{app:communication-examples} constructs the concrete fibre elements.

\paragraph{Witness introduction and elimination.}\label{par:witness-introduction-elimination}
The typed step family is the target fibre
\[
 \mathsf{Step}_\Pi(q,q')=\{e\in E_\Pi(q):t(e)=q'\}
\]
of~\eqref{eq:committed-fibre}. Its introduction supplies the selected medium witness, matching offers and target equation. Dependent elimination applies the specified handler to this retained witness; beta and eta are the dependent-sum equations. Each medium specializes this one family, with no new logical connective or independent transition specification.

\begin{theorem}[Operational typing and concurrent witnesses]\label{thm:operational-typing-concurrent-witnesses}
The witnessed communication rule preserves the scoped and sorted runtime state. Its resource charts give commuting conversions with the same selected and residual witnesses. Network interpretation commutes with substitution and the typed structural operations.
\end{theorem}
\begin{proof}\vmode
\leavevmode
\begin{enumerate}[label=(\roman*)]
\item \textbf{Typed committed updates.}
A committed witness is formed in the dependent fibre of Definition~\ref{def:formal-committed-witness-fibres}: each premise has its declared subject and payload sort, and the medium output is in its declared sorted store and continuation fibre.
Input elimination substitutes this payload into the typed continuation; the substitution theorem gives its target scope and sort.
Output continuation, unchanged framed components, and restriction preserve their existing types.

\item \textbf{Capture avoidance and target typing.}
Tagged fresh occurrence names avoid capture.
Thus dependent introduction and elimination construct a target in the runtime-state type.

\item \textbf{Finite derivations.}
Recursive transitions have finite derivations, so induction on their last rule handles unfolding and structural transport.
The concrete operations construct their required typed fibre elements in Appendix~\ref{app:communication-examples}; the theorem uses that common introduction rule, not a list of media.

\item \textbf{Independent charts.}
A chart retains the selected instructions and equalizes their evaluated states at equal subsets of executed directions.
Its two-direction faces therefore give the same composite target and transported witnesses; repeated restriction gives every higher face equation.
This is the resource-coefficient construction of Proposition~\ref{prop:resource-charts}.

\item \textbf{Compositional interpretation.}
Finally the interpretation on declarations and each network constructor is that of Theorem~\ref{thm:abstract-communication-interpretation}; substitution on declarations and induction on constructors prove its stated compatibility.
\end{enumerate}
\end{proof}


\begin{construction}[The semantic meaning of a completed step]\vmode\label{con:completed-step-interpretation}
\leavevmode
\begin{enumerate}[label=(\roman*)]
\item \textbf{Execution continuation.}
Let $D$ be the constructed category of witnessed execution paths modulo its resource squares.
At a current state $q$, the continuation family is the covariant representable $D(q,-)$.

\item \textbf{Precomposition by a completed step.}
A step $d:q\to q'$ induces
\[
 d^*:D(q',-)\Rightarrow D(q,-),\qquad p\longmapsto pd.
\]
For $u:r\to r'$, $u(pd)=(up)d$, proving naturality.

\item \textbf{Square compatibility.}
For a resource square, equality of its composite arrows gives equality of these transformations.
This construction works directly for the full audit interface, whether or not a particular public ordinary-strategy presentation has been constructed.

\item \textbf{Event residual.}
For an event-structure history $h\in\Conf_f(U)$, use the already constructed residual $U/h$.
Its configurations are exactly the $x\subseteq|U|\setminus h$ for which $h\cup x\in\Conf(U)$.

\item \textbf{Residual iteration.}
If $e$ is enabled after $h$, both $(U/h)/\{e\}$ and $U/(h\cup\{e\})$ have the same surviving events, inherited order, and finite consistency test $h\cup\{e\}\cup\downarrow X\in\Conf(U)$.
The identity on the surviving events is consequently an isomorphism.
\end{enumerate}
\end{construction}

\subsubsection{The communication parameter and its operational rule}\label{subsec:pc-protocol-parameter}
The syntax, offers, binding, and occurrence typing of $\pcLang$ are those of Section~\ref{subsec:pc-named}. We now form $\pcLang[\Pi]$ by supplying a store and witnessed operations for the communication medium $\Pi$. The parameter supplies the typed operation witnesses, their store updates, and the selected offer premises of Definition~\ref{def:pc-medium}. Each instance in Appendix~\ref{app:communication-examples} constructs its causal interface and its comparison with the physical operation histories.

Definition~\ref{def:pc-medium} fixes the common witness format; the media of Appendix~\ref{app:communication-examples} instantiate its store and operation sets.

The scoped elaboration of Section~\ref{subsec:pc-named-proof} constructs this rule with changing active components. The calculations in Appendix~\ref{app:communication-examples} establish the causal interface and history comparison for each medium directly from its selected offers and store updates.

\paragraph{The object of committed witnesses.}
The edge carrier and its selected-premise pullback are Definition~\ref{def:formal-committed-witness-fibres}; equation~\eqref{eq:committed-fibre} applies to each declared medium.

\begin{proposition}[Named operational elaboration]\label{prop:pc-named-elaboration}
At a declared medium, elaborating active prefixes into offers preserves and reflects committed transitions, selected branches, substitutions and medium witnesses, modulo transported addresses and scope congruence. Forgetting witness multiplicity gives the corresponding transition relation. Retaining the specified residual cells gives the same concurrent families.
\end{proposition}
\begin{proof}\vmode
\leavevmode
\begin{enumerate}[label=(\roman*)]
\item \textbf{Elaboration comparisons.}
The explicit inverse maps between retained prefix records and offer records, including their source--target and cubical comparison squares, are constructed in Section~\ref{subsec:pc-named-proof}.

\item \textbf{Transition and cube preservation.}
Proposition~\ref{prop:pc-finite-support} proves the transition and cell bijections; Lemma~\ref{lem:pc-substitution} proves their compatibility with scope and substitution.

\item \textbf{Occurrence erasure.}
Taking the image after erasing the retained-record fibres gives the stated transition relation.
\end{enumerate}
\end{proof}

\subsection{Distributed observation and control}\label{sec:distributed-games}
Fix the retained execution presentation of a communication protocol. We ask which histories can be generated by decisions based on the declared local observations. Permission functions describe those decisions; their positive commitments are then realized in a control arena. The configuration comparison will show how this implementation preserves execution, concurrency and the specified winning conditions.\subsection{Observations and distributed strategies}\label{subsec:observations-distributed-strategies}
\begin{definition}[Plant action and veto interface]\label{def:formal-plant-action-and-veto-interface}
Let $U$ be a protocol presentation of Definition~\ref{def:abstract-communication-protocol}. An occurrence $e$ retains its action type $a(e)$, finite participant support $\dom(e)$, and veto set $V(e)\subseteq\dom(e)$. For each $p$, a local control-name map $\nu_p$ sends the types vetoable by $p$ to its exposed alphabet $L_p$. The identity name map gives the global-action alphabet $A_p^v$. Uncontrollable occurrences have $V(e)=\varnothing$. Enabling means the witnessed configuration extension $e\in\En_U(h)$.
\end{definition}
\begin{definition}[Local observation and permission controller]\label{def:causal-permission-controller}
Let $\mathsf H_p=\{v_p(h):h\in\Conf_f(U)\}$ with the views of Definition~\ref{def:protocol-causal-views}. A declared observation is $o_p:\mathsf H_p\to O_p$, and a controller is $k_p:O_p\to\Pow(L_p)$. Write
\[
 g_k(e)\ \Longleftrightarrow\
 \forall p\in V(e),\quad\nu_p(a(e))\in k_p(o_p(D_{e,p})).
\]
The controlled event structure is the downward-closed restriction
\[
 U_k=\{e\in U:\forall d\leq e,\ g_k(d)\},\qquad G(k)=\Conf_f(U_k).
\]
At a controlled history the enabled extensions are
\begin{equation}\label{eq:controlled}
 \En_k(h)=\{e\in\En_U(h):g_k(e)\}.
\end{equation}
The event $e$ retains its full witness even when different occurrences have the same exposed action name. Empty veto sets add no test.
\end{definition}

The decision argument is the causal past of the action, not an arbitrary scheduling prefix. Lemma~\ref{lem:sequential-worker-decision-view} identifies it with the current local view for the existing sequential-worker implementations. A participant with concurrent actions can have several causal decision views; information from another event enters a decision through its actual causal dependency. Full causal observation and payload observation are different choices of $o_p$.

\begin{proposition}[Observation descent]\label{prop:observation-descent}
A view-based permission map $\kappa_p:\mathsf H_p\to\Pow(L_p)$ descends uniquely to the attained image of $o_p$ exactly when $\kappa_p\pi_1=\kappa_p\pi_2$ on $\mathsf H_p\times_{O_p}\mathsf H_p$. The same assertion holds internally with the power object.
\end{proposition}
\begin{proof}\vmode
\leavevmode
\begin{enumerate}[label=(\roman*)]
\item \textbf{Descent through the attained image.}
The effective epimorphism onto the attained image is the coequalizer of its kernel pair.
The displayed equation is exactly its factorization condition.
\end{enumerate}
\end{proof}

\subsubsection{Required permissions and generated histories}\label{subsec:permission-realization}
A history language determines the permissions needed to realize its steps, while a controller generates the histories permitted by its decisions. The adjunction below relates these two operations and identifies the languages realizable at the specified observations.

\begin{definition}[History languages and controller order]\label{def:formal-prefix-closed-history-families}
Let $\mathsf{HistSub}(U)$ be the inclusion poset of subsets of $\Conf_f(U)$ containing the empty configuration and closed under causal prefixes. They need not be closed under combining compatible alternatives. The controller order is pointwise inclusion on $\mathsf{Ctrl}=\prod_p\Pow(L_p)^{O_p}$.
\end{definition}
\begin{definition}[Required permissions]\label{def:requirements}
For $L\in\mathsf{HistSub}(U)$ put
\begin{equation}\label{eq:required}
 R_L(p,o)=\{\nu_p(a(e)):e\in h\in L,\ p\in V(e),\ o_p(D_{e,p})=o\}.
\end{equation}
\end{definition}
\begin{lemma}[Generation respects causal prefixes]\label{lem:controlled-prefix}
The configurations of $U_k$ are exactly the histories generated by~\eqref{eq:controlled}, and $G(k)\in\mathsf{HistSub}(U)$.
\end{lemma}
\begin{proof}\vmode
\leavevmode
\begin{enumerate}[label=(\roman*)]
\item \textbf{Downward closure.}
The restriction defining $U_k$ is downward closed.

\item \textbf{Generation and its converse.}
A finite configuration of it can be linearized by Theorem~\ref{thm:generic-history-exchange}; each event passes its fixed principal-past guards.
Conversely, every generated extension passes those guards and retains its already accepted predecessors, so every generated history is a configuration of $U_k$.
Its causal prefixes use the same guards.
\end{enumerate}
\end{proof}
\begin{proposition}[Permissions and histories form an adjunction]\label{prop:permission-adjunction}
The monotone maps $R,G$ satisfy
\begin{equation}\label{eq:permission-adjunction}
 R_L\leq k\quad\Longleftrightarrow\quad L\subseteq G(k).
\end{equation}
Thus $R\dashv G$, $C=GR$ is a closure operation, $RGR=R$, and its fixed points are precisely the controller-realizable languages. The closure $C(L)$ is the least such language containing $L$.
\end{proposition}
\begin{proof}\vmode
\leavevmode
\begin{enumerate}[label=(\roman*)]
\item \textbf{Permissions imply realization.}
If $R_L\leq k$, every event in every history of $L$, including all its predecessors, passes its guard.
Hence $L\subseteq\Conf_f(U_k)$.

\item \textbf{Realization implies required permissions.}
Conversely, every permission in~\eqref{eq:required} is witnessed in a controlled history when $L\subseteq G(k)$, so belongs to $k$.

\item \textbf{Monotonicity.}
Inclusion of languages and of permission sets makes both constructions monotone.

\item \textbf{Closure and fixed points.}
The unit $L\subseteq GRL$ and counit $RGk\leq k$ give $RGR=R$ and $GRG=G$ by the two opposite inequalities.
These prove idempotence, the fixed-point assertion, and leastness.
\end{enumerate}
\end{proof}
A permission can depend on positive evidence about a payload. For a local action name $a$ and open $U_a$ of a declared view domain, set $a\in k_p(y)$ exactly when $y\in U_a$. Pullback along a Scott-continuous local observation gives an open on decision histories. Each permitted occurrence thus has finite evidence in its own causal view. Finite intersection over its vetoing participants expresses their joint permission. This instance supplies the same $g_k$ to the following realization theorem and retains the causal view at which each commitment is made.

\begin{theorem}[Observation-sensitive realization]\label{thm:observation-sensitive-realization}
A candidate language $L$ is generated by a controller exactly when every enabled extension $h\to h\cup\{e\}$ leaving $L$ satisfies
\begin{equation}\label{eq:refusal}
 \exists p\in V(e):\quad \nu_p(a(e))\notin R_L(p,o_p(D_{e,p})).
\end{equation}
Equivalently $GRL=L$, and the least realizing controller is $R_L$.
\end{theorem}
\begin{proof}\vmode
\leavevmode
\begin{enumerate}[label=(\roman*)]
\item \textbf{Necessity of the refusal condition.}
For a realizable language the adjunction gives $GRL=L$, so an omitted extension cannot pass all its guards.

\item \textbf{First-exit sufficiency.}
Conversely a finite history of $GRL$ outside $L$ has a first extension leaving $L$ in any linearization; all its guards pass, contradicting~\eqref{eq:refusal}.

\item \textbf{Equality and leastness.}
The unit supplies the other inclusion.
The adjunction proves leastness, including empty permission sets at unused observations.
\end{enumerate}
\end{proof}
\begin{corollary}[Occurrence subprocesses]\label{cor:occurrence-permissions}
For a downward-closed occurrence subset $T\subseteq U$ with inherited consistency, its configurations are controller-realizable exactly when they pass~\eqref{eq:refusal}. Their required permissions are already witnessed by $[e)\to[e]$, $e\in T$.
\end{corollary}
\begin{proof}\vmode
\leavevmode
\begin{enumerate}[label=(\roman*)]
\item \textbf{Principal configurations.}
Principal configurations belong to $T$, and each decision view depends only on that principal past.

\item \textbf{Recovery of the occurrence subset.}
An event belongs to a downward-closed subset exactly when its principal configuration does, so equality of configuration families gives equality of the occurrence subprocesses.
\end{enumerate}
\end{proof}
\begin{proposition}[Locally exposed control names]\label{prop:localnames}
The preceding results hold for arbitrary $\nu_p$, including noninjective exposed names. With identity names they recover the global-action formulas. No additional authority to distinguish hidden witnesses is introduced by compilation.
\end{proposition}
\begin{proof}\vmode
\leavevmode
\begin{enumerate}[label=(\roman*)]
\item \textbf{Dependence on exposed names.}
Every guard and required union uses only $\nu_p(a(e))$.
The two implications of the adjunction and the first-exit argument do not use injectivity.
The compiler below tests precisely these same membership statements.
\end{enumerate}
\end{proof}

\subsubsection{The causal-view commitment compiler}
The permission calculation describes which physical histories a controller allows. We now realize those decisions as positive commitments in a control arena. The comparison must account for the additional commitment events while preserving the physical histories and their independent extensions.

\begin{definition}[Permission skeleton and compiled arena]\label{def:causal-view-control-arena}
For $p\in P$ and $v\in\mathsf H_p$ introduce a positive slot $c_{p,v}$; for each $e\in U$ introduce a negative physical event $e^-$. On this skeleton $B_U$ generate order by
\[
 d^-<e^-\ (d<e),\qquad d^-<c_{p,v}\ (d\in v),\qquad
 c_{p,D_{e,p}}<e^-\ (p\in V(e)).
\]
Put $r(e^-)=[e]_U$ and $r(c_{p,v})=v$. A finite skeleton family $F$ is consistent when $\bigcup_{b\in F}r(b)\in\mathsf{Con}_U$.

An event of the compiled arena $\mathcal C(U)$ is $(b,\lambda)$, where $\lambda$ assigns a set $K\subseteq L_p$ to every slot $c_{p,v}\leq b$ in the skeleton past and satisfies
\[
 \nu_p(a(e))\in\lambda(c_{p,D_{e,p}})
 \quad\text{for every }e^-\leq b\text{ and }p\in V(e).
\]
Order is restriction of decorated principal pasts. Finite consistency means skeleton consistency and agreement of the assignments on overlaps. Polarity is that of $b$.
\end{definition}
\begin{lemma}[The compiled arena is an event structure]\label{lem:permission-arena-event-structure}
Definition~\ref{def:causal-view-control-arena} constructs a polarized event structure with finite principal pasts. Its induced order on physical events is exactly that of $U$.
\end{lemma}
\begin{proof}\vmode
\leavevmode
\begin{enumerate}[label=(\roman*)]
\item \textbf{Acyclic skeleton order.}
A segment $d^-<c_{p,v}<e^-$ has $v=D_{e,p}$ and $d\in[e)$, so projects to $d<e$ in $U$.
Any putative skeleton cycle would therefore project to a strict physical cycle; a cycle cannot consist only of slots.
The order is acyclic and induces exactly the original physical order.

\item \textbf{Finite skeleton pasts.}
A physical event has only finitely many physical ancestors and finitely many required slots per ancestor.
A slot has finite support $v$ and the finite pasts of its members.
Thus all skeleton pasts are finite.

\item \textbf{Skeleton consistency.}
Predecessor supports are contained in successor supports, so the consistency test is hereditary and closed under adjoining predecessors; singleton supports are configurations.

\item \textbf{Existence of decorations.}
Every principal past admits a decoration by assigning each slot its full alphabet.
Restriction preserves the guards.

\item \textbf{Decorated principal pasts.}
The decorated principal past of $(b,\lambda)$ consists exactly of these restrictions, and is finite.

\item \textbf{Consistency of decorations.}
Consistency is hereditary, and adjoining a predecessor adds a restriction already agreeing with every compatible assignment.

\item \textbf{Separation and event-structure laws.}
Two distinct decorations of one skeleton event disagree at a slot in their common domain and cannot coexist.
These statements verify all event-structure laws and the claimed physical-order calculation.
\end{enumerate}
\end{proof}
\begin{definition}[Selection and erasure]\label{def:formal-commitment-erasure-and-selection}
Select the slot value
\begin{equation}\label{eq:selectedcommit}
 \lambda(c_{p,v})=k_p(o_p(v)).
\end{equation}
Let $S_k\subseteq\mathcal C(U)$ contain precisely the events whose entire decorated past has these selected values. Erasure forgets slots and decorations and retains the physical occurrences. Its history map is denoted $\epsilon$.
\end{definition}
\begin{lemma}[Selected controller display]\label{lem:restriction}
The subset $S_k$ is downward closed and its inclusion in $\mathcal C(U)$ is receptive and courteous.
\end{lemma}
\begin{proof}\vmode
\leavevmode
\begin{enumerate}[label=(\roman*)]
\item \textbf{Downward closure.}
Restrictions of selected assignments are selected.

\item \textbf{Receptivity.}
An enabled negative event has all its slot choices among its already present predecessors, so every such extension remains selected uniquely.

\item \textbf{Courtesy.}
The inclusion adds no immediate causal edge, since every intervening predecessor would also be selected.
These are receptivity and courtesy.
\end{enumerate}
\end{proof}
\begin{theorem}[Selected configuration classification]\label{thm:permission-configuration-classification}
For $h\in\Conf_f(U)$ put
\[
 \mathsf{Req}(h)=\{(p,D_{e,p}):e\in h,\ p\in V(e)\},\qquad
 \mathsf{Avail}(h)=\{(p,v):v\in\mathsf H_p,\ v\subseteq h\}.
\]
There is an order isomorphism
\begin{equation}\label{eq:permission-configurations}
 \Conf_f(S_k)\cong
 \{(h,J):h\in\Conf_f(U_k),\ J\text{ finite},\
                  \mathsf{Req}(h)\subseteq J\subseteq\mathsf{Avail}(h)\},
\end{equation}
where the right-hand order is inclusion in both coordinates and $\epsilon(h,J)=h$.
\end{theorem}
\begin{proof}\vmode
\leavevmode
\begin{enumerate}[label=(\roman*)]
\item \textbf{Projection to histories and slots.}
Project a selected configuration to its physical events $h$ and slots $J$.
Local uniqueness of decorations makes these sets.

\item \textbf{Required and available slots.}
Downward closure gives physical downward closure, every required slot, and the support of each retained slot inside $h$.
Consistency gives consistency of $h$; the guards imply that every event of $h$ belongs to $U_k$.

\item \textbf{Reconstruction of decorated events.}
Conversely, for such $(h,J)$, decorate each retained skeleton event by the selected assignment on its past.
Its physical ancestors are in $h$ and its required slots are in $\mathsf{Req}(h)$; the same is true for the physical ancestors of any slot.
Thus every principal past is present and all guards hold.

\item \textbf{Consistency.}
Assignments agree, and the union of supports is contained in the consistent configuration $h$.
This constructs a selected configuration.

\item \textbf{Inverse order correspondence.}
Projection and reconstruction are inverse and preserve inclusion, giving the stated order isomorphism.
\end{enumerate}
\end{proof}
\begin{theorem}[Execution and concurrency preservation]\label{thm:compiler}
Erasure relates selected compiled execution and controlled execution by a rooted weak bisimulation. It preserves physical causal order and lifts every finite controlled history and every finite enabled concurrent family from each existing selected lift. Preparing $e$ takes at most $|V(e)|$ positive steps; preparing a path $e_1\cdots e_n$ takes at most $\sum_i|V(e_i)|$.
\end{theorem}
\begin{proof}\vmode
\leavevmode
\begin{enumerate}[label=(\roman*)]
\item \textbf{Forward execution erasure.}
A slot addition erases to an identity.
A negative extension projects, by~\eqref{eq:permission-configurations}, to a controlled physical extension.

\item \textbf{Preparation of a controlled step.}
Given $e\in\En_k(h)$ at a lift $(h,J)$, add the missing slots $(p,D_{e,p})$ for $p\in V(e)$.
Their supports lie in $[e)\subseteq h$, and their other predecessors are already among the requirements of $h$.
They are mutually consistent and need no new physical event.

\item \textbf{Weak simulation and path lifting.}
The resulting configuration admits $e^-$, proving the reverse weak simulation and the bound.
Iteration gives path lifting.

\item \textbf{Preparation of a concurrent family.}
For a finite jointly enabled family $F$, prepare the union of its required slots.

\item \textbf{Lifted cube and physical order.}
With this fixed enlarged slot set $J'$, every $(h\cup F',J')$, $F'\subseteq F$, satisfies~\eqref{eq:permission-configurations}.
These are the faces of the lifted cube.
Shared control owners merely supply shared predecessors, and do not order its physical events.
Lemma~\ref{lem:permission-arena-event-structure} gives exact preservation of physical causal order.
\end{enumerate}
\end{proof}

\begin{construction}[Erasure bisimulation and prepared cubes]\vmode\label{con:compiler-bisimulation-cubes}
\leavevmode
\begin{enumerate}[label=(\roman*)]
\item \textbf{Transition labels and comparison relation.}
Label a selected commitment extension by $\tau$ and a plant extension
by its retained physical event.
Controlled histories have only the
latter transitions.
The relation is
\[
 \mathcal R=\{((h,J),h):
   \mathsf{Req}(h)\subseteq J\subseteq\mathsf{Avail}(h),\
   J\text{ finite}\}.
\]
A $\tau$-extension has unchanged second projection.

\item \textbf{Forward physical step.}
A compiled
plant extension has exactly the corresponding controlled extension.

\item \textbf{Reverse weak lift.}
Conversely, for a controlled $e$ define
$J_e=J\cup\{(p,D_{e,p}):p\in V(e)\}$.
Adding the missing slots and then $e$ is the required
$\tau^*e$ lift; the initial pair is $(\varnothing,\varnothing)$.
These are the rooted weak-bisimulation clauses with their labels.

\item \textbf{Prepared concurrent vertices.}
For a finite jointly enabled family $F$, put
\[
 J_F=J\cup\bigcup_{e\in F}\{(p,D_{e,p}):p\in V(e)\},\qquad
 \widehat c_F(G)=(h\cup G,J_F)\quad(G\subseteq F).
\]
All new slots are supported in $h$, and the configuration
classification proves every displayed vertex exists.

\item \textbf{Edges and erasure.}
The edge from
$G$ to $G\cup\{e\}$ is the selected occurrence of $e$.
Erasure satisfies $\epsilon\widehat c_F(G)=h\cup G$.

\item \textbf{Faces and preparation bound.}
For disjoint $B,G\subseteq F$, the face vertex is
$(h\cup B\cup G,J_F)$; restriction and successive execution
therefore obey the face identities by associativity of union.
The preparation precedes the cube and adds at most the stated sum
of veto-support sizes.
Shared slots are added once.
\end{enumerate}
\end{construction}

The comparison is an execution erasure. Plant events are negative in the control arena; erasure is not asserted to be a polarity-preserving strategy map to the original message interface. The reusable-register implementation and its global finite-participant bound are proved in Appendix~\ref{subsec:register-control-implementation}. Arbitrary set-valued policies are semantic parameters; an effective implementation uses an actual coded controller and the partial-computation interpretation of Section~\ref{sec:effective-foundations}.

\subsubsection{The kernel of a distributed controller}
\begin{definition}[Controller configuration kernel]\label{def:controller-configuration-kernel}
For $x\in\Conf_f(\mathcal C(U))$, put $\Phi_k(x)=\{*\mid x\in\Conf_f(S_k)\}$, with root at the empty configuration.
\end{definition}
\begin{theorem}[Distributed strategy kernel]\label{thm:controllerkernel}
The displayed subterminal fibres define a pointed reached Scott kernel, with cubical unstraightening $K(S_k)\to K(\mathcal C(U))$.
\end{theorem}
\begin{proof}\vmode
\leavevmode
\begin{enumerate}[label=(\roman*)]
\item \textbf{Controller fibre.}
Lemma~\ref{lem:restriction} and Theorem~\ref{thm:scott-discrete-fibrations} give exactly these fibres and their restriction maps.

\item \textbf{Kernel restrictions.}
Removing positive events leaves selected choices, and any allowed negative extension remains selected, so the support is closed under Scott restriction.

\item \textbf{Reachability.}
Every finite selected configuration has a linearization of selected prefixes, proving reachability.
\end{enumerate}
\end{proof}
\begin{proposition}[Decoding the constructed image]\label{prop:constructed-image-decoding}
A downward-closed arena restriction containing every enabled negative extension, choosing exactly one permission set at each reachable slot, and choosing the same set at equal declared observations, determines a controller on its reached observations. Recompilation reconstructs that restriction and its kernel; unreached values may be assigned the empty set.
\end{proposition}
\begin{proof}\vmode
\leavevmode
\begin{enumerate}[label=(\roman*)]
\item \textbf{Uniform decoding.}
Observation uniformity makes the read-off value a function.

\item \textbf{Equality of configurations.}
Induction along a finite linearization shows that the original restriction and the recompiled one have the same extensions: positive ones are the same unique selected choices and negative ones are all arena-enabled extensions.
Thus their finite configurations coincide, and each event is determined by its finite principal configuration.

\item \textbf{Stability under recompilation.}
No newly reached slot appears under recompilation, by the same induction.
\end{enumerate}
\end{proof}

\subsection{Winning conditions and information-preserving implementation}\label{subsec:winning-information-preserving-implementation}
The compiler comparison lifts physical executions. Winning transfer also needs to account for commitment completion and maximality. The following statements retain those requirements explicitly, so that the implementation comparison supports the stated winning conclusions.

\begin{definition}[Safety and controlled deadlock]\label{def:formal-safety-and-controlled-deadlock}
Let $B\subseteq\Conf_f(U)$ be the bad histories. A controller is safe when every history in $G(k)$ avoids $B$; causal prefixes already belong to $G(k)$. It creates no new deadlock when, for every $h\in G(k)$,
\begin{equation}\label{eq:deadlock}
 \En_U(h)\ne\varnothing\quad\Longrightarrow\quad\En_k(h)\ne\varnothing.
\end{equation}
For a concrete state evaluation $q$, a bad-state predicate is pulled back along $q$. Existing plant termination is allowed.
\end{definition}
\begin{definition}[Positive maximality and physical winning]\label{def:formal-positive-maximality-and-winning}
A configuration is positive-maximal when it has no positive extension. A strategy is winning when the images of all its positive-maximal configurations satisfy the arena winning condition, as in~\cite{winning}. Erasure on arbitrary configurations is the union of the erasures of finite principal histories. Let $\mathsf{Complete}(x)$ mean that every slot available at its physical history is present with a choice, and let $\En_\mathcal C(x)$ be its enabled negative extensions. Define, independently of a controller,
\begin{equation}\label{eq:winning}
 \begin{aligned}
 x\in W_{\mathrm{phys}}\quad\Longleftrightarrow\quad&
 \text{every finite causal prefix of }\epsilon x\text{ avoids }B,\\
 &\forall h\in\Conf_f(U),\quad\epsilon x=h\ \Longrightarrow\\
 &\qquad\bigl(\mathsf{Complete}(x)\ \land\
   (\En_U(h)\ne\varnothing\Rightarrow\En_\mathcal C(x)\ne\varnothing)\bigr).
 \end{aligned}
\end{equation}
\end{definition}
\begin{definition}[Positive saturation]\label{def:positive-saturation}
For a finite selected lift $(h,J)$ set $\mathsf{Sat}(h,J)=(h,\mathsf{Avail}(h))$, interpreted as the directed union of all finite additions to $J$ when necessary.
\end{definition}
\begin{lemma}[Canonical positive saturation]\label{lem:controller-saturation}
The saturation is a selected positive-maximal configuration with the same erasure. Every selected positive-maximal configuration with finite erasure $h$ is complete, and its negative extensions correspond bijectively to $\En_k(h)$.
\end{lemma}
\begin{proof}\vmode
\leavevmode
\begin{enumerate}[label=(\roman*)]
\item \textbf{Completion of available slots.}
Every finite set of available slots can be adjoined by~\eqref{eq:permission-configurations}; their directed union is a configuration by finite consistency and downward closure.
No further slot has its physical support inside $h$, so no positive extension remains.

\item \textbf{Maximality criterion.}
Conversely, any missing available slot has its physical past in $h$ and all the slots required by those physical events already present.
Its selected decoration is therefore an enabled positive extension, contradicting maximality.

\item \textbf{Physical extension at completeness.}
At completeness all requirements of a physically enabled $e$ are present; its negative decoration exists exactly when its controller guard passes.

\item \textbf{Uniqueness and infinite configurations.}
Uniqueness follows from the existing slot choices.
This argument also applies to an infinite selected configuration with finite physical erasure.
\end{enumerate}
\end{proof}
\begin{theorem}[Winning transfer]\label{thm:winning-transfer}
A controller is safe and creates no new deadlock exactly when $S_k$, equivalently its kernel $\Phi_k$, is winning for $W_{\mathrm{phys}}$.
\end{theorem}
\begin{proof}\vmode
\leavevmode
\begin{enumerate}[label=(\roman*)]
\item \textbf{Transfer from physical safety.}
Every finite physical prefix of a selected configuration belongs to $G(k)$ by the classification.
Safety therefore gives the first clause of~\eqref{eq:winning}; saturation completeness and the enabled-event bijection give its finite-erasure clause from~\eqref{eq:deadlock}.

\item \textbf{Recovery from compiled winning.}
Conversely lift each $h\in G(k)$ by $(h,\mathsf{Req}(h))$ and saturate it.
Winning at this configuration gives safety of $h$ and, by the bijection, no new deadlock.

\item \textbf{Infinite erasures.}
Infinite physical erasures are tested through their finite prefixes, exactly as specified by the arena condition.
\end{enumerate}
\end{proof}
For finite $P$ and finite $h$, there are only finitely many pairs $(p,v)$ with $v\subseteq h$, so saturation is finite. Its bound is the number of missing available slots; the separate register implementation has the sharper current-register bound $|P|$. Unique permission selection does not require a single allowed physical action: per-participant action determinism is the additional winning predicate $|\{e\in\En_\mathcal C(x):p\in V(e)\}|\leq1$, transported by the enabled-event bijection.

\subsubsection{Realizing full causal memory}
\begin{definition}[Causal-history instrumentation]\label{def:generic-causal-certificates}
Retain the original occurrence identifiers and labels. For each $e$, construct a finite certificate by taking the union of the certificates at its immediate predecessors, identifying shared occurrence identifiers, and adjoining $e$ with its original incoming edges. The order is the transitive closure. A participant's accumulated view certificate is the union at its participating occurrences, with the inherited order. The enriched presentation attaches these certificates to its events and dependency records; erasure forgets them.
\end{definition}
\begin{proposition}[Exact causal certificates]\label{prop:generic-causal-memory}
The certificate at $e$ is canonically $[e]_U$, and the accumulated certificate at $p$ in $h$ is $v_p(h)$. Erasure of the enriched presentation is an event-structure isomorphism and preserves all configuration cubes. Full-causal-view controllers are observation-based on the enriched presentation.
\end{proposition}
\begin{proof}\vmode
\leavevmode
\begin{enumerate}[label=(\roman*)]
\item \textbf{Inductive reconstruction of the past.}
Induct on the finite causal depth of $e$.
The strict past of $e$ is the union of the principal pasts of its immediate predecessors.
Their occurrence identifiers identify exactly their shared subhistories, so the union and the new top reconstruct $[e]$ with its order.

\item \textbf{Participant view.}
The union over $p$-events is exactly the defining downclosure of $v_p(h)$.

\item \textbf{Enrichment and controller realization.}
Enrichment adds a uniquely determined label without changing events, order or consistency, proving the isomorphism and cube statement.
Reading the reconstructed view supplies the controller argument.
\end{enumerate}
\end{proof}
An operational implementation transmits this metadata along its retained dependency or resource-supply records. The concrete send/receive and joint-action implementations are in Appendix~\ref{subsec:register-control-implementation}. This instrumentation changes the information exposed by the protocol; it does not identify full causal observation with the original payload observation.

\subsubsection{Independent composition of controllers}
\begin{proposition}\label{prop:parcontrol}
For independent presentations with disjoint participant tags, the compiled arenas and selected strategies split by tensor, and $\Phi_{k\amalg k'}(x,x')\cong\Phi_k(x)\times\Phi_{k'}(x')$. Conjunctive safety and componentwise absence of new deadlocks are preserved.
\end{proposition}
\begin{proof}\vmode
\leavevmode
\begin{enumerate}[label=(\roman*)]
\item \textbf{Componentwise configuration structure.}
Local views, slots, physical orders and consistency split into the two components.
Decorations and guards therefore split as well.
Formula~\eqref{eq:permission-configurations} gives the configuration product, and the kernel formula follows.

\item \textbf{Enabling and safety.}
Each enabled product event belongs to one component, where the controlled extension exists.
Safety is checked in each component.

\item \textbf{Tensor coherence.}
Tagged rebracketing gives the tensor coherences.
\end{enumerate}
\end{proof}
Connecting ports instead changes the interaction and observations. Its strategy interpretation is the secured composition of Theorem~\ref{thm:abstract-communication-interpretation}; winning composition uses the appropriate game winning specification~\cite{winning}. Independent safety alone does not prove safety of arbitrary wiring.

\subsection{Calculus: observation-typed controllers and commitments}\label{subsec:calculus-control}
The permission functions and compiled slots now have process rules. Their parameters record the available observations, and their evidence records the selected permissions; the earlier compiler theorem supplies the resulting history and cube comparison.

\begin{definition}[Observation-typed control rules]\label{def:observation-typed-control-rules}
The parameter type is $\mathsf{Ctrl}_p=\Set(O_p,\Pow(L_p))$, with evaluation $\mathsf{permissions}(\mathsf{ctrl}(k_p),o)=k_p(o)$. At a slot $(p,v)$, its positive command commits the whole selected set $k_p(o_p(v))$. The derived offer $\mathsf{allow}_p(o,\ell)$ carries membership evidence $\ell\in k_p(o)$. For an enabled occurrence $e$ after $h$, the permission-evidence type is
\[
 \mathsf{Permit}_k(h,e)=\prod_{p\in V(e)}
 \{*\mid\nu_p(a(e))\in k_p(o_p(D_{e,p}))\}.
\]
A controlled step retains its plant witness and an element of this product. Its compiled rule reads the required slots, executes the physical event, and retains their causal records. A semantic controller parameter need not have an effective program; coded realizations use the existing lift interpretation for computations that may not return.
\end{definition}
\begin{proposition}[Process reading of observation descent]\label{prop:pc-controller}
A view-indexed controller expression defines an observation-indexed one exactly when its decisions are constant on observation fibres. Its compiled process realizes~\eqref{eq:controlled}, preserves uncontrollable enabling, and has the history and cube comparisons of Theorem~\ref{thm:compiler}.
\end{proposition}
\begin{proof}\vmode
\leavevmode
\begin{enumerate}[label=(\roman*)]
\item \textbf{Observation descent and guards.}
Proposition~\ref{prop:observation-descent} constructs the unique function on attained observations.
Evaluation and membership elimination give exactly the guard product, with the singleton empty product for an uncontrollable event.

\item \textbf{Selected strategy and readback.}
The selected arena has these positive choices and all its negative extensions by Lemma~\ref{lem:restriction}.
The occurrence representation theorem gives its process declaration, and Theorem~\ref{thm:compiler} gives the physical comparison.

\item \textbf{Substitution compatibility.}
Evaluation, product elimination and substitution commute componentwise, so these are also the computation and substitution rules of the declared control terms.
\end{enumerate}
\end{proof}

\paragraph{Permissions and observations.}
For a history language $L$, the permission adjunction $R\dashv G$ gives $G(R_L)$ as the least controller-realizable language containing $L$. The winning-transfer theorem uses the lifting of executions and its designated winning configurations. Contextual may tests compare which finite successes a completed implementation can exhibit.


\section{Witnessed execution, relational feedback, and time}\label{sec:execution-games}
A connected strategy records the causal past of its visible events. To interpret iteration and recursion, we also need the successive operational transitions through which execution proceeds. This section constructs chronological histories, finite return and delayed state closure, and compares their execution witnesses with the existing game and kernel operations. The offers and medium commitments of $\pcLang[\Pi]$ supply the transitions. Compact connection, routed return and delayed closure retain their distinct types throughout the comparison.\begin{definition}[Compact feedback]\label{def:formal-compact-feedback}
The compact structure of $\mathsf{CG}$ constructs a trace. For $f:A\parallel U\to B\parallel U$, define
\[
 \operatorname{Tr}^{U}_{\parallel}(f)=
 (1_B\parallel\epsilon_U)(f\parallel1_{U^\ast})
 (1_A\parallel\eta_U),
\]
with the canonical symmetries understood. This is a composite of secured composition, cup and cap. It does not count how many internal operational transitions occurred.
\end{definition}

\subsection{Chronological execution and generated time}\label{subsec:chronological-execution-generated-time}
Chronological paths record successive extensions of an execution. Their lengths and prefix maps generate the temporal indexing used below, while the original concurrent cells continue to record independent events.

\begin{definition}[Chronological path objects]\label{def:formal-chronological-path-objects}
Take the one-event extension graph $C=(S\xleftarrow{s}D\xrightarrow{t}S)$ of a game or strategy. Its length-$n$ history object $H_n(C)$ consists of $n$ composable witnessed extensions, with $H_0(C)=S$. Deleting the last extension gives its chronological prefixes. These constructions also make sense internally in a Grothendieck topos. Concurrent cells remain the cells of the game; counting chronological steps does not impose an execution order on them.
\end{definition}

\subsubsection{Duration, prefixes, and the generated temporal topos}
\begin{definition}[Duration and translation clock]\label{def:formal-duration-and-translation-clock}
The free category on one vertex and one loop has arrow object
$L=\coprod_{n\geq0}1$ with composition addition. Every free-path category has a duration functor to this category, sending $H_n(C)$ to $n$. Its translation category has objects $r\in L$ and arrows $(r,k):r\to r+k$. It is the internal ordinal $\omega_{\E}$: an arrow $r\to s$ exists exactly when $r\leq s$, and then $k$ is uniquely determined. This constructs the lower-d\'ecalage clock directly from chronological length.
\end{definition}

\begin{definition}[Temporal objects and later]
Write
\[
 \mathcal T(\E)=\Fun(\omega^{\op},\E).
\]
Write $\Delta S$ for the constant tower on $S$. A \textbf{temporal object} is a tower $X_0\xleftarrow{d_0}X_1\xleftarrow{d_1}\cdots$. Its generated later object and its next map are
\[
 (\triangleright X)_0=1,\qquad
 (\triangleright X)_{n+1}=X_n,\qquad
 (\mathsf{next}_X)_0={!},\quad
 (\mathsf{next}_X)_{n+1}=d_n.
\]
The restriction from degree one to zero is terminal, and the other restrictions are inherited from $X$. These definitions also apply in $\Fun(\omega^{\op},\E/S)$, replacing $1$ by $1_S$.
\end{definition}

\begin{definition}[History-tower maps]\label{def:history-tower-maps}
For the already constructed finite path objects $H_n(C)$, define the \textbf{prefix map} $\mathsf{prefix}_n:H_{n+1}(C)\to H_n(C)$ by deleting the final edge and the \textbf{initial-endpoint map} $s_n:H_n(C)\to S$ by the first vertex. The \textbf{history tower} $\mathsf{Hist}(C)$ has degree $n$ object $H_n(C)$ and these prefix restrictions. Its \textbf{coherent-execution object} is $\lim_nH_n(C)$, with its limit projections. Proposition~\ref{prop:historytower} verifies these maps and their universal property.
\end{definition}

\begin{proposition}[Generated delay and history tower]\label{prop:historytower}
The later functor is right adjoint to successor precomposition $X\mapsto(X_{n+1})_n$ and preserves limits. The histories of $C$ form a temporal object
\[
 \mathsf{Hist}(C)_n=H_n(C)
\]
by deleting the last transition. It lies over the constant object $S$ via the initial endpoint. Its internal object of coherent infinite executions is $\lim_nH_n(C)$. Its set of global sections is
\[
 \Gamma_{\mathcal T(\E)}\mathsf{Hist}(C)
 \cong\Hom_{\E}\bigl(1,\lim_nH_n(C)\bigr).
\] Finite histories ending in a deadlock are retained at their finite stages, whether or not they extend to a global section.

The initial endpoint is a natural transformation to the constant tower:
\[
\begin{tikzcd}[column sep=large]
H_{n+1}(C) \arrow[r,"{\mathsf{prefix}}"] \arrow[d,"{s_{n+1}}"']
 &H_n(C) \arrow[d,"{s_n}"]\\
S \arrow[r,equal] &S .
\end{tikzcd}
\]
Thus the clock forgets the last transition without forgetting the execution's entry state.

\end{proposition}
\begin{proof}\vmode
\leavevmode
\begin{enumerate}[label=(\roman*)]
\item \textbf{Successor adjunction.}
A tower map $Y\to\triangleright X$ has a forced degree-zero component and otherwise consists of compatible maps $Y_{n+1}\to X_n$.
These are exactly maps from the successor restriction of $Y$ to $X$, giving the adjunction.

\item \textbf{Prefix restriction.}
Deleting a last edge preserves the initial endpoint and iterated deletions satisfy the tower equations.

\item \textbf{Global history sections.}
A map from the terminal tower into this tower is, by the universal property of a limit, a compatible family of finite paths.
Compatibility says precisely that each is the prefix of the next.
No surjectivity of the restriction maps is used.
\end{enumerate}
\end{proof}

\subsubsection{Guarded equations derived from the clock}
The prefix clock supplies the delay operation. We now use its degree-by-degree structure to solve guarded equations and identify the resulting fixed histories.

\begin{definition}[Guarded temporal map]\label{def:guarded-map}
A temporal equation $e:Z\times V\to V$ is \textbf{generated guarded in $V$} when it factors through $1_Z\times\mathsf{next}_V$. This is a property of $e$; a factorization is not part of the retained equation. A fixed section is a map $u:Z\to V$ with $u=e\langle1_Z,u\rangle$.
The \textbf{fixed-section set} is
\[
 \mathsf{Fix}(e)=\{u\in\Hom_{\mathcal T(\E)}(Z,V):u=e\langle1_Z,u\rangle\}.
\]
The notation $e^{\ddagger}$ denotes its unique element when it is a singleton. For $f:Z\times V\to O\times V$ with $e=\pi_Vf$, its \textbf{visible execution}, once this singleton property holds, is $\mathsf{fb}_V(f)=\pi_Of\langle1_Z,e^{\ddagger}\rangle:Z\to O$. Theorem~\ref{thm:guarded-execution} constructs and proves uniqueness for every generated guarded equation.
\end{definition}

\begin{definition}[Hidden-state closure]\label{def:formal-hidden-state-closure}
For a map $f:Z\times V\to O\times V$ between temporal objects, define the \textbf{hidden-state closure of its discrete companion} to be the span with carrier $\{(z,v)\mid f_V(z,v)=v\}$, visible input $z$, and visible output $f_O(z,v)$. This is a span between discrete temporal interfaces and therefore has the response-game interpretation below.
\end{definition}

The following is the parameterized guarded fixed-point construction of \cite[Theorem~2.4]{birkedal-guarded}; we give its degreewise form to verify the present indexing and its interpretation in $\E$.
\begin{theorem}[Unique generated guarded execution]\label{thm:guarded-execution}
Every generated guarded equation has a unique fixed section $e^{\ddagger}:Z\to V$, natural under substitution of parameters. Consequently a temporal map
$f:Z\times V\to O\times V$ with guarded hidden component $e=\pi_V f$ has a canonical visible execution
\[
 \mathsf{fb}_V(f)=\pi_O f\langle1_Z,e^{\ddagger}\rangle:Z\to O.
\]
Closing the hidden state of its companion $f_!$ between discrete temporal interfaces gives precisely the companion of $\mathsf{fb}_V(f)$.
\end{theorem}
\begin{proof}\vmode
\leavevmode
\begin{enumerate}[label=(\roman*)]
\item \textbf{Degreewise solution.}
Temporarily display $e=\bar e(1\times\mathsf{next})$.
Define
\[
 u_0=\bar e_0\langle1,!\rangle,\qquad
 u_{n+1}=\bar e_{n+1}\langle1,u_n d^Z_n\rangle.
\]
Naturality of $\bar e$ gives $d^V_0u_1=u_0d^Z_0$.

\item \textbf{Tower compatibility.}
Inductively it gives $d^V_nu_{n+1}=u_nd^Z_n$, because the corresponding equality one degree lower identifies the delayed arguments.
Thus $u$ is a tower map.

\item \textbf{Fixed-point equation.}
Substitution of the formula for $\mathsf{next}$ proves the fixed-section equation at every degree.

\item \textbf{Uniqueness.}
Any other solution has the same degree-zero component and, by induction, the same component in each subsequent degree.
Uniqueness proves independence of the temporary factorization and parameter naturality.

\item \textbf{Hidden-state equalizer.}
For the last assertion, the companion carrier is $Z\times V$.
Its hidden-state closure is the equalizer object
\[
 F=\{(z,v)\mid e(z,v)=v\}\longrightarrow Z.
\]
For any temporal object $T$ and map $a:T\to Z$, maps $T\to F$ over $a$ are exactly fixed sections of $e(a\times1)$, hence form a singleton by the result just proved.

\item \textbf{Visible feedback.}
The projection $F\to Z$ is therefore an isomorphism by Yoneda.
Its visible value is $\pi_Of(z,e^{\ddagger}(z))$, which proves the companion identification.

\item \textbf{Universal property.}
The universal property used here can be recorded as
\[
\begin{tikzcd}[column sep=large]
Z \arrow[r,"{\langle1,e^{\ddagger}\rangle}"] \arrow[dr,"{e^{\ddagger}}"']
 &Z\times V \arrow[d,"{e}"]\\
 &V
\end{tikzcd}
\qquad
\begin{tikzcd}[column sep=large]
F \arrow[r,hook,"{j}"] \arrow[d,"{p}","{\cong}"']
 &Z\times V \arrow[r,shift left=1.5,"{e}"] \arrow[r,shift right=1.5,"{\pi_V}"'] &V\\
Z \arrow[ur,"{\langle1,e^{\ddagger}\rangle}"'] &&
\end{tikzcd}
\]
Here $j$ is the equalizer of $e$ and $\pi_V$, and $p=\pi_Zj$.
Its inverse sends $z$ to the unique fixed pair over $z$.
\end{enumerate}
\end{proof}

\begin{lemma}[Constructed simultaneous elimination]\label{lem:guardedelimination}
If $e=(e_U,e_V):Z\times U\times V\to U\times V$ is generated guarded in $U\times V$, solving the $V$ equation first constructs a generated guarded equation in $U$. Its iterated solution is the unique simultaneous solution. The construction commutes with parameter substitution.
\end{lemma}
\begin{proof}\vmode
\leavevmode
\begin{enumerate}[label=(\roman*)]
\item \textbf{Second-component solution.}
Display $e(z,u,v)=\bar e(z,\mathsf{next}_Uu,\mathsf{next}_Vv)$, using the product comparison for later.
With $(z,u)$ as parameter, the second component is guarded in $v$, so has its unique solution $v^\ddagger:Z\times U\to V$.

\item \textbf{Factorization after substitution.}
For any tower map $h:Z\times U\to V$, construct $k:Z\times\triangleright U\to\triangleright V$ by $k_0={!}$ and
\[
 k_{n+1}(z,u')=h_n(d^Z_nz,u').
\]
Naturality of $h$ verifies the tower equations and gives
$k(1\times\mathsf{next}_U)=\mathsf{next}_Vh$.
Use this construction for $h=v^\ddagger$.
The substituted equation has the factorization
\[
 e_U(z,u,v^\ddagger(z,u))
 =\bar e_U\bigl(z,\mathsf{next}_Uu,k(z,\mathsf{next}_Uu)\bigr).
\]
It is therefore guarded in $u$ and has a unique solution $u^\ddagger$.

\item \textbf{Joint fixed point.}
The pair
$(u^\ddagger(z),v^\ddagger(z,u^\ddagger(z)))$ satisfies both original equations, hence agrees with their unique joint solution by Theorem~\ref{thm:guarded-execution} applied to $U\times V$.

\item \textbf{Uniqueness and finite iteration.}
Uniqueness removes dependence on the displayed factorization and proves parameter naturality.
Induction gives the same result for any finite jointly guarded family.
\end{enumerate}
\end{proof}

This constructs the guarded execution used here, rather than selecting a fixed-point operator independently. Its relationship with the general guarded fixpoint--trace correspondence is described in \cite{miliuslitak,goncharovschroeder}.

\subsubsection{First exit with all witnesses retained}
For routed execution, the relevant result is a first return to the external sector. The construction below retains each finite path leading to that exit, including its internal transitions, so that later comparisons can act on the execution witnesses.

\begin{definition}[Block correspondences and return witnesses]\label{def:formal-block-correspondences-and-return-witnesses}
A \textbf{correspondence} $R:X\rightsquigarrow Y$ means a span $X\leftarrow R\to Y$. We write $R;Q$ for pullback composition in chronological order and $\boxplus$ for the coproduct of carriers with fixed endpoints. Let
\[
 \alpha:X+U\rightsquigarrow Y+U
\]
be a correspondence. Extensivity decomposes it into four blocks
\[
 D:X\rightsquigarrow Y,\quad E:X\rightsquigarrow U,
 \quad C:U\rightsquigarrow U,\quad F:U\rightsquigarrow Y.
\]
Put $C^0=1_U$, $P_n=E;C^n$ and $R_n=E;C^n;F$.

The matching in a completed return is the pullback
\[
\begin{tikzcd}[column sep=large]
R_n \arrow[r] \arrow[d] &F \arrow[d,"{s_F}"]\\
P_n \arrow[r,"{t_{P_n}}"'] &U .
\end{tikzcd}
\]
Its point is the full tuple $(e,c_1,\ldots,c_n,f)$; composition never replaces this tuple by mere existence of an intermediate state.
\end{definition}

\begin{definition}[Finite return and its temporal refinement]\label{def:finite-return-temporal-refinement}
The finite-return correspondence and its first-exit tower over $X$ are
\begin{align}
 \mathsf{Ret}_U(\alpha)&=D\boxplus\coprod_{n\geq0}R_n,
       \label{eq:return-new}\\
 \mathsf{FB}_0(\alpha)&=X,\nonumber\\
 \mathsf{FB}_{m+1}(\alpha)&=D\boxplus
            \coprod_{0\leq n<m}R_n\boxplus P_m.\label{eq:firstexit-new}
\end{align}
Restrictions fix previously completed summands, delete the exit from the newest completed summand, and delete the last internal turn from the running summand. The first restriction is the source map to $X$. Define $\mathsf{Done}_1=D$ and $\mathsf{Done}_{n+2}=R_n$.
\end{definition}

\begin{theorem}[Finite completion]\label{thm:finitecompletion}
The displayed restrictions define a temporal object. There is a canonical isomorphism of correspondences
\[
 \coprod_{m\geq1}\mathsf{Done}_m(\alpha)
       \ \cong\ \mathsf{Ret}_U(\alpha)
       =D\boxplus E;\left(\coprod_{n\geq0}C^n\right);F.
\]
Its elements are exactly finite routed executions which enter from $X$ and first return to $Y$, retaining every transition witness. If $\alpha$ is a partial map, the source leg of this correspondence is monic, so it is its finite-return partial map. In general it is a genuinely proof-relevant correspondence.
\end{theorem}
\begin{proof}\vmode
\leavevmode
\begin{enumerate}[label=(\roman*)]
\item \textbf{Well-defined restrictions.}
Each restriction has a defined target summand: deleting an exit in $R_{m-1}$ gives $P_{m-1}$, and deleting a last internal turn in $P_m$ gives $P_{m-1}$.
Earlier completed summands persist.
Composites are ordinary repeated prefix deletion, giving the tower.

\item \textbf{First-return decomposition.}
A routed return is either a single $D$ witness or uniquely a word $e,c_1,\ldots,c_n,f$.
Matching endpoints gives exactly the pullback defining $R_n$.

\item \textbf{Finite completion.}
Coproduct distribution over pullback proves the formula.

\item \textbf{Uniqueness for partial maps.}
For a partial map, any two executions from the same source have equal first steps whenever both exist.
Inductively they have equal successive steps.
Two purported first exits of different lengths would therefore require the same partial-map value to lie both in $Y$ and in $U$, impossible by disjointness.
Exits of the same length have equal witnesses.

\item \textbf{Monicity and general spans.}
These statements express that the diagonal of the source map is an isomorphism, proving monicity internally.
For an arbitrary span, several steps may have the same source, and none of these witnesses is identified.
\end{enumerate}
\end{proof}

\begin{proposition}[Effective finite return]\label{prop:effective-finite-return}
For an effective deterministic partial transition on the coded carriers $X+U$ and $Y+U$, the finite-return partial map of Theorem~\ref{thm:finitecompletion} has an effective open domain. Its classifier is the increasing, uniformly effective limit of bounded return approximants in $\mathbf{Kl}(\mathsf{Lift})$.
\end{proposition}
\begin{proof}\vmode
\leavevmode
\begin{enumerate}[label=(\roman*)]
\item \textbf{Kleisli approximants.}
Write $\widehat\alpha:X+U\to\mathsf{Lift}(Y+U)$ for the transition classifier.
Put $h_0=\bot$ and
\[
 h_{n+1}=[\eta_Y,h_n\iota_U]^\dagger\widehat\alpha,
 \qquad r_n=h_n\iota_X,
 \qquad b^\dagger=\mu_Y\mathsf{Lift}(b).
\]
Induction shows that $r_n$ returns exactly on first exits using at most $n$ transitions, with the same returned value.

\item \textbf{Finite-path interpretation.}
Its successive Kleisli factors follow exactly the routed state history.

\item \textbf{Effective increasing chain.}
Each approximant is effective by the monad constructions.
First-exit uniqueness makes $r_n\leq r_{n+1}$.

\item \textbf{Enumeration of the return domain.}
Dovetailing their approximation relations enumerates the relation of $r=\bigvee_n r_n$; every returned compact approximation appears at a finite stage.
Thus its domain is the union of the effective open bounded-return domains.

\item \textbf{Returned value and classifier.}
The unique returned value is precisely that of the correspondence in Theorem~\ref{thm:finitecompletion}.
The classifier gives
\[
\begin{tikzcd}[column sep=large]
\mathsf{Ret}_U(\alpha)\arrow[r,"t"]\arrow[d,hook,"s"']
 \arrow[dr,phantom,"\lrcorner",very near start]
 &Y\arrow[d,hook,"\eta_Y"]\\
X\arrow[r,"\bigvee_n r_n"']&\mathsf{Lift}(Y).
\end{tikzcd}
\]
Uniformity is supplied by the displayed recursion on $n$.

\item \textbf{Scope of uniform effectiveness.}
This construction does not assert that a nonuniform family of effective points has an effective supremum.
The general branching return correspondence keeps all its witnesses and is still interpreted by the span and coend constructions below.
\end{enumerate}
\end{proof}

The family $C^n$ itself is not a family of disjoint partial-map domains: the same initial point may support many lengths. It is the \emph{first-exit} condition that makes the completed domains disjoint in the partial-map case. Infinite internal execution and a finite blocked execution both fail to produce a return; the history and first-exit towers retain their different finite-prefix behaviour.

\subsection{Relational feedback for witnessed paths}\label{subsec:relational-feedback}
A comparison of transitions extends to paths when each transition from a related pair of states has a matching witness. Finite induction then gives a comparison of first returns, including their intermediate states.

\begin{theorem}[Witnessed relational first return]\label{thm:relational-first-return}
For $i=0,1$, let $\alpha_i:X_i+U_i\rightsquigarrow Y_i+U_i$ have blocks $D_i,E_i,C_i,F_i$. Let $R_X,R_U,R_Y$ be relations on the corresponding pairs of carriers. A paired block witness, for example $L_C:R_U\rightsquigarrow R_U$, has maps to $C_0,C_1$ commuting with both endpoint projections. Require sourcewise coverage: each transition of either $C_i$ from a state in a related pair has a paired witness in $L_C$ over that pair; equivalently
\[
 L_C\longrightarrow R_U\times_{U_i}C_i
\]
is surjective for $i=0,1$. Define $L_D,L_E,L_F$ with the corresponding endpoint relations and the same coverage property. Then
\[
 L_{\mathrm{ret}}=L_D\boxplus\coprod_{n\geq0}L_E;L_C^n;L_F
\]
compares the two first-return correspondences and is sourcewise surjective onto each. It retains every intermediate state, transition and paired proof witness.
\end{theorem}
\begin{proof}\vmode
\leavevmode
\begin{enumerate}[label=(\roman*)]
\item \textbf{Direct and initial lifts.}
A direct return is lifted by coverage of $L_D$.
For an execution $(e,c_1,\ldots,c_n,f)$ and a related source pair, lift $e$ using $L_E$; its target is a related internal pair.

\item \textbf{Successive internal and exit lifts.}
Lift the successive $c_j$ at that pair, then lift $f$ using $L_F$.
Finite induction gives a matching tuple with related output, in either direction.

\item \textbf{Relational composition.}
Pullback composition records the intermediate paired endpoints, and coproduct tags record the number of turns, proving precisely the asserted surjectivity and endpoint equations.

\item \textbf{Coherence and scope.}
Rebracketing preserves the same tuples.
This is a constructive rule for a supplied lockstep invariant; complete observational comparisons, which can also match different internal schedules, are provided by Theorem~\ref{thm:complete-certificates}.
\end{enumerate}
\end{proof}

\subsection{Execution and geometry of interaction in concurrent games}\label{subsec:execution-interaction-geometry}
We now calculate execution in the Scott and kernel calculus of Section~\ref{sec:game-calculus}. There are three related levels: matrices of completed execution witnesses, initialized games exposing a return, and the cubical histories inside the interaction. Keeping these levels explicit supplies a translation even when the protocol has no returning run. In particular, a possibly empty set of completed returns is not silently given a strategy root.

\subsubsection{The discrete Scott calculation of finite return}
\begin{definition}[Discrete interfaces and span kernels]\label{def:formal-discrete-interfaces-and-span-kernels}
For a set $X$, let $\underline X$ be the polarized precubical set having $X$ as its vertices and no positive-dimensional cells. Directly from the Kan-extension definition,
\[
 \Sc(\underline X)\cong\operatorname{Disc}(X).
\]
For a span $R:X\rightsquigarrow Y$, define the unpointed kernel
\[
 W_R(x,y)=\{r\in R\mid s(r)=x,\ t(r)=y\},
\]
as a set-valued functor on the two discrete Scott categories. This unpointed calculation retains all carrier witnesses; initialization is constructed in the next subsection.
\end{definition}

\begin{theorem}[Finite return as a path sum of coends]\label{thm:game-return-matrix}
The Scott kernel of a matched span pipeline is its raw coend:
\[
 W_{R;Q}(x,z)\cong\int^{y\in\operatorname{Disc}(Y)}
       W_R(x,y)\times W_Q(y,z).
\]
For the block correspondence $\alpha$ of Definition~\ref{def:finite-return-temporal-refinement},
\begin{equation}\label{eq:game-return}
 W_{\mathsf{Ret}_U(\alpha)}
 \cong W_D\ \amalg\ \coprod_{n\geq0}
         W_F\star W_C^{\star n}\star W_E.
\end{equation}
These are natural isomorphisms of set-valued kernels retaining every carrier witness. They commute with the routed trace comparisons of Proposition~\ref{prop:local-route-normalization}.
\end{theorem}
\begin{proof}\vmode
\leavevmode
\begin{enumerate}[label=(\roman*)]
\item \textbf{Discrete coend calculation.}
The indexing category is discrete, so its ordinary coend is the coproduct
$\coprod_y W_R(x,y)\times W_Q(y,z)$.
An element is exactly a pair of carrier witnesses with common endpoint $y$, which is the fibre of the pipeline pullback.
No nonidentity indexing arrow imposes a further relation.

\item \textbf{Finite-return summands.}
Coproducts of spans are fibrewise coproducts of these kernels.
Apply the calculation to every $E;C^n;F$ summand and then take their coproduct.

\item \textbf{Witness identity and coherence.}
On both sides an element is the identical finite routed witness list.
Consequently flattening, sliding a cut, or decomposing a two-colour hidden word gives the same comparison.
\end{enumerate}
\end{proof}

Equation~\eqref{eq:game-return} is a calculation inside the Scott/coend semantics of Section~\ref{sec:game-calculus}. The disjoint-union feedback on its left is represented on the right by a \emph{path sum}, not by taking a single raw compact trace of $W_\alpha$. Disjoint union of routing sectors and parallel tensor of games have different meanings. This formula supplies the actual comparison without identifying those tensors.

\subsubsection{Initialized return strategies}
Completed returns may form an empty family. To expose return behaviour as a game strategy, we therefore construct an initialized interface and distinguish its root from an actual returning execution.

\begin{definition}[One-request arena and response strategy]\label{def:responsegame}
For a set $X$, let $\mathsf A_X$ be the all-positive event arena with one event for each $x\in X$, all distinct events in conflict and no nontrivial causal order. Its configurations are $\varnothing$ and the singletons $x$. For a span $R:X\rightsquigarrow Y$, its response strategy has negative events $i_x$ for all $x\in X$ and positive events $o_r$ for all $r\in R$, with
\[
 i_{s(r)}<o_r.
\]
Distinct requests conflict, and all distinct responses conflict. The display sends $i_x$ to $x$ in $\mathsf A_X^\ast$ and $o_r$ to $t(r)$ in $\mathsf A_Y$.
\end{definition}

This strategy accepts a request even when it can never respond. Its positive response occurrences retain the entire witness $r$, not merely its endpoint label.

\begin{proposition}[Response kernel]\label{prop:responsekernel}
The strategy is receptive and courteous. Its pointed, reached Scott kernel is
\[
 \Phi_R(\varnothing,\varnothing)=*,\quad
 \Phi_R(x,\varnothing)=*,\quad
 \Phi_R(\varnothing,y)=\varnothing,\quad
 \Phi_R(x,y)=W_R(x,y).
\]
Unstraightening recovers exactly the event display of Definition~\ref{def:responsegame}.
\end{proposition}
\begin{proof}\vmode
\leavevmode
\begin{enumerate}[label=(\roman*)]
\item \textbf{Configuration classification.}
The strategy configurations are the empty configuration, $\{i_x\}$, and $\{i_{s(r)},o_r\}$.

\item \textbf{Receptivity.}
Every arena-enabled negative request has its unique lift at the empty configuration.
After a request no other negative event is arena-enabled.

\item \textbf{Courtesy and local injectivity.}
The only causal dependency introduced by the display is negative-to-positive, which courtesy permits.
Local injectivity holds on every configuration.

\item \textbf{Kernel fibres and restrictions.}
The fibre descriptions follow from the listed configurations.
The nontrivial kernel restriction forgets a response witness; accepting a request and then adding a response reaches every element from the root.

\item \textbf{Unstraightening comparison.}
The Scott discrete-fibration and unstraightening calculation of Theorem~\ref{thm:scott-discrete-fibrations} now gives the asserted display.
Alternatively its one-edge restrictions are exactly the two displayed transitions, with no independent pairs.
\end{enumerate}
\end{proof}

\begin{theorem}[Response interaction and completed protocol GoI]\label{thm:responsecomposition}
There is a canonical isomorphism of pointed reached kernels
\[
 \Phi_Q\star_{\mathsf A_Y}\Phi_R\ \cong\ \Phi_{R;Q}.
\]
In particular $\Phi_{\mathsf{Ret}_U(\alpha)}$ is an initialized game interpretation of finite return. Its response events are exactly the finite return witnesses: an entry, a composable finite list of selected internal rule records, and an exit. The pullbacks in each summand retain every one of these records.
\end{theorem}
\begin{proof}\vmode
\leavevmode
\begin{enumerate}[label=(\roman*)]
\item \textbf{Completed-response fibre.}
The Scott category of $\mathsf A_Y$ has precisely the arrows $\varnothing\to y$.
At external pair $(x,z)$ with a response present, the coend has vertices $W_R(x,y)\times W_Q(y,z)$; the would-be terms at $\varnothing$ are empty because $Q$ cannot respond without a request.
No nonidentity coend relation remains.
Its value is therefore the discrete coproduct over $y$, namely $W_{R;Q}(x,z)$.

\item \textbf{Pending-request fibre.}
At $(x,\varnothing)$ there is one representative from the empty intermediate configuration and one for each $r\in R$ with source $x$.
Each latter representative is identified with the empty one by its intermediate restriction arrow, so the quotient is a singleton.

\item \textbf{Root and restrictions.}
At $(\varnothing,\varnothing)$ it is also a singleton, and at $(\varnothing,z)$ it is empty.
These identifications respect all restrictions and the root and give exactly Proposition~\ref{prop:responsekernel}; every resulting element is reached.

\item \textbf{Application to return.}
Apply the construction to the return span and use Theorem~\ref{thm:game-return-matrix} for its response fibres.
\end{enumerate}
\end{proof}

This theorem retains distinct completed realizations even when their observable outputs agree. It deliberately does not turn an infinite hidden run into a response. The next construction records the execution before such completion or divergence is forgotten.

\begin{proposition}[Process elaboration and finite return]\label{prop:pc-return-transport}
An isomorphism of committed execution graphs preserving the declared input and exit boundaries induces a bijection on their finite-return witnesses and an isomorphism of their initialized response strategies. It preserves chronological length. If the comparison also retains residual cells, the corresponding concurrent families have the same images at the return boundary.
\end{proposition}
\begin{proof}\vmode
\leavevmode
\begin{enumerate}[label=(\roman*)]
\item \textbf{Transport of witnessed paths.}
Apply the graph isomorphism to each edge of a finite path.
Matching endpoints make concatenation well-defined, and the inverse acts edgewise.

\item \textbf{First-return boundaries.}
The declared boundaries identify the same first-return decomposition, so each path lies in the corresponding summand of \eqref{eq:game-return}.

\item \textbf{Response and duration.}
The response construction has one negative event for each request and one positive event for each returned witness; the path bijection identifies these events and their dependencies.
Each edge remains one edge, preserving length.

\item \textbf{Concurrent cells and scope.}
A residual cell uses the same family of witness occurrences, and the given cell comparison transports all its faces.
This is the concurrent statement.
A weak observational equivalence alone need not preserve either path length or witness multiplicity.
\end{enumerate}
\end{proof}

The return construction also accepts the named elaborations of Proposition~\ref{prop:pc-named-elaboration}. Its kernel reads each process path as the same list of medium and premise witnesses. The positive responses are therefore the finite rule-witness lists in~\eqref{eq:game-return}. Every request is accepted, including one whose response fibre is empty, so the comparison also accounts for nonreturning requests.

\subsubsection{Concurrent execution inside the return boundary}
Fix the retained execution presentation $U$ of an abstract communication protocol, rooted at its empty configuration. The following explicit boundary construction prevents an arbitrary state predicate from accidentally destroying a commuting square. Entry and exit are boundary markers; the original internal actions retain exactly their previous residual independence.

\begin{definition}[Routed history arena]\label{def:routedhistoryarena}
Fix an external entry value $x$. Each witness $e\in E$ over $x$ retains an entry history $h_e\in\Conf_f(U)$ and starts a branch with an entry event followed by the residual $U/h_e$ of Proposition~\ref{prop:generic-protocol-residual}. Its empty configuration denotes that retained entry history. An initialized presentation is the case $h_e=\varnothing$. Different entry branches conflict. The entry precedes every event in its branch. For every finite internal configuration $h$ in that branch and every exit witness $f$ in its declared exit fibre, adjoin an exit event $o_{h,f}$ whose past is precisely the entry together with $h$. It conflicts with every internal event outside $h$ and with every other exit event. Each direct witness $d\in D$ over $x$ gives a separate one-event completed branch, conflicting with the other branches. Denote the resulting polarized event structure by $\mathsf R_x(\alpha)$. Give the internal events their previously declared polarities and give boundary markers their declared external polarities.
\end{definition}

The exit fibre is indexed by the actual residual configuration. For a concrete state-indexed exit, pull the fibre back along the implementation's proved state evaluation. The construction uses no chosen linear extension. Its finite configurations are exactly the empty configuration, a direct completed branch, an entry followed by an internal configuration, and such a configuration with its selected exit.

\begin{theorem}[Concurrent return adequacy]\label{thm:concurrentreturn}
The routed history arena is an event structure. Its completed configurations are in bijection with finite return witnesses modulo adjacent swaps of independent internal residuals. Before this quotient, their linear extensions are the individual routed execution witnesses. Every finite internal configuration cube lifts with all its faces.

Its game interpretation is precisely
\[
 \Sc\bigl(K(\mathsf R_x(\alpha))\bigr)
   =\operatorname{Lan}_{Y}(J_{\mathrm{Sc}})
           \bigl(K(\mathsf R_x(\alpha))\bigr).
\]
The constant singleton kernel on this arena, rooted at the empty configuration, is reached and has cubical unstraightening $K(\mathsf R_x(\alpha))$. Its ordinary Scott category is the configuration order \eqref{eq:scottorder}.
\end{theorem}
\begin{proof}\vmode
\leavevmode
\begin{enumerate}[label=(\roman*)]
\item \textbf{Entry and internal consistency.}
Entry adds a common predecessor to a branch and incompatibility separates branches.
Retain the internal unfolding's finite consistency relation.
Equivalently, declare a finite event set consistent precisely when it is contained in one of the listed configurations.

\item \textbf{Event-structure axioms.}
Every event has a finite past, every singleton is consistent, and taking subsets preserves consistency.
Adding a causal predecessor also preserves consistency because each listed configuration is downward closed.

\item \textbf{Exit configurations.}
A consistent finite downset containing an exit must contain its entire specified past and no event outside it; without an exit it is an internal configuration with its entry, or the empty configuration.
Thus the event-structure axioms hold and its configurations are precisely those listed.

\item \textbf{History correspondence.}
Theorem~\ref{thm:generic-history-exchange} identifies residual configurations with their witnessed extension paths modulo adjacent independent swaps.
Here $C$ is the residual configuration-extension correspondence; an operational return comparison uses its explicit witness map.

\item \textbf{First-return witnesses.}
The entry is first and the exit last in every linear extension of a completed branch.
Specifying its entry, its internal linear extension, and its exit is exactly a point of $E;C^n;F$.
Direct branches give $D$.
This establishes the first two assertions.

\item \textbf{Preservation of independence.}
Internal independent families are unchanged before an exit is selected, so the same cubical families and face maps lift.
The proof uses the actual configuration family and does not impose a disjoint-owner condition.

\item \textbf{Categorical history kernel.}
The categorical formula is the very definition of the Scott interpretation, rather than a replacement by a new event-game construction.
The Grothendieck construction of a terminal presheaf is the identity on its base category.
Pulling it back along the unit in \eqref{eq:cubical-unstraightening} therefore gives the original cubical arena.

\item \textbf{Reachability.}
Every finite configuration can be built in a linear extension of its events.
A negative addition is a restriction in the execution direction, while a positive addition reverses a positive restriction.
Hence all singleton fibres are reached.
The ordinary category calculation is Theorem~\ref{thm:eventscott}.
\end{enumerate}
\end{proof}

The terminal kernel here describes the complete open behaviour on the routed execution arena; it is not being asserted to be a strategy on an unrelated smaller port arena. The response strategy instead exposes only the outer request and return, retaining internal histories as response witnesses. These are two explicitly constructed observations of the same execution. To impose a distributed controller, first use its compiled event strategy $S_\kappa$ and kernel from Theorem~\ref{thm:controllerkernel}, and then apply the boundary construction to its configurations. The precise erasure comparison is the following quotient construction. Let $H$ be the controlled physical histories, $\widehat H$ the selected compiled histories, and $\varepsilon:\widehat H\twoheadrightarrow H$ the map of Theorem~\ref{thm:compiler}. Declare entry and exit tests on $H$ and pull them back along $\varepsilon$. Thus a compiled boundary witness retains the same physical boundary witness. Erasure induces a surjection $e_R:\widehat R\twoheadrightarrow R$ on completed chronological return witnesses: it deletes commitments and retains all physical steps and both boundary witnesses. Every physical return lifts from each compiled entry representative by filling the finitely many missing veto commitments before each step. Hence
\begin{equation}\label{eq:compiled-return-quotient}
 \widehat H/(\widehat H\times_H\widehat H)\cong H,
 \qquad
 \widehat R/(\widehat R\times_R\widehat R)\cong R.
\end{equation}
Here a quotient means the coequalizer of the two kernel-pair projections. In $\Set$ a surjection has this coequalizer: a function constant on its fibres factors uniquely by assigning its common value to each image point. This proves both isomorphisms without choosing a section. Erasure also induces the corresponding surjection on completed history configurations after independent exchanges are identified. It preserves physical causal order; each enabled physical residual cube lifts after commitment completion by Theorem~\ref{thm:compiler}. These are weak execution and quotient comparisons, not a dimension-preserving isomorphism of the compiled and physical cubical carriers. A declared physical observer $o$ becomes $o\varepsilon$, with the same attained observations.

\subsubsection{Guarded state closure as a game-kernel calculation}
The return calculation treats completed paths. Guarded state closure instead organizes successive state approximations. We now identify its delay and comparison maps in the same game-kernel interpretation.

\begin{definition}[Discrete companion and state closure]\label{def:formal-discrete-companion-and-state-closure}
There is a second exact comparison, using the cartesian state interface rather than a return port. Perform the following construction internally in the temporal topos, or equivalently degree by degree with compatible restriction maps. Let $f:Z\times V\to O\times V$ be a temporal map. Its companion kernel on discrete Scott interfaces is
\[
 K_f((z,v),(o,w))=
 \begin{cases}*,&f(z,v)=(o,w),\\ \varnothing,&\text{otherwise}.
 \end{cases}
\]
The notation is a membership description of the graph subobject, so does not require a decision procedure for equality. Define its discrete-state \textbf{closure} by
\[
 \mathsf{cl}_V K_f(z,o)
   =\int^{v\in\operatorname{Disc}(V)}K_f((z,v),(o,v)).
\]
\end{definition}

\begin{theorem}[Guarded execution commutes with Scott closure]\label{thm:guardedgame}
If the hidden component of $f$ is generated guarded, there is a canonical isomorphism of temporal kernels
\[
 \mathsf{cl}_V K_f\ \cong\ K_{\mathsf{fb}_V(f)}.
\]
The carrier interpretation is the companion comparison of Theorem~\ref{thm:guarded-execution}. It is natural under parameter substitution and compatible with finite simultaneous guarded elimination. Before restricting to companions, the same discrete closure retains the entire set of hidden-state carrier witnesses.
\end{theorem}
\begin{proof}\vmode
\leavevmode
\begin{enumerate}[label=(\roman*)]
\item \textbf{Fixed-point fibre.}
The coend is over a discrete category.
Its fibre at $(z,o)$ consists exactly of values $v$ with
\[
 f_V(z,v)=v,\qquad f_O(z,v)=o.
\]
In internal terms it is the visible fibre of the equalizer object $\{(z,v)\mid f_V(z,v)=v\}$.

\item \textbf{Unique guarded solution.}
Theorem~\ref{thm:guarded-execution} identifies that object over $Z$ with $Z$ itself, by its unique fixed section.
Its visible value is $\mathsf{fb}_V(f)$, giving the isomorphism.

\item \textbf{Naturality and substitution.}
This is an equality of the graph fibres, so respects the temporal restrictions.
Parameter substitution preserves the fixed-section equation and its unique solution.

\item \textbf{Simultaneous elimination.}
For a jointly guarded finite system, Lemma~\ref{lem:guardedelimination} constructs the successive guarded equations and identifies their solution with the simultaneous solution; the same equality therefore holds after passage to kernels.

\item \textbf{Scope of uniqueness.}
A general carrier correspondence has no such unique-solution assertion, and the discrete coend simply retains all pairs of a carrier realization and its matching hidden value, exactly as the diagonal closure does.
\end{enumerate}
\end{proof}

\begin{definition}[Companion profunctor and closing-arrow coend]\label{def:formal-companion-profunctor-and-closing-arrow-coend}
For small categories $\mathcal Z,\mathcal U,\mathcal O$ and a functor $f:\mathcal Z\times\mathcal U\to\mathcal O\times\mathcal U$, write $f_O=\pi_Of$ and $f_U=\pi_Uf$. Its \textbf{companion profunctor} has type
\[
 K_f:(\mathcal Z\times\mathcal U)\times(\mathcal O\times\mathcal U)^\op\to\Set,
 \qquad K_f((z,u),(o,w))=\mathcal O(o,f_O(z,u))\times\mathcal U(w,f_U(z,u)).
\]
Its \textbf{raw hidden-state closure} is the functor $\mathcal Z\times\mathcal O^\op\to\Set$ with fibre
\[
 \int^{u\in\mathcal U}\mathcal O(o,f_O(z,u))\times\mathcal U(u,f_U(z,u)).
\]
A representative is $(u,\kappa,\lambda)$ with $\kappa:o\to f_O(z,u)$ and $\lambda:u\to f_U(z,u)$. For $a:z\to z'$ and $b:o'\to o$, its image is represented by
\[
 (u,\ f_O(a,1_u)\kappa b,\ f_U(a,1_u)\lambda).
\]
Functoriality of $f$ makes this assignment respect the coend relations. Thus the \textbf{closure} retains closing arrows and their coend identifications; it is not an equality predicate on hidden states. For discrete $\mathcal U$, only the identity closing arrows and identity coend relations remain, recovering Definition~\ref{def:formal-discrete-companion-and-state-closure}.
\end{definition}

Proposition~\ref{prop:local-delayed-paths} gives the delayed realization of these witnessed paths, including their full witness fibres.

\begin{example}[Hiding and stopping observations]\label{ex:pc-hiding-stopping}
Let an internal arena have two conflicting positive moves $b_0,b_1$. One ordinary strategy chooses either; its partner accepts either and offers an external positive $c$ only after $b_0$. The hidden ordinary composite has the possible occurrence $c$, with its retained visible principal history. The execution that selected $b_1$ instead stops without $c$, but its purely internal branch contributes no visible event after hiding. Thus that stopping alternative cannot be reconstructed from the hidden display alone. A stopping or must observation is attached to the execution presentation that retains it, rather than inferred from ordinary display equality.
\end{example}

\subsection{Coherence and delayed realization of witnessed execution}
The preceding constructions provide several presentations of witnessed execution. Their comparison maps must also agree under composition and delay. The following coherence calculations identify the retained data along those routes.

\begin{proposition}[Routed normalization]\label{prop:local-route-normalization}
Finite return of execution correspondences has natural, coherent trace comparisons for the disjoint-union tensor. They restrict to the constructed partial-map locus.
\end{proposition}
\begin{proof}\vmode
\leavevmode
\begin{enumerate}[label=(\roman*)]
\item \textbf{Route normal form.}
Expand a finite wiring expression into its returning routes.
A route retains its ordered component visits, stops at its first external exit, and carries the iterated pullback of the visited witness carriers.
Its normal form is the coproduct of these pullbacks over all finite returning routes.

\item \textbf{Distribution over coproducts.}
Distributivity of pullback over coproduct sends a nested realization to its flattened route and witness tuple; cutting at the expression's interfaces is the inverse.

\item \textbf{Naturality and superposition.}
Naturality concatenates external routes; superposing adjoins an independent sector; vanishing at the empty sector removes no visits.

\item \textbf{Vanishing and sliding.}
Vanishing at $U+V$ decomposes the same route into successive excursions through the first eliminated sector.
Sliding moves a cut while retaining its component visits.

\item \textbf{Yanking.}
Tracing the symmetry has the single straight route and identity carrier.
These prove the trace equations~\cite{jsv}.

\item \textbf{Coherence and partial maps.}
Every coherence composite fixes the same flattened route and tuple, hence is equal.
The monic-source calculation of Theorem~\ref{thm:finitecompletion} gives the restriction to partial maps.
\end{enumerate}
\end{proof}

\begin{construction}[The typed route comparisons]\vmode\label{con:typed-route-comparisons}
\leavevmode
\begin{enumerate}[label=(\roman*)]
\item \textbf{Route-indexed carriers.}
For a finite routed wiring expression $W:X\rightsquigarrow Y$, let
$\mathsf{Routes}(W)$ be the finite component-visit words that enter
at $X$, follow the wired sectors, and stop at their first $Y$ exit.
For a route $\gamma$ let $R_\gamma$ be the iterated pullback of its
component witness carriers over the consecutive visited states.

\item \textbf{Flattening map.}
Define
\[
 \mathsf N(W)=\coprod_{\gamma\in\mathsf{Routes}(W)}R_\gamma .
\]
The map $F_W$ from the iterated realization of $W$ to $\mathsf N(W)$
concatenates its nested visit words and witness tuples.

\item \textbf{Inverse cutting map.}
Its inverse
cuts a word whenever it crosses a subexpression's declared interface
and places the corresponding subtuple in that summand.
A cut is
determined by the retained sector tags, even when a state value
appears more than once.

\item \textbf{Inverse equations.}
Concatenation followed by cutting, and
cutting followed by concatenation, are identities.

\item \textbf{Trace comparisons and coherence.}
The trace comparisons, with unitors and symmetry rebracketings
understood as the explicit disjoint-sum bijections, are:
\begin{center}\small
\begin{tabular}{@{}p{.19\linewidth}p{.73\linewidth}@{}}
\toprule Law & Types and carrier bijection\\\midrule
Naturality &
$f:X+U\rightsquigarrow Y+U$, $a:X'\rightsquigarrow X$,
$b:Y\rightsquigarrow Y'$:
$\operatorname{tr}^{U}((b+1)f(a+1))
\cong b\,\operatorname{tr}^{U}(f)\,a$.
Attach the entry and exit witnesses to the same returning route.\\
Sliding &
$f:X+U\rightsquigarrow Y+V$, $g:V\rightsquigarrow U$:
$\operatorname{tr}^{U}((1+g)f)
\cong\operatorname{tr}^{V}(f(1+g))$.
Move the internal cut across each retained $g$ visit.\\
Empty vanishing &
$f:X+0\rightsquigarrow Y+0$:
$\operatorname{tr}^{0}(f)\cong f$ under unitors.
There is no internal visit.\\
Double vanishing &
$f:X+U+V\rightsquigarrow Y+U+V$:
$\operatorname{tr}^{U+V}(f)
\cong\operatorname{tr}^{U}(\operatorname{tr}^{V}(f))$.
Partition the word into its successive $V$ excursions between
$U$ visits and its final exit.\\
Superposing &
For $g:Z\rightsquigarrow Z'$, the trace of $g+f$ is
$g+\operatorname{tr}^{U}(f)$ after reassociation.
Its routes lie in one of the two disjoint sectors.\\
Yanking &
$\operatorname{tr}^{U}(\sigma_{U,U})\cong1_U$.
A returning route enters the hidden copy and immediately exits;
its matched witness is the same element of $U$.\\
\bottomrule
\end{tabular}
\end{center}
Each comparison is $F_{W'}^{-1}F_W$ after identification of the
displayed common route carrier.
Therefore two composites with the
same endpoints both equal that formula.
This proves all coherence
equations on carriers and their endpoint maps.

\item \textbf{Partial-map specialization.}
Monic source legs
give the already constructed restriction to partial maps.
\end{enumerate}
\end{construction}

\begin{proposition}[Delayed realization of witnessed paths]\label{prop:local-delayed-paths}
For the execution graph $C=(S\leftarrow D\to S)$ put $V_n=S^{n+1}$ and $O_n=\Lambda^n$, with last-entry deletion. There is a temporal correspondence
\[
 R_C:\Delta S\times\triangleright V\rightsquigarrow O\times V
\]
with carrier $\mathsf{Hist}(C)$. Precomposing its hidden input with $\mathsf{next}_V$ and then imposing equality of hidden input and output gives precisely this history carrier, with its initial state and emitted word. It commutes with witnessed graph maps, prefix deletion, and the entry/exit pullbacks of finite return.
\end{proposition}
\begin{proof}\vmode
\leavevmode
\begin{enumerate}[label=(\roman*)]
\item \textbf{State sequence and anchors.}
A path $h=(d_1,\ldots,d_n)$ has the unique state sequence $(s_0,\ldots,s_n)$.
For $n>0$, give it input $(s_0,(s_0,\ldots,s_{n-1}))$ and output $((\ell d_1,\ldots,\ell d_n),(s_0,\ldots,s_n))$; at degree zero use $(s_0,*)$ and $(\eps,(s_0))$.

\item \textbf{Prefix compatibility.}
These anchors commute with deletion.

\item \textbf{Unique state closure.}
After precomposition a witness is $(h,v)$ whose delayed prefix agrees with that of $h$.
Closure forces $v$ to equal the full state sequence of $h$.

\item \textbf{Path correspondence and completion.}
Projection to $h$ and adjoining this unique sequence are inverse, with
\[
\begin{tikzcd}[column sep=large]
\mathsf{Cl}_{n+1}\arrow[r,"\cong"]\arrow[d,"d_n"']&H_{n+1}(C)\arrow[d,"d_n"]\\
\mathsf{Cl}_{n}\arrow[r,"\cong"']&H_n(C).
\end{tikzcd}
\]
An entry and exit attach the same witnesses to both carriers, so their pullbacks preserve this comparison.
Coproduct over return lengths completes it.

\item \textbf{Audit evaluation and nondeterminism.}
For administrative receipts evaluate the cumulative terminal history on this same tuple; parent pointers and evidence are unchanged.
Thus delay retains nondeterministic path witnesses and does not select a branch or turn nonreturn into a reply.
\end{enumerate}
\end{proof}

\subsection{Calculus: feedback, guarded recursion, and finite return}\label{subsec:calculus-execution}
We now interpret the recursive and routed process constructors using the execution operations just constructed. The interpretation keeps the connection, finite-return and temporal judgements separate, and each rule retains the witnesses appropriate to its type.

\subsubsection{Connection, finite return, and temporal unfolding}\label{subsec:pc-feedback}

Compact feedback connects dual ports using secured composition. Finite return routes a process through disjoint external and internal state sectors. A returned value retains one complete finite path and its transition witnesses. Temporal unfolding records prefixes at successive chronological lengths and constructs the later operator from that indexing. Their roles are summarized by
\begin{center}
\begin{tabular}{@{}p{.22\linewidth}p{.31\linewidth}p{.39\linewidth}@{}}
\toprule Construction & Process operation & Retained semantic object\\\midrule
Compact feedback & Connect dual ports & Ordinary strategy built by secured cut\\
Finite return & Route until a declared exit & A disjoint family of finite witnessed paths\\
Guarded unfolding & Advance a temporal stage & A compatible history tower with its generated delay\\
\bottomrule
\end{tabular}
\end{center}
\begin{definition}[Routed execution syntax]\label{def:formal-routed-execution-syntax}
For a process $P$ with a declared entry, exit, and internal sector $U$, write $\mathsf{run}_U(P)$ for the routed execution term: a completed response consists of a finite path of its committed transitions ending at the first declared exit, with its full witness list. The preceding semantic constructions and the following rules specify each interpretation. A response strategy accepts its request even when no finite returning continuation exists.
\end{definition}

\begin{definition}[Typed compact feedback rule]\label{def:calculus-feedback-rule}
For $A,B,U:\mathsf{Int}$, feedback is the derived interface constructor
\[
 \frac{F:A\otimes U\to B\otimes U}{\mathsf{feedback}_U(F):A\to B}.
\]
Its elimination is its expansion into the cup, $F$ and the cap of Definition~\ref{def:formal-compact-feedback}; thus it is an abbreviation in the existing interface syntax. Its computation equations include unit feedback and yanking:
\[
 \mathsf{feedback}_I(F)\equiv F,\qquad
 \mathsf{feedback}_U(\sigma_{U,U})\equiv\mathsf{fwd}_U,
\]
with unitors in the first equation. Expanding the definition proves the first by unit coherence and the second by one snake equation. Feedback therefore has the claimed type and computations by Proposition~\ref{prop:process-typing-structural-conversions}. The finite-return type below retains its separate operational path witnesses.
\end{definition}

\subsubsection{Recursive processes and replicated services}\label{subsec:pc-iteration}
Recursive unfolding extends the finite process syntax through successive stages. The service construction then uses that unfolding to repeat a state update after each completed command, retaining the command and reply occurrences.

\begin{definition}[Recursive and replicated execution]\label{def:formal-recursive-and-replicated-execution}
The recursive constructors are those of Definition~\ref{def:formal-sorted-value-contexts-and-process-constructors}. Recursive unfolding uses capture-avoiding substitution $P[\mu X.P/X]$. Input-guarded service recursion has a distinguished operational presentation: a persistent offer $\mathsf{srv}_a(x.P)$ accepts one concrete input, starts $P[v/x]$, and retains the server,
\[
 \mathsf{srv}_a(x.P)\xrightarrow{a?v}
 P[v/x]\mid\mathsf{srv}_a(x.P).
\]
Each accepted request gives its new service instance an occurrence address. This is the input-driven unfolding of $\mu X.a?(x).(P\mid X)$ with $x$ fresh for $X$. At the ordinary transition-relation level it has the replicated-input behavior of $!a?(x).P$ modulo structural congruence.

The replicated-input and asynchronous-output rules give the operational communication form used by $\pi_F$. Its categorical model in \cite[Definition~6]{pi-freyd} additionally includes a cartesian value category and an identity-on-objects monoidal functor with a right adjoint. The common categorical comparison here is compact wiring. Compact \textbf{closure} of $\mathsf{CG}$ alone is not a construction of that extra Freyd structure.

General replication uses the declared congruence $!P\equiv P\mid!P$. Its exact witnesses are defined by the finite-support occurrence construction of Proposition~\ref{prop:pc-finite-support}: a replication occurrence has virtual copies indexed by $\mathbb N$, and a commitment retains the selected body occurrences and medium record. Unused unfoldings contribute no events or additional witnesses. Copies are whole process bodies, so independent actions within one copy remain independent; there is no copy-activation token consumed by its first action. Forgetting addresses recovers the structural-congruence transition relation. General replication is not derived from the input-guarded service equation or from a compact cup. The established game-semantic treatment of copy indices and reindexing symmetries is \cite[Section~2.5.2]{ccw-thin}; the operational witness comparison used here is proved below.
\end{definition}

The temporal fixed-point construction applies when dependence on the recursive variable factors through the generated next map. This factorization connects operational guardedness to the temporal interpretation.

Equality has a different scope. Proposition~\ref{prop:pc-normalization} decides finite occurrence equality. It cannot decide equality of all recursive processes: a finite machine medium that replies to a request exactly upon halting would turn such a procedure into a halting decision.

\begin{definition}[Temporal formation and guarded computation]\label{def:temporal-formation-guarded-computation}
A \textbf{temporal type} $A:\mathsf{Time}$ is a tower $A\in\mathcal T(\E)=\E^{\omega^\op}$. A temporal context is a finite product of towers, and its judgment $\Gamma\vdash t:A$ denotes a tower map $\Gamma\to A$. This interpretation is distinct from an interface judgment in $\mathsf{CG}$.

In $\mathcal T(\mathcal E)$ the later type is $ (\triangleright A)_0=1$, $(\triangleright A)_{n+1}=A_n$, with its inherited restrictions. The next map has components $!$ and $d_n^A$. Formation and introduction are
\[
 \frac{A:\Time}{\triangleright A:\Time}\qquad
 \frac{\Gamma\vdash t:A}{\Gamma\vdash\mathsf{next}(t):\triangleright A}.
\]
The \textbf{guarded recursion rule} and its computation are
\[
 \frac{\Gamma,x:\triangleright A\vdash t:A}
 {\Gamma\vdash\mathsf{fix}(x.t):A},\qquad
 \mathsf{fix}(x.t)=t[\mathsf{next}(\mathsf{fix}(x.t))/x].
\]
This is a temporal judgment with $\llbracket t\rrbracket:\Gamma\times\triangleright A\to A$. There is no elimination rule $\triangleright A\to A$ at an arbitrary current stage. Observation one stage later uses $(\triangleright A)_{n+1}=A_n$.
\end{definition}

\begin{theorem}[Construction and uniqueness of the guarded rule]\label{thm:guarded-rule-construction-uniqueness}
The guarded rule is defined, satisfies its computation equation, is unique among solutions, and commutes with parameter substitution.
\end{theorem}
\begin{proof}\vmode
\leavevmode
\begin{enumerate}[label=(\roman*)]
\item \textbf{Guarded solution.}
Apply the degreewise construction of Theorem~\ref{thm:guarded-execution} to $\llbracket t\rrbracket:\Gamma\times\triangleright A\to A$.
Its unique solution is the interpretation of $\mathsf{fix}$, and its equation is exactly the substitution displayed in the rule.

\item \textbf{Substitution and interpretation.}
The same uniqueness identifies its pullback along every parameter map, proving substitution.
Thus the rule uses the constructed fixed point rather than another recursion principle.
\end{enumerate}
\end{proof}

\paragraph{Server computation.}\label{par:server-computation}
The recursive and replicated rules are those of Definition~\ref{def:formal-recursive-and-replicated-execution}. In particular one unfolding of $\mathsf{srv}_a(x.P)=\mu X.a?(x).(P\mid X)$ exposes its input; commitment installs $P[v/x]$ and retains the server. The separate temporal rule applies to bodies factoring through the generated delay, not to unrestricted recursive terms.

\begin{definition}[Return-type formation and path elimination]\label{def:return-types-path-elimination}
For the constructed routed correspondence $\alpha:X+U\rightsquigarrow Y+U$, write its four blocks as $D:X\rightsquigarrow Y$, $E:X\rightsquigarrow U$, $C:U\rightsquigarrow U$ and $F:U\rightsquigarrow Y$. Their endpoint fibres are sets. The return type at $(x,y)$ is
\[
 \mathsf{Ret}_U(\alpha)(x,y)=D(x,y)\; +
 \coprod_{n\geq0}\coprod_{u_0,\ldots,u_n\in U}
 E(x,u_0)\times\Bigl(\prod_{j=1}^n C(u_{j-1},u_j)\Bigr)\times F(u_n,y).
\]
Its introductions are $\mathsf{direct}(d)$ and
$\mathsf{path}(e,c_1,\ldots,c_n,f)$ with exactly these types and matching endpoints. Its elimination into any target family supplies one function on the direct summand and one on every displayed path summand. The computation rules apply the corresponding function to the retained tuple; eta reconstructs that tuple. Equivalently these are the coproduct and pullback universal properties already used by the finite-return theorem.

The process form $\mathsf{run}_U(P)$ is interpreted by the initialized response strategy for this return family. A request at $x$ has a possible positive reply $(y,r)$ for each $r\in\mathsf{Ret}_U(\alpha_P)(x,y)$. A request has no reply when all these fibres are empty; its negative acceptance remains present. Distinct paths remain distinct reply witnesses. Guarded temporal execution and this finite-return type have their separately constructed interpretations.
\end{definition}
\begin{proof}\vmode
\leavevmode
\begin{enumerate}[label=(\roman*)]
\item \textbf{First-return decomposition.}
Extensivity partitions a completed routed path into its direct case or its first entry, a unique finite number of internal edges, and its first exit.
Pullback matching gives the displayed tuple and its intermediate states.

\item \textbf{Inverse path assembly.}
Conversely such a tuple is a routed path with those endpoints and no earlier exit.
These two constructions are inverse.

\item \textbf{Elimination principle.}
Functions out of the coproduct are exactly the stated elimination branches, proving beta/eta.

\item \textbf{Response interpretation.}
The response-game construction attaches its response to the particular $r$, so gives the stated process interpretation even when no return exists.
\end{enumerate}
\end{proof}

\begin{definition}[Response interpretation of routed processes]\label{def:routed-process-response}
Writing $\alpha_P$ for the routed committed correspondence of $P$, the response interpretation of the term introduced in Section~\ref{subsec:pc-feedback} is
\[
 \llbracket\mathsf{run}_U(P)\rrbracket_{\mathrm{resp}}
   =\Phi_{\mathsf{Ret}_U(\alpha_P)}.
\]
\end{definition}


\subsubsection{Computational interpretation of communication}\label{subsec:effective-communication}
The comparison $U$ of Theorem~\ref{thm:effective-game-refinement} forgets the computational presentation of a strategy, whereas $\mathsf D$ of Proposition~\ref{prop:effective-configuration-comparison} acts on its event display and configuration evidence. The strong monad of Proposition~\ref{prop:effective-lift-monad} acts on $\mathsf{EDom}$. These constructions give the following interpretation of the communication operations.

Finite codes for commands, values, banks and audit nodes have discrete presentations in $\mathsf{EDom}$. Finite evaluators act by effective maps; potentially nonterminating computations act by Kleisli arrows into $\mathsf{Lift}$. Proposition~\ref{prop:effective-name-action} transports their enumerated finite configuration evidence.

Chronological execution and delayed histories retain the successive states of these computations. Extending a history extends its finite state sequence. We do not model a changing current state as a monotone map from growing histories into a discrete state set.

For the chronological histories of a fixed initialized machine, write $H_f$ for its finite witnessed paths and $V_f$ for their finite state sequences, both ordered by prefix. The execution map $e:H_f\to V_f$ computes the whole sequence and preserves prefixes. Lower closure of its finite images therefore extends it to an effective map on ideals. The evaluation and history comparisons are
\[
\begin{tikzcd}[column sep=large]
U\arrow[r,"f"]\arrow[d,hook]
 &Y\arrow[d,hook,"\eta_Y"]\\
X\arrow[r,"\widehat f"']&\mathsf{Lift}(Y)
\end{tikzcd}
\qquad
\begin{tikzcd}[column sep=large]
H_f\arrow[r,"e"]\arrow[d,hook,"\downarrow"']
 &V_f\arrow[d,hook,"\downarrow"]\\
\operatorname{Idl}(H_f)\arrow[r,"\operatorname{Idl}(e)"']
 &\operatorname{Idl}(V_f).
\end{tikzcd}
\]
Section~\ref{sec:execution-games} identifies deterministic finite return with this classifier; Section~\ref{subsec:crypto-categories} identifies strict public programs with its Kleisli composition and restriction structure; Section~\ref{subsec:crypto-games} identifies service comparison maps with the strict transducer simulations constructed there. General concurrent choices retain their span witnesses. Neither the computational refinement nor the classifier identifies different operational witnesses or makes arbitrary observation predicates decidable.


The game interpretation now accounts for both a protocol's exposed interaction and the transitions retained in its execution. A cryptographic service uses the same ports, composition, and observations. Its simulator must produce its replies from the information those ports supply.

\section{Public information and composable cryptographic security}\label{sec:cryptographic-communication-games}\label{sec:public-information-secure-refinement}
A simulator is a strategy that translates interaction at the ideal adversarial ports into interaction at the real adversarial ports. Correctness and security give opposite factorizations over the honest interface. To construct the simulator, we identify the coalition's public observations, determine which operations on private data descend to them, and interpret the resulting controller as a communication strategy.

\subsection{Resources, protocols, and simulation}\label{subsec:resource-simulation}
\begin{definition}[Simulator judgement]\label{def:simulator-realization}
For $r:I\to H\otimes A$, $s:I\to H\otimes B$ and an actual strategy $\sigma:B\to A$, write
\[
 \sigma\Vdash_H r\preccurlyeq s
 \quad\Longleftrightarrow\quad r\eqsc(1_H\otimes\sigma)s.
\]
Here $H$ is the common honest interface, $A$ is the real adversarial interface, and $B$ is its ideal counterpart. The test compares the two resulting states at the same boundary $H\otimes A$. More generally, an honest adapter $f:H\to K$ and a simulator $q:B\to A$ compare $(f\otimes1_A)r$ and $(1_K\otimes q)s$ at $K\otimes A$ when $s:I\to K\otimes B$.
\end{definition}

\begin{theorem}[Exact observational criterion for simulation]\label{thm:observational-simulation}
For a fixed strategy $\sigma$, the following statements are equivalent:
\leavevmode
\begin{enumerate}[label=(\roman*)]
\item $\sigma\Vdash_H r\preccurlyeq s$;
\item for every strategy $t:H\otimes A\to O$ and every Scott open $U\subseteq\D(O)$,
\[
 (t\circ r)\models U
 \quad\Longleftrightarrow\quad
 \bigl(t\circ(1_H\otimes\sigma)\circ s\bigr)\models U;
\]
\item $\Trace(r)=\Trace((1_H\otimes\sigma)s)$;
\item the real and simulated displays admit a trace certificate;
\item for every $U\in\Op(T_{H\otimes A})$, the real and simulated states have the same $\mathsf{May}(U)$ at their common serialized boundary.
\end{enumerate}
For the test condition it suffices to take the characteristic tests $\tau_w$ of Lemma~\ref{lem:characteristic-test} with their success opens. These judgements compose with witness $\sigma\rho$ and tensor with witness $\sigma\otimes\rho$, including in the presence of a joint adversary.
\end{theorem}
\begin{proof}\vmode
\leavevmode
\begin{enumerate}[label=(\roman*)]
\item \textbf{Equivalent realization criteria.}
Contextual equivalence implies agreement for every displayed test by postcomposition. Applying this agreement to each $\tau_w$ characterizes equality of trace images by Lemma~\ref{lem:characteristic-test}. Conversely, Theorem~\ref{thm:trace-full-abstraction} carries trace equality back to contextual equivalence, and Theorem~\ref{thm:complete-certificates} identifies it with the existence of a complete certificate. Theorem~\ref{thm:open-trace-testing} identifies the same trace comparison with (v).

\item \textbf{Sequential composition.}
If $r\eqsc(1\otimes\sigma)s$ and $s\eqsc(1\otimes\rho)t$, contextual congruence gives $r\eqsc(1\otimes\sigma\rho)t$.

\item \textbf{Parallel composition.}
For tensor, first replace one component in the context containing the other, then replace the other.
Thus the proof allows a subsequent observer interacting with both components together.

\item \textbf{Composition of witnesses.}
At the witness level transport certificates by these contexts and compose them by pullback.
\end{enumerate}
\end{proof}


A trace certificate verifies the action of a supplied online simulator. The simulator itself must be a strategy using the ideal ports, and this requirement fixes the quantifier order for private inputs: realization of a family indexed by $x$ means $\exists\sigma\,\forall x$, with $\sigma$ chosen before all tests and input fillings.

Pointwise choices $\forall x\,\exists\sigma_x$ do not suffice. If the ideal resource discloses nothing and the real resource emits the bit $x$, a constant-output simulator can match each fixed input separately by depending on $x$. No single simulator without that input produces precisely the two different constant behaviours. Keeping the input as an honest port enforces uniformity in the resource equation.

\begin{corollary}[Observing a simulated output]\label{cor:honest-readout}
If $\sigma\Vdash_H r\preccurlyeq s$ and an ordinary honest adapter $d:H\to H'$ observes the exposed output, the same $\sigma$ realizes the two adapted states over $H'$. In particular the adapter may query basis observations of an issued payload or apply a continuous observation at an exported query port.
\end{corollary}
\begin{proof}
Contextual congruence followed by tensor interchange gives
\[
 (d\otimes1_A)r\ \eqsc\
 (d\otimes1_A)(1_H\otimes\sigma)s
   = (1_{H'}\otimes\sigma)(d\otimes1_B)s .
\]
The query adapter of Proposition~\ref{prop:continuous-query-strategy} is an ordinary strategy of this kind at its stated type. Theorem~\ref{thm:open-trace-testing} applies at the resulting boundary.
\end{proof}

\begin{corollary}[Transfer of relational privacy]\label{cor:relational-privacy-transfer}
If the same $\sigma$ realizes $r_x$ from $s_x$ for every $x$, and an input relation $E$ satisfies $xEy\Rightarrow s_x\eqsc s_y$, then $xEy\Rightarrow r_x\eqsc r_y$.
\end{corollary}
\begin{proof}\vmode
\leavevmode
\begin{enumerate}[label=(\roman*)]
\item \textbf{Transport of privacy through the simulator.}
For $xEy$, contextual congruence gives
$r_x\eqsc(1\otimes\sigma)s_x\eqsc(1\otimes\sigma)s_y\eqsc r_y$.
The same strategy and the same exposed interfaces occur throughout.
\end{enumerate}
\end{proof}

\paragraph{Owned resources and conversions.}
Fix the finite participant set $P$ and the chosen coalitions $\mathcal K\subseteq\Pow(P)$. Work in $\mathcal C=\CG/{\eqsc}$. The same equations in $\CG$ give exact display comparisons. The constructions restrict to the declared cryptographic interface category $\Csc$ below. A resource is a strategy with owned interfaces, and an implementation is a tuple of local strategies. Ownership is retained when a coalition's ports are grouped.

\begin{definition}[Resources and correct party conversions]\label{def:formal-multipartite-resource-categories}
For $\mathbf A=(A_i)_{i\in P}$, put $A_P=\bigotimes_iA_i$. A resource is $r:I\to A_P$. A correct conversion $\mathbf f:r\to s$, where $s:I\to B_P$, is a tuple $f_i:A_i\to B_i$ satisfying $(\bigotimes_i f_i)r=s$. For $J\in\mathcal K$, group its interfaces as $A=A_J$, $B=B_J$ and put $H=B_{P\setminus J}$. Write
\[
 \widetilde r=(f_{P\setminus J}\otimes1_A)r:I\to H\otimes A,
 \qquad c=f_J:A\to B.
\]
Correctness is then $s=(1_H\otimes c)\widetilde r$. A coalition attacker is any strategy $a:A\to E$, including one acting jointly on all the coalition ports.
\end{definition}

\begin{definition}[Simulation squares]\label{def:crypto-simulation-category}
For $r:I\to H\otimes A$ and $s:I\to K\otimes B$, a comparison consists of an honest adapter $f:H\to K$ and a simulator $q:B\to A$ with $(f\otimes1_A)r=(1_K\otimes q)s$. Compact transposition assigns
\[
 \widehat r:A^\ast\to H,\qquad \widehat s:B^\ast\to K,
 \qquad f\widehat r=\widehat s q^\ast.
\]
These commuting squares, with their actual side strategies, are the arrows of $\mathsf{CompSim}(\mathcal C)$. A secure party conversion also retains its entire party tuple and correctness equation, and one such square for each tested coalition.
\end{definition}

\begin{proposition}[Simulation squares and honest-interface indexing]\label{prop:simulation-squares-codomain-opfibration}
Compact transposition gives a symmetric monoidal equivalence
\[
 \mathsf{CompSim}(\mathcal C)\simeq\operatorname{Arr}(\mathcal C).
\]
The projection to the honest interface is the split opfibration $\operatorname{cod}$; its fibre at $H$ is $\mathcal C/H$. The chosen pushforward along $f:H\to K$ is postcomposition $f_!(u)=fu$. In particular
\[
 \operatorname{Sim}_f(r,s)\cong
 \operatorname{Hom}_{\mathcal C/K}(f_!\widehat r,\widehat s).
\]
\end{proposition}
\begin{proof}\vmode
\leavevmode
\begin{enumerate}[label=(\roman*)]
\item \textbf{Transposition of the comparison.}
Transposition carries the comparison equation to the displayed square, and naming gives its inverse.

\item \textbf{Composition of simulation squares.}
Duality is a contravariant equivalence: $(qq')^\ast=(q')^\ast q^\ast$.
Hence composition becomes square composition, while tensor and its grouping comparisons become pointwise tensor.

\item \textbf{Chosen opcartesian lift.}
For $u:X\to H$, the chosen arrow over $f$ is $(1_X,f):u\to fu$.

\item \textbf{Universal factorization.}
Any square $(v,gf):u\to w$, where $w:Y\to L$ and $g:K\to L$, factors uniquely as this square followed by $(v,g):fu\to w$.
Its equation is the same $gfu=wv$.
This is the opcartesian universal property.

\item \textbf{Split indexing and fibres.}
Postcomposition preserves identities and composition strictly, proving splitness.
Squares over $1_H$ are exactly slice arrows.
Every comparison over $f$ factors through the chosen opcartesian arrow and a slice arrow, giving the final bijection.
\end{enumerate}
\end{proof}

\begin{theorem}[Correctness and security as mutual factorization]\label{thm:correctness-security-mutual-factorization}
For a correct conversion and a fixed coalition as above, the following are equivalent:
\leavevmode
\begin{enumerate}[label=(\roman*)]
\item For every $a:A\to E$ there is $b:B\to E$ with
\begin{equation}\label{eq:all-attacks}
 (1_H\otimes a)\widetilde r=(1_H\otimes b)s.
\end{equation}
\item There is a simulator $\sigma:B\to A$ with
\begin{equation}\label{eq:simulator-factorization}
 \widetilde r=(1_H\otimes\sigma)s.
\end{equation}
\item In $\mathcal C/H$ there is an arrow $\widehat{\widetilde r}\to\widehat s$, with upper map $\sigma^\ast$.
\end{enumerate}
Correctness supplies the reverse slice arrow, with upper map $c^\ast$:
\begin{equation}\label{eq:mutual-security-factorization}
 \widehat s=\widehat{\widetilde r}\,c^\ast,
 \qquad
 \widehat{\widetilde r}=\widehat s\,\sigma^\ast.
\end{equation}
Thus correct secure refinement is mutual factorization in this slice. The two arrows need not be inverse. All equations in $\Csc$ compare every contextual Scott observation with the same simulator.
\end{theorem}
\begin{proof}\vmode
\leavevmode
\begin{enumerate}[label=(\roman*)]
\item \textbf{Extraction of a uniform simulator.}
Evaluate (i) at $E=A$ and $a=1_A$ to obtain $\sigma$.

\item \textbf{Simulation of every attacker.}
Conversely, postcompose the equation in (ii) with $1_H\otimes a$ and set $b=a\sigma$.
This chooses one simulator before any later test; its quantifiers are $\forall a\,\exists b\,\forall\text{ tests}$.

\item \textbf{Transposed security equation.}
Transposition proves (ii)$\Leftrightarrow$(iii).

\item \textbf{Transposed correctness equation.}
Transposing the already specified correctness equation gives the other factorization.
Neither equation cancels the resource strategy, so neither asserts $c\sigma=1$ or $\sigma c=1$.

\item \textbf{Ambient contexts and joint attackers.}
The equations also hold in the ambient game quotient under its inclusion.
Consequently postcomposition covers an attacker whose final arena lies outside the chosen interface fragment as well.
There is no restriction to attackers expressible by a particular finite implementation language.
\end{enumerate}
\end{proof}

\begin{theorem}[Composition of secure party conversions]\label{thm:secure-conversion-composition}
Correct party conversions secure against all $J\in\mathcal K$ compose and tensor. Keeping their simulator witnesses gives a symmetric monoidal category. For sequential conversions the coalition witness is $\sigma_f\sigma_g$; for parallel conversions it is $\sigma_f\otimes\sigma_g$, with the specified grouping maps.
\end{theorem}
\begin{proof}\vmode
\leavevmode
\begin{enumerate}[label=(\roman*)]
\item \textbf{Composition of correctness.}
Correctness equations compose componentwise.

\item \textbf{Pasting simulation squares.}
Simulation squares paste: if $f\widehat r=\widehat s q^\ast$ and $g\widehat s=\widehat t p^\ast$, then
\[
 gf\widehat r=\widehat t p^\ast q^\ast
             =\widehat t(qp)^\ast.
\]
Tensoring the equations proves the parallel case.

\item \textbf{Joint attackers.}
A subsequent joint attacker $a$ is simulated by $a(\sigma_f\otimes\sigma_g)$, so no independence of parallel attacks is imposed.

\item \textbf{Identities and coherence.}
Identities have identity witnesses and structural isomorphisms have inverse coalition witnesses.
The game equations give associativity and interchange.

\item \textbf{Coalitionwise composition.}
Perform this construction for each coalition; forgetting the witnesses preserves the same operations.
\end{enumerate}
\end{proof}

\begin{remark}[The scope of the codomain construction]\label{rem:game-codomain-pullback-obstruction}
The codomain projection is not a fibration on the whole game category. Let $P_+$ have one positive event, and let $d_0,d_1:I\to P_+$ be respectively the silent and emitting strategies. A cone $u,v:Z\to I$ over $d_0,d_1$ would satisfy $d_0u=d_1v$. Every arena $Z$ has a state $t:I\to Z$: restrict to negative events whose entire past is negative. This inclusion is receptive, since the predecessors of an enabled negative move are already in that restriction, and courteous by inherited order. There is only one scalar $I\to I$, because its total display has empty target. Thus $ut=vt=1_I$, forcing $d_0=d_1$, contrary to the open observing the positive event. This argument also holds in $\Csc$. Therefore this cospan has no cone and no pullback.

The domain projection $\operatorname{Arr}(\mathcal C)\to\mathcal C$ is always a split fibration, with precomposition lifts and fibres $X/\mathcal C$. Compact duality gives $\mathcal C/H\simeq(H^\ast/\mathcal C)^\op$. Event-structure pullbacks used in secured interaction, configuration pullbacks, and the strict-service pullbacks below belong to their stated categories; none is a pullback of arbitrary game strategies.
\end{remark}

\paragraph{Attainable views.}\label{subsec:security-images}
For $\widetilde r:I\to H\otimes A$ and $s:I\to H\otimes B$, set
\[
 \Att_A(E)=\mathcal C(A,E),\quad
 \Att_B(E)=\mathcal C(B,E),\quad
 \Beh_H(E)=\mathcal C(I,H\otimes E).
\]
Here $\Beh_H$ is the exposed-state functor, and $\Views_r$ is an attainable-view image. Postcomposition makes these covariant functors. The natural maps $\Phi_{r,E}(a)=(1_H\otimes a)\widetilde r$ and $\Phi_{s,E}(b)=(1_H\otimes b)s$ have pointwise images $\Views_r,\Views_s\subseteq\Beh_H$. Compact transposition at $H$ identifies $\Beh_H$ with $\mathcal C(H^\ast,-)$ and these images with the principal cosieves of the transposed resources: all their postcomposites by game strategies.

\begin{corollary}[Images, correctness, and the uniform witness]\label{cor:security-images}
Without using correctness, $\Views_r\subseteq\Views_s$ is equivalent to a simulator equation, to factoring $\Phi_r$ through the image inclusion for $s$, and to a natural $T:\Att_A\to\Att_B$ with $\Phi_r=\Phi_sT$. For a correct conversion, security is equivalent to
\[
 \Views_r=\Views_s.
\]
The image inclusion is unique; the simulator need not be unique.
\end{corollary}
\begin{proof}\vmode
\leavevmode
\begin{enumerate}[label=(\roman*)]
\item \textbf{Simulator from image membership.}
Pointwise image membership at $E=A$, $1_A$ gives $\sigma:B\to A$.
Its postcomposites give $T_E(a)=a\sigma$ and the image inclusion.

\item \textbf{Simulator from naturality.}
Conversely naturality forces $T_E(a)=aT_A(1_A)$, recovering that same simulator equation.

\item \textbf{Image subfunctors.}
Pointwise images are subfunctors by postcomposition, and factoring through their inclusions is exactly the stated membership condition.

\item \textbf{Reverse inclusion from correctness.}
Correctness supplies $s=(1_H\otimes c)\widetilde r$, so every ideal attacker $b$ has real counterpart $bc$; this proves the reverse inclusion.
In the functor category $[\mathcal C,\Set]$, $\Phi_r$ factors as a pointwise surjection onto $\Views_r$ followed by the inclusion into $\Beh_H$. The same holds for $s$. The Yoneda calculation above recovers the simulator from the identity attacker~\cite{security-yoneda}. The two resource factorizations of Theorem~\ref{thm:correctness-security-mutual-factorization} are arrows in the slice $\mathcal C/H$.
\end{enumerate}
\end{proof}

\begin{remark}[Multipartite resource presentation]\label{rem:multipartite-resource-presentation}
The same states and correct tuples form
\[
 \Res_P(\mathcal C)=\int\mathcal C(I,T_P-),
 \qquad T_P(\mathbf A)=\bigotimes_iA_i.
\]
Tensor is the existing grouping permutation applied to states. Secure conversions form its wide symmetric monoidal subcategory; retaining witnesses gives the category of Theorem~\ref{thm:secure-conversion-composition}. Grouping a coalition is a strong monoidal functor $G_J:\mathcal C^P\to\mathcal C^{P\setminus J}\times\mathcal C$. With the remaining tensor $T_H$ and state functor $R=\mathcal C(I,-)$, the square from $\int RT_HG_J$ to $\int RT_H$ over $G_J$ is a pullback: its objects are precisely grouped states, and its arrows are the same correct party tuples. In a non-strict presentation this is the canonical equivalent pullback. Secure conversions are the inverse image of those tested against the unrestricted joint coalition component. This recovers the resource and attack constructions of~\cite[Sections~3--5]{crypto-bk}; their simulator witnesses and composition have already been constructed as game squares.
\end{remark}

\subsection{Scott observations, public frames, and private relations}\label{subsec:public-observations}
The public frame records assertions observable through the chosen participant boundary. For a finite participant set \(P\), a coalition \(J\subseteq P\), and \(r:I\to A_P=\bigotimes_{i\in P} A_i\), let \(\rho_r:S_r\to A_P\) represent \(r\), write \(D_r=\D(S_r)\), and put \(\partial_Jx=\rho_r[x]\cap A_J\). The boundary map transports configurations; observations are Scott-open sets. Here the domain order is inclusion, rather than the polarized Scott order used to calculate interaction. Proposition~\ref{prop:effective-configuration-domain} supplies its finite compact basis; see also~\cite{effective-domains,pub-escardo}.

\begin{proposition}[Generated public frame]\label{prop:generated-public-frame}
For finite \(k\in\Conf_f(A_J)\), put
\[
B_k^{r,J}=\{x\in D_r:k\subseteq\partial_Jx\},\qquad
\Pub_J(r)=\left\{\bigcup_{k\in L}B_k^{r,J}:L\subseteq\Conf_f(A_J)\right\}.
\]
Each \(B_k^{r,J}\) is Scott open and \(\Pub_J(r)\) is a subframe of \(\Op(D_r)\). It is exactly the image of
\[
\partial_J^*:\Op(\D(A_J))\longrightarrow\Op(D_r),\qquad U\longmapsto\partial_J^{-1}U.
\]
Moreover \(J\subseteq K\) implies \(\Pub_J(r)\subseteq\Pub_K(r)\).
\end{proposition}
\begin{proof}\vmode
\leavevmode
\begin{enumerate}[label=(\roman*)]
\item \textbf{Continuity and finite generators.}
Boundary images preserve inclusions and directed unions.
A finite \(k\) contained in the boundary of a directed union is already contained in one member, proving openness.

\item \textbf{Frame operations.}
The top is \(B_\varnothing\), and \(B_k\cap B_\ell=B_{k\cup\ell}\) when the union is a configuration, and is empty otherwise.
Arbitrary unions and finite intersections therefore remain in \(\Pub_J(r)\).

\item \textbf{Generation of all public opens.}
Every Scott open of \(\D(A_J)\) is a union of \(\uparrow k\), giving the image statement.

\item \textbf{Observer monotonicity and scope.}
Boundary projection proves observer monotonicity.
None of these unions is required to be enumerable.
\end{enumerate}
\end{proof}

The finite-generator calculation also applies to the cumulative execution and audit displays. Retain their internal history events and project to the actual exposed interface. Historical origin assertions are opens of this history domain, and their public approximation uses that projection.

After hiding, Lemma~\ref{lem:observation-descent-hiding} identifies exactly which opens descend to the visible domain. A private origin can remain part of the retained evidence without becoming an event of the hidden strategy.

\begin{proposition}[Private indistinguishability and public information order]\label{prop:private-information-order}
The relations induced by the public frame satisfy
\begin{align*}
x\preceq_Jy
&\quad\Longleftrightarrow\quad
\forall U\in\Pub_J(r)\;(x\in U\Rightarrow y\in U)
\quad\Longleftrightarrow\quad\partial_Jx\subseteq\partial_Jy,\\
x\equiv_Jy
&\quad\Longleftrightarrow\quad
\forall U\in\Pub_J(r)\;(x\in U\Longleftrightarrow y\in U)
\quad\Longleftrightarrow\quad\partial_Jx=\partial_Jy.
\end{align*}
The relation \(K_J(r)=\{(x,y):x\equiv_Jy\}\), ordered componentwise, is an internal equivalence relation in \(\mathbf{DCPO}\).
\end{proposition}
\begin{proof}\vmode
\leavevmode
\begin{enumerate}[label=(\roman*)]
\item \textbf{Boundary reconstruction.}
For any \(e\in\partial_Jx\), use the finite boundary configuration \(k=\downarrow e\).
Preservation of membership in \(B_k\) implies \(e\in\partial_Jy\).
This proves the nontrivial implication for the preorder; the reverse follows from upward closure.

\item \textbf{Observational equivalence.}
The equivalence is its symmetric part.

\item \textbf{Kernel-pair structure.}
Directed unions preserve equality of boundaries, so \(K_J(r)\) is a dcpo with componentwise joins.
Its projections, diagonal, symmetry, and the transitivity map on compatible triples are continuous coordinate maps.

\item \textbf{Observer enlargement.}
Observer enlargement refines the equivalence relation.
\end{enumerate}
\end{proof}

The public frame, rather than only the private equivalence relation, retains positive information order. Fibres of \(\partial_J\) describe alternative internal configurations with the same current public boundary; they do not assert equality of future interactions.

\begin{proposition}[Public approximation adjunction]\label{prop:interior}
The inclusion \(i_J:\Pub_J(r)\hookrightarrow\Op(D_r)\) has right adjoint
\[
R_J(U)=\bigcup\{V\in\Pub_J(r):V\subseteq U\}.
\]
The operator \(\Box_J=i_JR_J\) is monotone, deflationary, idempotent, preserves finite meets, and has precisely the public opens as fixed points.
\end{proposition}
\begin{proof}\vmode
\leavevmode
\begin{enumerate}[label=(\roman*)]
\item \textbf{Adjunction.}
For public \(V\), \(V\subseteq U\) iff \(V\subseteq R_J(U)\).
This is the adjunction.

\item \textbf{Deflation and idempotence.}
The greatest-public-open formula gives deflationarity and idempotence.

\item \textbf{Finite meets and top.}
The intersection of two greatest public subopens is public and contained in the intersection, proving finite-meet preservation in both directions.
Top is public.
\end{enumerate}
\end{proof}

\subsection{The public-observation fibration and hiding}
We next compare public assertions as the underlying display changes or internal events are hidden. Pullback transports assertions along display maps; the hiding calculation determines which assertions depend only on the visible history.

At a fixed participant interface, Section~\ref{subsec:strategy-levels} constructs the hom-category of displays \(\mathbb{CG}(I,A_P)\). A display map \(h:r\Rightarrow s\) satisfies \(\rho_sh=\rho_r\). Consequently
\[
\D(h)^{-1}B_k^{s,J}=B_k^{r,J},\qquad
\Pub_J:\mathbb{CG}(I,A_P)^{\mathrm{op}}\longrightarrow\Frm.
\]
The Grothendieck construction of its underlying poset functor is a split fibration, by the direct calculation below; see also \cite[Chapter 1]{doct-jacobs}.

\begin{proposition}[Public-observation fibration]\label{prop:fibration}
Over \(r\), the fibre is \(\Pub_J(r)\). There is an arrow \((r,U)\to(s,V)\) over \(h\) precisely when \(U\subseteq\D(h)^{-1}V\). The chosen cartesian lift is
\[
\begin{tikzcd}[column sep=large]
(r,h^*V)\arrow[r]\arrow[d]& (s,V)\arrow[d]\\
r\arrow[r,"h"']&s.
\end{tikzcd}
\]
These fibrations respect the parallel-interface comparisons.
\end{proposition}
\begin{proof}\vmode
\leavevmode
\begin{enumerate}[label=(\roman*)]
\item \textbf{Cartesian lifting.}
For \(g:t\Rightarrow r\), the condition \(W\subseteq(hg)^*V\) is exactly \(W\subseteq g^*h^*V\); this is the unique-factorization property of the displayed lift.

\item \textbf{Split functoriality.}
Identity and composition of inverse images are strict.

\item \textbf{Parallel public opens.}
Parallel configuration domains are products, and
\(B_{(k,\ell)}^{r\otimes s,J}=B_k^{r,J}\times B_\ell^{s,J}\).
The compact pairs form the basis of the product configuration domain.
Thus all parallel public opens are unions of these rectangles.

\item \textbf{Naturality and display isomorphisms.}
Inverse images preserve the rectangles; the coherence maps permute and rebracket the same finite configurations.
Display isomorphisms induce the same identification of public generators, so their public frames are independent of the chosen representative.
\end{enumerate}
\end{proof}

This fibration is on display/configuration comparisons. A simulation square is compared through its real and simulated interaction displays at the same exposed boundary. A strategy retains its witnessed interaction span; this construction uses configuration maps only for actual display comparisons.

\begin{lemma}[Constructed descent through hiding]\label{lem:observation-descent-hiding}
For an interaction event structure \(E\) and its visible restriction \(E|W\), let
\[
v:\D(E)\to\D(E|W),\quad x\mapsto x\cap W,
\qquad
j:\D(E|W)\to\D(E),\quad y\mapsto\downarrow_Ey.
\]
Both maps are continuous, \(vj=1\), and \(jv(x)\subseteq x\). A Scott open \(U\subseteq\D(E)\) descends through \(v\) iff it is constant on the fibres of \(v\).
\end{lemma}
\begin{proof}\vmode
\leavevmode
\begin{enumerate}[label=(\roman*)]
\item \textbf{Causal-closure section.}
Causal closure of a visible configuration is consistent: every finite collection of its events lies in the finite causal closure of a finite consistent collection of visible events.
It is down-closed by construction.

\item \textbf{Continuity and retraction.}
Both maps preserve directed unions; visible downclosure gives \(vj=1\).

\item \textbf{Descent of fibre-constant opens.}
For a fibre-constant \(U\), take \(V=j^{-1}U\), which is Scott open.
Since \(x\) and \(jv(x)\) have the same visible restriction, \(U=v^{-1}V\).

\item \textbf{Converse.}
The converse is immediate.
\end{enumerate}
\end{proof}

For external observations the relevant square is
\[
\begin{tikzcd}[column sep=large]
\D(E)\arrow[r,"v"]\arrow[d,"\partial_J^{\rm int}"']&
\D(E|W)\arrow[d,"\partial_J^{\rm comp}"]\\
\D(A_J)\arrow[r,equal]&\D(A_J).
\end{tikzcd}
\]
It gives \(\Pub_J(\mathrm{int})=v^{-1}\Pub_J(\mathrm{comp})\). The causal-past section proves descent for hiding; the external display still determines which retained occurrence distinctions are publicly accessible. Secured interaction remains the composition construction: the above argument does not replace it by an unrestricted pullback.

\begin{construction}[Configuration comparisons over a boundary]\vmode\label{con:boundary-dcpo}
\leavevmode
\begin{enumerate}[label=(\roman*)]
\item \textbf{Boundary slice.}
For the existing boundary domain $D_B=\D(B)$, form the slice
$\mathsf{Bd}_B=\mathbf{DCPO}/D_B$.
Its objects are continuous maps
$d:D\to D_B$ and its arrows $h:(D,d)\to(E,e)$ satisfy $eh=d$.

\item \textbf{Public frame.}
Set $\Pub(d)=\im(d^*:\Op(D_B)\to\Op(D))$.
This extends the public
frame construction from displays to their configuration domains; it does not
interpret an arbitrary game strategy as a function of configurations.
\end{enumerate}
\end{construction}

\begin{proposition}[Configuration observation fibration]\label{prop:configuration-fibration}
Inverse image defines $\Pub:\mathsf{Bd}_B^{\op}\to\Frm$. Each
$h^*:\Pub(e)\to\Pub(d)$ is surjective. The Grothendieck construction
of its underlying poset functor is a split fibration extending
Proposition~\ref{prop:fibration}. The private relation at $d$ is the kernel
pair $K_d=D\times_{D_B}D$, and
$K_d=(h\times h)^{-1}K_e$.
\end{proposition}
\begin{proof}\vmode
\leavevmode
\begin{enumerate}[label=(\roman*)]
\item \textbf{Inverse-image transport.}
For every boundary open $U$, $h^{-1}e^{-1}U=d^{-1}U$; these inverse
images exhaust $\Pub(d)$ and preserve arbitrary joins and finite meets.

\item \textbf{Split cartesian lifting.}
The cartesian lift of $V$ is $(D,h^*V)\to(E,V)$, since
$W\subseteq(hg)^*V$ is precisely $W\subseteq g^*h^*V$.
Inverse images respect identities and composition on the nose.

\item \textbf{Kernel characterization.}
The Scott topology of a dcpo is $T_0$: if $x\not\leq y$, the complement
of the Scott-closed principal ideal $\downarrow y$ contains $x$ and not
$y$.
Hence two points of $D$ agree on every public open exactly when
their images under $d$ coincide.

\item \textbf{Kernel-pair dcpo.}
The kernel pair is a dcpo with
coordinatewise directed joins; its equivalence-relation structure maps
are continuous coordinate maps.

\item \textbf{Naturality and restriction.}
Finally $eh=d$ proves the kernel
equation directly.
Restricting along the configuration-domain functor
recovers the display fibration.
\end{enumerate}
\end{proof}

\begin{proposition}[Public frames transported by finite comparisons]\label{prop:comparison-public-frames}
Let $S,T$ be two displays on the same exposed arena $B$ related by a
finite causal or complete trace certificate. Their configuration-domain pullback
\[
\begin{tikzcd}[column sep=large]
 Z=\D(S)\times_{\D(B)}\D(T)\arrow[r,"q"]\arrow[d,"p"']
 \arrow[dr,phantom,"\lrcorner",very near start]
 &\D(T)\arrow[d,"\partial_T"]\\
 \D(S)\arrow[r,"\partial_S"']&\D(B)
\end{tikzcd}
\]
induces frame isomorphisms
$\Pub(\partial_S)\xrightarrow{p^*}\Pub(\partial_Z)
\xleftarrow{q^*}\Pub(\partial_T)$. This transports current public
information and its private kernel through the actual comparison.
\end{proposition}
\begin{proof}\vmode
\leavevmode
\begin{enumerate}[label=(\roman*)]
\item \textbf{Configuration pullback.}
Equality of the two boundaries is closed under coordinatewise directed
joins, so $Z$ and its projections are objects and arrows of
Construction~\ref{con:boundary-dcpo}.
Proposition~\ref{prop:configuration-fibration}
gives surjectivity of both inverse-image maps.

\item \textbf{Finite separation of public opens.}
To prove injectivity of
$p^*$, take distinct public opens $U,V$ of $\D(S)$, and, interchanging
them if necessary, $x\in U\setminus V$.
A finite subconfiguration
$x_0\subseteq x$ already belongs to $U$.
It cannot belong to $V$, since
$V$ is upward closed.

\item \textbf{Injectivity from certificate coverage.}
The certificate supplies a finite configuration
$y_0$ of $T$ with the same displayed boundary.
Thus $(x_0,y_0)\in Z$
separates $p^*U$ from $p^*V$.

\item \textbf{Frame isomorphisms.}
The reverse coverage proves injectivity of
$q^*$.
Preservation of all joins and finite meets was already proved,
so these bijections are frame isomorphisms.

\item \textbf{Kernel comparison.}
The kernel comparison is
the kernel equation of Proposition~\ref{prop:configuration-fibration}.
\end{enumerate}
\end{proof}
This pullback compares configuration domains; it is not the secured
interaction used to compose games. The causal-order clause gives the stronger comparison of Lemma~\ref{lem:crypto-causal-congruence}; trace equality gives Theorem~\ref{thm:trace-full-abstraction}. For the frame isomorphism only equality of finite boundary images is used, which follows from either certificate. Thus current private alternatives, their fibrational
transport, and interactive security have separate proved connections.

The map $\partial_J$ records which events are already visible. To determine what a participant can learn by making another request, we must also specify the operations available on the values it holds.

\subsection{Public operation algebras and labelled presentations}\label{subsec:information-interfaces}
To construct a simulator, we must identify which operations are available through its interface. We therefore realize a diagram of private operations as a public reference service, then characterize the values accessible through its commands. A private set-theoretic map is available to a participant only through an actual command port. Products and images here are formed in $\Set$; their service interpretations are strategies.

\begin{construction}[Reference realization of an operation diagram]\vmode\label{con:reference-interface}
\leavevmode
\begin{enumerate}[label=(\roman*)]
\item \textbf{Private reference table.}
For sorted carriers $V_\tau$, realize a value by an entry in a
client-local reference table.

\item \textbf{Command lookup.}
A command names an operation and its
operand references; lookup checks ownership, issuance and sorts.

\item \textbf{Results and constants.}
A successful value is stored under the next reference, a public test
returns its literal result, and invalid lookup or an undefined
operation returns failure.
Public constants have nullary commands.

\item \textbf{Private imports and opacity.}
Private values enter only through declared import ports.
Thus neither
an identifier nor a private product representation supplies an
undeclared observation.
\end{enumerate}
\end{construction}

\begin{definition}[Operation signature and reference state]\label{def:reference-controller-signature}
Each operation $o$ has an input-sort list
$(\tau_1(o),\ldots,\tau_{a(o)}(o))$ and result kind
$\mathsf{value}(\tau(o))$ or $\mathsf{test}(O_o)$. Set
$Y_o=V_{\tau(o)}$ or $Y_o=O_o$, respectively.
Its private-state evaluator is
\[
 \mathsf{ev}_o:B\times\prod_jV_{\tau_j(o)}
 \longrightarrow B\times(Y_o+\{\mathsf{fail}\}).
\]
For a pure partial operation $B=1$, and failure leaves the state
unchanged. Stateful operations specify their failure-state update.
Possibly nonreturning evaluators use the lift of
Section~\ref{sec:effective-foundations}.
Reference states and public alphabets are
\[
 \begin{gathered}
 s=(b,(\rho_c)_c,(n_c)_c),\qquad
 \rho_c:\underline{n_c}\to\coprod_\tau V_\tau,\qquad
 \mathsf H=\mathsf{Client}\times\mathbb N,\\
 Q_{\rm pub}=\coprod_{c,o}\mathsf H^{a(o)},\qquad
 R_{\rm pub}=\mathsf H+
       \coprod_{o:\mathsf{test}(O_o)}O_o+\{\mathsf{fail}\}.
 \end{gathered}
\]
Only finitely many counters are nonzero; result tags are retained.
Write $\rho_{c,\tau}$ for the restriction to sort $\tau$.
The private field $b$ contains operation caches and, in generated-name
instances, the cumulative allocation world. Initial tables contain
the declared entries. Constants, copying and allocation are
respectively nullary value operations, identity on values, and the
specified stateful introductions. A private import has its own typed
port, outside $Q_{\rm pub}$.
\end{definition}

\begin{construction}[Total reference transition]\vmode\label{con:reference-transition-algorithm}
\leavevmode
\begin{enumerate}[label=(\roman*)]
\item \textbf{Operand validation.}
On $(c,o,\vec h)$, an operand $(c_j,i_j)$ is valid exactly when
$c_j=c$, $i_j<n_c$, and $\rho_c(i_j)$ has sort $\tau_j(o)$.
Invalid lookup returns $(s,\mathsf{fail})$.

\item \textbf{Operation evaluation.}
Otherwise evaluate $\mathsf{ev}_o(b,\vec v)=(b',z)$.

\item \textbf{Failure and test replies.}
Failure and tests retain the tables and counters, update $b$ to $b'$,
and return the tagged result.

\item \textbf{Value issuance.}
A value result of sort $\tau$ appends
\[
 \rho'_c=\rho_c[n_c\mapsto(\tau,z)],\qquad
 h_{\rm new}=(c,n_c),\qquad n'_c=n_c+1.
\]
Other tables are unchanged.

\item \textbf{Private imports.}
A private import first performs its
specified world transport and then this issuance update.

\item \textbf{Total transition.}
These clauses define $\delta:S\times Q_{\rm pub}\to S\times R_{\rm pub}$.
\end{enumerate}
\end{construction}

\begin{lemma}[Reference-state preservation]\label{lem:reference-state-preservation}
This transition preserves domains, ownership and sorts, and the
retained-world support invariant in generated-value instances.
Issued labels depend only on the client and its counter.
\end{lemma}
\begin{proof}\vmode
\leavevmode
\begin{enumerate}[label=(\roman*)]
\item \textbf{Table domains and sorts.}
Lookup failure changes nothing; tests and operation failure preserve
the tables.
Value issuance appends one correctly sorted entry at the
next index.

\item \textbf{World support.}
Operation and import updates retain its support in the
current world.

\item \textbf{Issued labels and termination.}
The formula for $h_{\rm new}$ proves the label claim.
Each transition uses finitely many lookups and one terminating
evaluation in the total instance.
\end{enumerate}
\end{proof}

\begin{remark}[Opacity is a property of the exposed arena]\label{rem:representation-opacity}
All strategies on the public arena remain available as opponents,
and Scott opens inspect its actual labels. Reference replies expose
issued identifiers; later observations are the declared commands.
Exporting a raw private tree instead would expose that tree to
principal opens. Equality of referenced values is an explicit command,
and comparisons below match the reference labels literally.
\end{remark}

\begin{definition}[The equality-exposing public diagram]\label{def:public-equality-diagram}
For the generated instances the public operation diagram is the concrete
constructor/destructor reference diagram instantiated in Appendix~\ref{app:cryptographic-examples}, including
its value-equality commands. Write
$\mathsf{eq}_\tau$ for that command at sort $\tau$ and
$\mathsf{true},\mathsf{false}$ for its literal result tags.
Its specified interpretation satisfies
\[
 \mathsf{eq}_\tau(v,w)=\mathsf{true}\ \Longleftrightarrow\ v=w.
\]
The command is applied by evaluating its two recipes into references
and then using the actual table lookup and value test. Thus for any
successful recipes $R,S$ of sort $\tau$ its positive result is
precisely $\crEval_\phi R=\crEval_\phi S$.
The notation below refers to this equality-exposing diagram; other
operation diagrams retain their declared observations.
\end{definition}

\paragraph{Generated presentations.}\label{subsec:crypto-categories}
\begin{construction}[Accessible algebra and its labelled generators]\vmode\label{con:accessible-labelled-algebra}
\leavevmode
\begin{enumerate}[label=(\roman*)]
\item \textbf{Available generators.}
Fix the actual available handles $W$, their values $\phi:W\to V$, and the public constants.
Allocations and selected receipts already performed are included among these sources.

\item \textbf{Public recipes.}
Interpret finite pure value programs in the existing calculus, using only the public introductions, strict operations and tests of the connected reference service.
A \emph{recipe} below is an abbreviation for such a finite value computation, not an additional language.

\item \textbf{Accessible image.}
Write $\mathsf{Recipes}^{\downarrow}_{\phi,\tau}(W)$ for its successful evaluation domain of result sort $\tau$, and form the image
\[
 A_{\phi,\tau}=\im\bigl(
 \mathsf{Recipes}^{\downarrow}_{\phi,\tau}(W)
 \xrightarrow{\crEval_\phi}V_\tau\bigr).
\]
Retain the generating map $\phi:W\to A_\phi$.

\item \textbf{Closure under public operations.}
A finite tuple of successful computations can be sequenced and followed by a public operation; hence its defined result again belongs to $A_\phi$.

\item \textbf{Observable results.}
Observable result sorts retain their prescribed carriers.
Data reading returns ordinary data, whose reintroduction uses the actual data command; it does not introduce arbitrary private values.

\item \textbf{Closure notation.}
For a sorted subset $K\subseteq V$, the same image construction with the inclusion of $K$ as source is denoted $\mathsf{DY}(K)$.
Let $\mathsf{Cl}_{\rm pub}(V)$ be the poset of subobjects containing the public constants and closed under the defined public operations, and $U$ its inclusion into $\mathsf{Sub}(V)$.

\item \textbf{Witnessed derivability.}
The judgement $\Gamma\vdash_{\rm rec}t$, when used, means membership in this image; a witnessed use retains the successful computation and evaluates it.
Image elimination supplies no choice of a computation from membership alone.

\item \textbf{Static equivalence.}
Two labelled frames are \textbf{statically equivalent} when their finite public computations have the same successes, equality tests and data-reading results.

\item \textbf{Algebra comparisons.}
A map of their generated public algebras commutes with the public operations and fixes constants and observable results; an isomorphism also reflects their domains.
These are comparisons of a fixed available frame, separately from future interactive behaviour~\cite[Section 4]{crypto-abf}.
\end{enumerate}
\end{construction}

\begin{proposition}[The accessible-algebra reflection]\label{prop:dy-reflection}
The construction $\mathsf{DY}$ is left adjoint to $U$. Its induced monad is idempotent and preserves directed unions. Its elements are exactly the successful finite public computations on the specified sources.
\end{proposition}
\begin{proof}\vmode
\leavevmode
\begin{enumerate}[label=(\roman*)]
\item \textbf{Least closed extension.}
If $K\subseteq U(L)$ and $L$ is closed, induction on a finite computation puts every returned value in $L$.
Conversely the source computations give $K\subseteq\mathsf{DY}(K)$.

\item \textbf{Reflection and closure laws.}
Thus
\[
 \mathsf{DY}(K)\subseteq L\quad\Longleftrightarrow\quad K\subseteq U(L).
\]
Closure was constructed by program composition, giving idempotence.

\item \textbf{Finitary generation.}
A finite computation uses finitely many sources, which lie together in one member of a directed union; this proves preservation of that union.
\end{enumerate}
\end{proof}
Joint knowledge is the reflected union $\mathsf{DY}(K\cup L)$: a ciphertext in one frame and its key in another can jointly reveal its plaintext. The generating labels must also be retained; frames $(x,y,z)\mapsto(k,k,\ell)$ and $(x,y,z)\mapsto(k,\ell,k)$ have the same generated values but different $x=y$ tests.

\begin{theorem}[Static equivalence as a canonical algebra isomorphism]\label{thm:frame-isomorphism}
For frames with the same sorted handle domain and public literal alphabet,
\[
 \phi\sim_s\psi\quad\Longleftrightarrow\quad
 \text{there is a unique public-algebra isomorphism }
 a:A_\phi\xrightarrow{\cong}A_\psi\text{ with }a\phi=\psi.
\]
This characterizes the constructor/destructor frame observations just defined; it does not identify a current frame with an entire interactive process.
\end{theorem}
\begin{proof}\vmode
\leavevmode
\begin{enumerate}[label=(\roman*)]
\item \textbf{Evaluation quotient.}
Static equivalence identifies the successful recipe domains and
their equality kernel pairs.
The evaluation images therefore give
the unique induced bijection
\[
 a(\crEval_\phi R)=\crEval_\psi R,
 \qquad
\begin{tikzcd}[column sep=large]
 \mathsf{Recipes}^{\downarrow}_\phi(W)
  \arrow[r,two heads,"\crEval_\phi"]
  \arrow[dr,two heads,"\crEval_\psi"']
 &A_\phi\arrow[d,"a","\cong"']\\
 &A_\psi .
\end{tikzcd}
\]
Appending an operation to operand recipes proves commutation and
domain reflection.

\item \textbf{Observable results and uniqueness.}
Test and reading observations fix their result
carriers; nullary and handle recipes fix the labelled generators.
Recipe induction proves uniqueness.
No representatives are chosen.

\item \textbf{Converse by recipe induction.}
Conversely that induction along an operation-preserving isomorphism
gives matching success and $a\crEval_\phi R=\crEval_\psi R$;
injectivity reflects equality and the fixed observable carriers
give the same literal results.
These are exactly the stated observations.
\end{enumerate}
\end{proof}

\begin{corollary}[Derived handles]\label{cor:derived-handles}
Adding corresponding successful computations as fresh labelled handles preserves static equivalence and leaves the generated public algebras unchanged.
\end{corollary}
\begin{proof}\vmode
\leavevmode
\begin{enumerate}[label=(\roman*)]
\item \textbf{Substitution of derived handles.}
Substitution replaces the new handle by its producing computation.
The canonical isomorphism of Theorem~\ref{thm:frame-isomorphism} therefore preserves its value and its labelled generating map.
\end{enumerate}
\end{proof}

\begin{definition}[Labelled public presentations]\label{def:crypto-presentations}
For a fixed sorted handle domain $W$, let $\mathsf{Pres}_{\rm pub}(W)$ be the groupoid whose objects are generated public algebras with their labelled generating maps $\phi:W\to A_\phi$, and whose arrows are public-algebra isomorphisms $a$ satisfying $a\phi=\psi$. Theorem~\ref{thm:frame-isomorphism} makes each hom-set empty or a singleton and identifies its connected components with static-equivalence classes. The observable result sorts and the public literal alphabet are part of this definition. Forgetting $W\to A_\phi$ would lose the correlations tested by equality of handle recipes.
\end{definition}

\begin{proposition}[Closed invariants and fixed values]\label{prop:crypto-information-invariants}
The knowledge reflection gives two constructions of nondeducibility certificates. A sorted public-operation-closed subobject $L\subseteq V$ containing the public constants and the frame satisfies $A_\phi\subseteq L$. For a group $G$ of public-algebra automorphisms fixing the labelled frame and public constants,
\[
 A_\phi\subseteq V^G
 =\{v:\alpha(v)=v\text{ for every }\alpha\in G\}.
\]
In the first construction a value outside $L$ is not derivable; in the second a value moved by one such automorphism is not derivable.
\end{proposition}
\begin{proof}\vmode
\leavevmode
\begin{enumerate}[label=(\roman*)]
\item \textbf{Closed-invariant bound.}
The first inclusion is the adjunction of Proposition~\ref{prop:dy-reflection}.

\item \textbf{Closure of fixed values.}
For the second, the fixed values contain the generators.
If an operation is defined on fixed operands, its value is fixed because every $\alpha$ commutes with that operation.
Thus $V^G$ is closed, and the first construction applies.

\item \textbf{Recipe induction and scope.}
Equivalently, induction on a successful recipe fixes its value.
The argument proves an inclusion; it does not identify all fixed values with recipe values.
\end{enumerate}
\end{proof}

\begin{definition}[Public comparison presentations]\label{def:crypto-public-comparison}
A public comparison presentation for two frames with the same handle domain is a span in $\mathsf{Pres}_{\rm pub}(W)$,
\[
 (A_\phi,\phi)\xleftarrow[\cong]{e_0}(N,\nu)
                       \xrightarrow[\cong]{e_1}(A_\psi,\psi).
\]
The middle presentation may be constructed by replacing protected values by typed markers, or by using an explicitly calculated algebra automorphism. A marker's operations include the unsuccessful destructor cases; it is not an additional operation for opening the protected value.
\end{definition}

\begin{proposition}[Public comparison and derived operations]\label{prop:crypto-public-comparison}
A public comparison presentation gives the canonical comparison $e_1e_0^{-1}:A_\phi\cong A_\psi$ and hence static equivalence. Every finite public program commutes with both expansion maps, including its definedness and observable results. Adding corresponding successful recipe handles preserves the comparison presentation.
\end{proposition}
\begin{proof}\vmode
\leavevmode
\begin{enumerate}[label=(\roman*)]
\item \textbf{Labelled algebra isomorphism.}
The composite $e_1e_0^{-1}$ preserves labelled generators, so
Theorem~\ref{thm:frame-isomorphism} gives static equivalence.

\item \textbf{Primitive operation equations.}
For each primitive operation and $i=0,1$,
\[
 f_N(\vec u)\downarrow\ \Longleftrightarrow\
 f_{A_i}(e_i\vec u)\downarrow,
 \qquad e_if_N(\vec u)=f_{A_i}(e_i\vec u)
\]
on that common domain.

\item \textbf{Strict-program induction.}
Strict composition preserves these equations,
including failure of an unused intermediate instruction.
Induction
gives the program statement, and Corollary~\ref{cor:derived-handles}
extends the generating triangles.
\end{enumerate}
\end{proof}

The free-constructor examples in Appendix~\ref{app:cryptographic-examples} construct the maps $e_i$ by normalization. Their images contain the frame and are closed under every defined public operation, so they contain every recipe value. Conversely, expansion sends generated normal forms to recipe values. Root separation and private-operand invariants prove injectivity, while the destructor cases verify preservation and reflection of definedness.

The scalar examples use field automorphisms instead. In each case the calculation constructs an arrow of Definition~\ref{def:crypto-public-comparison} and verifies its operation equations.

\subsection{Reference realization and comparison transport}\label{subsec:reference-comparison-transport}
A labelled public-algebra comparison already transports finite public programs. To compare future interaction, we extend it to controller states and verify the updates for allocation, import and failure. These equations give the comparison coalgebra. A distributed network additionally retains its causal certificate, so its communication dependencies remain part of the comparison.\paragraph{State coordinates of representation transport.}
For a stateful operation the operation diagram includes its state
coordinate $B$. Write $\beta:B\to B'$ for its representation
transport and $h_o:Y_o\to Y'_o$ for its result transport, which is
the identity on a public test carrier. The operation-commutation
equation then has its full type
\[
 (\beta\times(h_o+1))\mathsf{ev}_o
   =\mathsf{ev}'_o
       \left(\beta\times\prod_j h_{\tau_j(o)}\right).
\]
For pure operations $B=1$. For the constructed banks and workers,
$\beta$ transports stored value fields and selected-supplier
references and fixes control tags and counters. At allocation or
import its successor is the explicitly constructed extension of
that transport. Thus representation independence checks the
whole state update, including its failure branch.

\begin{proposition}[Evaluator representation independence]\label{prop:generic-evaluator-comparison}
A bijection between two accessible value representations that commutes
with every exposed operation, preserves and reflects its definedness,
and fixes public constants and test results lifts to a strict comparison
of their reference controllers. The lifted map preserves issued
identifiers, ownership and counters. The same construction applies to
an evolving comparison retained as part of the comparison state.
\end{proposition}
\begin{proof}\vmode
\leavevmode
\begin{enumerate}[label=(\roman*)]
\item \textbf{Transport of reference states.}
Send $\rho$ to $h\rho$, transport the private state by $\beta$, and
fix domains, ownership and counters.

\item \textbf{Comparison graph.}
For a fixed representation,
this is the graph instance of the paired update in
Construction~\ref{con:evolving-comparison-coalgebra} below.

\item \textbf{Strict transition equations.}
Lemma~\ref{lem:comparison-command-induction} verifies its clauses:
operation commutation gives corresponding new entries, domain
reflection gives identical failures, and tests fix their literals.
Thus the graph's projections give the strict square.

\item \textbf{Evolving comparison.}
Retaining
$(s,t,h)$ instead gives the evolving span, with the explicit
allocation and import extensions calculated in each instance.
\end{enumerate}
\end{proof}

\begin{construction}[The evolving comparison coalgebra]\vmode\label{con:evolving-comparison-coalgebra}
\leavevmode
\begin{enumerate}[label=(\roman*)]
\item \textbf{Comparison-state carrier.}
For the paired operation instance, let $M$ consist of triples
$(s,t,h)$ with equal table domains and counters, corresponding private
operation records, and sort-indexed isomorphisms $h_\tau$ between
the accessible value representations such that
$h_\tau\rho^s_{c,\tau}=\rho^t_{c,\tau}$.
Correspondence of private operation records means the explicit
paired-state relation constructed for that instance:
its operation, allocation and import clauses are part of the following
update, not a consequence inferred from table equality.

\item \textbf{Lookup and operation update.}
Given $m=(s,t,h)$ and a command, apply the lookup cases of
Construction~\ref{con:reference-transition-algorithm} simultaneously.
Failed lookup retains $h$.
A successful value operation extends the
tables by the corresponding results; a test has the same literal;
operation failure uses the paired failure-state clause.

\item \textbf{Import and allocation transport.}
A private
import or allocation extends $h$ by the instance's constructed
name/world or algebraic extension.

\item \textbf{Common reply and successor.}
Denote the resulting triple by
$m'$ and the common public reply by $r$.

\item \textbf{Strict coalgebra projections.}
This defines
\[
 \delta_M(m,a)=(m',r),\qquad
 (p\times1_R)\delta_M=\delta_S(p\times1_Q),\qquad
 (q\times1_R)\delta_M=\delta_T(q\times1_Q),
\]
where $p(s,t,h)=s$ and $q(s,t,h)=t$.
Thus
$\delta_M:M\times Q\to M\times R$ is a coalgebra span over the
two fixed command/reply alphabets.
\end{enumerate}
\end{construction}
\begin{lemma}[Command induction for comparison states]\label{lem:comparison-command-induction}
The preceding update lands in $M$ for the operation-preserving
representation map of Proposition~\ref{prop:generic-evaluator-comparison},
and for each explicitly extended comparison constructed below. Both
displayed projection equations hold. Hence every finite adaptive
command history has the same literal replies at the two projections.
\end{lemma}
\begin{proof}\vmode
\leavevmode
\begin{enumerate}[label=(\roman*)]
\item \textbf{Lookup correspondence.}
Lookup uses equal domains, owner tags and sorts.

\item \textbf{Values and tests.}
For a value operation,
$h f(\vec v)=f'(h\vec v)$ is exactly the new table-entry invariant;
preservation and reflection of the domain gives the same failure
case.
Tests fix their result carriers.

\item \textbf{Counters and world extension.}
Counters advance together.
For changing worlds the update is the actual extension map constructed
in the respective allocation/import lemma, whose restriction agrees
with $h$ on old values.

\item \textbf{Invariant and strictness.}
Thus the invariant is preserved in every
clause.
Projecting the clause gives exactly the corresponding
controller update.

\item \textbf{Adaptive command induction.}
Induct on command length.
Equal preceding replies
allow the same next adaptive request; the successor equation supplies
its equal reply.
\end{enumerate}
\end{proof}
This result concerns all future commands, including joint equality
tests. A comparison of two separately considered values does not
automatically compare the pair of them.

\paragraph{Game realization.}\label{subsec:crypto-games}
\begin{definition}[Total services and their simulation square]\label{def:formal-total-services-and-their-simulation-square}
For sets $Q,R$, a \textbf{total service} is $(S,s_0,\delta)$ with $s_0\in S$ and $\delta:S\times Q\to S\times R$. Its uninitialized carrier is equivalently a coalgebra $S\to(S\times R)^Q$. A \textbf{strict service simulation} from $\delta$ to $\delta'$ is a state map $a:S\to S'$ satisfying
\[
 (a\times\id_R)\delta=\delta'(a\times\id_Q).
\]
It compares initialized services when $a(s_0)=s'_0$. The equation is stated here before the word extension; Theorem~\ref{thm:service-behaviour-canonical-comparison} proves its behavioural consequence and records it as~\eqref{eq:extended-service-square}.
\end{definition}

\begin{definition}[Fixed call--return arena]\label{def:crypto-call-return}
For request set $Q$ and reply set $R$, let $A(Q,R)$ have as events the nonempty finite prefixes of alternating words
\[
 q_0^-r_0^+q_1^-r_1^+\cdots.
\]
Order is prefix order, and consistency means that the events lie on one branch. The request and reply sets are fixed before a particular service is chosen.
\end{definition}

\begin{lemma}[Stateful service strategy]\label{lem:crypto-stateful}
Let a service have a constructed state set $S$, initial state $s_0$, and total update function $\delta:S\times Q\to S\times R$. Its legal command histories form a receptive, courteous strategy on $A(Q,R)$.
\end{lemma}
\begin{proof}\vmode
\leavevmode
\begin{enumerate}[label=(\roman*)]
\item \textbf{Legal prefix display.}
Take the prefix subtree containing all requests and exactly the reply prescribed by $\delta$ after each request.
Map it into $A(Q,R)$ by inclusion.
Its finite configurations are branch prefixes and hence map injectively to configurations.

\item \textbf{Receptivity.}
At an even prefix every enabled negative request has exactly one extension.
At an odd prefix no negative move is enabled.
Thus the display is receptive.

\item \textbf{Courtesy.}
Every immediate causal edge is already an arena edge, so it is courteous.

\item \textbf{State evaluation and finite pasts.}
Internal state is calculated from the finite request history; it is not an extra public event label.
All branches have finite principal pasts.
\end{enumerate}
\end{proof}

\begin{lemma}[Controllers with private service calls]\label{lem:crypto-controllers}
A sequential client which accepts a request, makes finitely many sequential call--return service calls, and returns a reply has a receptive, courteous strategy on its fixed client arena and the duals of its fixed service arenas. It may also wait for a service reply indefinitely.
\end{lemma}
\begin{proof}\vmode
\leavevmode
\begin{enumerate}[label=(\roman*)]
\item \textbf{Control-history unfolding.}
Unfold its finite control histories, retaining request/reply values and branch choices.

\item \textbf{Added causal dependencies.}
A new service call is a positive event after the negative client request or the preceding negative service reply.
A return to the client is likewise positive after the last negative input.
These are the only added cross-port immediate dependencies, and all have sign $(-,+)$.

\item \textbf{Courtesy.}
Dependencies between a service request and reply, or between a client reply and the next request, are already in the corresponding arena.
Hence courtesy holds.

\item \textbf{Pending-call receptivity.}
While a call is pending, all possible replies to that call are accepted, with invalid data taking the failure branch.
There is no other outstanding call.

\item \textbf{Client-request receptivity.}
The next client request is arena-enabled only after the preceding client reply.
These observations identify every enabled negative move and its unique handling extension, proving receptivity.

\item \textbf{Finite histories and waiting.}
Unfolding records finite prefixes, so the principal pasts are finite.
Waiting is omission of a positive continuation, not refusal of an enabled negative move.
\end{enumerate}
\end{proof}

Local evaluation and private storage use these strategies. Transport uses any declared communication protocol strategy of Section~\ref{sec:protocol-games}, with its actual typed ports. A party's private-value transfer ports and evaluation internals are connected and hidden; its exposed port carries commands, references, and test results. An attacking client receives the same command interface with its own store.

\begin{definition}[Cryptographic communication games]\label{def:cryptographic-game-category}
Let $\mathcal A_{\rm crypt}$ extend $\mathcal A_{\rm comm}$ by the declared service arenas, including the one-event positive observation arena, closed under the empty arena, finite parallel product and polarity reversal. Define $\crCG$ to be the full subcategory of $\CG$ on these arenas. Its distinguished service strategies are the displays constructed from the actual operation diagrams and controllers above. Honest implementations are interpretations of network terms using these declarations; adversaries range over all ordinary strategies at their exposed interfaces.

The inclusion $U_{\rm crypt}:\crCG\to\CG$ retains the entire arena and strategy display. Thus a private representation supplies only the commands exposed by its service, not arbitrary additional public projections. The enhancement consists of these interpreted interfaces and service declarations inside ordinary concurrent games; it does not attach a security certificate to every game morphism.
Write $\mathsf{Net}_{\rm crypt}$ for the typed expression category of the paragraph ``\nameref{par:typed-network-interpretation}'' extended by these service declarations, with the same tensor, connection, forwarding and port-bending constructors and compact structural equations.
\end{definition}

\begin{theorem}[Compact closed cryptographic communication]\label{thm:compact-cryptographic-communication}
The category $\crCG$ is compact closed, and $U_{\rm crypt}$ is fully faithful and preserves the specified compact structure. The service and client declarations extend the existing network interpretation to a compact-preserving interpretation
\[
 \llbracket-\rrbracket_{\rm crypt}:\mathsf{Net}_{\rm crypt}\longrightarrow\crCG.
\]
Its restriction to the original communication declarations is their original interpretation. Independent networks retain their concurrent causal orders. Coded declarations have the effective interpretation of Section~\ref{subsec:effective-bicategory}.
\end{theorem}
\begin{proof}\vmode
\leavevmode
\begin{enumerate}[label=(\roman*)]
\item \textbf{Primitive strategies.}
Lemmas~\ref{lem:crypto-stateful} and~\ref{lem:crypto-controllers}
give the new receptive, courteous generators.

\item \textbf{Closed interface class.}
The selected arena
class is closed under tensor, unit and reversal, so the full
subcategory contains the copycat cups and caps and inherits their
snake and coherence equations from $\CG$.

\item \textbf{Extended interpretation.}
Extend Theorem~\ref{thm:abstract-communication-interpretation}
recursively on the network constructors: parallel is disjoint
parallel display, connection is secured interaction and hiding,
and bending uses these cups and caps.
Theorem~\ref{thm:ordinary-category}
supplies their equations and therefore the claimed interpretation.
It agrees with the original interpretation on its generators.

\item \textbf{Independent computations.}
Parallel events retain componentwise consistency; interaction
identifies matched occurrences and hiding restricts the resulting
causal order.
Independent computations are consequently retained.

\item \textbf{Effective presentation.}
The effective constructors of Section~\ref{subsec:effective-bicategory}
give the same recursion on coded declarations.
\end{enumerate}
\end{proof}

The contextual quotient of this declared interface category is
$\Csc=\crCG/{\eqsc}$, compact closed by Theorem~\ref{thm:contextual-congruence-calculus}. Equalities there retain all ambient tests at the exposed ports. Exact equations in $\crCG$ imply those equalities, but are kept separate when occurrence identity matters.

\paragraph{Canonical coalgebra comparisons.}\label{subsec:canonical-comparisons}
State comparisons are local proof witnesses for equations between service strategies. They are distinct from the simulator strategy in a resource comparison, which must itself use the actual ideal ports.

\begin{definition}[Behaviour and comparison carriers]\label{def:service-behaviour-comparison-carriers}
For fixed $Q,R$ put $B(X)=(X\times R)^Q$ and $\Omega=R^{Q^+}$. Define $\zeta(F)(q)=(F_q,F(q))$, where $F_q(w)=F(qw)$. The behaviour $\beta_S(s)$ of a service state sends a nonempty request word to its last reply, calculated by finite transition recursion. A witnessed service comparison is a span $S\xleftarrow p M\xrightarrow q T$ of strict $B$-coalgebra maps. A selected $m_0\in M$ compares its two projected initial states. The canonical comparison carrier is $S\times_\Omega T$.
\end{definition}

\begin{theorem}[Behaviour, replacement, and canonical comparison]\label{thm:service-behaviour-canonical-comparison}
The coalgebra $(\Omega,\zeta)$ is final. A strict service map
\begin{equation}\label{eq:extended-service-square}
 (a\times1_R)\delta_S=\delta_T(a\times1_Q)
\end{equation}
satisfies $\beta_Ta=\beta_S$. Corresponding initialized states have isomorphic fixed call--return strategy displays and can be replaced in every concurrent game context. The carrier $S\times_\Omega T$ has a canonical total service transition and is terminal among witnessed comparison spans from $S$ to $T$.
\end{theorem}
\begin{proof}\vmode
\leavevmode
\begin{enumerate}[label=(\roman*)]
\item \textbf{Final behaviour coalgebra.}
On a one-letter word the coalgebra equation determines the reply; on $qw$ it determines the derivative.
Finite induction therefore proves both that $\beta_S$ is a coalgebra map and that it is the unique such map to $\Omega$.
This includes $Q=\varnothing$, when the behaviour carrier is a singleton.

\item \textbf{Equality of initialized displays.}
Equation~\eqref{eq:extended-service-square} gives equal replies and corresponding successors, hence $\beta_Ta=\beta_S$ by the same induction.
The legal prefix subtrees of Lemma~\ref{lem:crypto-stateful} consequently coincide, including a last pending request.
Identity on these prefixes is the display isomorphism.

\item \textbf{Contextual replacement.}
Game composition and tensor transport it through any network context.

\item \textbf{Comparison-state transition.}
For $(s,t)\in S\times_\Omega T$, write $\delta_S(s,u)=(s',r)$ and $\delta_T(t,u)=(t',r')$.
Equality of behaviours at $u$ gives $r=r'$; equality at every $uw$ gives $\beta_S(s')=\beta_T(t')$.
Thus
\[
 \delta_M((s,t),u)=((s',t'),r)
\]
is the required total transition and both projections are strict.

\item \textbf{Terminal comparison span.}
For any comparison $S\xleftarrow a N\xrightarrow b T$, finality gives $\beta_Sa=\beta_Tb$.
The unique map $n\mapsto(a(n),b(n))$ to this carrier preserves replies and successors by the strict squares, proving terminality among comparisons.

\item \textbf{Initial-state criterion and scope.}
In particular an initial pair is comparable exactly when its behaviours agree.
This is a semantic carrier, not a decision procedure for behaviour equality; concrete invariants provide its computable witnesses when required.
\end{enumerate}
\end{proof}

\begin{proposition}[Constructing service certificates from comparison states]
\label{prop:service-history-certificate}
For the total services of Definition~\ref{def:service-behaviour-comparison-carriers}, let
$S\xleftarrow{p}M\xrightarrow{q}T$ be a strict comparison of their command/reply
transition systems with a chosen initial comparison state $m_0$.
For every finite adaptive public command history, iteration in $M$ constructs
two service paths with identical public requests and replies. Pairs obtained in this way form a complete trace certificate between the two initialized call--return displays.
The strict comparison is the pair of commuting squares
\[
\begin{tikzcd}[column sep=large,row sep=large]
 S\times Q \arrow[d,"\delta_S"']
 &M\times Q \arrow[l,"p\times1_Q"'] \arrow[r,"q\times1_Q"]
       \arrow[d,"\delta_M"]
 &T\times Q \arrow[d,"\delta_T"]\\
 S\times R
 &M\times R \arrow[l,"p\times1_R"] \arrow[r,"q\times1_R"']
 &T\times R.
\end{tikzcd}
\]
\end{proposition}
\begin{proof}\vmode
\leavevmode
\begin{enumerate}[label=(\roman*)]
\item \textbf{Iteration of strict comparison states.}
Use the total update functions already supplied by the service presentations.
If the
comparison state is $m$, strictness says that its transition on a request $a$
has one public reply $b$ and a successor $m'$ whose two projections are the
respective service successors on $a$, with that same reply $b$.
Starting at
$m_0$ and iterating gives the comparison for every finite request sequence.

\item \textbf{Adaptive and pending requests.}
At an adaptive step a context has received the same preceding public history,
so the same next request remains available on both sides.
Stop also at request
prefixes of the call/return presentation; the pending request is the same on
both sides.

\item \textbf{Surjective path projections.}
Every finite service path determines its request sequence and whether its
last request is pending, and is recovered
by this construction on the corresponding projection, since each reference
transition is deterministic.
The two projections onto path carriers are
therefore surjective.

\item \textbf{Trace certificate.}
Each path in an initialized call--return display specifies its sequence of requests and replies in order. The two projections therefore give a span of serialized configurations with surjective legs, and strictness makes their displayed words equal.

\item \textbf{Contextual equivalence.}
The equal-word span proves equality of service
trace languages; Theorems~\ref{thm:trace-full-abstraction} and~\ref{thm:complete-certificates} extend the comparison to arbitrary
ordinary contexts.
No inspection of the private representation is added to
the public command signature.
\end{enumerate}
\end{proof}

\begin{proposition}[Pullback composition of strict comparisons]\label{prop:service-pullbacks}
For strict maps $S\to T\leftarrow U$, the state pullback $S\times_TU$ has a unique total service making its projections strict. Comparison spans compose by these pullbacks. Their identities and associativity comparisons retain the same tuples of witnesses. Independent service products preserve the corresponding componentwise simulation equations.
\end{proposition}
\begin{proof}\vmode
\leavevmode
\begin{enumerate}[label=(\roman*)]
\item \textbf{Pullback transition.}
On a common request, the replies agree because both images are the reply in $T$; the successors have equal images for the same reason.
Their pair defines the pullback transition.

\item \textbf{Universal property.}
Any compatible pair of strict maps has a unique pairing, and the transition equation makes it strict.
This proves the universal property.

\item \textbf{Associativity.}
Threefold composition retains the same triples, proving associativity and coherence.

\item \textbf{Independent products and interpretation.}
Independent products use the separate request and reply components and their coordinate transitions.
Thus the fixed-interface service category has pullbacks and its codomain projection is a fibration.
Its interpretation proves game equations; it does not identify these pullbacks with pullbacks of arbitrary game strategies.
\end{enumerate}
\end{proof}

\begin{corollary}[Complete observed invariants]\label{cor:complete-observed-service-invariants}
For deterministic total services with current observations $o_S,o_T$ in the same carrier, augment $\beta_S(s)$ by $u\mapsto o_S(\delta_S^*(s,u))$ for $u\in Q^*$, and similarly for $T$. The equality pullback of these augmented behaviours is the greatest relation preserving current observations, replies and successors. It is $\bigcap_{n\geq0}F^n(R_0)$, where $R_0$ equates current observations and $F$ also requires equal next replies and successors in its argument. Its projections are strict.
\end{corollary}
\begin{proof}\vmode
\leavevmode
\begin{enumerate}[label=(\roman*)]
\item \textbf{Finite-depth agreement.}
The empty word tests the current observation and derivatives test each successor.
Finite request induction embeds every preserved relation in the equality pullback and identifies $F^n(R_0)$ with agreement through depth $n$.

\item \textbf{Strict projections.}
The transition calculation in the theorem makes the projections strict.
\end{enumerate}
\end{proof}

\begin{proposition}[Execution, family, and kernel transport]\label{prop:controller-execution-kernel-transport}
Let $\mathbb W_Q=BQ^*$ and $\mathbb W_R=BR^*$ be the free word categories. The service seed
\begin{equation}\label{eq:service-seed}
 \theta_\delta(s)(q)=([r],s')\quad\text{when }\delta(s,q)=(s',r)
\end{equation}
extends uniquely to an execution span $\mathbb W_Q\leftarrow E_\delta\to\mathbb W_R$ whose input leg is a discrete opfibration. Strict state simulations induce maps of these spans and of their family and kernel presentations, preserving pipeline and independent tensor.
\end{proposition}
\begin{proof}\vmode
\leavevmode
\begin{enumerate}[label=(\roman*)]
\item \textbf{Iterated service update.}
Define $\delta^*(s,\eps)=s$, $\omega^*(s,\eps)=\eps$ and extend by one request using $\delta$.

\item \textbf{Execution category.}
Induction gives
\[
 \delta^*(s,uv)=\delta^*(\delta^*(s,u),v),\qquad
 \omega^*(s,uv)=\omega^*(s,u)\omega^*(\delta^*(s,u),v).
\]
The execution category has states as objects and $(s,u):s\to\delta^*(s,u)$ as arrows; input and output are $u$ and $\omega^*(s,u)$.

\item \textbf{Unique lifting and strict maps.}
A request word has one lift from each state, proving the discrete-opfibration assertion and uniqueness.
The strict square extends by the same induction and conversely restricts to words of length one.

\item \textbf{Straightening and kernel.}
More generally a span $P=(A\xleftarrow\pi E\xrightarrow o B)$ with $\pi$ a discrete opfibration is reconstructed from its unique-lifting update $\delta(f,z)$ and output $\omega(f,z)$.
Straightening gives $A\to\operatorname{Fam}(B)$; its kernel is $K_P(a,b)=\coprod_{\pi z=a}B(b,oz)$.

\item \textbf{Composition comparison.}
For a second execution span $Q=(B\xleftarrow\rho F\xrightarrow v C)$, pullback composes the spans.
Write $\delta_Q,\omega_Q$ for its lifting update and output.

\item \textbf{Inverse coend calculation.}
The comparison from $K_Q\star_B K_P$ to the composite kernel is
\[
 [(z,k:b\to oz,w,l:c\to vw)]\longmapsto
 (z,\delta_Q(k,w),\omega_Q(k,w)l).
\]
The inverse uses $b=oz$ and $k=1$; unique lifting and the coend relation prove both inverse equations.

\item \textbf{Coherence and interpretation.}
Reassociation retains the flattened witness and tensor separates components, proving coherence.
This is the execution/kernel reading of the same controllers, not another cryptographic semantic base; compare~\cite{categorical-mealy,mealy-lenses}.
\end{enumerate}
\end{proof}

\begin{proposition}[Public operations and service replacement]\label{prop:primitive-service-comparison}
An operation-preserving public comparison commutes with every finite strict public program. Its reference-state lifting preserves adaptive call--return histories. The resulting comparison of initialized service displays is preserved by sequential connection and independent tensor.
\end{proposition}
\begin{proof}\vmode
\leavevmode
\begin{enumerate}[label=(\roman*)]
\item \textbf{Preservation of public programs.}
Theorem~\ref{thm:frame-isomorphism} and
Proposition~\ref{prop:crypto-public-comparison} give the program
equations, including failure.

\item \textbf{Adaptive service comparison.}
The paired transition of
Construction~\ref{con:evolving-comparison-coalgebra} and command
induction give the adaptive histories;
Theorem~\ref{thm:service-behaviour-canonical-comparison} interprets
them as game equations.

\item \textbf{Contextual replacement.}
The paired histories give a complete trace certificate by Proposition~\ref{prop:service-history-certificate}. Transport through a connection or tensor follows from Theorem~\ref{thm:complete-certificates}. Thus the same public program can be used inside either service network.
\end{enumerate}
\end{proof}
A comparison records correspondence between issued names as the reference tables grow. Each allocation instance supplies an extension map on this correspondence. This is the role of world-indexed name comparisons in~\cite{crypto-stark,benton-hofmann-nigam}.
\subsection{Descent and online simulation}\label{subsec:causal-descent}
Representation independence compares two given controllers. A simulator
has the additional obligation to produce its replies using the
information available from the ideal resource at that time. The next
construction proves an exact criterion for the initial-disclosure
case, without selecting a private representative from each fibre.

\begin{theorem}[Continuous descent of public behaviour]\label{thm:continuous-disclosure}
Let $\ell:X\twoheadrightarrow L$ be the continuous disclosure onto its declared dcpo of public values and let $b:X\to B$ be the continuous behaviour being disclosed indirectly. Write $h^*\Op(Y)=\{h^{-1}U:U\in\Op(Y)\}$ for a map into $Y$. There exists a unique Scott-continuous $j:L\to B$ with $b=j\ell$ exactly when
\begin{equation}\label{eq:continuous-disclosure}
                 b^*\Op(B)\ \subseteq\ \ell^*\Op(L).
\end{equation}
\end{theorem}
\begin{proof}
A factor sends $b^{-1}U$ to the form $\ell^{-1}(j^{-1}U)$. Conversely, the inclusion gives identical membership in every behaviour open for $b(x)$ and $b(y)$ when $\ell(x)=\ell(y)$. The Scott topology separates points: the complement of a principal lower set is Scott open and distinguishes $z$ from $z'$ whenever $z\not\le z'$. For algebraic behaviours the constructive argument compares membership in each principal open of a compact: the two values are suprema of the same compact lower approximations. For products of algebraic reply domains this argument applies coordinate by coordinate. Thus the graph $\{(\ell(x),b(x)):x\in X\}$ is single-valued and total on $L$. It defines the unique $j$ without a choice of private representative. For each $U$ the inclusion supplies $V\in\Op(L)$ with $\ell^{-1}V=b^{-1}U$. Surjectivity gives $j^{-1}U=V$. Open inverse images give Scott continuity: they preserve the specialization order, and an open containing the image of a directed supremum must already contain the image of a member. Separation then identifies the supremum of the images with that image.
\end{proof}

\begin{corollary}[Online replies from an observed disclosure]\label{cor:continuous-disclosure-controller}
For discrete request names $Q$ and reply dcpo $R$ use the behaviour domain $\Omega=R^{Q^+}$ with pointwise order. A continuously initialized deterministic service with continuous transitions has a continuous future-reply map $b:X\to\Omega$. Equation~\eqref{eq:continuous-disclosure} is necessary and sufficient for a continuous reply controller initialized by the disclosed value. Here a word in $Q^+$ records a sequence of calls; its last reply can carry information of any declared reply type.
\end{corollary}
\begin{proof}
For each nonempty request sequence $w$, finite iteration of the transitions and projection of the last reply compute $b(-)(w)$ continuously. Directed joins of behaviours are pointwise, so $b$ is continuous. Evaluation $F\mapsto F(q)$ and the derivative $F\mapsto F_q$, where $F_q(w)=F(qw)$, preserve directed joins. After the ideal initialization notification, start at $j(\ell(x))$ and answer $q$ by $(F_q,F(q))$. Induction on the calls gives the specified replies. Conversely a continuously initialized controller over $L$ has its own continuous future-reply map by the same coordinate argument and hence furnishes $j$.

For algebraic disclosure and reply domains with exported basis-query interfaces, there is also a direct implementation of the readout. For the actual call prefix $w$ and requested compact reply evidence $k$, the witnesses are precisely the compact $t$ with
\[
             t\le l,\qquad k\le j(t)(w).
\]
This is Proposition~\ref{prop:continuous-query-strategy} for the coordinate $l\mapsto j(l)(w)$. Index its private disclosure lanes by call occurrence and reply lane, activate each readout after that call and the initialization notification, and retain its prefix $w$ for subsequent queries on that reply. Different calls have disjoint lanes. Every finite run projects to the controller's requested successful observations by the displayed criterion. Conversely every finite sequence of those successes lifts by inserting the corresponding private query witnesses. The construction uses precisely the read access granted at the disclosure port.
\end{proof}

Fibre constancy suffices at a discrete disclosure. Indeed the set factor from the graph is then continuous, and its preimages supply the opens in~\eqref{eq:continuous-disclosure}. For a nondiscrete disclosure the distinction is observable. Take $X=\Delta\{0,1\}$, let $L=B=\{0<1\}$ and give $\ell$ the underlying identity and $b$ the values $b(0)=1$, $b(1)=0$. Both maps from $X$ are continuous. The only set factor sends the open $\{1\}$ back to $\{0\}$, which is not upper in $L$. The following set-level formulation recovers the discrete controller and its derivative-closed reachable behaviour carrier.

\begin{theorem}[Initial-disclosure descent]\label{thm:initial-disclosure-descent}
For a total controller $\delta$ initialized by $i:X\to S$, let an ideal
initialization disclose $\ell:X\to L$, where $L=\im\ell$. Public
requests begin after the initialization notification. There exists a
deterministic reply controller initialized solely by $\ell(x)$ and
having behaviour $\beta i(x)$ exactly when
\[
 \ell(x)=\ell(y)\quad\Longrightarrow\quad
 \beta i(x)=\beta i(y).
\]
In that case the behaviour image $V=\im\beta\subseteq\Omega$ carries
a canonical simulator controller, initialized by the unique map
$j:L\to V$ with $j\ell=\beta i$.
\end{theorem}
\begin{proof}\vmode
\leavevmode
\begin{enumerate}[label=(\roman*)]
\item \textbf{Derivative-closed behaviour image.}
For $F=\beta(s)$, the derivative $F_q$ is the behaviour of the
successor of $s$ on $q$.
Thus $V$ is closed under every derivative.

\item \textbf{Simulator transition.}
Its transition sends $(F,q)$ to $(F_q,F(q))$ and is independent of a
representative $s$.

\item \textbf{Initialization descent.}
Fibre constancy gives a unique $j$ by image
elimination: $j(l)$ is the unique behaviour $\beta i(x)$ for an $x$
with $\ell(x)=l$.
Its existence uses surjectivity onto $L$, and its
uniqueness uses the displayed implication, so no choice of $x$ is
retained.

\item \textbf{Behaviour equality.}
Induction on requests proves the simulator's equality with
$\beta i(x)$.

\item \textbf{Necessity.}
Conversely a deterministic controller initialized by
the same $l$ has the same behaviour on every request word, proving
fibre constancy.

\item \textbf{Online factorization.}
The factorization is
\[
\begin{tikzcd}[column sep=large]
 X\arrow[r,"i"]\arrow[d,two heads,"\ell"']&S\arrow[d,"\beta"]\\
 L\arrow[r,"j"']&V\arrow[r,hook]&\Omega.
\end{tikzcd}
\]
Lemma~\ref{lem:crypto-stateful} realizes the simulator as a game
strategy.
Its initialization waits for precisely the ideal notification,
so it creates no earlier reply.
This proves the online simulator
equation at the declared initialized interface.
\end{enumerate}
\end{proof}
The construction is semantic: equality in $\Omega$ is not asserted to
be decidable. In the concrete examples, finite controller programs
compute the required replies without storing arbitrary elements of
$\Omega$. Later disclosure is handled by actual later ideal calls and
their replies. A factorization of completed transcripts alone does not
construct such a controller: a real first reply containing a private
bit cannot be simulated from an ideal resource that first discloses
that bit only after a later request. Both complete transcripts could
contain the bit while their first-prefix behaviours differ.

\paragraph{Descent of public operations.}\label{subsec:operation-descent}
\begin{theorem}[Exact classifier for public operations]\label{thm:operation-descent}
Let $p:V\twoheadrightarrow P$ be the image projection to the public representation and let $k:V^n\to V+\{\mathsf{fail}\}$ be a specified operation. Extend $p$ by the identity on failure to $p_\bot$. There is a unique public operation $\bar k:P^n\to P+\{\mathsf{fail}\}$ satisfying
\[
 p_\bot k=\bar k\,p^n
\]
exactly when
\begin{equation}\label{eq:operation-fibre-criterion}
 p^n(v)=p^n(v')\quad\Longrightarrow\quad
 p_\bot k(v)=p_\bot k(v').
\end{equation}
Thus admissible public operations are exactly those preserving the kernel relation of the public projection, including their failure outcomes. The same construction applies to heterogeneous input and output carriers.
\end{theorem}
\begin{proof}\vmode
\leavevmode
\begin{enumerate}[label=(\roman*)]
\item \textbf{Necessity of fibre constancy.}
The factorization implies the criterion.

\item \textbf{Definition by image elimination.}
Conversely each public input tuple has a private preimage, since $p$ is surjective; its $p_\bot k$ value is independent of the preimage by the criterion.
Image elimination therefore defines $\bar k$, without retaining or choosing any private representative.

\item \textbf{Uniqueness and computation.}
Surjectivity gives uniqueness and the displayed equation.

\item \textbf{Nullary and heterogeneous operations.}
For nullary operations use the single empty tuple.
The componentwise argument proves the heterogeneous case.
\end{enumerate}
\end{proof}

\begin{corollary}[Continuous public operations, including failure]\label{cor:open-operation-descent}
For a finite tuple of continuous public projections $p_i:V_i\twoheadrightarrow P_i$, put $p_{\rm in}=\prod_i p_i$. Let $k$ be a continuous private operation with result sort $V_o+\{\mathsf{fail}\}$, the disjoint sum with its tagged failure. Put $b=(p_o+\mathrm{id})k$. A unique continuous public operation $\bar k$ with $b=\bar k p_{\rm in}$ exists exactly when
\[
 b^*\Op(P_o+\{\mathsf{fail}\})\ \subseteq\
 p_{\rm in}^*\Op\bigl(\textstyle\prod_i P_i\bigr).
\]
\end{corollary}
\begin{proof}
The finite product projection is surjective, including the nullary case. Apply Theorem~\ref{thm:continuous-disclosure} to this map and the projected outcome. In the sum, both the failure summand and every open of the success summand are open. Their pullbacks account for the actual failure reply as well as successful observations. A pending command has no reply occurrence. On discrete carriers all subsets are open, so the condition reduces to~\eqref{eq:operation-fibre-criterion}.
\end{proof}

A total verdict in the discrete Boolean sort requires both outcome fibres to be Scott open. Even success of exact equality can fail this test for evolving payloads: in $D=\{0<1\}$ the equality diagonal contains $(0,0)$ and omits the larger $(0,1)$. Exact equality on issued identifiers, bit tags and the specified discrete symbolic sorts therefore retains its original call--return rule. A request for a basis observation of a growing payload receives positive evidence as in Construction~\ref{con:positive-observation-port}.

\begin{corollary}[Reference-table realization of descent]\label{cor:reference-descent-realization}
For a reference service whose declared operations satisfy these descent equations, project each table entry by $p$ and retain the issued identifiers, counters, ownership and public command labels. Corresponding initialized services have a strict comparison and therefore the same call--return display. If initialization is supplied by an ideal public notification, the projected controller is a simulator using precisely that notification and the declared public operations.
\end{corollary}
\begin{proof}\vmode
\leavevmode
\begin{enumerate}[label=(\roman*)]
\item \textbf{Lookup correspondence.}
Lookup has the same success or failure on both tables because their issued domains, types and owners coincide.

\item \textbf{Operation descent.}
For a valid lookup, Theorem~\ref{thm:operation-descent} gives the same public result or failure.

\item \textbf{Issuance and public tests.}
Successful allocation installs the projected value under the same next identifier and increments the same counter.
Public tests use their specified output projection and hence have the same replies.

\item \textbf{Command induction.}
Induction on commands gives the strict transition square; Theorem~\ref{thm:service-behaviour-canonical-comparison} supplies the game comparison.

\item \textbf{Available ideal initialization.}
Initialization reads only its actual ideal notification, so the constructed simulator has no earlier private access.
\end{enumerate}
\end{proof}

There is a corresponding readout at a reference table. Issue a handle under the service's original freshness rule. At an authorized observation request attach Construction~\ref{con:positive-observation-port} to the public result, using separate lanes per handle and request. On related tables the projected result agrees, so exactly the same requested acknowledgements can occur by the same compact-membership test. Relate the extended states by the original table relation together with equal issued and answered public lanes. A request extends this relation on both sides, and a positive reply extends it precisely when its basis test succeeds. Induction on finite extensions gives the two trace inclusions, including prefixes with unanswered requests. This implements observations at the public projection without inspecting a private representative.

For proof services, let $V$ contain checked proof records and let $p$ extract their validated statements. Conjunction, refresh and the declared hiding operations are public exactly when their statement operations satisfy~\eqref{eq:operation-fibre-criterion}. The proof-erasure construction in Appendix~\ref{app:cryptographic-examples} implements this reference-table comparison, while private audit records remain evidence over the real execution.

The criterion also detects an invalid disclosure. Two valid records for the same statement can have hidden witness bits $0$ and $1$. An operation returning that bit would give different replies on one $p$-fibre and therefore fail the criterion. The descent calculation determines which proposed public operations respect the statement interface; it does not make an inaccessible private function into an available command.

\begin{proposition}[Descent of observation families]\label{prop:observation-family-descent}
For a map $C:X\to Y$ and a family $V:X\to\Pow(O)$, constancy
$C(x)=C(x')\Rightarrow V(x)=V(x')$ is equivalent to unique descent
along $q:X\twoheadrightarrow\im C$. Writing the descended relation
as $S\subseteq\im C\times O$, the equation is
$V=(q\times1_O)^{-1}S$.
\end{proposition}
\begin{proof}\vmode
\leavevmode
\begin{enumerate}[label=(\roman*)]
\item \textbf{Descent and uniqueness.}
Set $S(y)=V(x)$ for any preimage of $y$.
Fibre constancy gives
independence and uniqueness, without selecting representatives.

\item \textbf{Converse and scope.}
Conversely the pullback equation gives fibre constancy.
This is
subobject descent in $\Set$; an online resource realization additionally
supplies the actual initialization and subsequent command transitions.
\end{enumerate}
\end{proof}

\begin{theorem}[From a constructed comparison to secure refinement]\label{thm:categorical-crypto-bridge}
For a correct conversion $f:r\to s$ and coalition $J$, let $\sigma$ be a constructed strategy using the actual ideal ports. A strict initialized comparison of the real and simulated services, a finite causal certificate of their networks, or a complete trace certificate proves
\[
 (f_H\otimes1)r\eqsc(1\otimes\sigma)s.
\]
Together with correctness this gives the mutual slice factorizations of Theorem~\ref{thm:correctness-security-mutual-factorization}, and hence equality of attainable attacker-view images. The common trace image also gives the public-frame comparison at the exposed boundary, as in Proposition~\ref{prop:comparison-public-frames}. These witnessed refinements compose and tensor.
\end{theorem}
\begin{proof}\vmode
\leavevmode
\begin{enumerate}[label=(\roman*)]
\item \textbf{Comparison implies realization.}
A strict service comparison identifies the fixed call--return displays, including pending requests, by Theorem~\ref{thm:service-behaviour-canonical-comparison}.
A causal certificate gives contextual equality by Lemma~\ref{lem:crypto-causal-congruence}; a trace certificate does so by Theorem~\ref{thm:complete-certificates}.
Theorem~\ref{thm:observational-simulation} identifies the resulting comparison with the simulator judgement for this $\sigma$.

\item \textbf{Public information transport.}
Proposition~\ref{prop:comparison-public-frames} gives the public-frame comparison on their configuration pullback.

\item \textbf{Resource factorizations.}
Transposing the simulator and correctness equations gives the two slice arrows.
Theorems~\ref{thm:correctness-security-mutual-factorization} and~\ref{thm:secure-conversion-composition} prove security and composition; Corollary~\ref{cor:security-images} gives the image equality.

\item \textbf{Uniform attacker simulation.}
Every attacker $a$ is simulated by the existing composite $a\sigma$.
\end{enumerate}
\end{proof}

\subsection{Process presentation of reference services and simulators}\label{subsec:cryptographic-process-rules}
The reference-controller constructions now supply process implementations of operations and simulators. The rules below specify the command results and state transitions; their interpretation transports the service comparisons and resource judgements already proved.

\begin{definition}[Strict destructor rules]\label{def:calculus-destructor-rules}
The actual operation carriers extend the value rules of Definition~\ref{def:value-formation-constructors}. A constructed partial operation $d:V_{\vec\tau}\rightharpoonup V_\sigma$ with terminating definedness test has its value-or-failure evaluator $\widehat d:V_{\vec\tau}\to V_\sigma+\{\crFail\}$. Its process eliminator is
\[
 \frac{\Gamma\vdash\vec v:\vec\tau\quad
       \Gamma,x:\sigma;\Xi\vdash P\ \mathsf{proc}\quad
       \Gamma;\Xi\vdash Q\ \mathsf{proc}}
 {\Gamma;\Xi\vdash\mathsf{let}\ x=d(\vec v)\ \mathsf{in}\ P\ \mathsf{else}\ Q\ \mathsf{proc}}.
\]
Its computation is the selected success transition to $P[w/x]$ when $\widehat d(\vec v)=\mathsf{inl}(w)$ and the selected failure transition to $Q$ when $\widehat d(\vec v)=\mathsf{inr}(\crFail)$. The operands are evaluated at the current valuation before the branch is selected. Both branches retain the evaluation witness of the applied operational instance.
\end{definition}
\begin{proposition}[Strict evaluation and partial-operation elimination]\label{prop:strict-evaluation-partial-elimination}
These destructor rules preserve scope and sorts and have precisely the constructed evaluator's success and failure computations.
\end{proposition}
\begin{proof}\vmode
\leavevmode
\begin{enumerate}[label=(\roman*)]
\item \textbf{Typed strict elimination.}
A successful result has sort $\sigma$, so
Theorem~\ref{thm:substitution-indexed-interpretation} types $P[w/x]$;
failure uses the typed $Q$.

\item \textbf{Evaluation and substitution.}
Induction on the strict evaluation tree
preserves the selected branch and operands.

\item \textbf{Branch order and effectiveness.}
Evaluation precedes
either continuation even when $x$ is unused.
Sequential effective
evaluation uses Proposition~\ref{prop:calculus-effective}; its
undefined lift value remains distinct from a returned failure tag.
\end{enumerate}
\end{proof}

\paragraph{Dependent command results.}\label{par:dependent-command-results}
For command $a$, let $E_a$ be its tagged component of the existing committed-edge object. Its result family is the source fibre, with its derived target and reply projections:
\begin{equation}\label{eq:dependent-command-result-family}
 \mathsf{Result}_a(s)=(E_a)_s,\qquad
 e\longmapsto(t(e),\mathsf{reply}(e),e).
\end{equation}
This displayed tuple is the graph of those projections, not additional result data. Postconditions are maps or inclusions on this fibre. The client binds the runtime reply; the state and audit stay at their private interfaces.

\begin{lemma}[Evaluation and origin in a finite operation audit]\label{lem:generic-operation-audit}
An actual finite operation audit has a sorted input family $R_h$, node family $N_h$, and an operation tag and ordered predecessor tuple at every noninput node. There is a unique unfolding
\[
u_h:N_h\longrightarrow\mathsf{Term}_{\Sigma}(R_h),\qquad
u_h(r)=r,\quad u_h(n)=\sigma_n(u_h(n_1),\ldots,u_h(n_k)).
\]
For an interpretation of the inputs and operations, composition with the term-algebra fold is the unique audit evaluation. Unfolding and evaluation are natural under maps preserving inputs, tags and operand pointers, and under commuting maps of operation carriers. Strict partial operations use their value-or-failure carriers. A selected unary ancestry $\pi:P\to R\amalg P$ has the unique root map $r:P\to R$ with $r=[1_R,r]\pi$. These constructions are unchanged by independent exchanges.
\end{lemma}
\begin{proof}\vmode
\leavevmode
\begin{enumerate}[label=(\roman*)]
\item \textbf{Audit fold and uniqueness.}
Recursion in the finite creation order assigns each input its
variable and each node its operation applied to the already
assigned predecessors.
The same induction proves uniqueness and,
after composition with the term fold, the evaluation equations.

\item \textbf{Naturality of evaluation.}
For an audit map $i$, both $u_{h'}i_N$ and
$\mathsf{Term}_\Sigma(i_R)u_h$ satisfy those equations; uniqueness
proves naturality.
It likewise proves commutation with maps of
operation carriers and independence of linearization.

\item \textbf{Unary ancestry.}
For unary ancestry, decorate roots by themselves and each other
node by its parent's decoration.

\item \textbf{Termination and transport.}
Parent chains strictly decrease
creation order, so this recursion terminates and uniquely satisfies
$r=[1_R,r]\pi$.
Pointer-preserving maps and independent exchanges
preserve its equations.
\end{enumerate}
\end{proof}
The term object here is the free expression carrier of the audit's actual operation signature. It supplies an evaluation construction without exposing private operands at a public port. Appendix~\ref{app:symbolic-implementation} specializes it to the concrete symbolic evaluator.

\begin{definition}[Service-call formation and elimination]\label{def:service-call-formation-elimination}
For the constructed service $\delta:S\times Q\to S\times R$, its request and reply carriers are value types $Q,R$ and its private state carrier is $S$. The \textbf{client call} is the typed abbreviation
\[
 \mathsf{call}_\delta(q;x.P):=\mathsf{req}!q.\mathsf{reply}?(x).P,
\]
with $\mathsf{req}:\mathsf{Chan}(Q)$ and $\mathsf{reply}:\mathsf{Chan}(R)$. Its introduction and continuation elimination are the two prefix rules: from $\Gamma\vdash q:Q$ and $\Gamma,x:R;\Xi\vdash P\ \mathsf{proc}$ they derive $\Gamma;\Xi\vdash\mathsf{call}_\delta(q;x.P)\ \mathsf{proc}$. The selected service update is $\delta(s,q)=(s',r)$; receiving its reply computes $P[r/x]$ and retains $s'$ as the next private service state. The request, intervening evaluator commands, and reply keep their separate occurrences in the service history.

A strict simulation $a:S\to S'$ transports that update by $(a\times1_R)\delta=\delta'(a\times1_Q)$ and hence transports its request/reply witness. For a possibly nonreturning computation the semantic map has target $\mathsf{Lift}(S\times R)$, or $\mathsf{Lift}(S\times(R+\{\crFail\}))$ when failure is observable. These are respectively undefined computation and returned failure; the client continuation is installed only on a returned reply.
\end{definition}

\paragraph{Process form of the service.}
\begin{definition}[Recursive service process]\label{def:formal-recursive-service-process}
Write $\mathsf{Serve}(s)$ for the state-indexed equations
\[
 \begin{aligned}
 \mathsf{Serve}(s)&=\mathsf{req}?(q).\\
 &\quad\mathsf{let}\ (s',r)=\delta(s,q)\ \mathsf{in}\;
 \mathsf{reply}!r.\mathsf{Serve}(s').
 \end{aligned}
\]
Their meaning is the union of the finite unfoldings just constructed, starting at $s_0$. The alternating interface orders the next request after the preceding reply. Expanding the service state and evaluator yields the actual command rules, including failure. Lemma~\ref{lem:crypto-stateful} supplies the ordinary strategy, and Theorem~\ref{thm:pc-representation} supplies its occurrence declaration. Thus a stateful service has both a named recursive presentation and an exact typed occurrence presentation.
\end{definition}


\paragraph{Simulator rules.}\label{subsec:sim-calculus}
Use the existing judgement \(\Gamma;\Delta\vdash P:A\). Abbreviate a single-port term by \(P:A\to B\), write \(QP\) for the already typed cut, and use \(\langle P,Q\rangle\) for tensor. A resource is a state term \(\vdash R:A_P\); a protocol is a tuple \(\Pi=(P_i:A_i\to B_i)\). Its correctness judgement has premise
\[
\Pi R\eqd S.
\]
Here \(\eqd\) is equality of interpreted displays, kept distinct from structural congruence and operational reduction.

The simulator judgement is \(\Pi:R\Rightarrow_J^{\Sigma}S\), whose witness is an existing process term \(\Sigma:B_J\to A_J\). For exact display comparison its introduction rule is
\[
\frac{\Pi R\eqd S\qquad \Sigma:B_J\to A_J\qquad
 (\Pi_H\otimes1)R\eqd(1\otimes\Sigma)S}
 {\Pi:R\Rightarrow_J^{\Sigma}S}.
\tag{Sim-I}
\]
Its derived elimination and composition rules are
\[
\frac{\Pi:R\Rightarrow_J^{\Sigma}S\qquad a:A_J\to E}
 {(\Pi_H\otimes a)R\eqd(1\otimes a\Sigma)S},
\tag{Sim-E}
\]
\[
\frac{\Pi:R\Rightarrow_J^{\Sigma}S\qquad
       \Lambda:S\Rightarrow_J^{\Theta}T}
 {\Lambda\Pi:R\Rightarrow_J^{\Sigma\Theta}T},
\qquad
\frac{\Pi:R\Rightarrow_J^{\Sigma}S\qquad
       \Lambda:R'\Rightarrow_J^{\Theta}S'}
 {\Pi\otimes\Lambda:R\otimes R'\Rightarrow_J^{\Sigma\otimes\Theta}S\otimes S'}.
\tag{Sim-Cut/Tensor}
\]
Replacing every $\eqd$ in these rules by $\eqsc$ gives the observational rules. Their introduction premise is equivalently a trace certificate by Theorem~\ref{thm:observational-simulation}; in that version the conclusion records the simulator strategy $\llbracket\Sigma\rrbracket\Vdash_H\llbracket(\Pi_H\otimes1)R\rrbracket\preccurlyeq\llbracket S\rrbracket$. Exact equations imply this version, while retaining stronger occurrence information. These abbreviations include the inherited grouping maps. Soundness follows from Theorem~\ref{thm:correctness-security-mutual-factorization} and Theorem~\ref{thm:secure-conversion-composition}; substitution is inherited from the existing typed cut/substitution theorem. No new transition rule is required to execute a simulator: it is a process built from the existing constructors.

\paragraph{Readback and preservation.}
A simulator witness has an occurrence-process presentation by
the existing readback theorem. The same cut $a\Sigma$ then presents
its composite with each attacker.

\begin{proposition}[Simulator presentation and operational soundness]\label{prop:readback}
Every semantic simulator has a set-indexed occurrence-process presentation. A finite simulator display has a finite presentation. The protocol and simulator rules are sound for exact and Scott-observational simulation, and their observation conclusions are sound for the finite execution judgement of Theorem~\ref{thm:evidence}.
\end{proposition}
\begin{proof}\vmode
\leavevmode
\begin{enumerate}[label=(\roman*)]
\item \textbf{Simulator readback.}
Readback represents the simulator's ordinary display, with the arrow obtained by the existing port transposition.

\item \textbf{Interpretation of connection and tensor.}
Cut and tensor interpret as the corresponding strategy operations.
Hence the semantic simulator equation is the interpreted equation of the represented processes, and \textup{(Sim-E)} supplies each attacker equation.

\item \textbf{Agreement of tests.}
Theorem~\ref{thm:contextual-congruence-calculus} transports observational equality and Theorem~\ref{thm:evidence} supplies the finite-execution interpretation.

\item \textbf{Finite and unrestricted evidence.}
For a supplied finite exact display certificate, Proposition~\ref{prop:pc-normalization} checks its equation.
The unrestricted assertion layer supplies no algorithm for arbitrary observational equivalence.
\end{enumerate}
\end{proof}

The process construction distinguishes finite terms from set-indexed declarations: the readback theorem covers all ordinary strategies through the latter, while finite normalization covers supplied finite exact certificates. Secret inputs remain at their honest ports in the simulator equation. Composition with any subsequent input filling preserves that same witness.

\begin{proposition}[Process preservation of secure refinement]\label{prop:generic-process-preservation}
The service and strict-operation rules preserve runtime substitution, strict service squares, comparison spans and their induced public-frame maps. Their named realizations satisfy the same simulator equations as their game interpretations. With the constructed correctness converters they therefore give the same mutual slice factorizations and attainable-view-image equalities.
\end{proposition}
\begin{proof}\vmode
\leavevmode
\begin{enumerate}[label=(\roman*)]
\item \textbf{Substitution and strict evaluation.}
Substitution evaluates the same request operands and installs the returned value in the typed continuation.
Strict evaluation precedes either continuation, retaining failure even when its result variable is unused.

\item \textbf{Compositional interpretation.}
Theorem~\ref{thm:substitution-indexed-interpretation} gives typing; induction on finite service unfoldings transports the transition square, and applying both projections gives the span case.

\item \textbf{Reference and comparison transport.}
The reference comparison keeps identifiers and test results.
Connection and tensor transport the resulting causal certificates through secured interaction, so Theorem~\ref{thm:categorical-crypto-bridge} applies.

\item \textbf{Occurrence representation.}
Arbitrary semantic witnesses have the occurrence readback of Theorem~\ref{thm:pc-representation}; concrete finite presentations retain their stated effective implementations.
\end{enumerate}
\end{proof}
Appendix~\ref{app:cryptographic-examples} instantiates these rules. Recipes are finite expressions in the existing value language; Scott opens remain semantic assertions rather than new effective tests.


Public operations and controller descent supply the strategy at the ideal ports, and correctness gives the reverse resource factorization. The protocol history retains additional evidence about how a reply was produced. We now give this evidence its own history and view indices.

\section{Protocol evidence and policy}\label{sec:protocol-evidence-policy}
Section~\ref{subsec:public-observations} constructs the assertions that a coalition can observe through its boundary. A retained history can also carry a witness for an assertion, such as a signing occurrence in the causal past of a release. Let $U$ be a specified execution presentation or its controlled restriction and put $H=\Conf_f(U)$. We index evidence families by these histories and their decision views, and construct the maps that transport evidence through disclosure, execution, and policy changes.

\subsection{Evidence spans, quantification, and response games}\label{subsec:admin-equipment}
An assertion can carry a family of witnesses indexed by a history or view. Spans organize changes of those indices, and their adjunctions provide existential and universal evidence transport. Response games give an exact game interpretation of the resulting discrete interfaces.

\begin{definition}[Evidence slices and quantified families]\label{def:formal-evidence-slices-and-quantified-families}
Work in $\E=\Set$. For an actual index set $X$, write $\mathsf{Fam}(X)=\Set/X$. For $f:X\to Y$, pullback is $f^*$, dependent sum $\Sigma_f$ composes a family with $f$, and dependent product has fibre $(\Pi_fA)_y=\prod_{f(x)=y}A_x$. On subsets their adjoints are $\exists_f$ and $\forall_f$. The sum of witnesses and its image predicate remain different constructions.
\end{definition}
These are the usual family constructions~\cite{doct-jacobs}; their required universal properties are proved below at the actual index sets. They record evidence about the witnessed protocol operations and do not add operations to a cryptographic service.

\begin{definition}[Span squares and an equipment]\label{def:admin-equipment}
A square from $R:X\rightsquigarrow Y$ to $R':X'\rightsquigarrow Y'$ with vertical boundaries $f:X\to X'$ and $g:Y\to Y'$ is a carrier map $u:R\to R'$ satisfying $s'u=fs$ and $t'u=gt$. Horizontal composition uses the maps induced on pullbacks; vertical composition is composition of the three maps. An \textbf{equipment} is such a pseudo double category with cartesian restrictions of horizontal arrows along both vertical boundaries: every square into a restricted arrow factors uniquely through its defining square, including after further changes of endpoints. The following construction is the span instance of framed bicategories \cite{shulman}; we give its universal maps explicitly.

The defining square is equivalently the two endpoint triangles in
\[
\begin{tikzcd}[column sep=large]
X \arrow[d,"{f}"'] &R \arrow[l,"{s}"'] \arrow[r,"{t}"] \arrow[d,"{u}"] &Y \arrow[d,"{g}"]\\
X' &R' \arrow[l,"{s'}"'] \arrow[r,"{t'}"] &Y'.
\end{tikzcd}
\]
Horizontal composition sends $(r,q)$ to $(u(r),v(q))$ in the corresponding pullback; its existence follows from the shared endpoint equation.

\end{definition}

\begin{theorem}[Evidence-span equipment and response games]\label{thm:admin-equipment}
Spans and span maps form an equipment. Set-valued spans are interpreted by the response strategies of Definition~\ref{def:responsegame}. Its restriction of $R':X'\rightsquigarrow Y'$ along $f,g$ has carrier
\[
                   X\times_{X'}R'\times_{Y'}Y.
\]
The graph and reverse graph of $f:X\to Y$ are its companion and conjoint,
\[
 f_!=\bigl(X\xleftarrow{1}X\xrightarrow fY\bigr),\qquad
 f^\dagger=\bigl(Y\xleftarrow fX\xrightarrow{1}X\bigr),
 \qquad f_!\dashv f^\dagger.
\]
The construction holds in every slice $\E/H$ as well.
\end{theorem}
\begin{proof}\vmode
\leavevmode
\begin{enumerate}[label=(\roman*)]
\item \textbf{Span and response composition.}
Composition of spans is pullback and maps of spans act on their carrier witnesses.
For set-valued spans, Theorems~\ref{thm:game-return-matrix} and~\ref{thm:responsecomposition} identify this operation with the discrete kernel and initialized response-strategy calculations.

\item \textbf{Discrete execution interpretation.}
Regarding $S$ as the discrete category on $S$, a transducer between such interfaces is exactly a span: its only actions and outputs are identities.
Its pipeline is the same pullback.
The equipment structure is therefore the span construction of \cite[Examples~4.15]{shulman}.

\item \textbf{Equipment universal property.}
Explicitly, a square over $f,g$ is precisely a map into $X\times_{X'}R'\times_{Y'}Y$ preserving its three projections.
This pullback universal property gives its cartesian restriction and coherence under further endpoint changes.

\item \textbf{Companion and conjoint.}
Restricting an identity span at either end gives the displayed graph spans.

\item \textbf{Adjunction and internalization.}
The adjunction unit is the diagonal $X\to X\times_YX$, and the counit has carrier map $f:X\to Y$; their triangles are the projection--diagonal identities.
The same pullbacks and universal maps exist in every slice, proving the last assertion.
\end{enumerate}
\end{proof}

\begin{proposition}[Action on evidence and base change]\label{prop:evidence-action-base-change}
A span $R=(X\xleftarrow{s}R\xrightarrow{t}Y)$ acts on outgoing evidence by
\[
          R^\sharp=\Sigma_s t^*:\mathsf{Fam}(Y)\to\mathsf{Fam}(X).
\]
There are coherent isomorphisms $(R;Q)^\sharp\cong R^\sharp Q^\sharp$ and $1^\sharp\cong1$. For graph spans, $(f_!)^\sharp=f^*$ and $(f^\dagger)^\sharp=\Sigma_f$, recovering $\Sigma_f\dashv f^*$. Reindexing additionally has the right adjoint $\Pi_f$. Pullback squares give the usual dependent-sum and dependent-product base-change isomorphisms.
\end{proposition}
\begin{proof}\vmode
\leavevmode
\begin{enumerate}[label=(\roman*)]
\item \textbf{Fibrewise action.}
The fibre of $R^\sharp Z$ at $x$ consists of $(r,z)$ with $s(r)=x$ and $z\in Z(t(r))$.

\item \textbf{Composition and coherence.}
Applying $Q^\sharp$ first gives triples $(r,q,z)$ with matching middle endpoint; these are exactly the fibres for the pullback carrier of $R;Q$.
The identifications for longer composites retain the same ordered witness tuple, proving coherence.

\item \textbf{Graph specialization.}
The graph calculations follow by substituting their anchors.

\item \textbf{Evidence adjunctions.}
Transposing a family map $\Sigma_fU\to V$ amounts to a map $U_x\to V_{f(x)}$ for every $x$, proving the first adjunction; transposing a map $f^*V\to U$ gives the tuple of its components over each fibre, proving the second.

\item \textbf{Base change.}
For a pullback square, the corresponding index fibres are canonically the same; dependent sums and products over them give base change.
The internal proof is the identical slice universal-property calculation.
These are adjunctions on evidence categories; $\Pi_f$ is not being declared another horizontal adjoint of the graph transducer.

\item \textbf{Explicit pullback correspondence.}
Explicitly, for the pullback below, a point over $x'$ is the same pair on either route:
\[
\begin{tikzcd}[column sep=large]
X' \arrow[r,"{g'}"] \arrow[d,"{f'}"'] &X \arrow[d,"{f}"]\\
Y' \arrow[r,"{g}"'] &Y
\end{tikzcd}
\qquad
 g^*\Sigma_f\cong\Sigma_{f'}(g')^*,\qquad
 g^*\Pi_f\cong\Pi_{f'}(g')^*.
\]
The comparisons send a witness, respectively a compatible tuple over the fibre, to itself under that fibre bijection.
Iterated pullbacks retain the same tuple and hence prove their coherence.
\end{enumerate}
\end{proof}

The generic mode transpositions and parameter-family closed structure are constructed in Appendix~\ref{subsec:evidence-mode-calculus} and Appendix~\ref{subsec:parameter-families-closed-resources}. Here they are applied to the actual attained views and retained protocol histories.

\subsection{Current decision views and their evidence families}\label{subsec:admin-observation}
A decision view identifies histories that present the same current information. We now compute how evidence moves between histories and attained views, and when it descends to information available at the decision point.

\begin{definition}[Histories and attained observations]\label{def:admin-observation}
For current decisions, let $H$ be the set of finite full histories of a specified process or controlled subprocess. Its communication, service, or administrative interface determines which fields are visible. For participant $q$, compose its local causal-view map with its declared observation to obtain $o_q:H\to O_q$. Local observations may include received values, successful local tests, and local records, while omitting packet-source names and other participants' private data. The projection $o:H\to O$ retains the actual interface records exposed by this application. Replace each codomain by its image when referring to attained views. Predicates on an object $X$ are its subobjects, written $\operatorname{Sub}(X)$; in sets these are subsets.
\end{definition}

\begin{definition}[Observation and action operators]\label{def:observation-action-operators}
For the attained-observation map $o:H\to O$, define endomaps of $\operatorname{Sub}(H)$ by $\mathsf K_o=o^*\forall_o$ and $\mathsf{Poss}_o=o^*\exists_o$. A predicate is \textbf{observation-saturated} when it has the form $o^*Q$ for $Q\in\operatorname{Sub}(O)$. For an action-witness span $H\xleftarrow{s}E_a\xrightarrow{t}H$, define $\langle a\rangle=\exists_s t^*$ and $[a]=\forall_s t^*$, also as endomaps of $\operatorname{Sub}(H)$. Their corresponding evidence operations are $\Sigma_s t^*$ and $\Pi_s t^*:\E/H\to\E/H$.
\end{definition}

\begin{proposition}[Derived epistemic and action modalities]\label{prop:admin-modalities}
The observation map determines
\[
 \mathsf K_o(P)=o^*\forall_oP,
 \qquad \mathsf{Poss}_o(P)=o^*\exists_oP.
\]
The first is the greatest observation-saturated predicate contained in $P$; the second is the least observation-saturated predicate containing $P$. Both are idempotent, and their fixed predicates are exactly the inverse images of predicates on attained views. For an action-witness span $H\xleftarrow{s}E_a\xrightarrow{t}H$, the derived dynamic modalities are
\[
 \langle a\rangle P=\exists_s t^*P,\qquad
 [a]P=\forall_s t^*P.
\]
In sets their meanings are
\begin{align*}
 \mathsf K_o(P)(h)&\iff \forall h'\,(o(h')=o(h)\Rightarrow P(h')),\\
 \mathsf{Poss}_o(P)(h)&\iff \exists h'\,(o(h')=o(h)\ \text{and}\ P(h')),\\
 [a]P(h)&\iff\forall e\in E_a\,(s(e)=h\Rightarrow P(t(e))).
\end{align*}
\end{proposition}
\begin{proof}\vmode
\leavevmode
\begin{enumerate}[label=(\roman*)]
\item \textbf{Slice adjunctions.}
The slice adjunctions follow from the universal properties of dependent sum, pullback, and dependent product.

\item \textbf{Predicate operations.}
Factoring the dependent-sum map through its image gives the existential adjoint on subobjects.
Their internal element descriptions give the three formulas and the analogous existential formula for $\langle a\rangle$.

\item \textbf{Greatest public subpredicate.}
The counit of $o^*\dashv\forall_o$ gives $\mathsf K_o(P)\subseteq P$.
If $o^*Q\subseteq P$, adjoint transposition gives $Q\subseteq\forall_oP$, hence $o^*Q\subseteq\mathsf K_o(P)$.
Since $\mathsf K_o(P)$ is itself an inverse image, this proves the greatest-property and its idempotence.

\item \textbf{Least public superpredicate.}
The other adjunction gives the least-property and idempotence of $\mathsf{Poss}_o$.

\item \textbf{Fixed predicates and scope.}
Applying either property to an inverse image proves the assertion about fixed predicates.
These arguments use order and adjunction, not Boolean complementation, so they interpret internally.
\end{enumerate}
\end{proof}

These subset operators describe current decision alternatives; they are not asserted to preserve Scott openness. Historical public approximation remains $\Box_J$ of Proposition~\ref{prop:interior}. In particular current authorization can change after revocation, whereas a witnessed past validation is persistent.

The evidence-family operation $o^*\Pi_o$ has actual tuples of witnesses over an observation fibre. The preceding idempotence statement concerns predicates, whose fibres are subterminal; it does not identify iterated products of arbitrary evidence families. Nor does an arbitrary physical extension preserve a knowledge predicate: its observation fibres and current facts can change.

\begin{theorem}[Decision descent and compositional disclosure]\label{thm:admin-descent}
For $o:H\twoheadrightarrow O$ onto attained views and an outcome map $d:H\to D$, the following are equivalent:
\[
 d\pi_1=d\pi_2\text{ on }H\times_OH;
 \qquad d=\bar d\,o\text{ for a unique }\bar d:O\to D.
\]
Consequently $D^O\to D^H\rightrightarrows D^{H\times_OH}$ is an equalizer. View-indexed families of arrows are closed under pointwise identities and composition, and under pointwise tensor in a monoidal target category. For a decidable verdict $d:H\to\{0,1\}$ and $P=d^{-1}(1)$, descent is also equivalent to $\mathsf K_o(P)=P$ and to $\mathsf{Poss}_o(P)=P$.
\end{theorem}
\begin{proof}\vmode
\leavevmode
\begin{enumerate}[label=(\roman*)]
\item \textbf{Kernel-pair descent.}
An effective epimorphism is the coequalizer of its kernel pair.
Its universal property gives precisely the displayed equivalence and uniqueness.

\item \textbf{Parameterized descent.}
The same statement with a parameter object gives the equalizer of exponentials, since its product with $o$ is again an effective epimorphism.

\item \textbf{Closure under operations.}
If $d_h=\bar d_{o(h)}$ and $e_h=\bar e_{o(h)}$ are composable, their composite is $(\bar e\bar d)_{o(h)}$; identities are constant and a tensor is $(\bar d\otimes\bar f)_{o(h)}$.
These formulas prove the closure statements.

\item \textbf{Decidable verdicts.}
A decidable verdict is constant on each fibre exactly when its true predicate is a union of whole fibres.

\item \textbf{Predicate characterization.}
Proposition~\ref{prop:admin-modalities} characterizes those predicates by either fixed-point equation.
The kernel-pair assertion itself applies internally to every outcome object and does not manufacture a decision map from an arbitrary predicate.
\end{enumerate}
\end{proof}

\begin{theorem}[Observation operators and descent]\label{thm:observation-operators-descent}
On evidence over \(H\),
\[
 T_o=o^*\Sigma_o,\qquad K_o=o^*\Pi_o
\]
are respectively a monad and a comonad, and \(T_o\dashv K_o\).
There are equivalences
\[
 \mathcal E/O\simeq\operatorname{Alg}(T_o)
             \simeq\operatorname{Coalg}(K_o).
\]
On predicates, \(o^*\exists_o\) and \(o^*\forall_o\) are idempotent
closure and interior operators. The evidence operators are generally
not idempotent.
\end{theorem}
\begin{proof}\vmode
\leavevmode
\begin{enumerate}[label=(\roman*)]
\item \textbf{Observation monad.}
Internally, \((T_oA)_h\) has pairs \((h',a)\) with \(o(h')=o(h)\).
Its unit uses \(h'=h\), and its multiplication deletes the intermediate
history.

\item \textbf{Observation comonad.}
The object \((K_oA)_h\) consists of sections of \(A\) over the
same observation fibre.
Its counit evaluates at \(h\); its
comultiplication reindexes the same section from every history in the
fibre.

\item \textbf{Structural equations.}
Inserting or deleting indices proves the structure equations.

\item \textbf{Composite adjunction.}
The two original adjunctions give
\[
 \Hom_H(o^*\Sigma_oA,B)
 \cong\Hom_O(\Sigma_oA,\Pi_oB)
 \cong\Hom_H(A,o^*\Pi_oB).
\]
The monad of this additional adjunction is \(K_oT_o\), not \(T_o\).

\item \textbf{Transport from an algebra.}
A \(T_o\)-algebra gives transports
\(\alpha_{h,h'}:A_{h'}\to A_h\) on equal-observation pairs, satisfying
\(\alpha_{h,h}=\id\) and
\(\alpha_{h,h'}\alpha_{h',h''}=\alpha_{h,h''}\).
These equations imply
invertibility.

\item \textbf{Effective descent quotient.}
Quotient the total family by
\((h,a)\sim(h',a')\) iff
\(o(h)=o(h')\) and \(a'=\alpha_{h',h}(a)\).
Topos equivalence relations
have effective pullback-stable quotients.

\item \textbf{Recovery of the original family.}
The quotient lies over \(O\).
Its pullback to a fixed \(h\) is isomorphic to \(A_h\): the inverse sends
\([h',a']\) to \(\alpha_{h,h'}a'\), independently of representative.
This proves the first equivalence, including maps.

\item \textbf{Coalgebra descent.}
A \(K_o\)-coalgebra
gives the opposite transport notation \(c_{h,h'}:A_h\to A_{h'}\), with
the same identity and composition laws, proving the second equivalence.
No representative is chosen per observation.

\item \textbf{Predicate idempotence.}
On predicates the composites are the least saturated predicate above and
the greatest saturated predicate below.
This gives idempotence.

\item \textbf{Evidence nonidempotence.}
For
nonidempotence on evidence take \(H=\{0,1\}\), \(O=1\): a singleton
family has \(T_o\)-fibre sizes \(2,4\) after one and two applications;
a two-element family has \(K_o\)-fibre sizes \(4,16\).
\end{enumerate}
\end{proof}

\begin{proposition}[Correct action adjunction and composition]
For \(a=(X\xleftarrow{s}E_a\xrightarrow{t}Y)\), put
\(\Diamond_a=\Sigma_s t^*\) and \(\Box_a=\Pi_s t^*\). Then
\[
 \Diamond_a\dashv\Box_{a^{\op}},\qquad
 \Diamond_{a;b}\cong\Diamond_a\Diamond_b,\qquad
 \Box_{a;b}\cong\Box_a\Box_b.
\]
\end{proposition}
\begin{proof}\vmode
\leavevmode
\begin{enumerate}[label=(\roman*)]
\item \textbf{Action adjunction.}
\(\Hom_X(\Sigma_st^*B,A)\cong
\Hom_{E_a}(t^*B,s^*A)\cong\Hom_Y(B,\Pi_ts^*A)\).

\item \textbf{Composition by pullback.}
For composition, the witness carrier is \(E_a\times_YE_b\).

\item \textbf{Base-change comparison.}
Beck--Chevalley identifies the two dependent sums, respectively products,
over the same pairs of composable witnesses.

\item \textbf{Coherence and direction of adjunction.}
For three or more actions
the comparison retains the same tuple, proving coherence.
There is no
general adjunction between forward diamond and forward box.
\end{enumerate}
\end{proof}

An actual category of executions supplies a monad by summing over outgoing
executions: identity is the unit and composition is multiplication. A
single transition relation need not do so. For the concurrent execution,
the witnesses are its actual causal executions or its proved exchange
classes; this operation does not replace them by unexamined words.

\subsection{Policy categories, affirmation, and resource framing}\label{subsec:admin-policy}
Evidence about a history must also be organized by the permitted policy actions. Witnessed policy paths record those actions and their resource frames; the resulting free constructions supply the affirmation operations used in the resource calculus.

\begin{definition}[The policy category]\label{def:admin-policy}
Fix a service and one \textbf{decision view} $v$. Let $\mathsf R_v$ be its resource types and let $M_v$ be their finite-multiset commutative monoid. This forgets occurrence names for the purpose of resource typing; physical histories still retain them. A policy generator consists of a named rule with its selected credential witnesses at $v$, a replacement $u\to u'$, and an unchanged frame $z$, giving
\[
             u+z\xrightarrow{(\rho,z)}u'+z.
\]
Writing $a_\rho,b_\rho$ for the input and output multisets of the witnessed rule, its generator fibres are
\[
 \mathsf{Gen}_v(m,n)=\{(\rho,z):a_\rho+z=m,\ b_\rho+z=n\}.
\]
Define $\mathcal A_v$ to have objects $M_v$ and paths of these generators as arrows, modulo the congruence generated by the diamonds
\[
\begin{tikzcd}[column sep=large,row sep=large]
u_1+u_2+z \arrow[r,"{\rho_1}"] \arrow[d,"{\rho_2}"']
  & u'_1+u_2+z \arrow[d,"{\rho_2}"] \\
u_1+u'_2+z \arrow[r,"{\rho_1}"']
  & u'_1+u'_2+z
\end{tikzcd}
\]
Here the two instances act on the two displayed resource portions and read the fixed persistent view. Repeated types in a multiset do not identify distinct named rule or credential witnesses. A diamond witnesses two disjoint allocations before their occurrence names are forgotten.
\end{definition}

\begin{definition}[Policy paths with changing views]\label{def:varying-view-policy-category}
Let $\mathcal P_{\rm phys}$ be the category of physical decision
histories and their witnessed extension paths modulo physical
independent squares. A history $h$ has its actual consulted view
$v(h)$ and inventory multiset $m(h)\in M_{v(h)}$.
Define a graph with objects $(v,m)$, $m\in M_v$, and the following
generators:
\leavevmode
\begin{enumerate}[label=(\roman*)]
\item every fixed-view witnessed policy generator
$(v,m)\to(v,n)$ of Definition~\ref{def:admin-policy};
\item for each actual physical view-changing step $e:h\to h'$,
a generator
$\gamma_e:(v(h),m(h))\to(v(h'),m(h'))$ retaining that step's
view-update witness.
\end{enumerate}
An unchanged-view physical policy step maps to its witnessed
generator, and a step with unchanged view and no inventory/policy
effect maps to the identity. On a view-changing step use $\gamma_e$.
Quotient the free path category by the fixed-view policy diamonds
and by the images, under this assignment, of all physical independent
squares. Denote the quotient by $\mathcal A_{\rm hist}$.

Each endpoint inventory is typed from the actual owned occurrences
at that endpoint. Thus $\gamma_e$ does not posit a map
$M_v\to M_{v'}$ for arbitrary inventories or a monad morphism between
arbitrary views. It is a typed generator at the two recorded
inventories.
\end{definition}
\begin{lemma}[Typed historical readback]\label{lem:historical-readback-functor}
The assignments above define a functor
$R:\mathcal P_{\rm phys}\to\mathcal A_{\rm hist}$ with
$R(h)=(v(h),m(h))$. At a fixed view its policy part agrees with
the existing readback into $\mathcal A_v$.
\end{lemma}
\begin{proof}\vmode
\leavevmode
\begin{enumerate}[label=(\roman*)]
\item \textbf{Typed generators and squares.}
Every generating arrow has the displayed endpoints, so concatenation
defines a free-path functor.
A physical independent square is one
of the generating relations of the target.
The functor therefore
descends to the source quotient.

\item \textbf{Constant-view specialization.}
At a constant view a physical
policy step has exactly its consumed, produced and framed type
multisets and its actual credential witness, which is the generator
of $\mathcal A_v$.
Its independent square is the corresponding
fixed-view diamond.

\item \textbf{Functoriality.}
Identities and composition are inherited from
path concatenation.
\end{enumerate}
\end{proof}

\begin{proposition}[Category, framing, and execution readback]\label{prop:admin-policy-cat}
Concatenation defines $\mathcal A_v$. Adding a frame $z$ defines a functor $z+(-)$, with $0+(-)=\Id$ and $z+(w+(-))=(z+w)+(-)$. Physical decision paths whose consulted view is $v$ map to policy paths, preserving their independent exchanges. For varying views, the same readback is a functor into the category generated by view-tagged policy rules and the actual view-update steps.
\end{proposition}
\begin{proof}\vmode
\leavevmode
\begin{enumerate}[label=(\roman*)]
\item \textbf{Path quotient.}
Take the smallest equivalence relation on parallel paths containing the displayed diamonds and closed under left and right concatenation.
The endpoints of each diamond agree, so the quotient inherits identities and associative concatenation.

\item \textbf{Resource framing.}
Adding $z$ to every frame takes a generating diamond to the same diamond with a larger frame.
It therefore descends to the quotient and satisfies the two functor equations by addition of multisets.

\item \textbf{Physical readback.}
A physical policy step records the consumed and produced types, its actual credential witness, and the types of all unchanged resources.
Send it to that generator.

\item \textbf{Independent and silent steps.}
A physical independent square replaces disjoint occurrences and has unchanged guard reads, hence maps to the corresponding policy diamond.
A physical step with no policy effect maps to an identity when its view is unchanged.

\item \textbf{Changing views.}
For changing views, retain each actual view change as a generator between its tagged states, and quotient by the images of physical independent squares.
The same assignment on generators then respects every exchange and defines the stated functor.

\item \textbf{Scope of readback.}
This constructs a readback of executions; it does not declare every abstract policy path to be schedulable by a fixed sequential service.
\end{enumerate}
\end{proof}

\begin{definition}[Policy affirmation and free evidence]\label{def:policy-affirmation-operations}
For $X:M_v\to\Set$ define
\[
 (T_vX)(m)=\coprod_n\mathcal A_v(m,n)\times X(n),\quad
 \eta_X(x)=(1_m,x),\quad \mu_X(p,(q,x))=(qp,x).
\]
An element $(p,x)$ is a \textbf{policy affirmation}. On family maps $f:X\to Y$ use $T_vf(p,x)=(p,f(x))$. The \textbf{free policy-evidence functor} $F_v:\Set^{M_v}\to\Set^{\mathcal A_v^\op}$ has the same object formula and action $(F_vX)(r)(p,x)=(pr,x)$; $U_v$ restricts a presheaf to its object-indexed family. Its counit candidate is $(p,y)\mapsto Y(p)y$. The next results prove the monad and adjunction laws. For $m\in M_v$, the \textbf{concentrated family} $\mathbf m$ is $\mathbf m(n)=\{*\mid n=m\}$.
\end{definition}

\begin{theorem}[Policy evidence, Kan extension, and the derived monad]\label{thm:policy-evidence-adjunction-monad}
The arrow span of $\mathcal A_v$ is a monad in spans and in its response-game interpretation. For the discrete object inclusion $j_v:M_v\to\mathcal A_v$, the free-evidence adjunction is
\[
 F_v=\operatorname{Lan}_{j_v^\op}:\Set^{M_v}
 \rightleftarrows\Set^{\mathcal A_v^\op}:U_v=j_v^*,\qquad U_vF_v=T_v,
\]
with
\begin{equation}\label{eq:admin-aff}
 (T_vX)(m)=\coprod_n\mathcal A_v(m,n)\times X(n).
\end{equation}
Its unit is $x\mapsto(1_m,x)$, its multiplication is $(p,(q,x))\mapsto(qp,x)$, and the counit is $(p,y)\mapsto Y(p)y$. The algebra category is $\Set^{\mathcal A_v^\op}$. Thus a policy affirmation retains a path and evidence at its endpoint.
\end{theorem}
\begin{proof}\vmode
\leavevmode
\begin{enumerate}[label=(\roman*)]
\item \textbf{Path-span monoid.}
The span $M_v\leftarrow(\mathcal A_v)_1\to M_v$ has identity-arrow unit and composition multiplication on the pullback of composable arrows.
Its laws are the category laws, and Theorem~\ref{thm:responsecomposition} transports them to response games.

\item \textbf{Family monad and its laws.}
Its backward action on families is the displayed sum over arrows with source $m$.
Both unit equations remove an identity path; both associative composites send $(p,(q,(r,x)))$ to $(rqp,x)$.

\item \textbf{Free presheaf action.}
Precomposition gives $(F_vX)(u)(p,x)=(pu,x)$.

\item \textbf{Free-forgetful adjunction.}
A family map $g:X\to U_vY$ extends uniquely by $(p:m\to n,x)\mapsto Y(p)g_n(x)$.
Restriction to $(1_m,-)$ is inverse, by naturality.
These natural bijections prove the adjunction and its triangles.

\item \textbf{Kan-extension identification.}
Since $M_v$ is discrete, its pointwise left Kan formula is precisely this sum over paths into the argument in $\mathcal A_v^\op$, proving the claimed Kan extension.
The induced multiplication composes the retained paths, so its monad is the same $T_v$.

\item \textbf{Algebras and presheaves.}
An algebra assigns maps $X(n)\to X(m)$ to paths $p:m\to n$.
The unit and multiplication equations say that identities act identically and composition acts contravariantly.
Algebra morphisms satisfy exactly naturality, proving the algebra-category identification.
\end{enumerate}
\end{proof}


\begin{definition}[Resource tensor]\label{def:admin-tensor}
For families on $M_v$, put
\[
 (X\otimes_vY)(m)=\coprod_{m_1+m_2=m}X(m_1)\times Y(m_2),
 \qquad I_v(m)=\{*\mid m=0\}.
\]
The index of a tensor element is a decomposition of the type multiset. For a named finite inventory $I$ with type map $\tau:I\to\mathsf R_v$, define
\[
 |I|_\tau=\sum_{i\in I}[\tau(i)],\qquad
 \mathsf{Split}(I)=\{(I_1,I_2):I_1\amalg I_2=I\}.
\]
Its type-forgetting map sends $(I_1,I_2)$ to $(|I_1|_\tau,|I_2|_\tau)$. It is surjective onto the decompositions $m_1+m_2=|I|_\tau$: induct on a finite enumeration of $I$, placing each occurrence into a part with a remaining count of its type. This proves existence, without selecting a distinguished partition. The map need not be injective; two equal-typed occurrences admit two different allocations of one occurrence to each part. A physical realization of a tensor term retains the selected element of $\mathsf{Split}(I)$, while the graded tensor retains its image. Reassociation of physical allocations and of tensor indices both sends a partition into three parts to the same ordered triple.
\end{definition}

\begin{proposition}[Tensor and policy strength]\label{prop:admin-strength}
The displayed tensor is symmetric monoidal. Framing constructs a strength
\[
 X\otimes_v T_vY\longrightarrow T_v(X\otimes_vY)
\]
compatible with the unit and multiplication of $T_v$. The operation never creates an additional owned occurrence.
\end{proposition}
\begin{proof}\vmode
\leavevmode
\begin{enumerate}[label=(\roman*)]
\item \textbf{Tensor coherence.}
Either bracketing of three tensors is indexed by the same triples $m_1+m_2+m_3=m$ with the same three entries.
These identifications give the associator, its pentagon, and the unitors; swapping the two summands gives symmetry and its equations.

\item \textbf{Framed strength.}
For a tensor element with $x\in X(m_1)$ and $(p:m_2\to n_2,y)$, use the framed path $m_1+p:m\to m_1+n_2$ and return $(x,y)$ in its target tensor summand.

\item \textbf{Unit and multiplication laws.}
Framing the identity gives the strength unit law.
Framing a composite is the composite of framed paths, giving compatibility with multiplication.

\item \textbf{Associativity and resource retention.}
Iterated framing adds the same multisets in either order, giving the associativity equation for strength.
The $X$ resource allocation is carried unchanged as a frame; it is not inserted twice into the inventory.

\item \textbf{Strength equations.}
Writing $\mathsf{st}$ for the displayed strength and omitting the subscript $v$, its multiplication square is
\[
\begin{tikzcd}[column sep=large,row sep=large]
X\otimes T^2Y \arrow[r,"{\mathsf{st}_{X,TY}}"] \arrow[d,"{1\otimes\mu_Y}"']
 &T(X\otimes TY) \arrow[r,"{T\mathsf{st}_{X,Y}}"] &T^2(X\otimes Y) \arrow[d,"{\mu_{X\otimes Y}}"]\\
X\otimes TY \arrow[rr,"{\mathsf{st}_{X,Y}}"'] &&T(X\otimes Y).
\end{tikzcd}
\]
It commutes because $m_1+(qp)=(m_1+q)(m_1+p)$ in the policy category.
\end{enumerate}
\end{proof}

\begin{lemma}[Parallel policy programs]\label{lem:admin-parallelterms}
For $p:m_1\to n_1$ and $q:m_2\to n_2$, the paths
\[
 (n_1+q)(p+m_2),\qquad (p+n_2)(m_1+q)
\]
are equal in $\mathcal A_v$. They define a natural map
$T_vX\otimes_vT_vY\to T_v(X\otimes_vY)$, coherent with tensor symmetry, associativity, identity programs, and sequencing on disjoint allocations.
\end{lemma}
\begin{proof}\vmode
\leavevmode
\begin{enumerate}[label=(\roman*)]
\item \textbf{Exchange of disjoint paths.}
For two generators the equality is the defining disjoint-allocation diamond.
For arbitrary finite paths, move each generator of $p$ past each generator of $q$, retaining their respective frames.
This gives the displayed equality.

\item \textbf{Parallel evidence.}
Append the two endpoint evidence entries in their tensor summand to define the natural map.

\item \textbf{Associativity and symmetry.}
Three paths on three allocations can be executed in any order by the same adjacent exchanges; hence either associativity map has the same path class and evidence tuple.
Swapping two paths gives symmetry, an empty path gives the unit, and concatenating within either allocation gives the sequencing equation.

\item \textbf{Monoidal compatibility.}
These equalities also show that the map is well-defined when either input path is replaced by an exchange-equivalent path.
\end{enumerate}
\end{proof}

\begin{proposition}[Convolution, parallel policy, and Kleisli readback]\label{prop:admin-policy-kleisli}
The resource tensor is the Day convolution of families on the discrete commutative monoid $M_v$. With this tensor, $T_v$ is a commutative strong monad: the two orders of its strength-induced parallel computation agree and equal the map of Lemma~\ref{lem:admin-parallelterms}. Its Kleisli category is symmetric monoidal, and the concentrated families give a fully faithful functor
\[
 \mathcal A_v\longrightarrow\operatorname{Kl}(T_v),\qquad
 m\longmapsto\mathbf m,\qquad
 \Hom_{\operatorname{Kl}(T_v)}(\mathbf m,\mathbf n)
 \cong\mathcal A_v(m,n).
\]
\end{proposition}
\begin{proof}\vmode
\leavevmode
\begin{enumerate}[label=(\roman*)]
\item \textbf{Discrete convolution.}
For the discrete indexing category, convolution along $+$ is the sum over $m_1+m_2=m$ in Definition~\ref{def:admin-tensor}; there are no nonidentity index arrows to quotient.

\item \textbf{Parallel path tensor.}
For $p:m_1\to n_1$ and $q:m_2\to n_2$, the two composites of strength execute $(n_1+q)(p+m_2)$ and $(p+n_2)(m_1+q)$.
Lemma~\ref{lem:admin-parallelterms} identifies them.
Write their common path as $p\otimes q$.

\item \textbf{Unit and multiplication compatibility.}
The resulting map $TX\otimes TY\to T(X\otimes Y)$ sends $((p,x),(q,y))$ to $(p\otimes q,(x,y))$.
Identity paths give its unit compatibility.
For two successive paths in each factor, exchanging the middle disjoint steps gives
\[
 (p'p)\otimes(q'q)=(p'\otimes q')(p\otimes q),
\]
which is compatibility with monad multiplication.

\item \textbf{Coherence and Kleisli tensor.}
Three disjoint paths and exchange of two factors give associativity and symmetry; the same path classes prove all coherences.
Thus tensor of Kleisli maps is their ordinary tensor followed by this comparison, with the unit and composition just verified.

\item \textbf{Concentrated-family readback.}
A family map $\mathbf m\to T\mathbf n$ is an element of $(T\mathbf n)(m)$, and this fibre is $\mathcal A_v(m,n)$.
The Kleisli identity is the identity path and Kleisli composition is concatenation, proving full faithfulness.

\item \textbf{Physical interpretation.}
This readback concerns policy derivations.
The physical execution functor of Proposition~\ref{prop:admin-policy-cat} retains the worker discipline and does not identify every abstract policy exchange with a physical independent square.
\end{enumerate}
\end{proof}

\begin{proposition}[Monoidal Kan-extension construction]\label{prop:admin-monoidal-kan}
The framed policy diamonds construct a symmetric monoidal structure on $\mathcal A_v$, with tensor on objects $+$ and tensor on arrows
\[
 p\otimes q=(n_1+q)(p+m_2)=(p+n_2)(m_1+q)
 \quad(p:m_1\to n_1,\ q:m_2\to n_2).
\]
The identity-on-objects functor $j_v:M_v\to\mathcal A_v$ is strong symmetric monoidal. With Day convolution on the two presheaf categories, the policy-evidence adjunction is a monoidal adjunction. Its induced monoidal monad is exactly the commutative strong monad already calculated above.
\end{proposition}
\begin{proof}\vmode
\leavevmode
\begin{enumerate}[label=(\roman*)]
\item \textbf{Monoidal policy category.}
Lemma~\ref{lem:admin-parallelterms} proves that the displayed tensor is well-defined.
The interchange equation in Proposition~\ref{prop:admin-policy-kleisli} proves bifunctoriality.
Iterated framing gives strict association and unit, and the diamond equality gives symmetry, with coherence on every path.

\item \textbf{Monoidal inclusion.}
The inclusion $j_v$ preserves $+$, $0$ and their structural arrows, so is strong symmetric monoidal by construction.

\item \textbf{Kan tensor comparison.}
Write $F=\operatorname{Lan}_{j_v^\op}$ and $\star$ for convolution on $\Set^{\mathcal A_v^\op}$.
Distributing products over sums and applying co-Yoneda twice gives
\begin{align*}
 (FX\star FY)(m)
 &\cong\int^{a,b}\mathcal A_v(m,a+b)\times
       \left(\coprod_u\mathcal A_v(a,u)\times X(u)\right)\\
 &\hspace{25mm}{}\times\left(\coprod_w\mathcal A_v(b,w)\times Y(w)\right)\\
 &\cong\coprod_{u,w}\mathcal A_v(m,u+w)\times X(u)\times Y(w)\\
 &\cong F(X\otimes_vY)(m).
\end{align*}
Explicitly the middle comparison sends a path $r:m\to a+b$ and paths $p:a\to u$, $q:b\to w$ to $(p\otimes q)r$.

\item \textbf{Inverse and unit.}
Its inverse takes $a=u$, $b=w$ and identity paths; the coend relations prove the inverse equations.
Also $FI_v=\mathcal A_v(-,0)$, the convolution unit.

\item \textbf{Strong monoidal coherence.}
Associativity, unit and symmetry commute because all routes compose the same tensor of paths.
Thus $F$ is strong symmetric monoidal.
This is the applicable instance of monoidal Kan extension along a strong promonoidal functor~\cite[Proposition~1]{day-street-kan}.

\item \textbf{Monoidal adjunction.}
Restriction $U_v=j_v^*$ has the lax monoidal comparison sending $(y,z)$ in the summand $m=a+b$ to the convolution class $[1_m,y,z]$.
The unit and counit of Theorem~\ref{thm:policy-evidence-adjunction-monad} commute with these comparisons by their identity-path and path-action formulas.
Hence the adjunction is monoidal.

\item \textbf{Induced monoidal monad.}
Applying $U_v$ to the preceding strong comparison yields the monad's map
\[
 ((p,x),(q,y))\longmapsto(p\otimes q,(x,y)).
\]
This is the map of Lemma~\ref{lem:admin-parallelterms}; composing its first input with the monad unit yields exactly the strength of Proposition~\ref{prop:admin-strength}.
The monoidal unit and multiplication equations therefore coincide with the explicit equations already proved, and their existing diagrams are the corresponding coherence diagrams of this adjunction.
\end{enumerate}
\end{proof}

The monad $T_v$ is a policy execution monad on evidence families. Its arrows are policy paths acting on evidence, whereas the service output-word monad of Section~\ref{subsec:crypto-games} concatenates replies and the lift monad classifies return domains. Their actions and comparisons are given at their respective interfaces.

\subsection{Evidence transport over retained histories}\label{subsec:admin-audit-modes}
The span and evidence adjunctions of Section~\ref{subsec:admin-equipment} must retain the history to which their witnesses belong. We now perform their set-family constructions over the chosen history coordinate, including an inherited entry history under prefix insertion. The following proposition identifies the resulting fibrewise comparisons.

\begin{proposition}[Transport at an audited past]\label{prop:admin-audittransport}
For a map of evidence indices $f:X\to Y$ over $H$, the preceding span adjunction and the evidence adjunctions preserve the audit coordinate. Their units, counits, and mates constructed from maps over $H$ do so as well. Pulling them back to a history $h$ recovers the constructions on its evidence fibres, coherently with any further base change. Thus the evidence operations retain every event of the selected history presentation.
\end{proposition}
\begin{proof}\vmode
\leavevmode
\begin{enumerate}[label=(\roman*)]
\item \textbf{Construction over the history base.}
Perform Theorem~\ref{thm:admin-equipment} in $\E/H$.
Every carrier map, diagonal, and projection then lies over the identity of $H$.

\item \textbf{Fibrewise evidence operations.}
The family formulas in Proposition~\ref{prop:evidence-action-base-change} sum or multiply over $x$ with fixed image $y$; since $f$ is over $H$, these elements all have the same history coordinate.

\item \textbf{Units, counits and mates.}
The units and counits are the same fibre maps, and the mate formula uses only their composition and functorial images, so it also preserves that coordinate.

\item \textbf{Base change.}
Pullback stability and the base-change isomorphisms identify their fibres at $h$.

\item \textbf{History-extension transport.}
If evidence has specified transport along a history extension, naturality of these maps gives its commuting square.
For evidence represented by finite causal witnesses, prefix insertion retains those witnesses and composes by union, giving that transport directly.
Current decision predicates instead use their specified change-of-view maps.
\end{enumerate}
\end{proof}

\subsection{Game interpretation of evidence interfaces and modes}\label{subsec:admin-game-modes}
The evidence adjunctions now receive a game interpretation. Their discrete response interfaces admit an exact calculation, while the underlying protocol history game retains its concurrent cells. A one-request interface for a pair of values therefore remains distinct from two independently arriving requests; the comparison below keeps those types explicit.

\begin{theorem}[Evidence spans and policy monads in games]\label{thm:admin-matrices}
For a span $R:X\rightsquigarrow Y$, its discrete Scott kernel $W_R$ and its initialized response kernel $\Phi_R$ of Section~\ref{subsec:execution-interaction-geometry} preserve composition, identity spans, and maps of spans. They therefore transport the graph adjunctions of Theorem~\ref{thm:admin-equipment} and the horizontal policy monad of Theorem~\ref{thm:policy-evidence-adjunction-monad}. On the full categories of response tests $X\rightsquigarrow\{*\}$, their action is exactly $\Sigma_s t^*$, with the same evidence adjunctions, locks, and keys.

Resource tensor is represented by the explicit splitting span $+^\dagger:M_v\rightsquigarrow M_v\times M_v$ followed by the external product of the two evidence families. Its game interpretation has fibre $(X\otimes_vY)(m)$, the sum over resource allocations in Definition~\ref{def:admin-tensor}, and preserves the tensor and strength equations proved above.
\end{theorem}
\begin{proof}\vmode
\leavevmode
\begin{enumerate}[label=(\roman*)]
\item \textbf{Response interpretation of spans.}
Theorem~\ref{thm:game-return-matrix} identifies discrete kernel composition with the pullback of witness carriers; Theorem~\ref{thm:responsecomposition} gives its initialized strategy comparison.
Both send a composite to the same ordered witness tuple, naturally in span maps.

\item \textbf{Structural comparisons.}
Identity, associativity, graph units and counits, and policy multiplication therefore transport with their proved equations.

\item \textbf{Evidence-family category.}
A response test $X\rightsquigarrow\{*\}$ is exactly a family over $X$.
Its strategy maps fix the request and map reply witnesses within that fibre.
Hence the full response-test category is $\mathsf{Fam}(X)$ and precomposition is $\Sigma_s t^*$, as in Proposition~\ref{prop:evidence-action-base-change}; its evidence adjunctions are unchanged.

\item \textbf{Resource splitting.}
The fibre of the splitting span $+^\dagger$ over $m$ consists of pairs $m=m_1+m_2$.
Adding the two evidence entries gives precisely $(X\otimes_vY)(m)$.

\item \textbf{Coherence and interface distinction.}
Reassociation retains the same three-part allocation, and framing retains the same path and unchanged resource part, proving transport of tensor and strength coherence.
A single request for a pair is still distinct from two independent request ports.
\end{enumerate}
\end{proof}

\begin{proposition}[Representable kernel modes]\label{prop:admin-representables}
For a functor $F:C\to D$ of small ordinary categories, the set-valued kernels
\[
 F_!(c,d)=\Hom_D(d,Fc),\qquad
 F^\dagger(d,c)=\Hom_D(Fc,d)
\]
form an adjunction in the ordinary unpointed profunctor bicategory. For rooted categories, root-preserving $F$ gives pointed kernels and pointed unit and counit, rooted at identities. A transformation $\alpha:F\Rightarrow G$ gives unpointed cells $F_!\Rightarrow G_!$ and $G^\dagger\Rightarrow F^\dagger$, related by mates; their root elements have image $\alpha_i$. For occurrence-induced Scott functors, all these cells are canonically pointed.
\end{proposition}
\begin{proof}\vmode
\leavevmode
\begin{enumerate}[label=(\roman*)]
\item \textbf{Representable unit.}
Co-Yoneda identifies $(F^\dagger\star F_!)(c,c')$ with $\Hom_D(Fc',Fc)$.
Applying $F$ to arrows gives its unit from $\Hom_C(c',c)$.

\item \textbf{Representable counit.}
Composition in $D$ gives the counit
\[
 \int^c\Hom_D(Fc,d)\times
                  \Hom_D(d',Fc)
       \longrightarrow\Hom_D(d',d).
\]
In either triangle a representative arrow is composed with its identity and evaluated by co-Yoneda, leaving that arrow unchanged.

\item \textbf{Triangle equations.}
The same evaluation formulas commute with reassociation of the coend quotients.

\item \textbf{Comparison cells and mates.}
A component $Fc\to Gc$ acts by postcomposition in the first kernel and precomposition in the second.
Naturality and the counit formula give the mate equation of Lemma~\ref{lem:admin-key}.

\item \textbf{Root preservation.}
Units and counits preserve the root because they apply $F$ to an identity or compose two identities.

\item \textbf{Reached occurrence modes.}
For an occurrence map, $F(\varnothing)=\varnothing$, and thinness of the Scott configuration category forces $\alpha_{\varnothing}=\id_{\varnothing}$.
Thus its comparison cells and mates preserve the root as well.
This last conclusion is derived for these functors, not asserted for arbitrary transformations of rooted functors.
\end{enumerate}
\end{proof}

An arbitrary transformation of root-preserving functors need not preserve a kernel point: for $C=\mathbf 1$ and $D=B\mathbb N$, the transformation with component $1$ sends the root $0$ of the representable to $1$. Proposition~\ref{prop:admin-representables} separates this unpointed statement from the canonically pointed occurrence case used below.

\begin{definition}[Occurrence interface maps]\label{def:admin-occmap}
An occurrence interface map $f:A\rightharpoonup B$ is a polarity-preserving partial function on events which sends configurations to configurations and is injective on every configuration where defined. It induces $F(x)=f[x]$ on finite configurations. This map is monotone for the Scott orders: a negative removal or positive addition maps to an extension of the same sign, possibly an identity. It preserves compatible unions and nonempty compatible intersections. For the latter statement, local injectivity in a common bounding configuration identifies the preimages of any common image event.
\end{definition}

The principal interface constructions satisfy these conditions directly. For a visible event set $V\subseteq A$, restrict the arena to $V$ with inherited order and consistency and use the partial identity $A\rightharpoonup A\mathord\downarrow V$. It sends a configuration to its intersection with $V$, which is consistent and downward closed in the restricted order. Local injectivity and polarity preservation follow from the partial identity. Disjoint interface reindexing is an isomorphism on the selected factors. Thus these maps require no further admissibility condition.

To apply the functional construction of~\cite[Definition~5.9]{winskel-functional}, extend $F(x)=f[x]$ to all configurations. Images of infinite configurations are configurations: each finite subset lies in the image of a finite subconfiguration. Direct image preserves compatible unions, including the empty union.

For a nonempty compatible family $(x_i)$, take $b\in\bigcap_i f[x_i]$ and a preimage in one member. Local injectivity in a common bounding configuration forces that same preimage to lie in every member. Hence $f[\bigcap_i x_i]=\bigcap_i f[x_i]$. Finite configurations have finite images, and polarity is preserved. These properties prove that $F$ is additive-stable. The following finite principal histories give its local realization, and the established functional construction supplies the strategy adjunction.

\begin{theorem}[Occurrence maps give reached strategy modes]\label{thm:admin-occurrence-lifts}
An occurrence interface map has receptive, courteous lifts
$f_!:A\to B$ and $f^\dagger:B\to A$ with configuration families
\begin{align*}
 \mathcal F_!&=\{(x,y):y\sqsubseteq_B f[x]\},\\
 \mathcal F_\dagger&=\{(y,x):f[x]\sqsubseteq_B y\}.
\end{align*}
Their Scott kernels are the subterminal representable kernels of $F$ in Proposition~\ref{prop:admin-representables}. Cubical unstraightening directly recovers the constructed strategies, and these kernels are reached. The lifts form an adjunction for ordinary secured strategy composition. Pointwise Scott comparisons between two such $F,G$ give strategy keys and their mates.
\end{theorem}
\begin{proof}\vmode
\leavevmode
\begin{enumerate}[label=(\roman*)]
\item \textbf{Functional strategy construction.}
The preceding verification makes $F$ additive-stable on all configurations.
Theorems~5.2 and~5.11 of the extended version of \cite{winskel-functional} therefore construct the two receptive, courteous strategies with the displayed configuration families and give their adjunction for ordinary secured interaction and hiding.

\item \textbf{Finite principal-history presentation.}
Restricting to finite configurations gives the families used here.
Their prime events are the principal histories of Lemma~\ref{lem:admin-family}; thus no separate strategy presentation is supplied.

\item \textbf{Kernel identification.}
Each family has at most one configuration over a given interface configuration.
Theorem~\ref{thm:scott-discrete-fibrations} consequently identifies its Scott kernel with the corresponding subterminal representable and its cubical unstraightening with the constructed configuration carrier.

\item \textbf{Reachability and adjunction.}
Every finite strategy configuration has a linear extension from the empty configuration, so these kernels are reached.
Proposition~\ref{prop:admin-representables} proves pointedness of the representable unit, counit, comparisons and mates; Theorem~\ref{thm:ordinarycomp} identifies the reached kernel composition with the cited ordinary strategy composition.

\item \textbf{Pointwise comparison cells.}
For completeness, a pointwise Scott comparison $F\Rightarrow G$ gives the family inclusions $\mathcal F_{f,!}\subseteq\mathcal F_{g,!}$ and $\mathcal F_{g,\dagger}\subseteq\mathcal F_{f,\dagger}$.

\item \textbf{Well-defined principal-history map.}
Map a principal source occurrence to the target principal history of its label.
If two bounding configurations contain the same source principal history, that history itself belongs to the target family and is an ideal in each bound; Lemma~\ref{lem:admin-family} therefore makes the assigned target occurrence independent of the bound.

\item \textbf{Keys and mate equation.}
This defines the strategy keys and gives the stated kernel maps.
Their mate equation is the representable mate equation, or equivalently the unit--counit equation in the constructed adjunction.
\end{enumerate}
\end{proof}

\begin{example}[A directed key from erasure]\label{ex:admin-key}
Let $A$ have one positive event and let $f$ erase it. On configurations, $f[-]\sqsubseteq\Id$. The lift $f_!$ accepts the negative request and produces no positive event, whereas $\Id_!$ is copycat. Inclusion is a key $f_!\Rightarrow\Id_!$, and its mate is $\Id^\dagger\Rightarrow f^\dagger$. Both are obtained from the partial identity on occurrences and the Scott construction. A map of arbitrary observation values is instead handled by Theorem~\ref{thm:admin-matrices}; it is not silently promoted to an occurrence interface map.
\end{example}

\begin{definition}[Outgoing strategy-test modes]\label{def:formal-outgoing-strategy-test-modes}
At each test interface $X$, an ordinary adjunction $\sigma\dashv\rho$ gives the mode $(-)\rho\dashv(-)\sigma$ on outgoing strategy tests. Its hom bijection sends $u:a\rho\to b$ to $(u\sigma)(a\eta)$, with inverse $v\mapsto(b\epsilon)(v\rho)$. The two triangles prove inverseness and whiskering proves naturality in the test. Hence the \textbf{lock} and key calculus of Definitions~\ref{def:admin-mode} and~\ref{def:evidence-transposition-rules} applies to the \textbf{reached} ordinary mode lifts just constructed. Proposition~\ref{prop:admin-representables} supplies the corresponding raw kernel modes. The configuration construction and Theorem~\ref{thm:ordinarycomp} identify the reached strategy modes used here.
\end{definition}

\subsection{Protocol histories, evidence descent, and execution}\label{subsec:admin-game-execution}
The causal-view compiler can have several commitment histories representing one physical history. Apply it to a protocol presentation and its declared control interface, and pull observations back along its witness-preserving erasure. We now identify when evidence is independent of the compiled representative and how it descends through that erasure.

\begin{definition}[Compiler descent category]\label{def:compiler-descent-category}
For commitment erasure $\varepsilon:\widehat H\to H$, the \textbf{kernel-pair groupoid} has objects $\widehat H$, an arrow $(\hat h,\hat k):\hat h\to\hat k$ exactly when $\varepsilon\hat h=\varepsilon\hat k$, identity $(\hat h,\hat h)$ and composition $(\hat k,\hat l)(\hat h,\hat k)=(\hat h,\hat l)$. The \textbf{descent category} $\mathsf{Desc}(\varepsilon)$ has objects families $X\to\widehat H$ with bijections $\theta_{\hat h,\hat k}:X_{\hat h}\to X_{\hat k}$ satisfying identity and composition; arrows are family maps commuting with these bijections. Define $K_\varepsilon:\Set/H\to\mathsf{Desc}(\varepsilon)$ by pullback and identity transport between copies of one fibre. The lemma constructs its inverse by a quotient and proves the equivalence.
\end{definition}

\begin{definition}[Descent on history-extension categories]\label{def:history-functor-descent}
Let $\varepsilon:\widehat{\mathcal H}\to\mathcal H$ be the functor
on finite history-extension categories induced by commitment erasure.
Define
$\mathcal R_1=\widehat{\mathcal H}\times_{\mathcal H}
\widehat{\mathcal H}$: its objects are pairs of histories with equal
erasure and its arrows are pairs $(\hat u,\hat v)$ with
$\varepsilon\hat u=\varepsilon\hat v$. Its projection functors are
$p_1,p_2$. A descent object is a functor
$X:\widehat{\mathcal H}\to\Set$ and a natural isomorphism
$\theta:Xp_1\Rightarrow Xp_2$ satisfying
\[
 \theta_{\hat h,\hat h}=1,\qquad
 \theta_{\hat k,\hat l}\theta_{\hat h,\hat k}
       =\theta_{\hat h,\hat l}.
\]
For a pair of lifts of one physical extension its naturality is
\[
\begin{tikzcd}[column sep=large]
 X(\hat h)\arrow[r,"\theta_{\hat h,\hat h'}"]
   \arrow[d,"X(\hat u)"']
 &X(\hat h')\arrow[d,"X(\hat v)"]\\
 X(\hat k)\arrow[r,"\theta_{\hat k,\hat k'}"']
 &X(\hat k').
\end{tikzcd}
\]
Morphisms are natural transformations commuting with $\theta$.
\end{definition}
\begin{lemma}[Quotient extension maps]\label{lem:descent-extension-maps}
For such a descent object, the quotient fibres
\[
 E(h)=
 \left(\coprod_{\varepsilon\hat h=h}X(\hat h)\right)/\theta
\]
have a unique functor structure for which
\[
 E(u)[\hat h,x]=[\hat k,X(\hat u)x]
 \quad\text{whenever }\varepsilon\hat u=u,\quad
 \hat u:\hat h\to\hat k .
\]
This gives an equivalence between functors on $\mathcal H$ and the
displayed history descent category.
\end{lemma}
\begin{proof}\vmode
\leavevmode
\begin{enumerate}[label=(\roman*)]
\item \textbf{Definition on physical extensions.}
Every finite physical extension lifts from every compiled source
by the missing-slot construction of
Construction~\ref{con:compiler-bisimulation-cubes}.

\item \textbf{Independence of representatives.}
For two such
lifts, the naturality square identifies the two target classes.
For different representatives of the source class, the same square
uses their source $\theta$ and again identifies the target classes.
Thus the displayed prescription is a total single-valued relation,
and hence a function, without choosing a global family of lifts.

\item \textbf{Functoriality.}
Identity lifts prove the identity law.
Composing two lifts proves
the composition law; independence of lifts makes it the required
composite for arbitrary representatives.

\item \textbf{Pullback and quotient equivalence.}
Pulling back a functor $E$ gives identity transport between copies
of the same fibre.
Conversely
$x\mapsto[\hat h,x]$ is bijective by the identity and cocycle
equations and is natural by the displayed formula.
These maps give
the inverse equivalence on objects and on equivariant natural
transformations.
Their composites are identities after the canonical
fibre identifications, which also proves the unit and counit equations.
\end{enumerate}
\end{proof}

\begin{lemma}[Constructed evidence descent for compiler erasure]\label{lem:admin-evidence-descent}
Let $\varepsilon:\widehat H\to H$ erase the selected commitments, and let $\mathcal R=\widehat H\times_H\widehat H\rightrightarrows\widehat H$ be its kernel-pair groupoid. A physical evidence family $E\to H$ has a canonical action on its pullback: if $\varepsilon\hat h=\varepsilon\hat k$, the transport from $E_{\varepsilon\hat h}$ to $E_{\varepsilon\hat k}$ is the identity. This constructs an equivalence
\begin{equation}\label{eq:admin-evidence-descent}
 K_\varepsilon:\Set/H\simeq\mathsf{Desc}(\varepsilon).
\end{equation}
It also descends functors on the history-extension categories equipped with natural kernel-pair actions. For $o:H\to O$, write $\hat o=o\varepsilon$. The evidence quantifiers are
\begin{align}
 (\Pi_oE)_v&\cong
 \left\{(e_{\hat h})_{\hat o(\hat h)=v}:
 \begin{aligned}
 &e_{\hat h}\in E_{\varepsilon\hat h},\\[-2pt]
 &e_{\hat h}=e_{\hat k}\text{ if }\varepsilon\hat h=\varepsilon\hat k
 \end{aligned}\right\},\label{eq:admin-desc-product}\\
 (\Sigma_oE)_v&\cong
 \left(\coprod_{\hat o(\hat h)=v}E_{\varepsilon\hat h}\right)\big/\sim,
 \label{eq:admin-desc-sum}
\end{align}
where $\sim$ identifies the same evidence in copies over the same physical history. In the descent category these give the adjunctions
\begin{equation}\label{eq:admin-desc-adjunctions}
 \Sigma_oK_\varepsilon^{-1}\dashv K_\varepsilon o^*
       \dashv\Pi_oK_\varepsilon^{-1}.
\end{equation}
\end{lemma}
\begin{proof}\vmode
\leavevmode
\begin{enumerate}[label=(\roman*)]
\item \textbf{Kernel-pair descent equivalence.}
The identity transports satisfy the identity and cocycle equations.
To construct the inverse of \eqref{eq:admin-evidence-descent}, take a family $X_{\hat h}$ with action $\theta$ and put
\[
 E_h=\left(\coprod_{\varepsilon\hat h=h}X_{\hat h}\right)\big/\sim,
 \qquad
 (\hat h,x)\sim(\hat k,\theta_{\hat h,\hat k}x).
\]
For each $\hat h$, the map $x\mapsto[\hat h,x]$ is bijective: transport supplies surjectivity onto every class, and the identity and cocycle equations supply injectivity.
Equivariant maps descend to these quotients.
This proves the equivalence, the usual effective descent of an epimorphism in a topos; see \cite[Introduction]{prezado-descent}.

\item \textbf{Lifting physical extensions.}
In this application all actions on physical evidence are the constructed identity actions, not additional protocol fields.
Every finite physical extension lifts from every compiled representative: linearize its new events and use Theorem~\ref{thm:compiler} to fill only the missing selected veto commitments before each plant step.
The bound is the sum of the finite veto-support sizes of those steps.

\item \textbf{Quotient extension maps.}
For a functor $X:\widehat{\mathcal H}\to\Set$ with natural kernel-pair action, define its descended extension map by
\[
 E(u)[\hat h,x]=[\hat k,X(\hat u)x],
 \qquad \varepsilon\hat u=u:h\to k.
\]
Such a lift exists from each $\hat h$.

\item \textbf{Well-definedness.}
Two lifts over $u$ form an arrow in the kernel-pair category; naturality of $\theta$ identifies their outputs in the quotient.
The same argument handles different representatives.

\item \textbf{Functoriality and natural transformations.}
Successive lifts prove composition, and identity lifts prove identities.
Natural transformations descend by the same calculation.

\item \textbf{Compatible sections.}
A section over an observation fibre repeats its value on every compiled representative, giving \eqref{eq:admin-desc-product}.
Conversely, compatibility gives a unique value at each physical history; surjectivity of $\varepsilon$ establishes existence without choosing a global section.

\item \textbf{Witness quotients and adjunctions.}
The corresponding quotient of witnesses gives \eqref{eq:admin-desc-sum}.
Transport of the two slice adjunctions across the equivalence gives \eqref{eq:admin-desc-adjunctions}, including their units, counits and triangle equations.

\item \textbf{Predicate specialization.}
For subterminal evidence the compatibility in the product is automatic, and the sum has the same existential image.
Thus the existing predicate formulas are recovered.

\item \textbf{Quotient and pullback distinction.}
Unrestricted products over duplicate compiled representatives need not agree for general evidence, which is why the descent interpretation is used.
The base quotient and its evidence pullback are
\[
\begin{tikzcd}[column sep=large]
\widehat H\times_H\widehat H
 \arrow[r,shift left=1.5,"{p_1}"] \arrow[r,shift right=1.5,"{p_2}"']
 &\widehat H \arrow[r,two heads,"{\varepsilon}"] &H
\end{tikzcd}
\qquad
\begin{tikzcd}[column sep=large]
\varepsilon^*E \arrow[r] \arrow[d] &E \arrow[d]\\
\widehat H \arrow[r,"{\varepsilon}"'] &H .
\end{tikzcd}
\]
The first row is a coequalizer.
The second is a pullback.
Their distinct universal properties are what construct descent; neither requires a selected representative of each fibre.
\end{enumerate}
\end{proof}

\begin{theorem}[Protocol histories and evidence judgements in games]\label{thm:protocol-evidence-game-history}
Let $\kappa$ be an observation-based controller of the protocol presentation $U$, let $H=\Conf_f(U_\kappa)$, and let $\widehat H$ be the histories with the selected control commitments inserted. The existing compiler and Scott construction give maps
\begin{equation}\label{eq:admin-historycomparison}
\begin{tikzcd}[column sep=large]
H & \widehat H \arrow[l,two heads,"{\varepsilon}"']
              \arrow[r,"{\theta}","{\cong}"'] & \Conf_f(S_\kappa)
\end{tikzcd}
\end{equation}
Here $\varepsilon$ is surjective, erases only control commitments, and retains the causal order, physical occurrence witnesses and their declared labels. Independent controlled families lift to compiled cubes. The kernel of $S_\kappa$ is its reached Scott kernel and its unstraightening is $K(S_\kappa)$.

For an observer $o:H\to O$, interpret the game observer through $o\varepsilon\theta^{-1}$. Predicates transported through \eqref{eq:admin-historycomparison}, and evidence interpreted by the canonical descent equivalence of Lemma~\ref{lem:admin-evidence-descent}, have the same history-predicate, knowledge, possibility, and decision-descent interpretations. Any further specified family of physical evidence is transported by the same equivalence. Evidence-valued quantifiers use \eqref{eq:admin-desc-product}--\eqref{eq:admin-desc-adjunctions}. For predicates, specifically,
\begin{equation}\label{eq:admin-knowledge-transfer}
 \varepsilon^*\mathsf K_o(P)=\mathsf K_{o\varepsilon}(\varepsilon^*P),
 \qquad
 \varepsilon^*\mathsf{Poss}_o(P)=\mathsf{Poss}_{o\varepsilon}(\varepsilon^*P).
\end{equation}
Action modalities commute with predicate pullback and with the corresponding evidence descent when a physical action in $\widehat H$ includes the finite intervening control commitments, retaining its actual rule witness. The resulting equivalence transports the displayed evidence, policy and observation maps and their coherence equations.
\end{theorem}
\begin{proof}\vmode
\leavevmode
\begin{enumerate}[label=(\roman*)]
\item \textbf{Compiled history comparison.}
Theorems~\ref{thm:permission-configuration-classification} and~\ref{thm:compiler} give the history identification, surjective erasure, extension lifting from every representative, and cube lifting.
Lemma~\ref{lem:restriction} and Theorem~\ref{thm:scott-discrete-fibrations} give the strategy kernel and unstraightening.

\item \textbf{Physical evidence pullback.}
Each physical event retains its original witness and label; a positive slot changes neither.
A physical evidence family therefore pulls back with identity transport between representatives, exactly as in Lemma~\ref{lem:admin-evidence-descent}.

\item \textbf{Knowledge transfer.}
For knowledge, the forward implication in \eqref{eq:admin-knowledge-transfer} substitutes $\varepsilon\hat h'$ for an alternative history.
Conversely, take $h'$ with $o(h')=o(\varepsilon\hat h)$.
Surjectivity supplies a lift $\hat h'$; the right-hand universal statement implies $P(h')$.
This is existential elimination using surjectivity, not a global choice of lifts.

\item \textbf{Possibility transfer.}
For possibility, a compiled witness projects to a physical witness, and a physical witness has a lift.
These prove both equalities.

\item \textbf{Decision descent.}
They also show that constancy of $d\varepsilon$ on $o\varepsilon$-fibres is equivalent to constancy of $d$ on $o$-fibres; Theorem~\ref{thm:admin-descent} then transfers descent.

\item \textbf{Physical action relation.}
For each physical action $a$, define $\hat h\mathrel{\widehat a}\hat h'$ by finitely many selected commitments, one physical $a$ step, and finitely many selected commitments.
Every such extension projects to an $a$ extension.

\item \textbf{Quantification over lifted actions.}
Conversely the compiler proof lifts every permitted $a$ extension from every current compiled representative by preparing its missing causal-view slots.
Existential and universal quantification over these two extension relations therefore commute with inverse image of predicates along $\varepsilon$.
The finite commitment bound rules out replacing a physical response by an infinite sequence of these bookkeeping steps.

\item \textbf{Evidence-action carrier.}
For evidence, let $\widehat E_a$ be the set of these weak executions, with maps $\widehat s,\widehat t$ to $\widehat H$, and let $E_a$ be the physical $a$-extension witnesses with source and target $s,t$.
Erasure $\varepsilon_a$ induces
\begin{equation}\label{eq:admin-action-comparison}
 \lambda_a:\widehat E_a\twoheadrightarrow
           \widehat H\times_{H,s}E_a,\qquad
 \lambda_a(\widehat e)=(\widehat s(\widehat e),\varepsilon_a(\widehat e)),
 \qquad \varepsilon\widehat t=t\varepsilon_a.
\end{equation}
The source equation $\varepsilon\widehat s=s\varepsilon_a$ makes $\lambda_a$ well-typed.

\item \textbf{Surjective extension lifting.}
Given $(\widehat h,e)$ in the pullback, the preceding commitment completion followed by the selected physical witness $e$ constructs a preimage.
Hence $\lambda_a$ is surjective, expressing liftability from each compiled source.

\item \textbf{Descent of duplicate realizations.}
Apply the kernel-pair quotient and compatible-product constructions of Lemma~\ref{lem:admin-evidence-descent} to these witness fibres.
They identify duplicate compiled realizations of one physical witness, while keeping different physical witnesses distinct even when their public labels agree.
This proves the evidence-valued action comparison.

\item \textbf{Endpoint compatibility.}
The source and target equations used above are the commuting diagram
\[
\begin{tikzcd}[column sep=large]
\widehat H \arrow[d,"{\varepsilon}"'] &\widehat E_a
 \arrow[l,"{\widehat s}"'] \arrow[r,"{\widehat t}"] \arrow[d,"{\varepsilon_a}"]
 &\widehat H \arrow[d,"{\varepsilon}"]\\
H &E_a \arrow[l,"{s}"'] \arrow[r,"{t}"] &H .
\end{tikzcd}
\]
The additional surjection $\lambda_a$ in \eqref{eq:admin-action-comparison}, rather than commutativity alone, proves that an action can be lifted from every compiled source.
Quotienting its fibres retains the actual physical witness $e\in E_a$.

\item \textbf{Compound-rule soundness.}
Soundness of compound predicate rules follows by induction.

\item \textbf{Policy and mode transport.}
Physical evidence and policy maps have the canonical identity descent actions, and their naturality under extension is covered by Lemma~\ref{lem:admin-evidence-descent}.
Its equivalence transports adjunction units, counits, monad equations and mates; Theorem~\ref{thm:admin-matrices} gives their response-test interpretation.
The occurrence-mode locks are interpreted by Theorem~\ref{thm:admin-occurrence-lifts}.
Thus all interpretations have their stated judgement sorts and derived coherence.
\end{enumerate}
\end{proof}

The control arena exposes positive commitments and negative plant execution, as in Section~\ref{subsec:observations-distributed-strategies}. The response arena exposes negative requests and positive replies. These are different interfaces of the same witnessed operational construction. The proof uses the compiled commitment interface to implement authorization; it does not delete an enabled negative request from a fixed message arena. Every original protocol occurrence remains in $H$, and erasure removes only the explicitly compiled commitments.

The observation comparison is the triangle
\[
\begin{tikzcd}[column sep=large]
\widehat H \arrow[r,two heads,"{\varepsilon}"] \arrow[dr,"{\widehat o}"']
 &H \arrow[d,"{o}"]\\
 &O .
\end{tikzcd}
\]
Hence its kernel relation is computed by the actual equation $o\varepsilon(\hat h)=o\varepsilon(\hat k)$; no private compilation annotation becomes observable through this comparison.

\subsection{Calculus: resource programs and evidence modes}\label{subsec:administrative-process-calculus}
The process syntax remains the existing named and occurrence calculus. A selected resource command is a witnessed protocol step: its premises retain its typed inputs, persistent reads, owned occurrences, and local control. The following rule notation records the resource and evidence calculations for those steps~\cite{garg-linear,bowers-consumable}.

\begin{definition}[Resource types and programs]\label{def:resource-types-programs}
At a fixed decision view $v$ and parameter context $\Gamma$, a resource type is the family $A(\gamma,m)$ constructed above. The concentrated family $\mathbf m$, the unit, tensor, and internal hom have the meanings of Theorem~\ref{thm:parameter-substitution-closed-resources}. An owned context denotes their tensor, including its selected split. A pure judgement $\Gamma;D\vdash_{v,\mathrm{lin}}f:A$ names a constructed family map $D\to A$. Currying and application are the transpose and evaluation of that theorem. Constants name actual resource maps, and policy constants name selected elements of $\mathsf{Gen}_v(m,n)$.

A program judgement $\Gamma;D\vdash_v t:A$ has the constructors
\[
 \frac{\Gamma;D\vdash_{v,\mathrm{lin}} f:A}{\Gamma;D\vdash_v\mathsf{return}(f):A},
 \qquad
 \frac{\Gamma;D\vdash_v t:A\quad\Gamma;A\vdash_v u:B}
 {\Gamma;D\vdash_v\mathsf{bind}(t,u):B},
\]
\[
 \frac{\Gamma;D\vdash_v t:A}{\Gamma;Z\otimes_vD\vdash_v\mathsf{frame}_Z(t):Z\otimes_v A},
 \quad
 \frac{\Gamma;D\vdash_vt:A\quad\Gamma;D'\vdash_vu:B}
 {\Gamma;D\otimes_vD'\vdash_vt\parallel u:A\otimes_vB}.
\]
There is no structural copy or discard of an owned occurrence. Parameter substitution precomposes every family and constant and acts recursively on the finite term tree. Runtime calls still require their actual typed operands and command witnesses.
\end{definition}
\begin{definition}[Interpretation and historical evidence]\label{def:formal-resource-program-and-evidence-interpretation}
Recursion gives
\begin{equation}\label{eq:admin-term-interpretation}
 \llbracket t\rrbracket:D\longrightarrow T_vA.
\end{equation}
Return uses $\eta$, bind uses $\mu T_v\llbracket u\rrbracket\llbracket t\rrbracket$, framing uses strength, and parallel uses the equal framed orders of Lemma~\ref{lem:admin-parallelterms}. A policy generator inserts its retained path. A completed physical command retains its typed result and its history coordinate. Maps between such evidence families preserve that coordinate; concrete origin and receipt projections are constructed with the signed-grant instance. Historical public knowledge uses $\Box_J$ on Scott opens. The separately defined current-decision operators on subsets and evidence families use their displayed sum, product, and image formulas.
\end{definition}
\begin{proposition}[Resource implication computation]\label{prop:resource-implication-rules}
Pure abstraction and application satisfy beta and eta and commute with parameter substitution. Currying an effectful map $D\otimes A\to T_vB$ has codomain $A\multimap_v T_vB$.
\end{proposition}
\begin{proof}\vmode
\leavevmode
\begin{enumerate}[label=(\roman*)]
\item \textbf{Evaluation and currying.}
Evaluation and currying are the inverse componentwise maps in Theorem~\ref{thm:parameter-substitution-closed-resources}.

\item \textbf{Substitution and effectful codomain.}
Applying them at $\theta(\gamma')$ gives the substitution equation.
The codomain of the transposed effectful map follows by substituting $T_vB$ for $B$ in that bijection.
\end{enumerate}
\end{proof}
\begin{theorem}[Resource-program soundness and substitution]\label{thm:admin-soundness}
Every displayed program has the interpretation of its declared type, with
\begin{equation}\label{eq:admin-substitution}
 \llbracket t[\theta]\rrbracket=\theta^*\llbracket t\rrbracket.
\end{equation}
Policy sequencing satisfies the monad equations; parallel and framing satisfy the proved exchange and strength equations. The constructed evidence adjunctions satisfy the transposition and mate equations. Interpretation therefore descends through the typed congruence generated by these equations. Structural evidence and resource operations do not duplicate an atomic owned resource.
\end{theorem}
\begin{proof}\vmode
\leavevmode
\begin{enumerate}[label=(\roman*)]
\item \textbf{Typing induction.}
Induction on terms proves typing.
A generator has its declared path endpoints; return inserts an identity, bind concatenates paths, and framing and parallel retain their grade splits.

\item \textbf{Parameter substitution.}
The formulas are componentwise in the parameter, proving substitution by the same induction.

\item \textbf{Semantic equations.}
The monad, strength, and parallel equations are Theorem~\ref{thm:policy-evidence-adjunction-monad}, Proposition~\ref{prop:admin-strength}, and Lemma~\ref{lem:admin-parallelterms}; beta, eta, and mates are Proposition~\ref{prop:transposition-computation-substitution}.
Each generating equation consequently has equal interpretations, and equality is closed under the constructors.

\item \textbf{Obstruction to duplication.}
For an atom $r$, $\mathbf r([r])$ is inhabited but $(\mathbf r\otimes_v\mathbf r)([r])$ is empty.
A structural duplication would give a function from a singleton to that empty fibre.

\item \textbf{Operational resource invariants.}
Operational invariants of concrete rule constants are proved at their actual witness maps; structural soundness does not create additional resource occurrences.
\end{enumerate}
\end{proof}
\begin{corollary}[Transport to the constructed game histories]\label{cor:calculus-admin-transport}
The resource and evidence rules have the same interpretation on the compiled histories and response interfaces constructed in Theorem~\ref{thm:protocol-evidence-game-history}.
\end{corollary}
\begin{proof}\vmode
\leavevmode
\begin{enumerate}[label=(\roman*)]
\item \textbf{History-preserving descent.}
The evidence descent retains the physical witness, the inventory, and every path and provenance projection.

\item \textbf{Transport of rules and programs.}
It therefore transports each generator and each defining equation.
Induction on a program transports its composite interpretation; the response and occurrence comparisons are Theorems~\ref{thm:admin-matrices} and~\ref{thm:admin-occurrence-lifts}.
\end{enumerate}
\end{proof}
\subsubsection{The constructed evidence modes}\label{subsec:multimodal-structure}
Evidence transport uses $\Sigma_f\dashv f^*\dashv\Pi_f$ on actual sets of indices. Policy evidence uses $F_v\dashv U_v$ on the witnessed path category. The optional runtime-key parameter has the reader monad and the explicit law $T_vS\to ST_v$ constructed above. Occurrence maps give their reached strategy adjunctions; history-generated time gives its earlier delay construction. These modes are constructed over the specified protocol history and policy carriers. They neither enlarge the public command alphabet nor replace the Scott-open observation and behaviour-image security definitions.


Public frames describe observable assertions, evidence families retain their witnesses, and policy operations act with their specified history and resource indices. The comparison maps explain how this information changes under disclosure, hiding and execution. These are semantic constructions; the remaining task is to identify the computational presentations and finite procedures available for their realizers and comparisons.

\section{Effective presentations and complete finite verification}\label{sec:effective-foundations}
The preceding constructions establish semantic comparisons and construct their witnesses. We now ask how finite evidence can be represented computationally and when a comparison can be decided or a simulator synthesized. The configuration domains already have finite compact bases by Section~\ref{subsec:configuration-domains}. Effective presentations equip these bases and partial computations with enumerations while preserving the game operations. The final results distinguish procedures on supplied finite tables from the obstruction for unrestricted enumerated presentations. The domain-theoretic background is~\cite{effective-domains}.

\subsection{Effective configuration presentations}\label{subsec:effective-configurations}
We begin with presentations of compact evidence and the maps acting on it. These presentations will support the strategy constructors and the partial computations used by services, before the finite verification procedures are considered.



\begin{definition}[Effective domains and maps]\label{def:formal-effective-domains-and-maps}
Write $\mathsf{EDom}$ for the category of algebraic dcpos with an enumerable compact basis and computably enumerable basis order, and effective Scott-continuous maps. Empty and unpointed objects are included. A presentation of a map $f:X\to Y$ enumerates its \textbf{approximation relation}
\begin{equation}\label{eq:effective-approximable-map}
 R_f(k,\ell)\quad\Longleftrightarrow\quad \ell\leq f(k),
 \qquad k\in K_X,\quad \ell\in K_Y.
\end{equation}
Morphisms are the denoted maps; an enumeration is their computational structure and different enumerations may denote the same map. Products use pairs of compact codes and coproducts use tagged codes. A point $x\in X$ is effective when
$K_x=\{k\in K_X:k\leq x\}$ has an enumeration. No enumeration of $K_x$ is selected merely from a statement of its existence.
\end{definition}


\begin{definition}[Effective event presentation]\label{def:effective-event-presentation}
An \textbf{effective presentation} of an event structure $E$ consists of an enumerable code set $C_E\subseteq\mathbb N$, a surjection $\operatorname{decode}:C_E\to|E|$, and total computable realizers
\[
 \mathsf{eq},\mathsf{le}:C_E^2\to\{0,1\},\quad
 \mathsf{con}:C_E^*\to\{0,1\},\quad
 \mathsf{past}:C_E\to C_E^*,\quad
 \mathsf{pol}:C_E\to\{-,+\}.
\]
They realize respectively equality and causal order of decoded events, consistency of the decoded finite list, a list of all events in the decoded principal past, and polarity. Thus the equations specifying a realizer compare its decoded result with the corresponding event-structure operation. An \textbf{effective display} $\sigma:E\to A$ is accompanied by a computable $\widehat\sigma:C_E\to C_A$ with $\operatorname{decode}_A\widehat\sigma=\sigma\operatorname{decode}_E$. Codes may retain occurrence addresses, finite histories and matching witnesses. These algorithms are presentation data, not consequences of countability or of the finiteness of each individual past.
\end{definition}

For such a presentation, enumerate finite lists of event codes, remove repetitions, test consistency and test closure under the computed principal pasts. This enumerates exactly $K_E$, with decidable inclusion. The next proposition records how these presentations are obtained for the game operations used here.

\begin{proposition}[Effective game constructors]\label{prop:effective-game-constructors}
The effective presentations are preserved by parallel composition, polarity reversal, copycat, secured interaction and hiding. The occurrence constructions generated by finite coded commands, queues, banks and finite verification algorithms have these presentations. Their finite configuration maps for displays and hiding act uniformly on effective points.
\end{proposition}
\begin{proof}\vmode
\leavevmode
\begin{enumerate}[label=(\roman*)]
\item \textbf{Parallel and dual presentations.}
Parallel composition tags event codes and tests consistency componentwise; duality changes the sign algorithm.

\item \textbf{Effective copycat.}
For copycat, the two copies of an event's finite arena past contain its entire copycat past.
Closure under the arena and copy edges is a finite calculation in that set.
Causal comparison and consistency are consequently decidable by the same finite closure test.

\item \textbf{Enumeration of secured matchings.}
Enumerate finite configurations of both strategies, compute their displayed events, and retain pairs with equal middle-interface images.
Local injectivity determines the matching.
Compute the finite joint graph and test acyclicity.

\item \textbf{Principal interaction events.}
An interaction event is a rooted principal downset of such a graph, with its source occurrence identities retained.
This enumerates the events of Section~\ref{subsec:game-secured}.

\item \textbf{Equality and causal comparison.}
Finite labelled histories have decidable equality and inclusion, and the principal past of an event is the finite list of rooted principal histories inside its code.

\item \textbf{Consistency test.}
For joint consistency, form the union of the finitely many proposed histories.
Check that the projections give configurations of both sources, that the matching is well defined, that its joint order is acyclic, and that each proposed history is the principal history of its root in that union.
These tests are finite.

\item \textbf{Completeness of the test.}
Success gives a common matching witness.
Conversely, if a common witness exists, the union is an ideal in it by Lemma~\ref{lem:matching-past} and passes these tests.
Thus the test decides consistency without searching over possible larger witnesses.

\item \textbf{Hiding and finite configuration maps.}
Retain the interaction events with visible roots, together with their full matching-history codes.
Restrict the computed principal past to those visible roots; consistency is inherited.
On finite configurations, visibility is the map $v$ of Proposition~\ref{prop:finite-secured-witnesses}; its matching presentation is the finite pruning operation.
Displays compute finite images.

\item \textbf{Continuous extension.}
These operations preserve directed unions, and their finite images enumerate the compact approximations to the resulting configurations.

\item \textbf{Finite operational records.}
For a finite core declaration, reachability, surviving events, causal closure and forbidden-cohort tests are finite calculations in the declaration.
The scoped construction of Section~\ref{subsec:pc-named-proof} uses finite proof trees, virtual addresses, binding pointers and finite environments.
A coded medium or service rule adds its premises, selected suppliers, residual store and result; a completed partial evaluation retains its finite computation witness.

\item \textbf{Record checking and resource charts.}
Checking a proposed finite record and its principal histories is therefore a finite computation.
The resource charts retain these records when independent steps are exchanged.
The later protocol and service constructions specify their particular record checkers; the preceding interaction and hiding algorithms then compute their composed ordinary presentations.
\end{enumerate}
\end{proof}

\begin{construction}[Typed finite-code procedures]\vmode\label{con:effective-procedure-contracts}
\leavevmode
\begin{enumerate}[label=(\roman*)]
\item \textbf{Valid input codes.}
Codes in the following procedures are valid codes of the effective
presentations; an enumerable code set is not thereby asserted to
have decidable membership.

\item \textbf{Configuration test.}
$\mathsf{CfgTest}_E:C_E^*\to\{0,1\}$ removes decoded
duplicates using $\mathsf{eq}$, applies $\mathsf{con}$, and for
each retained code tests membership of every element of
$\mathsf{past}(c)$ in the list.
It returns $1$ exactly on finite
configurations.

\item \textbf{Secured matching procedure.}
$\mathsf{Match}_{S,T}$ takes two finite configuration lists,
tests equality of their displayed middle-interface images, and
constructs the unique locally injective matching.
A finite
topological-sort test decides securedness.
Its output is either
failure or the finite joint graph.

\item \textbf{Principal-history and hiding procedures.}
$\mathsf{Principal}(m,r)$ computes the predecessor closure of
the indicated root in a finite secured graph and returns its rooted
code.
$\mathsf{Hide}$ retains visible-root codes with their complete
matching records.

\item \textbf{Horizontal map computation.}
For horizontal maps first compute component
images and then $\mathsf{Principal}$, exactly as in
Lemma~\ref{lem:horizontal-principal-map}.

\item \textbf{Joint consistency test.}
$\mathsf{ConTest}$ for a finite list of interaction-event codes
forms their union, tests component configurations and securedness,
and checks that each proposed principal code is the root past in
that union.
It does not enumerate larger candidate extensions.

\item \textbf{Receptive-lift search.}
A required negative lift enumerates source event codes and
tests the display equation and enabledness at the given finite
configuration.
Receptivity proves that this search terminates with
one decoded event.
Duplicate codes do not affect that event.

\item \textbf{Termination and correctness.}
The first four procedures terminate by finite-list recursion and
finite graph closure.
Their correctness is the configuration,
matching and principal-history characterization already proved.
The fifth is an enumerative search with termination supplied by the
negative-extension theorem, not a finite a priori search bound.

\item \textbf{Operational certificate checking.}
For operational records, a certificate consists of the finite exposure
derivation, selected medium record and, for a completed partial
evaluation, its finite machine computation.
Checking each inference
and each computation step is finite in the declared coded instance.
Lemma~\ref{lem:retained-reading-bijection} proves that forgetting proof
bureaucracy preserves its denoted record.

\item \textbf{Use of the supplied presentation.}
Every procedure above uses
the supplied decoded-equality tests; replacing an input code by an
equal decoded code therefore gives an equal decoded output.
\end{enumerate}
\end{construction}

This constructs the presentations of the coded instances, while the set-valued interpretation of arbitrary external arenas and value sets remains available. The algorithms come from their occurrence records and finite verification procedures. Neither countability alone nor the compact closed equations select them.

\subsection{The effective strategy bicategory and its comparison}\label{subsec:effective-bicategory}
The computational presentations refine the two categorical levels of Section~\ref{subsec:strategy-levels}. We retain the finite-history representatives until the comparison maps have been constructed.

\begin{definition}[Effectively presented strategies]\label{def:effective-strategy-bicategory}
The effective hom categories and their objects are specified as follows.
\leavevmode
\begin{enumerate}[label=(\roman*)]
\item An object is an arena with an effective event presentation.
\item An object of $\mathbb{CG}_{\mathrm{eff}}(A,B)$ is an ordinary receptive, courteous display $\sigma:S\to A^\ast\parallel B$ whose source has an effective presentation and whose display map computes event codes.
\item A morphism $h:\sigma\Rightarrow\sigma'$ is a computable total event-structure map $h:S\to S'$ satisfying $\sigma'h=\sigma$. The morphism is the denoted map; different programs computing that map present the same morphism. Identities and vertical composition are computed on event codes.
\end{enumerate}
\end{definition}

\begin{theorem}[Effective refinement of the strategy bicategory]\label{thm:effective-game-refinement}
The constructors of Proposition~\ref{prop:effective-game-constructors} extend to a bicategory $\mathbb{CG}_{\mathrm{eff}}$ and a locally faithful forgetful pseudofunctor
\[
U:\mathbb{CG}_{\mathrm{eff}}\longrightarrow\mathbb{CG}.
\]
Its composition and identity comparisons, parallel and symmetry comparisons, and compact cups and caps have computable representatives. Consequently
\[
\mathsf{CG}_{\mathrm{eff}}(A,B)
 =\operatorname{Ob}\mathbb{CG}_{\mathrm{eff}}(A,B)/\cong_{\mathrm{eff}}
\]
defines a compact closed category, and $U$ induces a strong symmetric monoidal functor $\overline U:\mathsf{CG}_{\mathrm{eff}}\to\mathsf{CG}$ preserving the specified compact structure. Here $\cong_{\mathrm{eff}}$ means isomorphism in the effective hom category.
\end{theorem}
\begin{proof}\vmode
\leavevmode
\begin{enumerate}[label=(\roman*)]
\item \textbf{Horizontal map algorithm.}
Proposition~\ref{prop:effective-game-constructors} constructs the source and display of a composite.
For computable display maps $h:S\to S'$ and $k:T\to T'$, apply $h,k$ to a finite matching history and take the principal past of its mapped root in the resulting matching.
This is a finite algorithm.

\item \textbf{Securedness of mapped histories.}
To verify securedness, observe that a total event map reflects causal order within each configuration: if $h(s)\leq h(t)$ with $s,t\in x$, then $h([t])$ is a configuration containing $h(t)$; local injectivity on $x$ forces $s\in[t]$.
Hence every target causal edge lifts to a source causal relation.
The image matching is locally injective and acyclic.
Its principal history is therefore defined; taking that principal history is essential because a display map may forget a waiting dependency.

\item \textbf{Functoriality.}
This algorithm is the map induced by secured interaction, followed by visible restriction.
The universal property used in Section~\ref{subsec:strategy-levels} proves that it preserves configurations and gives the required map of displays.
It respects identities and vertical composition by that same uniqueness, so horizontal composition is a functor and interchange holds.

\item \textbf{Effective associator.}
Expand a nested principal history to the finite joint matching graph of its constituent strategies.
Take the principal past of the ultimate visible root and regroup its component histories in the other bracketing.
This computes the usual associator and its inverse.

\item \textbf{Configuration-level unit adjustment.}
For the unit comparisons, remove the copycat copies and perform the finite configuration adjustments determined by their matching boundary.
These adjustments consist of maximal positive removals and the unique receptive negative extensions, as in~\cite[Section~4.4]{ccrw}.

\item \textbf{Computing the adjustments.}
A removal is computed in the current finite configuration.
A required negative extension is found by enumerating source events and testing its display and enabledness; receptivity proves termination and uniqueness.

\item \textbf{Unitor on events.}
Apply this configuration computation to an event's visible principal past and return the unique maximal event of its image.
This computes the canonical unitor on event codes.

\item \textbf{Inverse unitor.}
Its inverse is computable by enumerating domain event codes until their image equals the given code, since the unitor is bijective and event equality is decidable.

\item \textbf{Tensor and compact structure.}
Rebracketing and permutation of parallel tags are computable.
The copycat construction computes the cups and caps by changing the boundary assignments.
The snake comparisons are the preceding unit comparisons read on those assignments.
Thus the tensor, symmetry, dual objects and compact comparisons lift those of~\cite[Section~5]{ccrw}.

\item \textbf{Coherence.}
All comparison maps are the underlying ordinary comparisons of Section~\ref{subsec:strategy-levels}.
In a pentagon every route expands to the same finite joint history and regroups it with the same visible root.
Copycat removal gives the corresponding triangle and snake comparisons.
More generally the ordinary coherence equations of~\cite[Sections~4--5]{ccrw} hold between the constructed maps.

\item \textbf{Forgetful pseudofunctor.}
Since $U$ forgets algorithms but retains their denoted event maps, it is locally faithful, so these equations also hold in the effective hom categories.
The canonical comparisons from the chosen coded constructors to the ordinary constructors give the stated pseudofunctor.

\item \textbf{Quotient category.}
Quotienting by its computable invertible 2-cells gives the asserted category and compact-preserving functor.
\end{enumerate}
\end{proof}

Writing $U_{A,B}$ for the hom functors, compatibility with composition is the natural comparison
\[
\begin{tikzcd}[column sep=large,row sep=large]
\mathbb{CG}_{\mathrm{eff}}(B,C)\times\mathbb{CG}_{\mathrm{eff}}(A,B)
 \arrow[r,"\odot_{\mathrm{eff}}"] \arrow[d,"U_{B,C}\times U_{A,B}"']
 &\mathbb{CG}_{\mathrm{eff}}(A,C) \arrow[d,"U_{A,C}"]\\
\mathbb{CG}(UB,UC)\times\mathbb{CG}(UA,UB)
 \arrow[r,"\odot"']
 &\mathbb{CG}(UA,UC).
\end{tikzcd}
\]
The square commutes up to the specified natural isomorphism; on the quotient categories it commutes as an equality. Tensor has the corresponding comparison. Presentations remain part of the effective objects and strategies; no identification by an arbitrary uncomputed display isomorphism is made.

\begin{corollary}[Effective interpretation of typed networks]\label{cor:effective-process-interpretation}
Let $\mathsf{Net}_{\mathrm{code}}$ be the typed expression category of the paragraph ``\nameref{par:typed-network-interpretation}'' on the coded component presentations of Proposition~\ref{prop:effective-game-constructors}. Forgetting the code annotations gives $J:\mathsf{Net}_{\mathrm{code}}\to\mathsf{Net}$. There is a compact-preserving interpretation $\mathcal I_{\mathrm{eff}}$ and a natural comparison
\[
\begin{tikzcd}[column sep=large,row sep=large]
\mathsf{Net}_{\mathrm{code}}
 \arrow[r,"\mathcal I_{\mathrm{eff}}"] \arrow[d,"J"']
 &\mathsf{CG}_{\mathrm{eff}} \arrow[d,"\overline U"]\\
\mathsf{Net} \arrow[r,"\mathcal I"'] &\mathsf{CG}.
\end{tikzcd}
\qquad
\overline U\mathcal I_{\mathrm{eff}}\cong\mathcal I J .
\]
In particular, for a coded term $P$, forgetting its effective interpretation recovers its ordinary displayed interpretation up to the constructed display isomorphism.
\end{corollary}
\begin{proof}\vmode
\leavevmode
\begin{enumerate}[label=(\roman*)]
\item \textbf{Effective component interpretation.}
On a component, use its constructed effective display; forgetting its codes recovers that component's ordinary display.
Forwarding and the cup and cap use the computed copycat presentation.

\item \textbf{Induction on network expressions.}
Induct on the finite expression tree.
Tensor uses tagged parallel and connection uses finite secured interaction and hiding.

\item \textbf{Structural comparison and coherence.}
The comparisons of Theorem~\ref{thm:effective-game-refinement} extend the component comparisons at each node.
Its coherence proves the symmetric monoidal and compact equations, so the interpretation descends to the structural quotient defining $\mathsf{Net}_{\mathrm{code}}$.

\item \textbf{Naturality and named elaboration.}
The same induction proves naturality of the comparison.
On named components, the finite elaboration records of Section~\ref{subsec:pc-named-proof} supply their codes; the later medium and service interpretations use their displayed finite record checkers.
\end{enumerate}
\end{proof}

\subsection{Uniform maps and their computational structure}\label{subsec:effective-maps}
Effective computation must act uniformly on the compact evidence of its input. Using the approximation relations and effective points of Definition~\ref{def:formal-effective-domains-and-maps}, we now construct the map operations needed by the subsequent service and process interpretations.

\begin{proposition}[Uniform action on names]\label{prop:effective-name-action}
An enumeration of $R_f$ transforms any enumeration of $K_x$ uniformly into one of $K_{f(x)}$. Identities and composition are effective, with
\begin{align}
 K_{f(x)}&=\{\ell:\exists k\in K_x\ R_f(k,\ell)\},\label{eq:effective-point-action}\\
 R_{gf}(k,m)&\Longleftrightarrow
     \exists\ell\in K_Y\bigl(R_f(k,\ell)\wedge R_g(\ell,m)\bigr).
     \label{eq:effective-composition}
\end{align}
\end{proposition}
\begin{proof}\vmode
\leavevmode
\begin{enumerate}[label=(\roman*)]
\item \textbf{Compact approximation criterion.}
Continuity gives $f(x)=\bigvee_{k\leq x}f(k)$.
Compactness of $\ell$ gives~\eqref{eq:effective-point-action}; a dovetailed enumeration of the two relations realizes it.

\item \textbf{Composition relation.}
Apply the same argument to $g(f(k))$ to obtain~\eqref{eq:effective-composition}.

\item \textbf{Identities and associativity.}
The identity relation is the enumerable basis order.
Substitution of existential witnesses proves associativity and the identity equations; the corresponding enumerators are obtained by finite pairing and dovetailing.

\item \textbf{Products and coproducts.}
Product and coproduct maps use the same componentwise procedures.
\end{enumerate}
\end{proof}

\begin{definition}[Enumerator-name representation]\label{def:formal-enumerator-name-representation}
Writing $\mathsf N(X)$ for enumerator names of effective points and $\nu_X$ for their denotation, the constructed transformer $\widetilde f$ gives
\[
\begin{tikzcd}[column sep=large]
\mathsf N(X)\arrow[r,"\widetilde f"]\arrow[d,"\nu_X"']
 &\mathsf N(Y)\arrow[d,"\nu_Y"]\\
X\arrow[r,"f"']&Y .
\end{tikzcd}
\]
\end{definition}
The lower objects remain the complete dcpos. Effective points alone need not contain the supremum of a nonuniform directed family. Likewise preservation of effective points alone is insufficient to specify an algorithm: every individual natural number is effective, but an arbitrary function on natural numbers need not have a uniform name transformer. Retaining the computational structure of a presentation makes the distinction precise; see~\cite[Section~3]{effective-dominances}.

\subsection{Configuration domains of effective displays}\label{subsec:effective-display-domains}
The effective event presentation must recover the configuration domain used by observation. The following comparison identifies its compact approximations and the maps induced by effective displays.

\begin{definition}[Effective event-map category]\label{def:formal-effective-event-map-category}
Let $\mathsf{ES}_{\mathrm{eff}}$ have effectively presented event structures and computable total event-structure maps. Polarity is forgotten in this category. The following construction connects their event presentations to the domain category just defined.
\end{definition}

\begin{proposition}[Effective configuration comparison]\label{prop:effective-configuration-comparison}
\leavevmode
\begin{enumerate}[label=(\roman*)]
\item Direct image gives a functor
\[
\mathsf D:\mathsf{ES}_{\mathrm{eff}}\longrightarrow\mathsf{EDom},
\qquad
\mathsf D(h)(x)=h[x],
\]
with effective canonical isomorphisms
$\mathsf D(E\parallel F)\cong\mathsf D(E)\times\mathsf D(F)$,
$\mathsf D(I)\cong1$, and $\mathsf D(E^\ast)\cong\mathsf D(E)$.
\item An effective display $\sigma:S\to A^\ast\parallel B$ determines the map
$d_\sigma:\mathsf D(S)\to\mathsf D(A)\times\mathsf D(B)$, or equivalently the witnessed span
\[
\begin{tikzcd}[column sep=large]
\mathsf D(A)&
\mathsf D(S)\arrow[l,"\pi_A d_\sigma"']\arrow[r,"\pi_B d_\sigma"]
&\mathsf D(B).
\end{tikzcd}
\]
On display maps this gives a functor
\[
\mathbb{CG}_{\mathrm{eff}}(A,B)
 \longrightarrow
\mathsf{EDom}/\bigl(\mathsf D(A)\times\mathsf D(B)\bigr).
\]
\item Interaction and hiding have the effective configuration comparison
\[
\begin{tikzcd}[column sep=large,row sep=large]
\mathsf D(T\mathbin{\bowtie_B}S)
 \arrow[r,"\mathsf{vis}"] \arrow[d,"\cong"']
 &\mathsf D(T\odot S)\arrow[d,"\cong"]\\
\operatorname{Idl}\mathsf M_f(\sigma,\tau)
 \arrow[r,"\operatorname{Idl}(\mathsf{prune})"']
 &\operatorname{Idl}\mathsf M_f^{\mathrm{vis}}(\sigma,\tau).
\end{tikzcd}
\]
Here $\mathsf{vis}$ restricts a configuration to its visible events, and the lower bases are those of Proposition~\ref{prop:finite-secured-witnesses}.
\end{enumerate}
\end{proposition}
\begin{proof}\vmode
\leavevmode
\begin{enumerate}[label=(\roman*)]
\item \textbf{Effective configuration map.}
A total event map sends configurations to configurations, is monotone on their inclusions, and preserves directed unions by direct image.
Its finite images are finite configurations.
Consequently
\[
R_{\mathsf D(h)}(k,\ell)
 \quad\Longleftrightarrow\quad
\ell\subseteq h[k]
\]
is decidable on compact codes, hence enumerable.

\item \textbf{Functoriality.}
Direct images preserve identities and composition.

\item \textbf{Tensor and polarity.}
Splitting and joining the tags of a finite parallel configuration computes the product isomorphism; the empty arena has its single empty configuration.
Reversing polarity leaves the inclusion domain unchanged.

\item \textbf{Extension to all configurations.}
These maps extend to all configurations by directed unions.

\item \textbf{Display comparison.}
Apply (i) to the display and then the product isomorphism.
For $h:\sigma\Rightarrow\sigma'$, the display equation gives
\[
\begin{tikzcd}[column sep=large,row sep=large]
\mathsf D(S)\arrow[r,"\mathsf D(h)"]\arrow[d,"d_\sigma"']
 &\mathsf D(S')\arrow[d,"d_{\sigma'}"]\\
\mathsf D(A)\times\mathsf D(B)
 \arrow[r,equal]
 &\mathsf D(A)\times\mathsf D(B).
\end{tikzcd}
\]
Functoriality in the slice follows from (i).

\item \textbf{Finite secured basis.}
Proposition~\ref{prop:finite-secured-witnesses} gives the isomorphisms and the commuting square on finite configurations.
Proposition~\ref{prop:effective-game-constructors} computes those finite maps and their inverses.

\item \textbf{Ideal-completion comparison.}
Taking ideal completions gives the displayed square, because visibility and finite pruning are monotone and their extensions preserve directed unions.
Proposition~\ref{prop:effective-name-action} supplies the uniform action on effective points.
\end{enumerate}
\end{proof}

A strategy therefore retains configuration witnesses over both interfaces; the resulting span need not be the graph of a function $\mathsf D(A)\to\mathsf D(B)$. Composition must first test secured matching and then project to the visible configuration. An unrestricted pullback of interface configuration domains does not perform that test.

The two functors constructed here consequently have different domains: one forgets computational presentations of strategies, while the other sends event maps to configuration maps. The displayed homwise comparison connects them. Polarity reversal leaves the inclusion domain unchanged under the configuration functor, but sends the Scott configuration category of Theorem~\ref{thm:eventscott} to its opposite.

\subsection{Effective open domains and the lift classifier}\label{subsec:effective-lifting}
A partial computation needs a domain on which it returns. We now classify effective open domains by the lift construction, providing the semantic operation used for potentially nonreturning service updates.

\begin{definition}[Effective open domains and lifting]\label{def:formal-effective-open-domains-and-lifting}
An \textbf{effective open subdomain} $U\hookrightarrow X$ is presented by an enumeration of compact generators with
\[
 U=\bigcup_{i\in W}\mathord\uparrow k_i.
\]
Its compact basis consists of the compact $k\in K_X$ for which some enumerated $k_i\leq k$; a code retains this finite positive witness. Repeated witnessed codes denote the same compact and are compared by their underlying basis order. Directed joins are inherited from $X$. This supplies an algebraic dcpo presentation of $U$, including the empty case. It is necessary to include unpointed objects: an open part of a pointed configuration domain need not have a least configuration.

Construct $\mathsf{Lift}(X)$ by ideal completion of the basis
\begin{equation}\label{eq:effective-lift-basis}
 K_X^\bot=\{\bot\}+\{\mathsf{ret}(k):k\in K_X\},\qquad
 \bot\leq\mathsf{ret}(k),\quad
 \mathsf{ret}(k)\leq\mathsf{ret}(\ell)\Longleftrightarrow k\leq\ell.
\end{equation}
The unit $\eta_X$ sends $x$ to the ideal generated by its returned compact approximations. If $X=\mathsf D(E)$, the event presentation adds a fresh return event before every original event; a finite set is consistent exactly when its old-event part is consistent. Its empty configuration is nonreturn; the singleton return configuration represents $\eta_X(\varnothing)$. Thus return of the empty configuration and nonreturn are distinct.
\end{definition}

\begin{theorem}[Effective partial-map classifier]\label{thm:effective-partial-classifier}
Effective open inclusions are stable under pullback along effective maps and under composition. The map $\eta_Y:Y\hookrightarrow\mathsf{Lift}(Y)$ classifies partial maps with such domains: for an effective map $f:U\to Y$ there is a unique effective map $\widehat f:X\to\mathsf{Lift}(Y)$ making the following square a pullback:
\[
\begin{tikzcd}[column sep=large]
U\arrow[r,"f"]\arrow[d,hook,"j"']\arrow[dr,phantom,"\lrcorner",very near start]
 &Y\arrow[d,hook,"\eta_Y"]\\
X\arrow[r,"\widehat f"']&\mathsf{Lift}(Y).
\end{tikzcd}
\]
Consequently, naturally in $X,Y$,
\begin{equation}\label{eq:effective-partial-bijection}
 \mathbf{Par}_{\rm eo}(\mathsf{EDom})(X,Y)
 \cong\mathsf{EDom}(X,\mathsf{Lift}(Y)).
\end{equation}
Here partial maps are spans with an effective open left leg, identified by isomorphism of their domains.
\end{theorem}
\begin{proof}\vmode
\leavevmode
\begin{enumerate}[label=(\roman*)]
\item \textbf{Effective inverse-image domain.}
For $g:Z\to X$, the preimage of $U$ is generated by compact $c$ for which $R_g(c,k_i)$ holds for an enumerated generator $k_i$.
Continuity and compactness prove equality with $g^{-1}U$.
This gives its enumeration and the pullback presentation.

\item \textbf{Open subdomains.}
An open in an open subdomain is enumerated by its compact codes regarded in the ambient basis, proving closure under composition.

\item \textbf{Construction of the partial classifier.}
Construct $\widehat f$ on compact approximations by always enumerating $\bot$ and by the rule
\begin{equation}\label{eq:effective-classifier-enumeration}
 R_{\widehat f}(k,\mathsf{ret}(\ell))
 \Longleftrightarrow
 \exists u\in K_U\bigl(u\leq k\ \wedge\ R_f(u,\ell)\bigr).
\end{equation}
For $x\in U$, the compacts of $U$ below $x$ are cofinal among the compacts of $X$ below $x$: combine a compact witnessing openness with any other compact below $x$.

\item \textbf{Returned values and nonreturn.}
Thus the rule yields exactly $\eta_Yf(x)$.
Outside $U$ it yields only the bottom approximation.

\item \textbf{Continuity and effectiveness.}
The rule also proves continuity and effectiveness, without a test for nonmembership in $U$.

\item \textbf{Recovery of the return domain.}
Conversely, for $h:X\to\mathsf{Lift}(Y)$ enumerate the compact $k$ for which $R_h(k,\mathsf{ret}(\ell))$ holds for some $\ell$.
Their upper closures form its return domain.

\item \textbf{Inverse correspondence.}
Returned approximations give an effective map on that domain, and the preceding rule recovers $h$.
It also proves the asserted pullback and uniqueness, including $Y=\varnothing$.

\item \textbf{Naturality.}
Pulling back once more or postcomposing on the returned values proves naturality.
\end{enumerate}
\end{proof}

\begin{definition}[Effective open classifier]\label{def:formal-effective-open-classifier}
For $Y=1$, put $\Sigma=\mathsf{Lift}(1)$ and $\top=\eta_1$. The characteristic map of an effective open is $\chi_U=\widehat{!_U}$. The classifier and its pullback stability are displayed by
\[
\begin{tikzcd}
g^{-1}U\arrow[r]\arrow[d,hook]\arrow[dr,phantom,"\lrcorner",very near start]
 &U\arrow[r]\arrow[d,hook]&1\arrow[d,hook,"\top"]\\
Z\arrow[r,"g"']&X\arrow[r,"\chi_U"']&\Sigma .
\end{tikzcd}
\]
This is the effective open dominance used here: truth has the displayed stable class of pullbacks, closed under composition. The enumerators supply its computational content. The Scott-open universal property of lifting is developed in~\cite[Theorem~30]{effective-lifting}; the explicit basis construction above proves the effective version needed here. It does not infer a choice of enumerators from an unstructured assertion of semidecidability.
\end{definition}

\begin{definition}[Lift operations and Kleisli arrows]\label{def:lift-operations}
For $X,Y\in\mathsf{EDom}$, define the \textbf{lift multiplication} $\mu_X:\mathsf{Lift}^2X\to\mathsf{Lift}X$ and \textbf{strength} $\mathsf{st}_{X,Y}:X\times\mathsf{Lift}Y\to\mathsf{Lift}(X\times Y)$ by
\[
\begin{aligned}
\mu_X(\bot)&=\bot,&\mu_X(\mathsf{ret}(\bot))&=\bot,&
\mu_X(\mathsf{ret}(\mathsf{ret}(k)))&=\mathsf{ret}(k),\\
\mathsf{st}_{X,Y}(k,\bot)&=\bot,&
\mathsf{st}_{X,Y}(k,\mathsf{ret}(\ell))&=\mathsf{ret}(k,\ell),
\end{aligned}
\]
extending by directed joins. The \textbf{bind operation} is $u\mathbin{\mathsf{bind}}h=\mu_Y\mathsf{Lift}(h)(u)$ for $h:X\to\mathsf{Lift}Y$. For a monad $(T,\eta,\mu)$ on $C$, its \textbf{Kleisli category} $\operatorname{Kl}(T)$ has objects those of $C$ and $\operatorname{Kl}(T)(X,Y)=C(X,TY)$, identity $\eta_X$ and composition $g\mathbin{\star}f=\mu_ZT(g)f$. The proposition proves the monad laws for this particular lift construction.
\end{definition}

\begin{proposition}[Constructed strong lift monad]\label{prop:effective-lift-monad}
The classifier determines a strong monad on $\mathsf{EDom}$. On compact generators its multiplication and strength are
\begin{align*}
 \mu_X(\bot)&=\bot,&
 \mu_X(\mathsf{ret}(\bot))&=\bot,&
 \mu_X(\mathsf{ret}(\mathsf{ret}(k)))&=\mathsf{ret}(k),\\
 \mathsf{st}_{X,Y}(k,\bot)&=\bot,&
 \mathsf{st}_{X,Y}(k,\mathsf{ret}(\ell))&=\mathsf{ret}(k,\ell).
\end{align*}
Continuous extension gives the maps on all ideals. Write $u\mathbin{\mathsf{bind}}h=\mu\,\mathsf{Lift}(h)(u)$. Under~\eqref{eq:effective-partial-bijection}, Kleisli composition is composition by pullback of partial domains.
\end{proposition}
\begin{proof}\vmode
\leavevmode
\begin{enumerate}[label=(\roman*)]
\item \textbf{Effective generator maps.}
The generator rules are monotone and effective; extend by directed joins.

\item \textbf{Naturality.}
Naturality follows on compact generators and hence on their joins.

\item \textbf{Monad and strength laws.}
Both unit laws remove one return layer and associativity removes two layers, propagating the first nonreturn.
The strength equations retain the same product coordinates and remove the same layers, so hold on generators and then everywhere.

\item \textbf{Partial-map composition.}
For partial maps $f:U\to Y$ and $g:V\to Z$, the returned domain of $\mu_Z\mathsf{Lift}(\widehat g)\widehat f$ is $U\times_YV$ and its returned value is $g(f(x))$.
The classifier's uniqueness proves the composition assertion and its coherence.
\end{enumerate}
\end{proof}
\[
\begin{tikzcd}[column sep=large]
\mathsf{Lift}^3X\arrow[r,"\mathsf{Lift}(\mu_X)"]\arrow[d,"\mu_{\mathsf{Lift}X}"']
 &\mathsf{Lift}^2X\arrow[d,"\mu_X"]\\
\mathsf{Lift}^2X\arrow[r,"\mu_X"']&\mathsf{Lift}X
\end{tikzcd}
\qquad
\begin{tikzcd}
U\times_YV\arrow[r]\arrow[d,hook]\arrow[dr,phantom,"\lrcorner",very near start]
 &V\arrow[d,hook]\\
U\arrow[r,"f"']&Y .
\end{tikzcd}
\]

\subsection{Effective partial controllers and return domains}
The lift classifier can now be applied to controller updates. Their return domains determine the successful transition graph, whose completed paths retain the finite witnesses needed by the execution interpretation.

\subsubsection{Effective partial services and completed paths}\label{subsec:effective-service-paths}
\begin{definition}[Effective partial service]\label{def:formal-effective-partial-service}
The total service instances of Section~\ref{subsec:crypto-games} reply to invalid requests with the explicit tag $\crFail$. A computation which may instead fail to return is an effective map
\[
 \widehat\delta:S\times Q\longrightarrow\mathsf{Lift}(S\times R).
\]
Let $U_\delta\hookrightarrow S\times Q$ be its classified return domain. Its successful-step graph has objects $S$, edges $U_\delta$, source $(s,q)\mapsto s$, and target the returned state. Its edge decoration is the returned reply. The free path category of this graph carries the input request word and output reply word as functors; each generator retains its actual return-domain witness. This retains the enabling witness instead of asserting a total action of every public request at every state.
\end{definition}

\begin{proposition}[Partial simulation and return-domain reflection]\label{prop:effective-partial-simulation}
For an effective state map $a:S\to S'$, the equation
\begin{equation}\label{eq:effective-service-square}
 \mathsf{Lift}(a\times1_R)\widehat\delta
   =\widehat\delta'(a\times1_Q)
\end{equation}
is equivalent to an output-preserving map $\bar a:U_\delta\to U_{\delta'}$ of the successful-step graphs whose source-domain square is a pullback:
\[
\begin{tikzcd}[column sep=large]
U_\delta\arrow[r,"\bar a"]\arrow[d,hook]
 \arrow[dr,phantom,"\lrcorner",very near start]
 &U_{\delta'}\arrow[d,hook]\\
S\times Q\arrow[r,"a\times1_Q"']&S'\times Q .
\end{tikzcd}
\]
The same property holds for completed finite request words. Their path functor preserves returned states and output words, and it preserves and reflects the domain of completed execution.
\end{proposition}
\begin{proof}\vmode
\leavevmode
\begin{enumerate}[label=(\roman*)]
\item \textbf{Return-domain pullback.}
Pull back the return inclusion along both sides of~\eqref{eq:effective-service-square}.
Theorem~\ref{thm:effective-partial-classifier} identifies the resulting domain with both $U_\delta$ and $(a\times1_Q)^*U_{\delta'}$.

\item \textbf{Transition comparison and converse.}
Equality of returned values gives the target and output equations.
Conversely the pullback identifies the return domains, and the graph equations identify their values; classifier uniqueness gives~\eqref{eq:effective-service-square}.

\item \textbf{Iteration over request words.}
For words put $\widehat\delta^*(s,\eps)=\eta(s,\eps)$ and extend by Kleisli composition, concatenating the next reply.
Naturality of multiplication and strength proves the corresponding square by induction on word length.

\item \textbf{Completed-path carrier.}
Applying the classifier once more gives
\[
\begin{tikzcd}[column sep=large]
U_{\delta^*}\arrow[r,"\bar a^*"]\arrow[d,hook]
 \arrow[dr,phantom,"\lrcorner",very near start]
 &U_{\delta'^*}\arrow[d,hook]\\
S\times Q^*\arrow[r,"a\times1_{Q^*}"']&S'\times Q^* .
\end{tikzcd}
\]
A completed path is uniquely determined by its initial state and request word on this deterministic domain.
Thus $U_{\delta^*}$ is precisely the arrow carrier of the free path category of the successful-step graph.

\item \textbf{Concatenation and effectiveness.}
Path concatenation is the same Kleisli recursion, so $\bar a^*$ is the asserted functor.
Its effective presentation follows by finite composition and enumeration over word lengths.
\end{enumerate}
\end{proof}

The pullback condition here is derived from the computational simulation equation. An arbitrary output-preserving functor on completed paths would only preserve successful runs and would not by itself reflect their existence. On the total service fragment, the return-domain square is automatically a pullback and the equation reduces to the strict service simulation of Definition~\ref{def:formal-total-services-and-their-simulation-square}, whose word extension is Proposition~\ref{prop:controller-execution-kernel-transport}.

When replies include explicit failure, the computational codomain is $\mathsf{Lift}(S\times(R+\{\crFail\}))$. It distinguishes nonreturn, a returned failure, and a successful reply. Definition~\ref{def:responsegame} accepts a request even when its return fibre is empty, so the initialized response strategy preserves this distinction.

Proposition~\ref{prop:effective-finite-return} identifies the deterministic return fibre with the lift interpretation. The witnessed paths and temporal tower retain the internal execution, including the prefixes of a computation that has not returned.


\subsection{Calculus: effective computations and sequential elimination}\label{subsec:calculus-effective}
The effective maps and lift operation now interpret computational process terms. Sequential elimination uses a returned value only on its classified domain, preserving the distinction between a returned failure tag and a computation that has not returned.

\begin{definition}[Computational formation and rules]\label{def:computational-formation-rules}
A \textbf{computational type} $X:\mathsf{Ty}_{\rm eff}$ is an object of $\mathsf{EDom}$. A \textbf{computational context} is a finite product of these objects; $\Gamma\vdash v:X$ denotes an effective map $\Gamma\to X$, while $\Gamma\vdash c\mathbin{\triangleright}X$ denotes an effective map $\Gamma\to\mathsf{Lift}X$. Thus the atomic formation and judgment interpretations are
\[
 \frac{X\in\operatorname{Ob}(\mathsf{EDom})}{X:\mathsf{Ty}_{\rm eff}},\qquad
 \llbracket\Gamma\vdash c\mathbin{\triangleright}X\rrbracket:
 \llbracket\Gamma\rrbracket\to\mathsf{Lift}X.
\]

For an effective domain $X$, $\Lift X$ is the classifier constructed above. At the computational judgment sort use
\[
 \frac{X:\Ty_{\rm eff}}{\Lift X:\Ty_{\rm eff}},\qquad
 \frac{\Gamma\vdash v:X}{\Gamma\vdash\ret(v)\mathbin{\triangleright}X},\qquad
 \frac{\Gamma\vdash c\mathbin{\triangleright}X\quad\Gamma,x:X\vdash d\mathbin{\triangleright}Y}
 {\Gamma\vdash\mathsf{let}\ x\leftarrow c\ \mathsf{in}\ d\mathbin{\triangleright}Y}.
\]
This is a cartesian computational context, not the owned interface context $\Delta$. The interpretations are $\eta_Xv$ and
\[
 \llbracket\mathsf{let}\rrbracket
   =\mu_Y\Lift(\llbracket d\rrbracket)
      \mathsf{st}_{\Gamma,X}\langle\id,\llbracket c\rrbracket\rangle.
\]
The computation equations are left and right unit for bind and associativity of sequential bind. These are the equations of the constructed monad and its strength.
\end{definition}

\begin{proposition}[Computational interpretation and equations]\label{prop:calculus-effective}
The return and sequential-elimination rules are well typed and satisfy
\[
 \mathsf{let}\ x\leftarrow\mathsf{return}(v)\ \mathsf{in}\ d\equiv d[v/x],\qquad
 \mathsf{let}\ x\leftarrow c\ \mathsf{in}\mathsf{return}(x)\equiv c,
\]
and the associative equation for two successive binds, with capture-avoiding substitution.
\end{proposition}
\begin{proof}\vmode
\leavevmode
\begin{enumerate}[label=(\roman*)]
\item \textbf{Sequential interpretation.}
The strength in Definition~\ref{def:computational-formation-rules} has target $\mathsf{Lift}(\Gamma\times X)$.
Applying $\mathsf{Lift}(\llbracket d\rrbracket)$ and then $\mu_Y$ gives $\mathsf{Lift}Y$, as required.

\item \textbf{Unit equations.}
The left and right monad units prove the two displayed equations.

\item \textbf{Associativity.}
The associativity and strength equations of Proposition~\ref{prop:effective-lift-monad} prove associativity with the same cartesian parameter in both continuations.

\item \textbf{Parameter substitution.}
Naturality of strength proves parameter substitution.
These are equations of the explicitly constructed lift.
\end{enumerate}
\end{proof}
\subsection{Finite display normalization}
With a complete finite display, all occurrence data can be inspected. The normalization procedure compares that data up to display isomorphism; the following subsection uses traces to decide the separate observational comparison.

\begin{proposition}[Finite process equality]\label{prop:pc-normalization}
Finite core declarations and their finite typed cuts have a terminating construction of their displayed event structures. Equality up to display isomorphism is decidable. At a fixed finite event set, a live declaration has the canonical presentation given by its causal covers and its inclusion-minimal inconsistent downsets, relative to the declared interface.
\end{proposition}
\begin{proof}\vmode
\leavevmode
\begin{enumerate}[label=(\roman*)]
\item \textbf{Causal and exclusion normalization.}
Compute the generating causal closure, remove the dead principal histories described in Proposition~\ref{prop:pc-core-sound}, and retain the transitive reduction of the live order.
Saturate each written exclusion by its causal past and the inherited arena consistency test.

\item \textbf{Minimal inconsistent ideals.}
Enumerate the inconsistent finite downsets and retain their inclusion-minimal members; each is represented by its maximal cohort.
Their antichain determines all inconsistent downsets.

\item \textbf{Canonical finite data.}
Configuration families determine the finite event structure: order is forced precedence in all configurations, and a finite set is consistent exactly when its causal closure is a configuration.
Thus the normalized data are canonical at fixed occurrence names.

\item \textbf{Finite isomorphism search.}
Permitting display-preserving renaming reduces equality to a finite bijection search.

\item \textbf{Forwarding and bending.}
A finite tagged forwarding adapter is computed from finite copycat by relabelling its external boundary.
Port bending and unbending are finite connections with the corresponding copycat cup and cap.

\item \textbf{Secured-cut normalization.}
At a cut, enumerate the finite secured matchings and retain the visible principal histories as in Section~\ref{subsec:pc-secured-proof}.
This gives another finite display, to which the same procedure applies.

\item \textbf{Term induction and scope.}
Structural induction on the finite term gives termination.
This algorithm supplies equality of constructed displays; structural $\beta\eta$ conversions alone are not asserted to be a complete presentation of all leaf equations.
\end{enumerate}
\end{proof}

\begin{construction}[Finite display normalization algorithm]\vmode\label{con:finite-display-algorithm}
\leavevmode
\begin{enumerate}[label=(\roman*)]
\item \textbf{Finite input tables.}
Represent a finite arena or source by a numbered finite carrier,
its order and consistency tables, polarity labels, and display table.
A finite core term additionally gives its finite waiting/exclusion
lists and finite typed syntax tree.

\item \textbf{Core normalization.}
At a core leaf compute causal transitive closure, test each
principal past, and remove precisely the dead events.
Enumerate
subsets of the remaining finite carrier to obtain its configuration
table.

\item \textbf{Tensor construction.}
At a tensor take tagged disjoint union.

\item \textbf{Cut construction.}
At a cut enumerate pairs
of configurations, apply $\mathsf{Match}$, and enumerate their
visible rooted principal histories, deduplicating by their retained
labelled graph codes.
Compute the finite source tables from these
histories.

\item \textbf{Copycat and bending.}
Copycat and port bending use the same finite operations.

\item \textbf{Canonical serialization.}
For a completed display, enumerate display-preserving bijections
of its finite carrier and serialize the transported order,
consistency, polarity and display tables.
The least serialization
in the fixed code order is its isomorphism normal form.

\item \textbf{Termination.}
Every loop is over a finite explicitly bounded carrier obtained at a
strict subtree, so structural induction proves termination.

\item \textbf{Correctness and scope.}
The game-constructor lemmas prove the denotation of each computed
table.
Equality of the final serializations is equivalent to existence
of a display-preserving bijection preserving and reflecting those
tables, which is exactly display isomorphism.
The procedure concerns
these finite presentations, not equality of arbitrary recursively
executing processes.
\end{enumerate}
\end{construction}


\subsection{Complete verification on supplied finite displays}\label{subsec:finite-observational-verification}
Finite traces give a terminating observational decision procedure when complete finite tables are supplied. A bound on equivalent strategy representatives then turns that procedure into a complete search for simulators over finite arenas.

\begin{theorem}[Finite observational decision and separating tests]\label{thm:finite-observational-decision}
For supplied finite arena and source tables, contextual may equivalence is decidable. The decision constructs either a complete trace certificate or an ordinary separating test.
\end{theorem}
\begin{proof}\vmode
\leavevmode
\begin{enumerate}[label=(\roman*)]
\item \textbf{Enumeration of trace witnesses.}
Enumerate all source subsets, retain the down-closed consistent ones, enumerate their finite linear extensions, and apply the display to each word.
This calculates both trace sets and the finite serialization carriers.

\item \textbf{Certificate or separating test.}
Equal sets give the equality pullback of Definition~\ref{def:trace-certificate}; otherwise choose a word in the symmetric difference and construct $\tau_w$ by Lemma~\ref{lem:characteristic-test}.
Its two success outcomes differ.

\item \textbf{Completeness and termination.}
Theorem~\ref{thm:trace-full-abstraction} proves completeness, while all loops are finite.
This decides observation, separately from display isomorphism in Proposition~\ref{prop:pc-normalization}.
\end{enumerate}
\end{proof}

\begin{theorem}[Bounded finite representatives and simulator synthesis]\label{thm:bounded-simulator-synthesis}
For an arena with $N$ events, put
\[
 L_N=\sum_{j=0}^N\frac{N!}{(N-j)!},\qquad K_N=NL_N,\qquad
 B_N=K_N+N(K_N+1)^N.
\]
Every ordinary strategy on this arena has a contextually equivalent ordinary strategy with at most $B_N$ source events. Consequently, for finite resources with supplied complete finite tables and finite honest and adversarial arenas, existence of a simulator is decidable, and a successful decision returns a finite simulator and a complete trace certificate.
\end{theorem}
\begin{proof}\vmode
\leavevmode
\begin{enumerate}[label=(\roman*)]
\item \textbf{Choice of trace representatives.}
There are at most $L_N$ distinct-event boundary words.
For each realized word select one finite realizing configuration; finite induction suffices for these finitely many selections.

\item \textbf{Retained positive occurrences.}
Let $P$ be all positive occurrences in the selected configurations, so $|P|\leq K_N$ and $P$ contains their positive ancestors.

\item \textbf{Restricted strategy.}
Restrict the source to $P$ and every negative occurrence all of whose positive ancestors lie in $P$.
This is down-closed and inherits consistency and courtesy.

\item \textbf{Receptivity.}
It is receptive: an enabled negative has its unique lift in the original source; by courtesy every immediate predecessor has an arena-predecessor label already present in the given configuration, so local injectivity puts those source predecessors there.
Its positive ancestors therefore lie in $P$, and the lift remains in the restriction.

\item \textbf{Trace preservation.}
Every selected configuration remains, and no new trace is introduced.

\item \textbf{Bound on negative occurrences.}
Lemma~\ref{lem:pc-injective} identifies a negative occurrence by its arena label and tuple of positive arena-ancestor occurrences.
With at most $K_N$ selected positives, the number of such tuples is bounded by $(K_N+1)^N$, including absent coordinates.
Thus there are at most $N(K_N+1)^N$ negative occurrences.

\item \textbf{Finite representative.}
The restriction is finite with the stated bound and the same trace set.

\item \textbf{Finite simulator search.}
For a simulator $\sigma:B\to A$ use $N=|B|+|A|$.
Enumerate all numbered source tables of size at most $B_N$, all display maps to $B^\ast\parallel A$, and test the finite event-structure, display, receptivity and courtesy conditions.

\item \textbf{Observational decision.}
Construct each finite simulated composite and apply Theorem~\ref{thm:finite-observational-decision}.

\item \textbf{Completeness and returned witness.}
If any simulator exists, its bounded representative occurs in this search and contextual congruence preserves its resource equation.
Otherwise the finite search rejects.
A successful candidate supplies both the strategy and the canonical comparison certificate.
This is a completeness bound, not a practical complexity estimate.
\end{enumerate}
\end{proof}

For an effectively presented algebraic payload, represent an effective open by an enumeration of basic opens and a supplied value by its compact evidence. Dovetail the two enumerations to witness membership: a common compact generator suffices when the lower evidence of the value is enumerated. For an effectively presented continuous processor the corresponding pairs of input and output basis evidence can likewise be enumerated. Two enumerations of opens need not admit a total extensional equality test. On $\Delta\mathbb N$, compare the empty enumeration with one that lists $0$ just when a specified machine halts. Deciding equality would decide whether that machine halts. The finite decision and synthesis algorithms above continue to use their supplied complete interface and display tables; singleton relabelling preserves every entry of such a table.

\subsection{The exact limit for arbitrary effective presentations}\label{subsec:effective-incompleteness}
An enumerated presentation need not reveal whether further events will appear. The following reduction shows why the finite procedures cannot extend to a sound and complete enumerable certificate system for all such presentations, even when the underlying resources are small.

\begin{theorem}[No complete enumerable finite certificate system]\label{thm:no-effective-complete-certificates}
For the effective event presentations of Definition~\ref{def:effective-event-presentation}, there is no sound and complete system of finite security certificates with an effective checker. This remains true for correct conversions whose underlying resources have at most two events, when their complete finite tables are not supplied.
\end{theorem}
\begin{proof}\vmode
\leavevmode
\begin{enumerate}[label=(\roman*)]
\item \textbf{Machine-indexed resource.}
Given a machine index $e$, take $H$ with one positive event $p$ and $A$ with one negative event $n$.
Enumerate a resource source containing $n$ immediately and containing $p$, with $n<p$, exactly when machine $e$ halts.

\item \textbf{Effective strategy presentation.}
This is an effective presentation: use fixed codes $0,1$, enumerate $1$ upon halting, and give total computable equality, polarity, order, consistency and finite-past operations on valid codes by their fixed two-event tables.
The display is receptive and courteous in either case.

\item \textbf{Correct conversion.}
Let the ideal $s:I\to H$ be silent.
The converter $c:A\to I$ with empty source supplies no $n$, so $(1_H\otimes c)r_e=s$ irrespective of halting.

\item \textbf{Security and nonhalting.}
A simulator $I\to A$ is necessarily the unique receptive negative-event strategy.
The real and simulated trace sets agree exactly when $e$ does not halt: both then have only $\eps,n$; if it halts the real resource additionally has $np$.
Thus the correct conversion is secure exactly on the nonhalting indices.

\item \textbf{Certificate-enumeration contradiction.}
Enumerating finite certificates and running an effective checker would enumerate these indices if the system were sound and complete, contradicting the non-enumerability of the complement of the halting set.
\end{enumerate}
\end{proof}

The distinction is the input presentation. A supplied finite table certifies that its listed carrier is complete. An enumerator which has so far produced one event does not certify that it will never produce another. Likewise a finite program can generate unboundedly many occurrences or payloads. The finite synthesis theorem, the general semantic comparison calculus, and this effective obstruction therefore have different, precisely stated domains.

\subsection*{What the constructions make available}
The manuscript connects occurrence implementations, their exposed behaviour and simulator-witnessed realization. The process representation retains the causal and consistency data of ordinary strategies. Complete finite-trace testing then supplies an observational criterion, with certificates that preserve witnesses and compose through contexts. An actual simulator makes that criterion a resource realization statement, and the correctness and security witnesses compose with the protocols they compare.

Retained histories supply the views used by distributed controllers, the transition witnesses for finite return, and the indices for evidence and policy. The public observation map sends each history to the configuration visible at a coalition's ports. The appendices use these maps to prove the stated privacy, disclosure, and origin properties at each protocol interface.

The effective results delimit what can be checked uniformly. Complete finite tables permit observational decision and bounded simulator synthesis; an unrestricted enumerator does not certify that no further occurrence will appear. This difference concerns the input presentation, even when the underlying resource is finite. Likewise the cryptographic comparisons concern the exact symbolic operations and qualitative observations specified here. Their relation to quantitative interpretations is set out in Appendix~\ref{app:quantitative-comparison}.

\appendix

\section{Communication protocol examples}\label{app:communication-examples}
The following media instantiate the communication and control constructions of Section~\ref{sec:protocol-games}. Each example specifies the actual interface, occurrence realizer and causal execution witnesses. Its guards, exclusions and positive tags use the process representation of Theorem~\ref{thm:pc-representation}. The examples identify what their comparisons preserve: physical history maps transport controller and concurrency claims, while trace or causal certificates transport observational claims. Equality in a testing quotient alone does not establish equality of concurrent cells.

\subsection{Relay, join, and finite exclusion}
\begin{example}[Relay, join, and finite exclusion]\label{ex:pc-generators}
On a one-request/one-response interface with independent $r^-$ and $s^+$, the relay is $\pcCore[A;\{(r,s)\};\varnothing]$. Its histories are $\varnothing$, $\{r\}$, and $\{r,s\}$. With message labels $a?v$ and $b!v$ and the declared matching value branches, its named spelling is $a?(x).b!x.0$.

For independent requests $r^-,s^-$ and a response $p^+$, the declaration with $J=\{(r,p),(s,p)\}$ derives a join: both receipts are required, and neither receipt is consumed by the test. On three independent positives $p,q,r$, the clause $K=\{\{p,q,r\}\}$ permits every pair but excludes the triple. It cannot be replaced by pairwise conflicts on these occurrences, since every pair is consistent. The same constructors will describe supplying packets, joint readiness, and competing resource uses.
\end{example}
The named spelling of the relay in Example~\ref{ex:pc-generators} therefore has one acceptance for each input-value branch and a matching output waiting on that acceptance. Its positive action is enabled after receipt, and its input branches remain those of the declared arena. This supplies the connection between named syntax and the exact occurrence semantics before choosing any communication discipline.

\subsection{Concrete medium updates and interchange}
\begin{lemma}[Packet stores over a payload domain]\label{lem:open-queue-store}
The finite FIFO store is $\coprod_{n\ge0}D^n$, ordered componentwise within each length. Appending a packet is continuous, and on the open of nonempty stores taking the head and tail is continuous. For a receive guard $U\in\Op(D)$, the stores permitting receipt are the open $\bigcup_{n>0}\pi_1^{-1}U$ over those summands. Removing an already present head commutes with appending a packet, with the same head occurrence in both orders.
\end{lemma}
\begin{proof}
A directed family stays in one summand, where the operations are tupling and projections. Writing $xw$ for head and tail, both orders of operations return $x$ and leave $\mathsf{snoc}(w,y)$ after appending $y$. This calculation keeps the two designated occurrences even when $x=y$. For $D=\Delta M$ and $U=\{m\}$ it is the restricted dequeue on $mM^*$ below.
\end{proof}
The order on stores compares payload evidence at fixed length. The stream order of Section~\ref{subsec:open-streams} also compares lengths. In $0<1$, the observation $(0)\preceq(1)$ of one packet admits no suffix $w$ with $(1)=(0)w$. The buffer contents are recovered from the actual input and delivery occurrences; their observations alone can be strictly coarser.

\begin{definition}[Queue update maps]\label{def:queue-update-maps}
On the word monoid $M^*$, put $\mathsf{enq}_n(v)=vn$ and $\mathsf{deq}_m(mv)=v$, the latter with domain $mM^*$. These are the store projections of the FIFO commitments of Definition~\ref{def:rendezvous-fifo-computation}; for occurrence-valued $M$ they retain the selected packets.
\end{definition}

\begin{lemma}[FIFO interchange]\label{lem:fifo-square}
Two initially enabled FIFO operations with disjoint participating owners commute exactly when their retained resource guards permit both orders. For an enqueue and an earlier head removal the square is
\[
 \mathsf{deq}_m(\mathsf{enq}_n(mv))=vn
 =\mathsf{enq}_n(\mathsf{deq}_m(mv)).
\]
The residuals retain the same local transitions and the same packet occurrences.
\end{lemma}
\begin{proof}\vmode
\leavevmode
\begin{enumerate}[label=(\roman*)]
\item \textbf{Commuting queue updates.}
The enqueue changes neither the selected head nor the receiver's control; removal changes neither the sender's control nor its produced name.
Both orders therefore have the displayed word and the same controls.

\item \textbf{Disjoint owners and channels.}
On a point-to-point channel participant disjointness allows at most one selected send and one selected receive.
Different channels update distinct coordinates.

\item \textbf{Higher interchange and exclusion.}
This also proves the corresponding finite-family statement.
A receive of the new packet is not initially enabled and supplies a causal edge instead of a square.
\end{enumerate}
\end{proof}
For a bag, distinct selected removals and fresh insertions commute by finite-set deletion and union. Loss, duplication, exchange, and replay use respectively deletion, replacement by named children, guarded adjacent permutation, and fresh append with a retained source. Bounded buffers retain their actual branch: an initially full drop-new send and a receive do not form a square because the branch changes after the receive. Credit operations preserve each slot's selected location and the two word guards. Atomic joint operations are single directions; cancellation and delivery cannot consume one packet twice. Allocation keeps its selected name and permission and rechecks unusedness. These computations define every coefficient for these media, including mixed networks, rather than adding concurrency to their transition graph by a filling convention.

\subsection{Participant and resource histories}\label{subsec:concrete-participant-histories}
\begin{definition}[Participant diagrams and local-cut evaluation]\label{def:formal-participant-diagrams-and-local-cut-evaluation}
The protocol examples use finite causal diagrams with labelled input and output occurrences. A participant contributes its own chain of local actions; a resource use retains the supplying occurrence, and a communication retains its supplying send. Initial occurrences have no producing event. Each diagram has its finite causal order, and its executable prefixes are the order ideals. The constructions in Appendix~\ref{app:communication-examples} determine the order and verify that every ideal has the prescribed resource interpretation. Two schedules of one diagram differ by adjacent exchanges of incomparable events, so its simultaneous moves are derived from that order.

Lemma~\ref{lem:admin-family} realizes these locally ordered configuration families by their principal histories. Its ideal description is verified in Appendix~\ref{app:communication-examples} for each protocol; Section~\ref{subsec:disjunctive-causal-interfaces} retains compatible alternative causes when a single prime-history presentation does not apply.

For participant diagrams write $\dom(e)$ for the participants of an event. A participant's events in a configuration form a chain, and two incomparable events have disjoint participant supports. The current local cut, pending messages, and live resources are evaluated from any linearization. The examples prove that adjacent independent exchanges preserve this evaluation. We denote the resulting graph of evaluated cuts by $N$ when discussing distributed control; $q(h)$ is this evaluation, and $\En_N(q(h))$ means the witnessed extensions available at that cut. Section~\ref{subsec:histories} proves that the witnessed construction yields the same diagrams for rendezvous and FIFO networks.
Explicitly the \textbf{cut graph} has vertex set $N_0$ of reachable evaluated cuts and arrow set $N_1$ of witnessed operations between them, with endpoint maps $s_N,t_N:N_1\to N_0$. The \textbf{enabled-operation set} is $\En_N(q)=s_N^{-1}(q)$. Writing $\mathcal F_N$ for the family of finite participant diagrams just specified, evaluation has type $q:\mathcal F_N\to N_0$; a witnessed extension determines its corresponding arrow in $N_1$. The distinction between a history and its evaluated cut is retained by this map.
\end{definition}

\subsection{Rendezvous and FIFO occurrence histories}\label{subsec:histories}
\begin{definition}[Operational executions and finite histories]\label{def:formal-operational-executions-and-finite-histories}
An execution is a finite path in the plant graph. For rendezvous, associate to it the partial order generated by the chronological order of events involving each participant. For FIFO add an edge from each enqueue occurrence to its matching dequeue occurrence. Fix the initial plant state, including its queue words. If a channel initially contains $r$ packet occurrences, its $k$th dequeue consumes the $k$th initial packet when $k\leq r$, and matches its $(k-r)$th new enqueue when $k>r$, with the same payload. Initial packets are distinguished by their boundary positions; they are not additional internal send events. Add a causal edge from an enqueue only in the second case. Retain the \textbf{initial state}, local transition witnesses, participant names, and these matchings. A \textbf{finite history} is such a labelled partial order, up to label-preserving isomorphism. Empty initial queues give the usual $k$th-send/$k$th-receive matching.
\end{definition}

\begin{lemma}[History adequacy]\label{lem:protocol-history-adequacy}
Every down-closed subset of a finite history has a legal linearization. Any two linearizations of a history differ by swaps of adjacent independent witnesses from Proposition~\ref{prop:resource-charts}. Consequently a history determines a unique final plant state.
\end{lemma}
\begin{proof}\vmode
\leavevmode
\begin{enumerate}[label=(\roman*)]
\item \textbf{Execution of participant chains.}
Take a linear extension of the finite partial order.
A local transition occurs after all previous transitions of its participant, so its local source state is correct.

\item \textbf{FIFO suppliers.}
Receives are ordered by their consumer.
The first $r$ receives on a channel therefore consume the initial packets in boundary order.
Every subsequent receive follows its matched new send and all earlier receives; new sends are ordered by their producer.
Its matched packet is consequently at the head when it executes.

\item \textbf{Downward-closed prefixes.}
The same argument applies to a downset, since it retains every required local predecessor and every required new send.

\item \textbf{Independent exchanges.}
Finite linear extensions are connected by adjacent swaps of incomparable elements: move the first element of one extension leftwards in the other, then induct.
In rendezvous such elements have disjoint supports.
In FIFO they have different owners; if they share a channel, they are an enqueue and a dequeue of an earlier message, so the dequeue is initially enabled at the swapped position.
Lemma~\ref{lem:fifo-square} applies.

\item \textbf{Converse and state preservation.}
The converse swaps are precisely the generated squares.
Each swap preserves the final state and residual witnesses.
\end{enumerate}
\end{proof}

\begin{definition}[Occurrence unfolding]\label{def:unfolding}
For a \textbf{finite history} $h$ and $e\in h$, put $\iota_h(e)=[(\downarrow_h e,e)]$, the label-preserving isomorphism class of its rooted principal history. Occurrence events are these classes. Define their order by
\[
 [(p,e)]\leq[(q,f)]
 \quad\Longleftrightarrow\quad
 \exists d\in q:\ [(p,e)]=[(\downarrow_qd,d)].
\]
Here $p=\downarrow_pe$ and $q=\downarrow_qf$. A finite set $X$ of occurrences is consistent when $X\subseteq\iota_h[h]$ for a finite legal history $h$. Denote the resulting ordered consistency structure by $U_N$.
\end{definition}

Event structures and configurations have the meanings fixed in Definition~\ref{def:event-structure}; see also~\cite[Section~2]{profunctors}. In Definition~\ref{def:unfolding}, a principal subhistory retains all witnesses and matchings, so an event's identity includes its causal history, not merely a global counter.

\begin{lemma}[Occurrence rigidity]\label{lem:occurrence-rigidity}
The order in Definition~\ref{def:unfolding} is well-defined on classes and is a partial order. For every finite history $h$, the map $\iota_h$ is injective and preserves and reflects order. A down-closed restriction of $h$ preserves the principal-history class of each retained event.
\end{lemma}
\begin{proof}\vmode
\leavevmode
\begin{enumerate}[label=(\roman*)]
\item \textbf{Principal-history invariance.}
A labelled isomorphism carries rooted principal downsets to rooted principal downsets, proving independence of representatives.

\item \textbf{Order on occurrence classes.}
Reflexivity and transitivity follow by principal restriction.
Mutual comparison gives embeddings between two finite principal histories in both directions.
Their cardinalities agree, so each embedding is surjective, proving antisymmetry.

\item \textbf{Participant ranks.}
Every event in these rendezvous or point-to-point FIFO histories has a participant.
For $p\in\dom(e)$ put
\[
 r_p(e)=|\{d\leq_h e:p\in\dom(d)\}|.
\]
The $p$-events of $h$ form a chain.

\item \textbf{Injectivity.}
Isomorphic rooted principal histories have the same participant labels and the same ranks.
If $\iota_h(e)=\iota_h(f)$, choose $p\in\dom(e)=\dom(f)$.
Then $r_p(e)=r_p(f)$, so $e=f$ in that chain.
This proves injectivity.

\item \textbf{Order reflection.}
If $e\leq_h f$, principal restriction gives $\iota_h(e)\leq\iota_h(f)$.
Conversely this comparison supplies $d\leq_h f$ with $\iota_h(d)=\iota_h(e)$; injectivity gives $d=e$.

\item \textbf{Restriction and administrative histories.}
Down-closed restriction leaves every retained principal history unchanged.
The same proof applies to administrative histories, whose enriched labels retain their nonempty participant sets.
\end{enumerate}
\end{proof}

\begin{proposition}[Unfolding configurations]\label{prop:unfolding}
The occurrence construction is an event structure. Its finite configurations are exactly the finite histories, and its cubical configuration carrier maps to $X_N$ by taking final plant states and residual witnesses.
\end{proposition}
\begin{proof}\vmode
\leavevmode
\begin{enumerate}[label=(\roman*)]
\item \textbf{Order and finite pasts.}
Lemma~\ref{lem:occurrence-rigidity} gives the partial order.
The predecessors of $[(p,e)]$ are exactly $\iota_p[p]$, so every principal downset is finite.

\item \textbf{Consistency axioms.}
The empty history witnesses empty consistency, each singleton is represented by its principal history, and consistency is hereditary.
If a finite set is represented in $h$, every predecessor of one of its occurrences is also represented in $h$; this proves closure under adding predecessors.

\item \textbf{Recovery of a finite history.}
A finite configuration $X$ is represented in some $h$.
Order reflection makes $J=\iota_h^{-1}(X)$ a downset of $h$.
Lemma~\ref{lem:protocol-history-adequacy} makes $h|_J$ a legal history, and restriction preserves its principal classes, so $\iota_{h|_J}[J]=X$.

\item \textbf{Inverse configuration map.}
Conversely $\iota_h[h]$ is a configuration.
The map $e\mapsto\iota_h(e)$ identifies its order and labels with those of $h$.
These assignments are inverse up to labelled history isomorphism; in particular the reconstruction from $X$ is independent of its chosen representing history.

\item \textbf{Concurrent cubes.}
A configuration cube consists of a history and a finite consistent set of pairwise concurrent extensions.
Such extensions have no causal predecessors among one another.
Lemma~\ref{lem:protocol-history-adequacy} identifies their successive updates with a residual family, giving the map.

\item \textbf{Face compatibility.}
Executing or omitting an extension commutes with taking its residual, so it preserves faces.
\end{enumerate}
\end{proof}

\begin{definition}[Residual event structure and history graph]\label{def:residual-event-structure}
Use the residual $U/h$ of Definition~\ref{def:event-residual-carrier} and its configuration description in Lemma~\ref{lem:event-residual-configurations}. The \textbf{history-extension graph} $\widetilde C_U$ has vertices $\Conf_f(U)$ and edges $(h,e)$ with $e$ enabled at $h$, with source $h$ and target $h\cup\{e\}$. Its extension order is $\mathcal H_U=(\Conf_f(U),\subseteq)$. The \textbf{prefix insertion} is $j_h(x)=h\cup x$. The \textbf{rooted upper cube carrier} $K(U)_{\geq h}$ consists of cells whose initial configuration contains $h$, rooted at $h$.
\end{definition}

\begin{proposition}[Residual histories and execution lifting]\label{prop:history-residual}
For an occurrence unfolding $U$ of the initialized rendezvous or point-to-point FIFO plant and $h\in\Conf_f(U)$, define
\begin{align*}
 |U/h|&=\{e\in U\setminus h:h\cup\downarrow_Ue\in\Conf_f(U)\},\\
 X\in\mathsf{Con}_{U/h}
 &\Longleftrightarrow h\cup\downarrow_UX\in\Conf_f(U)
 \qquad(X\subseteq |U/h|\text{ finite}).
\end{align*}
With inherited order and polarity this is an event structure, and $j_h(x)=h\cup x$ gives
\begin{equation}\label{eq:history-residual}
 \Conf_f(U/h)\cong\{k\in\Conf_f(U):h\subseteq k\},\qquad
 K(U/h)\cong K(U)_{\geq h}.
\end{equation}
The right-hand cubical arena is rooted at $h$. For a residual configuration $x$, there is a canonical identity-on-events isomorphism
\[
 (U/h)/x\cong U/(h\cup x).
\]
Write $\widetilde C$ for the graph of one-event extensions of $H=\Conf_f(U)$. Terminal-state evaluation $q:\widetilde C\to C$ uniquely lifts every enabled rule witness from each $h\in H$ and uniquely lifts the plant's independent residual families.
\end{proposition}
\begin{proof}\vmode
\leavevmode
\begin{enumerate}[label=(\roman*)]
\item \textbf{Residual event structure.}
Residual principal downsets are finite.
Singletons are consistent by the definition of $|U/h|$; taking subsets or adjoining a causal predecessor preserves the defining consistency condition.

\item \textbf{Configuration bijection.}
A residual configuration $x$ has $h\cup x$ downward closed and consistent.
Conversely, removing $h$ from a configuration extending it gives exactly a residual configuration.
This proves the first bijection in \eqref{eq:history-residual}.

\item \textbf{Cube correspondence.}
The same bijection sends a cube $(x,V)$ to $(h\cup x,V)$ and commutes with all faces, proving the second.

\item \textbf{Iteration.}
Iterated residuals have the same remaining events and the same test of consistency with $h\cup x$, so their comparison and its iteration coherences are identities on events.
Testing consistency with the fixed past is essential; restricting the original consistency relation alone would not perform that test.

\item \textbf{Lifting an enabled operation.}
Given $h$ and an enabled plant witness at $q(h)$, append it to a linearization of $h$.
Its past is the union of its participants' preceding views, together with the supplying send's past for a receive of a newly sent packet.
The fixed initial-packet designation handles a boundary packet.
This determines its principal occurrence history independently of the linearization, by Lemma~\ref{lem:protocol-history-adequacy}, and gives the unique one-event lift.

\item \textbf{Lifting independent families.}
For an independent family, disjoint participant updates preserve these prerequisites; the possible shared FIFO update is exactly Lemma~\ref{lem:fifo-square}.
All residual permutations therefore lift to the same occurrence family and the same cube.
This constructs the graph and cubical comparison without a scheduler.
\end{enumerate}
\end{proof}

\subsection{Participant-free joint records and their occurrence chains}\label{subsec:nullary-translation}
\begin{definition}[Participant-free record carriers]\label{def:nullary-joint-carriers}
Let $R_0$ be the set of joint operation witnesses with empty participant support and identity state update. Write $\mathsf{Inj}(F,R_0)=\{u:F\to R_0:u\text{ injective}\}$. Define $T_{R_0}(F)=\mathsf{Inj}(F,R_0)$ with restriction faces, and $M_0=\mathbb N^{(R_0)}$, the finite-support functions under addition. The \textbf{record-occurrence chain structure} $E_0^{\mathrm{joint}}$ has events $(r,n)\in R_0\times\mathbb N_{>0}$, order $(r,n)\leq(s,m)$ iff $r=s$ and $n\leq m$, and all finite sets consistent. The suffix $\mathrm{ne}$ denotes restriction to nonempty-participant operations.
\end{definition}

\begin{proposition}[Occurrence chains for participant-free joint records]\label{prop:nullary-joint}
Let $R_0$ be the state-identity joint records with empty participant block $S\cup R=\varnothing$. Their tags remain distinct. Their cubical factor is
\[
 T_{R_0}(F)=\mathsf{Inj}(F,R_0),
\]
with faces given by restriction, and its execution category is the one-object category of the free commutative monoid $M_0=\mathbb N^{(R_0)}$. Its root-representable history display is the cubical configuration carrier of the event structure
\[
 |E_0^{\mathrm{joint}}|=R_0\times\mathbb N_{>0},\qquad
 (r,k)\le(s,l)\ \Longleftrightarrow\ r=s\text{ and }k\le l,
\]
with every finite subset consistent. For the pure joint plant there are canonical factorizations
\[
 X_{\mathrm{joint}}\cong X_{\mathrm{joint},\mathrm{ne}}\boxtimes T_{R_0},
 \qquad
 U_{\mathrm{joint}}\cong U_{\mathrm{joint},\mathrm{ne}}\parallel E_0^{\mathrm{joint}},
\]
where the subscript $\mathrm{ne}$ denotes the nonempty-participant part.
\end{proposition}
\begin{proof}\vmode
\leavevmode
\begin{enumerate}[label=(\roman*)]
\item \textbf{Participant-free record transitions.}
A record with empty participant set has no premises, leaves every local coordinate fixed, and leaves the medium index $b$ fixed.
It is therefore an everywhere-enabled loop independent of every other instruction.
Definition~\ref{def:resource-coefficient} requires distinct directions, giving exactly the injections in $T_{R_0}$.

\item \textbf{Commutative execution monoid.}
Paths in this factor are words in its records and square relations exchange adjacent distinct records.
Such exchanges preserve multiplicities.
Conversely, words with equal multiplicities are related by these exchanges: move the first letter of one word leftwards from its first occurrence in the other word and induct on length.
Thus the execution monoid is $M_0$, without choosing an ordering of $R_0$.

\item \textbf{Rooted history cover.}
By Theorem~\ref{thm:protocol-audit}, its history display has cells $(n,j)$ with $n\in M_0$ and $j:F\hookrightarrow R_0$.
Its lower and upper faces are respectively
\[
 d_f^0(n,j)=(n,j|_{F\setminus\{f\}}),\qquad
 d_f^1(n,j)=(n+\delta_{j(f)},j|_{F\setminus\{f\}}).
\]
The principal past of $(r,k)$ in $E_0^{\mathrm{joint}}$ has $k$ elements.

\item \textbf{Configuration and cube classification.}
Its finite configurations are exactly
\[
 x_n=\{(r,k):1\le k\le n(r)\}.
\]
At $x_n$ the next event on the $r$-chain is $(r,n(r)+1)$.
Next events of distinct records are concurrent.

\item \textbf{Event-chain correspondence.}
Hence
\[
 (n,j)\longmapsto
 \bigl(x_n,\{(j(f),n(j(f))+1):f\in F\},
             f\mapsto(j(f),n(j(f))+1)\bigr)
\]
is a bijection to its configuration cubes and commutes with the displayed faces.
This identifies the history cover with $K(E_0^{\mathrm{joint}})$.

\item \textbf{Separation of joint charts.}
Every joint chart partitions uniquely into its empty-support and nonempty-support directions.
The former are state identities, so this partition and the Day tensor formula of Section~\ref{subsec:concurrent-cells-scott-categories} prove the cubical product.

\item \textbf{Product of history presentations.}
Execution categories and root representables split as products.
For the nonempty-support part, the participant-chain proof of Lemmas~\ref{lem:protocol-history-adequacy} and~\ref{lem:occurrence-rigidity} applies: a principal occurrence has a participant rank, and disjoint blocks give exactly the exchanges.
Tensoring that unfolding with the explicit chains gives the second factorization.

\item \textbf{Occurrence interpretation and scope.}
Equivalently, in a full joint history order repeated occurrences of each empty-support record chronologically and leave distinct such records independent of each other and of all nonempty blocks.
Its configurations give the same product.
No empty-support record is removed and no additional causal relation between distinct records is introduced.
This event-structure calculation is external and uses no axiom of choice.
\end{enumerate}
\end{proof}

\subsection{Communication kernels and their composition}\label{subsec:communication-kernel-composition}
The first kernel calculations use streams and synchronous communication to identify the witnesses carried through composition. They also establish the readiness comparisons used by the later implementation examples.

\subsubsection{Example: streams and FIFO copycat}
Take $D=\Delta M$ in Section~\ref{subsec:open-streams}. The following arena realizes its finite compact positions with one input or delivery event per letter. Proposition~\ref{prop:open-stream-processor} gives its identity support; the queue proof below calculates the occurrence display realizing this support.

\begin{definition}[Stream arena and prefix category]\label{def:formal-stream-arena-and-prefix-category}
For a message set $M$, let $A_M$ be the rooted all-positive stream arena. Its vertices are words $u\in M^*$, its edges are $u\to um$, and it has no cells of dimension greater than one. Its Scott category is the prefix category $C_M$: there is one arrow $u\to v$ exactly when $u$ is a prefix of $v$. Write $u\leq v$ for this relation.
\end{definition}

\begin{definition}[FIFO history kernel]
The FIFO kernel from $A_M$ to itself is
\begin{equation}\label{eq:fifokernel}
Y_M(u,v)=C_M(v,u)=
\begin{cases}*,&v\leq u,\\\varnothing,&\text{otherwise}.
\end{cases}
\end{equation}
Its root is the unique element at $(\eps,\eps)$.
\end{definition}
Its input occurrences are negative in $A_M^\ast\boxtimes A_M$, and its delivered outputs are positive.

\begin{theorem}[FIFO as asynchronous copycat]\label{thm:fifocopycat}
The cubical display of $Y_M$ has exactly the following data:
\leavevmode
\begin{enumerate}
\item vertices $(u,v)$ with $v\leq u$;
\item input edges $(u,v)\to(um,v)$ for every $m$;
\item output edges $(u,v)\to(u,vn)$ exactly when $vn\leq u$;
\item input/output squares at $(u,v)$ exactly when $vn\leq u$.
\end{enumerate}
It is reached. The map $\pi(u,v)=w$, where $u=vw$, is a rooted execution covering of the queue execution object. This kernel is the identity kernel of the stream arena and represents ordinary asynchronous copycat.
\end{theorem}
\begin{proof}\vmode
\leavevmode
\begin{enumerate}[label=(\roman*)]
\item \textbf{Presheaf support.}
The support of \eqref{eq:fifokernel} is covariant in $u$ and contravariant in $v$, so it is a kernel.

\item \textbf{Terminal-corner cells.}
On an input edge the Scott-terminal vertex is its execution source, already supported; hence the input always lifts.
On an output edge it is the execution target, supported precisely when $vn\leq u$.
On a mixed square it is $(u,vn)$, giving the same test by Lemma~\ref{lem:terminalcorner}.
There are no other possible cells in the tensor of two one-dimensional stream arenas.

\item \textbf{Reachability.}
To reach $(u,v)$, first accept all letters of $u$ and then deliver the prefix $v$.
Every intermediate vertex is supported, so the kernel is reached.

\item \textbf{Queue-state covering.}
At a supported vertex write $u=vw$ uniquely.
An input appends to $w$, an output removes its first letter, and the square condition says exactly that this removal was already enabled.
Thus each edge and cell at $w$ has a unique lift from the chosen history vertex.
This proves the covering assertion.

\item \textbf{Copycat kernel.}
Finally $Y_M$ is literally the hom kernel $Y_{A_M}$ defined alongside \eqref{eq:coend}.

\item \textbf{Ordinary event presentation.}
Its event presentation has input and output chains, with the $k$th output dependent on the $k$th matching input.
Its configurations are exactly prefix pairs $(u,v)$ with $v\leq u$.
Receptivity and courtesy follow either from this presentation or from the preceding lift calculation and Theorem~\ref{thm:scott-discrete-fibrations}.
Hence it is the ordinary asynchronous copycat strategy under that comparison.
\end{enumerate}
\end{proof}

The queue-state graph and the history display are different objects: many history pairs can have the same residual word. The covering remembers the transmitted and delivered prefixes. This is why a cyclic state implementation can yield an acyclic occurrence strategy.

\begin{example}[FIFO moves and their concurrency]\label{ex:native-fifo-rules}
At a supported FIFO position $(u,v)$, input of $m$ extends the first prefix and delivery of $n$ extends the second:
\[
 (u,v)\xrightarrow{m?}(um,v),\qquad
 \frac{vn\leq u}{(u,v)\xrightarrow{n!}(u,vn)}.
\]
The input and delivery form a square exactly when $vn\leq u$ already holds at the source; their common target is $(um,vn)$. If $u=v$, input of $m$ may enable its delivery, but the two moves do not form a square. These assertions follow by substituting the hom kernel of Theorem~\ref{thm:fifocopycat} into the two lifting rules. Removing the delivered prefix afterwards gives the residual queue, whose automata comparison is the covering in that theorem.
\end{example}

\subsubsection{Example: prefix processors and composition}\label{subsec:native-prefix}
\begin{definition}[Prefix-processor kernel]\label{def:formal-prefix-processor-kernel}
Let $f:C_M\to C_N$ preserve the root. Define $H_f(u,v)=C_N(v,f(u))$. This has an intrinsic game interpretation: at $(u,v)$, an input $m$ extends to $(um,v)$, and an output $n$ extends to $(u,vn)$ exactly when $vn\leq f(u)$. An input and an old output commute under this same source condition.
\end{definition}
\begin{proposition}[Prefix-processor strategies]\label{prop:native-prefix}
The kernel $H_f$ is reached and is represented by a receptive, courteous strategy. For rooted prefix functors $f:C_M\to C_N$ and $g:C_N\to C_L$, there is a coherent rooted isomorphism
\[
 H_g\star H_f\cong H_{gf}.
\]
\end{proposition}
\begin{proof}\vmode
\leavevmode
\begin{enumerate}[label=(\roman*)]
\item \textbf{Variance and reachability.}
Prefix preservation defines input covariance and output restriction.
Reach $(u,v)$ by accepting $u$ and delivering $v$.

\item \textbf{Causal suppliers.}
In a supported finite pair, each output position has a first input prefix at which it appears in $f(u)$: the eligible prefix lengths form a nonempty final segment of a finite ordinal.
Join that input to the output, and retain the two port chains.

\item \textbf{Principal-history strategy.}
Ideals of this finite order are exactly the supported prefixes.
Principal histories therefore construct the strategy, with unique negative extensions and only additional negative-to-positive dependencies.

\item \textbf{Kernel cells.}
The terminal-corner calculation gives the displayed move and square rules.

\item \textbf{Composition by co-Yoneda.}
Finally co-Yoneda gives
\[
 \int^v C_N(v,f(u))\times C_L(z,g(v))\cong C_L(z,g(f(u))).
\]
The map composes the two arrows; its inverse uses $v=f(u)$.
The coend equation proves both inverses.

\item \textbf{Coherence.}
Iterating either map evaluates the same string of arrows, giving coherence and root preservation.
\end{enumerate}
\end{proof}
The case $f=\id$ is FIFO copycat. Section~\ref{subsec:buffered-translation} identifies this construction with deterministic processors equipped with asynchronous output buffers, including their visible concurrent cells.

\subsubsection{Example: synchronous communication}

A rendezvous is a single joint event with participant support $\{p,r\}$. Independent rendezvous events have disjoint supports. Its open game interface exposes two readiness signals and one completed communication.

\begin{definition}[One-handshake arena]
Fix matching local witnesses for one message. Let $R$ be the product of three single-event arenas: left readiness $l$ and right readiness $r$ are negative, and commit $c$ is positive. Vertices are triples $(l,r,c)\in\{0,1\}^3$. Define a kernel from the tensor unit by
\[
\Phi_{\mathrm{rv}}(l,r,c)=
\begin{cases}*,&c\leq l\text{ and }c\leq r,\\\varnothing,&\text{otherwise}.
\end{cases}
\]
\end{definition}

\begin{proposition}[Readiness and commit]\label{prop:rendezvouskernel}
$\Phi_{\mathrm{rv}}$ is a reached kernel. Its display is the event structure with incomparable negative readiness events $l,r$ and a positive event $c$ causally following both. Erasing readiness from completed observations gives exactly one atomic rendezvous event. Independent handshakes tensor to the corresponding independent joint-action cells.
\end{proposition}
\begin{proof}\vmode
\leavevmode
\begin{enumerate}[label=(\roman*)]
\item \textbf{Presheaf restrictions.}
Restrictions add negative readiness or remove positive commit, so they preserve the support condition.
This defines the presheaf.

\item \textbf{Readiness and commit edges.}
The two readiness edges lift and form a square.
A commit edge lifts precisely at $(1,1,0)$.

\item \textbf{Exclusion of mixed squares.}
A readiness/commit square would have a Scott-terminal vertex with commit present and that readiness absent, which is unsupported.
Thus no such square lifts.

\item \textbf{Reached configurations.}
These are exactly the configurations and cells of the displayed event structure.
Every supported vertex has the execution that performs its present readiness events and then its commit if present, proving reachability.

\item \textbf{Erasure and tensor.}
Erasing $l,r$ from the causal histories leaves either the empty observation or the singleton event $c$, exactly the prefixes of the atomic rendezvous.
For independent copies the history and erasure maps split as products; Proposition~\ref{prop:scottmonoidal} identifies their tensor and its joint-action cubes.
\end{enumerate}
\end{proof}

Erasure in this proposition is an execution observation; it collapses edges and is not a dimension-preserving map of the precubical category. The fixed matching witnesses also matter: exposing premature local choices among several potential partners can add a reservation protocol with its own deadlocks. The atomic rendezvous plant and its control interpretation below do not make such choices prematurely.

\paragraph{Shared readiness for a whole network}
The local calculation extends to the entire synchronous plant by attaching readiness to local conditions, shared by every competing consumer of a condition. This avoids selecting a communication partner before the physical joint event. The construction below also applies to the participant histories of the point-to-point FIFO and administrative plants, although FIFO communication already has its copycat presentation.

\begin{definition}[Conditions and shared readiness]\label{def:shared-readiness}
Let $U$ be the principal-history event structure of participant diagrams from Section~\ref{subsec:native-resource-histories}. For the rendezvous and FIFO diagrams its equality with the operational unfolding $U_N$ is proved in Section~\ref{subsec:histories}. Construct its local condition set
\[
 B=P\amalg\coprod_{e\in|U|}\dom(e),
\]
writing $b_p^0$ for the initial condition of $p$ and $b_{e,p}$ for the condition produced by $e$ at $p$. For $p\in\dom(e)$, the $p$-events strictly preceding $e$ form a finite chain by their participant-chain construction. Define $\pre_p(e)=b_p^0$ when this chain is empty, and $\pre_p(e)=b_{d,p}$ when $d$ is its greatest element. Put
\[
 \pre_B(e)=\{\pre_p(e):p\in\dom(e)\},\qquad
 B_{\mathrm{use}}=\bigcup_{e\in|U|}\pre_B(e).
\]
The latter is the image of the event--precondition incidence relation; no consuming event is selected. For a \textbf{finite history} $h$, define
\begin{align*}
 B_h&=\{b_p^0:p\in P\}\cup
       \{b_{e,p}:e\in h,\ p\in\dom(e)\},\\
 C_h&=\bigcup_{e\in h}\pre_B(e),\qquad
 A(h)=(B_h\setminus C_h)\cap B_{\mathrm{use}}.
\end{align*}
Listing each participant's event chain shows that $B_h\setminus C_h$ contains exactly one current condition per participant.

Define $S^{\mathrm{rd}}_U$ to have events $|U|\amalg\{\rho_b:b\in B_{\mathrm{use}}\}$ and order generated by the physical order of $U$ and
\[
 d<\rho_{b_{d,p}},\qquad
 \rho_b<e\quad(b\in\pre_B(e)).
\]
Set $\supp(e)=\downarrow_Ue$, $\supp(\rho_{b_p^0})=\varnothing$, and
$\supp(\rho_{b_{d,p}})=\downarrow_Ud$. A finite set $X$ is consistent when
$\bigcup_{z\in X}\supp(z)\in\mathsf{Con}_U$.
\end{definition}

\begin{theorem}[Readiness histories and physical recovery]\label{thm:readiness-history}
The preceding construction is an event structure. Its finite configurations have the unique normal form
\begin{equation}\label{eq:readiness-normal}
 x(h,J)=h\cup\{\rho_b:b\in C_h\cup J\},
 \qquad h\in\Conf_f(U),\quad J\in\Pf(A(h)).
\end{equation}
Erasure $\varepsilon_{\mathrm{rd}}x(h,J)=h$ is surjective. Every finite physical extension lifts from every representative after finitely many readiness steps, and every finite independent physical family lifts to a cube after such preparation. Preparing an enabled event $e$ requires exactly the missing readies $\pre_B(e)\setminus J$, at most $|\dom(e)|$ steps. Hiding readiness gives
\[
 S^{\mathrm{rd}}_U\mathord\downarrow |U|=U
\]
with its original order and consistency.
\end{theorem}
\begin{proof}\vmode
\leavevmode
\begin{enumerate}[label=(\roman*)]
\item \textbf{Acyclic physical order.}
A physical--readiness--physical segment $d<\rho_b<e$ compresses to $d<e$ in $U$.
Thus the generated order is acyclic and adds no order between physical events.

\item \textbf{Finite principal pasts.}
Principal downsets are finite: each physical event has finite physical past, with finitely many readiness events at its participant conditions; a readiness event has its producer's finite past.

\item \textbf{Consistency.}
Supports increase along each generating relation.
Consistency is consequently hereditary and closed under adjoining predecessors.
Singleton consistency follows from consistency of principal histories, and the empty history supplies the empty consistency case.

\item \textbf{Construction of normal configurations.}
For \eqref{eq:readiness-normal}, every physical predecessor lies in $h$, every readiness needed by an event of $h$ belongs to $C_h$, and every readiness in $J$ has its producer in $h$.
The displayed set is therefore downward closed and its supporting physical history is $h$, proving consistency.

\item \textbf{Recovery of readiness data.}
Conversely, the physical part $h$ of a configuration is a configuration of $U$.
Readiness at every condition consumed by $h$ is mandatory.
Any additional readiness has its producer in $h$ and its condition unconsumed there, so belongs to $A(h)$.
This gives the unique $J$, which is finite precisely when the configuration is finite, since $h$ and $C_h$ are finite.

\item \textbf{Preparation of one event.}
In these coordinates the one-event extensions are exactly
\begin{align}
 (h,J)&\xrightarrow{\rho_b}(h,J\cup\{b\}),
       &&b\in A(h)\setminus J,\label{eq:readiness-step}\\
 (h,J)&\xrightarrow{e}
       (h\cup\{e\},J\setminus\pre_B(e)),
       &&h\cup\{e\}\in\Conf_f(U),\quad\pre_B(e)\subseteq J.
       \label{eq:readiness-physical}
\end{align}
If $e$ is physically enabled, all of $\pre_B(e)$ are current conditions.
Add the missing readies $\pre_B(e)\setminus J$ and then execute $e$.
These readies have distinct owners and commute.

\item \textbf{Independent families and path bounds.}
A finite independent family has disjoint participant conditions; prepare the finite union of its missing readies, then its physical events form the original cube with the same face updates.
A finite physical extension lifts by applying this argument to a linearization of its new events.
Thus preparation for $e_1,\ldots,e_n$ takes at most $\sum_i|\dom(e_i)|$ readiness steps.
This proves finite preparation and absence of added deadlock over a physically enabled history, for arbitrary $P$.

\item \textbf{Physical hiding comparison.}
The preceding order calculation shows that restriction to physical events has exactly the order of $U$.
For a finite set of physical events, the consistency test is consistency of its causal closure, equivalent to its original consistency in $U$.
This proves the hiding equality, including its configuration cubes.
In particular erasure is a history observation and a partial event map; it is not asserted to preserve the dimensions of readiness cubes.
\end{enumerate}
\end{proof}

\begin{definition}[Readiness extension categories]\label{def:readiness-extension-categories}
Define $\mathcal R_f(h)=\Pf(A(h))$ and $\overline{\mathcal R}(h)=\Pow(A(h))$, both ordered by inclusion; on $h\subseteq k$ use $J\mapsto J\cap A(k)$. Their \textbf{Grothendieck extension orders} have objects $(h,J)$ and relation
\[
 (h,J)\leq(k,K)\quad\Longleftrightarrow\quad h\subseteq k\text{ and }J\cap A(k)\subseteq K.
\]
Erasure is $(h,J)\mapsto h$, the empty section is $\ell(h)=(h,\varnothing)$, and the completed section is $r(h)=(h,A(h))$. The \textbf{vertical arrows} have $h=k$. The next propositions verify functoriality, the adjunctions and localization.
\end{definition}

\begin{proposition}[Readiness as a Grothendieck construction]\label{prop:readiness-opfibration}
On the poset $\mathcal H_U$ of finite histories under inclusion, there is a strict functor
\[
 \mathcal R_f:\mathcal H_U\longrightarrow\mathbf{Pos},\qquad
 \mathcal R_f(h)=\Pf(A(h)),\qquad
 \mathcal R_f(h\subseteq k)(J)=J\cap A(k).
\]
There is an isomorphism $\mathcal H_{S^{\mathrm{rd}}_U}\cong\int\mathcal R_f$ under which $\varepsilon_{\mathrm{rd}}$ is a split opfibration. The section $\ell(h)=(h,\varnothing)$ satisfies $\ell\dashv\varepsilon_{\mathrm{rd}}$ on history-extension posets.
\end{proposition}
\begin{proof}\vmode
\leavevmode
\begin{enumerate}[label=(\roman*)]
\item \textbf{Persistence of current conditions.}
Condition names are never recycled.
A condition current at $h$ and at a later $l$ is current at every intermediate $k$, proving the composition equation for $\mathcal R_f$.
Intersection preserves finiteness.

\item \textbf{Grothendieck order.}
The normal form gives
\[
 x(h,J)\subseteq x(k,K)
 \iff h\subseteq k\ \text{and}\ J\cap A(k)\subseteq K.
\]
A ready in $J$ is either consumed by $k$, and hence mandatory there, or survives in its current cut.
This is exactly the order of the Grothendieck construction.

\item \textbf{Opcartesian lift.}
The arrow $(h,J)\to(k,J\cap A(k))$ is opcartesian, and these arrows compose strictly.

\item \textbf{Left adjoint section.}
The comparison $\ell(h)\leq(k,K)$ holds precisely when $h\subseteq k$, proving the adjunction.

\item \textbf{Right-section boundary.}
A right section with value $(h,A(h))$ belongs to this finite-history category only when $A(h)$ is finite.
The following construction supplies the right adjoint without imposing that restriction on $P$.
\end{enumerate}
\end{proof}

\begin{proposition}[Fibrewise completion and readiness localization]\label{prop:readiness-completion}
The fibrewise ideal completion of $\mathcal R_f$ is the strict functor
\[
 \overline{\mathcal R}(h)=\Pow(A(h)),\qquad
 \overline{\mathcal R}(h\subseteq k)(J)=J\cap A(k).
\]
Its Grothendieck construction is the poset of all readiness configurations whose physical erasure is finite. On this completed poset the sections
$\ell(h)=(h,\varnothing)$ and $r(h)=(h,A(h))$ give
$\ell\dashv\overline\varepsilon_{\mathrm{rd}}\dashv r$.
Already on finite histories, erasure is a categorical localization at the vertical readiness extensions.
\end{proposition}
\begin{proof}\vmode
\leavevmode
\begin{enumerate}[label=(\roman*)]
\item \textbf{Ideal completion of finite subsets.}
An ideal here is a nonempty directed lower subset.
The maps
\[
 \operatorname{Idl}(\Pf(A))\longleftrightarrow\Pow(A),\qquad
 \mathcal J\longmapsto\bigcup\mathcal J,\qquad
 J\longmapsto\Pf(J)
\]
are inverse order isomorphisms.
For the nontrivial direction, every finite subset of $\bigcup\mathcal J$ has a bound in $\mathcal J$ by finitely many uses of directedness, and hence lies in $\mathcal J$.
This uses no simultaneous choice.

\item \textbf{Extension of transport.}
Intersection with $A(k)$ preserves directed unions, so these isomorphisms identify the completed transition maps with the displayed intersections.
This is the usual ideal-completion construction; see~\cite{effective-domains}.

\item \textbf{Completed configurations.}
The proof of the normal form applies to arbitrary $J\subseteq A(h)$: every finite subset has consistent finite physical support, so $x(h,J)$ is a configuration.
Conversely every configuration with finite physical part has that form.

\item \textbf{Adjoint sections.}
The same order formula proves the two adjunctions and the split opfibration.
In particular this is a completion of the fibres over finite physical histories, not a claim that every physical history remains finite in the whole configuration domain.

\item \textbf{Vertical arrows.}
Let $W$ consist of the arrows $(h,J)\to(h,K)$ with $J\subseteq K$.

\item \textbf{Localization factorization.}
On finite histories there is a natural transformation
\[
 \alpha:\ell\varepsilon_{\mathrm{rd}}\Longrightarrow\Id,
 \qquad \alpha_{(h,J)}:(h,\varnothing)\longrightarrow(h,J),
 \qquad \varepsilon_{\mathrm{rd}}\ell=\Id.
\]
Every component of $\alpha$ belongs to $W$.
If a functor $G$ to any category sends $W$ to isomorphisms, $G\alpha$ is the natural isomorphism
$(G\ell)\varepsilon_{\mathrm{rd}}\cong G$.

\item \textbf{Universal property.}
Conversely every functor through erasure inverts $W$.
Restriction along $\ell$ and precomposition with erasure are inverse equivalences also on natural transformations, by the same $\alpha$.
This is precisely the localization universal property.
\[
\begin{tikzcd}[column sep=large]
\int\mathcal R_f\arrow[r,hook]\arrow[d,"\varepsilon_{\mathrm{rd}}"']
 &\int\overline{\mathcal R}\arrow[d,"\overline\varepsilon_{\mathrm{rd}}"]\\
\mathcal H_U\arrow[r,equal]&\mathcal H_U
\end{tikzcd}
\]
\end{enumerate}
\end{proof}

\begin{remark}[Preparation and infinite administrative motion]\label{rem:readiness-divergence}
Finite preparation does not imply termination of every readiness-only execution. For $P=\mathbb N$ with one independent enabled one-shot event at each participant, all $b_p^0$ belong to $A(\varnothing)$ and $\rho_{b_0^0},\rho_{b_1^0},\ldots$ is an infinite readiness-only execution. Its union is the completed configuration $x(\varnothing,A(\varnothing))$. In a specified finite-participant instance, the integer $|A(h)\setminus J|\leq|P|$ strictly decreases at each readiness step with fixed physical history. Thus that instance has the finite bound; the general construction instead has finite local preparation, fibrewise saturation and the localization just proved.
\end{remark}

\begin{definition}[Readiness strategy arena]\label{def:readiness-strategy-arena}
Give the readiness events of $S_U^{\mathrm{rd}}$ negative polarity, rendezvous completions positive polarity, and other physical events their specified signs. Define $\mathcal A_U^{\mathrm{rd}}$ on the same events and consistency, with order generated by physical causality, producer-to-readiness edges, and readiness-to-physical edges whose target is negative. The \textbf{readiness display} is the identity function $\sigma_U^{\mathrm{rd}}:|S_U^{\mathrm{rd}}|\to|\mathcal A_U^{\mathrm{rd}}|$. Its candidate configuration fibres are $\Phi_U^{\mathrm{rd}}(x)=\{*\mid x\in\Conf_f(S_U^{\mathrm{rd}})\}$. The theorem proves the strategy and kernel properties.
\end{definition}

\begin{theorem}[Global readiness strategy]\label{thm:readiness-strategy}
Give readiness negative polarity and rendezvous completion positive polarity, retaining the chosen interface signs on other physical events. Let $\mathcal A_U^{\mathrm{rd}}$ have the events and consistency of $S_U^{\mathrm{rd}}$, with order generated by the physical order, producer-to-readiness relations, and readiness-to-physical relations whose target is negative. The identity event map
\[
 \sigma_U^{\mathrm{rd}}:S_U^{\mathrm{rd}}\longrightarrow\mathcal A_U^{\mathrm{rd}}
\]
is receptive and courteous. Its reached Scott kernel is
\[
 \Phi_U^{\mathrm{rd}}(x)=
 \begin{cases}
 *,&x\in\Conf_f(S_U^{\mathrm{rd}}),\\
 \varnothing,&\text{otherwise},
 \end{cases}
 \qquad x\in\Conf_f(\mathcal A_U^{\mathrm{rd}}),
\]
and its cubical unstraightening is $K(S_U^{\mathrm{rd}})$. For an isolated rendezvous this is Proposition~\ref{prop:rendezvouskernel}. Independent networks give the corresponding tensor, coherently with association and symmetry.
\end{theorem}
\begin{proof}\vmode
\leavevmode
\begin{enumerate}[label=(\roman*)]
\item \textbf{Arena and display.}
The arena order is contained in the strategy order, so has finite principal downsets and its consistency is closed under predecessors.
Each strategy configuration is an arena configuration, making the identity event map a prestrategy.

\item \textbf{Receptivity.}
If a negative event is enabled in the arena at a strategy configuration, all generating predecessors into it were retained in the arena.
They lie in that configuration, whose strategy down-closure supplies all their predecessors.
Consistency is unchanged, so the event has its unique strategy lift.

\item \textbf{Courtesy.}
Every immediate strategy dependency absent from the arena must be a removed generating relation, hence goes from negative readiness to a positive physical event.
These prove receptivity and courtesy.

\item \textbf{Kernel comparison.}
Apply \cite[Theorem~3.20]{ccrw}, in the Scott form of Theorem~\ref{thm:scott-discrete-fibrations}.
The identity event map has at most one configuration over each arena configuration, giving the displayed subterminal fibres and the unstraightening.

\item \textbf{Reachability.}
Every finite supported configuration has a linearization from the empty configuration, so the kernel is reached.

\item \textbf{Handshake specialization.}
For an isolated handshake, $B_{\mathrm{use}}$ consists of its two initial conditions; the arena has independent $l^-,r^-,c^+$ and the strategy adds $l<c$, $r<c$.

\item \textbf{Independent tensor and coherence.}
For independent networks, conditions and configuration pairs $(h,J)$ split by component.
All iterated comparisons retain the same component tuple, proving tensor coherence.
\end{enumerate}
\end{proof}

Preparation acts uniformly on local conditions, including conditions consumed by private steps. Competing messages or partners share each participant's current readiness; preparing it is internal bookkeeping. The message interface and the control arena of Section~\ref{subsec:observations-distributed-strategies} retain their own polarities and are compared through physical histories.

With the physical signs fixed, readiness erasure is intersection with $|U|$. It therefore preserves configurations, compatible unions, nonempty compatible intersections and finiteness. This proves additive stability in the sense of~\cite[Definition~5.9]{winskel-functional}, so the occurrence-map adjunction of Theorem~\ref{thm:admin-occurrence-lifts} applies.

\begin{proposition}[Why the bare diagonal is not a channel kernel]\label{prop:diagonal}
For nonempty $M$, the relation $u=v$ on unrestricted input/output stream histories is not a rooted subterminal kernel $A_M\to A_M$.
\end{proposition}
\begin{proof}\vmode
\leavevmode
\begin{enumerate}[label=(\roman*)]
\item \textbf{Root restriction contradiction.}
Its root fibre at $(\eps,\eps)$ would be inhabited.
Covariance in the input sends that element along $\eps\to m$ to an element at $(m,\eps)$.
The diagonal requires this fibre to be empty, a contradiction.
\end{enumerate}
\end{proof}

Thus synchronous completion cannot be imposed on an unrestricted receptive asynchronous interface simply by deleting all unmatched prefixes. Atomic joint moves and explicit readiness/commit interfaces give well-defined constructions. Session-game models address related interface changes using synchronous games and call-return encodings \cite{cy}; the above calculation specifies the interface used here rather than invoking a general encoding theorem without its typing data.

\subsection{Buffered implementations of prefix strategies}\label{subsec:buffered-translation}
A prefix strategy can be implemented through buffered communication only after the supplied occurrences and readiness conditions have been accounted for. The following comparison records those implementation histories and their projection to the prefix behaviour.

\begin{definition}[Rooted stream transducer]\label{def:stream-transducer}
A rooted deterministic stream transducer has a state set $Q$, \textbf{initial state} $q_0$, update $\delta:Q\times M\to Q$, and finite output $\omega:Q\times M\to N^*$. Define $q_u$ and $f(u)$ recursively by
\[
q_\eps=q_0,\quad f(\eps)=\eps,\quad
q_{um}=\delta(q_u,m),\quad f(um)=f(u)\omega(q_u,m).
\]
Its buffered implementation has states $(q,w)\in Q\times N^*$, input transitions
\[
(q,w)\xrightarrow{m?}(\delta(q,m),w\omega(q,m)),
\]
and output transitions $(q,nw)\xrightarrow{n!}(q,w)$.
\end{definition}
The map $f:C_M\to C_N$ is a functor because it preserves prefixes. No output-production step is confused with the later delivery of the produced letters.

\begin{theorem}[Buffered realization]\label{thm:buffered}
The kernel
\[
H_f(u,v)=C_N(v,f(u))
\]
is reached. Its unstraightening is the history carrier of the buffered implementation, with projection
\[
(u,v)\longmapsto(q_u,w),\qquad f(u)=vw.
\]
An input and a delivery form a square precisely when that delivery was enabled before the input.
\end{theorem}
\begin{proof}\vmode
\leavevmode
\begin{enumerate}[label=(\roman*)]
\item \textbf{Presheaf variance.}
Prefix preservation makes the formula covariant in $u$ and contravariant in $v$.

\item \textbf{Input and delivery lifts.}
At $(u,v)$ an input $m$ always lifts, since $v\leq f(u)\leq f(um)$.
A delivery of $n$ lifts exactly when $vn\leq f(u)$.

\item \textbf{Concurrent cells.}
The terminal-corner test gives the same condition for the mixed square.

\item \textbf{Queue-state projection.}
Projecting sends these operations to appending the block $\omega(q_u,m)$ and removing a previously available head.
The word identity
$\mathsf{deq}_n(nwz)=wz$ proves the commuting square, including $z=\eps$.

\item \textbf{Reachable histories.}
Every supported history pair is executable by first accepting $u$ and then delivering $v$.
All legal executions have such a pair because delivery never exceeds production.

\item \textbf{Causal orders and exchange.}
Their partial orders contain the input chain, output chain, and the dependence of each output on the input that first produced it.
Adjacent exchanges of independent input/delivery steps identify precisely the linearizations of that partial order.
Thus the display is its history carrier and is reached.
\end{enumerate}
\end{proof}

\begin{theorem}[Composition of buffered transducers]\label{thm:buffercomposition}
For stream transducers inducing $f:C_M\to C_N$ and $g:C_N\to C_L$,
\begin{equation}\label{eq:buffercomposition}
H_g\star_{A_N}H_f\cong H_{gf}.
\end{equation}
The raw composite is already reached. The isomorphism preserves roots, all visible execution cells, and the visible finite histories of an operational pipeline with an internal FIFO channel. In particular $Y_M\star Y_M\cong Y_M$.
\end{theorem}
\begin{proof}\vmode
\leavevmode
\begin{enumerate}[label=(\roman*)]
\item \textbf{Co-Yoneda comparison.}
For external prefixes $u,z$, co-Yoneda gives
\[
\int^{v\in C_N} C_N(v,f(u))\times C_L(z,g(v))
\cong C_L(z,g(f(u))).
\]
The bijection sends the class of $(v\to f(u),z\to g(v))$ to the composite $z\to g(v)\to g(f(u))$.

\item \textbf{Inverse and normalization.}
Its inverse uses $v=f(u)$ and its identity arrow.
The coend relation identifies every representative with this one, proving both inverse equations.
Equivalently, supported intermediate prefixes all map to the largest allowed prefix $f(u)$; if $z\nleq g(f(u))$, no representative exists.

\item \textbf{Naturality.}
These formulas are natural, preserve the empty-prefix roots, and show directly that the result is reached by Theorem~\ref{thm:buffered}.

\item \textbf{Forward operational adequacy.}
For operational adequacy, every pipeline execution has an intermediate prefix $v\leq f(u)$ that the second machine has consumed, and its delivered output satisfies $z\leq g(v)$.
Thus its visible pair is supported on the right.

\item \textbf{Reverse operational adequacy.}
Conversely, if $z\leq g(f(u))$, accept $u$, transfer the finite word $f(u)$ internally, and deliver $z$.
This realizes every supported visible pair.

\item \textbf{Lifting visible steps.}
More precisely any legal visible step in the buffered composite can be lifted after finitely many internal transfers: transfer only enough of the outstanding middle word to produce the requested output.
An input step requires none.

\item \textbf{Visible concurrency.}
If a visible input and output are independent in the composite, pre-transfer enough to enable the output before the input; their two residual updates then commute.
Hence the visible cells also agree after internal saturation.

\item \textbf{Hidden forwarding and coherence.}
Co-Yoneda identifies the different amounts of outstanding internal transfer that have the same visible observation.
The last equation is the case $f=g=\id$.
The composition comparisons are coherent because their maps concatenate/evaluate the same finite sequence of intermediate prefixes, or equivalently by the coherence of co-Yoneda.
\end{enumerate}
\end{proof}

This theorem concerns finite visible behavior after internal saturation. It does not identify scheduler fairness or claim that an unfair execution must eventually forward every queued message. A liveness specification can observe that distinction even when the finite-history kernel does not.

\subsection{Occurrence and resource protocol strategies}\label{subsec:occurrence-resource-protocol-strategies}
Apply the causal-interface construction of Lemma~\ref{lem:causal-interface} where indicated below; audit and public communication interfaces retain their distinct types.
\subsubsection{Example: bag and lossy-FIFO matching kernels}
Payload equality does not identify the occurrence that supplied a delivery. These kernels retain supplier injections, making the difference between bag matching and order-preserving FIFO matching explicit.


\begin{proposition}[Open delivery tests with supplier witnesses]\label{prop:open-supplier-matching}
Fix $n$ input occurrences with payloads in $D$ and $m$ requested output observations $U_1,\ldots,U_m\in\Op(D)$. For a bag the successful input tuples form
\[
 \bigcup_{\alpha:[m]\hookrightarrow[n]} W_\alpha,\qquad
 W_\alpha=\bigcap_{j=1}^{m}\pi_{\alpha(j)}^{-1}U_j
       \ \subseteq D^n.
\]
The witness family is $\coprod_\alpha W_\alpha\to D^n$. For lossy FIFO retain the increasing injections; for reliable FIFO retain just $\alpha(j)=j$ when $m\le n$.
\end{proposition}
\begin{proof}
Fixing the match identifies the $j$th delivered payload with the $\alpha(j)$th input. A tuple admits that witness exactly when it lies in the indicated finite intersection, which is open by continuity of projections. Union over legal matches gives the success open. Input extension retains the indices, and restriction to an output prefix restricts the injection; both preserve the required memberships and compose by restriction of functions. The empty output has its unique empty supplier. Substituting $D=\Delta M$ and $U_j=\{v_j\}$ identifies the fibre over $u$ with the original label-preserving injections, including their distinct suppliers when labels coincide.
\end{proof}
For example, one request for $\{a\}$ against $(a,a)$ has two bag witnesses. Their common success is a single truth value. The same pullback formula along the source map of a fixed duplication, adjacent exchange or replay history tests its deliveries. Duplication may reuse a source; exchange permutes positions; replay selects the packet named by its command. The parent, transposition and replay command remain part of the respective witnesses.

\begin{definition}[Label-preserving supplier injections]\label{def:supplier-injections}
For finite words $u\in M^n$ and $v\in M^m$, put
\[
 \mathsf{Inj}_M(v,u)=
 \{\alpha:[m]\hookrightarrow[n]:v_j=u_{\alpha(j)}\},
 \qquad [n]=\{1,\ldots,n\}.
\]
An injection records the supplying input occurrence, even when several input payloads are equal.
\end{definition}

\begin{definition}[Bag and lossy-FIFO fibre families]\label{def:matching-fibre-families}
For words $u,v\in M^*$, define $\Phi_{\mathrm{bag}}(u,v)=\mathsf{Inj}_M(v,u)$ and $\Phi_{\mathrm{loss}}(u,v)=\{\alpha\in\mathsf{Inj}_M(v,u):\alpha\text{ increasing}\}$. Input-prefix extension includes the supplier positions into the longer input; output-prefix restriction restricts the supplier function. The empty injection is the root. These specify the candidate functors $C_M\times C_M^\op\to\Set$; Theorem~\ref{thm:matching-kernels} verifies their laws and strategy interpretation.
\end{definition}

\begin{theorem}[Occurrence matching strategies]\label{thm:matching-kernels}
The formulas
\begin{align}
 \Phi_{\mathrm{bag}}(u,v)&=\mathsf{Inj}_M(v,u),\label{eq:bag-kernel}\\
 \Phi_{\mathrm{loss}}(u,v)&=
       \{\alpha\in\mathsf{Inj}_M(v,u):\alpha\text{ is increasing}\}
       \label{eq:loss-kernel}
\end{align}
define reached kernels $A_M\to A_M$. They are the configuration kernels of the reliable occurrence bag and the lossy FIFO after hiding loss events and retaining visible causal witnesses. Reliable FIFO is the subcase $\alpha(j)=j$. Input extension preserves an injection, and output restriction restricts its domain.
\end{theorem}
\begin{proof}\vmode
\leavevmode
\begin{enumerate}[label=(\roman*)]
\item \textbf{Presheaf restrictions.}
Inclusion of the old input positions into a longer input word preserves label matching and injectivity, and preserves increasingness in \eqref{eq:loss-kernel}.
Restriction to an output prefix preserves the same conditions.
These operations commute and compose strictly, giving the kernel laws and the empty matching as root.

\item \textbf{Bag causal diagrams.}
For a bag, form the finite diagram with input chain $i_1<\cdots<i_n$, output chain $o_1<\cdots<o_m$, and supply edges $i_{\alpha(j)}<o_j$.

\item \textbf{Execution of ideals.}
Every ideal is executable: an output's supplying input has occurred, and injectivity says that no previous output consumed that occurrence.
Conversely a bag execution determines its unique supplying injection.

\item \textbf{Independent exchange.}
Any two linearizations differ by adjacent exchanges of incomparable events.
Such an exchange moves an input past an output consuming an already present distinct occurrence, exactly a bag chart.

\item \textbf{Ordinary bag strategy.}
Thus principal diagrams give an occurrence strategy.
Its only additional cross-interface immediate dependencies are negative-to-positive supply dependencies; negative input extensions are unique.
It is receptive and courteous and has fibre \eqref{eq:bag-kernel}.

\item \textbf{Lossy FIFO suppliers.}
For lossy FIFO, deliveries select distinct inputs in increasing order.

\item \textbf{Necessary deletions.}
Before delivering the selected input $\alpha(j)$, delete all still-pending earlier inputs not selected by previous deliveries.
These deletions are necessary, and their occurrences are determined by $\alpha$.
Deletions of distinct packets commute, as do deletions and operations on other packets whenever their head guards persist.

\item \textbf{Visible normalization.}
After hiding, the necessary deletions contribute their supplying-input ancestry, already contained in the input chain through $i_{\alpha(j)}$.
Deletions outside the visible causal past disappear.

\item \textbf{Kernel identification.}
The resulting visible principal diagrams are exactly the increasing matching diagrams.
This proves the fibre identification and the same strategy conditions.

\item \textbf{Reachability.}
Finally every bag matching is reached by accepting all of $u$ and delivering its selected occurrences.
Every increasing matching is reached by the preceding deletion construction, proving reachability.
\end{enumerate}
\end{proof}

At $(u,v)=(aa,a)$ both kernels have two elements; replacing the fibre by a truth value loses the supplying occurrence. The full loss audit retains also deletions of packets with no subsequent visible consequence, so it is related to \eqref{eq:loss-kernel} by hiding, not by an identity of full histories.

\begin{corollary}[Several bag ports]\label{cor:multiport-bag}
For a bag with several input and output ports, replace the input and output position sets by their disjoint unions, retain each port's chain, and use one label-preserving injection from all output occurrences to all input occurrences. This gives its reached communication kernel. Outputs on distinct ports form a square precisely when their selected occurrences and participating controls are distinct and both outputs are already enabled.
\end{corollary}
\begin{proof}\vmode
\leavevmode
\begin{enumerate}[label=(\roman*)]
\item \textbf{Joint supplier diagrams.}
The preceding ideal and linearization argument uses only the local port chains and global single consumption.

\item \textbf{Independent deletions and injectivity.}
The simultaneous deletions of distinct selected occurrences commute, including when their values coincide.
A repeated selected occurrence fails injectivity and cannot form the square.
\end{enumerate}
\end{proof}

\begin{example}[Move rules from a matching fibre]\label{ex:matching-moves}
A position of the bag or lossy channel is $(u,v,\alpha)$, including its supplying injection. Input of $a$ gives $(ua,v,\iota\alpha)$, where $\iota$ includes the old input positions. Delivery of $b$ chooses an unused position $i$ with $u_i=b$ and gives
\[
 (u,v,\alpha)\xrightarrow{b!,i}(u,vb,\alpha[i]),
 \qquad \alpha[i](|v|+1)=i.
\]
For loss, $i$ must additionally exceed every previous selected position. For reliable FIFO it must equal $|v|+1$. These conditions are obtained by taking the preimages of $\alpha$ under the output-restriction map in the respective kernel, so they are exactly the positive move rule. Input covariance gives the unique negative move. The selected position is retained even if several inputs have the same payload. An input and a delivery form a square precisely when the delivery's selected position was already present.
\end{example}

\subsubsection{Example: duplication genealogies and exchange histories}
Duplication introduces new occurrences with a shared ancestor. The forest presentation records that genealogy through delivery and exchange, so that hiding can remove unused detail without losing the suppliers of visible outputs.


\begin{definition}[Normalized duplication forest]\label{def:dup-forest}
Over $(u,v)$ take an ordered finite binary forest with one root for each position of $u$. An internal node is a duplication occurrence with its two ordered, named children. Leaves are marked delivered or pending; the delivered leaves form an initial segment of the total left-to-right frontier, with payload word $v$. Normalize by contracting every copied subtree with no delivered descendant to one pending leaf. The root position followed by its binary child address is the occurrence name, using the disjoint constructors of Definition~\ref{def:anchored-instruction}. Write $\mathsf{Forest}_M(u,v)$ for the normalized forests, up to the unique order- and address-preserving isomorphisms.
\end{definition}

\begin{theorem}[Duplication strategy]\label{thm:dup-kernel}
For the unbounded duplication FIFO without losses or exchanges,
\begin{equation}\label{eq:dup-kernel}
 \Phi_{\mathrm{dup}}(u,v)=\mathsf{Forest}_M(u,v)
\end{equation}
is its reached communication kernel after hiding duplications. Input extension adds a pending root. Output restriction changes the removed delivered leaves to pending and normalizes. Its witnesses retain all duplication occurrences in the causal past of the visible deliveries.
\end{theorem}
\begin{proof}\vmode
\leavevmode
\begin{enumerate}[label=(\roman*)]
\item \textbf{Forest normalization.}
Normalization is determined solely by the roots, the surviving delivered leaves, and their ancestor trees.
Normalizing twice gives the same forest; deleting output suffixes in stages therefore gives the same restriction as deleting their union.
Adding later roots commutes with this operation, proving the presheaf laws.

\item \textbf{Duplication causality.}
A duplication consumes exactly its parent and produces the two indicated children.
Its causal diagram contains the input and output chains, creation before use, and the frontier constraints on delivery.

\item \textbf{Independent frontier updates.}
Actions on distinct live leaves commute, and a parent and its descendants cannot be concurrent.
Every ideal therefore has a legal frontier execution; conversely the named parent pointers reconstruct the forest of every execution.

\item \textbf{Visible-past pruning.}
After hiding, a copied subtree outside all delivered ancestry contributes no visible event identity and is contracted.
This is exactly the stated normalization.

\item \textbf{Strategy and reachability.}
Principal visible histories give a receptive, courteous event strategy: unbounded input extension is unique, and all new cross-port causes run from input to output.
Every normalized forest is reached by accepting its roots, duplicating in ancestral order, and delivering the indicated frontier prefix.
\end{enumerate}
\end{proof}

Erasing the genealogy gives a nondecreasing root map $\alpha:[m]\to[n]$ whose image is an initial segment. In the absence of loss, a later root cannot pass all descendants of an earlier root before their delivery. The root map nevertheless forgets execution witnesses: at $(a,a)$, arbitrarily many duplications along the leftmost branch before the first delivery yield different forests with the same map.

When loss is also present, the trees retain deleted leaves and their deletion witnesses. Theorem~\ref{thm:protocol-audit} records these combined histories before hiding.

\begin{definition}[Exchange diagram]\label{def:exchange-diagram}
An exchange diagram is a finite legal occurrence execution drawn as packet strands. Each input creates its named strand; each swap acts on two currently adjacent strands and retains its ordered-pair witness; each delivery terminates the current head strand. Along a strand retain the order of all operations using it, and retain the local input and output chains. A diagram is identified with another by interchanging adjacent operations on disjoint packet strands and distinct participant controls when both operations are enabled. The unchanged strand and control data supply the same residual operations. Its visible reduction retains the causal past of its visible occurrences, with the inherited internal strand histories in their identities. Write $\mathsf{Ex}_M(u,v)$ for these reduced diagrams.
\end{definition}

\begin{proposition}[Exchange strategy and bag observation]\label{prop:exchange-kernel}
The exchange-only unbounded FIFO has the reached communication kernel
\[
 \Phi_{\mathrm{ex}}(u,v)=\mathsf{Ex}_M(u,v).
\]
Input extension adds a strand; output restriction deletes the output suffix and its now-unused internal causal past. Forgetting the swaps defines a surjection
\[
 \Phi_{\mathrm{ex}}(u,v)\longrightarrow\Phi_{\mathrm{bag}}(u,v),
\]
which is generally not injective.
\end{proposition}
\begin{proof}\vmode
\leavevmode
\begin{enumerate}[label=(\roman*)]
\item \textbf{Adjacency and head invariants.}
The immediate predecessors of an operation are its preceding strand uses, relevant port predecessor, and the preceding head-removal history when it delivers.
Moving a packet out of a gap between two prospective adjacent strands changes one of those strands, so the preceding swap that establishes their adjacency is included in this causal past.
Thus ideals retain the required adjacency and head witnesses.

\item \textbf{Independent operations.}
Incomparable operations have disjoint strand uses and preserve each other's guards; their adjacent exchange is precisely a resource square.

\item \textbf{Principal diagrams and kernel laws.}
The principal-diagram construction consequently gives the occurrence histories.
Taking the visible causal past is idempotent and commutes with successive output restrictions; adding a later input leaves existing adjacency and head guards unchanged.
These facts give the kernel laws.

\item \textbf{Strategy and reachability.}
The visible strategy is receptive on inputs and adds only input-to-output cross-interface causes, hence is courteous.
Each diagram has a finite legal linearization, proving reachability.

\item \textbf{Bag observation.}
Each delivered strand has a unique input root and is consumed once, giving an injection.
Conversely, move each selected pending packet to the head by finitely many adjacent swaps and deliver it; this realizes every bag injection.

\item \textbf{Distinct retained exchange histories.}
With input $ab$ and output $a$, direct delivery and the swaps $ab\to ba\to ab$ followed by delivery have the same injection.
The latter swaps use the delivered strand and remain in its principal history.
They are not removed by disjoint-operation exchanges, proving noninjectivity.
\end{enumerate}
\end{proof}

\begin{example}[Copy and exchange as resource events]\label{ex:genealogy-moves}
In a full duplication history, the pending frontier records live named leaves. Copying a selected leaf $p$ replaces it by its ordered children $p0,p1$ and retains $p$ as their parent. Delivery marks the first pending leaf as delivered. These are the one-event extensions of the forest diagrams: a child cannot be used before its creation, and two operations using the same live leaf cannot be concurrent. Hiding internal copies contracts precisely the subtrees with no visible delivered descendant, giving the normalization in Theorem~\ref{thm:dup-kernel}.

For exchange, a selected adjacent pair of strands gives the move $wpqz\to wqpz$, with that ordered-pair witness retained. The delivery move terminates the first pending strand. Operations on disjoint strands and distinct participant controls commute when both are enabled, because neither changes the other's adjacency or head test. These are exactly the exchanges used in Proposition~\ref{prop:exchange-kernel}. Repeating the same swaps is a history, not an identity generated by independence.
\end{example}

\subsubsection{Example: replay through its command interface}

\begin{definition}[Replay command histories]\label{def:formal-replay-command-histories}
Replay reads a retained source without consuming it, whereas duplication consumes a live parent. Let $C_{\mathrm{cmd}}$ be the rooted category of lawful witnessed command histories, including record creation and selected replay records. A record-creation command installs a named immutable source; a replay command selects an installed source and appends a fresh attempt carrying its value without consuming that source. These clauses generate the command histories and their resource diagrams. Its arrows are history extensions; independent command steps are equated by their actual charts. Give its cubical history carrier positive command polarity, so the dual input interface has negative command moves and Scott category $C_{\mathrm{cmd}}^\op$. Let $f(c)$ be the FIFO word of attempts appended by the command history $c$, in their enqueue order. Independent command exchanges preserve this word: two appends to the same named FIFO do not commute, while a command producing no attempt leaves the word unchanged. Thus $f:C_{\mathrm{cmd}}\to C_M$ is a prefix-preserving functor.
\end{definition}

\begin{proposition}[Retained-record replay kernel]\label{prop:replay-kernel}
On the command-history interface the buffered replay processor has kernel
\begin{equation}\label{eq:replay-kernel}
 \Phi_{\mathrm{replay}}(c,v)=C_M(v,f(c)).
\end{equation}
The command $c$ retains the selected record, its source occurrence, and the controller history. Composing with a controller and hiding their shared command interface internalizes replay choice while retaining its visible causal witnesses.
\end{proposition}
\begin{proof}\vmode
\leavevmode
\begin{enumerate}[label=(\roman*)]
\item \textbf{Kernel variance.}
Postcomposition along $f(c)\leq f(c')$ and precomposition along output prefixes give the two kernel actions.

\item \textbf{Ordinary strategy.}
The same commuting diagrams as in Proposition~\ref{prop:native-prefix} give their laws.

\item \textbf{Reachability.}
A supported pair is reached by executing $c$ and then delivering $v$.

\item \textbf{Command and delivery updates.}
A legal command appends its fresh attempt and keeps the selected retained record; the selected source record is read without being consumed.
Existing head delivery commutes with append, giving the mixed cells.

\item \textbf{Retained command histories.}
The full command-history arena retains its other command cells; the formula neither imposes an order on independent zero-output commands nor forgets which record a replay selected.
Interaction and hiding give the stated controller construction by Theorem~\ref{thm:ordinarycomp} whenever the interacting interfaces are presented by the ordinary strategies above; the full audit of the composite is also given directly by Theorem~\ref{thm:protocol-audit}.
\end{enumerate}
\end{proof}

After an input $a$ has been delivered and no live descendants remain, duplication cannot produce another packet. Replay can do so from a retained record. Root-value repetition alone therefore does not distinguish these two resource disciplines; the command and genealogy witnesses do.

\subsubsection{Example: capacity arenas and overflow histories}\label{subsec:capacity-strategies}
A capacity bound changes which inputs can be accepted. We first show why the unrestricted stream arena cannot express the required behaviour, then construct interfaces that retain capacity and overflow histories.


\begin{proposition}[Unrestricted-stream obstruction]\label{prop:capacity-interface}
A blocking buffer with positive finite capacity cannot be represented exactly by a receptive strategy on the unrestricted independent input/output stream arena. For positive capacity, the drop-new and drop-oldest presentations also fail to have an exact witnessed display on that arena.
\end{proposition}
\begin{proof}\vmode
\leavevmode
\begin{enumerate}[label=(\roman*)]
\item \textbf{Blocking-buffer obstruction.}
At a full blocking buffer the next negative stream input is arena-enabled but has no protocol lift.
Equivalently, the proposed support $v\leq u$, $|u|-|v|\leq k$ is not closed under input extension.

\item \textbf{Drop-new obstruction.}
For drop-new, start with a full queue and a coenabled head delivery and arrival.
Arrival then delivery discards the new packet, whereas delivery then arrival stores it.
These are distinct witnessed target states.

\item \textbf{Drop-oldest obstruction.}
For drop-oldest, arrival consumes the head selected by the competing delivery, disabling that selected witness.
Neither pair forms a resource square.

\item \textbf{Forced-square contradiction.}
In an unrestricted product arena the mixed square exists; its Scott-terminal corner is the configuration after the old delivery and before the new input.
An element there forces a lifted square by Lemma~\ref{lem:terminalcorner}.
Hence no exact display over that arena can have precisely the protocol cells.
\end{enumerate}
\end{proof}

\begin{definition}[Blocking-capacity arena]\label{def:capacity-arena}
For $k\geq0$, let $B_k$ have the two message-labelled stream chains and the causal constraints
\begin{equation}\label{eq:capacity-arena-order}
 o_j^+<i_{j+k}^-\qquad(j\geq1).
\end{equation}
More explicitly, its finite \textbf{configuration family} consists of compatible word pairs $(u,v)$ with $|u|\leq |v|+k$, without imposing matching payloads or $v\leq u$. Inside a common pair these prefixes are ideals of the two chains with \eqref{eq:capacity-arena-order}; taking principal histories constructs the arena event structure. Its root is $(\eps,\eps)$.
\end{definition}

\begin{theorem}[Blocking-buffer strategy]\label{thm:bounded-kernel}
On $\Sc(B_k)$ the subterminal presheaf
\begin{equation}\label{eq:bounded-kernel}
 \Phi_k(u,v)=\begin{cases}*,&v\leq u,\\\varnothing,&\text{otherwise}\end{cases}
\end{equation}
has reached display equal to the blocking FIFO history carrier. For $k>0$ it is already reached. Its strategy adds the supply dependencies $i_j^-<o_j^+$. At occupancy $q$ its mixed input/output square exists exactly when $0<q<k$. For $k=0$, its reached part is the root only.
\end{theorem}
\begin{proof}\vmode
\leavevmode
\begin{enumerate}[label=(\roman*)]
\item \textbf{Presheaf restrictions.}
An arena-allowed negative input extension preserves $v\leq u$.
Removing a maximal positive output also preserves it.
These generating restrictions define the presheaf.

\item \textbf{Positive-capacity schedules.}
For $k>0$, a supported pair can be executed by accepting and delivering its matched prefix, then accepting its remaining at most $k$ pending letters.
More generally its legal schedules are exactly the linearizations of the two port chains with supply edges and \eqref{eq:capacity-arena-order}: supply ensures a packet exists and the arena edge ensures a slot has been released.

\item \textbf{Exchange and strategy.}
Adjacent incomparable input and output steps have an old head and a free slot, so are exactly the resource squares.
Principal histories give the ordinary strategy, with the stated kernel by Theorem~\ref{thm:scott-discrete-fibrations}.

\item \textbf{Mixed cells.}
A mixed pair is initially enabled precisely at $0<q<k$, and its two updates commute there.

\item \textbf{Zero-capacity reachability.}
At $k=0$, nonempty supported pairs have equal input and output lengths, but the first delivery cannot execute before supply and the first input cannot execute before capacity release.
No progress word from the root reaches such a pair.
Thus $\Reach(\Phi_0)$ is the one-configuration strategy, agreeing with the zero-capacity blocking rules rather than the separate rendezvous rule.
\end{enumerate}
\end{proof}

\begin{definition}[Overflow branch diagrams]\label{def:overflow-diagram}
Fix $k>0$ and one input and one output port. A finite diagram records inputs $i_1,\ldots,i_n$, outputs $o_1,\ldots,o_m$, and the selected space/full branch of every input. Set $s_0=0$ and let $s_l$ count the space-branch inputs among $i_1,\ldots,i_l$. Let $t_j$ be the position of the $j$th space-branch input. Require $m\leq s_n$ and include the port chains and availability edges $i_{t_j}<o_j$. A space input $i_l$ has the additional predecessor $o_{s_l-k}$ when $s_l>k$. A full input requires $s_l\geq k$, and, with $r=s_l-k$, has
\begin{equation}\label{eq:overflow-cut}
 o_r<i_l<o_{r+1}.
\end{equation}
The marker $o_0$ is omitted; a required lower output must be present; the upper relation applies whenever that output is present, including in later extensions. Retain exactly the resulting acyclic diagrams. Payload evaluation uses ordinary append on a space branch, and respectively unchanged queue or head replacement on a full drop-new or drop-oldest branch.
\end{definition}

\begin{theorem}[Overflow-policy causal realization]\label{thm:overflow-strategy}
Every diagram of Definition~\ref{def:overflow-diagram} has legal linearizations with the recorded branches. All those linearizations give the same delivered payloads and residual queue. Conversely the diagrams are precisely the finite witnessed histories of the corresponding overflow policy modulo its resource squares. Principal diagrams construct an occurrence event structure $S_{\mathrm{new},k}$ or $S_{\mathrm{old},k}$, and Lemma~\ref{lem:causal-interface} gives its reached ordinary strategy on a protocol-aware arena. Every arrival retains its original payload label and its capacity-branch history.
\end{theorem}
\begin{proof}\vmode
\leavevmode
\begin{enumerate}[label=(\roman*)]
\item \textbf{Occupancy invariant.}
After $l$ inputs and $d$ deliveries, occupancy is $s_l-d$: the space branch increases it by one and either full branch preserves it.

\item \textbf{Space and full input bounds.}
At a space input, the lower output bound gives $d\geq s_l-k$, hence pre-input occupancy $s_l-1-d<k$.
At a full input, \eqref{eq:overflow-cut} gives exactly $d=s_l-k$, hence pre-input occupancy $k$.

\item \textbf{Delivery availability.}
Before $o_j$, its availability predecessor has occurred, so $s_l\geq j$ and occupancy is positive.
The inequalities also keep occupancy nonnegative.
This proves legality of every linearization by induction.

\item \textbf{Exchange of linearizations.}
Two linearizations of a finite partial order differ by adjacent exchanges of incomparable events.
Port chains exclude two inputs or two outputs from such a pair.

\item \textbf{Full and space branches.}
A full input is ordered against every output by \eqref{eq:overflow-cut}.
An incomparable space input and output consequently meet at an occupancy strictly between zero and $k$.
Appending and removing the old head commute there, with identical occurrence names and payload results.
Thus evaluation is independent of linearization.

\item \textbf{Recovery from a legal run.}
In a legal run, the branch at each input is fixed.
A full branch has exactly $s_l-k$ preceding deliveries; a space branch needs the displayed lower bound; and the $j$th delivery needs the $j$th net occupancy increase.
Therefore the run is a linearization of its constructed diagram.

\item \textbf{Well-defined evaluation.}
The preceding exchange calculation shows that all its linearizations lie in its resource-chart history class and that every permitted adjacent exchange stays in that class.

\item \textbf{Principal-history realization.}
Ideals retain the branch bounds and availability witnesses; their diagrams are legal prefixes.
Principal histories and consistency by common extension therefore construct the stated event structures.

\item \textbf{Causal interface.}
Apply Lemma~\ref{lem:causal-interface} to their inherited signs.
\end{enumerate}
\end{proof}

For $k=0$, drop-new has the negative input chain alone and discards every arrival, giving the stream kernel supported exactly at empty output. The paper's drop-oldest policy is defined only for $k>0$. The all-space case of the diagrams recovers the supply and release orders of Theorem~\ref{thm:bounded-kernel}.

\begin{example}[Capacity branches at a game configuration]\label{ex:capacity-moves}
Evaluate a supported history at its pending word $w$. A space-branch arrival of $a$ replaces $w$ by $wa$ when $|w|<k$. At $|w|=k>0$, drop-new keeps $w$, while drop-oldest replaces $bz$ by $za$; both retain the arrival and its selected full-branch witness. A delivery replaces $bz$ by $z$. The order bounds in Definition~\ref{def:overflow-diagram} put a full arrival between the deliveries immediately before and after its full cut, and permit precisely these extensions. A space arrival and an old delivery form a square only at $0<|w|<k$. These formulas evaluate game histories; they do not identify distinct histories having the same pending word.
\end{example}

\subsubsection{Example: credit circulation as an interface}

\begin{definition}[Slot-labelled credit words]\label{def:formal-slot-labelled-credit-words}
For the finite slot set $C$, retain slot-labelled inputs and deliveries. Let $u,v\in(C\times M)^*$ and $z\in C^*$ record accepted sends, deliveries, and completed returns. Write $\#_a$ for the count of slot $a$ and $\operatorname{slots}$ for the payload-forgetting word map.
\end{definition}

\begin{definition}[Credit arena]\label{def:credit-arena}
The arena $B_C^{\mathrm{cr}}$ has three port chains, with sends negative and deliveries and completed returns positive. Its compatible finite configuration triples satisfy
\begin{equation}\label{eq:credit-arena}
 z\leq\operatorname{slots}(v),\qquad
 \#_a(u)\leq\#_a(z)+1\quad(a\in C).
\end{equation}
Equivalently, the $j$th delivery precedes the $j$th return, with its slot value, and the $t$th return of slot $a$ precedes its $(t+1)$st send. Principal histories of these compatible prefix triples give the arena. The strategy additionally requires $v\leq u$.
\end{definition}

\begin{theorem}[Credit-return strategy]\label{thm:credit-kernel}
The indicator of $v\leq u$ on $B_C^{\mathrm{cr}}$ is a reached Scott kernel. Its moves accept a send on a currently available slot, deliver the next supplied packet, and return the next delivered slot. The updates move the same slot from the free set to the forward word, from its head to the reverse word, and from that head back to the free set. The per-slot causal cycle is
\begin{equation}\label{eq:credit-cycle}
 s_{a,t}^-<d_{a,t}^+<r_{a,t}^+<s_{a,t+1}^-.
\end{equation}
Here $Q$ is the undelivered suffix of $u$, $R$ is the unreturned suffix of $\operatorname{slots}(v)$, and $F$ is the set of free slots. At a supported triple,
\[
 |Q|=|u|-|v|,\qquad |R|=|v|-|z|,\qquad
 |F|=|C|-|u|+|z|.
\]
\end{theorem}
\begin{proof}\vmode
\leavevmode
\begin{enumerate}[label=(\roman*)]
\item \textbf{Prefix-ideal arena.}
Within a common triple, the arena prefixes are ideals of the stated delivery-to-return and return-to-send order together with the port chains.
They are closed under compatible unions and intersections, so principal histories construct the arena.

\item \textbf{Supply dependencies and strategy.}
Adding $s_{a,t}<d_{a,t}$ supplies the remaining edges of \eqref{eq:credit-cycle}.
These are negative-to-positive edges, and an arena-enabled negative send has precisely its slot available; it has a unique lift.
This proves receptivity and courtesy, hence the indicator is a Scott presheaf.

\item \textbf{State recovery.}
The three words recover the actual state: $Q$ is the suffix of $u$ after $v$, $R$ is the suffix of $\operatorname{slots}(v)$ after $z$, and $F$ is the set of slots not occurring in those suffixes.
The inequalities and $v\leq u$ ensure that no slot occurs more than once in their disjoint union.

\item \textbf{Soundness of updates.}
Conversely every rule updates those prefixes exactly as specified.

\item \textbf{Execution of supported triples.}
A supported triple is reached by accepting and delivering the prefix $v$, returning the required prefix $z$ early enough to permit later uses of its slots, then leaving the remaining sends and returns in their indicated queues.
More explicitly, process sends in their order and, whenever a slot is needed again, execute the forced earlier deliveries and return prefix through its preceding return.
The arena constraints ensure that these lie within $v,z$; FIFO order makes the procedure legal.
Complete the remaining specified deliveries and returns.
Thus every supported triple is executable.

\item \textbf{Independent resource squares.}
The state recovery and distinct-resource commutation identify its cubes, and counting the three locations gives the formulas.
\end{enumerate}
\end{proof}

If free-slot choice is made by the medium, expose a positive slot grant before the corresponding negative data action; if the sender supplies a slot, retain it in the negative action. Erasing slot names is a subsequent observation map. Collapsing several free-slot witnesses into several lifts of one bare negative input would violate receptivity. A delivered packet whose credit has not returned still prevents reuse, as the middle two positions of \eqref{eq:credit-cycle} show.

\begin{example}[The three credit moves]\label{ex:credit-moves}
At a supported credit configuration the legal extensions are
\[
 \frac{\#_a(u)=\#_a(z)}{(u,v,z)\xrightarrow{\mathsf{send}(a,m)}(u(a,m),v,z)},
 \quad
 \frac{v(a,m)\leq u}{(u,v,z)\xrightarrow{\mathsf{deliver}(a,m)}(u,v(a,m),z)},
\]
\[
 \frac{za\leq\operatorname{slots}(v)}
      {(u,v,z)\xrightarrow{\mathsf{return}(a)}(u,v,za)}.
\]
The first condition says that the slot has no outstanding send. The second is the supply dependency, and the third is the reverse-channel prefix dependency. Substitution in \eqref{eq:credit-arena} and $v\leq u$ proves that they are precisely the one-event configuration extensions. In particular delivery alone does not re-enable the slot's next send.
\end{example}

\subsubsection{Example: asynchronous multicast and atomic joint communication}
Multicast can share an accepted source while allowing recipient deliveries to proceed separately. The kernels below identify which supplier and dispatch choices remain common, and distinguish those histories from an atomic joint communication.


For one source payload and a finite set of recipient predicates $U_j$, the common payload test is the inverse image of $\prod_j U_j$ along the diagonal $D\to D^J$, namely $\bigcap_jU_j$ (here $J$ indexes the tested recipients). For several sources, fix the supplier in each lane and take the same pullbacks along its coordinate projections. Union over allowed supplier choices gives the success open by Proposition~\ref{prop:open-supplier-matching}. In multicast the witnesses retain the separate dispatch and delivery events. In atomic joint communication the tested participants share a single readiness-dependent commitment. The equality of payload tests in these two cases therefore preserves their distinct scheduling rules.

\begin{theorem}[Multicast kernels]\label{thm:multicast-kernels}
For accept-and-fanout with FIFO recipient lanes, the reached kernel is
\begin{equation}\label{eq:multicast-fifo}
 \Phi_{\mathrm{mc}}(u,(v_j)_{j\in J})
       =\prod_{j\in J}C_M(v_j,u).
\end{equation}
For bag recipient collections it is
\begin{equation}\label{eq:multicast-bag}
 \Phi_{\mathrm{mc,bag}}(u,(v_j)_{j\in J})
       =\prod_{j\in J}\mathsf{Inj}_M(v_j,u).
\end{equation}
For separate selectable pending dispatches followed by FIFO recipient lanes, put $L=\max_j|v_j|$, with $L=0$ when $J$ is empty. After hiding fanout the reached kernel is
\begin{equation}\label{eq:multicast-dispatch}
 \Phi_{\mathrm{disp}}(u,(v_j))=
 \{d:[L]\hookrightarrow[|u|]: (v_j)_r=u_{d(r)}
                 \text{ for all }j,\ r\leq|v_j|\}.
\end{equation}
The witness $d$ is a single common dispatch order, not a separate choice at each recipient.
\end{theorem}
\begin{proof}\vmode
\leavevmode
\begin{enumerate}[label=(\roman*)]
\item \textbf{Recipient supplier choices.}
Accept-and-fanout produces one child per recipient.
Each FIFO recipient consumes its children in input order, giving \eqref{eq:multicast-fifo}; each bag recipient chooses an injection, giving \eqref{eq:multicast-bag}.
Injectivity is required within each recipient, while different recipients may share a source because their children are distinct.

\item \textbf{Restrictions and reachability.}
Input extension and output-prefix restriction give the kernel actions.
Accept all inputs and then perform the prescribed recipient deliveries to prove reachability.

\item \textbf{Ordinary causal diagrams.}
Their causal diagrams have each source before its designated children and each recipient's local chain; distinct recipient children can be consumed independently.
These diagrams give receptive, courteous strategies by the matching argument of Theorem~\ref{thm:matching-kernels}.

\item \textbf{Common dispatch order.}
In the separate-dispatch presentation, two fanouts into the same FIFO lanes have an actual common append order.
Every lane sees a prefix of that order.

\item \textbf{Visible dispatch prefix.}
The longest delivered lane determines all dispatches in the visible causal past, giving the injection $d$ of length $L$.
A dispatch after that prefix has no visible descendant and disappears under hiding.

\item \textbf{Converse and output restriction.}
Conversely accept $u$, select dispatches in order $d$, and consume the indicated lane prefixes.
Output restriction truncates $d$ when the maximum consumed prefix decreases; successive truncations and input inclusions commute.

\item \textbf{Shared fanout causality.}
Its diagrams retain one source-to-fanout edge and the common fanout order needed by the FIFO lanes, so the same ideal argument identifies the reached strategy.

\item \textbf{Bag lanes and audit distinction.}
For bag lanes, independent fanouts commute and their hidden presentation has the fibre \eqref{eq:multicast-bag}.
Full audits still distinguish a separate fanout from atomic accept-and-fanout.
\end{enumerate}
\end{proof}

\begin{example}[Joint events with no participant]
For the set $R_0$ of participant-free joint records, introduce events $(r,n)$, $r\in R_0$, $n>0$, ordered by $n$ within each fixed $r$ and with all finite subsets consistent. A finite configuration records the number of occurrences of each record, and different records are independent. Tensor this factor with the participant-bearing strategy. Proposition~\ref{prop:nullary-joint} proves the exact comparison with the automata's repeated state-identity events and their cubes.
\end{example}

\begin{proposition}[Finite joint strategies]\label{prop:joint-strategy}
For every atomic joint record with participant set $I=S\cup R$, shared readiness gives negative occurrences $\rho_p$ for its current participant conditions and one positive commit $c$, with
\[
 \rho_p<c\quad(p\in I).
\]
For empty $I$, commits are the positive occurrence chains of Proposition~\ref{prop:nullary-joint}, with no readiness events. The resulting subterminal strategy is the indicator of configurations in which each present commit has its readies. Competing nonempty commits share current readies and retain their resource conflicts. Hiding readiness recovers precisely the atomic joint execution and its independent-family cubes, including all participant-free repetitions. For unreliable atomic broadcast, each permitted receiver subset gives one alternative commit.
\end{proposition}
\begin{proof}\vmode
\leavevmode
\begin{enumerate}[label=(\roman*)]
\item \textbf{Nonempty-support histories.}
First use the nonempty-support factor of Proposition~\ref{prop:nullary-joint}.
Its histories have participant chains and one event per selected joint record.
Every ideal executes by a linearization, since the local predecessors are present and disjoint blocks commute; occurrence rigidity follows from a rank in any participating chain.

\item \textbf{Readiness construction.}
Apply Definition~\ref{def:shared-readiness}.
The normal-form, lifting, and hiding arguments of Theorems~\ref{thm:readiness-history} and~\ref{thm:readiness-strategy} apply to each finite block.

\item \textbf{Receiver alternatives.}
Each receiver-subset alternative shares its current sender condition, so the required conflict is retained without an earlier partner reservation.

\item \textbf{Participant-free factor.}
Give $E_0^{\mathrm{joint}}$ positive polarity and use its identity event map as a strategy on that arena.
It is receptive because there are no negative moves, and courteous because its causality is already the arena's causality.

\item \textbf{Tensor strategy.}
Tensor it with the nonempty-support readiness strategy.
Receptivity holds factorwise with unique lifts; every immediate dependency lies in one factor, so courtesy holds factorwise as well.

\item \textbf{Hiding and Scott characterization.}
Both event maps are inclusions on configurations, giving a subterminal kernel.
Hiding only readiness acts on the first factor, and the product isomorphism of Proposition~\ref{prop:nullary-joint} proves exact recovery of full joint histories and cubes.
The standard Scott characterization is \cite[Theorem~3.20]{ccrw}.

\item \textbf{Resource-sensitive scope.}
Resource-sensitive joint records use their resolved resource histories and Lemma~\ref{lem:causal-interface}; their unrefined mixed audit remains Theorem~\ref{thm:protocol-audit}.
The empty-support factorization above concerns the state-identity joint records of the joint communication construction above, not autonomous losses or other medium operations that consume resources.
\end{enumerate}
\end{proof}

Atomic joint delivery is one event. In contrast, multicast's different child deliveries remain independent after their common production. The two constructions consequently retain different cubes even when a coarse observation reports the same recipients.

\subsubsection{Example: cancellation, service status, and allocation}

\begin{definition}[Cancellable service slot]\label{def:cancellable-service-slot}
The \textbf{cancellable slot display} has events $\{r,c_0,d,c_1\}$ with signs $r^-,c_0^-,d^+,c_1^-$, labelled respectively $\mathsf{request}$, $\mathsf{preCancel}$, $\mathsf{complete}$ and $\mathsf{pendingCancel}$. Its strict order is generated by $r<d,r<c_1$; its conflict is the symmetric upward closure of $c_0\#r$ and $c_1\#d$. Finite sets are consistent iff conflict-free. The \textbf{slot arena} retains $r<c_1$ and the same conflicts. The variant without pending cancellation restricts both carriers to $\{r,c_0,d\}$.
\end{definition}

\begin{proposition}[A service-slot strategy]\label{prop:service-strategy}
For one service slot with pending cancellation, let
\[
 r^- =\mathsf{request},\quad c_0^-=\mathsf{preCancel},\quad
 d^+=\mathsf{complete},\quad c_1^-=\mathsf{pendingCancel}.
\]
Give these events order $r<d$, $r<c_1$, and conflict generated by
\[
 c_0\mathrel{\#}r,\qquad c_1\mathrel{\#}d.
\]
Their principal histories give exactly the service-slot executions. The arena retains $r<c_1$ and all conflicts; the strategy supplies $r<d$. Its kernel is the indicator of its configurations. The policy without pending cancellation uses the same construction without the event $c_1$.
\end{proposition}
\begin{proof}\vmode
\leavevmode
\begin{enumerate}[label=(\roman*)]
\item \textbf{Slot configurations.}
The configurations are the empty fresh state, the requested pending state, request followed by completion, pre-cancellation, and request followed by pending cancellation.
The inherited conflicts exclude every illegal simultaneous terminal status.

\item \textbf{Extensions and causal interface.}
Their one-event extensions are exactly the displayed service rules.
The arena is the causal-interface construction of Lemma~\ref{lem:causal-interface}, proving the strategy and kernel assertions.
\end{enumerate}
\end{proof}

Named packet cancellation and delivery consume the same occurrence, so they conflict. Their protocol arena retains the named-occurrence availability and the terminal-status conflict. A successful cancellation is consequently a guarded negative action. A request/answer interface that also accepts unsuccessful cancellation attempts can be constructed separately, but its additional request and response events are not identified with the paper's successful-cancellation record. Distinct slots tensor when their participating controls and resources are disjoint; a shared sequential controller retains its chain.

\begin{definition}[Allocation and use interface]\label{def:allocation-use-interface}
For each witnessed allocation record, the \textbf{allocation interface} consists of shared negative participant readies followed by its positive commit $\mathsf{alloc}(j)$, retaining the selected name $j$, consumed permission and initial component. A subsequent negative operation of that component is labelled $\mathsf{use}(j)$ and has $\mathsf{alloc}(j)$ as an arena predecessor. Competing consumptions of one permission or one available name are inconsistent. The events are the principal resource histories of these records, so their labels retain the actual suppliers. Proposition~\ref{prop:allocation-strategy} proves their strategy and erasure properties.
\end{definition}

\begin{proposition}[Allocation strategy]\label{prop:allocation-strategy}
The allocation rules have a name-indexed protocol strategy. Use shared request readiness followed by a positive allocation commit carrying its selected unused name, permission, and initial component. For every later negative use of the created component the arena retains
\[
 \mathsf{alloc}(j)^+<\mathsf{use}(j)^-.
\]
Finite configurations activate finitely many names. Conflicts retain competition for a permission or a name, and disjoint allocations with disjoint participating controls give product cubes. Erasing preparation gives the original allocation records and states.
\end{proposition}
\begin{proof}\vmode
\leavevmode
\begin{enumerate}[label=(\roman*)]
\item \textbf{Allocation resources.}
Represent each available allocation permission and each as-yet-unused name by its initial resource occurrence.

\item \textbf{Consumption and local use.}
An allocation consumes its selected permission and unused-name occurrence, and produces the named component's initial local condition.
Subsequent uses advance that component's condition.
This records the permission, unused-name, and local-condition guards of the allocation rule; it introduces no choice of a preferred name.

\item \textbf{Resource diagrams and independence.}
The associated diagrams retain each condition's producer, all selected consumed occurrences, and participant chains.
Their ideals are executable, and disjoint updates commute, so principal histories give the occurrence presentation.

\item \textbf{Readiness and polarity.}
Insert shared readiness for the request and apply Lemma~\ref{lem:causal-interface}.
A positive creation preceding a negative use remains an arena edge.

\item \textbf{Finite pasts and consistency.}
Each finite diagram contains only finitely many allocations, proving finite activation even with an infinite name supply.
Competing consumptions conflict, while disjoint allocations preserve each other's guards and names.

\item \textbf{Erasure and state projection.}
The readiness erasure and state projection perform precisely the original atomic allocation update.
\end{enumerate}
\end{proof}

\begin{construction}[Resource-state evaluation of resolved diagrams]\vmode\label{con:resolved-diagram-evaluation}
\leavevmode
\begin{enumerate}[label=(\roman*)]
\item \textbf{Resolved history carrier.}
For each of the following concrete media, a finite resolved diagram
retains the selected consumed occurrences and its local port order.

\item \textbf{State-evaluation table.}
Define its state evaluation on an ideal $J$ by the following records.
\begin{center}\small
\begin{tabular}{@{}p{.18\linewidth}p{.74\linewidth}@{}}
\toprule Medium & State record $q(J)$\\\midrule
Duplication FIFO & The left-to-right list of produced, unconsumed live
leaves, with root/child addresses and payloads. A copied parent is
replaced by its two ordered children.\\
Exchange FIFO & The current word of named strands, obtained by the
recorded adjacent transpositions, with delivered strands removed.\\
Overflow & The pending occurrence word and the retained space/full
tag of each arrival. Full drop-new leaves the word unchanged;
full drop-oldest removes its head and appends the arriving occurrence.\\
Credit & The three locations $(F,Q,R)$: free slots, forward packet
word and reverse slot word. Each slot occurs in exactly one location.\\
Multicast & The pending dispatch set, where present, and each
recipient's word or bag of produced, unconsumed child occurrences.\\
Allocation & The finite allocated-name set, consumed permission
occurrences and current local conditions of created components.\\
\bottomrule
\end{tabular}
\end{center}
For duplication and exchange the record is evaluated by any
linearization; Lemma~\ref{lem:resolved-resource-invariants} proves
independence of that choice.

\item \textbf{Credit-state specialization.}
For credit it is equivalently recovered
from the three prefixes in Theorem~\ref{thm:credit-kernel}.

\item \textbf{Witnessed extensions.}
Each enabled extension carries its original rule witness, so evaluation
also specifies a map on witnessed edges.
\end{enumerate}
\end{construction}

\begin{lemma}[Local invariants for the concrete resource diagrams]\label{lem:resolved-resource-invariants}
For the concrete diagrams just listed, every ideal has a legal state
evaluation. Independent adjacent exchanges preserve that evaluation
and the residual selected witnesses. Principal restriction preserves
all selected supplier and occurrence identifiers. Consequently
evaluation gives a graph map on histories and a cubical map on their
constructed resource charts.
\end{lemma}
\begin{proof}\vmode
\leavevmode
\begin{enumerate}[label=(\roman*)]
\item \textbf{Induction invariant.}
The invariant is proved by induction on an executable prefix, using
the following guard calculations.

\item \textbf{Duplication updates.}
For duplication, initially the live list consists of the input roots
already accepted.
A copy deletes its selected live parent and inserts
its two children in that position.

\item \textbf{Creation-before-use invariant.}
Creation-before-use prevents a
child from being used earlier; the frontier constraint on a delivery
puts every preceding live leaf's removal in its causal past.
Thus
each ideal retains the operations needed to make its selected parent
live or its selected delivery the head.

\item \textbf{Independent leaf updates.}
Replacement at distinct live
leaves commutes.
Deletion of an old head commutes with a later
replacement, whereas a parent and its child cannot both be initially
selected.

\item \textbf{Exchange adjacency.}
For exchange, the two selected strands are adjacent when their
recorded swap is performed.
Any earlier operation establishing their
adjacency uses one of those strands, except removal of a head
predecessor, whose head-removal dependency is retained.
An operation
on disjoint strands cannot insert a strand between an already
adjacent selected pair without moving an endpoint.

\item \textbf{Exchange legality.}
Thus the strand
and head predecessors retain its guard in every ideal.

\item \textbf{Independent swaps.}
Disjoint
adjacent transpositions act on disjoint positions and commute;
head deletion commutes precisely with a disjoint swap still enabled
after that deletion.
A swap using the head fails that independence
test.
This proves legality and the exact exchange square.

\item \textbf{Overflow occupancy.}
For overflow, write $s_l$ for the number of space branches and $d$
for the number of deliveries in the ideal.
Its occupancy is $s_l-d$.

\item \textbf{Branch and availability bounds.}
The bounds in~\eqref{eq:overflow-cut} force a full arrival to occur
at occupancy $k$; the space predecessor bound forces its occupancy
to be less than $k$.
The availability predecessor forces a delivery
to have a head.

\item \textbf{Overflow independence.}
Full arrivals are ordered against every delivery.
The only mixed independent case is therefore a space append and an
old head removal with $0<|Q|<k$, where both orders give the same
named suffix followed by the new occurrence.

\item \textbf{Credit partition.}
For credit, send moves one slot from $F$ to the tail of $Q$,
delivery moves the head slot of $Q$ to the tail of $R$, and return
moves the head slot of $R$ to $F$.
These updates preserve the
disjoint partition of the slot set.

\item \textbf{Credit causality and independence.}
The supply and return-to-next-send
edges ensure the selected location exists in every ideal.
Updates
on different slots commute exactly when the two word-head guards
survive; an append and an earlier head removal give the FIFO square.
The same slot cannot be used simultaneously at two locations.

\item \textbf{Multicast children and suppliers.}
For multicast, a production creates distinct children $(i,j)$ for
recipients $j$.
A child delivery retains its production predecessor
and its recipient's local predecessors.
FIFO recipients therefore
remove their actual heads, while bag recipients remove distinct
selected children.

\item \textbf{Dispatch order.}
Separate dispatches into common FIFO lanes
retain their common append order.

\item \textbf{Independent lane updates.}
Incomparable updates either act
on different lanes or are an append and an older head removal.
They commute on the explicit lane records.
No independent pair
reverses the retained common dispatch order.

\item \textbf{Allocation resources.}
For allocation, each commit deletes its selected permission and
unused-name resources and creates its named initial condition.
All later uses retain that condition's producer.

\item \textbf{Allocation independence.}
Distinct commits
can commute only when their permissions, names and participant
conditions are disjoint; the two finite-set insertions and
deletions then commute.
Competition for either resource is excluded
by the recorded consistency relation.

\item \textbf{Evaluation by linearization.}
These calculations permit any linear extension of an ideal:
every selected resource and every guard has its stated predecessors.
Two such linear extensions are connected by swaps of incomparable
events, and the calculations prove identical residual records at
each swap.
Hence evaluation is well-defined.

\item \textbf{Restriction and cube compatibility.}
A downclosed restriction
does not change a retained event's supplier or principal past.
At a chart, evaluate each subset; the equal-subset equations are
exactly those just proved.
Executing a subset and then restricting
directions therefore commutes with evaluation, proving the cubical
claim.

\item \textbf{Event-structure realization.}
Applying Lemma~\ref{lem:admin-family} to these resolved
families uses this verified ideal description.
\end{enumerate}
\end{proof}

\subsubsection{Example: alternating-bit communication as a composite strategy}
The alternating-bit example combines processes with lossy data and acknowledgement media. Its retained occurrences allow the proof to track retries, acknowledgements and delivered prefixes through the composite strategy.


\begin{construction}[Alternating-bit processes and their composite]\vmode\label{con:alternating-bit-composite}
\leavevmode
\begin{enumerate}[label=(\roman*)]
\item \textbf{Network and process equations.}
Connect a sender and receiver by lossy FIFO data and acknowledgement media $d,a$.
Their process equations, for $b\in\{0,1\}$, are
\[
 \begin{aligned}
 \mathsf{Idle}_b&=\mathsf{input}?(m).\mathsf{Send}_b(m),\\
 \mathsf{Send}_b(m)&=d!(b,m).\mathsf{Wait}_b(m),\\
 \mathsf{Wait}_b(m)&=a?(c).
    \begin{cases}\mathsf{Idle}_{1-b},&c=b,\\
                  \mathsf{Wait}_b(m),&c\ne b\end{cases}
       +\mathsf{retry}.\mathsf{Send}_b(m),\\
 \mathsf{Expect}_b&=d?(c,m).
    \begin{cases}\mathsf{output}!m.a!b.\mathsf{Expect}_{1-b},&c=b,\\
                  a!(1-b).\mathsf{Expect}_b,&c\ne b.
    \end{cases}
 \end{aligned}
\]
The cases index the continuations of the finite input branches.

\item \textbf{Retry and initialization.}
Retry is the selected local action, with no timing constraint.
Start at $\mathsf{Idle}_0,\mathsf{Expect}_0$ with empty media.

\item \textbf{Retained outputs and acknowledgements.}
The emitted-output, acknowledgement and duplicate branches give the residual prefixes formerly denoted $\mathsf{Emit},\mathsf{Ack},\mathsf{Repeat}$; they are derived process positions.

\item \textbf{Operational elaboration.}
Elaborate the terms by the existing witnessed-medium construction and connect the two loss strategies, retaining their audit occurrences:
\begin{equation}\label{eq:abp-strategy}
 \Sigma_{\rm ABP}=\operatorname{Hide}_{d,a,\rm ctrl}
 (\Sigma_{\rm sender}\parallel\Sigma_{\rm loss,d}
  \parallel\Sigma_{\rm receiver}\parallel\Sigma_{\rm loss,a}).
\end{equation}
Its history carrier $S_{\rm ABP}$ is the principal-history realization of this elaboration.
The general history theorem supplies participant chains, packet suppliers, and removal of earlier FIFO packets; its resource squares supply exactly their independent exchanges.

\item \textbf{Ordinary interface.}
At the ordinary interface expose current conditions by shared readiness, including successful acknowledgement before the next application input.
This retains the protocol's actual enabling information.

\item \textbf{Residual graph.}
The residual process graph is also its lossy-channel-system presentation~\cite{abdulla-jonsson-1996}.
\end{enumerate}
\end{construction}

\begin{definition}[Alternating-bit boundary labels]\label{def:abp-boundary-labels}
For epoch $n$, the \textbf{application labels} $i_n^-$ and $o_n^+$ denote its input and output occurrences. Its \textbf{readiness label} $\mathsf{ready}_{n+1}^+$ is the successful matching-acknowledgement transition of the sender back to its next idle state, retained at the readiness boundary. Thus readiness labels name specified transitions of Construction~\ref{con:alternating-bit-composite}; they add no extra transition to the network.
\end{definition}

\begin{theorem}[Alternating-bit safety and finite application behavior]\label{thm:abp-strategy}
Before hiding, the full audited realization of \eqref{eq:abp-strategy} is the protocol kernel of Theorem~\ref{thm:protocol-audit}. At every finite execution prefix, the application words satisfy
\begin{equation}\label{eq:abp-prefix-safety}
 v\leq u,\qquad 0\leq |u|-|v|\leq1.
\end{equation}
After internal saturation, its finite application traces are exactly the capacity-one FIFO traces. Its ordinary readiness interface retains
\[
 i_n^-<o_n^+<\mathsf{ready}_{n+1}^+<i_{n+1}^-.
\]
The exact strategy retains retry, loss, duplicate-reception, and acknowledgement histories and is not identified with the hom kernel of an unrestricted stream arena.
\end{theorem}
\begin{proof}\vmode
\leavevmode
\begin{enumerate}[label=(\roman*)]
\item \textbf{Epoch labels and invariant.}
Label a transmission attempt by the application-input occurrence that produced it.
Number these epochs from zero, with bit $b_n=n\bmod2$.
We prove the following invariant whenever the sender enters $\mathsf{Idle}_{b_n}$: exactly the first $n$ application inputs have been output; the receiver's logical next epoch is $n$; any surviving older data and acknowledgement packets belong only to epoch $n-1$.

\item \textbf{Initial epoch.}
For $n=0$ the channels are empty.
The receiver may be processing a duplicate from the previous epoch, but its duplicate branch returns to the same logical expectation.

\item \textbf{FIFO data attempts.}
After input $n$, all new data attempts have bit $b_n$ and its fixed payload.
They are enqueued behind any surviving attempts from epoch $n-1$.
FIFO loss cannot reverse that order.

\item \textbf{Unique output in an epoch.}
The old attempts take the receiver's duplicate branch and produce only old acknowledgements.
The first surviving epoch-$n$ attempt produces precisely its application output; all further attempts of that epoch take the duplicate branch after the receiver's acknowledgement step.
Thus there is at most one output for this input, with its correct payload.

\item \textbf{Acknowledgement ordering.}
Before the first epoch-$n$ acknowledgement is produced, all surviving older data ahead of that attempt have been processed; all their acknowledgements were enqueued earlier.
The sender discards acknowledgements with the other bit.
By the induction invariant there is no still older acknowledgement with bit $b_n$.

\item \textbf{Advance to the next epoch.}
It can therefore leave its wait state on bit $b_n$ only after the receiver's epoch-$n$ output.
Consuming such an acknowledgement removes, or follows losses of, all earlier acknowledgements in that FIFO.
The remaining data and acknowledgements can only be duplicates of epoch $n$.
This proves the invariant at the next idle state.

\item \textbf{Application-prefix safety.}
During an epoch the sender accepts no additional input and the receiver produces at most its one output, proving \eqref{eq:abp-prefix-safety} at every intermediate prefix.

\item \textbf{Completeness of visible histories.}
Conversely, realize any finite capacity-one FIFO trace by taking one successful data transfer for each requested output and completing the acknowledgement transfer before the next input.
Every step is a displayed rule of Construction~\ref{con:alternating-bit-composite}; a finite trace may stop with one pending input.
This proves equality of the finite application trace languages after internal saturation.

\item \textbf{Ordinary application display.}
In $S_{\mathrm{ABP}}$, retain the application inputs, application outputs, and successful sender acknowledgement receptions, hiding the other events while retaining the visible events' principal-history identities.
Map this structure to the word-labelled stop-and-wait arena with the displayed output-to-ready and ready-to-next-input dependencies.

\item \textbf{Receptivity and courtesy.}
After a ready, the sender is actually idle, so every next input value has its unique lift.
All remaining extra immediate cross-interface causes run from an input to a positive output or ready; positive-to-negative causes and the visible positive order are already in the arena.
Thus the map is receptive and courteous.

\item \textbf{Kernel witnesses.}
Its fibre at $x$ is the set of these visible strategy configurations over $x$, and Theorem~\ref{thm:scott-discrete-fibrations} supplies its reached kernel.

\item \textbf{Retained network histories.}
The operational interaction retains the original local witnesses and FIFO matches, so before hiding its full chart kernel is the direct kernel of the whole plant.
The corresponding matched ordinary interfaces use the reached-coend comparison of Theorem~\ref{thm:ordinarycomp}.
Repeated losses or retries remain distinct audited histories and can continue forever; no eventual-delivery assertion follows from the finite-trace calculation.
\end{enumerate}
\end{proof}

For packets in $D$, the epoch proof of Theorem~\ref{thm:abp-strategy} transports the payload along each retransmission unchanged. The bit tests and acknowledgement matches examine only control data. At the first accepted delivery in epoch $i$ the input and output payload maps from that history consequently agree. Their inverse images of each $U\in\Op(D)$ agree, including under simultaneous observations at finitely many epochs. The proof of the length bound counts the actual pending application input, and the readiness transition still requires the matching acknowledgement. In particular the universal safety statement ranges over every prefix, while a particular observed delivery is a positive event.

\subsection{Protocol composition and retained witnesses}\label{subsec:pc-protocol-equations}
\begin{proposition}[Two bags and the intermediate-word normal form]\label{prop:bag-composition}
The raw composite of two bag kernels is reached and has fibre
\[
 (\Phi_{\mathrm{bag}}\star\Phi_{\mathrm{bag}})(u,w)
 =\int^{v\in C_M}\mathsf{Inj}_M(v,u)\times\mathsf{Inj}_M(w,v).
\]
Writing a representative as $\alpha:[k]\hookrightarrow[n]$ and $\beta:[m]\hookrightarrow[k]$, its unique normal form deletes the intermediate suffix after $\ell=\max\im\beta$, with $\ell=0$ for $m=0$. The map to one bag sends the class to $\alpha\beta$ and need not be injective.
\end{proposition}
\begin{proof}\vmode
\leavevmode
\begin{enumerate}[label=(\roman*)]
\item \textbf{Intermediate-word normalization.}
An arrow of $C_M$ extends an intermediate word by a suffix.
The coend relation adds or removes that suffix precisely when $\beta$ uses only the earlier positions; it restricts or extends $\alpha$ accordingly.
Therefore the prefix through the last used position, together with its two restricted injections, is invariant under each generating relation.

\item \textbf{Canonical representative.}
Every representative reduces to this prefix, proving existence and uniqueness of the normal form.

\item \textbf{Reachability.}
Accept all of $u$, forward its selected intermediate word in order $\alpha$, then deliver the selected positions $\beta$.
This executes every representative, so its class is reached.

\item \textbf{Noninjective bag comparison.}
For $u=ab,w=a$, the intermediate words $a$ and $ba$ give distinct normal forms: respectively $\alpha=(1),\beta=(1)$ and $\alpha=(2,1),\beta=(2)$.
They have the same composite injection.
In the second history the intermediate forwarding of $b$ precedes the forwarding of $a$, so it belongs to the visible output's internal causal past.
It is not an unused suffix.
Hence the comparison to a single bag is not an isomorphism.
\end{enumerate}
\end{proof}

Similarly, two lossy FIFOs distinguish the representative $u=ab,v=b,w=b$ from $u=ab,v=ab,w=b$: deletion occurred before or after the intermediate transfer. Co-Yoneda collapses the FIFO hom kernels in Theorem~\ref{thm:buffercomposition}, while these matching kernels retain their intermediate witnesses. A quotient to possible payload words therefore requires its own compatibility proof before it is transported through interaction.

The following table collects the intrinsic protocol constructions. Theorem~\ref{thm:protocol-audit} gives their common full-history audit, including mixed networks; the individual proofs calculate the indicated public interfaces.

\begingroup\small
\begin{longtable}{@{}p{.23\textwidth}p{.46\textwidth}p{.23\textwidth}@{}}
\caption{Communication protocols and their strategy realizations.}\label{tab:protocol-strategies}\\
\toprule
Protocol & Strategy witnesses and interface & Construction \\
\midrule
\endfirsthead
\toprule
Protocol & Strategy witnesses and interface & Construction \\
\midrule
\endhead
FIFO & Prefix matching; asynchronous copycat & Theorem~\ref{thm:fifocopycat} \\[3pt]
Bag & Injective occurrence matching & Theorem~\ref{thm:matching-kernels} \\[3pt]
Lossy FIFO & Increasing occurrence matching after hiding drops & Theorem~\ref{thm:matching-kernels} \\[3pt]
Duplication & Ordered binary packet genealogies & Theorem~\ref{thm:dup-kernel} \\[3pt]
Exchange & Legal adjacent-swap histories & Proposition~\ref{prop:exchange-kernel} \\[3pt]
Replay & Retained-source commands and buffered attempts & Proposition~\ref{prop:replay-kernel} \\[3pt]
Blocking capacity & Prefix strategy on a capacity arena & Theorem~\ref{thm:bounded-kernel} \\[3pt]
Overflow policies & Capacity-branch causal diagrams & Theorem~\ref{thm:overflow-strategy} \\[3pt]
Credit return & Slot cycles through forward and return channels & Theorem~\ref{thm:credit-kernel} \\[3pt]
Multicast & Common production with recipient children; shared dispatch order for separate FIFO fanout & Theorem~\ref{thm:multicast-kernels} \\[3pt]
Joint communication & Shared readiness and one atomic commit & Proposition~\ref{prop:joint-strategy} \\[3pt]
Cancellation & Named consumption and status conflicts & Proposition~\ref{prop:service-strategy} \\[3pt]
Allocation & Permission and name consumption; created interfaces & Proposition~\ref{prop:allocation-strategy} \\[3pt]
Alternating bit & Controller/channel interaction and stop-and-wait readiness & Theorem~\ref{thm:abp-strategy} \\[3pt]
Mixed networks & Exact chart histories; explicit resource refinements & Theorem~\ref{thm:protocol-audit}, Proposition~\ref{prop:parallel-capacity} \\
\bottomrule
\end{longtable}
\endgroup

\subsection{Loss stages and alternative enabling causes}
\begin{definition}[Stage-labelled loss kernel]\label{def:loss-stage-kernel}
For $r\in\mathbb N$, define $L_r(u,v)$ to be the set of pairs $(\gamma,c)$, where $\gamma\in\mathsf{Inj}_M(v,u)$ is increasing and
\[
 c:\{1,\ldots,\gamma(|v|)\}\setminus\im\gamma\to[r].
\]
For $v=\eps$ this domain is empty. Input extension retains the pair and output restriction restricts $\gamma$ and $c$ to the prefix ending at the last retained supplier. These are the \textbf{stage-labelled loss fibres}; their composition is calculated in Proposition~\ref{prop:pc-loss-stages}.
\end{definition}

\begin{proposition}[Composition counts loss stages]\label{prop:pc-loss-stages}
Let $L_r$ have, over input $u\in M^n$ and output $v\in M^m$, an increasing supplier injection $\gamma:[m]\hookrightarrow[n]$ and a stage map
\[
 c:\{1,\ldots,\gamma(m)\}\setminus\im\gamma\longrightarrow[r].
\]
For $m=0$ the displayed domain is empty. These are the visible causal witnesses of $r$ successive lossy FIFO stages. There are coherent isomorphisms
\[
 L_s\star L_r\cong L_{r+s},\qquad L_0=Y_M.
\]
For a fixed $\gamma$, the fibre has $r^{\gamma(m)-m}$ elements, with the empty-domain convention.
\end{proposition}
\begin{proof}\vmode
\leavevmode
\begin{enumerate}[label=(\roman*)]
\item \textbf{Stage-labelled skipped positions.}
Normalize the intermediate word to the shortest prefix supplying the last visible output, as in the preceding intermediate-word calculation.
A skipped source position disappears either in the first block or in the second; retain its stage in that block and shift second-block stages by $r$.

\item \textbf{Inverse execution construction.}
Conversely, delete the first-block losses to reconstruct the normalized intermediate word, its first matching, and the second matching with its stage labels.

\item \textbf{Normalization and inverse equations.}
Positions beyond the last visible supplier disappear under visible-past normalization on both sides.
The two constructions are inverse.

\item \textbf{Composition and counting.}
Three blocks partition the same stage set into three successive blocks, proving coherence.
With no stages, the skipped set must be empty and $\gamma(j)=j$, giving reliable FIFO.
Counting maps from the skipped set proves the fibre formula.
\end{enumerate}
\end{proof}
For input $abc$ and output $c$, one loss stage has one such witness while two stages have four. The process equation therefore adds stage indices; it is not an idempotent loss equation. Together with Proposition~\ref{prop:bag-composition}, this gives explicit composition laws for the protocol language that retain intermediate communication witnesses.

\subsubsection{Parallel enabling causes in mixed media}\label{subsec:parallel-causes}
The exact cubical kernel remains available when a mixed resource guard has several compatible causes. This can already happen with capacity and loss, despite canonical packet naming.

\begin{proposition}[Capacity released by either of two losses]\label{prop:parallel-capacity}
In a capacity-two FIFO initially containing distinct occurrences $a,b$, let $L_a,L_b$ be their losses and let $E$ enqueue $c$. The histories modulo resource squares do not have a prime-event presentation retaining the same residual occurrence $E$ and the same history identifications. There is an explicit slot-witness refinement with prime causal alternatives and a surjective comparison back to those executions and cubes.
\end{proposition}
\begin{proof}\vmode
\leavevmode
\begin{enumerate}[label=(\roman*)]
\item \textbf{Physical episode.}
Initially $E$ is disabled.
All four executions
\[
 L_aL_bE,\qquad L_aEL_b,\qquad L_bL_aE,\qquad L_bEL_a
\]
are legal.

\item \textbf{Independent squares.}
The losses commute initially.
After either loss, the other loss and $E$ commute: both orders end at the queue containing only $c$, with identical selected instructions and names.
Thus all four have one endpoint history under the square equations.

\item \textbf{Failure of prime intersection closure.}
That history has prefixes $\{L_a,E\}$ and $\{L_b,E\}$, but their intersection $\{E\}$ is not executable.
Within a common prime-event configuration, configurations are ideals and their intersection is a configuration.
This proves the obstruction, relative to any common earlier history that produced the full queue as well.

\item \textbf{Slot-refined enqueue alternatives.}
Give the two occupied slots separate names.
Losing $a$ frees its slot and losing $b$ frees the other.
Retain the slot chosen by enqueue, giving alternatives $E_a$ and $E_b$ with causes $L_a<E_a$ and $L_b<E_b$.
Each alternative uses its actual released resource; all alternatives are retained.

\item \textbf{Forgetting the slot choice.}
Forget the slot choice by $E_a,E_b\mapsto E$.

\item \textbf{Edge and cube lifting.}
Every original edge lifts from a state over its source.
A coenabled enqueue/loss square lifts by using the slot already free before that loss, and the two-loss square lifts unchanged.
These are the relevant nontrivial cells, so the comparison is surjective on executions and charts.

\item \textbf{Refined causal realization.}
The refined resource diagrams have unique selected producers and their principal-history construction gives the causal alternatives.
Their forgetful map recovers the original cubical history identifications without choosing one loss as the cause of every enqueue.
\end{enumerate}
\end{proof}

We now identify the disjunctive causal presentation of this episode and its universal factorization property. The exact cubical history display remains the physical comparison target.

\subsubsection{The capacity episode and its comparison maps}
\begin{definition}[Capacity-episode carriers]\label{def:formal-capacity-episode-carriers}
Restrict Proposition~\ref{prop:parallel-capacity} to the episode consisting of its two specified losses and its one specified enqueue, each used at most once. Abbreviate $L_a,L_b,E$ by $a,b,e$. Its physical \textbf{configuration family} is
\[
 \mathcal F_{\mathrm{cap}}=\{\varnothing,\{a\},\{b\},\{a,b\},
                         \{a,e\},\{b,e\},\{a,b,e\}\},
\]
with equality as its equivalence. Let $P_{\mathrm{cap}}$ have prime events $a,b,e_a,e_b$, order $a<e_a$, $b<e_b$, all finite subsets consistent, and $e_a\equiv e_b$ as its only nontrivial equivalence. Write $\pi(a)=a$, $\pi(b)=b$, and $\pi(e_a)=\pi(e_b)=e$.
\end{definition}

\begin{theorem}[Exact disjunctive presentation of the episode]\label{thm:capacity-edc}
The structure $P_{\mathrm{cap}}$ is an EDC. The map $\pi:I(P_{\mathrm{cap}})\to\mathcal F_{\mathrm{cap}}$ is surjective on configurations and has the universal factorization property of Proposition~\ref{prop:edc-extremal-universal}. Its unambiguous configurations form the ordinary prime event structure $P_{\mathrm{slot}}$ obtained by making $e_a,e_b$ conflicting. The induced map
\[
 K(P_{\mathrm{slot}})\longrightarrow X_{\mathrm{cap}}
\]
to the physical episode's chart automaton is surjective on vertices, edges, and squares, and on finite executions. It is the named-slot refinement of Proposition~\ref{prop:parallel-capacity}.
\end{theorem}
\begin{proof}\vmode
\leavevmode
\begin{enumerate}[label=(\roman*)]
\item \textbf{Causal presentation and equivalence.}
The principal pasts are $\{a\}$, $\{b\}$, $\{a,e_a\}$, and $\{b,e_b\}$. Each is unambiguous.
Every image of a configuration containing $e$ also contains $a$ or $b$, and the seven displayed physical configurations all occur as images.
Reflection of equivalence holds because only $e_a,e_b$ share a physical image.

\item \textbf{Definition of a factor.}
For a map $f:I(Q)\to\mathcal F_{\mathrm{cap}}$ from an EDC, define $\bar f$ on exactly the domain of $f$ by the finite case formula
\[
 \bar f(q)=\begin{cases}
 a,&f(q)=a,\\
 b,&f(q)=b,\\
 e_a,&f(q)=e\text{ and }a\in f([q]),\\
 e_b,&f(q)=e\text{ and }a\notin f([q]).
 \end{cases}
\]
The image $f([q])$ belongs to the seven displayed configurations.
Membership of the named event $a$ is decidable on this explicitly finite configuration object, so the formula is finite case analysis, not an application of a choice principle or a decision of an arbitrary predicate.
The partial domain itself is left unchanged; there is no test deciding whether $f$ is defined.

\item \textbf{Configuration preservation.}
If $\bar f(q)=e_a$, there is $p\le q$ with $f(p)=a$, hence $\bar f(p)=a$.
If $\bar f(q)=e_b$, the configuration $f([q])$ contains $e$ and does not contain $a$; the displayed list then forces $b\in f([q])$, giving the other required predecessor.
It follows that every configuration has a downward-closed image, and consistency is automatic.

\item \textbf{Equivalence preservation and reflection.}
Equivalent source events have equivalent images, with the same saturated partial domain, because $f$ preserves equivalence and $e_a\equiv e_b$.
If two images within a configuration are equivalent, their $\pi$-images agree; reflection of equivalence for $f$ gives equivalence of their source events.

\item \textbf{Universal factorization.}
Thus $\bar f$ is an allowed map and $\pi\bar f=f$.
Every other factor has the same domain and pointwise equivalent values, since equality of physical images is exactly equivalence in $P_{\mathrm{cap}}$.
This proves the universal property directly, without Proposition~\ref{prop:edc-extremal-universal}.

\item \textbf{Unambiguous configurations.}
Both causal alternatives remain in the object.
An unambiguous configuration cannot contain both $e_a,e_b$.
Declaring that pair conflicting gives exactly those configurations, with the same causal order.

\item \textbf{Physical edges and squares.}
Their map to physical configurations preserves one-event extensions.
The physical squares are the initial $a/b$ square, the $b/e$ square after $a$, and the $a/e$ square after $b$.
They lift respectively unchanged, using $e_a$, and using $e_b$.
There is no initially enabled three-dimensional chart.

\item \textbf{Execution surjectivity.}
Every physical edge lifts from every source configuration above its source: before $e$ choose an already freed slot; losses retain their names.
Induction gives execution surjectivity.

\item \textbf{Terminal alternatives.}
The two refined terminal configurations with a single chosen enqueue cause have the same physical image, so the comparison is not injective.
\end{enumerate}
\end{proof}

The capacity example separates three carriers: $P_{\mathrm{cap}}$ has nine configurations, of which eight are unambiguous, and there are seven physical configurations. The ninth records both compatible causes of one physical enqueue, not two enqueue events.

Accordingly, $I(P_{\mathrm{cap}})\to\mathcal F_{\mathrm{cap}}$ is a map of equivalence families, whereas the cubical comparison uses the unambiguous ordinary structure. Their different types explain why deleting the equivalence relation does not turn a causal unfolding into the physical cubical carrier.

\subsection{Concrete medium process rules}
\begin{definition}[Rendezvous and FIFO computation]\label{def:rendezvous-fifo-computation}
For $a:\mathsf{Chan}(\tau)$ and $v:\tau$, the selected local and rendezvous witnesses give
\[
N_b(\tau.P\mid R)\longrightarrow N_b(P\mid R),\qquad
N_b(a!v.P\mid a?(x).Q\mid R)
 \xrightarrow{d}_{\rm rv}N_b(P\mid Q[v/x]\mid R).
\]
The rendezvous uses distinct active participants and leaves the medium store unchanged. For an occurrence-valued FIFO word $w$, commitment instead gives
\begin{align*}
N_{B[a:=w]}(a!v.P\mid R)&\xrightarrow{d_s}
 N_{B[a:=w(\kappa,v)]}(P\mid R),\\
N_{B[a:=(\kappa,v)w]}(a?(x).Q\mid R)&\xrightarrow{d_r}
 N_{B[a:=w]}(Q[v/x]\mid R).
\end{align*}
The send constructs $\kappa$ from its producer and output position; the receive retains it as supplier. Capacity $k>0$ additionally tests $|w|<k$ at append. Scope includes the store and received names. These rules instantiate the single commitment rule; the remaining media change its witness carrier and store update, not input substitution.
\end{definition}

The following table gives the operational reading of the strategy families developed in this section. Each preceding semantic construction calculates its actual interface and witness fibre.
\begin{longtable}{@{}p{.22\textwidth}p{.71\textwidth}@{}}
\toprule Discipline & Operation on the medium and selected offers\\\midrule\endhead
Rendezvous & Select a matching send and receive at distinct participants and advance both in one commitment.\\
FIFO & Append a sent packet; receive and consume the oldest packet with its actual supplying occurrence.\\
Bag and loss & A bag receive selects one pending packet; loss selects an explicit deletion witness. Increasing suppliers characterize visible lossy FIFO.\\
Capacity and credit & Capacity tests the actual residual word; credit tests an owned slot and its return history. Overflow selects its declared full-state branch.\\
Replay and duplication & Replay reads a retained source; duplication creates named descendants with their genealogy.\\
Multicast and joint communication & Retain child recipients or select the required offers in one joint operation, with the specified shared readiness.\\
Cancellation and allocation & Select and consume the live occurrence or create a fresh allocated occurrence; update the associated status.\\
Alternating-bit composition & Advance the explicit sender, receiver, retry and acknowledgement controls with the selected packet witnesses.\\
\bottomrule
\end{longtable}

For the rules of Definition~\ref{def:rendezvous-fifo-computation}, an already enabled receive of the old head and a send at a distinct participant leave the same residual word and continuations in either order. Lemma~\ref{lem:fifo-square} proves this square with its supplier retained. At an empty queue the new send instead supplies the receive, giving a dependency.

\begin{example}[A fresh channel sent by rendezvous]\label{ex:pc-name-passing}
With $b$ fresh for the receiver's surrounding context, scope extrusion and one rendezvous record derive
\[
 (\mathsf{new}\ b)(a!b.P)\mid a?(x).Q
 \longrightarrow (\mathsf{new}\ b)(P\mid Q[b/x]).
\]
The FIFO instance first stores $b$ within its scope, advances the sender, and later supplies that same name to the receiver. The intermediate queued name is part of the scoped store. This is a change of medium rules for the same offered actions.
\end{example}

\paragraph{Process realizations of the media.}
The preceding discipline table specializes the single commitment rule. FIFO copycat has the waiting edges $i_e^-<o_e^+$ on its stream interface; supplier injections and genealogies are the witnesses of the bag/loss and duplication instances. Shared readiness gives the ordinary rendezvous interface. These are the already constructed process displays and history comparisons. A sequential implementation that adds order between independent commitments requires its own comparison to their cells.

\paragraph{Asynchronous messages as pending occurrences.}\label{par:pc-asynchronous}
For the ordinary asynchronous message sublanguage, outputs have no continuation. Reify each active output $a!v$ as one bag packet. A selected-message input consumes exactly the corresponding packet and substitutes the same $v$. This bijection preserves receive commitments and the residual message family. It starts with the message already present; adding a separate enqueue step changes the step comparison.

Replay retains its persistent source, capacity retains the selected full-state branch, and credit consumes its actual owned slot. Joint communication retains its participating offer family; cancellation competes with delivery for the selected live occurrence. The alternating-bit term composes these media with its sender and receiver. Each instance's existing witness fibre and chart construction determine these dependencies; none duplicates an arbitrary capability or supplies a fairness guarantee.

\paragraph{The buffered-$\pi$ instance.}
\begin{definition}[Buffered channel instance]\label{def:formal-buffered-channel-instance}
For the buffered-$\pi$ instance of \cite[Section~2]{buffered-pi}, retain a store $B(a)=(k,w)$ for each buffered subject, with positive capacity $k$. The two internal rules append when $|w|<k$ and remove the head when $w=vw'$. Unbuffered subjects use the rendezvous rule. Buffer creation is a separate allocation transition. Public buffer endpoints have the corresponding environment input/output transitions. Local-name markers, restriction, and bound-output rules are transported with the store; dropping the scope information would not give this instance.
\end{definition}

\begin{proposition}[Buffered operational comparison]\label{prop:pc-buffered-comparison}
The translation sending a buffered process and its store to its named-offer elaboration with the same store preserves and reflects the structural labelled transition relation, including buffer creation and scoped name transport. At fixed selected prefix and medium witnesses it retains the commitment record.
\end{proposition}
\begin{proof}\vmode
\leavevmode
\begin{enumerate}[label=(\roman*)]
\item \textbf{Communication-rule correspondence.}
An unbuffered rule selects precisely the matching pair of offers.
Buffered output selects the unique append operation for its value and available capacity; buffered input selects the actual head-removal operation and substitutes that head in its continuation.

\item \textbf{Buffers and creation.}
The environment buffer rules select the same store boundary operation, with no local process premise.
Creation allocates the same fresh buffered name and empty word.

\item \textbf{Structural cases.}
Sum, parallel, restriction and structural-congruence cases are those of Proposition~\ref{prop:pc-named-elaboration}, using the target's local-name markers and open/close conventions.

\item \textbf{Converse and erasure.}
Conversely each selected operation is one of these cases and recovers its source rule.
Erasing occurrence addresses gives the ordinary transition relation.
\end{enumerate}
\end{proof}

The positive-capacity parameter in this comparison is distinct from an unbuffered channel. Setting the blocking-store capacity to zero disables its send rule; it does not derive rendezvous.


\subsection{Disjunctive causes and canonical causal interfaces}\label{subsec:disjunctive-causal-interfaces}
Some media admit compatible alternative causes for the same action. A single prime-history description cannot identify those alternatives without losing information. We therefore retain their causal equivalence and construct the comparison to a canonical causal interface.

Resource enabling determines the type of comparison required by the process language. A physical commitment may have compatible alternative enabling histories. The following general constructions retain those alternatives and distinguish the exact physical audit from its ordinary causal presentation; a nondeterministic prefix notation alone does not supply that comparison.

\subsubsection{Causal equivalence and universal unfolding}
\begin{definition}[Disjunctive-cause event structures]\label{def:formal-disjunctive-cause-event-structures}
An \textbf{event structure with equivalence} is a prime event structure $(P,\leq,\mathsf{Con})$ with an equivalence $\equiv$ on events. Its prime events represent causal realizations; their equivalence classes represent physical events. An \textbf{event structure with disjunctive causes} (an EDC) is such a structure in which every principal past $[p]$ is unambiguous:
\begin{equation}\label{eq:edc-unambiguous}
 p_1,p_2\leq p\text{ and }p_1\equiv p_2\quad\Longrightarrow\quad p_1=p_2.
\end{equation}
An arbitrary configuration may still contain distinct equivalent prime causes. Equivalence is different from a permutation symmetry of fresh names.
\end{definition}

\begin{definition}[Equivalence families and their maps]\label{def:formal-equivalence-families-and-their-maps}
Write $\mathsf{EqFam}$ for configuration families with equivalence, avoiding confusion with the coproduct completion $\operatorname{Fam}(B)$. A subfamily $\mathcal X\subseteq\mathcal F$ is compatible when some $z\in\mathcal F$ contains all its members, and \textbf{finitely compatible} when every finite subfamily has such a bound. A \textbf{configuration family} is a nonempty family closed under unions of finitely compatible subfamilies, and each $e\in x\in\mathcal F$ has a finite securing chain $e_1,\ldots,e_n=e$ inside $x$, with every prefix set in $\mathcal F$. Maps are partial functions preserving equivalence and configurations and reflecting equivalence within each configuration; preservation of equivalence includes equidefinedness on equivalence classes. \textbf{Equivalent maps} have the same domain and pointwise equivalent values. The functor $I$ sends an \textbf{event structure with equivalence} to its family of all configurations. These are the categories of \cite[Sections~2.2 and~3]{devisme-winskel}.
\end{definition}

\begin{definition}[EDC maps and realization codes]\label{def:edc-map-realization-types}
\leavevmode
\begin{enumerate}[label=(\roman*)]
\item The category $\mathsf{EDC}$ has the EDCs defined above as objects. Its arrows $f:P\rightharpoonup Q$ send configurations to configurations, have equivalence-saturated domains, preserve equivalence, and satisfy $f(p)\equiv f(p')\Rightarrow p\equiv p'$ whenever $p,p'\in x\cap\dom f$ for a configuration $x$. Composition is partial-function composition. The configuration functor has type $I:\mathsf{EDC}\to\mathsf{EqFam}$, $I(P)=(\Conf(P),\equiv_P)$.
\item A \textbf{causal realization} of $\mathcal F$ is a poset $R$ with finite principal pasts and a function $\rho:R\to|\mathcal F|$ sending each downset to a member of $\mathcal F$. A \textbf{total coarsening} is a surjection $f:R\to R'$ carrying downsets to downsets with $\rho=\rho'f$. A realization is \textbf{extremal} when every such coarsening is a label-preserving order isomorphism, and \textbf{prime extremal} when it also has a greatest element. This fixes the convention of~\cite[Section~4]{devisme-winskel}.
\item The \textbf{prime-extremal unfolding} $\operatorname{er}(\mathcal F)$ has the finite prime-extremal realization codes modulo label-preserving order isomorphism, ordered by principal restriction. A finite family is consistent when the union of its realization images belongs to $\mathcal F$; equivalence compares the labels of the tops. Its counit returns that top label. For an event structure with equivalence $Q$, $KQ$ restricts to the events with unambiguous principal past. The subsequent finite-code argument and the cited unfolding construction verify these operations.
\item $\mathsf{AC}_{\mathrm{fin}}$ names the assertion that $\prod_{i\in I}A_i\ne\varnothing$ for every set-indexed family of nonempty finite sets. It is the assertion analysed by Proposition~\ref{prop:edc-choice}, not an added ambient assumption.
\end{enumerate}
\end{definition}

\begin{definition}[Finite-configuration completion]\label{def:formal-finite-configuration-completion}
For a rooted family initially given only with bounded compatible-union \textbf{closure} and finite securing, the needed completeness is constructed, rather than supplied, by
\begin{equation}\label{eq:family-completion}
 \widehat{\mathcal F}=\left\{x\subseteq\bigcup\mathcal F:
 \forall u\subseteq_{\mathrm{fin}}x\ \exists v\in\mathcal F\;
 (v\text{ finite and }u\subseteq v\subseteq x)\right\}.
\end{equation}
Proposition~\ref{prop:family-completion} proves its universal property and equality of its finite configurations with those of $\mathcal F$. The unfolding counit for this input targets $\widehat{\mathcal F}$; the completion can add infinite configurations. It does not alter the finite capacity episode in Appendix~\ref{app:communication-examples}.
\end{definition}

Section~\ref{subsec:general-causal-translation} constructs a choice-free universal causal presentation by retaining all finite unambiguous realizations. It also compares this presentation with the prime-extremal unfolding of~\cite{devisme-winskel}: actual factorization into that smaller target has precisely finite-choice strength, while a span retains all its finite local factors without a global selection. The capacity construction in Appendix~\ref{app:communication-examples} gives its smaller causal carrier and its factorization directly.

\subsubsection{Canonical causal interfaces}
The event-structure causal-interface construction also has a minimality property and extends to the disjunctive presentation.
\begin{proposition}[Least causal interface]\label{prop:least-causal-interface}
For a polarized prime event structure $S$, the arena $A_S$ of Lemma~\ref{lem:causal-interface} has the least causal order, among orders contained in that of $S$ on the same events and consistency, for which the identity event map is a receptive courteous strategy. For a polarized EDC $S$, retaining its equivalence and applying the same deletion of immediate negative-to-positive dependencies constructs a polarized EDC $A_S$ and an EDC strategy $S\to A_S$, with the analogous least-order property.
\end{proposition}
\begin{proof}\vmode
\leavevmode
\begin{enumerate}[label=(\roman*)]
\item \textbf{Least forced arena order.}
Courtesy forces every immediate dependency of $S$ other than a negative-to-positive one to occur in the arena.
It therefore forces their transitive closure, exactly the order of $A_S$.

\item \textbf{Ordinary strategy.}
Lemma~\ref{lem:causal-interface} proves that this least order already gives a strategy.

\item \textbf{Unambiguous principal pasts.}
For an EDC, the new principal pasts are subsets of the original unambiguous pasts, so \eqref{eq:edc-unambiguous} holds; equivalence still preserves polarity.
Consider the ordinary structures of unambiguous configurations.

\item \textbf{Receptivity.}
If a negative event $n$ extends an unambiguous $S$-configuration $x$ in $A_S$, every immediate $S$-predecessor of $n$ belongs to $x$, because every edge into $n$ was retained.
Hence $x\cup\{n\}$ is $S$-downward closed; it is consistent for the unchanged underlying consistency relation and unambiguous as an event set.
This proves receptivity directly.

\item \textbf{Courtesy and consistency scope.}
The immediate dependencies of the ordinary structures are the original prime dependencies, so the deletion rule proves courtesy.
This argument does not identify the two ordinary consistency relations: deleting order can create additional unambiguous configurations of the arena.

\item \textbf{Equivalence and compatibility.}
The EDC identity map preserves and reflects equivalence.
If two one-event extensions of an $S$-configuration have compatible arena images, their union is $S$-downward closed and consistent for the unchanged underlying consistency relation, hence is a common $S$-extension.
In particular this gives the compatibility condition when one extension is negative.

\item \textbf{EDC strategy and leastness.}
Together with ordinary receptivity and courtesy, these are the strategy conditions characterized by copycat invariance in \cite[Theorem~4]{parallel-causes}.
They prove the EDC assertion.
Applying courtesy to an immediate dependency, which lies in an unambiguous principal past, again proves leastness.
\end{enumerate}
\end{proof}

The interface is canonical relative to its causal domain, polarity and physical-event equivalence. The public communication interfaces retain their specified ports and hiding maps. The following results locate the comparisons between these presentations.

Proposition~\ref{prop:edc-universal} constructs the general universal causal presentation without choice. Proposition~\ref{prop:edc-extremal-universal} treats the distinct prime-extremal target, and Proposition~\ref{prop:edc-choice} characterizes its finite-choice content. The capacity instance has the independent finite proof of Theorem~\ref{thm:capacity-edc}. Theorem~\ref{thm:history-universal} instead constructs the exact rooted display of the given cubical execution geometry. The capacity calculation compares these carriers explicitly. EDC interaction uses the equivalence-sensitive pullback and hiding of~\cite{parallel-causes}; the ordinary reached-coend theorem applies to the ordinary strategy presentations constructed here.

\subsection{General causal unfolding and its foundational scope}\label{subsec:general-causal-translation}
The general unfolding construction must distinguish retaining all realizations from selecting representatives. The following results construct the former, identify the choice content of extremal selection, and supply local comparisons that retain their witnesses.

Use the configuration families and event structures with equivalence defined in Section~\ref{subsec:disjunctive-causal-interfaces}. We first construct the general universal target from all finite unambiguous realizations, then compare it with the smaller prime-extremal target. The foundational comparison is over classical ZF; no choice principle is used in the all-realizations construction.

\begin{proposition}[Finite-configuration completion]\label{prop:family-completion}
Let $\mathcal F$ contain $\varnothing$, be closed under unions of compatible subfamilies, and have the finite securing chains specified above. Formula~\eqref{eq:family-completion}, with the same event equivalence, constructs a configuration family. The inclusion $\mathcal F\subseteq\widehat{\mathcal F}$ is the identity on events and preserves exactly the finite configurations. Every configuration-preserving partial map from $\mathcal F$ to a configuration family $\mathcal G\in\mathsf{EqFam}$ extends uniquely, by the same partial function, to $\widehat{\mathcal F}$. The construction is therefore an idempotent reflection from the bounded-union families just described to $\mathsf{EqFam}$. Finite labelled realizations and prime extremal histories are unchanged by completion.
\end{proposition}
\begin{proof}\vmode
\leavevmode
\begin{enumerate}[label=(\roman*)]
\item \textbf{Finite securing inside original configurations.}
If $x\in\mathcal F$ and $u\subseteq x$ is finite, take a finite securing configuration in $x$ for each member of $u$.
Their union belongs to $\mathcal F$, since they have the common bound $x$, and witnesses $x\in\widehat{\mathcal F}$.
These are finitely many existential witnesses, so no choice axiom is used.

\item \textbf{Equality of finite configurations.}
If $x\in\widehat{\mathcal F}$ is finite, applying its defining condition to $u=x$ gives $x\in\mathcal F$.
Thus the finite configurations coincide, and the underlying event sets coincide.

\item \textbf{Finite securing after completion.}
Every event in a completed configuration has a finite original securing configuration inside it, so securing is preserved.

\item \textbf{Finitely compatible union.}
Let $\mathcal X\subseteq\widehat{\mathcal F}$ be finitely compatible and put $x=\bigcup\mathcal X$.
A finite $u\subseteq x$ is covered by finitely many $x_j\in\mathcal X$.
Within these take finite original configurations $v_j\subseteq x_j$ covering $u$.

\item \textbf{Original finite upper witness.}
There is a common bound $z\in\widehat{\mathcal F}$ for the $x_j$.
Since $\bigcup_jv_j$ is finite, the defining condition for $z$ gives an original configuration containing this union.
The $v_j$ are therefore compatible in $\mathcal F$, and $v=\bigcup_jv_j$ belongs to $\mathcal F$.

\item \textbf{Completeness.}
As $u\subseteq v\subseteq x$, the defining condition holds for $x$.
The empty subfamily is handled by $\varnothing\in\mathcal F$.

\item \textbf{Finite approximation family.}
For $x\in\widehat{\mathcal F}$, let $\mathcal V_x$ be all finite members of $\mathcal F$ contained in $x$.
They cover $x$ and are finitely compatible: any finite union of them has a finite original bound inside $x$.

\item \textbf{Extension of maps.}
For a map $f:\mathcal F\to\mathcal G$, the family $\{f[v]:v\in\mathcal V_x\}$ is consequently finitely compatible in $\mathcal G$.
Its union is $f[x]$, so $f[x]\in\mathcal G$.

\item \textbf{Equivalence and uniqueness.}
Reflection of equivalence within $x$ is checked in an original finite subconfiguration containing the two events concerned.
Preservation of equivalence and domains is unchanged.
This proves the unique extension.

\item \textbf{Least completion.}
Taking $f$ to be an inclusion proves leastness of the completion, and applying the result to a complete family proves idempotence and the reflection.

\item \textbf{Preservation of finite extremality.}
A finite realization tests only finite images of its downsets; these images belong to the same finite families before and after completion.
Its finite total coarsenings likewise have finite carriers.
Hence its realization and extremality predicates are unchanged.
\end{enumerate}
\end{proof}

Completion is needed because bounded-union closure and finite securing do not ensure preservation of infinite configurations. For example, $\Pf(\mathbb N)$ has those properties, but its singleton prime histories form a discrete event structure whose configurations include $\mathbb N$. A proposed counit into $\Pf(\mathbb N)$ would fail on that configuration.

The completion is $\Pow(\mathbb N)$, and the constructed counit targets it. The proved comparison therefore accounts for the difference between compatible and finitely compatible unions.

\begin{definition}[All finite unambiguous causal realizations]\label{def:edc-all-realizations}
For $\mathcal F\in\mathsf{EqFam}$ on the event set $A$, form $\mathsf{Caus}(\mathcal F)$ as follows.
\leavevmode
\begin{enumerate}
\item An event is a triple $e=(X_e,\preceq_e,a_e)$, where $X_e\subseteq A$ is finite and unambiguous, $\preceq_e$ is a partial order with greatest element $a_e$, and every $\preceq_e$-downset belongs to $\mathcal F$.
\item For $b\in X_e$, put $e|b=(\mathord\downarrow_e b,\preceq_e|_{\mathord\downarrow_e b},b)$. Define $d\leq e$ when $d=e|b$ for some $b\in X_e$, and $d\equiv e$ when $a_d\equiv_A a_e$.
\item A finite set $C$ of these events is consistent when $\bigcup_{e\in C}X_e\in\mathcal F$. Define $\epsilon_{\mathcal F}(e)=a_e$.
\end{enumerate}
All the triples are formed by separation over finite subsets of $A$ and their binary relations. No representatives, minimal causes or global family of local factors are selected.
\end{definition}

\begin{proposition}[Choice-free universal causal presentation]\label{prop:edc-universal}
The preceding construction is an EDC and $\epsilon_{\mathcal F}:I(\mathsf{Caus}(\mathcal F))\to\mathcal F$ is a map. Every $f:I(P)\to\mathcal F$ from an EDC has an explicitly constructed factor $\widehat f:P\to\mathsf{Caus}(\mathcal F)$ with the same partial domain and $\epsilon_{\mathcal F}I(\widehat f)=f$. Any factor is equivalent to $\widehat f$. Consequently
\[
 \mathsf{EDC}(P,\mathsf{Caus}(\mathcal F))/\!\equiv
 \ \cong\ \mathsf{EqFam}(I(P),\mathcal F)/\!\equiv
\]
naturally in both variables. This gives an adjunction on the categories obtained by quotienting maps by equivalence, in ZF without choice.
\end{proposition}
\begin{proof}\vmode
\leavevmode
\begin{enumerate}[label=(\roman*)]
\item \textbf{Restriction order.}
Restrictions compose: $(e|b)|c=e|c$ whenever $c\preceq_e b$.
This proves transitivity and reflexivity of the order.
Mutual restriction forces equality of carriers and of the restricted orders and tops, proving antisymmetry.

\item \textbf{Finite pasts and consistency.}
The principal past is $[e]=\{e|b:b\in X_e\}$, hence finite and unambiguous.
Each $X_e$ belongs to $\mathcal F$, so singletons are consistent.

\item \textbf{Event-structure axioms.}
If $D\subseteq C$ and $C$ is consistent, the family $(X_d)_{d\in D}$ has the common bound $\bigcup_{e\in C}X_e$; compatible-union closure makes $D$ consistent.
Adjoining a predecessor changes no carrier union.
This proves all event-structure axioms, including empty consistency.

\item \textbf{Counit on configurations.}
If $x$ is a configuration, downward closure gives
\[
 \epsilon_{\mathcal F}[x]=\bigcup_{e\in x}X_e.
\]
Every finite subfamily of the carrier configurations has its union in $\mathcal F$, by consistency.
Completeness of $\mathcal F$ therefore gives $\epsilon_{\mathcal F}[x]\in\mathcal F$, also for infinite $x$.

\item \textbf{Counit equivalence conditions.}
Equivalence preservation and reflection follow from the definition by top labels.
Thus the counit is a map.

\item \textbf{Finite source history.}
For $p\in\dom f$, put $D_p=[p]\cap\dom f$ and $X_p=f[D_p]$.
Unambiguity of $[p]$ and reflection of equivalence by $f$ make $f|_{D_p}$ injective.

\item \textbf{Transported realization.}
Transport the order of $D_p$ along this bijection:
\[
 f(q)\preceq_p f(r)\quad\Longleftrightarrow\quad q\leq r
 \qquad(q,r\in D_p),\qquad
 \widehat f(p)=(X_p,\preceq_p,f(p)).
\]
The carrier is finite and unambiguous and has the displayed top.
If $Z$ is a downset in $X_p$, let $Y$ be the $P$-downclosure of its inverse image in $D_p$.
Then $Y\subseteq[p]$ is a configuration and $f[Y]=Z$: any defined predecessor maps back into $Z$ by its downward closure.
Hence $Z\in\mathcal F$, proving that the displayed triple is an event.

\item \textbf{Compatibility of restrictions.}
For defined $q\leq p$ one has the restriction equation
\[
 \widehat f(q)=\widehat f(p)|_{f(q)}.
\]
It proves downward closure of $\widehat f[x]$ for any configuration $x$.

\item \textbf{Configuration preservation.}
A finite union of its carriers is a union of the configurations $f[[p]]$ with common bound $f[x]$, so is in $\mathcal F$.
Thus the image is consistent.

\item \textbf{Factor-map conditions.}
Equidefinedness and equivalence preservation and reflection are inherited from $f$, since output equivalence compares exactly its top labels.
This proves that $\widehat f$ is a map and that its triangle commutes strictly.

\item \textbf{Uniqueness up to equivalence.}
If $g$ is any factor, equality of its composite with $f$ gives the same domain and top labels, whence $g\equiv\widehat f$.
The argument also applies when that triangle commutes up to equivalence.

\item \textbf{Functor on map classes.}
Postcomposition with $\epsilon_{\mathcal F}$ therefore gives the stated bijection of map classes.
For a map class $[u]:\mathcal F\to\mathcal G$, define
\[
 \mathsf{Caus}[u]=[\widehat{u\epsilon_{\mathcal F}}].
\]
Equivalent inputs give equivalent lifts.

\item \textbf{Functoriality and adjunction.}
The identity and composite equations follow by comparing their domains and top labels, or by uniqueness of factor classes.
The same calculation with precomposition proves naturality in $P$, and with postcomposition proves naturality in $\mathcal F$.
These are the functor and adjunction equations; strict naturality of the selected representatives is not needed.
\[
\begin{tikzcd}[column sep=large]
I(P)\arrow[r,"I(\widehat f)"]\arrow[dr,"f"']
 &I(\mathsf{Caus}(\mathcal F))\arrow[d,"\epsilon_{\mathcal F}"]\\
 &\mathcal F
\end{tikzcd}
\]
\end{enumerate}
\end{proof}

The construction retains nonminimal dependencies. In particular it is not the prime-extremal unfolding of~\cite{devisme-winskel}. Keeping those additional causal events is what permits the explicit factor above; the following result records the distinct foundational strength of requiring an extremal target.

\begin{proposition}[The prime-extremal factorization assertion]\label{prop:edc-extremal-universal}
For $\mathcal F\in\mathsf{EqFam}$, let $\operatorname{er}(\mathcal F)$ be the causal unfolding by prime extremal realizations of~\cite{devisme-winskel}. Let $KQ$ restrict an event structure with equivalence $Q$ to the events whose principal pasts are unambiguous. This constructs an EDC and a counit
\[
 \epsilon:I(K\operatorname{er}(\mathcal F))\longrightarrow\mathcal F
\]
without choice. The general assertion that every map $f:I(P)\to\mathcal F$ from an EDC factors through this counit is equivalent over ZF to $\mathsf{AC}_{\mathrm{fin}}$, as proved in Proposition~\ref{prop:edc-choice}. Factors, whenever they exist, are unique up to equivalence. Thus the extremal universal-factorization assertion includes a natural bijection of map classes, but it is distinct from the choice-free universal property of $\mathsf{Caus}(\mathcal F)$.
\end{proposition}
\begin{proof}\vmode
\leavevmode
\begin{enumerate}[label=(\roman*)]
\item \textbf{Prime-extremal counit.}
Remark~\ref{rem:unfolding-foundations} constructs the finite-code carrier of the extremal unfolding without representative selection.
Its order, consistency and top-label counit are the realization constructions of~\cite[Proposition~5.1]{devisme-winskel}.
Configuration preservation of the counit uses completeness of $\mathcal F$, by the same compatible-union argument as in Proposition~\ref{prop:edc-universal}; no factor map is needed to form it.

\item \textbf{Unambiguous coreflection.}
The good events defining $KQ$ are downward closed, so restriction gives an event structure satisfying~\eqref{eq:edc-unambiguous}.

\item \textbf{Factorization through the core.}
Given an actual map $h:P\to Q$ from an EDC, its image $h[p]$ is a configuration containing $h(p)$ whenever that value is defined, and therefore contains $[h(p)]$.
Equivalent events in this image have equivalent preimages in $[p]$, which agree by~\eqref{eq:edc-unambiguous}.
Hence $[h(p)]$ is unambiguous and $h$ restricts uniquely through $KQ$.
This proves the coreflection without selecting a factor into $Q$.

\item \textbf{Extremal comparison.}
Equivalence in the extremal unfolding is equivalence of top labels, so any two factors through the counit have the same partial domain and pointwise equivalent values.
Theorem~5.2 of~\cite{devisme-winskel} supplies general factors using choice; Lemma~\ref{lem:edc-finite-extension} and Proposition~\ref{prop:edc-choice} prove the precise equivalence with finite choice for EDC sources.
They concern this extremal target.

\item \textbf{Scope of the target.}
The general unfolding need not itself satisfy~\eqref{eq:edc-unambiguous}, as explained in~\cite[Section~7.1]{devisme-winskel}.
\end{enumerate}
\end{proof}

\begin{remark}[Finite codes and the scope of the unfolding theorem]\label{rem:unfolding-foundations}
The carrier of the prime-extremal unfolding can be formed without selecting representatives of isomorphism classes. Every prime history has a top and hence finite carrier. Form the set of finite codes $([n],\leq,\rho:[n]\to|\mathcal F|)$ satisfying the realization and prime-extremality predicates, with their tops, and quotient it by label-preserving order isomorphism. Top labels, the relation of principal restriction, and the consistency and equivalence predicates are invariant under these isomorphisms, so they descend to this quotient. This removes representative selection from the object construction.

The general existence of a factor map in Proposition~\ref{prop:edc-extremal-universal} is the external theorem of \cite[Theorem~5.2]{devisme-winskel}. Its proof explicitly uses choice, and its Lemma~4.1 uses Zorn's lemma for arbitrary extremal realizations. For EDC sources the following finite extension calculation sharpens the requirement to precisely $\mathsf{AC}_{\mathrm{fin}}$, choice for arbitrary families of nonempty finite sets. The foundation analysis here is over classical ZF; absence of a choice axiom does not assert an intuitionistic decision procedure for arbitrary configuration-membership predicates. The capacity factorization of Theorem~\ref{thm:capacity-edc} has its own finite-case proof. Proposition~\ref{prop:edc-local-cover} separately retains all finite local lifts without a global selection.
\end{remark}

\begin{lemma}[Finite extension of a causal factor]\label{lem:edc-finite-extension}
Put $E=K\operatorname{er}(\mathcal F)$ for $\mathcal F\in\mathsf{EqFam}$ and let $f:I(P)\to\mathcal F$ with $P$ an EDC. If $x$ is a finite configuration of $P$, $p$ is enabled at $x$, and $\lambda:x\rightharpoonup E$ is a factor of $f|_x$ through $\epsilon$, then $\lambda$ extends to a factor on $x\cup\{p\}$. This assertion is provable in ZF without choice. In particular every finite configuration admits a factor.
\end{lemma}
\begin{proof}\vmode
\leavevmode
\begin{enumerate}[label=(\roman*)]
\item \textbf{Undefined case and predecessor image.}
If $f(p)$ is undefined, extend the domain convention by leaving $p$ undefined.
Otherwise put $R=\lambda[[p)]$.
This is a finite configuration of $E$, and consider the nonempty finite collection
\[
 \mathcal D(R,f(p))=
 \{D\subseteq R:D\text{ is down-closed and }
       \epsilon[D]\cup\{f(p)\}\in\mathcal F\}.
\]
The collection contains $R$, since $f$ preserves $[p]$.

\item \textbf{Minimal predecessor selection.}
Take an inclusion-minimal member $D$ and adjoin a new top $t$ above all of $D$, labelled $f(p)$.

\item \textbf{Realization and extremality.}
Downsets excluding $t$ are realized by $D$, and the whole labelled order is a configuration by the defining condition on $D$.
The labels of $R$ are pairwise inequivalent, and none is equivalent to $f(p)$: use unambiguity of $[p]$ and reflection of equivalence by $f$.
Thus a label-preserving total coarsening of $D\cup\{t\}$ is bijective on its carrier and can only remove order relations.

\item \textbf{Rigidity of the old order.}
The finite configuration/extremal-realization correspondence of \cite[Proposition~5.1]{devisme-winskel} identifies $D$ with an extremal realization; this finite case uses no selection of an infinite extremal.
Consequently no order relation within $D$ can be removed.

\item \textbf{Minimality at the new root.}
If a dependency into $t$ were removed, its new strict past would be a proper down-closed $D'\subsetneq D$ with $\epsilon[D']\cup\{f(p)\}\in\mathcal F$, contrary to minimality.
The new history is therefore prime extremal and unambiguous.

\item \textbf{Extension event.}
Let $e_D\in E$ be its finite-code isomorphism class.
Its proper predecessor set is exactly $D$ under the principal-history identification.
Define $\lambda'(p)=e_D$ and retain $\lambda$ elsewhere.

\item \textbf{Configuration preservation.}
For a configuration $y\subseteq x\cup\{p\}$ containing $p$, its image is down-closed because $D\subseteq\lambda[[p)]$.
Its old image $\lambda[y\setminus\{p\}]$ is a configuration.
The labels of that image and of $D\cup\{e_D\}$ are configurations with the common bound $f[y]\in\mathcal F$.
Their compatible union supplies the consistency test in $\operatorname{er}(\mathcal F)$.
Thus $\lambda'[y]$ is a configuration.

\item \textbf{Equivalence and factorization.}
Preservation and reflection of equivalence follow from $f$, since equivalence in $E$ is equivalence of top labels; the same observation treats the equidefinedness condition.
Hence $\lambda'$ is the required factor.

\item \textbf{Finite induction.}
Induction on the size of a finite configuration proves the final assertion; only finitely many existential witnesses are used.
\end{enumerate}
\end{proof}

\begin{proposition}[Finite-choice strength of general factorization]\label{prop:edc-choice}
Over ZF, the existence-of-factor assertion of Proposition~\ref{prop:edc-extremal-universal}, for all $\mathcal F\in\mathsf{EqFam}$ and all EDC sources, is equivalent to $\mathsf{AC}_{\mathrm{fin}}$. Quotienting factor maps by pointwise equivalence does not remove this requirement when the morphisms being quotiented are actual maps.
\end{proposition}
\begin{proof}\vmode
\leavevmode
\begin{enumerate}[label=(\roman*)]
\item \textbf{Finite-set family encoding.}
Let $(A_i)_{i\in I}$ be a family of nonempty finite sets.
Form distinct root events $r_{i,a}$ for $a\in A_i$ and distinct events $d_i$.
Let $\mathcal F$ consist of all subsets $x$ of these events for which
\[
 d_i\in x\quad\Longrightarrow\quad
      (\exists a\in A_i)\ r_{i,a}\in x,
\]
with equality equivalence.

\item \textbf{Completeness of the family.}
Unions preserve this property; a root has a singleton securing chain and a $d_i$ has a two-event securing chain.
Thus this is a configuration family, without simultaneously selecting any of the witnesses $a$.

\item \textbf{Classification of extremal top histories.}
The prime extremal histories with top label $d_i$ are exactly $r_{i,a}<d_i$, for $a\in A_i$.
To check exhaustiveness, the top's realization contains some such root label.

\item \textbf{Elimination of extra predecessors.}
An event with a root label has no necessary predecessor: removing its incoming order constraints still gives a realization, so extremality makes it minimal.
Retaining only this selected root as a strict predecessor of the top also gives a realization.
Any other predecessors of that top would therefore give a nonisomorphic total coarsening, contrary to extremality.

\item \textbf{Remaining extremal histories.}
The resulting two-event histories are extremal and unambiguous.
Singleton root histories are the other prime extremals; consequently $K$ does not remove these alternatives.

\item \textbf{Prime source and family map.}
Construct a prime event structure $Q$ with these root events and an event $t_i$ above all $r_{i,a}$ for $a\in A_i$, with no further strict order, all finite subsets consistent, and equality equivalence.
Its principal pasts are finite, so $Q$ is an EDC.
The map
\[
 f(r_{i,a})=r_{i,a},\qquad f(t_i)=d_i
\]
is total, injective, and preserves configurations.

\item \textbf{Extraction of a choice function.}
A factor $h:Q\to K\operatorname{er}(\mathcal F)$ must send $t_i$ to one of the two-event histories $r_{i,a_i}<d_i$.
The unique root label in that history defines $a_i\in A_i$.
Therefore $i\mapsto a_i$ is a choice function.

\item \textbf{Equivalence and actual factors.}
Equivalence can identify the different possible values of $h(t_i)$ after a factor exists, but every actual factor already supplies this choice function.
An inhabited quotient of a set of maps also asserts that there is such a map, so taking map classes does not evade the argument.

\item \textbf{Finite-choice selector.}
Conversely, $\mathsf{AC}_{\mathrm{fin}}$ gives a selector on all nonempty finite collections of finite subsets of the set $E=K\operatorname{er}(\mathcal F)$.
Use this single selector on the inclusion-minimal members of $\mathcal D(R,f(p))$ in Lemma~\ref{lem:edc-finite-extension}.

\item \textbf{Depth recursion.}
Define $h$ by recursion on the finite causal depth of $p$, with $R=h[[p)]$ and with domain exactly $\dom f$.
Each principal past is finite.

\item \textbf{Top labels and configurations.}
The selected history has top label $f(p)$ and all its predecessors in $h[[p)]$.
For any configuration $x$ of $P$, its image is consequently down-closed.
The labels of a finite union of principal histories in $h[x]$ are configurations with the common bound $f[x]$; compatible-union closure proves the consistency test for $E$.
Hence $h[x]$ is a configuration, including when $x$ is infinite.

\item \textbf{Factor-map conditions.}
Equivalence preservation and reflection follow from $f$ and top-label equivalence.
Thus $h$ is a map and $\epsilon h=f$.

\item \textbf{Uniform selection and uniqueness.}
The selector is fixed before the depth recursion, so no dependent-choice step is used.
Any two factors are equidefined and have equivalent top labels at every event, proving the required uniqueness up to equivalence.
\end{enumerate}
\end{proof}

\begin{proposition}[The cover of all finite local factors]\label{prop:edc-local-cover}
For $f:I(P)\to\mathcal F$ as above there is, in ZF without choice, a constructed span of EDC maps
\[
 P\xleftarrow{\ \pi\ }\widetilde P
   \xrightarrow{\ H\ }K\operatorname{er}(\mathcal F),
 \qquad fI(\pi)=\epsilon I(H).
\]
The map $\pi$ is total, is an order isomorphism on each principal past, and lifts every finite configuration and every finite configuration extension from each of its lifts. It also lifts each finite configuration cube. These are finite-history statements; no global factor or lift of an arbitrary infinite configuration is selected.
\end{proposition}
\begin{proof}\vmode
\leavevmode
\begin{enumerate}[label=(\roman*)]
\item \textbf{Inhabited local factors.}
Write $E=K\operatorname{er}(\mathcal F)$ and construct the sets
\[
 L(x)=\{\lambda:x\rightharpoonup E:
       \lambda\text{ is an EDC map and }\epsilon\lambda=f|_x\}
 \qquad(x\in\Conf_f(P)).
\]
Restriction defines $L(y)\to L(x)$ whenever $x\subseteq y$.
Lemma~\ref{lem:edc-finite-extension} proves that these sets are inhabited, and proves surjectivity of each restriction along a one-event extension, without choosing a simultaneous family of their members.

\item \textbf{Events and order of the cover.}
The events of $\widetilde P$ are pairs $(p,\lambda)$ with $\lambda\in L([p])$.
Put
\[
 (q,\mu)\leq(p,\lambda)
 \quad\Longleftrightarrow\quad q\leq p\ \text{ and }\
 \mu=\lambda|_{[q]}.
\]
A finite set is consistent when its projected events are consistent in $P$ and its local maps agree on all overlaps of their principal pasts.

\item \textbf{Consistency and finite pasts.}
This condition is hereditary and is preserved by adjoining predecessors, since predecessors carry restrictions of the same maps.
Every singleton is consistent.
The principal past of $(p,\lambda)$ is precisely $\{(q,\lambda|_{[q]}):q\leq p\}$, so it is finite and projects isomorphically to $[p]$.

\item \textbf{Event equivalence.}
Define equivalence by $(p,\lambda)\equiv(q,\mu)$ iff $p\equiv q$.
Unambiguity of each principal past follows from that of $P$.

\item \textbf{Projection and factor maps.}
Define $\pi(p,\lambda)=p$ and $H(p,\lambda)=\lambda(p)$ where $f(p)$ is defined.

\item \textbf{Local injectivity and gluing.}
Two consistent events over the same $p$ have the same local map, so projection is locally injective.
Compatible local factors glue by union.

\item \textbf{Configuration preservation.}
Their image under $H$ is down-closed, and the top labels of every finite union of principal histories lie in a common image configuration under $f$; this proves consistency.
Preservation and reflection of equivalence follow from the corresponding properties of $f$ and $\pi$.
These arguments apply also to infinite configurations, by checking finite consistency.

\item \textbf{Commuting factorization.}
Thus both displayed arrows are EDC maps and their displayed equation holds.

\item \textbf{Finite extension lifting.}
More precisely, gluing and restriction are inverse in
\[
 \Conf_f(\widetilde P)\cong
 \{(x,\lambda):x\in\Conf_f(P),\ \lambda\in L(x)\}.
\]
The gluing argument proves the forward direction; the inverse sends $(x,\lambda)$ to $\{(p,\lambda|_{[p]}):p\in x\}$.
Lemma~\ref{lem:edc-finite-extension} therefore gives extension lifting from every existing lift.

\item \textbf{Concurrent cube lifting.}
A finite configuration cube has a finite family of independent events enabled at its initial configuration.
Extend its given factor to each of their principal pasts.
Overlaps lie in the initial configuration, so the extensions agree there and are jointly consistent.
Their principal orders create no new relation among the independent events.
They therefore define a lifted cube with exactly the projected faces.
\end{enumerate}
\end{proof}

The sets $L(x)$ form a presheaf on finite configurations. The configuration calculation realizes its compatible local factors as an event structure, and the finite extension property gives the open-map form of history comparison~\cite{jnw-open-maps}, including the cube lifting proved above.

A global splitting $s$ of $\pi$ would produce the actual factor $Hs$. Such splittings cannot be asserted in ZF in general. The constructed span instead retains the alternatives themselves, so it supplies the local comparison without selecting an unfolding map.

\begin{corollary}[Comparison with the extremal presentation]\label{cor:edc-all-extremal}
In ZF there is a span
\[
 \mathsf{Caus}(\mathcal F)\xleftarrow{\pi}
 \widetilde{\mathsf{Caus}(\mathcal F)}
 \xrightarrow{H}K\operatorname{er}(\mathcal F)
\]
commuting with the two counits into $\mathcal F$. Its left leg lifts finite histories, extensions and cubes from every existing lift. No global selection of extremal causes is involved.
\end{corollary}
\begin{proof}\vmode
\leavevmode
\begin{enumerate}[label=(\roman*)]
\item \textbf{Application of the local cover.}
Apply Proposition~\ref{prop:edc-local-cover} to the constructed map $\epsilon_{\mathcal F}:I(\mathsf{Caus}(\mathcal F))\to\mathcal F$.
All its premises were proved in Proposition~\ref{prop:edc-universal}.
The resulting span has exactly the stated equations and lifting properties.

\item \textbf{Retention of all alternatives.}
In the family used in Proposition~\ref{prop:edc-choice}, the canonical all-realizations lift sends $t_i$ to the finite star with all $r_{i,a}$ below its top $d_i$.
It therefore makes no selection of $a\in A_i$, whereas every actual factor into the extremal target does.
This also explains why the new universal property does not contradict that choice characterization.
\end{enumerate}
\end{proof}

\begin{remark}[Causal comparison types and universal properties]\label{rem:causal-comparison-register}
The constructions above have the following different codomains and
universal properties. Write a bar for quotienting parallel maps
by the specified pointwise equivalence.
\begin{center}\small
\begin{tabular}{@{}p{.22\linewidth}p{.69\linewidth}@{}}
\toprule Construction & Actual comparison\\\midrule
Completion &
$\mathcal F\hookrightarrow\widehat{\mathcal F}$ is identity on events
and finite configurations; extension of partial
configuration-family maps gives the reflection into
$\mathsf{EqFam}$.\\
All realizations &
$\epsilon_{\mathcal F}:I\mathsf{Caus}(\mathcal F)\to\mathcal F$,
with the explicit partial factor $\widehat f$:
$\overline{\mathsf{EDC}}(P,\mathsf{Caus}(\mathcal F))
\cong\overline{\mathsf{EqFam}}(IP,\mathcal F)$ in ZF.\\
Extremal realizations &
$\epsilon:I(K\operatorname{er}(\mathcal F))\to\mathcal F$
exists in ZF. The assertion of factors for all EDC-source maps
has the finite-choice strength proved above.\\
All local factors &
$P\xleftarrow{\pi}\widetilde P\xrightarrow{H}
K\operatorname{er}(\mathcal F)$, with $\pi$ total and $H$ partial
on the inverse image of $\dom f$; $fI\pi=\epsilon IH$.
This gives finite history, extension and cube lifting in ZF,
without a selected splitting of $\pi$.\\
\bottomrule
\end{tabular}
\end{center}
In particular the functor on quotient map classes in the second row
is explicitly
\[
 \overline{\mathsf{Caus}}([g])
      =[\widehat{g\epsilon_{\mathcal F}}],
 \qquad g:\mathcal F\rightharpoonup\mathcal G.
\]
Both its identity comparison and its composite comparison have the
same counit composite and the same partial domain; uniqueness up
to equivalence proves the functor equations in the quotient category.
The counit equation proves naturality there. No corresponding
unselected extremal factor is used to define this functor.
The finite top-labelled realization codes, total downset-preserving
coarsenings, and top-label equivalence give the dictionary to the
external extremal-realization terminology. The all-realizations
construction deliberately retains the nonminimal orders absent
from that extremal target.
\end{remark}


\subsection{Sequential-worker control implementations}\label{subsec:register-control-implementation}
The abstract commitment compiler can now be compared with sequential workers and their registers. The implementation proof follows the same physical histories while accounting for the local commitments needed before each step.

The following implementations instantiate Section~\ref{sec:distributed-games} on the sequential participant histories constructed above. Here the current local view equals the action's causal decision view by Lemma~\ref{lem:sequential-worker-decision-view}. For display simplicity use identity control names; replacing each membership test by its exposed $\nu_p$ gives the local-name version. The register implementation and the causal-view-slot implementation both erase onto the same $\Conf_f(U_k)$ and lift its steps and cubes. They need not have identical administrative events.
\begin{definition}[Compiled plant $C(N)$]\label{def:register-compiler}
Add to each participant a register
$\eta_p\in\{\bot\}\amalg\Pow(A_p^v)$, initially $\bot$. There are two kinds of transition:
\leavevmode
\begin{enumerate}
\item a positive local transition $\mathsf{commit}_p(K)$, enabled at $\eta_p=\bot$, sets $\eta_p=K$ and leaves the plant state fixed;
\item a negative transition of plant type $a$, enabled when $a$ is plant-enabled and every $p\in V(a)$ has $\eta_p\ne\bot$ and $a\in\eta_p$, performs the plant update and resets $\eta_p=\bot$ for all $p\in\dom(a)$.
\end{enumerate}
The support of a commitment is $\{p\}$. Plant transitions retain their original supports and residual updates. In particular an uncontrollable transition never waits for a commitment.
\end{definition}

\begin{definition}[Commitment erasure and selection]\label{def:register-formal-commitment-erasure-and-selection}
Commitments on different participants commute. A commitment and a plant transition commute when the participant is outside the transition's support. Independent plant transitions commute by Proposition~\ref{prop:resource-charts} and change disjoint sets of commitment registers. This gives the compiled cubical execution and its occurrence unfolding $U_{C(N)}$.

Let $\epsilon$ erase commitments from a compiled execution and retain the induced order on plant events. The \textbf{controller} selects at an uncommitted participant the unique commitment
\begin{equation}\label{eq:register-selectedcommit}
K=k_p(o_p(v_p(\epsilon h))).
\end{equation}
Control commitments are erased before taking the observation; their scheduling therefore discloses no extra information.
\end{definition}

\begin{lemma}[Register invariant]\label{lem:register-register}
In every execution using the selected commitments, each register with $\eta_p\ne\bot$ equals $k_p(o_p(v_p(\epsilon h)))$ at the current erased history.
\end{lemma}
\begin{proof}\vmode
\leavevmode
\begin{enumerate}[label=(\roman*)]
\item \textbf{Register invariant through execution.}
The equality holds at the commitment that sets the register.
Until the next plant action involving $p$, its local view does not change, so neither does its observation.
That next action resets the register.
Induction over the execution proves the claim, including when an uncontrollable action occurs before a possible commitment.
\end{enumerate}
\end{proof}

\begin{theorem}[Execution and concurrency preservation]\label{thm:register-compiler}
Selected compiled execution and controlled execution of $N$ are related by a rooted weak bisimulation after erasing commitments. Every finite controlled history lifts, and erasure preserves the causal order on plant events. Every finite independent enabled family of controlled plant transitions lifts, after finitely many commitments, to a compiled cube. Preparing an enabled controlled action $a$ requires at most $|V(a)|$ commitments; a finite controlled path $a_1\cdots a_n$ requires at most $\sum_i|V(a_i)|$.
\end{theorem}
\begin{proof}\vmode
\leavevmode
\begin{enumerate}[label=(\roman*)]
\item \textbf{Forward erasure.}
A commitment erases to an identity.
If a compiled plant action executes, its guard and Lemma~\ref{lem:register-register} give exactly \eqref{eq:controlled}, so its erasure is allowed.

\item \textbf{Preparation of an allowed action.}
Conversely suppose $a$ is allowed after an erased history and choose any selected compiled lift.
Make the selected commitments of the currently uncommitted participants in $V(a)$.
These moves do not change any plant view.

\item \textbf{Reverse weak simulation.}
The guard now holds, so $a$ executes and extends the lift.
Lemma~\ref{lem:sequential-worker-decision-view} identifies these current-view guards with the abstract causal guards.
This is the step-by-step weak simulation in the other direction.

\item \textbf{Path lifting.}
Iteration proves finite-history lifting and the rooted weak bisimulation relation.

\item \textbf{Physical order preservation.}
Each commitment is local: its plant causal predecessors lie in its owner's preceding view, and its plant successors involving that owner already depend on that view by local order.
Removing the commitment therefore introduces no new order between plant events.
Synchronization edges and FIFO matching edges are unchanged.
Conversely all generating plant causal edges remain, so the induced order is exactly the source history order.

\item \textbf{Concurrent-family lifting.}
For a finite independent family, first fill all missing veto commitments.
Its members have disjoint participant supports; the possible shared queue calculation remains the one in Lemma~\ref{lem:fifo-square}.
Executing one member resets no register needed by another.
Their residual witnesses and guards therefore persist, giving the compiled cube with the required erased faces.

\item \textbf{Finite preparation bounds.}
Only currently missing registers in $V(a)$ need to be filled before $a$.
This set is finite because $V(a)\subseteq\dom(a)$.
Summing these bounds gives the path bound.

\item \textbf{Locality of preparation.}
For a finite independent family the preparation set is contained in the finite union of its veto sets.
No global completion of all participant registers is used in these lifting arguments.
\end{enumerate}
\end{proof}

\begin{remark}[Local bounds and commitment-only execution]\label{rem:register-compiler-divergence}
For $P=\mathbb N$, the initially enabled selected commitments at distinct participants give an infinite commitment-only execution, even for a plant of independent one-shot events. Thus the general compiler has finite local preparation, not global termination of administrative motion. For a specified finite participant network, each commitment with fixed plant history decreases the number of uncommitted registers, giving the bound $|P|$. The winning construction below handles arbitrary $P$ by a directed union of all finite preparations, without selecting an infinite execution.
\end{remark}

The theorem is an execution comparison. Erasing an edge and changing the interpretation of the remaining polarities is not asserted to be a morphism of polarized precubical presheaves.


\begin{proposition}[Register and causal-view comparison]\label{prop:register-causal-view-comparison}
The selected register implementation and $S_k\to\mathcal C(U)$ realize the same controlled physical histories, steps and concurrent families. Filling all current registers in a finite participant network takes at most $|P|$ commitments. Its positive-maximal configurations have the same physical winning test as the general compiler.
\end{proposition}
\begin{proof}\vmode
\leavevmode
\begin{enumerate}[label=(\roman*)]
\item \textbf{Identification of guards.}
The register invariant and Lemma~\ref{lem:sequential-worker-decision-view} identify its guard with $g_k(e)$.

\item \textbf{Common-history comparison.}
The two execution-lifting theorems therefore give surjections to the same controlled histories, each lifting extensions and cubes from every representative.
This supplies the comparison as the fibre product of their history maps, with the identical physical coordinate.

\item \textbf{Commitment completion.}
At fixed physical history, each register commitment reduces the number of uncommitted participants and no commitment resets one, proving the finite bound.
With arbitrary participant sets, the directed union of finite preparations fills all current registers, and each missing one has an enabled selected commitment.

\item \textbf{Maximality and winning transfer.}
Hence positive maximality is exactly completion; the register invariant identifies its enabled physical moves with $\En_k(h)$.
The proof of Theorem~\ref{thm:winning-transfer} applies with this saturation.
\end{enumerate}
\end{proof}

\subsubsection{Certificates in sequential communication implementations}
\begin{definition}[Causal-memory certificates]\label{def:formal-causal-memory-certificates}
Example~\ref{ex:private-bit-fixed-ping} has a constructive repair when the communication mechanism is allowed to carry metadata. Give each participant a finite labelled history register. Occurrences receive names formed from owner and local counter; a joint action uses the finite tuple of participating counters. Names are unique within an execution and identify common subhistories during merges.

Certificates are finite directed acyclic graphs whose event labels are the original plant labels and payloads. They do not recursively contain the enriched messages carrying them. Thus including the current send event in its certificate is a finite construction, not a circular definition. Event-kind tags distinguish the naming formats, and participating counters advance at every local or joint event.

For a local event append that event to the register. At a rendezvous merge the participating registers and append the joint event. At a send, transmit the payload and the sender's history including that send. At the matching receive, merge the transmitted history with the receiver's register and append the receive. Initially each register and each certificate on an initial boundary packet is empty. A certificate is stored with occurrence names and the induced causal order. The register invariant below supplies embeddings into one common full history; merge is the union of their images with the inherited order. Thus compatibility is proved from the construction, rather than supplied as an additional input.
\end{definition}

\begin{proposition}[Causal-memory implementation]\label{prop:causalmemory}
After every finite execution $h$, each register $C_p(h)$ has a name-preserving order embedding
\[
 j_{p,h}:C_p(h)\hookrightarrow h,\qquad
 \im(j_{p,h})=v_p(h).
\]
Each in-transit certificate similarly embeds as the view at its sending event, or as the empty history for an initial boundary packet. The embeddings agree on common occurrence names. In particular the register at $p$ is exactly $v_p(h)$. Erasing registers and certificates recovers the original plant's executions and independent cubes. Every full-causal-view controller is therefore an observation-based controller on this enriched protocol.
\end{proposition}
\begin{proof}\vmode
\leavevmode
\begin{enumerate}[label=(\roman*)]
\item \textbf{Initial and local certificates.}
Proceed by induction on an execution.
Initially the empty graphs have their unique embeddings.
A new local event extends its owner's previous downset; map its fresh name to that event.

\item \textbf{Rendezvous amalgamation.}
At a rendezvous, the existing embeddings identify the participant registers with downsets of the same preceding history.
Their union is a downset and their orders agree on intersections.
Append the new joint event above this union.

\item \textbf{Send and receive certificates.}
A send constructs its certificate after appending the send, so its image is the send's principal history.
At reception, the stored certificate and the receiver's current register again embed as downsets of one history; union and append the receive.
For an initial packet only the receiver's preceding view is required.

\item \textbf{Exact causal image.}
These are exactly the generating predecessor relations of Section~\ref{subsec:histories}.
Hence each embedding has the asserted image.

\item \textbf{Uniqueness and event orders.}
Unique names make its action forced on every event, proving agreement on overlaps.
Downsets inherit all relations between their events, so graph union with inherited order computes precisely the claimed merge.

\item \textbf{Absence of new guards.}
Certificates impose no new plant guard and change no payload value.

\item \textbf{Independent squares.}
Independent events affect disjoint participant registers.
If they also act on one queue, one appends a new certified message and the other removes an earlier certified message; these updates commute by the same word calculation as Lemma~\ref{lem:fifo-square}.
Thus erasure preserves and lifts the execution cells.

\item \textbf{Controller interpretation.}
Applying a causal-view policy to the reconstructed register gives the stated controller.
\end{enumerate}
\end{proof}

The implementation can use unbounded messages and memory. Its purpose is exact information realization. It does not prove finite-state implementability or identify the enriched observations with the observations of the original fixed-payload protocol.


\begin{example}[Private bit and fixed ping]\label{ex:private-bit-fixed-ping}
The environment supplies a private bit $b$ to $p$. Participant $p$ sends a fixed payload \texttt{ping} to $r$, which must then guess $b$. The two payload observations of $r$ coincide. A deterministic observation-based guess must consequently be the same in both branches and fails in one. The causal views differ: the receive depends on the send, which depends on the private-bit event. A full-causal-memory controller can read that event's label and guess correctly. Thus a causal-memory winning strategy cannot be transferred to this fixed-payload protocol by identifying dependence with knowledge.
\end{example}
\begin{example}[Required unions at indistinguishable views]
Two local views may have the same observation although one occurs before an included $a$-extension and the other only before omitted $a$-extensions. The required union permits $a$ at both views. The candidate can still be realizable if another participant refuses every omitted synchronization. Formula~\eqref{eq:refusal} checks exactly those other vetoes. Equality of the outgoing-label sets at the two representatives would incorrectly reject such a candidate.
\end{example}
\begin{example}[Returning from a FIFO queue]
Take one channel, an initial queue $mv$, an enabled enqueue of $n$ by its producer, and an enabled dequeue of $m$ by its consumer. An entry marker precedes these two actions and an exit marker is placed after both. The two completed routed words
\[
 e\,\mathsf{enq}_n\,\mathsf{deq}_m\,f,
 \qquad
 e\,\mathsf{deq}_m\,\mathsf{enq}_n\,f
\]
are distinct points of the return correspondence, with the same final queue $vn$. They are the two linear extensions of one internal two-event configuration and bound its execution square. The Scott functor orients that square according to the events' declared polarities. If the queue is initially empty, a dequeue of the newly enqueued message instead depends on the enqueue, and there is no such square. For rendezvous, the matching send and receive are one internal joint witness; independent rendezvous on disjoint participants produce the corresponding cubes. Thus the return construction preserves the operational distinction between synchronization, causal message delivery, and independent progress.
\end{example}


\subsection{Typing of the concrete medium operations}
\begin{proposition}[Operational typing and concurrent witnesses]\label{prop:concrete-medium-typing}
The concrete operational rules retain well-scoped, sorted runtime states and their selected witnesses. A constructed resource chart gives the corresponding commuting conversions of execution paths with the same residual witnesses.
\end{proposition}
\begin{proof}\vmode
\leavevmode
\begin{enumerate}[label=(\roman*)]
\item \textbf{Rendezvous typing.}
Rendezvous specializes the input body by a value of its declared sort, so Theorem~\ref{thm:substitution-indexed-interpretation} types the target continuation.
The output continuation already has that scope.

\item \textbf{FIFO and scope.}
FIFO append inserts a packet of precisely the queue's payload sort, and head removal returns a value of that same sort.
Restriction keeps its scope around the remaining store and continuations.

\item \textbf{Fresh names.}
Fresh-name creation uses the tagged producer/output-position names of Definition~\ref{def:formal-resource-occurrence-names}, so it cannot capture an existing occurrence.

\item \textbf{Removal and loss.}
Bag removal and cancellation remove a selected occurrence without changing the sorts of remaining packets.
Loss has the same property.

\item \textbf{Duplication and dispatch.}
Duplication and dispatch construct children with the selected parent's payload sort.

\item \textbf{Replay and exchange.}
Replay selects a retained record and constructs a new delivery occurrence of its recorded sort.
Exchange reorders the selected packets.

\item \textbf{Capacity and resource updates.}
Capacity and status branches retain the same sorted update on their tested source subset.
Credit and allocation updates remove and insert the explicitly declared resource sorts.

\item \textbf{Nullary operations.}
Nullary updates have no substituted input body.
These exhaust the medium operation families of Appendix~\ref{app:communication-examples}.

\item \textbf{Joint typed substitution.}
For a joint witness the selected participants are distinct; apply the preceding checks to its typed premises and leave the other components unchanged.
Substitution in all receiving continuations is sort preserving.

\item \textbf{Service and administrative operations.}
The later symbolic, controller and administrative rule extensions have their own typed interpretations in Sections~\ref{subsec:calculus-control}, \ref{subsec:cryptographic-process-rules}, and~\ref{subsec:administrative-process-calculus}.

\item \textbf{Structural and recursive transport.}
Structural transports preserve scopes by their binder equations.
Recursive unfolding substitutes a well-scoped recursive closure for its bound process variable; a finite transition derivation uses only finitely many such unfolding rules.
Induction on that derivation, rather than an attempted recursion through infinite execution, proves the recursive case.

\item \textbf{Resource charts.}
A resource chart is constructed by testing all its finite permutations and equalizing the states at equal subsets of executed directions.
In a two-direction chart, its two composites therefore have the same source, target, selected instructions and transported supplier records.

\item \textbf{Higher cells and scope.}
The same equal-subset condition proves every higher face equation.
This is exactly the earlier resource-coefficient construction, now used as a typed operational commuting conversion.
No global progress statement is needed for this subject-reduction proof.
Deadlocks, refusal, termination and infinite internal execution continue to have the separate definitions and results already present above.
\end{enumerate}
\end{proof}
\begin{corollary}[Communication instances of the abstract theory]\label{cor:communication-protocol-instances}
The ordinary displays constructed in this appendix instantiate Definition~\ref{def:abstract-communication-protocol}. Their networks are interpreted by Theorem~\ref{thm:abstract-communication-interpretation}; their retained causal presentations have the histories and residuals of Section~\ref{subsec:abstract-protocol-histories}, and their declared controls have the compiler and winning interpretation of Section~\ref{sec:distributed-games}.
\end{corollary}
\begin{proof}\vmode
\leavevmode
\begin{enumerate}[label=(\roman*)]
\item \textbf{Ordinary medium strategies.}
Each displayed construction verifies receptivity and courtesy at its stated interface.

\item \textbf{Network closure.}
The network theorem closes these declarations under composition and tensor.

\item \textbf{History and control applications.}
Applying the generic history, residual and control theorems to their constructed event structures gives the assertions.

\item \textbf{Physical audit comparisons.}
Physical audit carriers enter through the comparison maps proved for each instance; this does not identify compatible alternative physical causes with one prime event identity.
\end{enumerate}
\end{proof}

\section{Cryptographic protocol structures}\label{app:cryptographic-examples}
The examples below construct the witnesses required by Section~\ref{sec:cryptographic-communication-games}. Five algebraic families organize the development: protected presentations, actions and threshold symmetries, witness families, free evaluation algebras, and phased disclosure. Each instance specifies its operations, reference-state updates and strategy judgement. Where an actual simulator is constructed, it proves $\sigma\Vdash_H r\preccurlyeq s$, with correctness supplying the reverse factorization. Origin, nondeducibility, local privacy and ideal-resource constructions retain their separately stated conclusions.

\subsection{Protected presentations and channel realization}\label{subsec:protected-presentations}
The free introductions determine private operation algebras. Image
elimination constructs their decoders and checks; labelled public
presentations determine which representations the interface can distinguish.
Reference realization then supplies the strategies and comparison witnesses
of Proposition~\ref{prop:primitive-service-comparison}.

\begin{construction}[Protected records from typed introductions]\vmode\label{con:protected-records}
\leavevmode
\begin{enumerate}[label=(\roman*)]
\item \textbf{Name worlds and key interfaces.}
Work over the finite name worlds of Definition~\ref{def:calculus-indexed-carriers}.
Private key and nonce inputs are allocated objects; a public-key object is a tagged copy of its private-key object, with the constructed map $p:K\to J$.
A shared-key interface instead uses $J=K$ and $p=1_K$.

\item \textbf{Participant key access.}
Whether a participant obtains a key is determined by the connected allocation and publication ports.

\item \textbf{Free record introductions.}
There are two record introductions.
Sealing adjoins a tagged record of a message, an encryption key and a nonce; attestation adjoins a tagged record of a message and an issuing key:
\[
 \mathsf{seal}:M\times J\times N\longrightarrow M,
 \qquad \mathsf{attest}:M\times K\longrightarrow M.
\]
For the finite network of such introductions, form its message carrier by the free value construction of Definition~\ref{def:formal-sorted-value-contexts-and-process-constructors}: start with ordinary data and the typed envelopes of its allocated and public-key objects, and close under pairing and these record introductions.
Equivalently its polynomial is the coproduct of these actual introduction domains.

\item \textbf{Free carrier and disjoint images.}
The union of the finite-depth stages is initial: interpreting a root after its finitely many children constructs the unique fold.
Thus record images are disjoint and each introduction is injective.
The carrier and its fold are derived from the introductions, rather than supplied as further operation data.

\item \textbf{Successful decoder image.}
For sealing form the successful-record graph $E=K\times M\times N$, with
\[
 a(k,m,r)=(k,\mathsf{seal}(m,p(k),r)),\qquad b(k,m,r)=m.
\]
Factor $a$ as $E\xrightarrow{q}D\xhookrightarrow{i}K\times M$.

\item \textbf{Kernel-pair descent.}
Equality of two $a$-images identifies the message operands by injectivity of the sealing introduction.
Hence $b$ coequalizes the kernel pair of $q$, and image elimination constructs the unique $d:D\to M$ with $dq=b$.
This is the partial decoder, including its exact successful domain; no decoder or retraction law is independently supplied.

\item \textbf{Attestation verification.}
For attestation take the image of
\[
 (k,m)\longmapsto(p(k),m,\mathsf{attest}(m,k))
\]
in $J\times M\times M$.
Verification is its checked membership operation.

\item \textbf{Root and key checks.}
Root inspection, followed by comparison of the retained key image and message, computes this check.
On the tagged public-key copy $p$ is injective; on a shared key it is the identity.

\item \textbf{Other image eliminators.}
Pair projections, ordinary-data reading and typed coercions are likewise the partial projections or inverses on their introduction images.
Outside those domains they return failure.

\item \textbf{Concrete equality.}
These constructions use the actual equality operation of the value instance; for arbitrary external ordinary data this is a set-valued operation, not an asserted decision procedure.
Finite bitstrings give its terminating implementation.

\item \textbf{Public interface.}
Expose only these specified introductions and checks through Construction~\ref{con:reference-interface}.
In particular the private product representation does not supply public projections of an encryption's key, nonce or plaintext.

\item \textbf{Strict evaluation and substitution.}
Strict evaluation is the existing process elimination for the resulting partial maps.
Its substitution law is Theorem~\ref{thm:substitution-indexed-interpretation}; naturality under name injections follows because they preserve and reflect tags and operand equality.
\end{enumerate}
\end{construction}

The decoder equation is the image-factorization equation $dq=b$; successful verification is membership in the constructed issuance image. An operation audit additionally retains the issuing occurrence and its selected operands, so its fold supplies an origin map. These are correctness and provenance statements. Confidentiality requires the public representation comparison below: the decoding equation also holds for an implementation with a publicly readable plaintext port.

\subsubsection{Generated sealing and authentication services}\label{subsec:crypto-values}
The channel constructions require exact introduction, decoding and checking operations. We first build these services on private value carriers; their subsequent reference realization determines which operations a participant can actually invoke.

\begin{construction}[Sealing and authentication instances]\vmode\label{con:sealing-authentication-instances}
\leavevmode
\begin{enumerate}[label=(\roman*)]
\item \textbf{Instantiation of protected records.}
Instantiate Construction~\ref{con:protected-records} on four separate key-allocation ports: encryption, symmetric encryption, signing, and MAC.
The two sealing instances and two attestation instances are
\[
 \begin{array}{c|cc}
 &\text{public-key port}&\text{shared-key port}\\\hline
 \mathsf{seal}&(\crAEnc,\crPK)&(\crSEnc,1)\\
 \mathsf{attest}&(\crSig,\crVK)&(\crMac,1).
 \end{array}
\]
The pairs record the introduction and its key map; their input objects are supplied by those ports.

\item \textbf{Key and message carriers.}
Write $T_{\mathsf K_e},T_{\mathsf K_s},T_{\mathsf K_\sigma},T_{\mathsf K_m}$ for the allocated key carriers, $T_{\mathsf N}$ for nonces, and $T_{\mathsf{Ch}}$ for the existing scoped channel names.
The public-key carriers $T_{\mathsf P_e},T_{\mathsf P_\sigma}$ are the tagged images of $\crPK,\crVK$.
The free message carrier is $T_{\mathsf M}$, abbreviated $T$; its envelope inclusions are $\iota_\tau$.

\item \textbf{Notation and derived checks.}
The notation records the objects produced by these four instances.
The derived decoders are $\crADec,\crSDec$; the attestation checks are $\mathsf{vfy},\mathsf{mcheck}$.
Thus, for example,
\[
 \crADec(\crAEnc(m,\crPK(k),r),k)=m,
 \qquad \mathsf{vfy}(\crSig(m,s),m,\crVK(s))=\crOK.
\]
The shared-key equations are the identity-key instances of the same construction.

\item \textbf{Failure and other eliminators.}
All unmatched checks fail.
Pair projections, typed coercion and data reading are its other image eliminators.

\item \textbf{Typed operation replies.}
Successful value replies have their declared value sort, verification replies have type $\{\crOK,\crFail\}$, and data reading has type $D+\{\crFail\}$.
Operations are invoked by the strict process rules.

\item \textbf{Name allocation.}
Allocation is tagged extension of the cumulative name world.
Its concrete realization is constructed in Proposition~\ref{prop:crypto-freshness}; neither that representation nor an arbitrary-name import is a public operation.
\end{enumerate}
\end{construction}

\begin{corollary}[Evaluation and name transport]\label{cor:evaluation-name-transport}
Every finite strict value program has its unique value or failure. The decoder and verification equations above hold, and evaluation commutes with sort-preserving injections of allocated names, including failure and equality tests.
\end{corollary}
\begin{proof}\vmode
\leavevmode
\begin{enumerate}[label=(\roman*)]
\item \textbf{Unique strict evaluation.}
The free value fold evaluates introductions.
The image eliminators of Construction~\ref{con:protected-records} compute the unique checked continuation, and strict process composition propagates failure.

\item \textbf{Naturality under name injections.}
Name injections preserve and reflect every image-membership test; induction over the existing finite program therefore proves the naturality equation.
Its bijective case is renaming.

\item \textbf{Symbolic randomness and extension scope.}
Encryption randomness here is an allocated nonce, with no probability distribution.
Additional algebraic operations require their own constructed carriers and comparison maps.
\end{enumerate}
\end{proof}

\subsubsection{Reference processes and allocation}\label{subsec:reference-processes-allocation}
\begin{construction}[Cryptographic reference processes]\vmode\label{con:cryptographic-reference-processes}
\leavevmode
\begin{enumerate}[label=(\roman*)]
\item \textbf{Reference-bank realization.}
Apply Construction~\ref{con:reference-interface} to the generated services.
A bank is the process $\mathsf{Serve}(s)$ of Definition~\ref{def:formal-recursive-service-process}; its private residual $s$ retains the finite immutable reference table $H$, allocation state and audit.

\item \textbf{Command evaluation.}
A command looks up its typed operands, executes the corresponding strict operation, and issues a fresh handle for a successful value; tests return their observable result and failure creates no entry.

\item \textbf{Data and copy operations.}
Data introduction takes an element of $D$.
Copy is the identity operation with a newly issued reference.

\item \textbf{Import and send operations.}
An import is a matched private-wire receipt, and send is the existing medium process applied to the selected entry.
The command and reply carriers are the coproducts of these actual process interfaces.

\item \textbf{Independent banks and handles.}
Use one bank for each issuing worker and tensor independent banks.
A handle is its bank tag and issuance ordinal; the client frame is their coproduct.

\item \textbf{Ownership checks.}
Wrong-owner, missing and wrongly sorted handles fail at lookup.

\item \textbf{Shared-bank request lanes.}
Several clients of one bank connect by separate FIFO request/reply lanes.
Its worker selects a ready lane, performs one service update, replies on that lane, and recurs.
This is the shared-readiness and service construction already given; accesses to one bank retain its chain, whereas distinct banks have independent controls and counters.

\item \textbf{Gateway routing.}
Gateways receive their public call before routing these private requests.

\item \textbf{Operation audit.}
Attach the operation audit of Lemma~\ref{lem:generic-operation-audit}: a successful result records its operation and selected operand occurrences; copy and transport retain their selected predecessor.
Thus the audit, its stored values, and the residual states are obtained by elaborating the bank processes.

\item \textbf{Value equality and resource use.}
Public equality compares values; it does not identify their issuing occurrences.
A reusable key is a private immutable entry, while a one-use authorization is a consumed resource token.
\end{enumerate}
\end{construction}

\paragraph{Initialization.}
Before executing the reference processes, prepare their finite initialized frame by Proposition~\ref{prop:crypto-initialization} and retain its audit. Each allocator extends the cumulative name world, including names no longer publicly held. Independent workers use occurrence-tagged extensions; Proposition~\ref{prop:crypto-freshness} realizes these extensions and their coherent exchanges, including imports.

\paragraph{Representation by dependency diagrams.}\label{par:dependency-diagram-representation}
Finite sorted acyclic node tables are the dependency diagrams of Lemma~\ref{lem:generic-operation-audit}, with sharing permitted. Its unique unfolding gives $\alpha:S_N\to S_T$ on reference-process states, fixing public references, counters and provenance. Constructor nodes append their selected predecessors; image eliminators inspect the unfolded root and compare the indicated operands.
\begin{corollary}[Representation independence of symbolic storage]\label{cor:symbolic-storage-representation-independence}
The unfolding satisfies
\begin{equation}\label{eq:crypto-coalgebra-square}
 (\alpha\times1_R)\delta_N=\delta_T(\alpha\times1_Q).
\end{equation}
Corresponding initialized reference strategies are isomorphic and interchangeable in concurrent contexts.
\end{corollary}
\begin{proof}\vmode
\leavevmode
\begin{enumerate}[label=(\roman*)]
\item \textbf{Evaluator square.}
The audit fold commutes with introductions; its root and operand equations give the same checked image eliminations, including failure and equality.

\item \textbf{Reference transition square.}
Lookup and issuance retain their labels.
These are the operation squares for the reference construction, hence give the displayed coalgebra square.

\item \textbf{Contextual replacement.}
Apply Theorem~\ref{thm:service-behaviour-canonical-comparison} to the graph of $\alpha$.

\item \textbf{Effective unfolding.}
Recursive unfolding is effective on finite tables, giving Proposition~\ref{prop:effective-name-action}'s effective comparison.
Sharing need not be injectively represented by trees.
\end{enumerate}
\end{proof}

\subsubsection{Examples: authentication and confidentiality protocols}\label{subsec:crypto-channels}
The channel proofs below construct the simulator first and compare the resulting open networks, retaining secret inputs at their honest ports. Their causal certificates convert to complete trace certificates by Construction~\ref{con:certificate-interaction}; Theorem~\ref{thm:observational-simulation} then supplies the uniform simulation criterion against arbitrary joint tests.

\paragraph{Channel boundaries.}\label{par:channel-boundaries}
The four wrappers below use the record $\mathsf{rec}(d)=(\mathsf{channel},\mathsf{sid},A,B,d)$ and check its exact session and role fields. Each source resource includes its private key allocation, registered public-image connections, reference evaluator and communication media. Secret-key ports connect only to the designated participants.

Define the source and target resources by the following process networks. Their common boundary has one honest input, one application output, an adversarial notification request and a FIFO instruction/submission port with transport receipts. The targets are distinguished by the notification and instruction objects:
\[
 \begin{array}{c|cc}
 &\mathsf{leak}(d)&\mathsf{Sel}\\\hline
 \mathcal F_{\rm auth}&d&\{\mathsf{stored}\}\\
 \mathcal F_{\rm conf}&*&\{\mathsf{stored}\}+D\\
 \mathcal F_{\rm secret\text{-}auth}&*&\{\mathsf{stored}\}.
 \end{array}
\]
The $D$ summand is written $\mathsf{put}(d)$. The one-use selection and application processes determine delivery; submission receipts do not reveal validation. The adversary has the declared public reference commands and can retain, copy, drop and replay packets before submission. Submitted packets follow their FIFO order. Notification, input and application dependencies are those of the displayed terms, with no fairness or absence-test operation added.

\paragraph{Channel processes and their shared components}\label{par:channel-process-components}
\begin{definition}[Channel component processes]\label{def:formal-typed-worker-state-carriers}
Use the FIFO, service-slot and bank processes already constructed. For a value sort $X$, an empty or filled cell has carrier $X_\bot=1+X$. A one-shot request prefix and its reply continuation have the three residual positions unreceived, pending and answered; no separate request-flag implementation is added. Request identifiers are the addresses of their actual request occurrences.

For each FIFO $Q$, the abbreviation $\mathsf{post}_Q(i,z;P)$ appends the occurrence $(i,z)$ and enables its transport-receipt continuation $P$. Its medium is constructed from the FIFO and a finite family of receipt tokens: the enqueue commitment sends $(Q,A)$ to $(Q(i,z),A\cup\{i\})$, and the receipt consumes $i\in A$. The FIFO consumer and receipt use distinct participant controls, so neither waits for the other. A client that resumes only after receipt uses that receipt as its continuation prefix. These two updates are instances of the witnessed medium rule, with the packet and receipt sharing the selected enqueue as producer. They define the abbreviation also when the consumer is blocked.

The first-selection cell $\ell$ starts empty. Its selector consumes one queued candidate at a time, fills $\ell$ with the first valid datum and leaves a filled cell unchanged. Its application process receives its request independently and offers its one reply only after $\ell$ is filled. The one-use output occurrence prevents a second reply. A cell read tests the filled value and retains its producing occurrence without consuming the cell. Distinct readers have their own controls; their unchanged-cell updates commute by the finite-coefficient construction of Definition~\ref{def:resource-coefficient}. Thus the cell supplies a common cause without serializing its readers. The notation $\mathsf{await}_c(x;P)$ denotes this private cell-read continuation, and $\mathsf{fill}_c(x)$ its one-use producer. Their histories are the service-slot histories of Section~\ref{subsec:protocol-history-causal-interfaces}.

All definitions below are finite process terms or state-indexed recursive service equations of Definition~\ref{def:formal-recursive-service-process}. Composition hides the displayed private cell, queue and bank ports. Elaborating these terms by~\eqref{eq:committed-fibre} constructs their residual positions, intermediate calls and concurrent cells. Thus a register means a captured value in that residual syntax, and a control position means its retained prefix address.
\end{definition}

\begin{definition}[Ideal channel process]\label{def:formal-ideal-channel-transition-system}
The private input cell $a$ and selected-result cell $\ell$ start empty; the instruction FIFO $Q$ and its receipt-token family start empty. The instruction sort is $\mathsf{Sel}=\{\mathsf{stored}\}$ for authentication and secret authentication, and $\mathsf{Sel}=\{\mathsf{stored}\}+\{\mathsf{put}(d):d\in D\}$ for confidentiality. Define the independent input, notification and application processes by
\begin{align}
\mathsf{In}&=\mathsf{input}?(d).\mathsf{fill}_a(d),\label{eq:ideal-input}\\
\mathsf{Notify}&=\mathsf{notifyReq}?().\mathsf{await}_a(d;\mathsf{notifyRep}!\mathsf{leak}(d).0),\label{eq:ideal-notify}\\
\mathsf{App}&=\mathsf{appReq}?().\mathsf{await}_\ell(d;\mathsf{appRep}!d.0),\label{eq:ideal-app}\\
\mathsf{Ingress}&=\mathsf{instruct}?(i,z).
 \mathsf{post}_Q(i,z;\mathsf{receipt}!i.\mathsf{Ingress}).\label{eq:ideal-instruction}
\end{align}
Here $\mathsf{leak}(d)=d$ for authentication and $*$ otherwise. The FIFO selector has state $\ell\in D_\bot$ and the following strict computation on its retained head:
\begin{equation}\label{eq:ideal-select}
\begin{array}{c|c|c}
\text{head}&\text{cell}&\text{continuation}\\\hline
(i,z)&\ell\ne\bot&\text{consume head; retain }\ell\\
(i,\mathsf{put}(d))&\ell=\bot&\text{consume head; fill }\ell\text{ with }d\\
(i,\mathsf{stored})&\ell=\bot&\mathsf{await}_a(d;\text{consume head; fill }\ell\text{ with }d).
\end{array}
\end{equation}
The selected head remains unavailable to later iterations until its continuation completes; a stored instruction therefore blocks later instructions while awaiting input. Its receipt is already independently enabled. A direct instruction can select before input, and later input still enables notification. Consuming a head at a filled cell never produces another selection.

Let $\mathsf{Ideal}_W$ be the connection of these processes and cells. Its state carrier is derived, rather than separately specified:
\begin{equation}\label{eq:ideal-channel-state}
 Z_W=\mathsf{Pos}(\mathsf{Ideal}_W),
\end{equation}
where $\mathsf{Pos}$ is the reachable residual-ledger carrier of Definition~\ref{def:formal-residual-ledgers-and-enabled-offers}. Its coordinates retain the two cell values, queued occurrences, pending request prefixes and receipt tokens. Negative requests are accepted at their declared ports; a blocked cell read delays only its positive continuation. Theorem~\ref{thm:compact-cryptographic-communication} constructs the ordinary strategy. The generating dependencies are exactly the request, selected supplier and local continuation dependencies of these terms.
\end{definition}

\begin{definition}[Real channel processes]\label{def:real-channel-processes}
Let $\mathsf{rec}(d)$ be the fixed tagged record and let $\mathsf{parse}$ be the strict composite of projections, tag tests and data reading, with result in $D_\bot$. Thus $\mathsf{parse}(\mathsf{rec}(d))=d$. The wrapper operations are the following service programs; each displayed constructor, destructor and test is called through its typed bank interface:
\[
\begin{array}{c|l|l}
W&\mathsf{Encode}_W(d)&\mathsf{Check}_W(t)\\\hline
\mathsf{sig}&\langle\mathsf{rec}(d),\crSig(\mathsf{rec}(d),s)\rangle
 &\mathsf{vfy}(\mathsf{snd}(t),\mathsf{fst}(t),\crVK(s));\ \mathsf{parse}(\mathsf{fst}(t))\\
\mathsf{mac}&\langle\mathsf{rec}(d),\crMac(\mathsf{rec}(d),k)\rangle
 &\mathsf{mcheck}(\mathsf{snd}(t),\mathsf{fst}(t),k);\ \mathsf{parse}(\mathsf{fst}(t))\\
\mathsf{pk}&(\mathsf{new}\ r:\mathsf N)\,\crAEnc(\mathsf{rec}(d),\crPK(k),r)
 &\mathsf{parse}(\crADec(t,k))\\
\mathsf{sk}&(\mathsf{new}\ r:\mathsf N)\,\crSEnc(\mathsf{rec}(d),k,r)
 &\mathsf{parse}(\crSDec(t,k)).
\end{array}
\]
Repeated operands in a row denote the same bound intermediate result: first bind the record for encoding, or the projected record and tag for checking. These are finite strict let/call terms under Definitions~\ref{def:calculus-destructor-rules} and~\ref{def:service-call-formation-elimination}. A test must return $\crOK$ before its continuation; the first failed lookup, destructor or test returns $\bot$. The receiver result is $D_\bot$. Allocation uses the existing tagged name extension.

With a private outgoing cell $p$, submitted-packet FIFO $Q$, successful-result FIFO $C$, first-selection cell $\ell$, and separate private endpoint banks, define
\[
\begin{aligned}
S_W&=\mathsf{appIn}?(d).\mathsf{encode}_W(d;z.\mathsf{fill}_p(z)),\\
O_W&=\mathsf{packetReq}?().\mathsf{await}_p(z;
 \mathsf{import}_{\rm receipt}(z;h.\mathsf{packetRep}!h.0)),\\
I_W&=\mathsf{submit}?(i,h).\mathsf{lookup}(h;z.
 \mathsf{post}_Q(i,z;\mathsf{receipt}!i.I_W))\ \mathsf{else}\ \mathsf{failReply}.I_W,\\
R_W&=Q?(i,z).\mathsf{check}_W(z;d.
 \mathsf{post}_C(i,d;R_W))\ \mathsf{else}\ R_W,\\
B_W&=\mathsf{appReq}?().\mathsf{await}_\ell(d;\mathsf{appRep}!d.0).
\end{aligned}
\]
The $C$-consumer applies the first-selection operation to each head datum. In $R_W$, continuation after $\mathsf{post}_C$ waits for its private transport receipt, but not for selection. The public evaluator is the recursive service $E_W$ of Definition~\ref{def:formal-recursive-service-process}, using the constructed reference operations and routed lookups. Its result bank and the outgoing receipt bank have distinct issuer counters; actual accesses retain each bank's local order. The sender and receiver use their own private banks.

Connect these processes, queues and cells, hiding their private ports, to obtain $\mathsf{Real}_W$. All cells are initially empty, queues and receipt families empty, and processes at their initial prefixes. Its entire state carrier, including pending calls, is
\begin{equation}\label{eq:real-channel-state}
 X_W=\mathsf{Pos}(\mathsf{Real}_W).
\end{equation}
Proposition~\ref{prop:pc-finite-support} and~\eqref{eq:committed-fibre} construct its witnessed transitions and cells. In particular submission supplies both its receipt and its queued packet; the packet supplies checking and selection; application output depends on both its own request and that selection. There is no edge from the transport receipt to checking, selection or application output. This follows from the separate continuations in $\mathsf{post}$ and the independent application process, not from an identification of interleavings. Request identifiers, FIFO suppliers, bank operands and the one-use selection remain in the occurrence interpretation.
\end{definition}

\begin{lemma}[Opaque-ciphertext calculation]\label{lem:crypto-opaque}
Let $c=\crAEnc(u,\crPK(k),r)$ be the single honest ciphertext, where the privately allocated $k$ and $r$ have no exposed handles. In the stated wrapper, replace $c$ in the adversary's frame by a fresh opaque ciphertext marker $\kappa$. Public recipe operations have a corresponding evaluation on these markers: equality identifies the same marker; decryption of $\kappa$ with any available secret key fails; other constructors can contain it; other root destructors reject it. This correspondence preserves success, equality, and data-reading results. The same statement holds for symmetric encryption when its key has no adversarial handle.
\end{lemma}
\begin{proof}\vmode
\leavevmode
\begin{enumerate}[label=(\roman*)]
\item \textbf{Private key and nonce invariant.}
The recipient's key never occurs as an available bare key, by the provenance induction and the actual port connections.
The nonce $r$ likewise occurs outside the honest endpoints only under $c$.

\item \textbf{Inaccessibility of the sealed value.}
The only relevant destructor which could open $c$ requires its unavailable key.
No other operation extracts $r$.

\item \textbf{Marker expansion.}
Consider terms built from available atoms, ordinary constants, and the marker.

\item \textbf{Injectivity of expansion.}
Expansion $\kappa\mapsto c$ is injective on these accessible expressions modulo their evaluated values.
A term not using $\kappa$ cannot equal $c$: matching its encryption root would require a bare occurrence of $r$.
A proper expression containing $\kappa$ cannot expand to $c$ itself by finiteness of trees.
Root-by-root comparison reduces all remaining equalities to corresponding comparisons of operands.

\item \textbf{Preservation of eliminators.}
Decryption of a marker fails with every available key, while decryption of a constructed outer ciphertext has the same key comparison and returns corresponding operands.
Projection, signature/MAC checks, and data reading have the same root decisions.

\item \textbf{Recipe induction.}
Induction on recipes proves the result, including failed recipes and nested use of the marker in attacker-created terms.
\end{enumerate}
\end{proof}

The lemma's non-exposure facts are proved invariants of these closed wrappers. They are not a claim about a service whose caller is also given the secret key. Explicit disclosure updates the frame and can invalidate the abstraction; Section~\ref{subsec:admin-crypto} treats that update without asserting continued secrecy.

\paragraph{The simulator and its finite causal certificate}\label{par:channel-simulator-causal-certificate}
\begin{definition}[Simulator marker expansion]\label{def:formal-simulator-marker-expansion}
We now give the state relation and the maps used in the channel proof. All adversarial commands, including unsuccessful guesses of handles, use the same gateways, issuer counters, and routed bank protocol on the two sides. This retains their actual access order.

For signature and MAC, let $\widetilde T$ be a copy of the symbolic algebra with privately renamed honest keys and corresponding public images. Its value comparison $e:\widetilde T\to T$ is that atom renaming. The simulator obtains the honest datum from the ideal authentication notification, forms the indicated packet, and places it in its simulated outgoing store. For encryption, adjoin one fresh opaque message constant $\kappa$ to the simulator's accessible algebra. When the ideal confidential notification arrives, make that marker available as the outgoing packet. Expand it in the proof by
\begin{equation}\label{eq:marker-expansion}
e_d(\kappa)=\mathsf{Encode}_W(d),
\end{equation}
with the actual honest nonce, and extend $e_d$ through constructors and the private atom renaming. The simulator never evaluates the right side of \eqref{eq:marker-expansion}. Lemma~\ref{lem:crypto-opaque} proves that $e_d$ is injective on accessible values and commutes with all public successes, failures, equality tests, and data reading. Prior to honest input no marker is installed; the expansion then acts only on the values already available. Counters and handle identifiers are identical on the two sides.
\end{definition}

\paragraph{Executable receiver classification.}
\begin{definition}[Simulator classifier and worker transitions]\label{def:formal-simulator-classifier-and-worker-transitions}
Write $\mathsf{Instr}_W=\{\mathsf{reject},\mathsf{stored}\}+\{\mathsf{put}(d):d\in D\}$ for the classifier result set; only the public-key classifier uses its $\mathsf{put}$ summand. The simulator replaces the real private check by the following total map
\begin{equation}\label{eq:sim-classifier}
\chi_W:\widetilde T\longrightarrow
 \{\mathsf{reject},\mathsf{stored}\}+\{\mathsf{put}(d):d\in D\}:
\end{equation}
\[
\begin{array}{c|l}
W&\chi_W(t)\\\hline
\mathsf{sig},\mathsf{mac}&\begin{aligned}&\mathsf{stored}\text{ if the simulated tag and record check succeeds};\\&\mathsf{reject}\text{ otherwise}\end{aligned}\\
\mathsf{pk}&\begin{aligned}&\mathsf{stored}\text{ if }t=\kappa;\\&
 \mathsf{put}(d)\text{ if }t=\crAEnc(u,\widetilde{\crPK(k)},r_1)\text{ and }\mathsf{parse}(u)=d;\\&
 \mathsf{reject}\text{ otherwise}\end{aligned}\\
\mathsf{sk}&\mathsf{stored}\text{ if }t=\kappa;\quad\mathsf{reject}\text{ otherwise}.
\end{array}
\]
The middle test uses only the simulator's own node tags, public key, and available operands. In the first row, the origin lemma ensures that every successful check is for the one honest datum. In the last row no accessible constructor has the protected symmetric key as an argument. The classification therefore satisfies the commuting equation
\begin{equation}\label{eq:classification-square}
\mathsf{Check}_W(e_d(t))=\operatorname{val}_d(\chi_W(t)),\qquad
\operatorname{val}_d(\mathsf{reject})=\bot,\quad
\operatorname{val}_d(\mathsf{stored})=d,\quad
\operatorname{val}_d(\mathsf{put}(d'))=d'.
\end{equation}
For a pre-input public-key injection the marker case is absent, and the $\mathsf{put}$ calculation does not use $d$. In diagrams the computation is
\[
\begin{tikzcd}[column sep=large]
\widetilde T \arrow[r,"{\chi_W}"] \arrow[d,"{e_d}"']
 &\mathsf{Instr}_W \arrow[d,"{\operatorname{val}_d}"]\\
T_{\mathsf M} \arrow[r,"{\mathsf{Check}_W}"'] &D_\bot .
\end{tikzcd}
\]
The proof of the square is exhaustive. A marker has the honest encryption's successful check. A non-marker public-key ciphertext has its displayed operands in the table, so its decryption returns that $u$ and the parse decides the result. Every other root or wrong key fails. A signature/MAC passing the protected-key check has the honest origin and hence that exact record. An accessible symmetric ciphertext under the protected key is the marker. These are all root cases; nested marker occurrences are inspected by the same public destructor rules until one of these root cases is \textbf{reached}. In particular a marker in a purported ordinary-data field fails the data test.

The simulator uses the process network of \eqref{eq:real-channel-state}. Its outgoing, evaluator and ingress processes are the same terms over $\widetilde T$; the receiver consumes the same submitted queue, computes $\chi_W$, and discards $\mathsf{reject}$. For either other result, $R$ appends that instruction to the ideal selection FIFO and receives its private transport receipt. It does not wait for selection. The source FIFO $C$ corresponds to this ideal FIFO by $\operatorname{val}_d$ on its entries. The ideal cell, rather than the simulator, stores the selected datum. The simulator records its submitted symbolic instructions; it does not inspect the private ideal selection state. The public application worker is the ideal one. An initial private notification request waits for the honest input without blocking $E$ or $I$. Its result supplies the simulated packet; $O$ still waits for its own public request before returning a handle. The finite classifier is implemented by root inspections and tests in a private table, so it never needs the unknown encryption plaintext.
\end{definition}

\paragraph{Comparison of residual process positions}
\begin{construction}[Comparison of residual channel processes]\vmode\label{con:residual-channel-comparison}
\leavevmode
\begin{enumerate}[label=(\roman*)]
\item \textbf{Residual comparison carrier.}
Let $Y_W$ be the residual carrier of the ideal process connected to the simulator.

\item \textbf{Ledger and result projections.}
The residual-ledger projections of the displayed terms determine a comparison $\mathcal R_W$ over $X_W\times Y_W$: pull back the diagonals on corresponding public prefixes and occurrence identifiers, the graphs of $e_d$ on reference and submitted-packet entries, and the graphs of $\operatorname{val}_d$ on successful-result/instruction entries.
The selected-result cells use the diagonal on $1+D$.

\item \textbf{Product of slot comparisons.}
Products over the finitely occupied slots retain these witnesses; there is no separately specified state tuple.

\item \textbf{Public instruction positions.}
Request, queue and receipt positions are those of the same component processes.

\item \textbf{Private intervals and shared resources.}
At private encoding/checking intervals use the spans of their finite program prefixes: entries correspond, proper prefixes can correspond to the opposite entry or a proper prefix, and completing both programs relates their exits by~\eqref{eq:classification-square}.
Keep externally shared bank, queue and cell operations as individual occurrences.

\item \textbf{Publication and initialization.}
The outgoing publication is related to the ideal notification after the same input; neither requires the simulator to inspect the private selected cell.
Initialization is the comparison of prepared keys and reference banks.

\item \textbf{Successor invariant.}
Each common prefix preserves its diagonal.

\item \textbf{Evaluator and classifier updates.}
The evaluator comparison preserves the graph of $e_d$; queue insertion and removal preserve these graphs pointwise and keep the same supplying occurrence.
Classification transports the successful queue by $\operatorname{val}_d$; the first-selection and application rules preserve the cell diagonal.

\item \textbf{Failure and completeness.}
Failure leaves the same public continuation.
These are all constructors of the component terms, so their pullback comparison is preserved.
The following lemma proves that the private-prefix spans also preserve the visible causal boundary.
\end{enumerate}
\end{construction}

\begin{definition}[Channel dependency graphs and private blocks]\label{def:channel-contraction-graph}
For a finite witnessed channel history $H$, let $\mathsf G(H)$ have
its retained operation occurrences as vertices. Its generating edges
are: consecutive local worker or bank uses; matched call--return
and enqueue--dequeue edges; selected operand and handle-production
edges; and cell production/read and one-use output dependencies.
The boundary label is the actual public arena event. This is the
finite graph obtained by elaborating the component terms, before
transitive closure.

Mark every publicly shared bank access, FIFO operation,
selection-cell operation and public event as a retained vertex.
Keep other shared-cell incidences until the explicit input-read
elimination of Lemma~\ref{lem:channel-block-reachability}; the remaining
private notification and outgoing-cell paths have their actual
request/input entry dependencies and their single result exit.
Between these retained operations,
the sender's or receiver's sequential pure program, together with
its exclusively used private-bank calls, is a \emph{private block}.
Split at every external resource access. A block's record includes
its entry operands and their suppliers, its actual branch, and its
exits. Its internal calls use the bank's immutable initialized entries
or results produced within the block. A history may stop inside a
block. Define $\mathsf C(H)$ by replacing each completed block by
a summary vertex with its actual incoming and outgoing incidences.
For an incomplete block retain its entry and no new exit incidence.
When computing the causal past of public events, its unfinished
tail is absent.
\end{definition}

\begin{lemma}[Boundary incidence and contraction reachability]\label{lem:channel-block-reachability}
For these channel programs, contracting private blocks preserves and
reflects reachability between retained vertices. The additional
ideal read of the honest input in a stored selection contributes
no new order between retained visible events.
\end{lemma}
\begin{proof}\vmode
\leavevmode
\begin{enumerate}[label=(\roman*)]
\item \textbf{Private interval dependencies.}
Inspect the strict let/call programs for $\mathsf{Encode}_W$ and
$\mathsf{Check}_W$.
Each call follows the preceding reply.
Its operands
are its entry datum or packet, an initialized immutable key, or a
previous result.

\item \textbf{Exclusive bank access.}
The endpoint bank serves only that endpoint worker.
Consequently all external operand incidences enter before the first
use that needs them, all such operands are available at block entry,
and every exit follows that entry along the recorded sequential
program.
No other worker can read a proper private intermediate.

\item \textbf{Shared operations and failure branches.}
Splitting at shared operations gives the same property for each
maximal remaining block.
The actual failure branch has one failure
exit, and the success branch has its declared result exit.

\item \textbf{Reachability contraction.}
A path of $\mathsf G(H)$ contracts to a path of $\mathsf C(H)$.

\item \textbf{Reachability reflection.}
Conversely, whenever a contracted path traverses a summary vertex,
its incoming boundary source reaches the program entry and its
outgoing boundary target follows the selected program exit.
The
recorded finite call sequence connects them.
Replace each such
traversal by that sequence.
This expands the path and proves
reflection as well as preservation.

\item \textbf{Boundary-incidence requirement.}
Merely identifying arbitrary
subgraphs would not establish this reverse implication.

\item \textbf{Stored-instruction dependencies.}
A stored instruction is emitted only after the honest marker or
authenticated record has supplied its packet.
Honest input already
precedes that packet.
The selector's extra read has edges from that
same input and from the selected instruction, and its exit leads to
the selection.
The existing path from input through packet and
instruction to selection already supplies the former dependency.

\item \textbf{Removal of redundant reads.}
Removing the read and retaining this dependency changes no
reachability between the remaining vertices.

\item \textbf{Direct instructions and initialized keys.}
A direct
$\mathsf{put}(d')$ instruction performs no such read.
Initialized key
suppliers precede runtime on both sides.

\item \textbf{Unfinished private tails.}
Finally an unfinished
private tail has no outside consumer, so deleting it from the
causal past of the selected public events changes no such path.
\end{enumerate}
\end{proof}

\begin{lemma}[Exhaustive channel successor comparison]\label{lem:channel-successor-table}
The residual comparison is preserved by every retained operation and
has legal finite expansions in both directions. Its nontrivial
classifier cases are the following disjoint cases on accessible
values:
\begin{center}\small
\begin{tabular}{@{}p{.40\linewidth}p{.24\linewidth}p{.25\linewidth}@{}}
\toprule Case & Real check & Simulated instruction\\\midrule
Honest encryption marker & $d$ & $\mathsf{stored}$\\
Nonmarker public-key encryption under the registered key;
record parses as $d'$ & $d'$ & $\mathsf{put}(d')$\\
Same encryption root; record parse fails & $\bot$ & $\mathsf{reject}$\\
Wrong encryption key or wrong root & $\bot$ & $\mathsf{reject}$\\
Signature/MAC with successful protected check and record parse
& $d$ by origin & $\mathsf{stored}$\\
Signature/MAC check or parse fails & $\bot$ & $\mathsf{reject}$\\
Symmetric nonmarker accessible value & $\bot$ & $\mathsf{reject}$\\
\bottomrule
\end{tabular}
\end{center}
\end{lemma}
\begin{proof}\vmode
\leavevmode
\begin{enumerate}[label=(\roman*)]
\item \textbf{Receiver-classification equations.}
The marker case is its defining expansion.
In the public-key
constructor case the stored operands are known and the destructor
equation returns exactly the record tested by $\mathsf{parse}$.
All wrong-key and wrong-root cases fail by the strict destructor
domains.

\item \textbf{Protected-origin cases.}
The protected-origin lemma gives the signature/MAC success
case.
An accessible symmetric constructor cannot use the protected
bare key, so the only successful protected symmetric ciphertext is
the marker.
The cases are exhaustive by the sorted root constructors;
invalid handles fail before entering this table.

\item \textbf{Store and receipt updates.}
The remaining successor calculations are:
\[
\begin{array}{c|l}
\text{operation}&\text{updated comparison field}\\ \hline
\text{request or receipt}&\text{same occurrence identifier and continuation}\\
\text{evaluator}&\rho^T=e_d\rho^{\widetilde T}\text{ and equal counters}\\
\text{submission enqueue/dequeue}&\text{same FIFO supplier; entries related by }e_d\\
\text{successful check/post}&C\text{ related to the instruction FIFO by }\operatorname{val}_d\\
\text{first selection}&\ell_{\rm r}=\ell_{\rm i}\in1+D\\
\text{later selection}&\text{filled cell unchanged}\\
\text{application}&\text{same request and selected-cell producer}\\
\text{failure}&\text{same public failure or}\\[-2pt]
&\text{same private discard continuation}
\end{array}
\]
These are equations on the residual-ledger projections of the actual
terms.
FIFO removal selects the corresponding head, and insertion
appends the corresponding entry with the same supplier.
At an empty
cell the classifier equation gives the same stored datum; at a filled
cell both updates retain it.
Receipts consume the same receipt token
and impose no edge on the distinct consumer.

\item \textbf{Evaluator preservation.}
Evaluator preservation
is Lemma~\ref{lem:comparison-command-induction}.

\item \textbf{Construction of opposite histories.}
Induct along any finite linearization of the retained graph.
The
table supplies the enabled opposite operation, and a private block
expands by the finite strict program on its selected branch.
Its
operands and queue heads are supplied by the invariant just checked.
Thus induction constructs a legal opposite history, not merely an
equal endpoint.

\item \textbf{Independent exchange and converse.}
Swapping incomparable steps preserves the same
selected records by the resource-chart equations.
Reversing the
table argument constructs the opposite-direction expansion.

\item \textbf{Pending requests.}
Pending
requests need no reply; stopping inside a private block contributes
only its already matched public causal past.
\end{enumerate}
\end{proof}

\begin{lemma}[Private intervals and visible orders]\label{lem:private-intervals}
In the two worker networks just defined, replace each finite private pure checking or encoding interval by a single internal vertex with its entry and exit dependencies. Keep every externally shared bank event, queue operation, selection-cell operation, and public-boundary event as a separate vertex. Private-bank events used exclusively inside one interval are contracted with that interval. This replacement preserves the set of finite visible labelled partial orders. The real and simulated networks have the same such orders.
\end{lemma}
\begin{proof}\vmode
\leavevmode
\begin{enumerate}[label=(\roman*)]
\item \textbf{Private-block incidence.}
Use the graph and block maps of Definition~\ref{def:channel-contraction-graph}.
Lemma~\ref{lem:channel-block-reachability} proves their boundary-incidence property and both directions of reachability preservation: encoding uses the sender private bank and checking the receiver private bank, and neither serves another worker.
Their immutable key entries were initialized before runtime.
The ideal notification interval and private selection call have the same specified input and output connections.

\item \textbf{Path contraction.}
Replacing a directed finite path by its composite edge preserves reachability between all remaining vertices: a path through that interval enters at its entry and leaves at an exit, and conversely its contracted edge expands to that finite path.
At a branch the actual recorded outcome fixes the chosen finite path.

\item \textbf{Prefixes and expansion.}
A proper interval prefix without an exit has no new outside consequence and can be omitted when taking the causal past of visible events.
Conversely every contracted finite history expands by the corresponding finite program, whose branch correctness is the preceding process comparison.
Thus contraction and expansion give both inclusions of the visible finite orders.

\item \textbf{Correspondence of visible vertices.}
For the remaining vertices define the correspondence by worker name, its local ordinal, and the selected request/packet/operand identifiers.
Lemma~\ref{lem:channel-successor-table} gives a bijection on these retained records (with ordinals counting retained operations) and preserves their enabling data, including unsuccessful operations.

\item \textbf{Preservation of causal generators.}
It preserves four generating kinds of edges: local worker and bank chains; matched request/reply and enqueue/dequeue edges; selected operand/handle-production edges; and one-use cell and application-request edges.
These are precisely the history generators of the declared local machines and queues.
No edge is inserted from an ingress receipt to receiver processing.

\item \textbf{Contracted checks.}
Private checks have been contracted, so \eqref{eq:classification-square} gives the same receiver outcome at their exits.

\item \textbf{Stored and direct instructions.}
A stored instruction in the simulator can occur only after its marker or valid signed packet has become available; hence its additional ideal read of the honest input is already in that packet's causal past.
It adds no new visible ordering.
A direct public-key injection uses no such read.

\item \textbf{Visible-order isomorphism.}
Consequently the generating directed graphs are isomorphic after these contractions; their transitive closures and their restrictions to the displayed boundary are isomorphic.

\item \textbf{Legality of reverse expansion.}
Legality in the reverse direction is not inferred from equal endpoints: for every retained record the process comparison supplies the corresponding enabled operation with the same selected resource, and the interval calculation supplies its finite internal witnesses.
Linearize the finite graph and perform those operations.
The bank states and queue head equations are maintained by the state relation at every step.
This gives a legal opposite history; incomparability is retained by the same resource-chart tests.

\item \textbf{Visible-past restriction.}
Finally take the causal past of the selected visible events and hide the remaining internal vertices, as in Section~\ref{subsec:pc-secured-proof}.
This proves both inclusions, including prefixes with unperformed replies or private work pending.
\end{enumerate}
\end{proof}

Using the displayed-order carrier of Definition~\ref{def:crypto-causal-certificate}, the lemma constructs the equality
\begin{equation}\label{eq:channel-behaviour-certificate}
\mathcal B_A((f_H\otimes1)r)=\mathcal B_A((1\otimes q_W)s).
\end{equation}
\begin{definition}[Matched channel histories]\label{def:matched-channel-histories}
For the just-constructed real and ideal-with-simulator realizations, the \textbf{matched-history set} $\mathsf{MatchedHist}_W$ consists of the paired finite expansions $(h_{\rm r},h_{\rm i},\theta)$ constructed in the preceding proof, with their retained-record bijection $\theta$. Its projections have types $\pi_{\rm r}:\mathsf{MatchedHist}_W\to\Conf_f(S_{\rm real})$ and $\pi_{\rm i}:\mathsf{MatchedHist}_W\to\Conf_f(S_{\rm ideal+sim})$ and return the indicated expansion. The preceding lemma proves that these maps cover both sets of displayed finite orders.
\end{definition}

The resulting span is
\[
\begin{tikzcd}[column sep=large]
\Conf_f(S_{\rm real}) &\mathsf{MatchedHist}_W
 \arrow[l,"{\pi_{\rm r}}"'] \arrow[r,"{\pi_{\rm i}}"]
 &\Conf_f(S_{\rm ideal+sim}) .
\end{tikzcd}
\]
Different internal witnesses need not be identified as strategy events. Equation~\eqref{eq:channel-behaviour-certificate} is the exact sufficient criterion of Lemma~\ref{lem:crypto-causal-congruence}.

\paragraph{Correctness converters.}
\begin{definition}[Honest channel correctness converters]\label{def:channel-correctness-converters}
For authentication the honest adversarial converter requests the outgoing packet, parses its public record for the notification, and uses the same retained packet for stored delivery instructions. For confidentiality it exposes only $*$ on notification. A $\mathsf{put}(d)$ instruction uses its own result bank to form $\mathsf{Encode}_{\mathsf{pk}}(d)$ with fresh local randomness and the initialized public key. Symmetric confidentiality has only stored delivery. In each case a converter ingress immediately queues the target instruction and makes its receipt available; a distinct forwarding worker processes these instructions in order. A stored instruction waits for the outgoing packet, except that a known already-used selection token makes all subsequent instructions inert. A direct instruction constructs its ciphertext without waiting for that packet. Forwarding submits the selected source handle and consumes the converter's one-use planning token after the first valid forwarding. The source's receiver and application worker retain the actual selection and output.
\end{definition}

\begin{lemma}[Correctness-converter record comparison]\label{lem:channel-correctness-records}
The honest converters have a two-sided finite causal comparison with
their ideal channel, obtained from the following record correspondence:
\[
\begin{array}{c|l}
\text{ideal record}&\text{converter and source records}\\ \hline
\text{notification}&\text{outgoing packet request and its checked notification}\\
(i,\mathsf{stored})&\text{forward the retained honest packet at instruction }i\\
(i,\mathsf{put}(d'))&\text{construct and forward }\mathsf{Encode}_{\mathsf{pk}}(d')\\
\text{receipt }i&\text{the same independent ingress receipt}\\
\text{first selection}&\text{consume the planning token and select the decoded datum}\\
\text{later instruction}&\text{unchanged selected cell; no second delivery}\\
\text{application reply}&\text{same request and first selected datum}
\end{array}
\]
\end{lemma}
\begin{proof}\vmode
\leavevmode
\begin{enumerate}[label=(\roman*)]
\item \textbf{Stored-branch computation.}
On the stored branch the operation equations give
\[
 \mathsf{Check}_W(\mathsf{Encode}_W(d))=d.
\]
On a direct branch the same equation gives $d'$ and uses only
the registered public key, fresh local nonce and supplied $d'$.
Thus this branch has no honest-input predecessor.

\item \textbf{Separation of result banks.}
Its result bank
is separate from the outgoing receipt bank, so no such predecessor
is introduced by shared bank control.

\item \textbf{Ingress and forwarding.}
Ingress appends the instruction and creates its independent receipt
token.
The forwarding worker consumes the queued instructions in
order.

\item \textbf{One-use planning.}
Its one-use planning token chooses at most the first valid
instruction.

\item \textbf{Store readiness.}
A first stored instruction waits for the outgoing
packet; subsequent inert instructions do not acquire a fresh
input dependency.
These are exactly the empty/filled-cell branches
of~\eqref{eq:ideal-select}.

\item \textbf{Application worker.}
The application worker retains its own
request predecessor and the selected datum's producer.

\item \textbf{Inductive history pairing.}
Construct the history pairing by induction on these retained records.
For every ideal record the table gives a finite constructor/forwarding
expansion with the same selected packet, queue and receipt witnesses.
For every converter record it gives its ideal contraction; private
proper prefixes have no new public consequence.

\item \textbf{Private-block contraction.}
The private programs
use exclusive result banks and have only their indicated packet or
notification exits, so the path replacement proof of
Lemma~\ref{lem:channel-block-reachability} applies to these programs
as well.

\item \textbf{Independent receipts and deliveries.}
Independent ingress receipts and application requests remain
independent.
Hence the two expansions preserve the same visible
partial order, including histories ending with a pending request.

\item \textbf{Surjective causal certificate.}
The set of all such paired finite expansions, with their two
configuration projections, is the certificate of
Definition~\ref{def:crypto-causal-certificate}.
\end{enumerate}
\end{proof}

\begin{proposition}[The common channel comparison]\label{prop:crypto-channel-pasting}
For each displayed wrapper $W$, the constructed public expansion, receiver classifier, and matched-history span give the simulation comparison $(f_H,q_W)$ in the contextual Scott quotient, as in Theorem~\ref{thm:correctness-security-mutual-factorization}. Its value and control calculations are the two squares
\[
 e_d\widetilde f=f e_d\quad\text{for public operations},
 \qquad \mathsf{Check}_W e_d=\operatorname{val}_d\chi_W.
\]
The first equation includes the equality of partial domains and observable results. The actual worker graphs lift these squares to the displayed game boundary, and the correctness converters give the separate correct resource conversions.
\end{proposition}
\begin{proof}\vmode
\leavevmode
\begin{enumerate}[label=(\roman*)]
\item \textbf{Public evaluator comparison.}
Lemma~\ref{lem:crypto-opaque} and private-key renaming give the
public comparison presentations.

\item \textbf{Exhaustive classifier.}
The exhaustive classifier cases
give~\eqref{eq:classification-square}.

\item \textbf{Residual history pairing.}
Construction~\ref{con:residual-channel-comparison} lifts these maps
to the actual bank, queue and cell records.

\item \textbf{Visible causal comparison.}
Lemma~\ref{lem:private-intervals} gives the two-sided causal
certificate~\eqref{eq:channel-behaviour-certificate}, retaining
pending requests and selected suppliers.

\item \textbf{Correctness and realization.}
Lemma~\ref{lem:channel-correctness-records} separately gives the
correctness certificate.
Theorem~\ref{thm:categorical-crypto-bridge}
turns these into the simulator and correctness equations.
\end{enumerate}
\end{proof}

\begin{theorem}[Symbolic channel realizations]\label{thm:crypto-channels}
Against arbitrary strategies on the exposed adversarial network/evaluation interface, the signature and MAC wrappers realize $\mathcal F_{\rm auth}$. The public-key wrapper realizes $\mathcal F_{\rm conf}$. The specified free-constructor symmetric wrapper realizes $\mathcal F_{\rm secret\text{-}auth}$. Each realization has a simulator satisfying \eqref{eq:simulator-factorization} at the visible game boundary.
\end{theorem}
\begin{proof}\vmode
\leavevmode
\begin{enumerate}[label=(\roman*)]
\item \textbf{Channel instances.}
Apply Proposition~\ref{prop:crypto-channel-pasting}.

\item \textbf{Authentication classifiers.}
Protected-key origin makes the signature and MAC classifiers select
only the stored record.

\item \textbf{Public-key injection.}
Public-key encryption additionally admits
$\mathsf{put}(d)$, giving confidentiality; the symmetric classifier
admits only its honest marker, giving secret authentication.

\item \textbf{Executable simulation.}
These classifiers are executable without the encryption plaintext.

\item \textbf{Attacker contexts and concurrency.}
Postcomposition with any exposed attacker $a$ uses the single
witness $aq_W$, by Theorem~\ref{thm:correctness-security-mutual-factorization}.
The causal certificates retain the independent banks, FIFO suppliers,
receipts and application requests.
\end{enumerate}
\end{proof}
\begin{proposition}[A secure channel for positive observations]\label{prop:secure-open-channel}
Fix a Scott open $U$ of an algebraic payload domain $D$. A one-shot instance of the secret-authenticated channel carries a positive observation of membership in $U$. It has the simulator of Theorem~\ref{thm:crypto-channels} at the adversarial interface. The predicate can itself be a supplied payload $V\in\Op(D)$, in which case the reported observation is $x\in V$.
\end{proposition}
\begin{proof}
At the honest sender attach the observation strategy for the characteristic map $\chi_U:D\to\mathbb S$, where $\mathbb S=\{0<1\}$. This map is continuous since the only nontrivial open of $\mathbb S$ is $\{1\}$. Request evidence for $1$ once. On its positive receipt send a fixed public literal through the original wrapper. At the honest receiver turn a delivery of that literal into a single positive receipt. The sender and receiver adapters have separate local state. The sender obtains evidence exactly when a compact $k\le x$ lies in $U$, by Proposition~\ref{prop:continuous-query-strategy}. The ideal channel can deliver the literal after such a send and a valid forwarding instruction; every receipt requires that send. If the adversary withholds delivery the receipt can remain pending. For a supplied predicate use the continuous characteristic map of the membership open in Proposition~\ref{prop:open-payload-basis} at the honest sender. These are ordinary honest contexts for the real channel and its ideal counterpart. Applying the same contexts to its simulator equation proves the claim by Corollary~\ref{cor:honest-readout}; the simulator receives only the original ideal channel's notifications, including their timing.

For a finite family of predicates use distinct local query lanes and one channel instance per predicate. Parallel composition of the simulator equations covers joint attackers. Duplicate successful queries then have distinct receipt occurrences. A delivery observation on the original signature, MAC or public-key wrapper can likewise be made by an honest adapter at its existing application port. The stored-packet and injected-packet branches, where offered, retain the selection and public-command rules of the channel under consideration.
\end{proof}

\begin{corollary}[Transmitting an open at a query interface]\label{cor:transmit-observation}
A sender holding $V\in\Op(D)$ can transmit it at the receiver interface $\mathsf Q_{\Op(D)}(I)$ using an indexed family of the preceding channels. With honest forwarding the finite receiver traces are exactly those of $P_V^I$. At the joint adversarial interface a single simulator realizes the family of ideal channels with the same honest adapters.
\end{corollary}
\begin{proof}
For each receiver lane $a=(i,U_F)$ use a separate channel for the basic test $U_F\subseteq V$ of Proposition~\ref{prop:open-payload-basis}. The sender takes the evidence from its own lane on the supplied predicate and sends the marker when that test succeeds. The receiver retains delivery at that lane and acknowledges exactly after both its local request and the delivery. This last adapter accepts the request independently of the delivery. With honest forwarding, any visible acknowledgement entails $U_F\subseteq V$. Conversely every finite serialization of $P_V^I$ lifts by forwarding the markers for its successful lanes before their acknowledgements. Other requests may remain pending. This proves the trace assertion including the order of repeated requests.

For an arbitrary index set take disjoint, lane-tagged copies of the channel displays. Order is componentwise and finite consistency means consistency in each component. Principal pasts stay finite, each enabled negative event lifts in its own lane, and courtesy holds componentwise. Use the same construction on copies of the specified simulator. A finite serialization involves finitely many components, where the finite tensor simulator equation applies. A matching serialization on this finite subfamily embeds back into the whole family with empty configurations elsewhere. The reverse inclusion is identical. Theorem~\ref{thm:open-trace-testing} gives contextual equivalence on the joint boundary; attaching the honest adapters preserves it. At that boundary the ideal notifications expose channel activity indexed by lane. In particular an observed notification for $a=(i,U_F)$ reveals the positive fact $U_F\subseteq V$. The claim compares the real execution to this scheduled ideal service, whose adversary can delay or suppress particular deliveries.
\end{proof}

\begin{corollary}[Channel refinement and behaviour images]\label{cor:channel-images}
For each channel realization and coalition, its correctness converter and simulator give opposite slice arrows over the honest interface. The real and ideal attainable attacker-view images are equal by Corollary~\ref{cor:security-images}.
\end{corollary}
\begin{proof}\vmode
\leavevmode
\begin{enumerate}[label=(\roman*)]
\item \textbf{Transposition of channel equations.}
Transpose the two equations of Proposition~\ref{prop:crypto-channel-pasting}
and apply Corollary~\ref{cor:security-images}.
\end{enumerate}
\end{proof}

The symmetric wrapper's integrity conclusion follows from its specified operations. The free constructor provides no ciphertext modification operation, and constructing another ciphertext under the same key requires that exact key. This conclusion therefore concerns this symbolic wrapper. It does not assert authentication for arbitrary symmetric encryption schemes. For public-key encryption, the successful injection operation explicitly shows why encryption alone does not authenticate the sender.

Theorem~\ref{thm:correctness-security-mutual-factorization} composes parallel copies with disjoint session resources. Replacing a packet store by a specified queue, loss or replay medium also transports that medium's selection and delivery controls to the ideal interface. The target must retain observable scheduling information, and the comparison supplies only the delivery and fairness guarantees of the chosen medium.

Shared long-term key services, adaptive replacement of honest programs and further protocol formats need their own simulators. Their transitions can be represented, but they are outside the conclusion of Theorem~\ref{thm:crypto-channels}.

\subsubsection{Local operations as process macros}\label{subsec:pc-applied}
Value constructors and strict success/failure branches are the rule instances of Definitions~\ref{def:value-formation-constructors} and~\ref{def:calculus-destructor-rules}. A handle bank is the state of the service process, not a second process syntax. Scoped allocation uses the existing binding and occurrence constructors. Finite public value computations use this same language; Proposition~\ref{prop:finite-public-computation-adequacy} identifies their values with the accessible reference entries. Applied-$\pi$ provides a related symbolic operational presentation~\cite{crypto-abf}; the present security relation is the contextual Scott relation of Definition~\ref{def:formal-contextual-boundary-equivalence}.

\begin{example}[An encrypted relay]\label{ex:pc-encrypted-relay}
At a private key interface, let $c=\crSEnc(m,k,r)$ with fresh local nonce $r$. A send on $a$, receive on $a$, and successful evaluation derive
\[
 a!c.0\mid
 a?(z).\mathsf{let}\ y=\crSDec(z,k)\ \mathsf{in}\ b!y.0
 \longrightarrow^* b!m.0.
\]
The receiver's key is supplied by its private service connection. With an unrelated key the failure continuation is taken. Under FIFO, the communication derivation has a separate enqueue and dequeue; its local value calculation is the same. In the public handle presentation, output exports a ciphertext handle and evaluation uses the receiver's authorized key handle. No constructor exposes the key to the network observer.
\end{example}

\subsubsection{Process interpretation of the wrapper comparison}\label{subsec:pc-crypto-wrappers}
Definition~\ref{def:real-channel-processes} gives the actual wrapper terms. Their service macros expand through Definition~\ref{def:service-call-formation-elimination}; Proposition~\ref{prop:pc-named-elaboration} supplies their witnessed graph and cells. The classifier and marker-bank programs replace the corresponding private continuations in the simulator. Proposition~\ref{prop:generic-process-preservation} transports the comparison of Theorem~\ref{thm:crypto-channels} through these terms, with the same public interface and finite causal certificate.

\subsubsection{The strict decoder as a service process}
\begin{proposition}[Decoder computation]\label{prop:decoder-computation}
For an available runtime key $k$ and ciphertext $c$, the process
$\mathsf{let}\ x=d(k,c)\ \mathsf{in}\ P\ \mathsf{else}\ Q$
executes the displayed destructor. A successful return $m$ installs $P[m/x]$; failure installs $Q$. On a matching encryption the returned value is its plaintext.
\end{proposition}
\begin{proof}\vmode
\leavevmode
\begin{enumerate}[label=(\roman*)]
\item \textbf{Lookup and branch evaluation.}
Typed lookup supplies the actual operands.
The evaluator chooses its defined pattern or failure case, and the existing substitution theorem installs the indicated continuation.

\item \textbf{Decoder equations.}
The decryption equation proves the encrypted case.

\item \textbf{Sequential decoding.}
Sequential decoders repeat this calculation at each checked command, including every failure branch.
No evidence assertion is used as a runtime key.
\end{enumerate}
\end{proof}

\subsubsection{Derived key encapsulation and authenticated encryption}\label{subsec:crypto-kem-aead}
The generated services also support composed encryption mechanisms. The following constructions specify their algorithms and interfaces, then check correctness and the stated origin and secrecy comparisons through the actual operations.

The following workers compose encryption, checked coercion, and MAC verification. All branches are actual service transitions; the receiver classifiers connect them to ideal delivery resources.

\begin{definition}[Key transport as an encapsulation service]\label{def:constructed-kem}
Given an encryption public key $p$, an encapsulation worker allocates a fresh symmetric key $s$ and nonce $r$ and returns the pair consisting of a local handle for $s$ and
\[
 c=\crAEnc(\iota_{\mathsf K_s}s,p,r).
\]
The sender transports only $c$ to the recipient. Decapsulation applies $\crADec(c,k)$ and the checked $\mathsf K_s$ coercion, returning failure on any other case. The worker has no command revealing another session's retained key.
\end{definition}

\begin{proposition}[Encapsulation equations and non-derivability]\label{prop:constructed-kem}
For $p=\crPK(k)$, decapsulation of the constructed ciphertext returns $s$. Among valid generated ciphertexts, equality determines the plaintext key, public key, and nonce. In the closed one-session resource exporting $p,c$ but neither $k$ nor $s$, the observer does not derive $s$.
\end{proposition}
\begin{proof}\vmode
\leavevmode
\begin{enumerate}[label=(\roman*)]
\item \textbf{Correctness and key uniqueness.}
The decryption equation followed by checked coercion gives $s$.
Equality of encryption trees identifies their three operands, proving the binding statement for valid generated ciphertexts.

\item \textbf{Private-key invariant.}
The private $k$ is never exported and there is no destructor for public keys.

\item \textbf{Session-value nondeducibility.}
Every observer occurrence of $s$ starts beneath the honest ciphertext.
Opening it requires $k$; another constructor or a failed destructor cannot expose $s$.
The same invariant proof as Lemma~\ref{lem:crypto-opaque} therefore excludes a bare $s$.

\item \textbf{Handle-level receipt.}
Receipt of $c$ gives only its selected root handle, not its children.
\end{enumerate}
\end{proof}
This is an explicitly rejecting symbolic KEM: malformed ciphertexts fail. Implicit-rejection KEMs, which return a derived key on malformed ciphertexts, have a different transition. The distinctions among such interfaces and their binding properties are studied in~\cite{crypto-kems}; the proposition makes no claim about the post-quantum security of a bitstring implementation.

\begin{definition}[Associated-data packet and receiver]\label{def:associated-data-packet}
The parameterized packet operation at keys $k_e\in T_{\mathsf K_s}$ and
$k_m\in T_{\mathsf K_m}$ takes $(m,a,r)\in T_{\mathsf M}^2\times T_{\mathsf N}$
to
\[
 c=\crSEnc(m,k_e,r),\qquad
 t=\crMac(\langle a,c\rangle,k_m),\qquad
 p=\langle c,t\rangle.
\]
The receiver, given $a$ and $p$, projects $c,t$, checks
$\mathsf{mcheck}(t,\langle a,c\rangle,k_m)=\crOK$, and only on success
evaluates $\crSDec(c,k_e)$. Failed lookup, projection, checking, or decryption
returns failure without returning an intermediate plaintext.

For its \textbf{initialized closed service}, allocate $k_e,k_m$
in separate key-sort namespaces and connect their banks only to the
packet workers and receiver. A packet worker receives $(m,a)$,
allocates a nonce in its own increasing-counter namespace, executes
the displayed program once, and exports only $p$ under that request's
packet identifier. It retains the request, ciphertext, tag and tagging
occurrence. Different workers have disjoint allocation namespaces.
The receiver's successful plaintext goes only to its application port.
Public gateways use the reference commands of
Construction~\ref{con:cryptographic-reference-processes}; neither
keys, nonce registers nor raw-term import are public commands.
Transport uses actual matched-wire imports. The finite worker programs,
banks and private connections are interpreted by
Theorem~\ref{thm:compact-cryptographic-communication}, as for the channel
workers. Allocation follows its request, so no future local nonce is
a literal supplied by that request.

The \textbf{pre-decryption public frame} consists of the public initial
entries, the issued packets, their identifiers, any publicly retained
associated data, and values obtained by public evaluator commands and the
observer's own allocations. Plaintexts returned at the receiver's application
port are separate observations. In a comparison of two plaintext experiments,
the initial public entries, associated data, packet identifiers and public
literal alphabet are retained. This is the same frame observation of
Construction~\ref{con:accessible-labelled-algebra}.
\end{definition}

\begin{proposition}[Associated-data binding and confidentiality]\label{prop:constructed-aead}
The initialized service returns $m$ on an honest packet. Every accepted pair
$(a,c)$ has an honest tagging occurrence of exactly that pair in its supplying
past. For any finite family of issued honest packets, changing their
plaintexts while retaining their associated data and packet identifiers
preserves the pre-decryption public frame observations. Corresponding public
evaluator commands, including the observer's fresh allocations, preserve this
comparison through every finite interaction. No distribution on allocated
names is used.
\end{proposition}
\begin{proof}\vmode
\leavevmode
\begin{enumerate}[label=(\roman*)]
\item \textbf{Correctness.}
Correctness is the composite of the derived MAC check and decoder.

\item \textbf{Common marker carrier.}
For the privacy comparison retain the finite issued-packet set $J$ and write
\[
 c_j^b=\crSEnc(m_j^b,k_e,r_j),\qquad
 t_j^b=\crMac(\langle a_j,c_j^b\rangle,k_m),\quad b=0,1.
\]
The two experiments use corresponding allocation occurrences.
Let $B$ be their common stock of external atoms, disjoint from the service's key and nonce allocators.
On the free external value carrier adjoin message markers $\kappa_j,\theta_j$ and construct expansion by its unique fold:
\[
 e_b(\kappa_j)=c_j^b,\quad e_b(\theta_j)=t_j^b,\quad e_b|_B=1.
\]
At markers all public root eliminations fail; equality uses their distinct tags and packet identifiers.

\item \textbf{Expansion operations.}
Ordinary roots have the already derived image eliminators.
Here are the instance-specific checks for this public algebra map.

\item \textbf{Root separation.}
Distinct ciphertext expansions have distinct fresh $r_j$; equality of tag expansions forces equality of their ciphertext operands.
Their root tags distinguish ciphertexts from MACs.

\item \textbf{Injectivity of expansion.}
An external nonmarker term cannot equal either expansion, since constructing its root would require $k_e$ or $k_m$, whereas its key operands lie in $B$.
Root injectivity then proves injectivity of $e_b$ by induction.

\item \textbf{Marker eliminators.}
At a marker, projection, coercion and data reading fail by its root; a successful symmetric decryption or MAC check would require the unavailable corresponding key.
Other cryptographic root checks fail by their tags.

\item \textbf{Nonmarker operations.}
At nonmarker roots, image membership and returned operands commute with expansion by its injectivity.
Thus every exposed partial operation preserves and reflects its domain, and equality and observable results agree.
This includes a free hash introduction when present.

\item \textbf{Accessible-image invariant.}
The image of $e_b$ contains the external initial values and is closed under all exposed operations.
Issuance extends the source by the two new markers and publishes $\langle\kappa_j,\theta_j\rangle$; an external allocation extends $B$ by its new occurrence tag.
Copy and wire transport retain their sources.

\item \textbf{Audit induction and nondeducibility.}
Induction on the reference-process audit therefore puts every accessible value in this image, while none of $k_e,k_m,r_j$ belongs to its key/nonce components.
This proves the required non-derivability from the actual port connection.

\item \textbf{Plaintext interface.}
Receiver plaintext outputs are outside the pre-decryption observation being compared.

\item \textbf{Common associated-data values.}
Fixed associated-data values have common normal forms, even when they contain earlier packets.
Indeed $e_0(u)=e_1(v)$ forces the same marker at a protected-key root: ciphertext equality fixes $r_j$, and MAC equality fixes its ciphertext and hence $j$.
Such a root cannot match an external nonmarker root.

\item \textbf{Recursive comparison of roots.}
All remaining comparisons descend to operands, proving $u=v$.
This also permits later plaintexts containing earlier packets without a cycle: an issuance only uses earlier available operands.

\item \textbf{Labelled public-frame isomorphism.}
Let $\nu$ be this common labelled frame and $N_J=\mathsf{DY}(\im\nu)$ in the marker algebra.
Expansion commutes with finite public computations; its image is closed and contains the real frame.
Conversely every element of that image is the expansion of a computation on $\nu$.
Hence $e_b$ restricts to an isomorphism $N_J\cong A_{\phi_b}$.

\item \textbf{Public comparison and extension.}
Proposition~\ref{prop:crypto-public-comparison} gives static equivalence; Proposition~\ref{prop:generic-evaluator-comparison} lifts it to the reference processes with the same domains, counters and replies.
At issuance and external allocation the explicit extensions above give the successor comparison, so the result includes adaptive finite commands.

\item \textbf{Receiver authentication classifier.}
For authentication normalize a submitted packet.
A tag passing the receiver's protected-key check must be $\theta_j$, since no external MAC has key $k_m$.
Injectivity forces its associated data and ciphertext to be $\bar a_j$ and $\kappa_j$.
Thus the receiver classifier is the finite lookup
\[
 \mathsf{Accept}(u,v,w)\ \Longleftrightarrow\
       \exists j\in J:\ (u,v,w)=(\bar a_j,\kappa_j,\theta_j).
\]
The decoder then returns $m_j^b$.

\item \textbf{Origin and replay.}
The selected tag's supplying audit reaches the actual tagging occurrence by Corollary~\ref{cor:crypto-origin}, since no external operation has the protected key.
Replay may reuse that origin.
This proves the claimed authentication and completes the comparison at the stated observation boundary.
\end{enumerate}
\end{proof}

\begin{construction}[Delivery resource and packet simulator]\vmode\label{con:aead-ideal}
\leavevmode
\begin{enumerate}[label=(\roman*)]
\item \textbf{Ideal records and notifications.}
The ideal associated-data service stores an application request $(m,a)$ under its freshly issued packet identifier $j$.
It publishes the same request notification and associated-data fields as the packet interface.

\item \textbf{Adversarial commands.}
Its adversarial commands are $\mathsf{deliver}(j)$ and $\mathsf{reject}$: a delivery returns the stored $m$ at the receiving application port, and rejection returns failure.

\item \textbf{Replay and failure.}
A stored record remains available after delivery, so this resource permits replay.
Missing identifiers fail.

\item \textbf{Common interface.}
The notification, delivery, and reply ports have the same order and separate worker dependencies as the concrete network; the ideal worker is the finite register program that implements these instructions.

\item \textbf{Simulator marker bank.}
The simulator maintains the marker banks $N_J$ from Proposition~\ref{prop:constructed-aead}.
On notification of $j$, it issues the packet handle containing $\langle\kappa_j,\theta_j\rangle$ in its private representation.

\item \textbf{Public evaluation.}
Public evaluator calls use the normal-form evaluator, preserving handle identifiers and replies.

\item \textbf{Packet classification.}
A submitted packet is classified by the displayed predicate $\mathsf{Accept}$: a successful classification sends $\mathsf{deliver}(j)$, and every other submission sends $\mathsf{reject}$.

\item \textbf{Available information.}
The simulator needs no plaintext to perform these steps.

\item \textbf{Private representation and strategies.}
Its representation is internal; the outside interface still offers precisely handles, commands, and replies.
These finite programs define strategies by Lemma~\ref{lem:crypto-stateful}.
\end{enumerate}
\end{construction}

\begin{proposition}[Associated-data security as image inclusion]\label{prop:aead-image}
With the application ports kept open, the concrete packet network and the simulated ideal delivery service have the same contextual Scott observations. The constructed packet simulator therefore gives a slice arrow from the real strategy to the ideal strategy, and the attainable-view-image inclusion of Corollary~\ref{cor:security-images}.
\end{proposition}
\begin{proof}\vmode
\leavevmode
\begin{enumerate}[label=(\roman*)]
\item \textbf{Public-bank comparison.}
Relate the public banks by the expansion of
Proposition~\ref{prop:constructed-aead}, retaining identifiers,
associated data, counters and application-request records.
Its
operation squares preserve every public reply and failure.

\item \textbf{Receiver classification.}
Its
classifier sends an accepted packet exactly to its issued identifier
$j$, whose ideal record returns the same $m_j$; rejection and replay
use the same branches.

\item \textbf{Fresh requests.}
New requests extend the marker bank after
the same input occurrence.

\item \textbf{Causal certificate.}
Matching these record updates in both
directions retains their selected suppliers and worker dependencies;
independent ports and banks retain their orders.
This constructs the
finite causal certificate, so Theorem~\ref{thm:categorical-crypto-bridge}
gives the simulator equation and Corollary~\ref{cor:security-images}
its image inclusion.

\item \textbf{Correctness wrapper.}
Honest forwarding is the successful-delivery
case and supplies the application correctness equation.
\end{enumerate}
\end{proof}
This ideal resource exposes the actual associated-data and scheduling fields and permits the receiver's actual replay behaviour. An additional replay filter would be a different service program, with its own state transition and comparison.

\subsubsection{Keyed functions and an exact opaque-oracle comparison}\label{subsec:crypto-keyed}
A fresh-key service can be compared with a memoized opaque oracle by tracking its public replies. The comparison below is exact in the symbolic operation model and retains the allocation and repeated-query behaviour needed by the reference service.

Derived keys are values in the explicitly enlarged key carrier. Its exposed commands determine their available uses.

The command interface here illustrates the role of discrete query identity. A repeated query selects a previously issued entry. Singleton relabelling preserves exactly that lookup by Proposition~\ref{prop:discrete-observation-recovery}. An observation of an authorized reply acts after this lookup. Comparing two arbitrary received predicates for extensional equality would instead require a new operation; the equality calculation in Section~\ref{subsec:operation-descent} explains the distinction. Encapsulation, authenticated decryption and tag verification likewise retain the specified exact-key and failure branches of their evaluators.

\begin{definition}[Keyed-function service]\label{def:formal-keyed-function-service}
Use a key carrier freely generated by key atoms and a constructor $F:K\times M\to K$, with no component projection. Original symmetric encryption accepts this carrier as its key argument and compares complete key values. Constructor recursion gives the extended finite-tree carrier. A private worker allocates $k$, retains it, and offers the public query $x\mapsto F(k,x)$ through handles. Its other public operations are the generated value operations; its private key is not exported or used by another hidden worker.
\end{definition}

\begin{theorem}[Symbolic keyed function versus memoized opaque names]\label{thm:symbolic-keyed-oracle}
The hidden-key worker has the same public handle observations as a memoizing worker which stores a finite map from query values to fresh opaque key names, returning the stored value on repeated queries and allocating one on a new query. The comparison permits adaptive finite queries and subsequent public uses of returned keys. It is an exact symbolic comparison, not a PRF security theorem.
\end{theorem}
\begin{proof}\vmode
\leavevmode
\begin{enumerate}[label=(\roman*)]
\item \textbf{Query-ordered expansion.}
After a finite history let $x_1,\ldots,x_j$ be the distinct queried values in the memoizing world, with results $u_1,\ldots,u_j$.
Define expansion recursively in query order by
\[
 e(u_i)=F(k,e(x_i)),
\]
fixing all public source values and extending through public constructors.

\item \textbf{Previously available operands.}
A query argument contains only values already available, so this recursion has no cycle.

\item \textbf{Private-key invariant.}
The private $k$ is not derivable: it starts only as the first operand of a hidden-key $F$ result, and no public destructor extracts that operand.

\item \textbf{Fresh-root separation.}
A newly returned $F(k,e(x_i))$ cannot be manufactured by applying public $F$ with an available first operand.

\item \textbf{Equality of special roots.}
Equality between two such special roots reduces to equality of their inputs; the memoizer identified exactly those inputs.

\item \textbf{Recursive operation preservation.}
Other constructor roots descend to operand comparisons.
Induction on the query history and on compared finite values therefore proves injectivity of $e$ on generated values, and preservation and reflection of every public destructor.
In particular symmetric decryption compares entire key values, so its key matches are preserved.

\item \textbf{Repeated and fresh queries.}
A repeated query returns corresponding old values.
On a new query extend $e$ by the displayed equation; injectivity follows from its input being different from every previous input.

\item \textbf{Strict reference comparison.}
The explicit old-query and new-query updates extend $e$ at each step.
Proposition~\ref{prop:generic-evaluator-comparison} lifts this evolving operation comparison to the reference banks, issuing the same fresh public handles even for repeated values.
Its comparison carrier retains the reached query table and expansion; Theorem~\ref{thm:service-behaviour-canonical-comparison} gives equality of the public call--return strategies, including adaptive queries.
\end{enumerate}
\end{proof}

The retained query table and expansion make $M$ an actual comparison transducer. Its two projections are strict simulations, and further comparisons compose by the local pullback of Proposition~\ref{prop:service-pullbacks}. The finite expansion recursion also supplies their effective maps on the coded comparison states; no decision procedure for arbitrary behavioural equivalence is used.

\[
\begin{tikzcd}[column sep=large]
\text{opaque table at }h\arrow[r,"\text{next query}"]\arrow[d,"e_h"']
 &\text{opaque table at }hq\arrow[d,"e_{hq}"]\\
\text{hidden-}F\text{ values at }h\arrow[r,"\text{next query}"']
 &\text{hidden-}F\text{ values at }hq
\end{tikzcd}
\]
The query interface here does not disclose $k$. A later key-disclosure operation is a new interface case and is not covered by this comparison. \begin{definition}[Tagged key-derivation trees]\label{def:tagged-key-derivation}
The \textbf{derivation tags} $\mathsf{tag}_i=\mathsf{data}(\mathsf{inr}(i))$ for $i\in\{0,1\}$ are the two new public data leaves in the coproduct $D\amalg\{0,1\}$. The \textbf{binary key-derivation map} sends $k$ to $(F(k,\mathsf{tag}_0),F(k,\mathsf{tag}_1))$ in the twofold product of the keyed-constructor result sort. Iterating this map over a finite binary tree gives its \textbf{key-derivation tree}, retaining the actual dependency equations. The coproduct extension is the same conservative data extension used for public format tags below.
\end{definition} The more general interaction of encryption and pseudorandom derivation is studied in~\cite{crypto-micciancio}; its computational result is not inferred from constructor freeness.

\subsection{Actions, threshold symmetries, and linear computation}\label{subsec:actions-threshold}
For a monoid $A$, the monoid algebra action on its actual carrier $X$
is a functor $BA\to\Set$. Two implemented actions $u,v$ commute
exactly when they extend to an action of $A\times B$ by
$(a,b)\cdot x=u_a(v_b(x))$. Indeed the multiplication law follows by
commuting $v_b$ past $u_{a'}$; restriction to $(a,1)$ and $(1,b)$ is
the inverse construction. Thus the square
\[
\begin{tikzcd}
 X\arrow[r,"u_a"]\arrow[d,"v_b"']&X\arrow[d,"v_b"]\\
 X\arrow[r,"u_a"']&X
\end{tikzcd}
\]
calculates agreement: exchange $u_a(x_0)$ and $v_b(x_0)$ and apply
the other local action. The returned values coincide. The scalar
action in Section~\ref{subsec:crypto-algebraic} is an explicit instance.
Its accessible submodule calculation supplies a separate information
theorem; agreement does not supply authentication or secrecy by itself.

\subsubsection{Explicit scalar, group, and pairing models}\label{subsec:crypto-algebraic}
The scalar, group, and pairing carriers below have explicit finite evaluators. Their algebra maps prove correctness and their fixed-value invariants prove the stated information properties.

\begin{definition}[Scalar and pairing carriers]\label{def:formal-scalar-and-pairing-carriers}
For a finite allocated-name object $A$, the scalar polynomial algebra over a coefficient field $k$ is the free commutative $k$-algebra
\[
 k[A]=\operatorname{Sym}_k(k^{(A)}),\qquad
 \operatorname{Hom}_{k\text{-}\mathrm{Alg}}(k[A],C)
       \cong\operatorname{Hom}_{\Set}(A,U C).
\]
Its fraction field is the localization
$K_A=(k[A]\setminus\{0\})^{-1}k[A]$. The universal maps are
\[
\begin{tikzcd}[column sep=large]
k[A]\arrow[r,hook,"\iota_A"]\arrow[dr,"f"']
 &K_A\arrow[d,dashed,"\exists!\,\bar f"]\\
 &C ,
\end{tikzcd}
\qquad f(k[A]\setminus\{0\})\subseteq C^\times.
\]
Thus allocation is the free extension $k[A]\to k[A+1]$, followed by localization. Injections of allocated names induce embeddings and hence extend to the fraction fields. The unit condition displayed in the diagram specifies exactly when a more general algebra map extends; a specialization killing a denominator is not such a map.

For the computational interpretation take an effective coefficient field, presented by finite codes and terminating field operations and equality. Both $\mathbb Q$ and each prime field $\mathbb F_p$ have these presentations. Represent $k[A]$ by finitely supported coefficient maps on $\mathbb N^A$, with finite convolution as multiplication. A monomial order proves that a product of nonzero polynomials is nonzero. Represent $K_A$ by pairs $(p,q)$ with $q\ne0$, modulo $pq'=p'q$. Cross multiplication decides equality; inversion exchanges numerator and denominator on the effective open domain $p\ne0$. Theorem~\ref{thm:effective-partial-classifier} therefore classifies its effective open return domain.

The public scalar interface exposes coefficient constants and field operations. It does not expose coefficients of a polynomial, syntax, variable support, or a test for being a freshly allocated name. A coefficient-reading operation succeeds exactly on the image of $k\hookrightarrow K_A$. On $p/q$, scan the finitely many coefficients of $q$ to find a nonzero coefficient, divide the corresponding coefficient of $p$ by it to obtain $c$, and test $p=cq$. Return $c$ exactly on success. This is independent of the fraction representative. The rational-reading instance is recovered by $k=\mathbb Q$.

Let $G_A$ and $G_{T,A}$ be disjoint tagged copies of the additive group of $K_A$. Define
\[
 g(a)g(b)=g(a+b),\quad g(a)^{-1}=g(-a),\quad
 g(a)^s=g(as),\quad \mathsf g=g(1),
\]
and the same group operations on the target copy. Define the bilinear operation
\[
 e:G_A\times G_A\longrightarrow G_{T,A},\qquad e(g(a),g(b))=g_T(ab).
\]
The pairing is the bilinear map induced by multiplication of the scalar algebra:
\[
\begin{tikzcd}[column sep=large]
K_A\times K_A\arrow[r,"\otimes"]\arrow[dr,"m"']
 &K_A\otimes_k K_A\arrow[d,"\overline m"]\\
 &K_A .
\end{tikzcd}
\]
The private tags transport this map to the source and target group carriers. There is no operation $G_A\to K_A$ or $G_{T,A}\to G_A$. The tag $g$ describes the private representation; public group construction is exponentiation of the public generator by an available scalar.
The \textbf{scalar sort} is $T_{\mathsf{Scalar}}(A)=K_A$. Its transport envelope is the tagged constructor $\iota_{\mathsf{Scalar}}:K_A\to T_{\mathsf M}(A)$ of the extended transport algebra; checked coercion returns its scalar argument and fails at other tags. Scalar operations act on $K_A$, not on a public syntax decomposition of its elements.
\end{definition}

\begin{proposition}[Group and pairing equations]\label{prop:constructed-groups}
These carriers and operations give computable symbolic abelian groups, a scalar action, and a bilinear pairing. For example
\[
 (u^a)^b=u^{ab},\quad e(u^a,v^b)=e(u,v)^{ab},
 \quad e(uu',v)=e(u,v)e(u',v).
\]
For independently privately allocated scalars $a,b$, public source-group values $\mathsf g,g(a),g(b)$ do not yield $g(ab)$ using source-group operations and public scalars. Adding the displayed pairing does yield $g_T(ab)$, and distinguishes a candidate source-group Diffie--Hellman value by a pairing equality.
\end{proposition}
\begin{proof}\vmode
\leavevmode
\begin{enumerate}[label=(\roman*)]
\item \textbf{Algebraic equations.}
Every equation reduces to the field equations on the represented coefficients; for example the coefficients of $e(g(x)^a,g(y)^b)$ and $e(g(x),g(y))^{ab}$ are both $xyab$.

\item \textbf{Public coefficient invariant.}
Scalar values generated publicly belong to the field $L$ generated by public scalars and the observer's own allocations, since no group operation returns a scalar.
Source-group values therefore have coefficients in the $L$-linear span of $1,a,b$, by induction on products, inverses and scalar powers.

\item \textbf{Source-group nondeducibility.}
Algebraic independence of $a,b$ over $L$ implies $ab$ is not in this span: clearing denominators would equate a nonzero multiple of the monomial $ab$ with a polynomial containing no such monomial.
Constructors of other sorts do not add a scalar or source-group destructor, so they do not alter this argument.

\item \textbf{Pairing and public equality test.}
Pairing directly gives $e(g(a),g(b))=g_T(ab)$.
For a candidate $z=g(c)$, the test $e(z,\mathsf g)=e(g(a),g(b))$ holds exactly when $c=ab$.

\item \textbf{Distinction of conclusions.}
Thus nondeducibility in the source group is not confused with indistinguishability of group tuples after adding pairings.
\end{enumerate}
\end{proof}

\[
\begin{tikzcd}[column sep=large]
G_A\times G_A\arrow[r,"e"]\arrow[d,"(-)^a\times(-)^b"']
 &G_{T,A}\arrow[d,"(-)^{ab}"]\\
G_A\times G_A\arrow[r,"e"']&G_{T,A}
\end{tikzcd}
\]
These models are constructed from the free scalar algebra and its localization; the characteristic-zero instance is obtained with $k=\mathbb Q$. Even for finite $k$, the allocated names are indeterminates, so this is not a finite cyclic group implementation or a proof of computational discrete-log hardness. They realize the displayed algebraic operations by calculation. The established Diffie--Hellman, pairing and XOR message theories in~\cite{crypto-tamarin} provide further presentations; identifying different presentations requires their actual equation-preserving maps.

XOR itself has a similarly concrete carrier: the finite subsets of a set of generators, with symmetric difference as addition and the empty set as zero. Equivalently this is the free $\mathbb F_2$-vector space on those generators. Associativity, commutativity and $x\mathbin\oplus x=0$ follow coefficientwise. No injection of its elements into the scalar field above is declared; these are different sorts and different algebraic theories.

\subsubsection{Reconstruction and coordinate symmetry}\label{subsec:crypto-sharing}
Threshold sharing requires two algebraic facts: enough shares reconstruct the secret, while fewer shares admit the translations used by the privacy comparison. We establish both facts in a common coordinate calculation before constructing the adaptive simulator.

Interpolation supplies reconstruction. An affine coordinate change supplies the coalition comparison and its fresh-coordinate simulator.

\begin{definition}[Sharing coordinate algebra]\label{def:formal-sharing-coordinate-algebra}
The sharing construction has a coefficient-independent presentation. For integers $1\leq t\leq n$, form the universal coordinate algebra
\begin{equation}\label{eq:sharing-coordinate-algebra}
 \mathcal B_n=\mathbb Z[x_1,\ldots,x_n,\Delta^{-1}],
 \qquad \Delta=\prod_i x_i\prod_{i<j}(x_j-x_i).
\end{equation}
An algebra map $\mathcal B_n\to C$ supplies evaluation coordinates $\xi_i$ together with inverses for the coordinates and their pairwise differences. These inverses are the images of elements of the localization. Put $V_C=C[Z]_{<t}\cong C^t$, let $\pi:V_C\to C$ be evaluation at zero, and write $E_J:V_C\to C^J$ for evaluation at the coordinates indexed by $J$.
\end{definition}

\begin{proposition}[Interpolation, torsors, and hidden-share translations]\label{prop:sharing-module-structure}
For $|I|=t$, $E_I$ is an isomorphism and reconstruction is
$\pi E_I^{-1}$. For $|J|<t$, put $K_J=\ker E_J$ and
\[
 \sigma_J(c)=c\prod_{j\in J}(1-Z/\xi_j).
\]
Then $\pi\sigma_J=1_C$, $E_J\sigma_J=0$, and
\[
 K_J\cong C\oplus(K_J\cap\ker\pi),\qquad
 f\longmapsto(\pi f,f-\sigma_J(\pi f)).
\]
Every fibre of $E_J$ is a torsor for the additive group of $K_J$.
Translation by $\sigma_J(c)$ fixes observed shares and shifts the
secret by $c$. All these constructions commute with base change
of $\mathcal B_n$-algebras.
\end{proposition}
\begin{proof}\vmode
\leavevmode
\begin{enumerate}[label=(\roman*)]
\item \textbf{Interpolation isomorphism.}
The determinant of the coefficient matrix of $E_I$ is the product
of the invertible differences $\xi_j-\xi_i$.
Its inverse columns are
the Lagrange polynomials
$L_i(Z)=\prod_{j\in I\setminus\{i\}}(Z-\xi_j)/(\xi_i-\xi_j)$.

\item \textbf{Sections of partial evaluation.}
For $J$, the same polynomials of degree $|J|-1<t$ give a section
of $E_J$; for $J=\varnothing$ its target is the zero module.
Thus every fibre is inhabited.

\item \textbf{Torsor fibres.}
Subtracting two elements in a fibre
gives its unique translating element of $K_J$, proving the torsor
claim.

\item \textbf{Hidden-coordinate decomposition.}
Evaluation at zero and at $\xi_j$ gives the two equations
for $\sigma_J$.
The displayed decomposition has inverse
$(c,g)\mapsto\sigma_J(c)+g$, as direct substitution verifies.

\item \textbf{Translations.}
Translation has inverse translation by $-\sigma_J(c)$.

\item \textbf{Base change.}
The sections split the relevant exact sequences, so tensoring with
any $C'$ preserves their kernels and decompositions.
Their formulas
use only the units in~\eqref{eq:sharing-coordinate-algebra}; hence
these identifications are the stated base-change maps.
\end{enumerate}
\end{proof}
\[
\begin{tikzcd}[column sep=large]
V_C\arrow[r,"E_I","\cong"']\arrow[dr,"\pi"']
 &C^I\arrow[d,"\mathsf{rec}_I"]\\
 &C
\end{tikzcd}
\qquad
\begin{tikzcd}[column sep=large]
V_C\arrow[r,"+\sigma_J(c)"]\arrow[d,"{(\pi,E_J)}"']
 &V_C\arrow[d,"{(\pi,E_J)}"]\\
C\times C^J\arrow[r,"(+c)\times1"']&C\times C^J
\end{tikzcd}
\]
For $\varphi:C\to C'$ the base-change comparison is
\[
\begin{tikzcd}[column sep=large]
C'\otimes_C V_C\arrow[r,"1\otimes E_I"]\arrow[d,"\cong"']
 &C'\otimes_C C^I\arrow[d,"\cong"]\\
V_{C'}\arrow[r,"E'_I"']&(C')^I .
\end{tikzcd}
\]
This is the module-theoretic reconstruction and fibre symmetry. The symbolic privacy result below realizes that symmetry as an automorphism of the generated public scalar algebra. It does not infer a probability distribution from a module translation.

For an executable scalar instance, use the effective field presentation of Section~\ref{subsec:crypto-algebraic} with specified public coordinates $\xi_i$ obtained from a map $\mathcal B_n\to k$. There are explicit such instances: $k=\mathbb Q$, $\xi_i=i$; or search for a prime $p>n$, take $k=\mathbb F_p$, and use the residue classes of $1,\ldots,n$. Trial division and the unboundedness of primes make this a terminating construction. The allocated sharing coefficients remain independent indeterminates in $K_A$ in both instances.

Use these coordinates and the same integer parameters $1\leq t\leq n$. The sharing worker receives a scalar $m$, allocates new scalar indeterminates $a_1,\ldots,a_{t-1}$ after receiving it, and computes
\[
 f_m(Z)=m+\sum_{\ell=1}^{t-1}a_\ell Z^\ell,
 \qquad s_i=f_m(\xi_i).
\]
It sends $s_i$ on the designated private port to party $i$. The coefficients have no outgoing port. A reconstruction worker receiving $t$ shares at distinct prescribed points $I$ returns
\[
 \mathsf{rec}_I((s_i)_{i\in I})=\sum_{i\in I}\lambda_i s_i,
 \qquad \lambda_i=\prod_{j\in I\setminus\{i\}}\frac{-\xi_j}{\xi_i-\xi_j}.
\]
Distinct point labels are checked explicitly; missing or duplicate labels fail. This worker is a threshold reconstruction operation, not a verifier of arbitrary maliciously supplied shares.

\begin{theorem}[Threshold reconstruction and exact symbolic privacy]\label{thm:constructed-sharing}
The reconstruction worker returns $m$ on every prescribed set of $t$ shares. For a set $J$ of fewer than $t$ receiving ports, the joint frame of their shares is statically equivalent for any two base-field messages $m_0,m_1$, with the same public base field. If the secret is itself a private scalar indeterminate, that joint frame does not derive it.
\end{theorem}
\begin{proof}\vmode
\leavevmode
\begin{enumerate}[label=(\roman*)]
\item \textbf{Reconstruction.}
Proposition~\ref{prop:sharing-module-structure} gives $\mathsf{rec}_IE_I=\pi$ over the scalar algebra.
Applying this identity to $f_m$ gives the reconstruction formula and the returned value $m$.
In coordinates, the reconstruction coefficients are the values $L_i(0)$ of the inverse Vandermonde columns.

\item \textbf{Privacy coefficient field.}
For privacy write $L$ for the field before the new coefficients were allocated and $\delta=m_1-m_0\in L$.

\item \textbf{Compensating polynomial.}
Define the explicitly computed polynomial
\[
 b(Z)=\delta\left(1-\prod_{i\in J}(1-Z/\xi_i)\right)
       =\sum_{\ell=1}^{t-1}b_\ell Z^\ell,
\]
padding zero coefficients.
It has constant coefficient zero and $b(\xi_i)=\delta$ for $i\in J$.

\item \textbf{Coefficient translation.}
The assignment $a_\ell\mapsto a_\ell+b_\ell$ extends to an $L$-automorphism $\alpha$ of $L(a_1,\ldots,a_{t-1})$, with explicit inverse $a_\ell\mapsto a_\ell-b_\ell$.
Hence
\[
 \alpha(f_{m_0}(\xi_i))=f_{m_1}(\xi_i)\quad(i\in J).
\]
It preserves field operations, equality, nonzero tests and failure of inversion.

\item \textbf{Preservation of public operations.}
It fixes the coefficient field $k$, and therefore preserves and reflects coefficient reading: if $\alpha(u)=c\in k$, then $u=\alpha^{-1}(c)=c$.
Extend it through all free envelope constructors.

\item \textbf{Frame comparison.}
It gives an arrow in $\mathsf{Pres}_{\rm pub}(W_J)$, so Proposition~\ref{prop:crypto-public-comparison} identifies every joint coalition test.

\item \textbf{Secret-shifting automorphism.}
For nondeducibility let the secret be a new indeterminate $s$ over the public base field.
Apply $s\mapsto s+1$ and simultaneously translate the coefficients by those of $-\left(1-\prod_{i\in J}(1-Z/\xi_i)\right)$.
This is an invertible translation of the independent indeterminates and fixes each observed $s_i$.

\item \textbf{Nondeducibility.}
It fixes all public generators but moves $s$.
Proposition~\ref{prop:crypto-information-invariants} puts the generated algebra inside the fixed-value subalgebra, so the moved value $s$ is not derivable.

\item \textbf{Observer allocations.}
New observer allocations extend the automorphism by the identity, which proves the result for arbitrary finite public computations.
\end{enumerate}
\end{proof}

\[
\begin{tikzcd}[column sep=large]
W_J\arrow[r,"(f_{m_0}(\xi_i))_{i\in J}"]\arrow[dr,"(f_{m_1}(\xi_i))_{i\in J}"']
 &A_{m_0,J}\arrow[d,"\alpha","\cong"']\\
 &A_{m_1,J}
\end{tikzcd}
\]
The proof uses a field automorphism, not a probability distribution on coefficients. It is also why exposing the polynomial representation or a test for being a freshly allocated scalar would be incorrect: such a test need not be preserved by translation. The finite fraction evaluator provides exactly the public scalar interface stated above. The familiar polynomial sharing construction originates in~\cite{crypto-shamir}; the theorem here proves its specified symbolic observation property. Verifiable multi-secret sharing has an applied-pi presentation in~\cite{crypto-vss}; an explicit verifiable service is constructed in Section~\ref{subsec:crypto-zk-applications}.

\subsubsection{Adaptive threshold disclosure and its simulator}
A fixed subthreshold frame comparison does not yet provide replies to adaptively chosen requests. We now construct an online simulator that extends its comparison after each disclosure and uses the secret only when the threshold resource releases it.

\begin{construction}[Threshold view and disclosure resources]\vmode\label{con:adaptive-sharing-resources}
\leavevmode
\begin{enumerate}[label=(\roman*)]
\item \textbf{Scalar instance.}
Fix the concrete scalar interface of Section~\ref{subsec:crypto-algebraic},
$k=\mathbb Q$, and $\xi_i=i$ for $1\leq i\leq n$.
The same formulas
apply to the other coordinate instances already constructed.

\item \textbf{Real initialization.}
A real
share-view resource accepts a private scalar input $m$, allocates
$a_1,\ldots,a_{t-1}$ afterwards, and stores $f_m$ as above.
It then
emits an initialization notification.

\item \textbf{Share-query interface.}
At its public query port a valid
index $i$ returns a cached reference to $f_m(\xi_i)$; first use of an
index issues that reference, and repetition returns the same one.
The
public scalar operations act on these references and on the caller's
own values.
Bad indices fail.
The coefficients have no public port.

\item \textbf{Ideal disclosure state.}
The ideal disclosure resource stores the same private $m$ and a set
$J$ of queried indices.

\item \textbf{Threshold replies.}
After its initialization notification, a valid
query adds $i$ to $J$ and returns $\star$ while $|J|<t$, and $m$ when
$|J|\geq t$.

\item \textbf{Repeated and invalid requests.}
Repeated indices do not increase the cardinality.
It
returns failure on bad indices.

\item \textbf{Common interface.}
The two resources have the same
initialization order and scalar-operation interface.

\item \textbf{Correctness wrapper.}
A correct wrapper from the real resource to the ideal interface stores
the distinct queried indices and their returned share references.
It
returns $\star$ before the threshold, computes the Lagrange
reconstruction at the first $t$ distinct indices, and returns that
reconstructed scalar thereafter.

\item \textbf{Hidden calls and public references.}
Its internal calls are connected and
hidden.
Public result references belong to its client-local table,
so their labels do not reveal its internal temporary references.
\end{enumerate}
\end{construction}

\begin{definition}[Threshold comparison states]\label{def:threshold-comparison-states}
A simulator state records its queried-index set $J$, cached public
references, issuer counters, and one of two modes:
\[
 \mathsf{pre}(J,(z_i)_{i\in J})\quad(|J|<t),\qquad
 \mathsf{post}(J_0,J,\widetilde f)\quad(|J_0|=t-1).
\]
In the second mode $J_0$ is the set immediately before disclosure and
$\widetilde f$ is the polynomial constructed by the displayed
Lagrange formula below. The pre-mode contains no secret input.
A comparison state additionally retains the real polynomial $f_m$,
the common ambient coefficient field $L$, and the map
\[
 \phi_J:L(z_i:i\in J)\longrightarrow
          L(f_m(\xi_i):i\in J),\qquad
 \phi_J(z_i)=f_m(\xi_i),\quad \phi_J|_L=\id
\]
in pre-mode. In post-mode use $\phi_{J_0}$ and its equality
$\phi_{J_0}(\widetilde f)=f_m$. Tables are related entrywise by the
restriction of this map to their actually accessible values; cache
identifiers and public counters are equal. The auxiliary $L$, which
contains $m$, belongs to the comparison proof, not to the simulator's
pre-mode storage.
\end{definition}

\begin{lemma}[Threshold transition invariant]\label{lem:threshold-transition-invariant}
The comparison states have the following total successor cases, and
each preserves the entrywise table relation:
\[
\begin{array}{c|c}
\text{request case}&\text{comparison update}\\ \hline
\text{invalid index}&\text{same state; failure}\\
i\in J&\text{same state; cached identifier}\\
i\notin J,\ |J|+1<t&J'=J\cup\{i\},\ z_i\text{ fresh},\ \phi_{J'}\\
i\notin J,\ |J|=t-1&J_0=J,\ \mathsf{post}(J_0,J\cup\{i\},\widetilde f)\\
\text{new index after disclosure}&\text{issue }\widetilde f(\xi_i)
       \text{ and extend the cache}\\
\text{public scalar operation}&\text{apply the same operation under }\phi\\
\text{independent public allocation}&\text{extend }\phi\text{ by the same new generator}
\end{array}
\]
Before disclosure the comparison extensions satisfy
\[
 \begin{tikzcd}[column sep=large]
 L(z_j:j\in J)\arrow[r,"\phi_J"]\arrow[d,hook]
 &L(f_m(\xi_j):j\in J)\arrow[d,hook]\\
 L(z_j:j\in J')\arrow[r,"\phi_{J'}"']
 &L(f_m(\xi_j):j\in J').
 \end{tikzcd}
\]
\end{lemma}
\begin{proof}\vmode
\leavevmode
\begin{enumerate}[label=(\roman*)]
\item \textbf{Subthreshold field isomorphism.}
For $q=|J|<t$, the first $q$ coefficient columns give the matrix
$M=(\xi_i^\ell)_{i\in J,1\leq\ell\leq q}$ with nonzero determinant
$(\prod_i\xi_i)\prod_{i<j}(\xi_j-\xi_i)$, with the corresponding
ordering sign.
Solving
$z=m\mathbf1+Ma+Nb$ for $a$ gives an inverse field substitution.
Thus the queried values are algebraically independent over $L$ and
$\phi_J$ is an isomorphism.

\item \textbf{Extension compatibility.}
At a new subthreshold index the same
calculation applies; the two maps agree on $L$ and every old
generator, proving the square.

\item \textbf{Threshold interpolation.}
At disclosure, Lagrange interpolation on $0$ and the $t-1$ old
coordinates gives a degree-below-$t$ polynomial with constant $m$
and old values $z_j$.
Applying $\phi_{J_0}$ gives the same values
as $f_m$ at these $t$ distinct points.
Their difference has degree
below $t$ and $t$ roots; successive division by the distinct linear
factors, or the invertible evaluation matrix, makes it zero.

\item \textbf{Later evaluations.}
Consequently every later cached or new evaluation corresponds.
For $t=1$ this is the constant polynomial $m$.

\item \textbf{Public operations.}
Field isomorphisms preserve and reflect zero, inversion domains and
equality.
They fix the ground field, hence also the success and result
of coefficient reading.

\item \textbf{Allocation and query updates.}
Independent allocations adjoin a new
indeterminate on both sides, extending the same map.
Invalid and
repeated queries have the stated unchanged-state clauses.

\item \textbf{Reference-table invariant.}
Reference
issuance appends corresponding values at equal counters, proving the
table invariant and the strict successor equations of
Construction~\ref{con:evolving-comparison-coalgebra}.
\end{enumerate}
\end{proof}

\begin{theorem}[Adaptive symbolic threshold realization]\label{thm:adaptive-sharing}
The wrapper of Construction~\ref{con:adaptive-sharing-resources} is
correct and has an online simulator for the share-query interface.
The simulator learns $m$ only when its ideal resource releases it at
the threshold. The real and simulated resources are contextually
Scott equivalent at their common honest-input and public-command
boundary. The correctness wrapper and online simulator therefore give the mutual slice factorizations of Theorem~\ref{thm:correctness-security-mutual-factorization}, and equal attainable-view images.
\end{theorem}
\begin{proof}\vmode
\leavevmode
\begin{enumerate}[label=(\roman*)]
\item \textbf{Correctness wrapper.}
The correct wrapper counts distinct indices and, at the threshold,
uses $\pi E_I^{-1}$.
Invalid requests and repetitions leave its
count unchanged.

\item \textbf{Pre-disclosure simulation.}
The simulator waits for initialization and forwards each query to
the ideal resource.
Before disclosure it allocates a fresh scalar
$z_i$ at each new valid index, issues its reference, and caches it.

\item \textbf{Threshold interpolation.}
At the first reply revealing $m$, let $J$ be the $t-1$ old indices
and compute
\[
 \widetilde f(Z)=m\prod_{j\in J}\frac{Z-\xi_j}{-\xi_j}
  +\sum_{i\in J}z_i\frac{Z}{\xi_i}
       \prod_{j\in J\setminus\{i\}}
           \frac{Z-\xi_j}{\xi_i-\xi_j}.
\]
Return its value at the new index.

\item \textbf{Later and invalid queries.}
Thereafter use this polynomial
for new indices and the cache for repetitions.
At $t=1$ the formula
is $\widetilde f=m$.
Invalid indices fail without update.

\item \textbf{Available simulator operands.}
Every instruction is a finite scalar computation on values actually
available to the simulator; it never reads $m$ before disclosure.

\item \textbf{Preserved comparison.}
Lemma~\ref{lem:threshold-transition-invariant} constructs the
compatible field isomorphisms $\phi_J$, including their extensions
at new indices, public allocations, and threshold interpolation.
Its command cases preserve every public operation, failure and
literal test result, so an adaptive observer chooses the same next
query.

\item \textbf{Reference-state relation.}
Equal counters and cached identifiers give the two projections
of the evolving comparison coalgebra.
Its auxiliary field $L$ may
contain $m$; the simulator's pre-mode contains only $z_J$ and its cache.

\item \textbf{Causal history pairing.}
Match the input, initialization, and each public request and reply
at their identical process positions.
Private ideal calls and
interpolation are internal sequential blocks of that request; their
entries and exits have the same boundary dependencies.
Thus the
strict comparison gives a two-sided finite causal certificate.

\item \textbf{Realization and factorization.}
Theorem~\ref{thm:categorical-crypto-bridge} gives the simulator
equation, and the correct wrapper gives mutual factorization and
equal attainable-view images.
\end{enumerate}
\end{proof}


The theorem uses one simulator for every honest secret input. Before disclosure, its state contains the fresh $z_i$ and public cache. The comparison field may mention $m$, but does not supply it to that state.

Privacy is thus indexed by the ideal disclosure policy: subthreshold privacy does not assert indistinguishability after an authorized threshold disclosure. Restricting the view to fewer than $t$ indices gives the following input-independent instance.

At an observation port, a share reference still exposes the operations of its public scalar algebra. For any declared readout $h$ to an algebraic domain and related real and simulated states, the field comparison in the proof intertwines $h$ on their public results. For every open $U$, membership of $h$ in $U$ therefore agrees after each adaptive prefix. In particular the scalar sort with equality order recovers all the original literal replies. A subset of a private field presentation defines a public test only when its readout descends through the public projection as in Corollary~\ref{cor:open-operation-descent}. The interpolation at the threshold still uses the field operations and the earlier issued shares.

\begin{corollary}[Input-independent simulation below threshold]\label{cor:input-independent-threshold-simulation}
For a fixed set $J$ with $|J|<t$, the one-shot joint share-publication
view has a simulator which receives only the publication notification
and allocates $|J|$ fresh scalar indeterminates. Its view-resource
comparison gives a simulation square uniformly in the
secret supplied before coefficient allocation.
\end{corollary}
\begin{proof}\vmode
\leavevmode
\begin{enumerate}[label=(\roman*)]
\item \textbf{Subthreshold restriction.}
Restrict the simulator and comparison of Theorem~\ref{thm:adaptive-sharing}
to those indices and their subsequent evaluator calls.

\item \textbf{Input-independent simulation.}
The threshold
is never reached, so only the initialization notification and fresh
$z_i$ are used.
The accessible-field comparison proves the assertion
for their joint publication as well as their sequential issuance.
\end{enumerate}
\end{proof}
The theorem specifies a read-only share-query resource with adaptive
disclosure. It introduces neither adversarial writes to an honest
sharing polynomial nor an implementation of distributed verifiable
sharing. Those are separate interfaces considered in
Section~\ref{subsec:crypto-zk-applications}.

\subsubsection{The distributed linear computation}\label{subsec:linear-computation}
The sharing symmetries can also compare the views of a distributed linear computation. The proof must retain the actual message suppliers and collector waits as well as the scalar relation, because the conclusion concerns the specified concurrent network.

The maps $E_I$ and $\pi$ are module homomorphisms, so their
reconstruction square transports linear combinations. For its distributed realization, input worker $i$ shares $x_i$ with an independent coefficient family $a_{i,\ell}$. Collector $j$ receives the $j$-th share from every input worker and computes
\[
 z_j=\sum_{i=1}^d c_i f_{x_i}(\xi_j)
\]
for public coefficient-field elements $c_i$. It sends $z_j$ to a reconstruction worker, which uses any $t$ designated collectors to return $y=\sum_i c_ix_i$. Each collector waits for one receipt from each input worker, retains their supplier identities, and replies with failure to duplicates. Inputs at distinct collectors use separate queues and banks. A read-only observation port permits a designated coalition to inspect its collectors' values through the scalar evaluator; it grants no write command on another collector's state.

\begin{theorem}[Correctness and privacy of the linear network]\label{thm:linear-mpc}
This network computes $y$. For any fixed coalition $J$ of fewer than $t$ collector observation ports, two input tuples with the same output $y$ have the same finite visible causal observations at those ports and the public output port. This is a security statement for the constructed read-only observation interface, not malicious-write tolerance.
\end{theorem}
\begin{proof}\vmode
\leavevmode
\begin{enumerate}[label=(\roman*)]
\item \textbf{Correctness polynomial.}
The polynomial $g(Z)=\sum_i c_i f_{x_i}(Z)$ has degree below $t$ and constant term $y$; interpolation proves correctness.

\item \textbf{Compensating translations.}
Let $\delta_i=x'_i-x_i$ and put $b_i(Z)=\delta_i(1-\prod_{j\in J}(1-Z/\xi_j))$.
Translate the fresh coefficient family for each input by the coefficients of $b_i$.

\item \textbf{Joint field automorphism.}
These translations commute on their disjoint families and have the inverse translations.
They carry every observed share to its counterpart.

\item \textbf{Preservation of the output.}
Since $\sum_i c_i\delta_i=0$, we have $\sum_i c_ib_i(Z)=0$, so every final collector value $z_j$, including those outside $J$, is fixed.

\item \textbf{Matching instruction occurrences.}
The honest programs have the same finite instruction graphs for both input tuples: all scalar additions and multiplications are total, and their readiness depends only on the received point and sender labels.
Match each allocation, queue receipt, computation, and reply with its same instruction occurrence.

\item \textbf{Observer operations.}
At an observed collector, the field automorphism carries the incoming frame and every ensuing public evaluator result to its counterpart.
Equality and coefficient reading therefore choose the same observer branches.

\item \textbf{Output and causal order.}
The final output is literally the same $y$.
Chosen FIFO suppliers, collector waits and independent banks are unchanged, so this matching preserves each finite configuration order.

\item \textbf{Reverse comparison and contexts.}
In the reverse direction use the inverse translations.
Lemma~\ref{lem:crypto-causal-congruence} gives equality in every closing observation context at the declared boundary.
No interleaving replaces the collectors' independent executions.
\end{enumerate}
\end{proof}

\subsection{Witness families and verified disclosure}\label{subsec:witness-families}
The category $\Set/X$ retains a family of witnesses over each
statement. It therefore records more than the validity subobject
$\im(p)\hookrightarrow X$. Its products are conjunction over a
common statement; composition of display maps supplies dependent
sum. The controller below realizes these operations while keeping
the witness field private.

\begin{construction}[Certificate controller from a checked relation]\vmode\label{con:generic-certificate}
\leavevmode
\begin{enumerate}[label=(\roman*)]
\item \textbf{Checked record carrier.}
For a concrete checked relation $R\subseteq X\times W$, let
$p:R\to X$ be its statement projection.

\item \textbf{Private issuance.}
A private issuance operation
checks $(x,w)\in R$ and appends that record to a finite certificate
table $T:H\rightharpoonup R$, under a fresh identifier.
Invalid
issuance returns failure.

\item \textbf{Public commands.}
Public commands read $pT(h)$, check that
$h$ is issued, and compare issued identifiers; invalid identifiers
fail.

\item \textbf{Ideal issuance and erased table.}
The ideal issuance port performs the same private check and
notifies the simulator only of the validated statement and the issued
identifier.
Thus its public table is $\bar T=pT$, with the same
domain and counters.
\end{enumerate}
\end{construction}

\begin{proposition}[Certificate erasure and closure]\label{prop:generic-certificate-erasure}
The table map $T\mapsto pT$ is a strict public-controller comparison.
It yields a simulator using the validated issuance notifications.
Derived certificate operations are transported by commuting squares
\[
\begin{tikzcd}[column sep=large]
 \prod_{i=1}^n R_i\arrow[r,"k"]\arrow[d,"\prod_i p_i"']
 &R\arrow[d,"p"]\\
 \prod_{i=1}^n X_i\arrow[r,"k_0"']&X.
\end{tikzcd}
\]
Their composites and independent products satisfy the same erasure
equation. Every issued statement belongs to $\im p$, with its
constructed witness retained in the private audit.
\end{proposition}
\begin{proof}\vmode
\leavevmode
\begin{enumerate}[label=(\roman*)]
\item \textbf{Table and lookup correspondence.}
The erasure fixes table domains, identifiers and counters.
Reading
returns $pT(h)$ on both sides; checking and equality use the same
identifier, and invalid lookup fails on both.

\item \textbf{Issuance updates.}
A valid private issuance
notifies the ideal side of precisely $p(r)$; both extend at the same
fresh identifier.
Invalid issuance creates no entry.

\item \textbf{Derived commands.}
For a derived command, the commuting square says that computing
$k$ then erasing gives the same new entry as computing $k_0$ on
the erased operands.
Invalid lookups fail first.

\item \textbf{Strict transition equations.}
These are the
strict transition equations of Corollary~\ref{cor:reference-descent-realization}.

\item \textbf{Composition and products.}
Pasting and products of the squares prove closure.

\item \textbf{Audit extraction.}
The finite
issuance audit evaluates each node in $R$, proving validity and
private witness extraction.

\item \textbf{Service replacement.}
Theorem~\ref{thm:service-behaviour-canonical-comparison}
interprets the initialized square as the claimed simulator.
\end{enumerate}
\end{proof}
For families $p:R\to X$, $p':R'\to X$, conjunction is their
pullback $R\times_XR'\to X$. Independent statements instead use
$R\times R'\to X\times X'$. For a witness family
$p_Z:R\to X\times Z$ and projection $q:X\times Z\to X$, hiding
the $Z$ component is $\Sigma_q(p_Z)=qp_Z$. Explicitly,
\[
 \operatorname{Hom}_{\Set/X}(\Sigma_qR,S)
 \cong\operatorname{Hom}_{\Set/(X\times Z)}(R,q^*S).
\]
A map $a:R\to S$ over $X$ corresponds to
$r\mapsto(p_Z(r),a(r))$ into the pullback; projection is its inverse.
This proves $\Sigma_q\dashv q^*$ by construction. Rebracketing
$((x,z),w)\mapsto(x,(z,w))$ realizes it on checked witness records.
These pullbacks are in $\Set$, and their reference controllers
are subsequently interpreted in games. The validity image and the
attainable-behaviour image of a simulator remain different constructions.

\subsubsection{Proof objects, witness erasure, and symbolic zero knowledge}\label{subsec:crypto-zk}
This instantiates the public-operation descent criterion of Theorem~\ref{thm:operation-descent}, through Construction~\ref{con:generic-certificate} and Proposition~\ref{prop:generic-certificate-erasure}. Certificates are returned references to checked private records. Their audit supplies witnesses, while erasure preserves the public service. These are constructions of a proof resource, not an implementation of a distributed zero-knowledge protocol.

\begin{construction}[Checked certificate instance]\vmode\label{con:checked-certificate-service}
\leavevmode
\begin{enumerate}[label=(\roman*)]
\item \textbf{Relation checker.}
A relation code $R$ is a finite strict Boolean program of the existing value calculus, with statement carrier $X_R$ and witness carrier $W_R$.
Its successful domain and statement projection are
\[
 V_R=\{(x,w):\crEval_R(x,w)=\mathsf{true}\},\qquad
 p_R:V_R\to X_R.
\]
Failure rejects; codes have syntactic equality and never inspect private certificate records.

\item \textbf{Conjunction and existential carriers.}
In the slice over $X$, conjunction is $V_R\times_XV_S$; existential hiding is the composite $V_R\to X\times Z\to X$.

\item \textbf{Derived checker programs.}
The conjunction checker evaluates $R(x,w)$ and $S(x,v)$ strictly and returns their Boolean conjunction; the hiding checker on $(u,(z,w))$ evaluates $R((u,z),w)$.
These are finite programs in the existing calculus.

\item \textbf{Statement comparison maps.}
Their comparison maps are
\[
 \kappa((x,w),(x,v))=(x,(w,v)),\qquad
 \eta((u,z),w)=(u,(z,w)).
\]
These are isomorphisms onto $V_{R\wedge S}$ and $V_{\exists_ZR}$ respectively, by evaluating those checker programs.
Their statement squares commute by the product equations.

\item \textbf{Certificate-service realization.}
Apply Construction~\ref{con:generic-certificate} to the coproduct of the $p_R$, retaining code and statement in the public projection.
Proposition~\ref{prop:generic-certificate-erasure} lifts $\kappa,\eta$ to the commands $\mathsf{conjoin},\mathsf{hide}$.

\item \textbf{Refresh and public observations.}
Refresh is identity on the checked record followed by fresh issuance.
Statement and verification are the existing public projection and its equality test.

\item \textbf{Identifiers and lookup failure.}
Copy preserves the identifier; each issuance, including refresh, allocates a new one.
Lookup, code, statement and sort mismatches use the reference service's failure continuation.

\item \textbf{Honest and public issuance.}
The honest issuance process supplies its witness at a private port and clears its temporary input after checking.
A public prover supplies only values available in its own reference bank.

\item \textbf{Retained client banks and transport.}
These banks are retained by witness erasure: forgetting a certificate's private witness does not erase a witness already known by a client.
Actual transport carries the issued identifier and its parent occurrence.

\item \textbf{Scope of public constructors.}
A deterministic public constructor on $(R,x,w)$ would expose candidate-witness equality; fresh issuance is therefore an operative part of this service.
\end{enumerate}
\end{construction}

\begin{theorem}[Validity, audit extraction, and witness-erasure simulation]\label{thm:proof-erasure}
Every verified certificate of this service has a witness in $V_R$, extracted by a finite fold of its private issuance audit. Erasing private witness fields gives a strict comparison to the statement-record service. Their public reference strategies agree, including adaptive verification, refresh, conjunction, hiding, copying and transport.
\end{theorem}
\begin{proof}\vmode
\leavevmode
\begin{enumerate}[label=(\roman*)]
\item \textbf{Audit extraction.}
The audit fold interprets checked issuance by its accepted witness,
refresh by identity, conjunction by $\kappa$, and hiding by $\eta$.
Their checker equations put each result in the required fibre;
copy and transport retain the selected ancestry.
This constructs
extraction on the private audit.

\item \textbf{Witness-erasure comparison.}
Proposition~\ref{prop:generic-certificate-erasure} applies to these
same squares.

\item \textbf{Public prover interface.}
At a public prover port both sides check the same
supplied arguments; rejection creates no record and acceptance uses
the same fresh identifier.

\item \textbf{Fixed public data.}
Erasure fixes clients' banks, statements,
parent pointers and counters, hence also equality and failed lookups.

\item \textbf{Simulation and issuance dependencies.}
The resulting strict service comparison identifies the two public displays. The identity strategy on their common interface therefore witnesses replacement in the sense of Theorem~\ref{thm:observational-simulation}.
The private issuance call and return keep their dependencies; its
witness field and audit-extraction port remain private.
\end{enumerate}
\end{proof}

For an observation $U$ on a validated statement the corresponding private-record predicate is $p_R^{-1}U$. It is Scott open at a continuous statement projection, including the discrete instance above. Each command square $p'k=\bar k p$ gives $k^{-1}(p'^{-1}U)=p^{-1}(\bar k^{-1}U)$, with the failure summand included for a failing operation. Iteration proves the same identity after a command sequence. Refresh issues a new identifier and preserves the projected statement, so observing the statement after refresh preserves both that membership result and the fresh issuance occurrence. A private witness can still be recovered by the audit projection where that projection is declared.

\begin{corollary}[Certificate families and erasure descent]\label{cor:crypto-proof-descent}
At a finite history, the issued certificates of code $R$ form a finite family $c_h:C_h\to X_R$ with a privately extracted lift $w_h:C_h\to V_R$ satisfying $p_Rw_h=c_h$. These lifts commute with history inclusions. Erasure retains $C_h,c_h$ and forgets $w_h$. Public service behaviour factors through the reachable erased-state image and is constant on the kernel pair of the erasure map.
\end{corollary}
\begin{proof}\vmode
\leavevmode
\begin{enumerate}[label=(\roman*)]
\item \textbf{Issued witness family.}
Take $C_h$ to be the issued identifiers and $c_h$ their stored statements.

\item \textbf{Extraction and extension.}
The audit recursion in Theorem~\ref{thm:proof-erasure} supplies $w_h$.
An extension preserves the existing parent records, so the same recursion gives the same witness on every old identifier, proving naturality.

\item \textbf{Public behaviour descent.}
The service square gives $\beta_P=\beta_{P^0}a$; factor $a$ through its reachable image to obtain descent.
Equal erased states thus have equal public behaviour.

\item \textbf{Retention of witness multiplicity.}
The family $C_h\to X_R$ is retained: its identifiers distinguish repeated copying of one certificate from two fresh issuances.
It is not replaced by the subobject $\im(p_R)$.
\end{enumerate}
\end{proof}
\paragraph{Certificate validity and public Scott observations}
\begin{theorem}[Certificate soundness and observable validation]\label{thm:certificate-soundness-validation}
For the finite issued-reference family $c_h:C_h\to X_R$, the private audit defines $w_h:C_h\to V_R$ with $p_Rw_h=c_h$, naturally under cumulative-history extension. Its statement map consequently factors through $P_R=\im(p_R)$ in sets. At the actual public game boundary, the open $U_{\mathrm{verify}}$ witnessing successful verification is contained in the public approximation of the historical open $U_{\mathrm{valid}}$ witnessing a checked issuance ancestry:
\[
 U_{\mathrm{verify}}\subseteq\Box_J U_{\mathrm{valid}}.
\]
The service supplies no public witness-extraction operation by this factorization.
\end{theorem}
\begin{proof}\vmode
\leavevmode
\begin{enumerate}[label=(\roman*)]
\item \textbf{Semantic soundness factorization.}
Corollary~\ref{cor:crypto-proof-descent} constructs $w_h$ by the finite audit fold, including the actual parent witnesses for conjunction and hiding.
Its equation proves the set-theoretic image factorization.

\item \textbf{Historical validation open.}
A verification occurrence and its finite checked ancestry persist in cumulative histories, so their existence defines Scott opens: each is a union of principal finite-configuration opens.

\item \textbf{Origin implication.}
The private audit proof gives $U_{\mathrm{verify}}\subseteq U_{\mathrm{valid}}$.

\item \textbf{Public approximation.}
The successful reply is a public event, hence $U_{\mathrm{verify}}$ is public.
Proposition~\ref{prop:interior} gives the displayed inclusion.

\item \textbf{Scope of semantic evidence.}
The factorization is a semantic statement about these records; the constructed public eliminators still return only statements and verification results.
\end{enumerate}
\end{proof}
Use the erased controller $P^0$ as the ideal verification resource.
Its private input checks $(R,x,w)$ before notifying the simulator
of $(R,x)$ and the issuance event. The simulator maintains the
erased table; it cannot mint an unchecked certificate. Once the
private input wires and temporary buffers are closed, accepted
witnesses with equal $(R,x)$, client banks and call pattern give
the same erased state, and hence the same public interaction.
This realizes witness privacy at the constructed proof-resource
interface; its checking authority is still a private resource.

\subsubsection{Verifiable sharing, selective disclosure, and certified computation}\label{subsec:crypto-zk-applications}
These examples compose the actual commitment, certificate, and value commands. The private witness maps and the erasure comparison prove their validity and privacy.

\begin{proposition}[A consistency construction for verifiable sharing]\label{prop:verifiable-sharing}
Commit to each coefficient $a_0,\ldots,a_{t-1}$ of a sharing polynomial with an independently allocated private salt:
\[
 c_\ell=H(\langle\iota_{\mathsf{Scalar}}a_\ell,\iota_{\mathsf N}r_\ell\rangle).
\]
Publish $\vec c$. Privately to recipient $i$, deliver $s_i$ and a certificate of the relation
\[
 R_{t}(i,s_i,\vec c;\vec a,\vec r)\;\Longleftrightarrow\;
 \bigwedge_{\ell<t} c_\ell=H(\langle\iota_{\mathsf{Scalar}}a_\ell,\iota_{\mathsf N}r_\ell\rangle)
 \ \wedge\ s_i=\sum_{\ell<t}a_\ell \xi_i^\ell.
\]
A recipient checks the certificate's entire statement against its received $i,s_i,\vec c$. Every $t$ accepted shares at distinct prescribed labels, with the coordinates $\xi_i$ of Section~\ref{subsec:crypto-sharing}, reconstruct the same committed constant coefficient. For the honest private dealer, fewer than $t$ recipients' joint views, together with all public coefficient commitments, are statically independent of the secret. The assertion is consistency and privacy; acceptance does not force a dealer to deliver missing shares.
\end{proposition}
\begin{proof}\vmode
\leavevmode
\begin{enumerate}[label=(\roman*)]
\item \textbf{Finite consistency checker.}
The relation is a finite checker using multiplication, addition, hash equality and the known finite parameter $t$.

\item \textbf{Extraction and binding.}
Audit extraction gives an opening vector and the share equation for every accepted certificate.
Binding of each coefficient commitment identifies these vectors across recipients.

\item \textbf{Reconstruction.}
Lagrange reconstruction in Theorem~\ref{thm:constructed-sharing} therefore yields the same $a_0$ from any $t$ distinct accepted shares, including when the dealer selected its inputs adversarially.

\item \textbf{Privacy translation.}
For honest privacy use the coefficient translation in that theorem for the observed recipient set $J$.

\item \textbf{Commitment markers.}
Replace each coefficient commitment by a distinct marker.

\item \textbf{Salt invariant and expansion.}
No salt is derivable: every occurrence outside the dealer is under its hash, and proof-record witnesses have no public projection.
Expansion is injective by the argument of Theorem~\ref{thm:hash-commitment}, simultaneously for all fresh salts.

\item \textbf{Accessible-frame comparison.}
Thus the translation on the accessible scalar field and the bijection on these commitment markers define an isomorphism of the observed generated algebras.

\item \textbf{Certificate statement relation.}
Recipient statements $(i,s_i,\vec c)$ are preserved by this isomorphism.

\item \textbf{Witness erasure.}
Erase the certificates' witness fields by Theorem~\ref{thm:proof-erasure}, then pair corresponding issuance identifiers.

\item \textbf{Public witness operations.}
Publicly proposed witnesses pass the finite relation checker on one side exactly when their images do on the other, because every operation and equality in that checker is preserved.
Refresh, conjunction and hiding retain the same relation by their command equations.

\item \textbf{Declared recipient view.}
Other recipients' statements and shares stay on their private ports.
This gives privacy of precisely the declared joint view, with no public broadcast of every $s_i$.

\item \textbf{Concurrent delivery pattern.}
The dealer sends the same finite pattern of packets and proof requests for each secret, so the same argument also preserves its visible call and delivery orders.
\end{enumerate}
\end{proof}

The symbolic VSS literature includes an applied-$\pi$ equivalence treatment in~\cite{crypto-vss}. The construction above gives the coefficient commitment and certificate mechanism explicitly; a broadcast agreement or guaranteed-delivery protocol is a separate communication task.

For selective disclosure let a signed attribute tuple be $\vec m=(\vec d,\vec h)$, with disclosed fields $\vec d$ and hidden fields $\vec h$. Define the finite checker
\[
 R_{\mathrm{cred}}((v,\vec d,e);(\vec h,\sigma))
 \quad\Longleftrightarrow\quad
 \mathsf{vfy}(\sigma,\langle\mathsf{domain},(\vec d,\vec h)\rangle,v)=\crOK.
\]
Here $\mathsf{domain}$ is a fixed public message tag and $e$ is the recipient's current challenge, included in the certificate's statement. A presentation is a fresh certificate for this relation; the recipient verifies the complete statement, including $e$. Correctness is the signature equation. A valid presentation has a signed tuple by audit extraction; at a private issuer key its signature has a signing ancestor by Corollary~\ref{cor:crypto-origin}. For two accepted tuples with equal disclosed fields, key and challenge, proof-worker witness erasure gives identical public presentation interactions. A certificate for a different challenge fails the statement check. This constructs selective disclosure and challenge binding at the proof-service boundary. It does not add non-transferable ownership of a credential or prevent a recipient from forwarding an already disclosed field. The DAA analysis in~\cite{crypto-zk-applied} motivates this combination of signatures and symbolic proofs without identifying it with that entire protocol.

For a finite circuit $C$, the checker $R_C((x,y);w)\Longleftrightarrow C(x,w)=y$ similarly gives a certificate for a computation with private inputs. Extraction proves the output equation; erasure hides the witness beyond $(C,x,y)$. The checker executes the circuit, so this construction makes no assertion that verification is shorter than circuit evaluation. The equations and interface, rather than a name such as ``SNARK'', determine what has been constructed.

\subsection{Free evaluation algebras and output descent}\label{subsec:evaluation-descent}
Correctness is naturality of an algebra homomorphism with respect
to circuit substitution. Privacy is a separate relation on public
presentations. For garbling this relation descends along the circuit's
output map, with an explicit simulator on each output fibre.

\subsubsection{Homomorphic circuit evaluation by a private fold}\label{subsec:crypto-homomorphic}


\begin{construction}[Free ciphertext algebra and decryption fold]\vmode\label{con:ciphertext-algebra-decryption-fold}
\leavevmode
\begin{enumerate}[label=(\roman*)]
\item \textbf{Plaintext algebra and circuits.}
Use the Boolean algebra or the constructed scalar algebra $B=K_A$, with its displayed gates.
Its circuits form the free finite-product syntax on those operations.

\item \textbf{Free ciphertext carrier.}
At encryption key $p$, form the free term algebra for the gate arities on the encrypted-input object $B\times N$, with a distinct constant introduction for each $b\in B$.
Write its universal introductions as $\mathsf{leaf}_p(b,r)$, $\mathsf{const}_p(b)$ and $\mathsf{gate}_{p,\omega}$; the existing free value construction gives its carrier $C_p$ and induction principle.

\item \textbf{Universal decryption fold.}
Writing $F_{\Sigma_p}\dashv U$ for the free-algebra adjunction,
$C_p=F_{\Sigma_p}(B\times N)$; the constants form part of
$\Sigma_p$.
The projection $B\times N\to B$ and the plaintext
$\Sigma_p$-algebra determine the unique fold $d_k:C_p\to B$ at $p=\crPK(k)$.

\item \textbf{Fold equations.}
Its algebra equations are
\[
 d_k\mathsf{leaf}_p(b,r)=b,\quad d_k\mathsf{const}_p(b)=b,
 \quad d_k\mathsf{gate}_{p,\omega}(\vec c)=\omega(d_k\vec c).
\]
Wrong keys fail before this elimination.

\item \textbf{Public commands and key checks.}
Public commands expose leaf, constant and gate introduction and equality, but no projection of children, plaintext or nonce.
Gate arity and common key index are their typed operands.

\item \textbf{Syntactic equality.}
The carrier is free syntax, so equal decrypted values do not equate distinct ciphertext trees.

\item \textbf{Circuit interpretation.}
For a circuit $C$, interpretation in this free algebra gives $\mathsf{Eval}_C:C_p^n\to C_p$; interpretation in $B$ gives $C_B:B^n\to B$.
Their compatibility with composition is the free finite-product interpretation.

\item \textbf{Decryption logical relation.}
The graph $R_k=\{(c,b):d_k(c)=b\}$ is a relation preserved
by every pair of corresponding ciphertext and plaintext gates.
Circuit substitution therefore preserves $R_k$.
This is the
algebra-level logical relation for correctness; the protected
reference interface realizes its public operations as strategies.

\item \textbf{Process realization.}
Section~\ref{subsec:calculus-extensions} gives their type rules.
\end{enumerate}
\end{construction}

\begin{theorem}[Homomorphic correctness and symbolic ciphertext privacy]\label{thm:homomorphic-circuits}
Encrypted circuit evaluation satisfies
\[
 d_k\bigl(\mathsf{Eval}_C(\mathsf{leaf}_p(b_1,r_1),\ldots,
 \mathsf{leaf}_p(b_n,r_n))\bigr)=C_B(b_1,\ldots,b_n)
 \quad(p=\crPK(k)).
\]
In the closed resource which allocates $k$ privately and exports finitely many encrypted inputs with fresh private nonces, all public circuit computations preserve static secrecy of those inputs at the boundary without the decryption key.
\end{theorem}
\begin{proof}\vmode
\leavevmode
\begin{enumerate}[label=(\roman*)]
\item \textbf{Circuit correctness.}
At an input node the equation is the leaf rule; at a constant node it is the constant rule; at a gate it follows by applying $\omega$ to the induction equations for its children.
This also proves compatibility with circuit substitution.

\item \textbf{Decryption homomorphism.}
Equivalently the private decryption is the algebra homomorphism in
\[
\begin{tikzcd}[column sep=large]
C_p^{n_\omega}\arrow[r,"\mathsf{gate}_{p,\omega}"]\arrow[d,"d_k^{n_\omega}"']
 &C_p\arrow[d,"d_k"]\\
B^{n_\omega}\arrow[r,"\omega"']&B.
\end{tikzcd}
\]
For privacy replace the honest leaves by distinct markers of ciphertext sort and allow all the public gate and constant operations on them.

\item \textbf{Hidden-key invariant.}
The private key cannot be extracted from a public key or a leaf.
An available-key decryption of an honest-key ciphertext therefore fails.

\item \textbf{Fresh-leaf separation.}
Leaf nonces are fresh and are not outputs of any public destructor.
A marker-free accessible tree cannot reconstruct an honest leaf, and a proper tree containing a marker cannot expand to that leaf.

\item \textbf{Injectivity and operation preservation.}
Comparing remaining roots descends to their operands; distinct honest leaves are separated by their nonces.
This proves injectivity of expansion, including all attacker-created gates and encryptions.

\item \textbf{Available-key decryptions.}
All successful decryptions at other keys, equality tests, and public constructors commute with expansion.

\item \textbf{Public-frame comparison.}
These are the expansion arrows of a common labelled public presentation, so Proposition~\ref{prop:crypto-public-comparison} proves static secrecy.

\item \textbf{Reference-state lifting.}
Retaining the expansion with the handle counters gives the command-by-command coalgebra span of Theorem~\ref{thm:service-behaviour-canonical-comparison}, and hence the adaptive evaluator comparison.

\item \textbf{Authorized decryption boundary.}
A recipient's successful authorized decryption is intentionally outside this secrecy boundary.
\end{enumerate}
\end{proof}

The gate operations freely generate the ciphertext carrier, and $d_k$ folds it into the plaintext algebra. The preceding arguments establish functionality and privacy for this construction. They do not identify every published ciphertext theory with the free-tree theory, nor establish succinct ciphertexts, circuit privacy or security for approximate decryption.

Dahl--Damg\aa rd~\cite{crypto-dahl} give a different symbolic process presentation with homomorphic encryption, commitments and proofs, including two-party protocols and a computational UC comparison. Relating their operation equations to the present carrier requires a further comparison; those equations do not follow from the fold equation alone.

For an authorized decryption readout and a Scott-open predicate $U$ of the plaintext algebra, the fold equation gives
\[
 (d_k\circ\mathsf{Eval}_C)^{-1}U
     =(C_B\circ d_k^{\,n})^{-1}U .
\]
The two sides are open for continuous interpretations of the gates and decoder, in particular for the original discrete algebras. This follows by composing the gate equations along the finite circuit. The observation belongs to the boundary of the holder of the decoding key.

\subsubsection{An MPC resource and its information boundary}\label{subsec:crypto-mpc-resource}
This stateful game is an ideal resource. Its information boundary is proved directly from its command transitions.

\begin{definition}[Circuit-evaluation resource]\label{def:formal-circuit-evaluation-resource}
Fix a finite total circuit $F:\prod_{i=1}^n X_i\to\prod_{i=1}^n Y_i$, built from the typed operations whose gate evaluations are total on these carriers. A session worker stores a partial vector $\vec x$, a status in $\{\mathsf{collecting},\mathsf{ready},\mathsf{aborted}\}$, and a set of delivered output positions. The command $\mathsf{submit}_i(x)$ is available only at party $i$'s input port, fills that party's previously empty slot, and otherwise fails. Once the last slot is filled the worker computes $\vec y=F(\vec x)$ and becomes ready. It emits public activity notifications containing only the session and party labels. The scheduling port can request $\mathsf{deliver}_i$; when ready and not previously delivered this sends $y_i$ on party $i$'s private output port and records that delivery. Before readiness or after a repeated delivery it fails. The command $\mathsf{abort}$ sets the aborted status; later delivery requests return an abort indication. Every other malformed or unauthorized request fails without change. No command reads an unowned input slot.

This is a concrete state machine, with finite state data at each stage and finite circuit evaluation. Lemma~\ref{lem:crypto-stateful} constructs its game strategy; parallel session workers give independent sessions. The scheduling interface supplies no fairness promise. In particular the worker need not produce an output after a scheduling opponent aborts or declines delivery.
\end{definition}

\begin{definition}[Fixed-input coalition observation]\label{def:mpc-view-carriers}
Fix the submission activity pattern of the resource and a coalition $J$. Close each input port with its one-shot sender, each noncoalition output with a sink, and initially disclose $\vec x_J$; retain the coalition output and scheduling ports. Let $\mathcal D_J$ be the set of finite displayed causal histories at this visible interface, up to the displayed-order isomorphisms already defined. The \textbf{behaviour carrier} is $\mathsf{Beh}_J=\Pow(\mathcal D_J)$. Define
\[
 \lambda_J:\prod_iX_i\to(\prod_{i\in J}X_i)\times(\prod_{i\in J}Y_i),\quad
 \vec x\mapsto(\vec x_J,F(\vec x)_J),\qquad
 \mathsf{View}_J:\prod_iX_i\to\mathsf{Beh}_J
\]
by taking all finite visible histories of that closed-input worker. The \textbf{reduced observation map} $\overline{\mathsf{View}}_J$ has domain $(\prod_{i\in J}X_i)\times(\prod_{i\in J}Y_i)$ and the same codomain; it evaluates the worker's flag transitions with only the supplied coalition input and output values, using the same request ports and submission pattern. The following proof checks every request case and establishes $\mathsf{View}_J=\overline{\mathsf{View}}_J\lambda_J$.
\end{definition}

\begin{proposition}[Exact information released by the MPC resource]\label{prop:mpc-resource-view}
For the fixed-input view, close all input ports with one-shot senders for $\vec x$, initially disclose $\vec x_J$ to the coalition, and close noncoalition output ports with sinks. The coalition output ports and the scheduling port remain live. For a coalition $J$, write
\[
 \lambda_J(\vec x)=(\vec x_J,F(\vec x)_J).
\]
For fixed typed inputs and a fixed submission activity pattern, the entire visible behaviour at the coalition and scheduling ports factors through $\lambda_J$. The factorization includes adaptive delivery and abort requests.
\end{proposition}
\begin{proof}\vmode
\leavevmode
\begin{enumerate}[label=(\roman*)]
\item \textbf{Reduced controller.}
Use the reduced controller of Definition~\ref{def:mpc-view-carriers}.

\item \textbf{Submission and delivery replies.}
Submission changes a slot flag and emits only its party label;
readiness depends on filled slots.
Delivery to $J$ returns the
supplied output component, while other delivery changes only its
flag at the visible boundary.

\item \textbf{Failure cases.}
Abort, duplicates and premature or
repeated deliveries inspect these same flags.

\item \textbf{Information factorization.}
Thus every reply and
successor is determined by $\lambda_J$ and the command history.
This is the explicit factorization
$\mathsf{View}_J=\overline{\mathsf{View}}_J\lambda_J$ and its
strict command comparison.

\item \textbf{Closed input realization.}
The prescribed closed input processes
have the same activity dependencies, so the visible finite causal
histories also agree.
Noncoalition input values stay on the closed
wires, outside the compared boundary.
\end{enumerate}
\end{proof}

\[
\begin{tikzcd}[column sep=large]
\prod_iX_i\arrow[r,"\lambda_J"]\arrow[dr,"\mathsf{View}_J"']
 &\left(\prod_{i\in J}X_i\right)\times\left(\prod_{i\in J}Y_i\right)
   \arrow[d,"\overline{\mathsf{View}}_J"]\\
 &\mathsf{Beh}_J
\end{tikzcd}
\]
The selection circuit $F((m_0,m_1),b)=(\star,m_b)$ gives the input/output functionality of oblivious transfer, with private receiver choice. Its delivery-controlled worker is a sequential service presentation, whereas Proposition~\ref{prop:crypto-ot} supplies the event-level OT game. This identifies their functionality; it does not equate their different scheduling interfaces without a further adapter. The linear sharing protocol in Theorem~\ref{thm:linear-mpc} is an actual distributed computation for its stated observation ports. The generic worker just constructed is an ideal resource. It is not a proof that arbitrary circuits admit maliciously secure distributed realization from the preceding channels. The symbolic MPC abstraction and verification theorem of~\cite{crypto-mpc} explicitly distinguishes these roles as well.

\subsubsection{Garbled circuits: construction, evaluation, and symbolic privacy}\label{subsec:crypto-garbling}
The garbling construction has two tasks: evaluation must return the circuit output, and the declared view must admit its symbolic privacy comparison. We specify the wire labels and encrypted rows first, then use their accessible operations to construct and transport that comparison.

Garbling combines the encryption service with the free hash introduction
of Definition~\ref{def:formal-free-symbolic-hash}. The output-only
simulator is a process on the one-shot bundle interface.

\begin{definition}[Garbled-circuit bundle]\label{def:formal-garbled-circuit-bundle}
Let $C$ be a finite acyclic Boolean circuit with binary gates; unary gates can ignore their second input, and constants can be supplied as public input bits. For each wire $w$ allocate two distinct private symmetric-key atoms $K_w^0,K_w^1$. All wire pairs are disjoint. For a gate $v=f(u,z)$ form four rows
\[
 T_v^{a,b}=\crSEnc\!\left(
       \crSEnc(\iota_{\mathsf K_s}K_v^{f(a,b)},K_z^b,r_v^{a,b}),
       K_u^a,s_v^{a,b}\right),\qquad a,b\in\{0,1\},
\]
with distinct fresh private nonces for every encryption occurrence. At each gate choose a permutation of its four rows by a finite hidden choice, and publish the resulting ordered list. Such a choice is a finite positive exclusion of the core language; no distribution is assigned. For each output wire publish the ordered decoding pair
\[
       (H(\iota_{\mathsf K_s}K_w^0),H(\iota_{\mathsf K_s}K_w^1)).
\]
For input $\vec x$, publish exactly the active key of each input wire. The other keys, the row permutations, and the encryption nonces have no public projection. Initialization finishes before publishing this bundle. Its public interface subsequently supplies the ordinary handle evaluator, with no later key-disclosure or second-input-selection command. In particular the output decoding table does not publish both output keys.
The \textbf{input-wire set} $\mathsf{In}(C)$ and \textbf{output-wire set} $\mathsf{Out}(C)$ are the designated boundary wires of $C$, so its Boolean evaluation has type $C:\{0,1\}^{\mathsf{In}(C)}\to\{0,1\}^{\mathsf{Out}(C)}$. For a wire set $W$, put $\mathsf{ActiveKeys}_W=\prod_{w\in W}\{K_w^0,K_w^1\}$. Define the \textbf{input encoding} $\mathsf{encode}_{\vec K}(\vec x)=(K_w^{x_w})_w$, the \textbf{garbled evaluation} $\mathsf{evaluate}_{\vec T}:\mathsf{ActiveKeys}_{\mathsf{In}(C)}\to\mathsf{ActiveKeys}_{\mathsf{Out}(C)}$ by topological gate evaluation with the public decryption operations, and the \textbf{output decoding} $\mathsf{decode}_{\vec H}:\mathsf{ActiveKeys}_{\mathsf{Out}(C)}\to\{0,1\}^{\mathsf{Out}(C)}$ by the matching position in each ordered hash pair. Theorem~\ref{thm:constructed-garbling} proves that each gate evaluation has exactly one successful row and that these maps are total on the displayed carriers.
\end{definition}

\begin{definition}[The garbling observation family]\label{def:crypto-garbling-family}
Let $\mathcal P_C$ be the finite product of the four-row permutation sets of the gates of $C$, and let $\Phi_{\vec x,\pi}$ be the published frame constructed above. In the common public handle alphabet put
\[
 \mathsf{View}_C(\vec x)
   =\{[\Phi_{\vec x,\pi}]_s:\pi\in\mathcal P_C\}.
\]
Thus $\mathsf{View}_C$ records a family of static frame comparisons; Scott observations remain opens of the execution domain. Equivalently it is the relation of possible static observations obtained as the image of the span with witnesses $(\vec x,\pi)$ and endpoints $\vec x,[\Phi_{\vec x,\pi}]_s$.
Let $\mathcal S_C$ be the set of static-equivalence classes of frames on this common public handle alphabet. The \textbf{garbling view family} has type $\mathsf{View}_C:\{0,1\}^{\mathsf{In}(C)}\to\Pow(\mathcal S_C)$.
\end{definition}

\begin{definition}[Garbled simulator family]\label{def:garbled-simulator-family}
For $\vec y\in\{0,1\}^{\mathsf{Out}(C)}$, define a \textbf{simulator frame} as follows. Allocate private pairs $(L_w,U_w)$ for each wire. Its selected input keys are $L_w$. At every gate construct four rows with input-key indices in $\{L,U\}^2$, using the same nested-encryption pattern, but put $L_v$ as the plaintext key in every row. Choose an arbitrary hidden permutation of these rows. For output bit $y_w$, put $H(\iota L_w)$ in decoding position $y_w$ and $H(\iota U_w)$ in the other position. All its nonces are fresh and distinct.
Let $\mathsf{SimView}_C(\vec y)$ be the set of static-equivalence classes of these frames over all the finite row permutations. Thus $\mathsf{SimView}_C:\{0,1\}^{\mathsf{Out}(C)}\to\Pow(\mathcal S_C)$, with $\mathcal S_C$ as in Definition~\ref{def:crypto-garbling-family}.
\end{definition}

\begin{definition}[Normalized garbling values]\label{def:garbling-normal-values}
Fix an input $\vec x$ and a row choice $\pi$, and let $a_w$ be the
evaluated wire bits. Index a message marker $\kappa_c$ by each distinct
accessible honest ciphertext $c$ whose outer encryption key is an
inactive wire key. Index a hash-result marker $\eta_w$ by each
published $H(\iota K_w^{1-a_w})$. Repeated copies of one value use
the same marker. The normalized sorted carrier is the closure of
the normalized published values and public constants under the
existing evaluated public operations, with these extra sorted roots.
Normalization $N$ replaces such a ciphertext without traversing it,
replaces the specified hash, and otherwise recurses through its
constructor operands. Expansion $E$ substitutes the indexed original
value for each marker and is the identity on retained atoms and
public constants.

On normalized values, projection and checked coercion use their
ordinary root cases. Decryption at $\kappa_c$ fails; decryption at
a retained encryption uses its key equality and returns its retained
plaintext. Hash markers have no inverse operation. Equality compares
normalized values, including marker identity. These clauses define
the public-operation diagram on the normalized carrier.
For the simulator replace active keys by $L_w$, inactive keys by
$U_w$, and use the same definitions at its published positions.
\end{definition}

\begin{lemma}[Nonderivability and normalized expansion]\label{lem:garbling-normal-expansion}
No inactive wire key or honest row nonce is obtainable as a bare
value by a public recipe. Expansion is an isomorphism from the
normalized accessible algebra to the actual accessible algebra,
fixing public constants and preserving and reflecting every public
operation's domain and result.
\end{lemma}
\begin{proof}\vmode
\leavevmode
\begin{enumerate}[label=(\roman*)]
\item \textbf{First-production argument.}
Consider the first producing operation for a purported inactive key.
Public allocation has a distinct owner/occurrence tag; constructors
and key derivations have their own roots; hash has no inverse.
Projection or copying reduces to an earlier producing occurrence.

\item \textbf{Inactive-key obstruction.}
The remaining possibility is an inner row decryption.
A row reached
with both active input keys returns the active output key; every
other row requires an inactive decryption key before it can produce
its plaintext.
This contradicts first production.

\item \textbf{Repeated inputs and fan-out.}
If the two input
wires coincide, the only row with both available keys is still its
active diagonal row.
Fan-out repeats a known key and creates no new
key-producing operation.

\item \textbf{Nonce nonderivability.}
A nonce can only be introduced by its
allocator or retained as a constructor argument; no exposed row
destructor returns the encryption nonce.
The same first-production
argument proves its nonderivability.
It applies equally to $U_w$
in the simulator.

\item \textbf{Separation of marker expansions.}
Distinct honest ciphertext occurrences have distinct fresh nonce
tags and therefore distinct expanded values.
An accessible
nonmarker encryption cannot equal the expansion of an opaque marker:
forming that root would require its inactive key and its particular
honest nonce as operands.
An accessible nonmarker hash cannot equal
the inactive-key hash marker, since that would require the inactive
key.
Thus marker expansions are distinct from every competing
accessible root expression.

\item \textbf{Injectivity of expansion.}
Comparing two expanded values by their
roots now reduces equality either to equality of the marker indices
or to equality of corresponding operands.
Induction on finite
constructor comparisons proves injectivity of $E$.
This argument
uses the non-generation calculation, not merely tree size.

\item \textbf{Constructor and destructor equations.}
Introductions commute with expansion by its constructor definition.
A projection or coercion either has the same retained root and
result or rejects on both sides.

\item \textbf{Decryption domains.}
At a retained encryption, key
equality is reflected by injectivity, so decryption has the same
domain and result.
At a ciphertext marker its required key is
unavailable, so every public decryption fails on both sides.
Other protected operations use the same root and key tests.

\item \textbf{Hashing, data and equality tests.}
Hashing has no new elimination; data reading rejects markers of
nondata sort and returns the same ordinary data otherwise.
Equality is exactly injectivity of $E$.

\item \textbf{Surjectivity on generated values.}
Inducting on recipes proves
$E(N(v))=v$ for every accessible result and that every normalized
generated result expands to that result.
This proves surjectivity
as well as all the operation-domain equations.
\end{enumerate}
\end{proof}

\begin{lemma}[Row-position comparison]\label{lem:garbling-row-permutation}
View a gate's row permutation as a bijection from its four published
positions to $\{0,1\}^2$. Define
\[
 r_{v,\vec x}(a,b)=(a\mathbin{\oplus}a_u,b\mathbin{\oplus}a_z),
 \qquad
 (\Theta_{\vec x}\pi)_v=r_{v,\vec x}\circ\pi_v.
\]
Then $\Theta_{\vec x}:\mathcal P_C\to\mathcal P_C$ is a bijection
with inverse given by the same coordinatewise xor. There is an
isomorphism of normalized labelled accessible algebras
\[
 h_{\vec x,\pi}:
 N(A_{\Phi_{\vec x,\pi}})
 \xrightarrow{\cong}
 N(A_{\mathsf{Sim}(C(\vec x),\Theta_{\vec x}\pi)})
\]
fixing every public handle and position.
\end{lemma}
\begin{proof}\vmode
\leavevmode
\begin{enumerate}[label=(\roman*)]
\item \textbf{Row permutation.}
Xor with a fixed pair is an involution, proving the permutation claim.

\item \textbf{Assignments on normalized generators.}
Map each active wire key to $L_w$, each retained encryption nonce
to the nonce at its corresponding gate and row position, and each
opaque marker to the marker at that position and encryption depth.

\item \textbf{Active and opaque row pattern.}
At a gate, two outer rows are opaque; of the two opened outer rows,
one inner row is opaque and the other returns the active output key.
The xor map preserves exactly this pattern.
Inactive-row plaintext
differences occur below an opaque root and are not fields of the
normalized value.

\item \textbf{Output decoding positions.}
The output decoding position of the active
hash is $a_w=C(\vec x)_w$, and the other position has its corresponding
hash marker.
Thus all exported normalized values correspond.

\item \textbf{Extension and inverse equations.}
Extend through public constructors and operations; their domain
reflection is Lemma~\ref{lem:garbling-normal-expansion}.
The inverse
assignments on keys, nonces and markers give both inverse equations.
\end{enumerate}
\end{proof}

\begin{theorem}[Symbolic garbling correctness and output-only frame privacy]\label{thm:constructed-garbling}
The published bundle computes $C(\vec x)$ by public recipes. For any two inputs $\vec x,\vec x'$ with $C(\vec x)=C(\vec x')$, the families of published frames obtained by the hidden row permutations are equal up to static equivalence: each permutation choice on one side has a matching choice on the other. This matching extends through adaptive public evaluator calls. There is an explicit simulator family $\mathsf{SimView}_C$ with
\[
 \mathsf{View}_C=\mathsf{SimView}_C\circ C.
\]
Its construction uses only $C$ and the output $C(\vec x)$.
\end{theorem}
\begin{proof}\vmode
\leavevmode
\begin{enumerate}[label=(\roman*)]
\item \textbf{Evaluation correctness.}
In topological order, the two active input keys open exactly the
row indexed by their bits and return the active output key; this
also holds when the input wires coincide.
An output key's hash
matches exactly its decoding position.
This proves correctness
with arbitrary fan-out.

\item \textbf{Composite frame comparison.}
Lemma~\ref{lem:garbling-normal-expansion} proves nonderivability
of inactive keys and row nonces and constructs the two normalized
expansion isomorphisms.
Lemma~\ref{lem:garbling-row-permutation}
gives the intervening labelled isomorphism $h_{\vec x,\pi}$.
Their composite compares the real frame with
$\mathsf{Sim}(C(\vec x),\Theta_{\vec x}\pi)$, preserving every
exported handle, row position, operation domain and test result.

\item \textbf{Equality of observation families.}
Since $\Theta_{\vec x}$ is a bijection, both inclusions hold:
\[
 \mathsf{View}_C(\vec x)=\mathsf{SimView}_C(C(\vec x)).
\]
Proposition~\ref{prop:observation-family-descent} identifies this
equation with pullback of the simulator observation relation along
$C\times1$.

\item \textbf{Output-only simulation.}
The simulator's recipe in
Definition~\ref{def:garbled-simulator-family} uses only the circuit
and its output; $\vec x$ occurs only in the comparison bijection.

\item \textbf{Adaptive evaluator comparison.}
The entrywise action of this isomorphism lifts by
Construction~\ref{con:evolving-comparison-coalgebra};
Lemma~\ref{lem:comparison-command-induction} gives all adaptive
reference histories, including failed calls and equality tests.

\item \textbf{Initialization and causal certificate.}
Initialization exposes no intermediate value and has the same
allocation/publication pattern.
Matching those boundary events and
the subsequent call--return positions supplies the finite causal
certificate.

\item \textbf{One-shot interface boundary.}
The one-shot boundary excludes a second input-key
selection and private-key disclosure, exactly as specified in the
bundle construction.
\end{enumerate}
\end{proof}

\[
\begin{tikzcd}[column sep=large]
\{0,1\}^{\mathsf{In}(C)}\arrow[r,"\mathsf{encode}_{\vec K}"]
 \arrow[d,"C"']&\mathsf{ActiveKeys}_{\mathsf{In}(C)}
 \arrow[d,"\mathsf{evaluate}_{\vec T}"]\\
\{0,1\}^{\mathsf{Out}(C)}
 &\mathsf{ActiveKeys}_{\mathsf{Out}(C)}\arrow[l,"\mathsf{decode}_{\vec H}"]
\end{tikzcd}
\]
All wires in the diagram are finite tuples, and correctness proves its commutativity for each constructed garbling. An OT resource can supply an evaluator's selected input key from a pair; substituting those outputs into the same evaluator preserves the correctness equation. The theorem above proves privacy of the fixed-input garbling view. It does not by itself prove a malicious two-party protocol, fairness, or a composable realization of OT. Li--Micciancio~\cite{crypto-garbling} develop symbolic garbled circuits and computational soundness for a specified class of protocols; that additional comparison has its own content. In particular nondeterministic matching of our finite row choices is not a probability-preserving coupling theorem.

\begin{construction}[One-shot garbling view resource]\vmode\label{con:garbling-view-resource}
\leavevmode
\begin{enumerate}[label=(\roman*)]
\item \textbf{Real worker.}
The real worker receives a private input $\vec x$, computes the finite circuit, allocates its private garbling, and publishes the bundle of Definition~\ref{def:crypto-garbling-family}.
Thereafter it offers the public evaluator on that bundle.

\item \textbf{Ideal worker and simulator.}
The ideal worker receives the same private input and supplies only $C(\vec x)$ to its simulator after that input occurrence.
The simulator constructs the frame of Definition~\ref{def:garbled-simulator-family}, chooses a hidden row permutation at each gate, and supplies the same bundle-handle and evaluator interface.

\item \textbf{Initialization and evaluator scheduling.}
The private initialization programs are sequential before publication; after publication their evaluator commands have the same bank and scheduling discipline.
\end{enumerate}
\end{construction}
\begin{corollary}[Garbled view simulation and its image]\label{cor:garbling-image}
The one-shot garbling view resource has the constructed simulation square; Corollary~\ref{cor:security-images} gives its attainable-view-image inclusion. Its public evaluation returns the ideal output $C(\vec x)$.
\end{corollary}
\begin{proof}\vmode
\leavevmode
\begin{enumerate}[label=(\roman*)]
\item \textbf{Matching hidden choices.}
The bijection $\Theta_{\vec x}$ in
Lemma~\ref{lem:garbling-row-permutation} matches the real and simulated
hidden choices in both directions.

\item \textbf{Finite causal certificate.}
Theorem~\ref{thm:constructed-garbling}
lifts each match to the evaluator comparison, and both initialization
programs expose only the same input and publication dependencies.
These matches form a two-sided finite causal certificate for the
one-shot resources.

\item \textbf{Simulator information.}
The simulator program uses the ideal output;
only the certificate uses $\vec x$.

\item \textbf{Security image and correctness.}
Apply
Theorem~\ref{thm:observational-simulation} and
Corollary~\ref{cor:security-images}.
The public evaluation equation
gives correctness.
This is qualitative matching of choices.
\end{enumerate}
\end{proof}


The garbling simulator depends only on the actual ideal output and is uniform on the output fibres of $C$. Its comparison witness may additionally refer to the real hidden row choices and input. Theorem~\ref{thm:observational-simulation} turns the constructed two-sided causal comparison into a trace-certificate judgement. Corollary~\ref{cor:relational-privacy-transfer} then transports equality on output fibres to the one-shot real view, at the declared interface without additional corruption or input-selection ports.

Applying a permitted open test $U$ at the garbled circuit's released output observes $C^{-1}U$ on the input. The simulator uses that same released output. Postcomposing the comparison with the authorized readout therefore preserves it for every $U$ and preserves the single simulator used for all these tests. The row, key and decoding commands remain those of the stated public evaluator.

\subsection{Phased disclosure and private choice}\label{subsec:phased-disclosure}
\subsubsection{Hashing, one-wayness, and a commitment construction}\label{subsec:crypto-hash-commit}
A disclosure phase changes the public observation map. Commitment
binding is a monomorphism calculation; hiding is a comparison of
its closed-phase public presentations. Blindness instead preserves
observations under the permutation of two hidden assignments.
The phase transitions are explicit controller operations.

\begin{definition}[Free symbolic hash]\label{def:formal-free-symbolic-hash}
Adjoin the constructor $H:T_{\mathsf M}\to T_{\mathsf M}$ to the polynomial defining finite values. Its image is tagged and it has no root destructor. There is no fixed output bit length. The public equality test is equality of finite trees; therefore $H(u)=H(v)$ if and only if $u=v$. This is the free symbolic hash used, for example, in the message theories of~\cite{crypto-tamarin}. It is not the assertion that a compressing bitstring function is injective.
\end{definition}

\begin{proposition}[Exact symbolic non-inversion]\label{prop:hash-noninversion}
Allocate a nonce $n$ in a private worker, export only $H(\iota_{\mathsf N}n)$, and close its allocation and storage ports. No public recipe over this export and the opponent's own allocations returns $n$ or its injection. The statement remains true with the original encryption, pairing, signature and MAC operations.
\end{proposition}
\begin{proof}\vmode
\leavevmode
\begin{enumerate}[label=(\roman*)]
\item \textbf{Protection invariant.}
In every source value, each occurrence of $n$ lies strictly below an $H$-node.
The opponent's fresh allocations are different by their instance/owner coordinates and increasing counters.

\item \textbf{Preservation by public operations.}
A constructor cannot remove this protection.
A successful original destructor selects an operand of a pair or an encryption, or returns an observable test; it never removes an $H$-node.
Checked injections remove their own outer tag, not a hash.
These cases preserve the invariant by induction on successful recipes.

\item \textbf{Exclusion of the protected name.}
A bare $n$, and its injection, violate the invariant.
Tests may reveal an equality about a guessed value but cannot insert a value.
The protected values form the closed invariant of Proposition~\ref{prop:crypto-information-invariants}; hence neither is derivable.
\end{enumerate}
\end{proof}

The known-message test $w=H(\mathsf{data}(d_0))$ distinguishes hashes of distinct public data $d_0,d_1$. Non-inversion is consequently weaker than static secrecy. This distinction is central in~\cite{crypto-hash-soundness}; its computational interpretation uses random oracles and is not invoked here.

\begin{definition}[Commitment process]\label{def:hash-commitment}
The committer receives a message $m$, allocates a fresh private nonce $r$, computes $c=H(\langle m,\iota_{\mathsf N}r\rangle)$, exports $c$, and retains $(m,r)$. Opening exports that pair. The public checker has type
\[
 \mathsf{opencheck}:T_{\mathsf M}\times T_{\mathsf M}\times T_{\mathsf N}
 \longrightarrow\{\mathsf{true},\mathsf{false}\},
\]
and returns whether $c=H(\langle m,\iota_{\mathsf N}r\rangle)$. A session cell permits one initial commitment and one opening, rejecting all additional state-changing requests. Its state is $\mathsf{empty}$, $(c,m,r,\mathsf{closed})$, or $(c,m,r,\mathsf{opened})$; an invalid request leaves the state unchanged and replies with failure.
\end{definition}

\begin{theorem}[Symbolic commitment correctness, binding, and pre-opening hiding]\label{thm:hash-commitment}
The displayed process opens its commitment successfully. A commitment has at most one opening pair. For either of two fixed messages supplied before its private nonce is allocated, the public pre-opening service has the same observations, up to the unchanged interface handles. These statements require no distribution on nonces.
\end{theorem}
\begin{proof}\vmode
\leavevmode
\begin{enumerate}[label=(\roman*)]
\item \textbf{Correctness and binding.}
The opening map $e(m,r)=H(\langle m,\iota_{\mathsf N}r\rangle)$
is monic, since each introduction in its composite is injective.
The opening checker is its graph-membership test.
Thus honest
opening succeeds and each commitment has at most one opening,
including commitments chosen by an adversary.

\item \textbf{Commitment marker and nonce protection.}
For hiding, substitute a fresh message marker $\kappa$ for the exported commitment.
Every occurrence of the private $r$ lies under that hash, and no public operation opens the hash, so the non-inversion induction proves that $r$ is not available.

\item \textbf{Injective expansion.}
A public value built without the marker cannot equal its expansion: constructing that hash would require $r$; data introduction cannot introduce a hash tree.
A proper constructor expression containing the marker cannot expand to the commitment itself by finite tree size.
Thus expansion is injective on generated public values.

\item \textbf{Preservation of public operations.}
All original destructors and the hash operation commute with expansion.
The opening checker returns false for every available candidate nonce, on both sides.
Equality follows from injectivity.

\item \textbf{Static frame comparison.}
These calculations construct the two expansion arrows of a common public presentation.
Proposition~\ref{prop:crypto-public-comparison} gives static equivalence.

\item \textbf{Interactive state comparison.}
Each public command uses the same handle counters and the same instruction branch; replacing its calculated value by the expansion preserves its reply.
The state relation retaining this expansion and the identical flags and counters is a witnessed service comparison; Theorem~\ref{thm:service-behaviour-canonical-comparison} gives the interactive equality.

\item \textbf{Pre-opening boundary.}
The commit and opening flags and their call--return order are the same; the comparison is for the pre-opening interface, where no opening has yet been delivered.
Opening deliberately publishes $m,r$ and is not included in the hiding assertion.
\end{enumerate}
\end{proof}

\[
\begin{tikzcd}[column sep=large]
 &A_\kappa\arrow[dl,"e_{m_0}"',"\cong"]\arrow[dr,"e_{m_1}","\cong"']&\\
A_{c_{m_0}}\arrow[rr,"e_{m_1}e_{m_0}^{-1}"',"\cong"]&&A_{c_{m_1}}
\end{tikzcd}
\]
The construction specifies binding and hiding separately. It has not introduced an equivocation trapdoor, a public witness extractor, or a computational commitment scheme. Laud~\cite{crypto-laud} develops a symbolic implementation calculus and constructs encryption from hashing and XOR. His encryption interface exposes its formal randomness and includes a ciphertext test; it is therefore not silently substituted for the different wrapper of Section~\ref{subsec:crypto-channels}.

\subsubsection{Blind signatures as a typed service}\label{subsec:crypto-blind}
\begin{construction}[Blind signing and checked publication]\vmode\label{con:blind-signing-checked-publication}
\leavevmode
\begin{enumerate}[label=(\roman*)]
\item \textbf{Sorted carriers and introductions.}
Use distinct sorts $B_{\rm n},B_{\rm sk},B_{\rm msg},B_{\rm resp},B_{\rm sig},B_{\rm pk}$.
The introductions $\mathsf{blind}(m,b)$, $\mathsf{bsign}(u,k)$, $\mathsf{final}(m,k)$ and $\mathsf{bpk}(k)$ form respectively a blinded message, its response, a final attestation and a public-key image.

\item \textbf{Opaque carriers.}
The blinded-message and response carriers have no public component parser.

\item \textbf{Checked unblinding.}
Pulling back a response's retained blinding name along the supplied name constructs
\[
 \mathsf{unblind}(\mathsf{bsign}(\mathsf{blind}(m,b),k),b)
 =\mathsf{final}(m,k),
\]
with failure off that domain.

\item \textbf{Final verification.}
Final verification is the checked issuance image of Construction~\ref{con:protected-records}; the public introductions require their actual operands, including a signing-key operand for $\mathsf{final}$.

\item \textbf{Two-session client state.}
The two-session client is a concrete worker network.
It retains the designated public key $p=\mathsf{bpk}(k)$ and the two fixed public messages $m_0,m_1$.

\item \textbf{Session execution.}
Session $i$ allocates a fresh private blinding name $b_i$, publishes its request, accepts one response, and unblinds it privately.
It then verifies the result against its retained message and $p$.
A failed unblind or verification records failure; a success stores the result in a private completion cell.

\item \textbf{Joint publication.}
A publication worker waits for both successful cells and publishes the two final signatures in the fixed order $m_0,m_1$, with no session labels.
A failed or unfinished session yields no final publication.
The individual checks have no public reply port.

\item \textbf{Publication dependencies and realization.}
Publication has both completion events as prerequisites, so selective completion creates no earlier one-message publication.
The request, response and completion instructions are the existing call--return workers and private connections; Theorem~\ref{thm:compact-cryptographic-communication} realizes their network as a strategy.
\end{enumerate}
\end{construction}

The group $S_2$ acts on the two assignments of messages to request
positions. The following comparison proves invariance of the public
behaviour under that action. Its witness preserves session positions
and the joint-publication prerequisites; it is not a claim that
every private controller state is identified with its permutation.

\begin{proposition}[Blind-signature equations and two-session blindness]\label{prop:blind-signatures}
Blind, sign and unblind produces a valid final signature. Every observed final signature under a private key has a direct signing-key constructor ancestor or an unblinding ancestor from a response under that key. A response has at most one successful blinding-name argument. The checked two-session worker hides the assignment of $m_0,m_1$ to the two requests from an interacting signer, including the final publication.
\end{proposition}
\begin{proof}\vmode
\leavevmode
\begin{enumerate}[label=(\roman*)]
\item \textbf{Correctness.}
The introduction and checked-elimination equations prove correctness.

\item \textbf{Origin and unique unblinding argument.}
Follow the finite supplying audit through copy and transport to its first producing operation.
A final signature is introduced directly with its key or by a successful unblind, whose response fixes its key and blinding name.
This proves origin and uniqueness of the unblinding argument.

\item \textbf{Session markers.}
Replace the requests at their session positions by markers $u_i$, expanding them in experiment $b\in\{0,1\}$ as $\mathsf{blind}(m_{b\mathbin\oplus i},b_i)$.
Here $b\mathbin\oplus i$ is addition modulo two.

\item \textbf{Injective public expansion.}
The private names are neither exported nor returned by any public destructor.
Constructor injectivity and induction on public computations show that expansion is injective and commutes with exposed operations: an available ordinary blinding name cannot open a marker or its signed response, and a newly constructed blinded message cannot equal a marker.

\item \textbf{Successful response criterion.}
A response can succeed in session $i$ only if it is $\mathsf{bsign}(u_i,k')$ for some signing key $k'$.
The subsequent verification against $p$ succeeds precisely when $\mathsf{bpk}(k')=p$, hence $k'=k$.
This condition depends on the session marker and the designated key, not on its hidden message assignment.

\item \textbf{Failure and publication results.}
All other responses give the same failure state in the two experiments.
Once both checks succeed, publication is exactly $(\mathsf{final}(m_0,k),\mathsf{final}(m_1,k))$ in both.
Before that joint event neither final signature is public.

\item \textbf{Comparison state and transitions.}
Retain these marker expansions, response flags and private completion cells as the comparison state.
Each public command and each response extends it with the same observable result.

\item \textbf{Finite causal certificate.}
The publication dependencies are the same pair of completion events, including histories with one response pending or rejected.
The identity on these event positions gives a two-sided finite causal certificate, so Lemma~\ref{lem:crypto-causal-congruence} proves the interactive comparison.

\item \textbf{Operational implementation of the checks.}
The fixed-key checks and joint publication are actual worker transitions; fixing final outputs in an initial frame alone would not prove this result.
\end{enumerate}
\end{proof}
This is the typed symbolic comparison at the constructed publication interface, not a computational one-more unforgeability theorem; compare~\cite{crypto-interactive}.

\subsubsection{Example: an ideal oblivious-transfer game}\label{subsec:crypto-ot}

\begin{definition}[Ideal oblivious-transfer carrier]\label{def:formal-ideal-oblivious-transfer-carrier}
The private selection map is
$\mathsf{sel}:D^2\times\{0,1\}\to D$, $(\vec m,c)\mapsto m_c$.
Its graph supplies the positive output events below, while each
party's projection specifies its observation boundary. For message
set $D$, take independent negative sender alternatives $a_{m_0,m_1}$ and receiver alternatives $b_c$, $c\in\{0,1\}$. Distinct sender alternatives conflict, as do $b_0,b_1$. For each triple $(m_0,m_1,c)$ add a positive event $z_{m_0,m_1,c}$ with causes exactly $a_{m_0,m_1}$ and $b_c$. Conflicts are inherited from these inputs. Display its output at the receiver interface as $m_c$, not as the input pair or as the triple naming the strategy event.
\end{definition}

\begin{proposition}[Ideal OT strategy]\label{prop:crypto-ot}
This display is a receptive, courteous strategy on the product of the one-shot input interfaces and the output interface. The two inputs are genuinely concurrent. Sender observation contains its own pair and no choice event; receiver observation contains its own choice and selected result, not the other input component.
\end{proposition}
\begin{proof}\vmode
\leavevmode
\begin{enumerate}[label=(\roman*)]
\item \textbf{Receptivity and courtesy.}
Every consistent choice of a sender input and receiver input is accepted, with a unique negative lift.
The only new causal edges are negative-to-positive, proving courtesy.

\item \textbf{Output dependency and value.}
An output cannot occur before both causes, and its label is precisely the selected component.

\item \textbf{Party observations.}
There is no order between the two negative inputs.
Projection to each party's boundary gives the stated observations; the strategy-event identifier is not itself a public payload.

\item \textbf{One-use cell realization.}
Equivalently, a one-use cell stores the two inputs and returns the selected value once, rejecting subsequent use in a reusable command presentation.
\end{enumerate}
\end{proof}

\begin{proposition}[Ideal selection of an observable payload]\label{prop:open-ideal-selection}
For an algebraic payload domain $D$ and the discrete choice bit, the ideal selection operation is continuous. Its receiver can be given a handle for the selected payload after the two input occurrences. At the observation interface its positive replies are exactly the requested basis observations of that payload.
\end{proposition}
\begin{proof}
For $U\in\Op(D)$ the inverse image under selection is
\[
 (U\times D\times\{0\})\ \cup\ (D\times U\times\{1\}).
\]
Equivalently, on each clopen choice component the selector is a coordinate projection. Retain the concurrent one-shot inputs of Proposition~\ref{prop:crypto-ot}, carrying the private pair and the bit. Replace the selected-value output by a receiver-only ready notification with the same two causes. In the continuation after ready, attach the state $P_{m_c}$ of Construction~\ref{con:positive-observation-port}, with its distinct lanes. A request has its unique lift in the active continuation. A basis acknowledgement is available precisely after that request and exactly when its compact lies below $m_c$. The guards for the ready occurrence are the original two input occurrences, even for the constant test $U=D$. Projection to the sender still gives just its own input. Varying the unselected component preserves every receiver-lane event and all its ordering and consistency relations. Thus the receiver comparison also holds at the observation interface. The original discrete, single-output port is recovered by atomic delivery of the singleton token as in Proposition~\ref{prop:discrete-observation-recovery}.
\end{proof}

This proves existence of the ideal game and its information interfaces. It is not a construction of a two-party protocol using only communication and local computations.

\subsubsection{Oblivious transfer and the realization boundary}\label{subsec:crypto-ot-boundary}

Proposition~\ref{prop:crypto-ot} constructs an ideal OT game. To see why a concrete realization requires more, consider a receiver that generates two key pairs, is instructed to erase the key for its unchosen bit, and receives one encrypted message under each public key. A malicious receiver keeps both keys and uses both decryption equations. The erasure instruction therefore cannot justify security against arbitrary receiver strategies.

Ambient compact closure also has a specific limitation here. In the compact-closed communication-only two-party setting of Broadbent--Karvonen, secure realization forces a splitting equation that OT fails~\cite[Theorem 8.1 and Corollary 8.2]{crypto-bk}. Their obstruction applies to that resource basis. It neither excludes every additional cryptographic resource nor decides whether the symbolic services constructed above suffice.

Equational OT presentations can support the analysis of protocols that use OT~\cite{crypto-bursuc}. They do not themselves realize OT from the constructors of Section~\ref{subsec:crypto-values}; the extended report also leaves the connection to simulation-based security for further work.

Here we have constructed the OT functionality and the channel realizations of Theorem~\ref{thm:crypto-channels}. A distributed OT realization from these symbolic services remains open. Providing a trusted selector would add a resource, rather than derive OT from the displayed encryption, signature and MAC operations.

\subsection{Scope of the established comparisons}\label{subsec:crypto-probability-boundary}
All constructions here use exact operations, formal allocation, finite
audits and qualitative Scott observation. They introduce no probability
distribution or complexity bound. In particular the fresh memoized
oracle theorem is an exact symbolic comparison, and the certificate
theorem extracts from an actual private issuance audit. Computational
pseudorandomness and extraction from arbitrary proof bitstrings are
different realization questions. The concrete distinctions are recorded
in Appendix~\ref{app:quantitative-comparison}; they are not additional
foundations for the present examples.

\paragraph{Coverage and comparison.}
The following table identifies the conclusion established for each construction and the relevant literature comparison. The following discussion then relates these conclusions to the general simulator criterion.
\begingroup\renewcommand{\arraystretch}{1.08}
\begin{longtable}{@{}>{\raggedright\arraybackslash}p{0.23\textwidth}>{\raggedright\arraybackslash}p{0.49\textwidth}>{\raggedright\arraybackslash}p{0.22\textwidth}@{}}
\toprule
Structure & Constructed result in this manuscript & Literature comparison\\
\midrule\endfirsthead
\toprule Structure & Constructed result in this manuscript & Literature comparison\\\midrule\endhead
Public-key and symmetric encryption; signatures; MACs & Typed free operations, handles, provenance, and the four explicit channel simulations of Section~\ref{subsec:crypto-channels}. & \cite{crypto-abf,crypto-bu}\\
Hashes and commitments & Non-inversion, unique opening, and pre-opening hiding by an accessible-algebra isomorphism. & \cite{crypto-hash-soundness,crypto-laud}\\
Key encapsulation and AEAD & Explicit encryption-based encapsulation and encrypt-then-MAC algorithms, checked correctness and origin/secrecy properties at their declared interfaces. & \cite{crypto-kems}\\
Keyed functions and derivation & Memoized fresh-key oracle comparison and deterministic derived-key terms; no claim of computational pseudorandomness. & \cite{crypto-micciancio,crypto-ggm}\\
DH, pairings and XOR & Concrete rational-function groups, bilinear equations, a DH non-derivability calculation, and a separate finite-support XOR algebra. & \cite{crypto-tamarin}\\
Threshold and verifiable sharing & Lagrange reconstruction, online threshold-disclosure simulation, and commitment/certificate consistency. & \cite{crypto-shamir,crypto-vss}\\
Homomorphic encryption & A free circuit ciphertext carrier, decryption fold, evaluation correctness and hidden-key static privacy. & \cite{crypto-dahl}\\
Zero-knowledge proof resources & Checked issuance, audit extraction, witness-erasure coalgebra map, conjunction and existential hiding, with selective disclosure and computation certificates. & \cite{crypto-zk-applied,crypto-zk-types,crypto-zk-malleable}\\
Blind signatures & Typed unblinding, provenance, and two-session blindness at the specified unlinking interface. & \cite{crypto-interactive}\\
MPC, OT and garbling & Linear distributed computation at observation ports; generic circuit ideal resource; ideal OT functionality; concrete garbling evaluation and output-only frame privacy. & \cite{crypto-mpc,crypto-bursuc,crypto-garbling}\\
\bottomrule
\end{longtable}
\endgroup
\subsubsection{Refinement conclusions of the examples}
Simulation squares compare the specified game strategies at their actual boundaries. The channel and associated-data constructions supply a simulator equation for their delivery interfaces. The threshold and garbling constructions supply simulators for their defined view resources. The other local proofs have the following precise relation to that criterion.
\begin{corollary}[Images of exact service replacements]\label{cor:service-images}
An initialized service comparison of Theorem~\ref{thm:service-behaviour-canonical-comparison}, with its common public interface, gives equal attainable attacker-view images. In particular the proof service and its witness-erased service have equal such images at their public command boundary.
\end{corollary}
\begin{proof}\vmode
\leavevmode
\begin{enumerate}[label=(\roman*)]
\item \textbf{Contextual display comparison.}
The comparison identifies the two public strategy displays.
Tensor with any untouched honest interface and postcompose with any attacker; functoriality preserves that equality.

\item \textbf{Behaviour images and simulator.}
The two behaviour transformations are therefore equal, so their pointwise images are equal.
The identity interface strategy supplies the simulator in both directions.

\item \textbf{Proof-service instance.}
Theorem~\ref{thm:proof-erasure} supplies this comparison for the proof service, including its transformed certificates.
Its private witness-extraction port is outside this public comparison.
\end{enumerate}
\end{proof}
Commitment hiding uses the pre-opening interface; opening changes the public information. Hash non-inversion and key non-derivability establish membership exclusions in a generated value algebra, which are different statements from a resource simulator equation. Signature, MAC, and certificate validity give historical origin implications in the public Scott frame. The ideal OT and generic circuit workers specify target resources; their construction alone is not a realization from a particular communication resource. The linear network retains its read-only coalition interface. Thus the general image theorem is used exactly where a simulator or an exact replacement has been constructed, while each remaining local invariant supplies its stated information or correctness conclusion.

\subsection{Process instances for operations, circuits and certificates}\label{subsec:calculus-extensions}
The following instances connect the algebraic constructions to concrete process commands. Each uses the existing service and substitution interpretation, so its process presentation carries the same comparison witnesses and the same declared conclusion.

These are process presentations of the same realizers and comparison witnesses, not a second security semantics. The shared readback and substitution theorems transport each proved judgement to the named syntax.
\paragraph{Operation and circuit instances.}\label{par:operation-circuit-instances}
Every constructed operation in this section instantiates the existing value rule, strict eliminator, or service call according to its type. The resulting macros have the interpretation and substitution law of Proposition~\ref{prop:generic-process-preservation}; interpretation is inherited from the same calculus.

For the free ciphertext algebra $C_p$, the corresponding introductions are
\[
 \frac{\Gamma\vdash b:B\quad\Gamma\vdash r:\mathsf N}
      {\Gamma\vdash\mathsf{leaf}_p(b,r):C_p},\qquad
 \frac{\Gamma\vdash b:B}{\Gamma\vdash\mathsf{const}_p(b):C_p},
\]
\[
 \frac{(\Gamma\vdash c_i:C_p)_{i=1}^n}
      {\Gamma\vdash\mathsf{gate}_{p,\omega}(\vec c):C_p}
 \quad(\omega:B^n\to B).
\]
At the available matching key the strict eliminator is the unique fold $d_k:C_p\to B$; its computations are the algebra-morphism equations of Theorem~\ref{thm:homomorphic-circuits}. Wrong keys take failure. Circuit substitution is composition in the free finite-product syntax, so fold uniqueness gives its decryption square. This supplies the type-theoretic presentation of the same construction without redefining its evaluator.

\begin{definition}[Certificate and private-audit types]\label{def:certificate-private-audit-types}
For a code $R$ of Construction~\ref{con:checked-certificate-service}, write $W_R(x)=\{w\in W_R:(x,w)\in V_R\}$. At a state $s$ of the constructed proof service, the \textbf{verified certificate type}
\[
 \mathsf{Proof}_{R,s}(x)=
 \{p\in\mathsf{Handle}_s:\mathsf{verify}_s(p,R,x)=\mathsf{true}\}
\]
is the fibre of its public verification map. Here $\mathsf{Handle}_s$ is the chosen port's accessible proof-handle set; lookup, access and sort checks are those of Construction~\ref{con:checked-certificate-service}. Thus $R$, $x:X_R$ and $s$ form the type $\mathsf{Proof}_{R,s}(x):\mathsf{type}_{\Set}$. Its public values are handles, not pairs exposing $w$.

For $p:\mathsf{Proof}_{R,s}(x)$, the \textbf{private audit type} $\mathsf{Audit}_{R,s}(p,x)$ consists of the finite audit of issuance, refresh, conjunction, and hiding, rooted at the record resolved by $p$. Its extracted witness has type
\[
 \mathsf{extract}_{R,s,p,x}:\mathsf{Audit}_{R,s}(p,x)\longrightarrow W_R(x),
\]
constructed by the finite audit fold of Theorem~\ref{thm:proof-erasure}. This map's domain is private audit evidence, and is not the public handle type.
\end{definition}

\paragraph{Derived proof rules.}\label{par:derived-certificate-proof-rules}
Use the existing completed-command judgement $s\vdash c\Downarrow s';r:F(s')$. Applying the checked certificate service and its witness-family maps gives
\[
 \frac{x:X_R\quad w:W_R(x)}
 {s\vdash\mathsf{prove}_R(x,w)\Downarrow s';p:\mathsf{Proof}_{R,s'}(x)},
\]
\[
 \frac{p:\mathsf{Proof}_{R,s}(x)\quad q:\mathsf{Proof}_{S,s}(x)}
 {s\vdash\mathsf{conjoin}(p,q)\Downarrow s';r:\mathsf{Proof}_{R\wedge S,s'}(x)},
\]
\[
 \frac{p:\mathsf{Proof}_{R,s}(u,z)}
 {s\vdash\mathsf{hide}(p)\Downarrow s';r:\mathsf{Proof}_{\exists_ZR,s'}(u)}.
\]
Targets and returned identifiers are computed by those service calls. Rejected checking or lookup takes the existing failure continuation. Refresh is the identity witness map with fresh issuance. Public elimination returns the statement or its verification test; private audit elimination is the fold with node maps identity, $\kappa$ and $\eta$. Thus its computation returns the checked issuing witness, a parent's witness, a conjunction pair, or the pair of hidden statement component and parent witness, respectively.

\begin{corollary}[Sound certificate rules]\label{cor:certificate-rule-soundness}
These introductions and eliminations have their displayed types and computations, and their public interpretation is preserved by witness erasure.
\end{corollary}
\begin{proof}\vmode
\leavevmode
\begin{enumerate}[label=(\roman*)]
\item \textbf{Witness fibres and handles.}
The checker equations identify the target witness fibres; the reference construction supplies the returned handles.

\item \textbf{Audit and public comparison.}
Theorem~\ref{thm:proof-erasure} supplies both the private audit fold and the public service square.

\item \textbf{Transport of certificate rules.}
Apply Proposition~\ref{prop:generic-process-preservation} to transport the rules and their substitutions through that square.
\end{enumerate}
\end{proof}

The proof-service construction can be compared with several symbolic treatments. Backes--Maffei--Unruh~\cite{crypto-zk-applied} develop an applied-$\pi$ account and a Direct Anonymous Attestation analysis; Backes--Hri\c{t}cu--Maffei~\cite{crypto-zk-types} supply proof-oriented typing. Malleable proof operations are treated explicitly in~\cite{crypto-zk-malleable}.

The conjunction and hiding maps here are verified service operations. Any further proof transformation must satisfy its checker equation and erasure case. Computational comparisons such as~\cite{crypto-zk-soundness} concern particular cryptographic realizations and do not follow from the proof-object syntax alone.


\subsection{Example: administrative resource histories and move rules}\label{subsec:admin-states}
\begin{construction}[Signed grants and owned stages]\vmode\label{con:signed-grants-owned-stages}
\leavevmode
\begin{enumerate}[label=(\roman*)]
\item \textbf{Indexed record carrier.}
For the finite principal and initialized transaction-key sets $\mathsf{Prin},K$, and resource/action fields $X_r,X_a$, form the tagged sum of grant and delegation records
\[
 \mathsf{Rec}=(\mathsf{Prin}^2\times K\times X_r\times X_a)
       +(\mathsf{Prin}^2\times X_r\times X_a).
\]
The first projection is $\mathsf{issuer}$; write $\mathsf{Rec}_p$ for its fibre.
Denote the records by $(\mathsf{grant},p,q,k,r,a)$ and $(\mathsf{delegate},p,p',r,a)$.

\item \textbf{Protected issuing ports.}
Instantiate the checked attestation construction at one protected issuing port per $p$, with record domain $\mathsf{Rec}_p$.
Its result is written $\mathsf{sig}_p(t)$; the private authenticated-record carrier is $\mathsf{Auth}=\coprod_p\mathsf{Rec}_p$.
Only the protected signing command creates an occurrence, using that principal's signing capability.

\item \textbf{Verification interface.}
Verification checks the issuer and exact record.
This is the local authenticated-record resource; Section~\ref{subsec:admin-crypto} constructs its concrete signature interpretation.

\item \textbf{Public values and initial ledgers.}
The public value carrier $V$ is the free pairing extension on ordinary data $B$ and these authenticated records.
Thus data, pairing, checked projection, copying and transport are inherited from the common value and reference constructions.
Capabilities are owned resources, not data.
Initialize public ledgers and channels by finite data/pair preparation, so their audits contain neither signatures nor protected replies.

\item \textbf{Registry and authority graph.}
At a decision view $v$, the service registry supplies recognized roots $\mathsf{Root}_{q,v}(r,a)$ and revoked delegation-record values.
Verified, unrevoked delegation occurrences are edges of the authority multigraph $G_v$, from their verified signer to the delegated issuer; copies remain distinct edges.
Registry replacement constructs the next view.

\item \textbf{Selected authority paths.}
The authority family is the free-path family of~\eqref{eq:admin-authority-paths}, rooted at those recognized principals.
A release retains its actual consulted view and selected authority path.

\item \textbf{Owned-stage program.}
Initialize one named $\mathsf{Fresh}(k,q,r,a)$ occurrence for each transaction.
Its owned-stage program is
\[
 \mathsf{Fresh}\xrightarrow{\mathsf{begin}}\mathsf{Pending}
 \xrightarrow{\mathsf{accept}}\mathsf{Available}
 \xrightarrow{\mathsf{release}}\mathsf{Spent}.
\]
Acceptance consumes the pending occurrence with a selected verification of the matching grant.
Release consumes its available successor with an authority path at the current view and an actual protected-effect edge; their precise dependent operands are~\eqref{eq:admin-authority-paths} and the command fibres below.

\item \textbf{Protected release reply.}
Release also constructs the immutable protected reply $(e,t,b,v,\gamma)$, retaining its release parent, grant, computed result, view and authority path.
Its public result is $b$.

\item \textbf{Stage exclusivity and duplicate reports.}
No other command advances these stages or recreates an earlier one.
Duplicate acceptance may report its retained record but changes no inventory.

\item \textbf{Revocation update.}
Revoking a delegation removes all occurrences of that record from a later authority graph while preserving the cumulative audit and earlier frozen replies.
\end{enumerate}
\end{construction}

\begin{construction}[Administrative process interpretation]\vmode\label{con:administrative-process-interpretation}
\leavevmode
\begin{enumerate}[label=(\roman*)]
\item \textbf{Connected command network.}
Connect the signed-grant commands, local registries and owned-stage processes by the chosen FIFO or rendezvous media.
Use the existing reference banks for locally available values and the operation audit for their complete ancestry.

\item \textbf{Residual state and availability.}
The state $S_{\rm adm}$ is the residual carrier of this network: its projections give controls, live inventories, local availability and current registry versions.
Availability and cumulative history remain distinct.

\item \textbf{Wire payloads and occurrence names.}
Administrative wire payloads remain the declared public values $V$; the banks represent local availability and install the payload of each matched receipt.
Names are the producer/output-position names already constructed.

\item \textbf{Dependent command edges.}
For a command $\rho$, the operand fibre $G_\rho(s)$ below specializes the committed-witness family; its actual update is $u_\rho$.
Consequently its local edge object is
\begin{equation}\label{eq:admin-rule-carrier}
 D_{\rm loc}=\coprod_{s\in S_{\rm adm}}\coprod_\rho G_\rho(s),
 \qquad (s,t)(s,\rho,g)=(s,u_\rho(s,g)).
\end{equation}
Elaboration and connection give $C_{\rm adm}=(S_{\rm adm}\leftarrow D_{\rm adm}\to S_{\rm adm})$.

\item \textbf{Retained operands and outputs.}
Each command retains its selected credentials, local operands, consumed tokens and consulted registry records; remote evidence must first arrive on a wire.
Multiple outputs retain the same producing event.

\item \textbf{Relay and reception rules.}
The relay receives and records a packet, then offers forward or drop; replay sends another occurrence from a retained entry.
External reception is uncontrollable, with rejection or validation in its subsequent local continuation.
Their elaboration supplies the witnessed graph.
\end{enumerate}
\end{construction}

\begin{definition}[Administrative extension and survival maps]\label{def:administrative-survival}
Let $\mathcal H_{\mathrm{adm}}$ have finite witnessed administrative histories as objects and prefix extensions as arrows. Let $\mathbf{PInj}$ be the category of sets and partial injections, composed on their matching domains. Define $\mathsf{Live}(h)$ to be the named inventory at the terminal cut of $h$. For $h\subseteq k$, the \textbf{survival map} $\mathsf{Live}(h)\rightharpoonup\mathsf{Live}(k)$ is the identity on those occurrence names not consumed by $k\setminus h$. Theorem~\ref{thm:admin-residuals} verifies well-definedness under exchange and functoriality.
\end{definition}

\begin{theorem}[Administrative residuals and inventories]\label{thm:admin-residuals}
The administrative rules construct concurrent cells and a principal-history event structure. These agree with the operational residual cubes and unfolding of Sections~\ref{subsec:local-resource-charts} and~\ref{subsec:histories}. An independent family retains its rule witnesses, credential pointers, and resource names under every residual permutation. Its finite histories form an extension category $\mathcal H_{\mathrm{adm}}$, and live-resource survival defines a functor
\[
 \mathsf{Live}:\mathcal H_{\mathrm{adm}}\longrightarrow\mathbf{PInj},
\]
where $\mathbf{PInj}$ has sets and partially defined injective functions.
\end{theorem}
\begin{proof}\vmode
\leavevmode
\begin{enumerate}[label=(\roman*)]
\item \textbf{Independent local updates.}
Independent local rules have disjoint participating components.
They consequently change disjoint local inventories, controls, ledgers, and registries; each selected local guard and predecessor remains available after the other rule.
Distinct produced names remain distinct and unchanged by the exchange.

\item \textbf{Channel exchange.}
For a shared FIFO channel the only additional case is an enqueue and an already enabled dequeue.
The equality in Lemma~\ref{lem:fifo-square} holds for occurrence-valued packets, so it preserves their source pointers as well as their payloads.
All other channel updates commute componentwise.

\item \textbf{Residual cubes.}
Repeating these equalities for a finite family gives the same residual witness at every face and the cubical equations of Proposition~\ref{prop:resource-charts}.

\item \textbf{Participant and supply dependencies.}
The history construction records participant chains and packet matching, including the initial-packet convention of Section~\ref{subsec:histories}.
Input ledger occurrences selected by a rule were already present at its owner's preceding local state or arrived by a matched reception, so their producing events lie in this causal past.
Dependencies between outputs constructed within one step attach to that single producing event; they introduce no extra physical event.
A later rule selecting any such output therefore retains its producer and all required input predecessors.

\item \textbf{Legal history prefixes.}
Every downset has a legal linearization by Lemma~\ref{lem:protocol-history-adequacy}; the enriched labels and pointers are invariant under its exchanges.
The ideals of these finite participant and supply diagrams are therefore the legal prefixes.

\item \textbf{History event structure.}
Lemma~\ref{lem:admin-family} constructs their event structure.
The same diagrams are obtained from the operational unfolding by Proposition~\ref{prop:unfolding}.
Take finite configurations and inclusions as $\mathcal H_{\mathrm{adm}}$.

\item \textbf{Partial survival maps.}
For $h\subseteq h'$, define survival on those members of $\mathsf{Live}(h)$ not consumed in $h'\setminus h$, sending each to the same named occurrence in $\mathsf{Live}(h')$.
Names are never recycled, so this is injective on its domain.

\item \textbf{Identity and composition.}
An occurrence survives $h\subseteq h''$ exactly when it survives both $h\subseteq h'$ and $h'\subseteq h''$.
This proves the composition equation for partial maps; the empty extension gives the identity.
\end{enumerate}
\end{proof}

Disjoint resources alone do not make two actions of one sequential component concurrent. Physical independence also includes the participant discipline above. Two transaction workers with separate controls can redeem disjoint permissions concurrently. A single sequential worker orders its redemptions. The policy category constructed below describes resource derivations; its comparison from physical execution preserves independent exchanges without imposing every policy exchange on a sequential implementation.

\subsection{Example: a signed grant and its historical receipts}\label{subsec:admin-grant}
A signed grant connects cryptographic validation with resource actions. The example retains issuance, acceptance, release and delivery records so that the evidence conclusions can refer to the historical occurrences that justify them.


An observation asserting that a particular finite signing and delegation history has occurred is the principal open of that history in the retained configuration domain. Its witness includes the signature and the authority path. Its visibility to a participant is determined by pullback from that participant's declared observation frame as in Section~\ref{subsec:public-observations}.

\begin{construction}[Dependent grant commands]\vmode\label{con:dependent-grant-commands}
\leavevmode
\begin{enumerate}[label=(\roman*)]
\item \textbf{Value operations and verification evidence.}
Generic value and transport commands use the matching fibres of the reference construction.
In particular a successful verification selects an actual locally available entry $b$ with value $\mathsf{sig}_p(t)$ and $t\in\mathsf{Rec}_p$; retain it as $(p,t,b)\in\mathsf{VG}_t(s)$.

\item \textbf{Registry and duplicate-report updates.}
Registry insertion/removal is the finite-record update at its owner's control.
Duplicate reports append audit evidence without changing owned resources.

\item \textbf{Inventory and authority fibres.}
Let $F_k(s),P_k(s),A_{k,t}(s)$ be the named fresh, pending and available inventory fibres.
The authority family is constructed by the free category on the verified, unrevoked delegation multigraph:
\begin{equation}\label{eq:admin-authority-paths}
 \mathsf{Auth}_v(p)=\coprod_{p_0\in\mathsf{Root}_{q,v}(r,a)}
                         \Path(G_v)(p_0,p).
\end{equation}
Let $\mathsf{Eff}_{q,r,a}(s)$ be the fibre of the local process's protected-effect edges, retaining their target and result.

\item \textbf{Dependent stage commands.}
For the actual grant $t=(\mathsf{grant},p,q,k,r,a)$, form
\[
 \begin{aligned}
 G_{\mathsf{begin}_k}(s)&=F_k(s),\\
 G_{\mathsf{accept}_{k,t}}(s)&=P_k(s)\times\mathsf{VG}_t(s),\\
 G_{\mathsf{release}_{k,t}}(s)&=A_{k,t}(s)\times
                  \mathsf{Auth}_{v(s)}(p)\times\mathsf{Eff}_{q,r,a}(s).
 \end{aligned}
\]
Commitment consumes the selected stage, produces its named successor, and frames the remaining inventory.

\item \textbf{Release evidence and emission.}
Acceptance retains its verification; release performs its selected effect and records $(e,t,\operatorname{result}(e),v(s),\gamma)$.
Its view is computed from the source's consulted records, not chosen independently.
Protected emission transmits the ordinary result and retains this reply as parent.
These are dependent process commands with the result fibre of equation~\eqref{eq:dependent-command-result-family}.

\item \textbf{Finite authority witnesses.}
Cycles in $G_v$ can yield infinitely many distinct finite paths.
The path object is the coproduct of composable-edge pullbacks over finite lengths, so every release witness is finite without requiring finite branching.

\item \textbf{Reachability search.}
To execute a reachability decision on the finite graph, erase loops from a path: a reachable vertex has a path shorter than the number of vertices.
Finite search therefore constructs a witness when one exists; the retained command operand is the path, not its Boolean existence test.
\end{enumerate}
\end{construction}

\begin{lemma}[Signer binding and transport]\label{lem:admin-signer-binding}
Every operationally produced signature component has value $\mathsf{sig}_p(d)$ with $d\in\mathsf{Rec}_p$. Copy, projection, communication, and replay preserve that index along the selected predecessor derivation. An accepted grant naming $p$ consequently has a signing origin at $p$, and an authority edge $p\to p'$ has a delegation-signing origin at its source $p$.
\end{lemma}
\begin{proof}\vmode
\leavevmode
\begin{enumerate}[label=(\roman*)]
\item \textbf{Signer invariant.}
The initial audit has no signature node.
At a signing node the operand lies in $\mathsf{Rec}_p$, hence names $p$ as issuer.

\item \textbf{Induction on the operation audit.}
Every remaining constructor either retains a component or follows a selected operand or transport pointer.
Induction in the audit fold of Lemma~\ref{lem:generic-operation-audit} preserves that index and reaches the signing origin.

\item \textbf{Verification and transport.}
Verification checks this indexed value, so acceptance and each delegation edge retain the claimed source.
Pointer naturality gives preservation by independent exchanges and history extensions.
\end{enumerate}
\end{proof}

\begin{definition}[Administrative occurrence families]\label{def:administrative-occurrence-families}
For a finite administrative history $h$, let $\mathsf{Begin}(h)$, $\mathsf{Accept}(h)$, $\mathsf{Release}(h)$ and $\mathsf{Delivery}(h)$ be the subsets of its events with the respective rule tags. Along a prefix inclusion use the inclusion of the same occurrences. These are \textbf{history-indexed occurrence families} $\mathcal H_{\mathrm{adm}}\to\Set$. Their origin maps return the selected earlier supplier occurrence; Theorem~\ref{thm:admin-origin} and the protected-reply construction establish the indicated maps and their properties.
\end{definition}

\begin{theorem}[Origin and one-use authorization]\label{thm:admin-origin}
Every accepted grant has a signing occurrence at its named issuer of its exact record in its causal past. Every protected release has a preceding acceptance and a selected finite verified authority chain at the release view. Every signature occurrence selected in this chain has a signing origin in that past. For each initialized key there is at most one begin, one acceptance, and one protected release. There are natural injections on finite histories
\[
 \mathsf{Release}\lhook\joinrel\longrightarrow\mathsf{Accept}
 \lhook\joinrel\longrightarrow\mathsf{Begin}.
\]
The chosen signing-origin maps, and these injections, are invariant under independent exchanges.
\end{theorem}
\begin{proof}\vmode
\leavevmode
\begin{enumerate}[label=(\roman*)]
\item \textbf{Signature ancestry traversal.}
Follow the selected occurrence of a signature backwards through its constructor derivation.
A copy, transmission, or reception preserves its predecessor; a projection selects its recorded operand; a pairing keeps the derivations of both components.
A selected signature subterm therefore either has a strictly earlier supplying occurrence or was created by a signing step.

\item \textbf{Signing origin.}
Induction on the finite derivation reaches that signing step.
The initialized ledgers and boundary packets are built from ordinary data and pairing, and there is no other operational signature constructor.
Every predecessor is in the causal past, so the resulting signing precedes validation.

\item \textbf{Issuer and delegation origins.}
Lemma~\ref{lem:admin-signer-binding} identifies that signing principal with the issuer of the accepted record.
Applying the same argument to each selected delegation edge proves its origin at the source principal of that edge.

\item \textbf{Authorization witness.}
The authority path and its root membership are fields of the release witness, so inverting that witness proves the authorization statement at the recorded view.

\item \textbf{Single-occurrence inventory invariant.}
For a fixed $k$, the sum of the counts of its fresh, pending, available, and spent occurrences is initially one.
Each displayed rule replaces one stage by its successor, and every other rule preserves these counts.
The stage can advance only three times, with no reverse transition; hence each advance occurs at most once.

\item \textbf{Unique preceding stages.}
Following the consumed available occurrence gives the unique preceding acceptance, and following its pending occurrence gives the unique begin.
These assignments are injective because a key has only one occurrence of each advance.

\item \textbf{Naturality and occurrence distinction.}
Prefix extension preserves the producing events and pointers, proving naturality.
Independent exchanges change neither the relevant pointers nor the invariant.
The argument follows the actual chosen derivation, so equal public signatures from different signing events are not collapsed.
\end{enumerate}
\end{proof}

\begin{definition}[Protected reply genealogy]\label{def:admin-replies}
For a finite full history $h$, let $R_h$ be its protected release occurrences and let $P_h$ be its non-root genealogy nodes: immutable reply records and protected packet components. A reply record has its producing release as parent. Emission has its reply record as parent; forwarding, copying, and replay create a new packet component with its selected earlier protected component as parent. A packet carrying several protected records has one genealogy node for each selected component. Let $D_h$ be the protected delivery occurrences, each retaining its incoming component. These constructors give maps
\[
 \pi_h:P_h\longrightarrow R_h\amalg P_h,
 \qquad \delta_h:D_h\longrightarrow P_h.
\]
The public wire value may omit these names. A drop removes the held packet and creates no delivery; its earlier genealogy nodes remain in $h$. History inclusions preserve every node and its parent, defining functors $R,P,D:\mathcal H_{\mathrm{adm}}\to\Set$ and natural maps $\pi:P\to R\amalg P$ and $\delta:D\to P$. A fragment beginning with an already protected packet uses its producing history as the inherited audit of Proposition~\ref{prop:admin-inherited-audit}; its root remains in that full history.
\end{definition}

\begin{theorem}[Historical receipts as a derived decoration]\label{thm:admin-receipt}
There is a unique natural root map $r:P\to R$ satisfying
\[
 r=[1_R,r]\pi.
\]
Consequently $o=r\delta:\mathsf{Delivery}\to\mathsf{Release}$ is the unique delivery-origin decoration induced by fixing release roots and preserving parents. It assigns each delivery its originating release and hence that release's exact record, authority derivation, and decision view. Adding this derived receipt decoration to witnessed histories and erasing it are inverse operations. They preserve history extensions and residual cubes.
\end{theorem}
\begin{proof}\vmode
\leavevmode
\begin{enumerate}[label=(\roman*)]
\item \textbf{Well-founded parent recursion.}
The parent maps of Definition~\ref{def:admin-replies} strictly decrease finite creation order.
The unary ancestry instance of Lemma~\ref{lem:generic-operation-audit}, with release occurrences as roots, constructs the unique $r_h$ satisfying $r_h=[1,r_h]\pi_h$.

\item \textbf{Natural delivery-origin map.}
Its naturality gives $r:P\to R$ over history inclusions.
Composing with $\delta$ gives the claimed delivery origin and its record, view and authority-path projections.

\item \textbf{Uniqueness of the decoration.}
Any parent-preserving decoration fixing release roots satisfies the same equation, hence is this one.
Adding and forgetting it are therefore inverse and preserve extensions and residual cubes.
\end{enumerate}
\end{proof}

The map $o$ is generally not injective: several replays can deliver descendants of one release. In contrast the two maps in Theorem~\ref{thm:admin-origin} are injective. Equal-valued releases can have different roots, so receipt reconstruction is over witnessed histories, not over payload traces alone. A later registry revocation leaves the historical receipt intact while changing the guard of a new decision.

\begin{proposition}[Inherited administrative audit]\label{prop:admin-inherited-audit}
Let $H=\Conf_f(U_{\mathrm{adm}})$ be the initialized full-history carrier, and form $q:\widetilde C_{\mathrm{adm}}\to C_{\mathrm{adm}}$ as in Proposition~\ref{prop:history-residual}. It has the same unique witness and residual-family lifting. On the residual after $h_0\in H$, define
\[
 \mathsf{Delivery}_{h_0}(x)=\mathsf{Delivery}(h_0\cup x),\qquad
 \mathsf{Release}_{h_0}(x)=\mathsf{Release}(h_0\cup x),\qquad
 o^{h_0}_x=o_{h_0\cup x}.
\]
These give a natural receipt decoration, including releases in the inherited past. Resource survival, signing origins and retained decision views have the same restriction construction. All these decorations commute with residual cubes and iterated residuals.
\end{proposition}
\begin{proof}\vmode
\leavevmode
\begin{enumerate}[label=(\roman*)]
\item \textbf{Inherited predecessor relation.}
The proof of Theorem~\ref{thm:admin-residuals} places every selected ledger predecessor in the participating local past or in a matched packet past.
Thus the principal-history lifting proof of Proposition~\ref{prop:history-residual} applies to its enriched witnesses.

\item \textbf{Pullback of audit families.}
Pull back the functors and natural maps of Theorems~\ref{thm:admin-origin} and~\ref{thm:admin-receipt} along $j_{h_0}:x\mapsto h_0\cup x$.
Explicitly, for $x\subseteq y$ their naturality includes the commuting square
\[
\begin{tikzcd}[column sep=large,row sep=large]
\mathsf{Delivery}(h_0\cup x) \arrow[r,"{o_{h_0\cup x}}"] \arrow[d,hook]
 & \mathsf{Release}(h_0\cup x) \arrow[d,hook] \\
\mathsf{Delivery}(h_0\cup y) \arrow[r,"{o_{h_0\cup y}}"']
 & \mathsf{Release}(h_0\cup y).
\end{tikzcd}
\]
Parents and release records are unchanged by extension or independent exchange.

\item \textbf{Iteration coherence.}
The identity-on-events comparison $(U/h_0)/x\cong U/(h_0\cup x)$ proves the iteration coherence.

\item \textbf{Live-resource composition.}
The live-resource maps remain the same partially defined injections, with composition already proved in Theorem~\ref{thm:admin-residuals}.
\end{enumerate}
\end{proof}

\subsubsection{Cryptographic implementation of signed records}
\begin{construction}[Cryptographic interpretation of administrative records]\vmode\label{con:administrative-record-cryptographic-interpretation}
\leavevmode
\begin{enumerate}[label=(\roman*)]
\item \textbf{Tagged record representation.}
Use the existing $\mathsf{Rec}\to\mathsf{Prin}$ and its checked attestation resource.
Its wire representation is the tagged sum of the same field products.

\item \textbf{Record retraction.}
The map $\mathsf{encode}$ inserts those tags; checked field elimination gives $\mathsf{decode}\,\mathsf{encode}=1$.

\item \textbf{Concrete serialization and inverse.}
For a concrete signature realization, form
\[
 D_{\rm adm}=B+\mathsf{Prin}+K+X_r+X_a+
              \{\mathsf{grant},\mathsf{delegate}\}.
\]
Its coproduct injections and right-associated pairing give a serialization $j$ of wire records into that instance's generated message carrier.
Successive checked projections and tag tests construct its partial inverse.

\item \textbf{Signature interpretation.}
At principal $p$ the registered private key $s_p$ interprets the authenticated record by
\[
 [p,t]\longmapsto\crSig(j\mathsf{encode}(t),s_p).
\]
The key is supplied only to that protected issuing port.

\item \textbf{Verification and origin preservation.}
Verification uses its registered public image and checks the encoded issuer and fields.
Injectivity of $j\mathsf{encode}$ makes this precisely the original authentication test; copying and transport preserve the selected issuing occurrence.

\item \textbf{Free pairing extension.}
Extend record encoding to administrative values by the free pairing functor.
Its fold $E$ fixes ordinary data, commutes with pairing, and sends a signed record to its typed wire record with the authentication component.

\item \textbf{Tag separation and authentication interface.}
The wire signature tag is distinct from the pair tag, and the record fields remain readable through their declared interface.
Thus generic pair projection does not become a parser for signed values.
Replacing that authentication component by the displayed concrete signature is the service interpretation just constructed.
\end{enumerate}
\end{construction}

\begin{lemma}[Administrative representation retraction]\label{lem:admin-encoding-retraction}
The induced decoder satisfies $\mathsf{decode}_V E=1$. The interpretation is injective and commutes with administrative pairing, checked projection, signing, verification, copy and transport.
\end{lemma}
\begin{proof}\vmode
\leavevmode
\begin{enumerate}[label=(\roman*)]
\item \textbf{Generator equation and free extension.}
The record inverse and authentication test give the equation on generators.
Uniqueness of the free pairing fold extends it to all values.

\item \textbf{Operation squares.}
Distinct tags give the same projection failures, and the verification square was proved on the issuing generator.

\item \textbf{Copy and transport.}
Copy and transport preserve the interpreted value and its supplier.
These are the required operation squares.
\end{enumerate}
\end{proof}

\begin{theorem}[Cryptographic interpretation of signed grants]\label{thm:crypto-admin-comparison}
The local encoding of Construction~\ref{con:administrative-record-cryptographic-interpretation}
gives an injective interpretation of administrative operation witnesses,
resource inventories, authority paths and visible causal histories,
commuting with the compiled administrative transitions. With the
decoded administrative observation, the administrative judgements
retain their meanings.
\end{theorem}
\begin{proof}\vmode
\leavevmode
\begin{enumerate}[label=(\roman*)]
\item \textbf{Value-operation interpretation.}
Lemma~\ref{lem:admin-encoding-retraction} supplies the injective value
map and each value-operation square.

\item \textbf{Local command compilation.}
Compile a selected operation by
its finite record-access and authentication calls, owned by that
operation's participant, and hide those internal calls.
Its selected
operands and receipt occurrences are retained.

\item \textbf{Capability and origin preservation.}
The signing command
requires the same owner's available capability, and verification checks
the same indexed record.
Induction along copy, projection and actual
transport retains that origin; a public reference does not create a
new authentication token.
This is the local origin calculation of
Lemma~\ref{lem:admin-signer-binding}, independent of a global recipe
language.
The concrete signature implementation has the same property
by Corollary~\ref{cor:crypto-origin}.

\item \textbf{Inventories and registry guards.}
Fresh, Pending, Available and Spent inventories are unchanged, as are
the authority-graph vertices, edges and selected paths.
Registry
guards consult the same computed source view.
Each compiled operation
therefore has the same control and resource update, with its private
value representation related by $E$.

\item \textbf{Forward and inverse history comparisons.}
Induction on a finite history
constructs the forward comparison.
Decoding constructs its inverse
on completed macros in the image.

\item \textbf{Independent exchanges and causal order.}
Independent operations retain their
participants, operand suppliers and consumed inventory occurrences;
their added internal calls occur within their respective expansions.
Their exchange squares and visible causal orders are consequently
preserved.
This is a finite causal comparison of the completed-macro
displays, not an identification of intermediate private call positions
with completed administrative actions.

\item \textbf{Administrative judgement transport.}
The decoded observation square now identifies the original predicates
and action spans with their interpreted copies.
Pullback and
quantification along these identified maps give the same administrative
judgements.

\item \textbf{Public fields and later transitions.}
The public record retains its issuer and scope fields;
authentication does not conceal them.
Publication, later disclosure
and rotation have their own transitions and are not inserted into an
earlier history by this representation map.
\end{enumerate}
\end{proof}

\begin{definition}[Decoded administrative observation square]\label{def:decoded-administrative-observation}
On the completed-macro histories of Theorem~\ref{thm:crypto-admin-comparison}, let $u:H_{\rm crypt}\to H_{\rm adm}$ be its constructed erasure. Include the decoded administrative fields in the observation of cryptographic tests, giving $o_{\rm adm}u=v o_{\rm crypt}$.
Here $o_{\rm crypt}:H_{\rm crypt}\to O_{\rm crypt}$, $o_{\rm adm}:H_{\rm adm}\to O_{\rm adm}$ and the \textbf{observation decoding map} is $v:O_{\rm crypt}\to O_{\rm adm}$. The commuting equation defines its decoding on attained observations.
\end{definition} Then
\[
 u^*K_{\rm adm}P\subseteq K_{\rm crypt}(u^*P).
\]
Equal cryptographic observations give equal administrative observations, proving the inclusion directly from Proposition~\ref{prop:admin-modalities}. With the decoded administrative observation, the completed-macro comparison gives equality; recording additional cryptographic tests can provide more information. The publication and disclosure updates above also have their own action spans on the enlarged history space, independently of this erasure comparison.

\subsubsection{Example: administrative executions and modal reasoning}\label{subsec:admin-examples}
Use a direct grant $t=(\mathsf{grant},p,q,k,r,a)$ with $p$ initially in $\mathsf{Root}_{q,v}(r,a)$. The chronological executions below use the relay of Construction~\ref{con:administrative-process-interpretation}; their causal histories omit ordering between independent components.

\begin{example}[Signed-grant process and replay]\label{ex:signed-grant-replay}
The successful direct-grant term has selected action macros
\[
\mathsf{begin}_k;\ \mathsf{sign}_p(t);\ \mathsf{send}(h_t);
\ \mathsf{receive}(h'_t);\ \mathsf{verify}_q(h'_t);
\ \mathsf{accept}_k;\ \mathsf{release}_{k,v};\ \mathsf{deliver}.
\]
Semicolons record the stated local continuations or supplier dependencies; they impose no order between independent participants. Signing and verification expand into the service programs, communication uses the selected medium, and acceptance and release are the command fibres of Construction~\ref{con:dependent-grant-commands}. The token progresses through fresh, pending, available and spent, and release freezes its reply and decision view. The next proposition gives the complete dependent typing.

Replaying the grant can repeat reception and verification, but supplies neither a pending nor an available token. Duplicate reporting leaves the spent inventory unchanged. Replaying the reply instead gives another delivery with the same release root. Independent workers on distinct tokens use the tensor and residual square; the subsequent examples treat revocation, delegation and loss.
\end{example}

\begin{proposition}[A complete dependent release derivation]
\label{prop:dependent-release-derivation}
At an actual decision view \(v\), a validated grant, its selected
witnessed authority path, and its single available owned token type the
existing release command. Its completed result contains the frozen reply
and historical release witness. Delivery derives a receipt with that
release origin. Repeated validation or receipt copying does not produce
a second available token.
\end{proposition}
\begin{proof}\vmode
\leavevmode
\begin{enumerate}[label=(\roman*)]
\item \textbf{Verification and acceptance.}
The verification carrier supplies the selected signer-indexed record, and Corollary~\ref{cor:validation-justified-knowledge} gives its historical origin assertion.
The acceptance command consumes the pending occurrence and produces its available successor.

\item \textbf{Dependent release witness.}
At the actual view $v$, the selected authority path and protected-effect edge combine with that occurrence in $G_{\mathsf{release}_{k,t}}(s)$.

\item \textbf{Release update and result.}
Committing this witness consumes the available token, produces spent, performs the selected effect and returns the frozen reply with its release parent.
These are exactly the premises and result fibre of the typed command.

\item \textbf{Policy and authority evidence.}
The policy interpretation inserts the selected resource-rule path; bind and framing retain its endpoint evidence and unchanged inventory.
The authority derivation remains a selected field of the release witness.

\item \textbf{Natural historical projections.}
Projection gives $\mathsf{Release}\to\mathsf{Accept}$, while the root fold of Theorem~\ref{thm:admin-receipt} gives $\mathsf{Delivery}\to\mathsf{Release}$.
Their naturality preserves the recorded view under later extension, including revocation.

\item \textbf{Single-use resource obstruction.}
Theorem~\ref{thm:admin-origin} gives the operational single-release invariant.
The resource calculus also has no structural duplication: $\mathbf r([r])$ is terminal and $(\mathbf r\otimes_v\mathbf r)([r])$ is empty.
Reverification and copied evidence affect neither the spent inventory nor this typing, so they cannot supply the missing available premise.
\end{enumerate}
\end{proof}

\begin{example}[Revocation before and after release]\label{ex:admin-revoke}
Let the grant have been accepted. If the service removes $p$ from its root registry before release, with no other verified authority path to $p$, the available token remains but no release witness exists at the new view. A rejection report can record that result. If the removal follows a successful release and precedes delivery of its frozen reply, the delivered receipt still identifies the earlier release view and its valid authority derivation. It proves historical authorization, not current eligibility for another transaction. Registry update and decision at one service are locally ordered, so no spurious commuting square equates these two cases.
\end{example}

\begin{example}[A witnessed delegated release]\label{ex:admin-delegation}
Let $p_0$ be a recognized root and let the service have verified occurrences of
$\mathsf{sig}_{p_0}(\mathsf{delegate},p_0,p_1,r,a)$ and
$\mathsf{sig}_{p_1}(\mathsf{delegate},p_1,p_2,r,a)$, together with a grant signed by $p_2$. The two retained delegation occurrences give the selected path $p_0\to p_1\to p_2$. Release consumes the same single available token as a direct grant and records that path. Theorem~\ref{thm:admin-origin} supplies causal signing origins for both edges and for the grant. A second path to $p_2$ gives a distinct justification witness, but not an extra available token. Removing the root or a consulted delegation edge from a later decision view changes which later paths are enabled; it leaves the earlier path recorded in a historical receipt unchanged.
\end{example}

\begin{example}[Two workers and a lossy relay]\label{ex:admin-parallel-loss}
Take separate service workers with independent controls, initialized keys $k_1,k_2$, and their respective locally held decision views. Their release steps replace disjoint available tokens, retain disjoint witnesses, and commute. The two orders form a square and their game unstraightening contains that square. Two attempts to release the same key instead compete for one occurrence and have no such square. A relay may receive a protected reply and choose drop rather than forward. The full history then contains the release and the drop, but no delivery of that packet. These are distinct full histories even when a public observer sees the same absence of a reply. The observation modalities quantify over those actual alternatives.
\end{example}

\begin{proposition}[Synchronous and asynchronous administrative instances]\label{prop:admin-two-media}
Replacing the two relay links by the synchronous interfaces of Section~\ref{subsec:communication-kernel-composition} gives a synchronous administrative instance. The origin, single-release, receipt, and decision-descent results hold in both this instance and the FIFO instance. Their communication histories and enabled-state structures retain their respective operational differences.
\end{proposition}
\begin{proof}\vmode
\leavevmode
\begin{enumerate}[label=(\roman*)]
\item \textbf{Synchronous history construction.}
First form the synchronous administrative history strategy using the matching participant conditions of Section~\ref{subsec:communication-kernel-composition}.
The local signing, ledger, registry, and resource rules are unchanged.

\item \textbf{Transfer ancestry.}
A synchronous joint transfer places the received occurrence after its supplying local record; a FIFO reception instead follows its matching enqueue.
Both supply the predecessor relation used in the induction of Theorem~\ref{thm:admin-origin}.

\item \textbf{Resource safety and receipts.}
Neither channel operation creates a permission, so the single-release invariant is unchanged.
Parent recursion in Theorem~\ref{thm:admin-receipt} uses the actual transfer occurrences in each plant, giving its historical receipt.
Observation descent uses that plant's declared observer.

\item \textbf{Readiness realization.}
For the synchronous message interface, apply Definition~\ref{def:shared-readiness} to its constructed administrative unfolding.
Lemma~\ref{lem:occurrence-rigidity} already supplies its participant chains; Theorems~\ref{thm:readiness-history} and~\ref{thm:readiness-strategy} construct the global strategy and its reached kernel.

\item \textbf{Shared readiness and physical erasure.}
Readiness belongs to a current local condition and is shared by competing transfers, so it does not reserve a partner.
Erasure retains the full physical causal order, rule witnesses, parent pointers, and inventories.

\item \textbf{Extension lifting.}
Every physical extension lifts from every readiness representative by \eqref{eq:readiness-physical}, and independent physical families lift after preparation.
Its fibres and extension lifting are described by Proposition~\ref{prop:readiness-opfibration}.

\item \textbf{Observer transport and descent.}
For an observer $o$ use $o\varepsilon_{\mathrm{rd}}$.
Readiness leaves every physical inventory and evidence fibre unchanged, giving identity transports on the kernel-pair fibres.
These satisfy the identity and cocycle equations.
Surjectivity of erasure proves the predicate identities of \eqref{eq:admin-knowledge-transfer} with $\varepsilon_{\mathrm{rd}}$ in place of $\varepsilon$: physical alternatives lift and readiness alternatives project.
Lemma~\ref{lem:admin-evidence-descent} applies with the extension lifting just proved, so physical evidence and its natural extension maps descend as well.

\item \textbf{Action-witness comparison.}
For an action witness, perform its missing readies and then that exact physical event.
This gives the source-wise surjection of \eqref{eq:admin-action-comparison}; its kernel-pair quotient identifies preparations of the same physical witness and retains distinct physical witnesses.

\item \textbf{Administrative interpretation.}
Thus origin, receipt, resource safety, and decision descent have the same interpretation at the synchronous game interface.

\item \textbf{FIFO and control-arena instances.}
The FIFO message interface remains the copycat construction of Theorem~\ref{thm:fifocopycat}.
The distributed-controller and winning interpretation of either physical plant uses its control arena as in Theorem~\ref{thm:protocol-evidence-game-history}; no change of that arena's polarities is made.
A pending FIFO packet and a readiness event remain distinct intermediate histories.

\item \textbf{Physical duration.}
The physical clock and GoI calculations use the physical witness graph of Proposition~\ref{prop:local-delayed-paths}; if duration is read on readiness paths, give readiness weight zero and physical edges weight one.
The free-path universal property extends these weights uniquely, and the resulting duration is physical length after erasure.
\end{enumerate}
\end{proof}


\subsection{Validation and disclosure at administrative views}\label{subsec:admin-validation}
For the administrative plant, instantiate Definition~\ref{def:admin-observation} with its finite cumulative histories, local causal views, verification records, and consulted registry fields. Historical knowledge uses the public Scott frame; current action predicates use the local index-set operations just constructed.

\begin{corollary}[Knowledge justified by validation]\label{cor:validation-justified-knowledge}
At the validating interface $J$, let $V_t$ be the public Scott open witnessing successful validation of a selected record $t$, and $S_t$ the historical open witnessing a matching signing origin. Then
\begin{equation}\label{eq:validation-implies-origin-knowledge}
 V_t\subseteq\Box_J S_t.
\end{equation}
\end{corollary}
\begin{proof}\vmode
\leavevmode
\begin{enumerate}[label=(\roman*)]
\item \textbf{Historical witness opens.}
The validation and its finite signing ancestry persist in cumulative histories, so their witness sets are Scott open.

\item \textbf{Public origin inclusion.}
The result is public and Theorem~\ref{thm:admin-origin} gives $V_t\subseteq S_t$.
The greatest-public-subopen property of Proposition~\ref{prop:interior} proves the inclusion.

\item \textbf{Information retained by the interface.}
It does not reveal an origin identifier omitted by the interface.

\item \textbf{Current verdict descent.}
For a current decidable verdict, the separately constructed fibre-descent theorem determines whether that verdict depends only on the attained decision view.
\end{enumerate}
\end{proof}

The disclosure calculation observes the decision before its result is announced. Consider histories with the same public request and signed grant but different private root registries. If only one admits an authority path, the verdict varies within a public observation fibre. Descent then fails even though both services implement the resource rule correctly.

This failure identifies the information carried by the verdict. To test timing or later communication, include those outcomes in the chosen behaviour object; the Boolean calculation alone does not determine them.

\subsection{Cryptographic publication, evidence, and knowledge}\label{subsec:admin-crypto}

Use cumulative histories, retaining executed preparation and actual transport sources. For current decisions the attained records have the following concrete projections; persistent public observations continue to be Scott opens.
\begin{definition}[Coalition interface records]\label{def:crypto-coalition-observation}
For a finite history $h$, let $o_J(h)$ retain the actual joint interface record of coalition $J$: its exposed causal events, issued handle labels, completed test replies, and declared scheduling or administrative fields. Put $O_J=\im(o_J)$. Future commands are tested by contextual interaction; their hypothetical outcomes are not inserted as already performed events.
\end{definition}
\begin{proposition}[Restriction of coalition records]\label{prop:crypto-coalition-modes}
For $J\subseteq K$, restriction gives $o_J=\pi_{KJ}o_K$, coherently for three coalitions. The local subset and evidence-family operators apply to these actual record maps. A current outcome descends to $O_J$ exactly when it is constant on their fibres.
\end{proposition}
\begin{proof}\vmode
\leavevmode
\begin{enumerate}[label=(\roman*)]
\item \textbf{Restriction and composition.}
Restriction retains precisely the indicated event and reply fields.
Restricting twice equals restricting once.

\item \textbf{Fibre descent criterion.}
The fibre criterion is Theorem~\ref{thm:admin-descent}.

\item \textbf{Joint-interface tests.}
Joint tests appear as actual commands on the joint interface; they are not reconstructed from separate component tests.
\end{enumerate}
\end{proof}

\begin{proposition}[Publication, disclosure, and revocation]\label{prop:crypto-disclosure}
Publishing a public key adds a handle for its public image and enables the corresponding public operations. It does not add a secret-key handle. Disclosing the secret key permits decryption of previously retained matching ciphertexts. Revocation removes an authorization edge from subsequent decisions, while leaving existing frame entries and receipts intact.
\end{proposition}
\begin{proof}\vmode
\leavevmode
\begin{enumerate}[label=(\roman*)]
\item \textbf{Public-key publication.}
Publication computes the public-image operation at the protected key interface and inserts a reference to that result in the recipient's table.
The key entry remains in its original private table.
The command transitions have no public inverse of the public-image operation.

\item \textbf{Later key disclosure.}
Later disclosure instead imports an actual reference to the secret key through its authorized transport port; the keyed-decoding square then returns the plaintext from a retained matching ciphertext.

\item \textbf{Revocation update.}
Revocation deletes the designated authorization edge from the current finite registry, as in Construction~\ref{con:signed-grants-owned-stages}, and does not delete existing references or receipts.
Thus the enabled administrative decisions change, while those already held values remain available.

\item \textbf{Concrete encryption instance.}
In the concrete encryption instance these commands are exactly public-key extraction, key transport, and decryption.
The claims follow from their state updates and decoding equation, independently of a recipe syntax.
\end{enumerate}
\end{proof}

Publication is a value-bearing event. It cannot in general be replaced by restricting the current state space to a proposition that was already true. Owner authentication additionally requires the registration and credential paths of Section~\ref{subsec:admin-policy}: anyone can publish its own key while writing somebody else's name in an ordinary message.

\begin{proposition}[Public knowledge of cryptographic origin]\label{prop:crypto-knowledge}
For the selected verification record, let $V$ be the public open witnessing its successful reply and $O$ the historical open witnessing the matching constructor in its supplying past. Then $V\subseteq\Box_J O$. A MAC origin is an event using the matching shared key, not necessarily an event of one uniquely identified holder.
\end{proposition}
\begin{proof}\vmode
\leavevmode
\begin{enumerate}[label=(\roman*)]
\item \textbf{Authenticated-record origin.}
For local authenticated records, Lemma~\ref{lem:admin-signer-binding} and the root evaluation of Lemma~\ref{lem:generic-operation-audit} give $V\subseteq O$.

\item \textbf{Signature and MAC origin.}
For the concrete signature or MAC realization, Corollary~\ref{cor:crypto-origin} gives the same inclusion.

\item \textbf{Public knowledge.}
Both have finite historical witnesses and $V$ is public, so Proposition~\ref{prop:interior} gives the conclusion.

\item \textbf{Scope of MAC attribution.}
The MAC operation is available to every actual key holder, proving the stated scope of the origin claim.
\end{enumerate}
\end{proof}


\begin{corollary}[Signed-grant evidence in game histories]\label{cor:signed-grant-evidence-games}
The signed-grant construction instantiates Theorem~\ref{thm:protocol-evidence-game-history}. Its origin, receipt, inventory, and single-release evidence retain their physical meanings in the compiled game histories. Its sequential-worker register implementation has the additional finite commitment bound of Proposition~\ref{prop:register-causal-view-comparison}.
\end{corollary}
\begin{proof}\vmode
\leavevmode
\begin{enumerate}[label=(\roman*)]
\item \textbf{Causal presentation and controller compilation.}
Theorem~\ref{thm:admin-residuals} constructs the causal presentation and its physical witness comparison.
Apply the abstract controller compiler and pull back its inventory, origin and receipt families along erasure.

\item \textbf{Evidence transport through slots.}
A slot does not consume a token, update a registry, create a release or change a parent pointer.
Theorem~\ref{thm:protocol-evidence-game-history} therefore transports exactly the evidence maps of Theorems~\ref{thm:admin-origin} and~\ref{thm:admin-receipt}, with identity transport at slots.

\item \textbf{Sequential-worker comparison.}
The sequential-worker comparison identifies its register guards and gives the stated bound.
\end{enumerate}
\end{proof}
\begin{proposition}[Historical release projections]\label{prop:signed-grant-release-projections}
A completed physical command has the result family of equation~\eqref{eq:dependent-command-result-family}, with its private audit and operational state coordinates. The grant example constructs projections
\begin{equation}\label{eq:release-delivery-evidence-projections}
 \mathsf{Release}(h)\to\mathsf{Accept}(h),\qquad
 \mathsf{Delivery}(h)\to\mathsf{Release}(h),
\end{equation}
and its authority-path record at the decision view. These are maps of historical evidence, not operations creating owned tokens. 
These maps retain the underlying historical witness.
\end{proposition}
\begin{proof}\vmode
\leavevmode
\begin{enumerate}[label=(\roman*)]
\item \textbf{Retained parent pointers.}
Theorem~\ref{thm:admin-origin} associates a release with its retained acceptance, and Theorem~\ref{thm:admin-receipt} associates a delivery with its retained release.

\item \textbf{Natural evidence maps.}
Reading these pointers constructs the displayed maps, naturally under history extension.
Their interpretation as evidence maps uses Section~\ref{sec:protocol-evidence-policy}.
\end{enumerate}
\end{proof}
\subsection{Interpretation of the constructed extensions}\label{subsec:extension-interpretation}
\begin{corollary}[Substitution and preservation]\label{cor:extension-substitution-preservation}
The service, circuit, certificate and administrative macros preserve scope, typing and runtime substitution, and commute with their game interpretations. They transport the constructed service comparisons, causal certificates and contextual Scott equations. Each semantic simulator has its occurrence readback, and every concrete simulator here has its displayed named-process realization. Their witnessed secure conversions give the mutual slice factorizations of Theorem~\ref{thm:correctness-security-mutual-factorization} and the attainable-view-image equalities of Corollary~\ref{cor:security-images}, at the stated source, target and coalition interfaces.
\end{corollary}
\begin{proof}\vmode
\leavevmode
\begin{enumerate}[label=(\roman*)]
\item \textbf{Operation and family preservation.}
Proposition~\ref{prop:generic-process-preservation} applies to the constructed operation maps; Theorem~\ref{thm:admin-soundness} gives the administrative family interpretation and Lemma~\ref{lem:admin-encoding-retraction} its operation comparison.

\item \textbf{Contextual transport.}
Connection and tensor preserve the resulting causal certificates.

\item \textbf{Categorical interpretation and readback.}
Apply Theorem~\ref{thm:categorical-crypto-bridge}, and use Theorem~\ref{thm:pc-representation} for occurrence readback.
The other examples retain their separately stated correctness, origin or view-privacy conclusions; the transport does not change those interfaces or turn them into stronger realizations.
\end{enumerate}
\end{proof}


The examples have now supplied process realizations and comparison witnesses at their specified interfaces. Channel and view simulations use the general realization judgement; the remaining results retain their stated correctness, origin or information conclusions. The next appendix realizes the shared allocation and audit operations on which these examples depend.

\section{Concrete realization of public operation algebras}\label{app:symbolic-implementation}
The public-algebra comparisons require concrete allocation, import, initialization and audit operations. This appendix constructs those operations and shows how their implementations retain the comparisons of Section~\ref{subsec:information-interfaces}. The cumulative name world keeps track of support even after a value leaves a current reference table.

\subsection{Allocation worlds and evaluation}\label{par:world-indexed-evaluation}
Use the finite sorted name-world category $\mathbb I$ of Definition~\ref{def:calculus-indexed-carriers}, with a small coded presentation and its tagged extensions $A\hookrightarrow A+1_\tau$. The free value construction gives $T:\mathbb I\to\Set^{\mathcal S}$. Corollary~\ref{cor:evaluation-name-transport} gives $T(j)^\bot\crEval_A(R)=\crEval_B(jR)$, preserving and reflecting each success and observable result. Independent extensions compare in their common tagged sum; all associativity and exchange maps fix the old names and each new occurrence tag, proving coherence. These are disjoint-sum constructions in sorted sets, not a claim of coproducts in the injection category. Worlds retain the cumulative allocated and imported support, including discarded values. The counter implementation is the representation of these same extensions proved below; it does not make a shared allocator concurrent with itself~\cite{crypto-stark}.

\begin{construction}[Open inputs over cumulative name worlds]\vmode\label{con:open-inputs-name-worlds}
\leavevmode
\begin{enumerate}[label=(\roman*)]
\item \textbf{World-indexed input family.}
At $A\in\mathbb I$, the raw input family is
\[
 \mathsf{Raw}_\tau(A)=\coprod_{j:A\hookrightarrow B}T_\tau(B).
\]
Its transition transports the entire stored state and audit along $T(j)$, records the selected boundary source $t$, and continues at $B$.

\item \textbf{Finite-support representation.}
Every finite incoming value with its identifications with old names gives such a summand by adjoining its finite support.
The category is small, so this is a set.

\item \textbf{Raw-input boundary.}
The ordinary public reference interface has no such raw-input port.

\item \textbf{Concrete receipt tokens.}
For its concrete realization, transport $j(a)$ to the retained token of $a$ and a new $b\notin j[A]$ to $(\mathsf{boundary},e,b)$, where $e$ is this receipt occurrence; retain its sort.
The inverse of $j$ on its image is unique, and the other tags are disjoint from local allocator tokens.
Thus the construction extends the current name embedding without choosing representatives.

\item \textbf{Allocation and evaluation.}
Local allocation uses the next worker token of~\eqref{eq:crypto-initial-counter}.
Evaluation naturality proves equality of all subsequent checked operations under this embedding.

\item \textbf{Successive and independent transports.}
Successive transports compose by $T(kj)=T(k)T(j)$.
Two independent receipts or allocations compare in the world containing their old names and both sets of new occurrence tags; both execution orders induce the same embedding.
A value passed from one new receipt to another retains its wire dependency instead.
This proves freshness, coherent independent exchange, and preservation of the exact command replies.

\item \textbf{Fixed command carrier and closed wires.}
A fixed command carrier is obtained by the coproduct over source worlds, returning failure on a mismatched source-world tag.
In a closed network each wire acts in the common ambient world and its source continues into the sender's audit.
\end{enumerate}
\end{construction}

\subsection{Audit evaluation and accessible computations}\label{subsec:audit-evaluation-accessible-computations}
Apply Lemma~\ref{lem:generic-operation-audit} to the actual value operations. Its input family $R_h$ consists of initialized inputs and declared boundary sources; its node family $N_h$ retains each selected operand and operation. Write $u_h$ for the unique term unfolding and $r_h$ for its unary root fold. Successful command equations make evaluation of $u_h$ equal the stored value. No new audit object or evaluation principle is needed.

For an origin calculation retain the selected subtree address: projection or decryption follows its recorded operand address, and the matching constructor or explicit boundary is a root. For certificates the node maps act instead on witness fibres. Thus equal values need not have equal origins. The common fold is natural under history extension and independent exchange in every instance.

\begin{proposition}[Adequacy of finite public computations]\label{prop:finite-public-computation-adequacy}
The values issued by finite reference-process executions are exactly the successful finite public computations on their available sources. Substituting a producing computation for its handle preserves evaluation, including failure.
\end{proposition}
\begin{proof}\vmode
\leavevmode
\begin{enumerate}[label=(\roman*)]
\item \textbf{Evaluation of the audit fold.}
The audit fold labels every source by its available value, each operation by the corresponding strict call, and every copy or receipt by its selected source.
Its evaluation gives the issued result.

\item \textbf{Execution of a finite computation.}
Conversely the finite computation's dependency diagram is executed by those calls in dependency order, retaining their returned handles.
Success gives every required operand.

\item \textbf{Substitution compatibility.}
The substitution equation is the existing value-evaluation square of Theorem~\ref{thm:substitution-indexed-interpretation}.
\end{enumerate}
\end{proof}

\begin{proposition}[Finite initialization by constructor processes]\label{prop:crypto-initialization}
Every finite sorted initial value frame has a finite private preparation using the displayed allocation, data-introduction, constructor and transport operations. At completion its exported frame is the specified frame up to a sort-preserving renaming of private atoms and the fixed binding of interface handles. Its cumulative audit contains an actual constructor origin for each signature or MAC subtree. The residual evaluator has the prescribed handle-level behaviour with the counters computed in \eqref{eq:crypto-initial-counter}.
\end{proposition}
\begin{proof}\vmode
\leavevmode
\begin{enumerate}[label=(\roman*)]
\item \textbf{Finite subterm graph.}
Take all subterms of the finitely many initialized trees.
Create one allocation node for each distinct atomic name, one introduction node for each ordinary-data leaf value, and one constructor node for each compound subterm, with operand edges from its immediate subterms.
Atomic occurrences shared by different trees use the same allocated entry.

\item \textbf{Allocation and operand supply.}
Allocate in the indicated owner resources using their initialization workers; private preparation transport supplies each constructor with its selected operands.

\item \textbf{Acyclic dependency order.}
The graph is finite and acyclic, since constructor dependencies decrease term depth.
Within each allocator, use its increasing counter order.
The finite extra allocations needed to reserve unused lower indices have no exported entries.
These allocator chains create no cycle with the subterm order.

\item \textbf{Execution and subterm correctness.}
Induction over this finite order constructs the commands and proves that each result is the corresponding subterm under the allocation renaming.
A finite topological execution exists by induction on the number of remaining nodes.
All such legal executions may be retained as preparation witnesses; the construction does not select topological orders for an arbitrary family of inputs.

\item \textbf{Exported roots and private handles.}
Copy and transport export exactly the designated roots.
Temporary handles belong to private preparation interfaces.

\item \textbf{Explicit handle adapters.}
The exposed handle slots are bound bijectively to the designated original slots; for counter-indexed handle alphabets this finite binding extends to future fresh slots by the increasing enumerations of the unused indices.
Both its forward and inverse lookup adapters are explicit, so prescribed public handle names and invalid-handle replies are respected.

\item \textbf{Completed frame and counters.}
At the completed configuration the exported values have exactly the stated binding.
The next counters are those reached by preparation, as in \eqref{eq:crypto-initial-counter}.

\item \textbf{Residual evaluator comparison.}
Relate subsequent states by this frame binding and by the atom renaming extended at each fresh allocation.
Corollary~\ref{cor:evaluation-name-transport} makes every constructor, destructor, equality and data-reading decision agree.
Copy and transport preserve the relation, while fresh allocation extends the atom bijection because neither side reuses a reserved atom.
Induction on requests gives the claimed residual evaluator behaviour.

\item \textbf{Cryptographic constructor origins.}
Every signature or MAC node in the subterm graph was produced by its actual constructor using its matching secret-key operand.
Its outgoing operand and export edges remain in the preparation audit.

\item \textbf{Inherited initialization audit.}
Retain the completed preparation $h_0$ as an inherited initial configuration: a residual history $x$ carries the cumulative history $h_0\cup x$, and restriction and iterated residuals use $(U/h_0)/x\cong U/(h_0\cup x)$.
This is the prefix-embedding construction of Proposition~\ref{prop:history-residual}, also used for administrative audits in Proposition~\ref{prop:admin-inherited-audit}.
Thus the preparation is an executed part of the new resource's cumulative audit, not a reconstructed assertion about the past of an unexplained imported value.
Its private events are hidden at the external interface without deleting them from that audit.
\end{enumerate}
\end{proof}

\begin{corollary}[Unforgeable origin in the constructed model]\label{cor:crypto-origin}
In a constructed closed cryptographic network, every successfully checked signature or MAC has a matching constructor event in its cumulative audit, including its executed preparation. Following operand and transport edges reaches such an event, and that event used the matching secret-key value. At an explicitly open raw-value boundary, the same backward traversal instead records that selected boundary source when no producing computation is present in the network.
\end{corollary}
\begin{proof}\vmode
\leavevmode
\begin{enumerate}[label=(\roman*)]
\item \textbf{Verification root and key.}
The verification pattern identifies the root constructor and key.

\item \textbf{Backward ancestry traversal.}
Follow the selected value backwards through copying, projection, transport and successful destructor results.
If a constructor has produced a value containing the relevant subtree, follow its appropriate operand until either that root was itself constructed or an initialized entry is reached.

\item \textbf{Termination and possible sources.}
Each traversal follows strictly earlier operand occurrences.
Other constructors cannot create the required root because their tree images are disjoint.
The finite ancestry therefore reaches its signature or MAC constructor, a prepared initial entry, or an explicitly open boundary source.

\item \textbf{Prepared and transported origins.}
Proposition~\ref{prop:crypto-initialization} supplies the earlier constructor for a prepared entry; actual transport continues the traversal into the sending audit.
In a closed network these exhaust all sources, since external data introduction has domain $D$ and does not introduce arbitrary signature or MAC trees.

\item \textbf{Open-boundary source.}
An open raw-value source is retained as such and is not relabelled as a signing event.

\item \textbf{Natural root fold.}
With the selected subtree address retained, this traversal is the root fold of Lemma~\ref{lem:generic-operation-audit}.
Its naturality transports the selected origin through audit extension and independent exchange.
\end{enumerate}
\end{proof}

Origin is not automatically the identity of a named person. A disclosed signing key can be used by its new holder. A shared MAC key can be used by every holder. These facts follow from the same operational construction.

\begin{proposition}[Freshness and compositional allocation]\label{prop:crypto-freshness}
The concrete names realize the cumulative world extensions: allocation never repeats an earlier atom, independent workers allocate disjoint atoms, and their updates commute. Open inputs preserve these properties. Renaming internal worker and instance identifiers preserves the public reference behaviour, coherently with exchange and reassociation.
\end{proposition}
\begin{proof}\vmode\label{def:formal-initialized-allocation-counters}
\leavevmode
\begin{enumerate}[label=(\roman*)]
\item \textbf{Occurrence names and initial counters.}
Represent a local allocated occurrence by $(i,p,j,\tau,n)$: its resource instance, owner, issuing worker, sort and ordinal.
The initialized counter is computed, rather than supplied:
\begin{equation}\label{eq:crypto-initial-counter}
 c_{i,p,j,\tau}=1+\max\bigl(\{-1\}\cup
 \{n:(i,p,j,\tau,n)\in\supp(S_{\mathrm{init}})\}\bigr).
\end{equation}
For an already running allocator use the maximum with its current counter.
Finite preparation reserves unused lower ordinals privately and reaches this counter.

\item \textbf{Imported-name representation.}
Imported names use the boundary-receipt tags of Construction~\ref{con:open-inputs-name-worlds}.

\item \textbf{Cumulative freshness invariant.}
The prepared counter is above every ordinal of its namespace in the initialized support.
Induct over the operation audit: data adds no atom, an introduction or checked elimination uses the operands' support, a wire retains its actual source, and allocation adjoins its next unused token.
Tokens are never recycled.
Boundary receipts use the disjoint tags of Construction~\ref{con:open-inputs-name-worlds}.
This constructs an embedding of the cumulative world at every stage, including names discarded from current banks.

\item \textbf{Independent allocation square.}
Distinct workers update distinct counters and banks, and their new tokens retain their own occurrence tags.
Both orders extend the embedding by those same tags, proving the independent square.
Shared-bank access retains its existing chain.

\item \textbf{Renaming and coherence.}
Internal renaming transports this embedding and fixes public handle slots via their lookup binding; evaluation naturality gives equal replies.
The associativity and exchange equations are the world-extension equations, checked on each old and fresh tag.
\end{enumerate}
\end{proof}

\subsection{Transport of the generated instances}\label{subsec:generated-instance-transport}
Each extension in Appendix~\ref{app:cryptographic-examples} constructs its private carrier and actual operation maps. Apply the same typed introductions, strict eliminations, reference service and game interpretation; new transported sorts use their coproduct envelopes. The executable instances use finite bitstrings for ordinary data and the finite evaluators supplied in the examples. Extending the key carrier retains its old atoms as one summand, so the surrounding free value construction and its equality are rebuilt functorially over that embedding.

Proposition~\ref{prop:generic-evaluator-comparison} lifts an operation comparison to references; Theorem~\ref{thm:service-behaviour-canonical-comparison} gives contextual replacement; a distributed network instead uses its finite causal certificate. Theorem~\ref{thm:categorical-crypto-bridge} then gives the simulation square and, with its correctness converter, mutual slice factorization. These are applications of the existing transport results, not a further extension signature or security theorem. A finite returned test is a principal Scott observation at its actual boundary configuration.

\begin{remark}[Joint experiments and future interaction]\label{rem:crypto-correlations}
Retain labelled generators and their correlations: handles holding $(k,k)$ and $(k,\ell)$ differ by one equality request. The generated algebra describes a frozen frame. Future interaction additionally uses the evolving service or causal comparison, with the actual shared services and ports.
\end{remark}


\section{Comparison with quantitative interpretations}\label{app:quantitative-comparison}
The exact symbolic comparisons proved in this manuscript have a specific relation to the usual quantitative meanings of ``one-way'' and ``pseudorandom''. This appendix identifies the additional semantic data involved in those meanings and gives a concrete separation between fresh-name and sampled-function behaviour. A pseudorandom function family is compared with a sampled random function by a resource-bounded observer. Its conventional distinguishing quantity has the form
\[
 \left|\Pr[A^{F_k}(1^\lambda)=1]-\Pr[A^{\rho}(1^\lambda)=1]\right|.
\]
The finite-name key oracle in Theorem~\ref{thm:symbolic-keyed-oracle} proves an exact algebraic indistinguishability statement. It does not instantiate this expression: there is no security parameter, distribution on $k$, probability on $\rho$, or time bound in that theorem. For example, the fresh memoized oracle always returns different names on different queries, whereas a uniform random function to a finite set of size $N$ returns equal outputs on two distinct queries with probability $1/N$. This supplies a concrete separation between the two semantics. GGM~\cite{crypto-ggm} constructs the conventional notion; symbolic encryption with derived pseudorandom keys is studied in~\cite{crypto-micciancio}. A probabilistic interpretation of the latter does not make freshness itself a distribution.

Symbolic non-inversion likewise means that the declared operations cannot derive a value. Computational one-wayness instead measures the success probability of efficient algorithms on sampled inputs. An injective free hash is a symbolic model, but does not establish collision resistance for a fixed-length compressing function.

The same distinction applies to the other constructions. Homomorphic evaluation supplies the fold equation, without a succinctness or computational-security conclusion. Proof-service erasure supplies exact symbolic simulation and audit extraction, without a polynomial-time extractor for arbitrary external proof bitstrings. Information-theoretic and computational conclusions require their quantitative semantics and realization theorem. A probability-based proof logic such as Squirrel~\cite{crypto-squirrel} addresses that level.


These distinctions locate the symbolic results precisely. The established comparisons concern the declared exact operations and their public observations; a quantitative realization must supply the corresponding distributions, resource bounds and comparison theorem. The technical presentations that follow return to the exact structures already used in the manuscript.

\section{Detailed process, kernel, and evidence presentations}\label{app:technical-presentations}
This appendix discharges the presentation obligations used in the main development. Named elaboration connects scoped syntax to occurrence realizers; the kernel calculations recover their secured interaction; and the evidence calculations identify the modes acting on retained histories. Each block can be read from the comparison result that invokes it, and the required constructions are given here.

\subsection{Named process syntax, binding, and offers}\label{subsec:pc-named}
This block constructs the named presentation used by the operational examples. It first fixes scope and substitution, then derives offers and their interpretation. The later elaboration proof connects these named steps to the retained occurrence histories.

\begin{definition}[Finite named process syntax]\label{def:formal-finite-named-process-syntax}
The base language $\pcLang$ includes named actions as well as its occurrence declarations and typed connections. Names label communicating ports and may themselves be transmitted. Initially values are names and ordinary message data. Its finite named surface syntax is
\[
 \begin{aligned}
 P,Q&::=0\mid(P\mid Q)\mid(\mathsf{new}\ a)P
           \mid\sum_{i\in I}\alpha_i.P_i,\\
 \alpha&::=a!v\mid a?(x)\mid\tau .
 \end{aligned}
\]
Here $I$ is finite, input binds $x$, and $0$ is the empty sum. The declarations $\pcCore[A;J;K]$, forwarding, and typed connection remain available as interface terms. We retain the waiting and finite-cohort constructors explicitly so that a concurrent join or a collective exclusion has its own occurrence presentation. Recursion and replication extend this grammar in Section~\ref{subsec:pc-iteration}; service commands extend it in Section~\ref{subsec:cryptographic-process-rules}.
\end{definition}

\paragraph{Sorted contexts and capture-avoiding substitution}\label{par:sorted-contexts-substitution}
\begin{definition}[Sorted value contexts and process constructors]\label{def:formal-sorted-value-contexts-and-process-constructors}
Let $\mathcal S$ be the set of value sorts, including the channel and data sorts of the declared signature. A \textbf{value context} $\Gamma=(\Gamma_\tau)_{\tau\in\mathcal S}$ is a finite sorted set. The set $\mathsf{Val}_\tau(\Gamma)$ consists of finite constructor trees of result sort $\tau$ with free leaves in $\Gamma$, using the stated value signature. Channel types record the payload sort of their corresponding send and input; the untyped-message instance uses one message sort. Binding is represented by an additional distinguished leaf $*$ in $\Gamma+\tau$. Equivalently one can use sorted de Bruijn indices; thus bound-variable renaming is not additional data.

The \textbf{scoped process sets} $\mathsf{Proc}(\Gamma)$ are the least sets generated by the displayed grammar, together with the finite recursive syntax $X$, $\mu X.P$ and $!P$. Process variables range over a separate finite context $\Xi$: $X\in\Xi$ gives a term, $\mu X$ maps $\mathsf{Proc}(\Gamma;\Xi+\{X\})\to\mathsf{Proc}(\Gamma;\Xi)$, and $!:\mathsf{Proc}(\Gamma;\Xi)\to\mathsf{Proc}(\Gamma;\Xi)$. We suppress $\Xi$ when empty or fixed. The remaining constructor types are:
\[
\begin{array}{ll}
0\in\mathsf{Proc}(\Gamma),&
\mid:\mathsf{Proc}(\Gamma)^2\to\mathsf{Proc}(\Gamma),\\
\nu_\tau:\mathsf{Proc}(\Gamma+\tau)\to\mathsf{Proc}(\Gamma),&
\mathsf{out}_{a,v}:\mathsf{Proc}(\Gamma)\to\mathsf{Proc}(\Gamma),\\
\mathsf{in}_a:\mathsf{Proc}(\Gamma+\tau)\to\mathsf{Proc}(\Gamma),&
\mathsf{tau}:\mathsf{Proc}(\Gamma)\to\mathsf{Proc}(\Gamma).
\end{array}
\]
Here $a$ has channel sort with payload $\tau$ and $v\in\mathsf{Val}_\tau(\Gamma)$. A finite sum is a finite family of such prefixed bodies, retaining its index set; reindexing by a bijection transports that family. Recursion variables have a separate binding context and do not stand for values. The operators $\mu X$ and $!$ have the preceding finite syntax types; their transition rules are specified by finite derivations below and their temporal interpretation is constructed in Section~\ref{subsec:pc-iteration}. These sets are the union of the finite-depth stages of this binding signature; recursion on depth supplies their elimination principle.
\end{definition}

\begin{definition}[Process formation and introduction]\label{def:process-formation-introduction}
\textbf{Scoped process formation and introduction.}
For each well-formed pair $(\Gamma,\Xi)$ there is a syntax sort $\Proc(\Gamma;\Xi)$. Its finite terms have rules
\[
 \frac{}{\Gamma;\Xi\vdash0\ \mathsf{proc}}\qquad
 \frac{\Gamma;\Xi\vdash P\ \mathsf{proc}\quad\Gamma;\Xi\vdash Q\ \mathsf{proc}}
 {\Gamma;\Xi\vdash P\mid Q\ \mathsf{proc}},
\]
\[
 \frac{\Gamma\vdash a:\mathsf{Chan}(\tau)\quad\Gamma\vdash v:\tau
       \quad\Gamma;\Xi\vdash P\ \mathsf{proc}}
 {\Gamma;\Xi\vdash a!v.P\ \mathsf{proc}},
\qquad
 \frac{\Gamma\vdash a:\mathsf{Chan}(\tau)\quad\Gamma,x:\tau;\Xi\vdash P\ \mathsf{proc}}
 {\Gamma;\Xi\vdash a?(x).P\ \mathsf{proc}},
\]
\[
 \frac{\Gamma;\Xi\vdash P\ \mathsf{proc}}{\Gamma;\Xi\vdash\tau.P\ \mathsf{proc}}
 \qquad
 \frac{\Gamma,a:\nu;\Xi\vdash P\ \mathsf{proc}}
 {\Gamma;\Xi\vdash(\mathsf{new}\ a:\nu)P\ \mathsf{proc}}.
\]
Here $\nu$ ranges over the name sorts whose allocation operations have been constructed. For a finite family of typed prefixed bodies there is the finite guarded-sum constructor. Its branch index remains in an offer witness. Empty sum is $0$. Process-variable, recursion and replication rules are
\[
 \frac{X\in\Xi}{\Gamma;\Xi\vdash X\ \mathsf{proc}}\quad
 \frac{\Gamma;\Xi,X\vdash P\ \mathsf{proc}}
 {\Gamma;\Xi\vdash\mu X.P\ \mathsf{proc}}\quad
 \frac{\Gamma;\Xi\vdash P\ \mathsf{proc}}
 {\Gamma;\Xi\vdash !P\ \mathsf{proc}}.
\]
Bound occurrences use the binder pointers and captured environments. The displayed $!$ constructs a process term, not an exponential connective on arbitrary owned interfaces.
\end{definition}

\begin{definition}[Process elimination and computation]\label{def:process-elimination-computation}
\textbf{Structural elimination of process syntax.}
A recursive interpretation of raw syntax supplies a handler with the displayed source and target for each constructor in Definition~\ref{def:process-formation-introduction}. Its fold is characterized by
\[
 \mathsf{fold}(c(P_1,\ldots,P_n))
   =c^{\mathrm{alg}}(\mathsf{fold}(P_1),\ldots,\mathsf{fold}(P_n)),
\]
using the extended context for a binding argument. These are the elimination/computation rules of the syntax as an inductive object. Structural induction follows by taking the family of propositions established by the handlers. Operational elimination is different: it observes offers and commits selected witnesses as in Section~\ref{subsec:calculus-protocols}. Structural recursion on a $\mu$ syntax node does not recursively execute its unfolded body.
\end{definition}

\begin{definition}[Capture-avoiding substitution]\label{def:capture-avoiding-substitution}
A sorted value substitution $\theta:\Gamma\to\mathsf{Val}(\Delta)$ lifts to
$\theta^{+\tau}:\Gamma+\tau\to\mathsf{Val}(\Delta+\tau)$ by weakening old terms and fixing the distinguished new leaf. Define substitution by recursion on the binding syntax, with clauses
\[
(\nu_\tau P)[\theta]=\nu_\tau(P[\theta^{+\tau}]),\qquad
(\mathsf{in}_aP)[\theta]=\mathsf{in}_{a[\theta]}(P[\theta^{+\tau}]).
\]
Output substitutes in subject, payload and continuation; sum and parallel act componentwise. A recursive binder fixes its process variable and substitutes in its captured value environment. The equations proved below are
\begin{equation}\label{eq:scoped-substitution}
P[\eta]=P,\qquad P[\theta][\psi]=P[\psi^*\theta].
\end{equation}
They make scoped syntax a right module for the sorted value-term monad. Allocation retains its binding pointer and, at a finite name world, uses the tagged extension $A\hookrightarrow A+1_\nu$.
\end{definition}

\begin{definition}[Valuation-indexed evaluation]\label{def:formal-valuation-indexed-evaluation}
For a closed value algebra $T$, put $\mathsf{Env}(\Gamma)=\prod_\tau T_\tau^{\Gamma_\tau}$. Evaluation makes $\theta$ induce
$\theta^*:\mathsf{Env}(\Delta)\to\mathsf{Env}(\Gamma)$, with
$(\theta^*v)(x)=\mathsf{eval}_v(\theta(x))$. Evaluation and substitution satisfy
\[
\begin{tikzcd}[column sep=large]
\mathsf{Val}_\tau(\Gamma)\times\mathsf{Env}(\Delta)
 \arrow[r,"{(-)[\theta]\times1}"] \arrow[d,"{1\times\theta^*}"']
&\mathsf{Val}_\tau(\Delta)\times\mathsf{Env}(\Delta) \arrow[d,"{\mathsf{eval}}"]\\
\mathsf{Val}_\tau(\Gamma)\times\mathsf{Env}(\Gamma)
 \arrow[r,"{\mathsf{eval}}"'] &T_\tau .
\end{tikzcd}
\]
The proof is recursion on the value tree. Process templates obey the corresponding equation after instantiation. In particular a symbolic equality or destructor branch is evaluated after instantiation, on both sides of that equation. This does not assert that a noninjective substitution preserves a transition previously selected by an inequality test. Bijective renaming gives an operational isomorphism; general substitution is reindexing of the valuation-indexed semantics.
\end{definition}

\begin{lemma}[Syntactic substitution before operational elaboration]\label{lem:syntactic-substitution-foundation}
For sorted substitutions $\theta:\Gamma\to\Val(\Delta)$ and
$\psi:\Delta\to\Val(\Lambda)$, substitution preserves each constructor's
formation judgement and
\[
 t[\eta]=t,\qquad t[\theta][\psi]=t[\psi^*\theta],\qquad
 \mathsf{eval}_v(t[\theta])=\mathsf{eval}_{\theta^*v}(t).
\]
The first two equations also hold for process syntax, with the lifted
substitution at each value binder and the unchanged process-variable
binder at $\mu$.
\end{lemma}
\begin{proof}\vmode
\leavevmode
\begin{enumerate}[label=(\roman*)]
\item \textbf{Variable and constructor cases.}
Induct on the finite constructor tree.
A variable is the corresponding
substitution component; a constructor uses the induction equations on
its ordered operands.

\item \textbf{Binder lifting equation.}
At a binder, both composite lifts fix the new leaf,
and on an old leaf both are the weakening of the same substituted term.
This proves $(\psi^*\theta)^{+\tau}=(\psi^{+\tau})^*\theta^{+\tau}$.

\item \textbf{Process constructors.}
The process cases apply this identity to input and restriction, act
componentwise on sum and parallel, and leave the process-variable
coordinate fixed at recursion.

\item \textbf{Evaluation compatibility.}
Evaluation follows by the same value-tree
induction, including the definedness and result of an evaluated test.
No transition or semantic comparison is used in this proof.
\end{enumerate}
\end{proof}

\begin{definition}[Indexed interpretation carriers]\label{def:calculus-indexed-carriers}
For a category $C$ and the valuation set $\mathsf{Env}(\Gamma)$ of Definition~\ref{def:formal-valuation-indexed-evaluation}, put
\[
 \mathcal F_C(\Gamma)=\prod_{\rho\in\mathsf{Env}(\Gamma)}C.
\]
An object is a family of objects of $C$ and an arrow is a pointwise family of arrows. A value substitution acts by precomposition with its evaluation map $\theta^*$. The \textbf{finite name-world category} $\mathbb I$ has finite sorted name sets as objects and sort-preserving injections as arrows. At world $A$, $T_\tau(A)$ denotes constructor values over those names and $\mathsf{Env}(A,\Gamma)=\prod_{x:\tau\in\Gamma}T_\tau(A)$. Its \textbf{fresh extension} is the tagged inclusion $A\hookrightarrow A+1_\nu$. Constructor evaluation and injection of names define the world actions; the allocated service carriers later extend this same construction.
\end{definition}

\begin{theorem}[Substitution and the indexed interpretation]\label{thm:substitution-indexed-interpretation}
Value substitution preserves the process judgments and satisfies $P[\id]=P$ and $P[\theta][\psi]=P[\psi^*\theta]$. It reindexes the interpretation at valuations. It does not assert preservation of a branch already selected by an inequality test under a later noninjective value substitution.
\end{theorem}
\begin{proof}\vmode
\leavevmode
\begin{enumerate}[label=(\roman*)]
\item \textbf{Typing under substitution.}
Induct on the scoped typing derivation.
Values substitute at their declared sorts; input and restriction use the lifted substitution, and recursion preserves its separate process-variable binder.

\item \textbf{Identity and composite substitutions.}
The identity equation and the composite equation follow from
$(\psi^*\theta)^{+\tau}=(\psi^{+\tau})^*\theta^{+\tau}$, checked on old leaves and the distinguished leaf.
This proves typing and~\eqref{eq:scoped-substitution}.

\item \textbf{Evaluation and test branches.}
The value-evaluation square of Definition~\ref{def:formal-valuation-indexed-evaluation} gives identical evaluated operands after either instantiation.
Thus a test selects the same branch on both sides of that square, including failure.

\item \textbf{Binder and operational comparison.}
Under a name binder use the same tagged extension; under a recursive binder use the same captured environment.
Induction on the finite accessible walks of Section~\ref{subsec:pc-named-proof} then identifies offers, selected medium records and residual continuations.

\item \textbf{Finite-cut transport.}
A common finite cut contains every record of a transition or cell, so its elaborations agree up to the specified address transport.

\item \textbf{Indexed functoriality.}
On $\mathcal F_C(\Gamma)$ this is precomposition by $\theta^*$; evaluation of composite substitutions proves its functoriality.
Use $C=\CG$ on the constructed ordinary-interface families and the full-audit carriers on their stated instances.
This proves the indexed interpretation without asserting that a previously selected inequality branch survives a subsequent noninjective substitution.
\end{enumerate}
\end{proof}

\begin{proposition}[The basic interpretation principle]\label{prop:basic-interpretation-principle}
The raw syntax of Construction~\ref{con:section-signatures} has structural recursion and induction. The explicitly typed semantic operations determine a unique recursive interpretation. When the displayed computation equations have been checked, this interpretation factors through the presented syntax.
\end{proposition}
\begin{proof}\vmode
\leavevmode
\begin{enumerate}[label=(\roman*)]
\item \textbf{Recursive interpretation and uniqueness.}
Interpret a finite syntax tree by applying the constructed semantic operation to the interpretations of its immediate subtrees, using the extended context at a binder.
Induction on depth proves uniqueness and typing.

\item \textbf{Binding substitution.}
Theorem~\ref{thm:substitution-indexed-interpretation} gives compatibility with binding substitution.

\item \textbf{Factorization through computation equations.}
Equality of interpreted maps is a typed congruence; checking the generating computation equations therefore gives its unique factorization through the presented syntax.

\item \textbf{Initial-semantics identification and scope.}
This is the finite-tree instance of initial binding semantics in~\cite[Theorem~4.1 and Section~4]{calculus-fpt}.
It asserts neither termination of recursive execution nor completeness of semantic equality.
\end{enumerate}
\end{proof}

\begin{definition}[Named offers and scoped restriction]\label{def:formal-named-offers-and-scoped-restriction}
For a scoped term in $\mathsf{Proc}(\Gamma)$, a \textbf{named offer} is a tuple $(w,i,a,\zeta,P\prime)$: $w$ is a prefix address, $i$ a member of that prefix's finite sum index, $a$ a channel value, $\zeta$ a payload in the channel's payload sort or its distinguished input binder, and $P\prime$ a continuation in $\mathsf{Proc}(\Gamma)$ or $\mathsf{Proc}(\Gamma+\tau)$ respectively. A committed input specializes its continuation by capture-avoiding substitution. Parallel composition exposes offers of distinct active components. Choosing one branch discards the other branches of that sequential component. Rebracketing and permuting parallel syntax transport addresses; they do not introduce scheduling events. The offer $\tau$ is a private local action with its own selected occurrence.

The binder $\mathsf{new}\ a$ creates a fresh scoped name. Scope extrusion is permitted across a parallel component in which the name is not free, after alpha-renaming to avoid capture. A transmitted private name remains within the scope containing its uses, including any pending medium entry. By contrast, $\pcCut_B$ connects already typed dual interfaces by secured interaction. The two binders implement different constructions.

For a fixed named occurrence interface, $\mathsf{accept}(n,a,v)$ spells the negative skeleton offer at $n$, and $\mathsf{emit}(p,b,v\mid R)$ spells the positive offer at $p$ waiting on the receipt set $R$. Their typing is the core formation rule on that interface; a finite-cohort clause retains its specified positive occurrences. An input prefix denotes its family of concrete value alternatives, with their interface branches and continuation dependencies. A sequential prefix is typed on its corresponding prefix interface: the interface records same-polarity sequencing and dependencies into negative moves. Bijective name renaming transports this interpretation componentwise. A substitution identifying channel subjects changes which operational offers can match.
\end{definition}

\begin{definition}[State judgments and prefix introduction]\label{def:state-judgments-prefix-introduction}
The \textbf{input-prefix type} is $?_{a,T}B:=\mathsf{Pref}^{-}_T(B)$; the \textbf{output-prefix type} is $!_{a,T}B:=\mathsf{Pref}^{+}_T(B)$. The prefix symbols $?$ and $!$ here are type constructors with these explicit arena meanings. They are different from process replication $!P$ and from a linear-logic exponential. Their formation rule is
\[
 \frac{\tau:\Ty_{\Val}\quad T=\llbracket\tau\rrbracket\quad
       (B_v:\Int)_{v\in T}\quad a:\mathsf{Chan}(\tau)}
 {?_{a,T}B:\Int\quad !_{a,T}B:\Int}.
\]
For open values this is interpreted valuation by valuation, as in Theorem~\ref{thm:substitution-indexed-interpretation}. The family is an index of semantic carriers, not an infinite list of syntax subtrees.

Write $\Gamma\vdash_{\rm st}P:B$ for a constructed strategy state $I\to B$ at each value valuation. An open interface term $\Gamma;\Delta\vdash P:B$ has such a state judgment on $(\bigotimes\Delta)^\ast\otimes B$ by iterated port bending. Prefixing applies to this whole boundary, so a sequential prefix also guards its continuation's incoming moves.

The introduction rules are
\[
 \frac{\Gamma\vdash a:\mathsf{Chan}(\tau)\quad
       \Gamma,x:\tau\vdash_{\rm st}P:B(x)}
 {\Gamma\vdash_{\rm st}a?(x).P:?_{a,T}B},
\]
\[
 \frac{\Gamma\vdash a:\mathsf{Chan}(\tau)\quad\Gamma\vdash v:\tau\quad
       \Gamma\vdash_{\rm st}P:B(v)}
 {\Gamma\vdash_{\rm st}a!v.P:!_{a,T}B}.
\]
For input, attach a negative root to every instantiated continuation display $S_v\to B_v$. For output, attach the positive root $r_v$ to the selected continuation only, and display it in the corresponding branch of the full output-prefix arena. These instructions construct the strategy carriers of the two rules.
\end{definition}

\begin{theorem}[Prefix typing, elimination, and computation]\label{thm:prefix-typing-computation}
Construction~\ref{con:dependent-prefix-arenas} gives an arena. The two prefix introductions give ordinary receptive, courteous strategies. Their operational elimination at an actually committed prefix $r_v$ has type
\[
 \mathsf{Pref}^{\epsilon}_T(B)/\{r_v\}\ \cong\ B_v,
\]
and the residual display is the selected continuation display. In particular input elimination computes capture-avoiding substitution $P[v/x]$; output elimination computes its supplied continuation $P$.
\end{theorem}
\begin{proof}\vmode
\leavevmode
\begin{enumerate}[label=(\roman*)]
\item \textbf{Prefix carrier and display.}
Lemma~\ref{lem:prefix-carrier} proves the arena and configuration assertions directly from the tagged order and consistency construction.
The same construction on the source displays maps a root to the corresponding arena root and a continuation event by $\sigma_v$.
The configuration description and local injectivity of $\sigma_v$ prove that this is a display.

\item \textbf{Receptivity at the root.}
For input, at the empty configuration every negative arena extension is exactly one $r_v$, and that root has one lift.
For output there is no enabled negative move at the empty configuration.

\item \textbf{Receptivity after commitment.}
After a chosen root, a negative extension must lie in its branch and lifts uniquely by receptivity of $\sigma_v$.
Incompatible roots from other branches create no missing receptive lift.

\item \textbf{Courtesy.}
A causal cover internal to a continuation satisfies courtesy by courtesy of $\sigma_v$.
A new cover from the prefix root into a continuation is mapped to an arena dependency already present in the prefix order, regardless of its two signs.
Thus all required courtesy implications hold.
This also explains why retaining an unchanged unprefixed incoming interface would be incorrect for a general sequential prefix.

\item \textbf{Residual arena.}
A different branch is inconsistent with $r_v$ and disappears in the residual.
Within branch $v$, deleting $r_v$ leaves exactly $|B_v|$, its inherited order and its original consistency.
The identity on continuation tags is the required arena isomorphism.

\item \textbf{Residual display and computation.}
The same identity on source continuation tags gives the commuting square
\[
\begin{tikzcd}[column sep=large]
 S_{\rm pref}/\{r_v\}\arrow[r,"\cong"]\arrow[d,"\sigma_{\rm pref}/r_v"']
 &S_v\arrow[d,"\sigma_v"]\\
 \mathsf{Pref}^{\epsilon}_T(B)/\{r_v\}\arrow[r,"\cong"']&B_v.
\end{tikzcd}
\]
For input $S_v$ is the interpretation of the body instantiated at $v$, which is $P[v/x]$ by Theorem~\ref{thm:substitution-indexed-interpretation}.
For output it is the selected body from the introduction premise.
This proves both computation clauses.

\item \textbf{Occurrence-indexed eliminator.}
The eliminator is indexed by a committed occurrence.
It is not a total current-state map selecting an arbitrary $v$ before communication.
Its protocol declaration supplies the typed operation witness and its exposed event.
\end{enumerate}
\end{proof}

\begin{definition}[Witnessed medium commitment]\label{def:pc-medium}
Use the carrier of Definition~\ref{def:calculus-medium-carrier}, with each local alphabet instantiated by the sorted offer labels just defined. For $\omega\in W_\Pi(b)$, put $I=I(\omega)$, $a_i=\mathsf{req}(b,\omega,i)$, $\ell=\mathsf{lab}(b,\omega)$ and $b'=\mathsf{next}(b,\omega)$. Its committed transition rule is
\begin{equation}\label{eq:pc-medium}
 \frac{(P_i\xrightarrow{a_i,u_i}P_i')_{i\in I}}
 {N_b(\vec P)\xrightarrow{\ell,\,\omega,(u_i)_{i\in I}}
       N_{b'}(\vec Q)},
 \qquad Q_i=\begin{cases}P_i',&i\in I,\\P_i,&i\notin I.\end{cases}
\end{equation}
The witness retains the actual premise offers $u_i$. Residual parallel continuations expose fresh child addresses. Nullary medium operations are allowed. The residual cells of the concrete media are the finite-limit coefficient construction of Definition~\ref{def:resource-coefficient}; they determine which commitments coexist and their transported witnesses. The generic format admits a supplied chart presentation, but every medium used below has its charts constructed from its operations.
The \textbf{medium-indexed process language} $\pcLang[\Pi]$ has the same scoped terms as $\pcLang$ and the committed transition derivations generated by~\eqref{eq:pc-medium}, with the medium store included in the network state.
\end{definition}

\begin{proposition}[Named presentation of the polarized core]\label{prop:pc-named-core}
On a named occurrence interface $A$, replace each negative event by its $\mathsf{accept}$ offer and each positive event $p$ by its $\mathsf{emit}$ offer with receipt set $J^{-1}(p)$, retaining $K$. Reading the committed occurrence set gives an isomorphism between this offered presentation and the transition graph of $\pcCore[A;J;K]$. It preserves all anchored concurrent cells and the interface display.
\end{proposition}
\begin{proof}\vmode
\leavevmode
\begin{enumerate}[label=(\roman*)]
\item \textbf{Enabled acceptance and emission.}
At a committed set $C$, an acceptance is enabled exactly by the negative rule of Definition~\ref{def:pc-core}.
An emission is enabled exactly when its arena extension exists, every receipt in $J^{-1}(p)$ is present, and no clause of $K$ is completed.
Its successor is $C\cup\{p\}$.

\item \textbf{Inverse edge assignments.}
Keeping the occurrence address makes these assignments inverse on edges, including equally labelled alternatives.

\item \textbf{Concurrent families and labels.}
A family is concurrent precisely when each subset can be committed at its anchor, which is the same test on both sides.
The labels and polarities are unchanged.
\end{enumerate}
\end{proof}

Communication between named components is obtained by assigning a medium $\Pi$ to their offers. We write $\pcLang[\Pi]$ for this operational instance of the same language. Section~\ref{subsec:pc-protocol-parameter} defines the typed witnessed commitment rule and its protocol parameter. Consequently the syntax and binding rules are fixed here; the protocol sections supply the additional operations and their causal interfaces.

\paragraph{Typed networks and the scope of the interpretation}\label{par:typed-network-interpretation}
\begin{definition}[Typed network category and interpretation]\label{def:formal-typed-network-category-and-interpretation}
The interface category and the value-context category have different roles. At a fixed value environment, let $\mathsf{Net}$ have finite interface lists as objects. Its arrows are finite typed expression trees built from constructed ordinary component declarations, forwarding, tensor, connection, and the cup and cap, quotiented by the symmetric monoidal and compact structural equations. Each named component is first elaborated by Proposition~\ref{prop:pc-named-elaboration}; the particular medium or service construction then supplies its declared ordinary interface and display. Components with only a cubical full-audit interpretation retain that interpretation until their stated ordinary refinement is applied.

Define $\mathcal I:\mathsf{Net}\to\mathsf{CG}$ on generators by those constructed displays, on tensor by parallel, and on connection by secured interaction. Proposition~\ref{prop:process-typing-structural-conversions} proves that it descends through the structural congruence. The complete composition square is
\[
\begin{tikzcd}[column sep=large]
\mathsf{Net}(B,C)\times\mathsf{Net}(A,B)
 \arrow[r,"{\mathsf{connect}}"] \arrow[d,"{\mathcal I\times\mathcal I}"']
&\mathsf{Net}(A,C) \arrow[d,"{\mathcal I}"]\\
\mathsf{CG}(B,C)\times\mathsf{CG}(A,B)
 \arrow[r,"{\circ}"'] &\mathsf{CG}(A,C).
\end{tikzcd}
\]
Tensor has the identical square with parallel composition in its bottom row. For a value substitution, instantiate a component at $v\in\mathsf{Env}(\Delta)$ or instantiate the unsubstituted component at $\theta^*v$. Equation~\eqref{eq:scoped-substitution} and the evaluated offer rules identify the resulting operational records. Their constructed interfaces and displays are therefore canonically isomorphic. This gives the indexed interpretation of open terms, with reindexing composition proved by the same substitution equation. It does not identify a value binder with a compact interface connection.

The operational comparison in Section~\ref{subsec:pc-named-proof} precedes this interpretation: it is an isomorphism of the witnessed graphs and their specified chart presheaves. The ordinary history and core-readback comparisons are then applied exactly on the instances for which they were constructed. In particular the full audit on its own execution arena, an ordinary public-interface strategy, and a quotient of observable behaviours remain separately typed constructions.
\end{definition}

\subsection{Scoped named elaboration}\label{subsec:pc-named-proof}
We construct the carriers used in Proposition~\ref{prop:pc-named-elaboration}. All sets below are formed from finite syntax, finite words, and finite maps. The value substitutions are those of the paragraph ``\nameref{par:sorted-contexts-substitution}''. A state includes the medium store under the same scope as the names occurring in it.

\paragraph{Finite derivations and virtual addresses}
\begin{definition}[Transition derivations and virtual addresses]\label{def:formal-transition-derivations-and-virtual-addresses}
For the set $\mathcal J$ of scoped transition judgements, the finite-premise rules define the indexed polynomial
\[
(\mathcal PX)_j=\coprod_{\rho:\operatorname{concl}(\rho)=j}
 \prod_{k\in\operatorname{Prem}(\rho)}X_{\operatorname{judg}(k)}.
\]
Its initial algebra $\mathsf{Der}=\colim_{n<\omega}\mathcal P^n0$ consists of finite proof trees: adjoining and removing the root are inverse, since finitely many subtrees have a common finite height bound. Recursion on height proves initiality. Unfolding $\mu X.P$ is a proof constructor with premise a transition of $P[\mu X.P/X]$, not a physical event. In particular $\mu X.X$ has no finite transition proof.

A virtual address is a finite rooted walk in the syntax graph, with back edges from recursive variables to their binders and edges $(\mathsf{copy},n)$ from a replication node to its body. A recursive binder captures the environment of its free value variables; its back edge restores that closure. A restriction at address $w$ binds the pointer $\mathsf b_w$, distinct from every free-name constructor. Input extends the continuation environment by its received value. These are the substitution clauses of Definition~\ref{def:capture-avoiding-substitution} along the walk. Thus an outer restriction is shared across copies, while a restriction inside a copy has that copy's address. Extrusion transports the pointer with its scope; no unused-name selection is involved.
\end{definition}

\paragraph{Finite ledgers and enabled offers}
\begin{definition}[Residual ledgers and enabled offers]\label{def:formal-residual-ledgers-and-enabled-offers}
A residual ledger $L$ is a finite partial map on virtual addresses recording consumed prefixes and their input values, selected sum branches, and persistent-server request ordinals and continuations. A position $(L,b)$ also retains the scopes of the medium store. Accessible walks respect these consumed prefixes, branch selections and captured environments; parallel exposes both components, while recursion and replication expose their addressed bodies without a physical step. Before a sum is selected, its branches offer alternatives of one sequential component.

An offer retains its prefix address, component, branch, subject, payload, environment and continuation:
\[
u=(w,p,i,a,\zeta,\rho,\kappa,P').
\]
For input, $\zeta=v$ ranges over the payload sort and installs $\rho[*:=v]$; output evaluates its payload in $\rho$; $\tau$ has no subject or payload. Persistent-server offers retain their request ordinals. Tests are evaluated at $\rho$ before their continuation is exposed, retaining the evaluation witness when it is an operational step.

Commitment records consumption at $w$, fixes its selected branch, and installs precisely this continuation and the medium update of~\eqref{eq:committed-fibre}. It consumes no parallel ancestor or replication-activation token. The component address prevents one sequential component from supplying two premises of a joint commitment. Starting at the empty ledger, take the union of finite reachable stages; this constructs exactly the position carrier used below, including recursion and dynamic scopes.
\end{definition}

\begin{definition}[Exposure judgements and ledger updates]\label{def:exposure-ledger-rules}
Fix a scoped syntax tree $P_0$ and its initial valuation. A closure code
is a finite environment record with pointers to lexical recursive binders;
unfolding a pointer reads its recorded body and captured value environment.
The pointer's self-reference is a field of this finite record, not an
infinite substituted term. An exposure item is
$(w,P,\rho,\kappa)$, where $w$ is a virtual walk, $\rho$ assigns values of
the declared sorts to the free value leaves of $P$, and $\kappa$ assigns
closure pointers to its free process variables.

For a finite ledger $L$, the judgement
$L\vdash w:P[\rho,\kappa]\ \mathsf{exposed}$ is the least relation
containing the root item and closed under the following rules.
The displayed tags are disjoint word constructors; $j=0,1$ in the
parallel row.
\[
\begin{array}{c|c}
\text{exposed body at }w&\text{new exposed item}\\ \hline
P_0\mid P_1&(w\,\mathsf{par}_j,P_j,\rho,\kappa)\\
\nu_\tau P&(w\,\mathsf{nu},P,\rho[*\mapsto\mathsf b_w],\kappa)\\
!P&(w\,\mathsf{copy}(n),P,\rho,\kappa)\quad(n\in\mathbb N)\\
\mu X.P&(w\,\mathsf{body},P,\rho,\kappa[X\mapsto c_w])\\
X,\ \kappa(X)=c&(w\,\mathsf{back}(c),P_c,\rho_c,
                         \kappa_c[X_c\mapsto c])
\end{array}
\]
Here $c_w$ records $(X,P,\rho,\kappa)$ at that binder occurrence.
At an exposed sum $P=\sum_{i\in I}\alpha_i.P_i$, an unselected ledger
entry gives the offer rules
\[
\begin{array}{c|c|c}
\alpha_i&\text{evaluated action}&\text{continuation environment}\\ \hline
a!v&\mathsf{out}(\mathsf{eval}_\rho a,\mathsf{eval}_\rho v)&\rho\\
a?(*)&\mathsf{in}(\mathsf{eval}_\rho a,z)\quad(z\in T_\tau)&\rho[*\mapsto z]\\
\tau&\mathsf{tau}&\rho
\end{array}
\]
The offer's prefix address is $w\,\mathsf{branch}(i)$ and its component
is the sum address $w$. Its record also includes $i,\rho,\kappa,P_i$
and the evaluated operands. These rules define
$\mathsf{Off}_p(L,b)$ as a set of records, not as a set of derivation
trees. A selected sum has no further offer: if its entry is
$(i,\rho',\kappa)$, its sole continuation exposure is
$(w\,\mathsf{branch}(i)\,\mathsf{cont},P_i,\rho',\kappa)$.
The empty sum exposes nothing.

A persistent input server retains a finite list of accepted requests.
If its list has length $n$, its next offer carries ordinal $n$;
commitment appends its actual received value and exposes its body at
$w\,\mathsf{request}(n)\,\mathsf{cont}$ with the extended environment.
All earlier request bodies remain exposed, and the server offers again
with ordinal $n+1$. This is the addressed form of the server equation.
Pure tests expose the continuation selected by their evaluated result.
An operational test instead offers its declared operation and exposes
that continuation only after commitment.

For a compatible family $(u_p)_{p\in I_\omega}$ define
$\mathsf{upd}(L,(u_p))$ by inserting the selected-sum entries and
appending the selected server entries. Compatibility means precisely
that the component addresses are distinct, as required by the
participant product in~\eqref{eq:committed-fibre}. Thus the inserted
entries have disjoint keys. The target function is
\[
 t\bigl((L,b),\omega,(u_p)\bigr)
   =\bigl(\mathsf{upd}(L,(u_p)),b_\omega\bigr).
\]
The runtime-state set is the least subset of finite well-sorted ledger
and store records containing the initial state and closed under this
function on the displayed witness fibres. The scope of a state retains
every binding pointer referenced by its ledger, closures or store.
\end{definition}

\begin{lemma}[Exposure typing, finite support and update locality]\label{lem:exposure-typing-locality}
Every exposed item and offer has its displayed scope and sorts.
The target of a committed witness is a state in those sorts. Each
offer has a finite exposure derivation and finite support; no offer is
derived from an infinite exposure walk. Updates at distinct component
addresses leave each other's ledger entries unchanged.
\end{lemma}
\begin{proof}\vmode
\leavevmode
\begin{enumerate}[label=(\roman*)]
\item \textbf{Environment typing.}
Induct on an exposure derivation.
The parallel and copy rules reuse a
typed environment.
Restriction extends it by the pointer of the
declared sort; input extends it by the selected value of that sort.

\item \textbf{Recursive closures.}
The recursive-body rule records exactly the current environment, and
the back-edge rule restores that recorded environment, so neither rule
reinterprets a bound value in a later input environment.

\item \textbf{Continuation and target typing.}
The sum rows
use sorted evaluation and Lemma~\ref{lem:syntactic-substitution-foundation}.
They therefore produce a typed continuation record.
A medium witness
already has its declared store-target type, proving the target claim.

\item \textbf{Finite support.}
Every inference has finite depth and finitely many premises; each
environment used in it is a finite record.
The union of its walks,
lexical references and stored-value supports is finite.
Closure pointers
are retained as frontier references, so this union does not unfold a
recursive cycle.

\item \textbf{Local ledger updates.}
Finally finite-map insertion at distinct keys commutes.
This last statement concerns ledger updates; simultaneous enabledness
and the medium-store updates remain the actual resource-chart tests.
\end{enumerate}
\end{proof}

\begin{definition}[Committed-witness fibres]\label{def:formal-committed-witness-fibres}
For a scoped state $q$, let $\mathsf{Off}_p(q)$ be its offers at active component $p$, as constructed in Section~\ref{subsec:pc-named-proof}. Write $W_\Pi(b)$ for the operation witnesses at store $b$, with participant set $I_\omega$, required action projections, and target $b_\omega$. The exact edge fibre is the dependent sum of matching premise fibres
\begin{equation}\label{eq:committed-fibre}
E_\Pi(q)=\coprod_{\omega\in W_\Pi(b)}
 \prod_{p\in I_\omega}
 \{u\in\mathsf{Off}_p(q):\operatorname{act}(u)=\operatorname{req}_p(\omega)\}.
\end{equation}
An input entry includes its selected value and hence its substituted continuation. Each equality is an evaluated subject/payload equality. The edge object is $E_\Pi=\coprod_qE_\Pi(q)$; its source is the coproduct projection, its target installs precisely the selected continuations and $b_\omega$, and its label is $\ell(\omega)$. This is a construction from products, pullbacks, and coproducts, not a new enabling axiom. For a nullary witness the product is terminal. The selected-premise pullback, for a fixed $\omega$, has the form
\[
\begin{tikzcd}[column sep=large]
\prod_{p\in I_\omega}\mathsf{Off}_p(q)_{\operatorname{req}_p(\omega)}
 \arrow[r] \arrow[d] &\prod_{p\in I_\omega}\mathsf{Off}_p(q) \arrow[d,"{\operatorname{act}}"]\\
1 \arrow[r,"{(\operatorname{req}_p(\omega))_p}"']
 &\prod_{p\in I_\omega}\mathsf{Act}_p .
\end{tikzcd}
\]
\end{definition}

\paragraph{Finite cuts and their colimit}

\begin{definition}[Finite cuts and witnessed colimits]\label{def:formal-finite-cuts-and-witnessed-colimits}
A finite cut $K$ is a finite prefix-closed set of virtual walks, including their intervening binder references and retaining closures at unexpanded frontiers. Finite unions give a filtered inclusion poset. Let $S_K,O_K,E_K$ be the positions, offers and committed edges supported in $K$, requiring both endpoints of an edge to lie there. Inclusion of cuts preserves every ledger, pointer and selected record. The concrete media inspect selected premises and their store, never absence of an unexpanded offer; hence these are injective maps of witnessed graphs and
\[
S=\colim_KS_K,\qquad O=\colim_KO_K,\qquad E=\colim_KE_K.
\]
Two presentations agree exactly when their retained records agree in a common cut. This erases unused unfoldings, while retaining actual choices and medium witnesses.

For each finite direction set $F$, construct $X_K(F)$ by the anchored-instruction test of Definition~\ref{def:resource-coefficient}, with all vertices and records in $K$. A chart and each equality in its finite test have a common finite support. Inclusion preserves and reflects those equalities, so
\begin{equation}\label{eq:cut-cubes}
X(F)=\colim_KX_K(F).
\end{equation}
Equivalently filtered colimits of sets commute with the finite products and equalizers defining this coefficient. Faces restrict the same records in every cut, giving the commuting squares
\[
\begin{tikzcd}[column sep=large]
E_K\arrow[r,hook]\arrow[d,"{(s,t)}"']&E_{K'}\arrow[d,"{(s,t)}"]\\
S_K^2\arrow[r,hook]&S_{K'}^2
\end{tikzcd}
\qquad
\begin{tikzcd}
X_K(G)\arrow[r]\arrow[d,"X_K(u)"']&X(G)\arrow[d,"X(u)"]\\
X_K(F)\arrow[r]&X(F).
\end{tikzcd}
\]
The cubical identities hold in each finite chart and hence in the colimit.
\end{definition}

\begin{lemma}[Support closure and evaluation reindexing]\label{lem:cut-support-reindexing}
The finite cuts form a filtered poset, and every state, offer, edge and
finite chart has a cut supporting all its data. Evaluated elaboration
at $P[\theta],v$ and at $P,\theta^*v$ gives equal records under the
constructor-address correspondence. This correspondence preserves
targets and every face of a finite chart.
\end{lemma}
\begin{proof}\vmode
\leavevmode
\begin{enumerate}[label=(\roman*)]
\item \textbf{Finite support closure.}
A support contains the prefixes of each retained walk and the finite
closure codes and binder references read along its derivation.
A
recursive frontier retains its pointer rather than recursively adding
its unfolded body.
Hence closing finitely many derivations under these
operations is finite.

\item \textbf{Common supports for cuts and charts.}
The closure of the union of two supports is a
common upper bound; equivalently take the union of their already
closed finite records.
A chart has finitely many vertices and selected
edges, so the same argument supports it.

\item \textbf{Evaluation reindexing.}
For reindexing, induct on the finite exposure derivation.
Evaluated
operands agree by Lemma~\ref{lem:syntactic-substitution-foundation}.

\item \textbf{Binder and closure transport.}
Input uses the lifted substitution fixing its received leaf;
restriction transports its distinguished pointer.
A recursive closure
captures the reindexed environment and its back edge restores it.

\item \textbf{Ledger and medium updates.}
The ledger update inserts exactly the same chosen branch and value
records.
The medium parameter is reindexed with those operands.
Thus both endpoints and each selected instruction agree.

\item \textbf{Chart face compatibility.}
A chart is
the finite collection of these updates and equal-subset tests; applying
the same comparison to every subset proves face compatibility.
\end{enumerate}
\end{proof}

\paragraph{Structural transports and the comparison maps}
\begin{definition}[Address-presentation groupoid]\label{def:formal-address-presentation-groupoid}
Associativity, exchange, removal of inactive components, and scope extrusion transport the retained prefix and binding addresses. Their inverses transport them back. Structural replication uses the explicit bijection
\[
\mathbb N\longrightarrow\{*\}\amalg\mathbb N,
\qquad 0\mapsto*,\quad n+1\mapsto n.
\]
It identifies the copies of $!P$ with the exposed copy and the residual copies of $P\mid!P$. Finite compositions of these maps and finite permutations of copy indices form a groupoid of address presentations. Transport is defined on the entire ledger, scope pointers, and stored values. Its action is a functor because transport by a composite function is the composite transport. A nontrivial permutation can act nontrivially on occurrences; it is not asserted to be the identity. We do not quotient distinct copy commitments into one event. Erasure of addresses is the subsequent coarse transition-relation map.
\end{definition}

\begin{construction}[Generating structural address transports]\vmode\label{con:structural-address-transports}
\leavevmode
\begin{enumerate}[label=(\roman*)]
\item \textbf{Parallel address transport.}
Write $0,1$ for parallel child tags.
Reassociation sends the path
prefixes $00,01,1$ to $0,10,11$, respectively; symmetry exchanges
$0$ and $1$.
The unit transport deletes the child prefix of the
nonzero component of $P\mid0$.
These formulas act on the prefix
preceding an arbitrary suffix.

\item \textbf{Scope extrusion.}
For scope extrusion
$(\nu a.P)\mid Q\cong\nu a.(P\mid Q)$ with the declared freshness
side condition, send the $P$-prefix $0\,\mathsf{nu}$ to
$\mathsf{nu}\,0$ and the $Q$-prefix $1$ to $\mathsf{nu}\,1$.
Transport the binder pointer itself by this binder correspondence;
the environment of $Q$ ignores its new, unused leaf.

\item \textbf{Replication address transport.}
Replication
sends $\mathsf{copy}(0)$ to the exposed-body prefix and
$\mathsf{copy}(n+1)$ to the residual-replication prefix
$\mathsf{copy}(n)$, with the outer parallel tag included.

\item \textbf{Ledger, value and closure transport.}
For any such supported address bijection $a$, define
\[
 (a_*L)(a(w))=a_*(L(w)),\qquad
 a_*(\mathsf b_w)=\mathsf b_{a(w)}.
\]
On value constructors apply $a_*$ to their pointer leaves; on closures
apply it to the lexical pointer and captured environment; on a store
apply it to the same name occurrences.

\item \textbf{Supported inverses.}
Removed inactive syntax has no
supported offer; a scope containing a pending store use is retained.
The inverse formulas give the inverse supported transport.
\end{enumerate}
\end{construction}
\begin{lemma}[Equivariance of retained operations]\label{lem:address-equivariance}
Structural address transports induce inverse maps of witnessed graphs
and their chart presheaves. They satisfy
$(ba)_*=b_*a_*$, $1_*=\id$, and commute with valuation reindexing.
\end{lemma}
\begin{proof}\vmode
\leavevmode
\begin{enumerate}[label=(\roman*)]
\item \textbf{Exposure-rule transport.}
Each exposure rule is carried to the same rule with its constructor
prefix renamed.
Extrusion preserves evaluated subjects because its
additional leaf is unused in the other component.
The replication
bijection retains the selected whole body and every suffix inside it.

\item \textbf{Updates and chart faces.}
The update equation follows from transport of finite-map insertion.
Every chart retains the same selected medium records and equality
tests, so its faces are transported.

\item \textbf{Composition and inverses.}
Composition and inverses hold on
address words, pointers and finite records by structural recursion.

\item \textbf{Value-substitution compatibility.}
Value substitution acts at free leaves, whereas address transport acts
at their retained occurrence pointers; the two recursions commute.
\end{enumerate}
\end{proof}

\begin{definition}[Retained prefix records]\label{def:retained-prefix-records}
Fix an address presentation and a finite cut $K$. A prefix record
consists of its source and target ledger positions, one medium witness,
and, for each required component, its selected prefix walk, branch,
evaluated operands, closure and continuation environment. The record
is admitted when the exposure and update rules of
Definition~\ref{def:exposure-ledger-rules} derive those data. Proofs of
exposure themselves are not fields of the record.
Let $G_{\rm pre,K}$ be the graph of these records.

Reading a finite operational proof traverses its prefix and context
rules and retains exactly these fields. Two proofs are
$\sim_{\rm ret}$-equivalent when those syntactic fields agree in this
fixed presentation. This forgets repeated inactive unfolding and
different proofs of the same exposure, but retains the chosen branch,
copy address and medium witness. It does not quotient the address
groupoid action.
\end{definition}
\begin{lemma}[Reading and reconstruction on a finite cut]\label{lem:retained-reading-bijection}
Reading and reconstruction induce inverse graph maps
\[
 e_K:G_{\rm pre,K}\xrightarrow{\cong}G_{\rm off,K},
 \qquad r_K:G_{\rm off,K}\xrightarrow{\cong}G_{\rm pre,K}.
\]
On each source--target fibre, reading proof trees is surjective and
has kernel $\sim_{\rm ret}$. These maps induce a natural isomorphism
of the finite-chart presheaves and commute with cut inclusions.
\end{lemma}
\begin{proof}\vmode
\leavevmode
\begin{enumerate}[label=(\roman*)]
\item \textbf{Reading retained fields.}
Reading sends the indexed family of retained prefix fields to the
same family in~\eqref{eq:committed-fibre}.
Each field is an offer by
its exposure derivation, and its evaluated action is the required
one.

\item \textbf{Reconstruction from exposure derivations.}
Conversely an element of that fibre supplies finite exposure
derivations.
Induct on each such derivation: parallel inserts its
context, restriction its scope, a recursive back edge its finite
unfolding rule, and a selected sum its prefix rule.
A replication
address uses finitely many unfoldings exposing that selected copy;
unused copies make no retained field.

\item \textbf{Medium rule and ledger target.}
Apply the medium rule to this
indexed family of premises.
The result has precisely the specified
ledger target.
Thus every offer record reconstructs a finite proof.

\item \textbf{Inverse equations modulo retained fields.}
Reading this proof returns its input fields; reconstructing a read
proof returns the same retained fields.
Equality of these fields is
exactly $\sim_{\rm ret}$ by its independent syntactic definition.
This proves the fibre bijections and their endpoint equations.

\item \textbf{Cell and face maps.}
At each vertex of a chart, the maps retain every selected instruction
and subset-state equality.
They therefore induce inverse cell maps;
restriction to a face commutes with retaining those same fields.

\item \textbf{Cut inclusion and address equivariance.}
Cut inclusion changes no field.
Lemma~\ref{lem:address-equivariance}
supplies the separate equivariance under structural transport.
\end{enumerate}
\end{proof}

\begin{proposition}[Finite-support recursive occurrences]\label{prop:pc-finite-support}
The constructions above give the occurrence presentation of recursive and replicated named processes. Reading prefix derivations as offer records induces isomorphisms of witnessed graphs and of their finite-chart presheaves. These are compatible with cut refinement and the structural address groupoid. Erasing addresses has exactly the structural-congruence transition relation as its image. No guardedness condition is used.
\end{proposition}
\begin{proof}\vmode
\leavevmode
\begin{enumerate}[label=(\roman*)]
\item \textbf{Reading on a finite cut.}
A finite prefix derivation uses finitely many accessible walks; their union is a cut.
Reading its selected branches, values and medium witness gives an edge of $E_K$ with the same residual position.

\item \textbf{Reconstructing a finite derivation.}
Conversely, traverse an edge's retained walks, applying the prefix, context, substitution and unfolding rules and then its medium rule.
Accessibility supplies their premises.

\item \textbf{Inverse record maps.}
The two procedures are inverse on retained records in a common cut; they do not identify arbitrary proof trees.

\item \textbf{Chart comparison.}
Apply these procedures at all vertices and edges of an $F$-chart in one common finite cut.
Both sides use the same selected updates and equal-subset tests, hence have the same residual records and face maps.
This proves the graph and chart isomorphisms.

\item \textbf{Cut and address compatibility.}
Cut inclusion fixes the records, and the address groupoid transports all their pointers by the same bijection, proving both compatibilities.

\item \textbf{Image transition relation.}
Every structural-congruence transition has a finite proof and hence such a record; conversely every record reconstructs a proof.
Erasing addresses therefore has exactly that transition relation as image, while different copy commitments can remain distinct above it.
\end{enumerate}
\end{proof}

\begin{lemma}[Offers and substitution]\label{lem:pc-substitution}
At a valuation $v\in\mathsf{Env}(\Delta)$, elaborating $P[\theta]$ gives the same evaluated records as elaborating $P$ at $\theta^*v$. The comparison preserves residual continuations, scopes, and the corresponding store parameters. Bijective atom renaming induces isomorphisms of witnessed graphs and charts.
\end{lemma}
\begin{proof}\vmode
\leavevmode
\begin{enumerate}[label=(\roman*)]
\item \textbf{Reindexing accessible walks.}
Apply Lemmas~\ref{lem:syntactic-substitution-foundation} and~\ref{lem:cut-support-reindexing} along each finite accessible walk.

\item \textbf{Binder values and offer records.}
Input assigns the same received value to the distinguished leaf; restriction transports the same binding pointer; recursion restores the same captured environment.
The evaluated subject, payload and test results therefore determine the same offer and medium record.

\item \textbf{Inverse renaming transports.}
A bijective atom renaming preserves and reflects the equality and guard tests, so these transports are inverse on graphs and their finite charts.
\end{enumerate}
\end{proof}

\begin{proof}[Proof of Proposition~\ref{prop:pc-named-elaboration}]\vmode
\leavevmode
\begin{enumerate}[label=(\roman*)]
\item \textbf{Inverse state and edge maps.}
Let $G_{\rm pre}$ be the graph of retained prefix records and $G_{\rm off}$ the graph constructed from \eqref{eq:committed-fibre}.
The preceding proposition defines inverse maps $e$ and $r$ on states and edges.
In particular
\[
\begin{tikzcd}[column sep=large]
E_{\rm pre} \arrow[r,"{e_E}","{\cong}"'] \arrow[d,"{(s,t)}"']
 & E_{\rm off} \arrow[d,"{(s,t)}"]\\
S_{\rm pre}\times S_{\rm pre} \arrow[r,"{e_S\times e_S}","{\cong}"']
 & S_{\rm off}\times S_{\rm off} .
\end{tikzcd}
\]
The input and scope cases are Lemma~\ref{lem:pc-substitution}; nullary operations use the terminal premise product.

\item \textbf{Cell and execution comparisons.}
The cell bijections give a natural isomorphism $X_{\rm pre}\cong X_{\rm off}$.
Thus their execution categories, obtained by quotienting paths by these same squares, are isomorphic, and so are their root-representable audit displays.

\item \textbf{Principal histories and core readback.}
On the ordinary history instances the principal-history construction transports each pair $(a,[a]_h)$ to the identical pair under $e$, preserving order, consistency and boundary display.
Core readback then gives the same ordinary strategy.

\item \textbf{Scott and audit compatibility.}
These comparisons form
\[
\begin{tikzcd}[column sep=large]
X_{\rm pre} \arrow[r,"{\cong}"] \arrow[d,"{\Sc}"'] &X_{\rm off} \arrow[d,"{\Sc}"]\\
\Sc(X_{\rm pre}) \arrow[r,"{\cong}"'] &\Sc(X_{\rm off}),
\end{tikzcd}
\]
with the audit polarization used when appropriate.

\item \textbf{Image transition relation.}
Forgetting the retained-record fibre replaces this graph by its image transition relation, exactly as asserted.
\end{enumerate}
\end{proof}

\subsection{Detailed kernel presentations}\label{subsec:technical-kernels}
This block supplies the generators, relations and reachability calculations used in Section~\ref{subsec:kernel-calculation}. The comparison proceeds from configuration categories to set-valued kernels and then identifies ordinary secured interaction inside their reached part.

\subsection{Generators and relations for Scott categories}
The generators below are the single-event extensions and their Scott reversals. A residual square supplies the commuting conversion for two independent process commitments. Relations generated by those cells identify the corresponding compositional routes while retaining the underlying occurrences.

\begin{definition}[Scott path category]\label{def:scott-path-category}
For an arena $X$, let $\mathsf{ScPath}(X)$ be the category with vertices as objects, positive edges oriented forwards, negative edges oriented backwards, and the commutative-square relation for every oriented two-cell. The edges and faces retain their actual witnesses.
\end{definition}
\begin{lemma}[Path presentation]\label{lem:pathpresentation}
There is a canonical isomorphism $\Sc(X)\cong\mathsf{ScPath}(X)$.
\end{lemma}
\begin{proof}\vmode
\leavevmode
\begin{enumerate}[label=(\roman*)]
\item \textbf{Vertex and edge data.}
A functor from a product of oriented intervals is determined by its vertex objects and one-coordinate arrows satisfying its two-coordinate commutative squares.

\item \textbf{Construction of the cube functor.}
To see sufficiency, order any finite set of changed coordinates and compose their arrows.
Adjacent transpositions connect all such orders, and each transposition preserves the composite by its square.
Thus the assignment defines the unique functor on that cube.

\item \textbf{Agreement on cofaces.}
Agreement of these functors on cofaces says exactly that shared vertices and edges have the same assignment.

\item \textbf{Universal-property comparison.}
Consequently functors from either displayed category to any $C$ are the same assignments of vertex objects and oriented edge arrows satisfying the square relations.
This is natural in $C$, so the universal properties identify the categories.
\end{enumerate}
\end{proof}
\subsection{Configuration families and event-generated carriers}
\begin{definition}[Configuration-family cube carrier]\label{def:family-cube-carrier}
For a labelled set $E$ and $\varnothing\in\mathcal F\subseteq\Pf(E)$, define
\[
 X(\mathcal F)_F=\{(x,\phi):x\in\mathcal F,\ \phi:F\hookrightarrow E\setminus x,
       \ x\cup\phi[V]\in\mathcal F\ (V\subseteq F)\}.
\]
For a coface $(u,\alpha):F\to G$, its face map sends $(x,\phi)\in X(\mathcal F)_G$ to $(x\cup\phi[\alpha^{-1}(1)],\phi u)$. This is the \textbf{configuration-family cube carrier}; Proposition~\ref{prop:family-cubes} proves the presheaf and product assertions.
\end{definition}

\begin{proposition}[Configuration families give explicit concurrent carriers]\label{prop:family-cubes}
Let $\mathcal F$ be a family of finite subsets of a labelled event set, containing $\varnothing$. Define $X(\mathcal F)_F$ to consist of pairs $(x,\phi)$ with $\phi:F\hookrightarrow E\setminus x$ and
\[
       x\cup\phi[V]\in\mathcal F\qquad(V\subseteq F).
\]
This is a symmetric precubical carrier, with vertices $\mathcal F$, one-event-extension edges, and exactly the displayed concurrent cells. It preserves inclusions of configuration families and independent parallel products. For an event structure $E$, it is $K(E)$ of Section~\ref{sec:game-calculus}.
\end{proposition}
\begin{proof}\vmode
\leavevmode
\begin{enumerate}[label=(\roman*)]
\item \textbf{Well-defined faces.}
A face retaining $J\subseteq F$ and assigning values $\epsilon:F\setminus J\to\{0,1\}$ sends the pair to
\[
 (x\cup\phi[\epsilon^{-1}(1)],\phi|_J).
\]
For $V\subseteq J$ its vertex is one of the original vertices, so the pair is a cell.

\item \textbf{Face and reindexing equations.}
Two successive faces take the same union of the assigned upper coordinates and retain the same restriction of $\phi$; this proves the face equations.
Bijections reindex $\phi$ and commute with these operations.

\item \textbf{Low dimensions and inclusions.}
The zero- and one-dimensional descriptions follow by substitution.
Inclusion of families preserves precisely these data.

\item \textbf{Tensor splitting.}
In a tagged parallel family every cell splits uniquely into the directions and vertices in the two factors; this gives the tensor isomorphism, with coherence because all iterated splittings retain the same tags.

\item \textbf{Event-structure instance.}
For $\mathcal F=\Conf_f(E)$ the defining condition is exactly the definition of $K(E)$.
\end{enumerate}
\end{proof}

This construction retains the finite histories and cells of the specified family. An inclusion $\mathcal F\subseteq\mathcal G$ induces the corresponding carrier inclusion. It supplies a direct comparison between a process configuration family and an enlarged family while retaining which histories and concurrent cells each admits.

\subsubsection{Event-generated cubes}
The carrier $K(E)$ is Definition~\ref{def:formal-event-generated-configuration-cubes}.

We calculate the Scott configuration order of Definition~\ref{def:scottorder} from the oriented cells.

\begin{theorem}[Ordinary Scott calculation]\label{thm:eventscott}
The relation \eqref{eq:scottorder} is a partial order and
\[
\Sc(K(E))\cong(\Conf_f(E),\sqsubseteq_E).
\]
The construction preserves parallel product and turns polarity reversal into opposition.
\end{theorem}
\begin{proof}\vmode
\leavevmode
\begin{enumerate}[label=(\roman*)]
\item \textbf{Commutation of opposite-polarity moves.}
An oriented path removes negative events and adds positive events.
If a positive addition is immediately followed by a negative removal, the events are distinct.
The removed event cannot causally precede the added one, since the final set is a configuration; the added event cannot precede the removed one, since the latter was present beforehand.
Their consistency follows from the intermediate configuration containing both.
Hence the two moves form a square and may be interchanged.

\item \textbf{Normal form and necessary order condition.}
Repeated interchange moves all negative removals before all positive additions.
In the resulting path the intermediate configuration is $x\cap y$, so existence of an arrow implies \eqref{eq:scottorder}.

\item \textbf{Construction of a path.}
Conversely remove the negative events in reverse causal order and add the positive events in causal order; this gives a path whenever \eqref{eq:scottorder} holds.

\item \textbf{Uniqueness modulo squares.}
Within either part, two permissible orders differ by interchanges of incomparable events, hence by square relations.
Thus there is exactly one arrow when the relation holds and none otherwise.

\item \textbf{Order laws.}
Transitivity follows from composition of paths.
If arrows exist in both directions, an event could neither be removed positively nor added negatively, so $x=y$.
This proves antisymmetry and the categorical identification.

\item \textbf{Products and reversal.}
Configurations of a disjoint parallel product split as pairs, and the formula splits componentwise; reversal exchanges removal and addition.
\end{enumerate}
\end{proof}

For stream trees this is the ordinary prefix category: there are no higher-dimensional cells or alternative directed paths. The free-category calculation gives exactly the same category. Every kernel below is a presheaf on the ordinary Scott category just constructed.

\subsection{Set-valued kernels, displays, and interaction}\label{subsec:kernels-displays-interaction}
\subsubsection{Categories of elements and cubical displays}
\begin{definition}[Pointed kernel]\label{def:pointed-kernel}
A raw kernel from $A$ to $B$ is a functor
\[
 \Phi:C_A\times C_B^\op\longrightarrow\Set
\]
with a chosen element $\rho\in\Phi(i_A,i_B)$. Equivalently it is a pointed presheaf on $C=C_A^\op\times C_B$. Its category of elements $E_\Phi$ has objects $(c,z)$ with $z\in\Phi(c)$ and an arrow $(c,z)\to(d,w)$ over $u:c\to d$ exactly when $\Phi(u)(w)=z$.
\end{definition}
\begin{proposition}[Constructed discrete fibration]\label{prop:kernel-elements}
The projection $E_\Phi\to C$ is a discrete fibration. Conversely, the fibres of a discrete fibration and its unique target lifts define a presheaf, naturally inverse to the category-of-elements construction up to canonical isomorphism over $C$.
\end{proposition}
\begin{proof}\vmode
\leavevmode
\begin{enumerate}[label=(\roman*)]
\item \textbf{Unique lifts.}
Given $u:c\to d$ and $(d,w)$, its unique lifted arrow is $(c,\Phi(u)w)\to(d,w)$.

\item \textbf{Functor laws.}
Presheaf functoriality proves preservation of identities and composition.

\item \textbf{Reconstruction of restriction maps.}
Conversely, composing the unique lifts defines the restriction maps and proves their functor laws.

\item \textbf{Inverse constructions.}
The maps sending a lifted object to its fibre element and an arrow to its base arrow identify the two reconstructions.
\end{enumerate}
\end{proof}
\begin{definition}[Cubical unstraightening]
For $M=A^\ast\boxtimes B$, use the unit of the Scott adjunction to define
\begin{equation}\label{eq:cubical-unstraightening}
 \Un(\Phi)=M\times_{N_{\mathrm{Sc}}C}N_{\mathrm{Sc}}E_\Phi.
\end{equation}
Thus a cell is an arena cell and a functor lifting its Scott cube to $E_\Phi$. The root is $(i_A,i_B,\rho)$.

The defining square is a pullback in the presheaf category:
\[
\begin{tikzcd}[column sep=large]
\Un(\Phi) \arrow[r] \arrow[d] &N_{\mathrm{Sc}}E_\Phi \arrow[d,"{N_{\mathrm{Sc}}p_\Phi}"]\\
M \arrow[r,"{\eta_M}"'] &N_{\mathrm{Sc}}C .
\end{tikzcd}
\]
Its two projections remember the arena cell and its actual lift, respectively.

\end{definition}
\begin{lemma}[Terminal corner]\label{lem:terminalcorner}
A lift of an arena cell $m$ is uniquely determined by an element of the kernel over its Scott-terminal vertex $m_t$. Consequently
\[
 \Un(\Phi)_c\cong\coprod_{m\in M_c}\Phi(m_t).
\]
The terminal vertex has value $1$ in each positive direction and $0$ in each negative direction.
\end{lemma}
\begin{proof}\vmode
\leavevmode
\begin{enumerate}[label=(\roman*)]
\item \textbf{Restriction from the terminal vertex.}
Every vertex $v$ of $J_{\mathrm{Sc}}(c)$ has a unique arrow to its terminal vertex.
Restrict the selected terminal element along these arrows.

\item \textbf{Functorial lift.}
The presheaf equations make the resulting assignments a functor on the cube.

\item \textbf{Uniqueness.}
Conversely the discrete-fibration property forces these assignments for any lift.

\item \textbf{Coface compatibility.}
Restricting a lift to a coface gives the same construction on that face, so the bijection is compatible with cell restriction.
\end{enumerate}
\end{proof}
An arena cell occurs in a strategy display precisely when its terminal fibre has an element. This calculates the strategy cells directly from its kernel.

\subsubsection{Reachability from progress and its endpoint image}
Let $M$ be an initialized arena and let $\Phi$ be a pointed presheaf on $C=\Sc(M)$. The following construction applies in particular to a morphism arena.
\begin{definition}[Progress localization and histories]\label{def:kernel-progress}
For every positive generating arena edge $u:c\to d$ and $w\in\Phi(d)$, the opposite category $E_\Phi^\op$ has the restriction arrow
$(d,w)\to(c,\Phi(u)w)$. Let $W_+$ be the set of all these arrows and define the ordinary localization
\[
 \mathsf{Prog}(\Phi)=E_\Phi^\op[W_+^{-1}].
\]
It is constructed by adjoining formal inverses and quotienting finite composable words by the category and inverse equations. Define
\[
 \mathsf{Hist}_\Phi(c)=
 \coprod_{z\in\Phi(c)}\Hom_{\mathsf{Prog}(\Phi)}(\rho,(c,z)).
\]
For $u:c\to d$, restriction takes $(w,h)$ to $(\Phi(u)w,r_{u,w}h)$, where $r_{u,w}$ is the opposite lifted restriction arrow. The \textbf{endpoint natural transformation} $\mathsf{end}:\mathsf{Hist}_\Phi\Rightarrow\Phi$ has component $\mathsf{end}_c(z,h)=z$. Define $\Reach(\Phi)$ to be its image in the presheaf category $\Set^{C^\op}$.
\end{definition}
\begin{lemma}[History and finite progress words]\label{lem:progress-words}
The history assignment is a presheaf and its endpoint map is natural. An element lies in $\Reach(\Phi)$ exactly when it can be obtained from the root by finitely many presheaf restrictions and reverse restrictions over positive generating arena edges.
\end{lemma}
\begin{proof}\vmode
\leavevmode
\begin{enumerate}[label=(\roman*)]
\item \textbf{Restriction laws.}
The unique lift of a composite base arrow is the composite of its lifts, so the displayed restriction formula respects composition and identities.

\item \textbf{Endpoint naturality.}
Forgetting the history commutes with restriction, proving naturality.

\item \textbf{Finite-word representation.}
An arrow in the ordinary localization is represented by a finite word of arrows of $E_\Phi^\op$ and inverses of $W_+$.
The former are precisely restrictions, and the latter are the indicated positive steps.

\item \textbf{History reconstruction and endpoint image.}
Conversely every such word represents a history.
Taking its endpoint proves the characterization.
Different history words may have the same endpoint; the image construction records existence while retaining the original endpoint element.
\end{enumerate}
\end{proof}
\begin{proposition}[Constructed reachability]\label{prop:kernel-reachability}
Reachability is a functorial pointed subpresheaf, is idempotent, and is right adjoint to the inclusion of the pointed kernels all of whose elements are reached. It is the smallest rooted subpresheaf closed under reverse positive restrictions. For every positive generator $u:c\to d$, its square with $\Phi(u):\Phi(d)\to\Phi(c)$ is a pullback.
\end{proposition}
\begin{proof}\vmode
\leavevmode
\begin{enumerate}[label=(\roman*)]
\item \textbf{Endpoint image.}
Presheaf images are computed objectwise, so the endpoint image is a subpresheaf.

\item \textbf{Functoriality.}
A pointed natural transformation induces a functor on element categories, carries each specified positive lift to a specified positive lift, and therefore induces a functor on the localizations.
It carries histories to histories, proving functoriality.

\item \textbf{Idempotence and minimality.}
Each prefix of a progress word is itself reached.
Hence every reached element is reached within the image subpresheaf, proving idempotence.
The same finite-word induction proves minimality.

\item \textbf{Positive-step pullback.}
The positive restriction and its inverse belong to the progress calculus, so $w$ is reached exactly when $\Phi(u)w$ is reached.
This is precisely the pullback assertion.

\item \textbf{Universal factorization.}
Finally, a pointed map from a reached kernel sends every progress word to one and therefore factors uniquely through the image.
This is the stated right adjunction.
\end{enumerate}
\end{proof}
The operational direction of a negative move is a restriction; the direction of a positive move is a reverse restriction. In an event display every finite legal execution therefore gives a progress history. The full history presheaf and its endpoint image remain distinct constructions.

For a positive generator $u:c\to d$, the pullback and image factorizations are
\[
\begin{tikzcd}[column sep=large]
\Reach(\Phi)(d) \arrow[r,hook] \arrow[d,"{\Reach(\Phi)(u)}"']
 &\Phi(d) \arrow[d,"{\Phi(u)}"]\\
\Reach(\Phi)(c) \arrow[r,hook] &\Phi(c)
\end{tikzcd}
\qquad
\begin{tikzcd}[column sep=large]
\mathsf{Hist}_\Phi \arrow[r,two heads] \arrow[dr,"{\mathsf{end}}"']
 &\Reach(\Phi) \arrow[d,hook]\\
 &\Phi .
\end{tikzcd}
\]
The first square is the reverse-positive closure equation; the second is the objectwise image defining reachability.

\subsubsection{Raw composition and compact duality}
\begin{definition}[Raw interaction]
For kernels $\Phi:A\to B$ and $\Psi:B\to C$, define
\begin{equation}\label{eq:coend}
 (\Psi\star_B\Phi)(a,c)=\int^{b\in C_B}\Phi(a,b)\times\Psi(b,c).
\end{equation}
Explicitly, this is the quotient of $\coprod_b\Phi(a,b)\times\Psi(b,c)$ by
\[
 (\Phi(a,u)x,y)\sim(x,\Psi(u,c)y),\qquad u:b\to b',
\]
where $x\in\Phi(a,b')$ and $y\in\Psi(b,c)$. Its root is the class of the two roots at $i_B$. The identity kernel is
$Y_A(a,a')=\Hom_{C_A}(a',a)$, rooted at the identity.
\end{definition}
\begin{theorem}[Raw kernel calculus]\label{thm:rawcompact}
Raw kernels and pointed natural transformations form a bicategory with associative composition up to coherent isomorphism, identity kernels $Y_A$, and symmetric monoidal tensor induced by product. Its object dual is $A^\ast$. Isomorphism classes of kernels consequently form a compact closed category.
\end{theorem}
\begin{proof}\vmode
\leavevmode
\begin{enumerate}[label=(\roman*)]
\item \textbf{Associator and pentagon.}
Both bracketings of three kernels are the quotient of triples by the two intermediate-arrow relations.
The associator retains the same triple, and every route around the pentagon does likewise.

\item \textbf{Units and triangle.}
Co-Yoneda gives the unit bijections
\[
 \int^b\Hom(b,a)\times F(b)\cong F(a),
 \qquad [u,z]\longmapsto F(u)z,
\]
with inverse $z\mapsto[1_a,z]$.
These formulas prove the triangle.

\item \textbf{Tensor composition and coherence.}
Tensor uses products of fibres; its composition comparison partitions the independent intermediate variables and therefore is a bijection.
Coordinate reassociation proves its coherence.

\item \textbf{Evaluation and coevaluation.}
Since $C_{A^\ast}\cong C_A^\op$, the hom kernel defines both evaluation and coevaluation, with the two identity roots.

\item \textbf{Snake equations.}
A snake composite is a coend of two hom sets.
Composition of arrows and the displayed co-Yoneda bijection identify it with the identity hom kernel.
This proves the duality equations and their compatibility with roots.
\end{enumerate}
\end{proof}
\begin{definition}[Secured interaction]
Put $\Psi\odot\Phi=\Reach(\Psi\star\Phi)$. A flat multiple secured interaction applies $\Reach$ once to the full multiple coend.
\end{definition}
The inclusions into raw coends give canonical maps from nested secured expressions to their flat counterparts. For the actual event strategies, Section~\ref{subsec:pc-secured-proof} identifies this operation with interaction and hiding and proves its associative comparison. Appendix~\ref{app:communication-examples} computes these maps for the concrete channels. These proofs specify the domain of each composition assertion rather than treating raw compact duality as a proof about every reachability restriction.
\subsubsection{The ordinary event-strategy comparison}
\begin{lemma}[Order reflection in a configuration]\label{lem:orderreflection}
If $s,t$ lie in a common configuration of $S$ and $f(s)\leq f(t)$, then $s\leq t$.
\end{lemma}
\begin{proof}\vmode
\leavevmode
\begin{enumerate}[label=(\roman*)]
\item \textbf{Predecessor lift.}
The image of $\downarrow t$ is a configuration containing $f(t)$, hence contains $f(s)$.
Some $s'\leq t$ has image $f(s)$.

\item \textbf{Local injectivity.}
Both $s,s'$ lie in the common configuration; local injectivity forces $s=s'$.
\end{enumerate}
\end{proof}

\begin{theorem}[Scott discrete fibrations]\label{thm:scott-discrete-fibrations}
For a polarized prestrategy $f:S\to E$, the map of finite Scott configuration orders is a discrete fibration exactly when $f$ is receptive and courteous. Its classifying presheaf has fibre
\[
F_f(x)=\{q\in\Conf_f(S):f(q)=x\}.
\]
Its cubical unstraightening is $K(S)$ over $K(E)$, and it is reached from the empty configuration.
\end{theorem}
\begin{proof}\vmode
\leavevmode
\begin{enumerate}[label=(\roman*)]
\item \textbf{Scott fibration and fibre description.}
The equivalence of receptivity and courtesy with the Scott discrete-fibration condition is \cite[Theorem~3.20]{ccrw}; Theorem~\ref{thm:eventscott} identifies its Scott orders with the ordinary categories used here.
The category-of-elements construction of Proposition~\ref{prop:kernel-elements} then gives exactly the displayed fibres.

\item \textbf{Terminal cube lift.}
For the cubical assertion, Lemma~\ref{lem:terminalcorner} identifies a lifted arena cube with its terminal strategy configuration.

\item \textbf{Atomic restrictions.}
An atomic Scott restriction adds its unique negative occurrence or removes its unique positive occurrence.
In the removal case the occurrence is maximal: any immediate successor would, by courtesy, force a successor of its image, contradicting maximality in the arena configuration.

\item \textbf{Identification of lifted faces.}
Local injectivity and Lemma~\ref{lem:orderreflection} identify the named occurrences in all faces.
Thus the lifted cube consists precisely of the configurations obtained by the corresponding independent event additions, and is a cell of $K(S)$.
Conversely every cell of $K(S)$ supplies this lift.

\item \textbf{Reachability.}
Finally every finite strategy configuration has a linearization from the empty configuration; its negative steps are restrictions and its positive steps are reverse positive restrictions.
Hence every fibre element is reached.

\item \textbf{Explicit cubical comparison.}
Writing $p_f:E_{F_f}\to\Sc K(E)$ for the element projection, the comparison fits the pullback square
\[
\begin{tikzcd}[column sep=large]
K(S) \arrow[r,"{\cong}"] \arrow[d,"{K(f)}"']
 &\Un(F_f) \arrow[d]\\
K(E) \arrow[r,equal] &K(E).
\end{tikzcd}
\]
On an $F$-cell the upper map sends its configuration family to the unique lift determined by its Scott-terminal corner.
The inverse reads the events from the lifted one-coordinate faces, as in the proof.
\end{enumerate}
\end{proof}

Section~\ref{subsec:pc-secured-proof} constructs interaction and hiding from finite matching configurations, proves its comparison with the reached set-valued coend, and verifies associativity and compact duality. In particular Theorem~\ref{thm:ordinarycomp} identifies the composite kernel directly, and Theorem~\ref{thm:ordinary-category} proves closure of this strategy class under composition. These are the ordinary configuration constructions studied in~\cite{ccrw,profunctors}.


\begin{corollary}[Kernel of the causal-interface display]\label{cor:causal-interface-kernel}
The strategy of Lemma~\ref{lem:causal-interface} has the reached kernel and display
\begin{equation}\label{eq:causal-interface-kernel}
 \Phi_S(x)=\{*\mid x\in\Conf_f(S)\},\qquad \Un(\Phi_S)\cong K(S).
\end{equation}
\end{corollary}
\begin{proof}\vmode
\leavevmode
\begin{enumerate}[label=(\roman*)]
\item \textbf{Configuration fibres.}
The identity event map has one configuration lift or none.

\item \textbf{Unstraightening and reachability.}
Theorem~\ref{thm:scott-discrete-fibrations} identifies these fibres and their unstraightening; finite linearizations prove reachability.
\end{enumerate}
\end{proof}
\subsection{Resource charts and finite histories}\label{subsec:local-resource-charts}
The resource calculations must retain the witnesses that license independent updates. Resource charts record those choices and their exchanges, supplying the finite histories used by the rooted displays below.

The rules construct concurrency through the anchored instructions and finite coefficients of Definitions~\ref{def:anchored-instruction}--\ref{def:resource-coefficient}, using the face category of Definition~\ref{def:formal-labelled-cube-faces-and-cubical-arenas}. The coefficient test proves their cubical equations uniformly from the typed witness data of Section~\ref{sec:protocol-games}.

\begin{proposition}[Constructed resource cells]\label{prop:resource-charts}
The coefficients are a presheaf, with
\[
 (u,\alpha)^*(q,(a_g)_{g\in G})=(q_B,(a_{u(f)})_{f\in F}),
 \qquad B=\alpha^{-1}(1).
\]
Its vertices and edges are exactly the states and selected rule records. With $Q_G=\mathsf C_\square(-,G)$ there is a canonical isomorphism
\[
 X_N=\colim_{(G,c)\in\int\mathcal C_N}Q_G\cong\mathcal C_N.
\]
Independent medium factors give the Day tensor, coherently with all faces and witness maps.
\end{proposition}
\begin{proof}\vmode
\leavevmode
\begin{enumerate}[label=(\roman*)]
\item \textbf{Residual coefficient test.}
A remaining permutation after executing $B$ is a segment of a full permutation, so is defined.
Its subset-prefix state is $q_{A\cup B}$; omission and execution therefore preserve the coefficient test.

\item \textbf{Face equations.}
Successive cofaces combine the same executed subsets and reindex the same instructions, proving the face equations.

\item \textbf{Singleton rules.}
Every singleton passes exactly when its rule is enabled.

\item \textbf{Colimit comparison.}
A representative $(c,\phi:F\to G)$ in the colimit maps to $\phi^*c$; its inverse is $(c,1_F)$.
These maps are inverse by the colimit relation.

\item \textbf{Independent-factor decomposition.}
For independent factors a chart partitions uniquely between them; this is the displayed Day formula, and iterated partitions prove coherence.

\item \textbf{Subset-state naturality.}
\[
\begin{tikzcd}[column sep=large]
\mathcal C_N(G)\arrow[r,"{c\mapsto(q_A)}"]\arrow[d,"{(u,\alpha)^*}"']
 & S^{\mathcal P(G)}\arrow[d,"{(q_A)\mapsto(q_{B\cup uA})}"]\\
\mathcal C_N(F)\arrow[r,"{c\mapsto(q_A)}"']&S^{\mathcal P(F)} .
\end{tikzcd}
\]
The upper map records the actual vertices of a chart; all edge equations use the original selected witnesses.
\end{enumerate}
\end{proof}

\begin{definition}[Direction-label quotient]\label{def:direction-label-quotient}
For a symmetric precubical carrier $X$, let $L(X)=X_1/{\sim}$, where $\sim$ is the equivalence relation generated by pairs of opposite edges of two-cells. For a set $\Lambda$, define the \textbf{cofree direction-labelled carrier} $!\Lambda$ by $(!\Lambda)_F=\Lambda^F$, with face maps given by restriction along their direction injections. The notation $\operatorname{Nat}(X,Y)$ denotes natural transformations between the two presheaves.
\end{definition}

\begin{proposition}[Direction labels]\label{prop:direction-labels}
Let $L(X)$ be the quotient of $X_1$ generated by identifying opposite edges of each square. For $(!\Lambda)_F=\Lambda^F$ there is a natural bijection
\[
 \Hom(L(X),\Lambda)\cong\operatorname{Nat}(X,!\Lambda).
\]
In particular every resource chart has coherent direction labels, while its exact emitted operational observation remains attached to its edges.
\end{proposition}
\begin{proof}\vmode
\leavevmode
\begin{enumerate}[label=(\roman*)]
\item \textbf{Direction classes.}
Any two parallel edges of a cube are joined by successive square faces, so their quotient classes agree.

\item \textbf{Canonical cubical labelling.}
Labelling each direction by this class defines a cubical map $X\to!L(X)$.

\item \textbf{Universal factorization.}
Any cubical labelling is constant on opposite edges and factors uniquely through the quotient.

\item \textbf{Naturality and scope.}
The constructions are inverse and commute with cubical maps.
Position-dependent observations, such as a loss index changing after another deletion, need not factor; their original witnesses remain in $X_1$.
\end{enumerate}
\end{proof}

\subsection{Universal rooted history displays}\label{subsec:universal-rooted-history-displays}
A retained history presentation has a distinguished initial history. The following construction uses that root to obtain an exact history kernel and its comparison maps, including the concurrent cells carried by the presentation.

\subsubsection{The exact cubical history kernel}
\begin{definition}[Resource execution and audit categories]\label{def:formal-resource-execution-and-audit-categories}
Let $X_\Pi$ be the rooted occurrence-refined resource-chart automaton of Proposition~\ref{prop:resource-charts}, with \textbf{initial state} $q_0$. Let $D_\Pi$ be the category generated by its execution edges, subject to the equations between the two boundary paths of each square. Thus $D_\Pi(q,q')$ consists of witnessed finite paths modulo adjacent exchanges licensed by actual charts. Give the same carrier its derived direction labels from Proposition~\ref{prop:direction-labels} and the constant negative sign, obtaining the audit arena $X_\Pi^{\mathrm{aud}}$.
\end{definition}

\begin{theorem}[Uniform protocol audit strategy]\label{thm:protocol-audit}
The Scott category of the all-negative audit arena is canonically $D_\Pi^\op$. Its representable presheaf
\begin{equation}\label{eq:protocol-audit-kernel}
 \mathcal H_\Pi(q)=D_\Pi(q_0,q)
       =\Sc(X_\Pi^{\mathrm{aud}})(q,q_0),
 \qquad \rho=1_{q_0},
\end{equation}
is a reached kernel from the tensor unit. Its cubical display has the formula
\begin{equation}\label{eq:protocol-audit-cells}
 \Un(\mathcal H_\Pi)_F
 \cong\coprod_{c\in(X_\Pi)_F}D_\Pi(q_0,c_{\mathbf0}).
\end{equation}
Projection to $X_\Pi$ uniquely lifts each witnessed edge and each chart from a history over its execution-initial vertex. It retains all finite executions and all higher cells, with exactly the declared adjacent-exchange identifications of history endpoints. Independent protocol factors give a coherent tensor comparison of these kernels.
\end{theorem}
\begin{proof}\vmode
\leavevmode
\begin{enumerate}[label=(\roman*)]
\item \textbf{Reversed execution category.}
The generators-and-relations calculation of Section~\ref{sec:game-calculus} reverses every execution edge and reverses every square equation.
It therefore gives $D_\Pi^\op$.

\item \textbf{Audit restrictions and reachability.}
A restriction in \eqref{eq:protocol-audit-kernel} appends the corresponding execution path.
Every element is reached from the identity by such restrictions.

\item \textbf{Terminal-corner cell formula.}
All directions are negative, so the Scott-terminal vertex of a cube is its execution-initial corner.
Lemma~\ref{lem:terminalcorner} gives \eqref{eq:protocol-audit-cells}.

\item \textbf{Faces and unique lifting.}
At a lower face retain the history and restrict the chart.
At an upper face append that direction to the history before restricting.
Appending a subset of directions in different orders gives the same arrow of $D_\Pi$, by the chart squares, proving the face equations and the unique-lifting assertion.

\item \textbf{Retention of higher cells.}
The displayed formula retains the base higher cell itself, so using square equations in the category does not identify distinct higher cells.

\item \textbf{Independent-factor products.}
For independent factors, the execution category is the product: generators from separate factors commute, and a path is determined modulo these exchanges by its two projected paths.
Hom sets out of the root therefore split as products.
The cell comparison is the corresponding partition of direction sets, and all coherence maps preserve the same tuples.
\end{enumerate}
\end{proof}

Theorem~\ref{thm:protocol-audit} applies to every medium in the protocol instances in Appendix~\ref{app:communication-examples}, to their mixed networks, and to the alternating-bit network. The all-negative audit interface puts enabling and operational alternatives in its arena. Definition~\ref{def:abstract-communication-protocol} specifies the ordinary strategy and its retained presentation; each operational comparison identifies its declared exposed boundary.

\subsubsection{The universal root and naturality of realization}
\begin{theorem}[Universal rooted history display]\label{thm:history-universal}
In $[D_\Pi,\Set]$ there is a natural bijection
\begin{equation}\label{eq:root-yoneda}
 \operatorname{Nat}(D_\Pi(q_0,-),F)\cong F(q_0).
\end{equation}
Consequently $(\mathcal H_\Pi,1_{q_0})$ is initial among functors with a chosen element at $q_0$, with maps preserving that element. Equivalently, it is the universal rooted presheaf on $\Sc(X_\Pi^{\mathrm{aud}})$, and its unstraightening is the cubical history display of Theorem~\ref{thm:protocol-audit}.
\end{theorem}
\begin{proof}\vmode
\leavevmode
\begin{enumerate}[label=(\roman*)]
\item \textbf{Root-induced natural transformation.}
Send $r\in F(q_0)$ to the transformation with component $h\mapsto F(h)r$.
Naturality follows from $F(uh)=F(u)F(h)$.

\item \textbf{Inverse evaluation and uniqueness.}
Evaluating at $1_{q_0}$ recovers $r$.
Conversely naturality along $h:q_0\to q$ forces every transformation to have this formula.
Thus the bijection and pointed initiality hold.

\item \textbf{Scott presheaf identification.}
The reversal $\Sc(X_\Pi^{\mathrm{aud}})=D_\Pi^\op$ identifies these covariant functors with Scott presheaves.

\item \textbf{Unstraightened cell formula.}
Unstraightening then has the cell formula already proved, retaining the base cube as well as its initial history.

\item \textbf{Pointed comparison square.}
For $u:q\to q'$, the pointed universal map $\alpha^r_q(h)=F(h)r$ obeys
\[
\begin{tikzcd}[column sep=large]
D_\Pi(q_0,q) \arrow[r,"{\alpha^r_q}"] \arrow[d,"{u\circ(-)}"'] &F(q) \arrow[d,"{F(u)}"]\\
D_\Pi(q_0,q') \arrow[r,"{\alpha^r_{q'}}"'] &F(q').
\end{tikzcd}
\]
Its value at the identity is $r$.
These two requirements force every component and hence prove uniqueness of the map of pointed kernels.
\end{enumerate}
\end{proof}

\begin{definition}[Chart-preserving specification maps]\label{def:formal-chart-preserving-specification-maps}
For naturality, retain specification maps on component states, operation witnesses, and resource data. Such a map induces a rooted map of witness graphs. Call it \textbf{chart-preserving} when it sends every \textbf{anchored instruction} family passing the resource test to the transported family passing that test, with the same direction set. This condition is an explicit preservation condition on the computed charts. An arbitrary graph map can identify two consumed resources, so need not have this property. In particular, injective renaming of resource and participant identifiers, with commuting guards and updates, preserves the tests: disjointness remains disjointness, every executable permutation maps to that permutation, and the equalities at its subset prefixes are preserved.
\end{definition}

\begin{proposition}[Naturality on displayed specifications]\label{prop:spec-naturality}
A root-preserving, label-preserving, chart-preserving specification map induces a cubical map $g:X_\Pi\to X_{\Pi'}$, an execution functor $d:D_\Pi\to D_{\Pi'}$, and a pointed transformation
\[
 \mathcal H_\Pi\longrightarrow d^*\mathcal H_{\Pi'},
 \qquad h\longmapsto d(h).
\]
Its displayed cell map is $(c,h)\mapsto(g(c),d(h))$. These assignments preserve identity maps and composition and commute with the operational realization maps.
\end{proposition}
\begin{proof}\vmode
\leavevmode
\begin{enumerate}[label=(\roman*)]
\item \textbf{Cubical specification map.}
The transported direction family and its subset states define $g$; omission and execution of a direction commute with the specification map, so its faces agree.

\item \textbf{Square and direction-label preservation.}
Every square relation therefore maps to a square relation.
Proposition~\ref{prop:direction-labels} supplies the induced direction-alphabet map $L(g)$, preserving the constant audit polarity.

\item \textbf{Path functor and pointed transformation.}
The path functor descends to $d$.
It sends $1_{q_0}$ to $1_{q'_0}$, and $d(hu)=d(h)d(u)$ gives the pointed transformation.

\item \textbf{Cubical display comparison.}
In the display formula, lower faces retain the history and upper faces append the corresponding direction.
Both commute with $(g,d)$, proving the cubical statement.

\item \textbf{Functoriality and retained witnesses.}
Identities and composites act on the same witnesses and charts.
The selected offer and medium records are retained by these same componentwise maps.

\item \textbf{Explicit coface naturality.}
For any coface $u:F\to G$, the explicit cell comparison $\widehat d_F(c,h)=(g(c),d(h))$ satisfies
\[
\begin{tikzcd}[column sep=large,row sep=large]
\Un(\mathcal H_\Pi)_G \arrow[r,"{\widehat d_G}"] \arrow[d,"{u^*}"']
 &\Un(\mathcal H_{\Pi'})_G \arrow[d,"{u^*}"]\\
\Un(\mathcal H_\Pi)_F \arrow[r,"{\widehat d_F}"']
 &\Un(\mathcal H_{\Pi'})_F .
\end{tikzcd}
\]
At an upper face the equality is $d(uh)=d(u)d(h)$; at a lower face it is restriction of $g(c)$.
Repeated faces therefore give the same comparison in every dimension.
\end{enumerate}
\end{proof}

\begin{remark}[Audit and exposed communication]
The root universal property is relative to the execution arena of the specified protocol. The all-negative polarization places operational alternatives and enabling in that arena. Public communication arenas have their specified input/output polarities and may hide internal actions or forget witnesses. Their strategy kernels are calculated by Theorem~\ref{thm:scott-discrete-fibrations}. A universal rooted audit and a public communication strategy answer different interface questions; the comparison between them is the explicit hiding, refinement, or observation map stated for each protocol.
\end{remark}

\subsection{Modes, transposition, and composition}\label{subsec:evidence-mode-calculus}
The remaining calculations identify the evidence modes and their transposes on the specified carriers. Their units, counits and mate equations make explicit how these operations compose and how the resulting rules are interpreted.

\begin{definition}[Modes, locks, and keys]\label{def:admin-mode}
An \textbf{evidence mode} is a constructed adjunction $L_\mu:\mathcal C_m\rightleftarrows\mathcal C_n:R_\mu$ between specified categories of tests or evidence. The graph construction supplies $\Sigma_f\dashv f^*$; dependent products supply $f^*\dashv\Pi_f$ on evidence. Identity and composite modes use identity and composite adjunctions. The \textbf{lock} of a context object $\Gamma\in\mathcal C_m$ along $\mu$ is $L_\mu\Gamma$. A \textbf{key} between parallel modes is a natural transformation $\theta:L_\mu\Rightarrow L_\nu$. Such a transformation is constructed from a given comparison of index maps or tests; the definition does not supply comparisons that the operational data lack.
\end{definition}

\paragraph{Natural transposition and mates.}
\begin{definition}[Modal transposition operations]\label{def:formal-modal-transposition-operations}
For a constructed adjunction $L_\mu\dashv R_\mu$, define the natural hom-set bijection
\[
 \operatorname{tr}_{\mu,\Gamma,B}:\Hom_{\mathcal C_n}(L_\mu\Gamma,B)
 \longrightarrow\Hom_{\mathcal C_m}(\Gamma,R_\mu B),
 \qquad t\longmapsto R_\mu(t)\eta^\mu_\Gamma.
\]
Its inverse is $u\mapsto\epsilon^\mu_B L_\mu(u)$. The two composites are identities by naturality and the unit--counit triangles. Subscripts are suppressed when determined by the types. The associated term constructors are defined in Section~\ref{subsec:calculus-evidence}.
\end{definition}

\begin{lemma}[Key transposition and substitution]\label{lem:admin-key}
The unique mate $\theta^\flat:R_\nu\Rightarrow R_\mu$ satisfying
\[
 \operatorname{tr}_\mu(t\theta_\Gamma)
       =\theta^\flat_B\operatorname{tr}_\nu(t),
 \qquad t:L_\nu\Gamma\to B,
\]
is
\[
 \theta^\flat_B=
 R_\mu(\epsilon^\nu_B)\,
 R_\mu(\theta_{R_\nu B})\,
 \eta^\mu_{R_\nu B}.
\]
Mates reverse vertical composition, respect identities, and are compatible with whiskering and substitution.
\end{lemma}
\begin{proof}\vmode
\leavevmode
\begin{enumerate}[label=(\roman*)]
\item \textbf{Forced mate formula.}
Set $\Gamma=R_\nu B$ and $t=\epsilon^\nu_B$ in the required equation.
Its $\nu$-transpose is the identity, forcing the displayed formula.

\item \textbf{Transposition equation.}
Every $t$ is $\epsilon^\nu_B L_\nu(\bar t)$ for its transpose $\bar t$.
Naturality of $\theta$ and of the unit moves $\bar t$ through that formula, proving the required equation for every $t$.

\item \textbf{Invertibility.}
The triangle identities make transposition invertible.

\item \textbf{Composition and substitution.}
Stacking two such transposition equations proves reversal of composition by uniqueness.
Identity, whiskering, and precomposition with a substitution are proved by the same equation and naturality.

\item \textbf{Explicit hom-set square.}
The complete mate equation is the commuting hom-set square
\[
\begin{tikzcd}[column sep=large,row sep=large]
\Hom(L_\nu\Gamma,B) \arrow[r,"{(-)\theta_\Gamma}"] \arrow[d,"{\operatorname{tr}_\nu}"']
 &\Hom(L_\mu\Gamma,B) \arrow[d,"{\operatorname{tr}_\mu}"]\\
\Hom(\Gamma,R_\nu B) \arrow[r,"{\theta^\flat_B(-)}"']
 &\Hom(\Gamma,R_\mu B).
\end{tikzcd}
\]
Naturality in $\Gamma$ is exactly compatibility with substitution of evidence contexts.
\end{enumerate}
\end{proof}

The evidence and response adjunctions constructed above, together with the temporal adjunction already constructed in Section~\ref{sec:execution-games}, supply modes. The following subsections construct policy and occurrence-interface modes and adjoin them by the following same operation.

\begin{proposition}[Composition and keys of the constructed modes]\label{prop:multimodal-composition}
Closing these constructed adjunctions and keys under identities, composition, and whiskering gives a $2$-category of modes. For $L_\mu\dashv R_\mu$ and $L_\nu\dashv R_\nu$, the composite has left part $L_\nu L_\mu$, right part $R_\mu R_\nu$, and structure maps
\[
 \eta^{\nu\mu}=R_\mu\eta^\nu L_\mu\,\eta^\mu,
 \qquad
 \epsilon^{\nu\mu}=\epsilon^\nu\,L_\nu\epsilon^\mu R_\nu.
\]
Key transposition is compatible with this composition and with precomposition in the hom-set context variable.
\end{proposition}
\begin{proof}\vmode
\leavevmode
\begin{enumerate}[label=(\roman*)]
\item \textbf{Composite adjunction.}
The two hom-set adjunction bijections compose, yielding the displayed unit and counit.
Their inverse equations follow from those of the original bijections, and hence give both triangles.

\item \textbf{Associativity and units.}
With three modes either association composes the same three bijections; the expanded units and counits are equal by functoriality.
Identity adjunctions give the units.

\item \textbf{Key composition and mates.}
Keys compose vertically as natural transformations and horizontally by whiskering; naturality gives the interchange equation.
Lemma~\ref{lem:admin-key} proves the corresponding mate equations.

\item \textbf{Substitution compatibility.}
Naturality of each hom-set bijection in its context object proves compatibility with precomposition by a substitution; composing these natural bijections preserves that equation.

\item \textbf{Administrative instance.}
The administrative interpretation of Theorem~\ref{thm:admin-soundness} instantiates this calculation.
All these operations use the already constructed adjunctions and maps.
\end{enumerate}
\end{proof}

\subsubsection{Rules for the constructed evidence adjunctions}\label{subsec:calculus-evidence}
\begin{definition}[Transposition rules]\label{def:evidence-transposition-rules}
For a constructed adjunction $L_\mu\dashv R_\mu$, the rule $\mathsf{shut}_\mu$ transposes a map $t:L_\mu\Gamma\to B$ to $R_\mu(t)\eta_\Gamma:\Gamma\to R_\mu B$. The inverse rule $\mathsf{open}_\mu$ sends $u$ to $\epsilon_B L_\mu(u)$. A key $\theta:L_\mu\Rightarrow L_\nu$ acts on the left by precomposition and on the right by its mate from Lemma~\ref{lem:admin-key}. Context substitution is precomposition by an actual parameter map. These are rules for the constructed evidence maps, with their stated source and target categories.
\end{definition}
\begin{proposition}[Transposition computation and substitution]\label{prop:transposition-computation-substitution}
Open and shut are inverse, commute with context substitution, and satisfy the mate equation for keys and their composites.
\end{proposition}
\begin{proof}\vmode
\leavevmode
\begin{enumerate}[label=(\roman*)]
\item \textbf{Beta and eta equations.}
The two adjunction triangles prove beta and eta.

\item \textbf{Substitution.}
Naturality of the hom-set bijection proves substitution.

\item \textbf{Key equation and coherence.}
The mate square of Lemma~\ref{lem:admin-key} proves the key equation, and its uniqueness proves compatibility with identity, composition, and whiskering.
\end{enumerate}
\end{proof}
\subsection{Parameter families and closed resource structure}\label{subsec:parameter-families-closed-resources}
Parameter families let the resource constructions vary over specified inputs. The final calculations identify their closed structure and the interaction with runtime-key parameters, retaining the owned occurrences required by the resource rules.

Parameter families and closed resource structure are constructed from the declared parameter sets, inventories, and witnessed policy paths. For a parameter set $\Gamma$, use families $A(\gamma,m)$ indexed by $\Gamma\times M_v$, and write $\mathcal R_v(\Gamma)=\Set^{\Gamma\times M_v}$. Policy-evidence families are $\Gamma$-indexed presheaves on $\mathcal A_v$, denoted $\mathcal S_v(\Gamma)$. Parameter substitution along $\theta:\Gamma'\to\Gamma$ is literal precomposition. These families describe evidence over the protocol; they are not additional public programs.

\begin{theorem}[Parameter substitution and closed resources]\label{thm:parameter-substitution-closed-resources}
At each parameter $\gamma$, the formulas
\[
 (A\otimes_vB)(\gamma,m)=\coprod_{a+b=m}A(\gamma,a)\times B(\gamma,b),
 \qquad I_v(\gamma,m)=\{*\mid m=0\},
\]
\begin{equation}\label{eq:parameterized-resource-hom-policy-monad}
 (A\multimap_v B)(\gamma,m)=\prod_n B(\gamma,m+n)^{A(\gamma,n)},
 \quad (T_v A)(\gamma,m)=\coprod_n\mathcal A_v(m,n)\times A(\gamma,n)
\end{equation}
give a closed symmetric monoidal category and the existing policy monad. Parameter substitution preserves these displayed family constructions strictly. Free policy evidence, Day convolution, and their adjunctions are computed at each parameter and commute with the same substitution.
\end{theorem}
\begin{proof}\vmode
\leavevmode
\begin{enumerate}[label=(\roman*)]
\item \textbf{Tensor--hom adjunction.}
A map $C\otimes_vA\to B$ is exactly a family of maps
$C(\gamma,m)\times A(\gamma,n)\to B(\gamma,m+n)$.
Currying each component gives $C\to A\multimap_v B$, and evaluation gives the inverse.

\item \textbf{Monoidal coherence.}
Reassociation and symmetry retain the same grade partition and entries.

\item \textbf{Policy unit and multiplication.}
Unit and multiplication insert identity paths and concatenate retained paths as in Theorem~\ref{thm:policy-evidence-adjunction-monad}.

\item \textbf{Parameter substitution.}
All these definitions evaluate at the same $\gamma$, so precomposition by $\theta$ commutes with them and with successive substitutions.

\item \textbf{Presheaf comparison.}
The presheaf coends and ends are likewise calculated at each parameter; the same pointwise argument gives their comparison.
\end{enumerate}
\end{proof}
For the following calculations fix a parameter and abbreviate $\PubRes_v=\Set^{M_v}$, $\PolicyRes_v=\Set^{\mathcal A_v^{\op}}$, and $\PolicyCat_v=\mathcal A_v$. The parameterized versions are these formulas applied at each $\gamma$. For actual finite inventories, tensor additionally retains its chosen partition into owned occurrences; forgetting the occurrence names gives the displayed grade split.
\begin{theorem}[Policy adjunction and closed comparison]\label{thm:policy-adjunction-closed-comparison}
Let \(i:M_v\to\PolicyCat_v\) be the identity on objects. Then
\[
 F=\operatorname{Lan}_{i^{\op}}:\PubRes_v\rightleftarrows\PolicyRes_v:U=i^*,
 \qquad
 (FA)(m)=\coprod_n\PolicyCat_v(m,n)\cdot A(n).
\]
The induced monad \(T=UF\) has identity-path unit and path-composition
multiplication. On \(\PolicyRes_v\), Day convolution is closed, with
\[
 (P\star Q)(m)=\int^{a,b}\PolicyCat_v(m,a+b)\cdot(P(a)\times Q(b)),
 \quad
 (P\multimap_\star Q)(m)=\int_nQ(m+n)^{P(n)}.
\]
The functor \(F\) is strong symmetric monoidal, and
\[
 \boxed{U(FA\multimap_\star B)\cong A\multimap UB.}
\]
\end{theorem}
\begin{proof}\vmode
\leavevmode
\begin{enumerate}[label=(\roman*)]
\item \textbf{Policy adjunction and monoidal comparison.}
Theorem~\ref{thm:policy-evidence-adjunction-monad} and Proposition~\ref{prop:admin-monoidal-kan} give the displayed adjunction, its monad and its strong monoidal comparison, including coherence.

\item \textbf{Closed structure.}
For closedness, maps out of the convolution coend are natural families $P(a)\times Q(b)\to Z(a+b)$.
Currying and imposing naturality in $b$ by the end give the displayed internal hom, with evaluation as inverse.

\item \textbf{Closed comparison isomorphism.}
Finally
\[
\begin{aligned}
 \Hom(C,U(FA\multimap_\star B))
 &\cong\Hom(FC\star FA,B)\\
 &\cong\Hom(F(C\otimes A),B)
 \cong\Hom(C,A\multimap UB).
\end{aligned}
\]
The bijections are natural; Yoneda gives the boxed comparison.
\end{enumerate}
\end{proof}

The framing operation gives
\(A\otimes TB\to T(A\otimes B)\). Its two parallel orders agree by the
constructed policy exchange diamonds, so \(T\) is a commutative strong
monad for this tensor. These are not claims of preservation of ordinary
cartesian implication by \(F\) or \(T\). In particular,
\(F1(m)=\coprod_n\PolicyCat_v(m,n)\) need not be terminal.

\begin{theorem}[A distributive law for keys and policy paths]\label{thm:key-policy-distributive-law}
For a supplied runtime-key family \(K_\gamma\), work at each parameter \(\gamma\) and put \(K=K_\gamma\). Let
\(S(A)(m)=A(m)^K\). There is a distributive law
\[
 \lambda:TS\longrightarrow ST,
 \qquad
 \lambda(p,f)=\bigl(k\mapsto(p,f(k))\bigr).
\]
Give \(S=(-)^K\) its constant-function unit and diagonal-evaluation multiplication. Consequently \(ST\) is a monad. The map is compatible with resource
framing and the parallel policy comparison. Inverting it is not part of
the construction.
\end{theorem}
\begin{proof}\vmode
\leavevmode
\begin{enumerate}[label=(\roman*)]
\item \textbf{Construction and naturality of the law.}
For each path summand, curry the evaluation map
\(K\times A(n)^K\to A(n)\) followed by its injection into \(TA(m)\).
The coproduct universal property constructs \(\lambda\), and these maps
are natural.

\item \textbf{Distributive-law unit equations.}
The four distributive-law equations are checked as follows.
For the \(T\)-unit both sides give \(k\mapsto(\id,f(k))\); for the
\(S\)-unit both give the constant function \(k\mapsto(p,a)\).

\item \textbf{Policy-path multiplication.}
For the \(T\)-multiplication, an element represented by paths \(p,q\)
and a key function gives \(k\mapsto(qp,f(k))\) by either route.

\item \textbf{Key multiplication.}
For the \(S\)-multiplication, an element \((p,F)\) with two key arguments
gives \(k\mapsto(p,F(k)(k))\) by either route.

\item \textbf{Componentwise equality of arrows.}
These formulas are componentwise: equality after evaluation and coproduct injection proves
equality of the actual arrows, without global points or a choice of path
representatives.

\item \textbf{Composite monad.}
The composite multiplication is
\[
 STST\xrightarrow{S\lambda T}SSTT
       \xrightarrow{\mu^S_{TT}}STT
       \xrightarrow{S\mu^T}ST.
\]
The four equations just checked and the two monads' equations prove its
unit and associativity, equivalently the standard distributive-law
construction~\cite{doct-beck}.

\item \textbf{Resource strength and framing.}
Resource strength for \(S\) retains a chosen
split and evaluates the key-dependent entry at the supplied key.
Every
comparison with \(\lambda\) therefore retains the same split, framed
path, and values evaluated at the same key.
This proves compatibility
with framing and parallel policy execution.
\end{enumerate}
\end{proof}

The distinction between \(TS\) and \(ST\) has content: the former chooses
its policy path before the key argument is supplied; the latter can have
a key-dependent path. The construction supplies a direction of comparison,
not an unproved exchange isomorphism. Similarly, a change of decision
view uses its actual maps or spans; there is no automatic monad morphism
between every pair \(T_v,T_w\).

\paragraph{Owned occurrences.}
The graded resource tensor records types and multiplicities. For an actual
finite inventory \(I\), retain the witness
\((I_1,I_2)\in\mathsf{Split}(I)\). Its type-forgetting map sends it to
\((|I_1|,|I_2|)\). Reassociation retains the same ordered partition into
three parts, so this map commutes with tensor reassociation and symmetry.
The logical tensor does not select a physical partition from its type
grade, and a logical path exchange is used operationally only through the
constructed proved worker and causal comparisons.

The key--policy map is stable under substitution: every path injection,
evaluation and currying in its construction is preserved by the componentwise
parameter substitutions in Theorem~\ref{thm:parameter-substitution-closed-resources}. The family interpretation makes this equality literal at each parameter.


\begin{thebibliography}{99}
\bibitem{effective-domains}
S. Abramsky and A. Jung.
\emph{Domain Theory}.
In \emph{Handbook of Logic in Computer Science}, vol. 3, Clarendon Press, pp. 1--168, 1994.
\href{https://www.cs.ox.ac.uk/people/samson.abramsky/handbook.pdf}{Corrected and expanded author version}.

\bibitem{effective-dominances}
M. H. Escard\'o and C. M. Knapp.
\emph{Partial Elements and Recursion via Dominances in Univalent Type Theory}.
CSL 2017, LIPIcs 82, pp. 21:1--21:16, 2017.
\href{https://doi.org/10.4230/LIPIcs.CSL.2017.21}{doi:10.4230/LIPIcs.CSL.2017.21}.

\bibitem{effective-lifting}
J. Sterling.
\emph{Tensorial Structure of the Lifting Doctrine in Constructive Domain Theory}.
Accepted author manuscript, 2024.
\href{https://doi.org/10.17863/CAM.108169}{doi:10.17863/CAM.108169}.

\bibitem{ccrw}
S. Castellan, P. Clairambault, S. Rideau, and G. Winskel.
\emph{Games and Strategies as Event Structures}.
Logical Methods in Computer Science 13(3:35), 2017.
\href{https://doi.org/10.23638/LMCS-13(3:35)2017}{doi:10.23638/LMCS-13(3:35)2017}.
Author version: \url{https://pageperso.lis-lab.fr/pierre.clairambault/papers/cg1.pdf}.

\bibitem{profunctors}
G. Winskel.
\emph{Strategies as Profunctors}.
FoSSaCS, Lecture Notes in Computer Science 7794, pp. 418--433, 2013.
\href{https://doi.org/10.1007/978-3-642-37075-5_27}{doi:10.1007/978-3-642-37075-5\_27}.
Author version: \url{https://www.cl.cam.ac.uk/~gw104/fossacs2013_submission_12.pdf}.

\bibitem{cy}
S. Castellan and N. Yoshida.
\emph{Two Sides of the Same Coin: Session Types and Game Semantics. A Synchronous Side and an Asynchronous Side}.
Proceedings of the ACM on Programming Languages 3(POPL), Article 27, 2019.
\href{https://doi.org/10.1145/3290340}{doi:10.1145/3290340}.

\bibitem{winskel-functional}
G. Winskel.
\emph{Making Concurrency Functional}.
Proceedings of LICS, pp. 1--14, 2023.
\href{https://doi.org/10.1109/LICS56636.2023.10175727}{doi:10.1109/LICS56636.2023.10175727}.
Extended author version: \url{https://arxiv.org/abs/2202.13910}.
The theorem numbering used here is that of the extended text, also available at
\url{https://strathprints.strath.ac.uk/86947/1/Winskel_LICS_2023_Making_concurrency_functional.pdf}.

\bibitem{buffered-pi}
X. Deng, Y. Zhang, Y. Deng, and F. Zhong.
\emph{The Buffered $\pi$-Calculus: A Model for Concurrent Languages}.
Author report, 2012, Section 2 and Tables 1--3.
\url{https://arxiv.org/abs/1212.6183}.

\bibitem{abdulla-jonsson-1996}
P. A. Abdulla and B. Jonsson.
\emph{Verifying Programs with Unreliable Channels}.
Information and Computation 127(2), pp. 91--101, 1996.
\href{https://doi.org/10.1006/inco.1996.0053}{doi:10.1006/inco.1996.0053}.

\bibitem{devisme-winskel}
M. de Visme and G. Winskel.
\emph{Causal Unfoldings and Disjunctive Causes}.
Logical Methods in Computer Science 19(2:4), 2023.
\href{https://doi.org/10.46298/lmcs-19(2:4)2023}{doi:10.46298/lmcs-19(2:4)2023}.

\bibitem{parallel-causes}
M. de Visme and G. Winskel.
\emph{Strategies with Parallel Causes}.
CSL, LIPIcs 82, pp. 41:1--41:20, 2017.
\href{https://doi.org/10.4230/LIPIcs.CSL.2017.41}{doi:10.4230/LIPIcs.CSL.2017.41}.

\bibitem{jnw-open-maps}
A. Joyal, M. Nielsen, and G. Winskel.
\emph{Bisimulation from Open Maps}.
Information and Computation 127(2), pp. 164--185, 1996.
\href{https://doi.org/10.1006/inco.1996.0057}{doi:10.1006/inco.1996.0057}.
Report version: \url{https://www.brics.dk/RS/94/7/BRICS-RS-94-7.pdf}.

\bibitem{winning}
P. Clairambault, J. Gutierrez, and G. Winskel.
\emph{The Winning Ways of Concurrent Games}.
Proceedings of LICS, 2012.
Author version: \url{https://www.cl.cam.ac.uk/~gw104/lics2012_submission_14.pdf}.

\bibitem{pi-freyd}
K. Sakayori and T. Tsukada.
\emph{A Categorical Model of an i/o-typed $\pi$-calculus}.
ESOP 2019, LNCS 11423, pp. 640--667.
\href{https://doi.org/10.1007/978-3-030-17184-1_23}{doi:10.1007/978-3-030-17184-1\_23}.

\bibitem{ccw-thin}
S. Castellan, P. Clairambault, and G. Winskel.
\emph{Thin Games with Symmetry and Concurrent Hyland--Ong Games}.
Logical Methods in Computer Science 15(1:18), 2019.
\href{https://doi.org/10.23638/LMCS-15(1:18)2019}{doi:10.23638/LMCS-15(1:18)2019}.
\url{https://lmcs.episciences.org/5248}.

\bibitem{birkedal-guarded}
L. Birkedal, R. E. M\o gelberg, J. Schwinghammer, and K. St\o vring.
\emph{First Steps in Synthetic Guarded Domain Theory: Step-Indexing in the Topos of Trees}.
Logical Methods in Computer Science 8(4:1), pp. 1--45, 2012.
\href{https://doi.org/10.2168/LMCS-8(4:1)2012}{doi:10.2168/LMCS-8(4:1)2012}.

\bibitem{miliuslitak}
S. Milius and T. Litak.
\emph{Guard Your Daggers and Traces: Properties of Guarded (Co-)recursion}.
Fundamenta Informaticae 150(3--4), pp. 407--449, 2017.
\href{https://doi.org/10.3233/FI-2017-1475}{doi:10.3233/FI-2017-1475}.
Author version: \url{https://arxiv.org/abs/1603.05214}.

\bibitem{goncharovschroeder}
S. Goncharov and L. Schr\"oder.
\emph{Guarded Traced Categories}.
FoSSaCS, Lecture Notes in Computer Science 10803, pp. 313--330, 2018.
\href{https://doi.org/10.1007/978-3-319-89366-2_17}{doi:10.1007/978-3-319-89366-2\_17}.
Author version: \url{https://arxiv.org/abs/1802.08756}.

\bibitem{jsv}
A. Joyal, R. Street, and D. Verity.
\emph{Traced Monoidal Categories}.
Mathematical Proceedings of the Cambridge Philosophical Society 119(3), pp. 447--468, 1996.
\href{https://doi.org/10.1017/S0305004100074338}{doi:10.1017/S0305004100074338}.

\bibitem{shulman}
M. Shulman.
\emph{Framed Bicategories and Monoidal Fibrations}.
Theory and Applications of Categories 20(18), pp. 650--738, 2008.
\url{https://tac.mta.ca/tac/volumes/20/18/20-18abs.html}.

\bibitem{crypto-abf} M. Abadi, B. Blanchet, and C. Fournet.
\emph{The Applied Pi Calculus: Mobile Values, New Names, and Secure Communication}. Journal of the ACM 65(1), Article 1, 2017.
Extended report: \url{https://arxiv.org/abs/1609.03003}.
\bibitem{crypto-stark}
I. Stark.
\emph{Categorical Models for Local Names}.
LISP and Symbolic Computation 9(1), pp. 77--107, 1996.
\href{https://homepages.inf.ed.ac.uk/stark/catmln.pdf}{Author's full text}.

\bibitem{categorical-mealy}
G. Boccali, A. Laretto, F. Loregian, and S. Luneia.
\emph{Bicategories of Automata, Automata in Bicategories}.
Applied Category Theory 2023, EPTCS 397, pp. 1--19, 2023.
\href{https://doi.org/10.4204/EPTCS.397.1}{doi:10.4204/EPTCS.397.1}.

\bibitem{mealy-lenses}
B. Clarke.
\emph{A Diagrammatic Approach to Symmetric Lenses}.
Applied Category Theory 2020, EPTCS 333, pp. 79--91, 2021.
\href{https://doi.org/10.4204/EPTCS.333.6}{doi:10.4204/EPTCS.333.6}.

\bibitem{crypto-rutten}
J. J. M. M. Rutten.
\emph{Universal Coalgebra: A Theory of Systems}.
Theoretical Computer Science 249, pp. 3--80, 2000.
\href{https://ir.cwi.nl/pub/48}{CWI full text}.

\bibitem{crypto-bk} A. Broadbent and M. Karvonen.
\emph{Categorical composable cryptography: extended version}. Logical Methods in Computer Science 19(4:30), 2023; corrected journal version published 22 October 2024.
\url{https://lmcs.episciences.org/12718}.
Corrected author version: \url{https://arxiv.org/abs/2208.13232v5}.
The theorem numbering cited here is that of the corrected version.
\bibitem{crypto-dkp} S. Delaune, S. Kremer, and O. Pereira.
\emph{Simulation Based Security in the Applied Pi Calculus}.
FSTTCS 2009, LIPIcs 4, pp.~169--180.
\url{https://doi.org/10.4230/LIPIcs.FSTTCS.2009.2316}.

\bibitem{crypto-hayman} J. Hayman.
\emph{Interaction and Causality in Digital Signature Exchange Protocols}.
Trustworthy Global Computing, 2014.
\url{https://www.cs.le.ac.uk/events/tgc2014/pre-proceedings/hayman.pdf}.

\bibitem{crypto-bursuc} S. Bursuc.
\emph{Secure two-party computation in applied pi-calculus: models and verification}.
In \emph{Trustworthy Global Computing}, LNCS 9533, pp. 1--15, 2016.
\url{https://doi.org/10.1007/978-3-319-28766-9_1}.

\bibitem{crypto-tamarin}
The Tamarin Team.
\emph{Tamarin Prover Manual: Cryptographic Messages}.
Online manual, accessed 11 November 2026.
\url{https://tamarin-prover.com/manual/master/book/004_cryptographic-messages.html}.

\bibitem{crypto-hash-soundness}
V. Cortier, S. Kremer, R. K\"usters, and B. Warinschi.
\emph{Computationally Sound Symbolic Secrecy in the Presence of Hash Functions}.
FSTTCS 2006, Lecture Notes in Computer Science 4337, pp. 176--187.
\href{https://publ.sec.uni-stuttgart.de/cortierkremerkuesterswarinschi-eprint-218-2006.pdf}{Full version}.

\bibitem{crypto-laud}
P. Laud.
\emph{Implementing Cryptographic Primitives in the Symbolic Model}.
NASA Formal Methods 2011, Lecture Notes in Computer Science 6617, pp. 267--281.
\href{https://doi.org/10.1007/978-3-642-20398-5_20}{doi:10.1007/978-3-642-20398-5\_20}.

\bibitem{crypto-kems}
C. Cremers, A. Dax, and N. Medinger.
\emph{Keeping Up with the KEMs: Stronger Security Notions for KEMs and Automated Analysis of KEM-based Protocols}.
ACM CCS, 2024.
\href{https://doi.org/10.1145/3658644.3670283}{doi:10.1145/3658644.3670283}.

\bibitem{crypto-micciancio}
D. Micciancio.
\emph{Symbolic Encryption with Pseudorandom Keys}.
EUROCRYPT 2019, Lecture Notes in Computer Science 11478, pp. 64--93.
\href{https://eprint.iacr.org/2009/249}{Full version, Cryptology ePrint 2009/249}.

\bibitem{crypto-shamir}
A. Shamir.
\emph{How to Share a Secret}.
Communications of the ACM 22(11), pp. 612--613, 1979.
\href{https://doi.org/10.1145/359168.359176}{doi:10.1145/359168.359176}.

\bibitem{crypto-vss}
H. Zhao, M. Li, K. Sakurai, Y. Ren, J. Z. Sun, and F. Wang.
\emph{A Finite Equivalence of Verifiable Multi-secret Sharing}.
International Journal of Computational Intelligence Systems 5(1), pp. 1--12, 2012.
\href{https://doi.org/10.1080/18756891.2012.670517}{doi:10.1080/18756891.2012.670517}.

\bibitem{crypto-dahl}
M. Dahl and I. Damg\aa rd.
\emph{Universally Composable Symbolic Analysis for Two-Party Protocols Based on Homomorphic Encryption}.
EUROCRYPT 2014, Lecture Notes in Computer Science 8441, pp. 695--712.
\href{https://doi.org/10.1007/978-3-642-55220-5_38}{doi:10.1007/978-3-642-55220-5\_38}.

\bibitem{crypto-zk-applied}
M. Backes, M. Maffei, and D. Unruh.
\emph{Zero-Knowledge in the Applied Pi-calculus and Automated Verification of the Direct Anonymous Attestation Protocol}.
Dagstuhl Seminar Proceedings 07421, pp. 1--43, 2008.
\href{https://doi.org/10.4230/DagSemProc.07421.4}{doi:10.4230/DagSemProc.07421.4}.

\bibitem{crypto-zk-types}
M. Backes, C. Hri\c{t}cu, and M. Maffei.
\emph{Type-checking Zero-knowledge}.
ACM CCS, pp. 357--370, 2008.
\href{https://catalin-hritcu.github.io/publications/zk-types-ccs2008.pdf}{Author's full text}.

\bibitem{crypto-zk-malleable}
M. Backes, F. Bendun, M. Maffei, E. Mohammadi, and K. Pecina.
\emph{Symbolic Malleable Zero-knowledge Proofs}.
IEEE CSF, pp. 412--426, 2015.
\href{https://mohammadi.eu/paper/malleable_zk.pdf}{Author's full text}.

\bibitem{crypto-zk-soundness}
M. Backes and D. Unruh.
\emph{Computational Soundness of Symbolic Zero-Knowledge Proofs Against Active Attackers}.
Journal of Computer Security 18(6), pp. 1077--1155, 2010.
\href{https://doi.org/10.3233/JCS-2009-0392}{doi:10.3233/JCS-2009-0392}.

\bibitem{crypto-interactive}
M. Backes, E. Mohammadi, and T. Ruffing.
\emph{Computational Soundness for Interactive Primitives}.
ESORICS 2015, pp. 125--145.
\href{https://mohammadi.eu/paper/cosp-int.pdf}{Author's full text}.

\bibitem{crypto-mpc}
M. Backes, M. Maffei, and E. Mohammadi.
\emph{Computationally Sound Abstraction and Verification of Secure Multi-Party Computations}.
FSTTCS 2010, Leibniz International Proceedings in Informatics 8, pp. 352--363.
\href{https://mohammadi.eu/paper/smpc.pdf}{Extended full version}.

\bibitem{crypto-garbling}
B. Li and D. Micciancio.
\emph{Symbolic Security of Garbled Circuits}.
IEEE CSF, 2018.
\href{https://eprint.iacr.org/2018/141}{Full version, Cryptology ePrint 2018/141}.

\bibitem{crypto-ggm}
O. Goldreich, S. Goldwasser, and S. Micali.
\emph{How to Construct Random Functions}.
Journal of the ACM 33(4), pp. 792--807, 1986.
\href{https://doi.org/10.1145/6490.6503}{doi:10.1145/6490.6503}.

\bibitem{crypto-squirrel}
D. Baelde, S. Delaune, C. Jacomme, A. Koutsos, and J. Lallemand.
\emph{The Squirrel Prover and its Logic}.
ACM SIGLOG News 11, 2024.
\href{https://adrienkoutsos.fr/papers/siglog.pdf}{Author's full text}.

\bibitem{crypto-bu} F. B\"ohl and D. Unruh.
\emph{Symbolic Universal Composability}. IEEE Computer Security Foundations Symposium, 2013, pp. 257--271.
Extended report: \url{https://eprint.iacr.org/2013/062}.
\bibitem{garg-linear}
D. Garg, L. Bauer, K. D. Bowers, F. Pfenning, and M. K. Reiter.
\emph{A Linear Logic of Authorization and Knowledge}.
ESORICS, Lecture Notes in Computer Science 4189, pp. 297--312, 2006.
\href{https://doi.org/10.1007/11863908_19}{doi:10.1007/11863908\_19}.

\bibitem{bowers-consumable}
K. D. Bowers, L. Bauer, D. Garg, F. Pfenning, and M. K. Reiter.
\emph{Consumable Credentials in Logic-Based Access-Control Systems}.
Proceedings of NDSS, pp. 143--157, 2007.
Author version: \url{https://www.cs.cmu.edu/~fp/papers/ndss07.pdf}.

\bibitem{day-street-kan}
B. Day and R. Street.
\emph{Kan Extensions Along Promonoidal Functors}.
Theory and Applications of Categories 1(4), pp. 72--77, 1995.
\href{https://www.emis.de/ft/26705}{Journal full text (EMIS mirror)}.

\bibitem{prezado-descent}
R. Prezado.
\emph{On Effective Descent $\mathcal V$-Functors and Familial Descent Morphisms}.
Journal of Pure and Applied Algebra 228(5), Article 107597, 2024.
\href{https://doi.org/10.1016/j.jpaa.2023.107597}{doi:10.1016/j.jpaa.2023.107597}.
Revised author version: \href{https://arxiv.org/abs/2305.04042v2}{arXiv:2305.04042v2}.

\bibitem{calculus-fpt} M. Fiore, G. Plotkin, and D. Turi.
\emph{Abstract Syntax and Variable Binding}. LICS, 1999, pp.~193--202.
\url{https://www.dcs.ed.ac.uk/home/dt/FiorePlotkinTuri.pdf}.

\bibitem{calculus-cp} L. Caires and F. Pfenning.
\emph{Session Types as Intuitionistic Linear Propositions}. CONCUR, LNCS 6269, 2010, pp.~222--236.
\url{https://www.cs.cmu.edu/~fp/papers/concur10.pdf}.

\bibitem{calculus-wadler} P. Wadler. \emph{Propositions as Sessions}.
Journal of Functional Programming 24(2--3), 2014, pp.~384--418.
\url{https://doi.org/10.1017/S095679681400001X}.

\bibitem{doct-jacobs} B. Jacobs. \emph{Categorical Logic and Type Theory}.
North-Holland, 1999.
\url{https://www.cs.ru.nl/B.Jacobs/CLT/bookinfo.html}.
\bibitem{doct-beck} J. Beck. Distributive laws.
In \emph{Seminar on Triples and Categorical Homology Theory},
Lecture Notes in Mathematics 80, Springer, 1969, 119--140.
\url{https://doi.org/10.1007/BFb0083084}.
\bibitem{pub-escardo} M. H. Escard\'o. Synthetic topology of data types and classical spaces. \emph{Electronic Notes in Theoretical Computer Science} 87 (2004), 21--156. \url{https://martinescardo.github.io/papers/syntop.pdf}.
\bibitem{pub-observable} S. Abramsky. Domain theory in logical form. \emph{Annals of Pure and Applied Logic} 51 (1991), 1--77. \url{https://doi.org/10.1016/0168-0072(91)90065-T}.
\bibitem{security-riehl-fs} E. Riehl. \emph{Factorization systems}. Notes, 2008. \url{https://math.jhu.edu/~eriehl/factorization.pdf}.
\bibitem{security-yoneda} E. Riehl. \emph{Category Theory in Context}. Dover, 2016. \url{https://math.jhu.edu/~eriehl/context.pdf}.

\bibitem{ccw-mays-musts}
S. Castellan, P. Clairambault, and G. Winskel.
\emph{The Mays and Musts of Concurrent Strategies}.
In \emph{Samson Abramsky on Logic and Structure in Computer Science and Beyond}, 2023.
Author version, 2021: \url{https://arxiv.org/abs/2108.10558}.
The theorem numbering cited here is that of this author version.

\bibitem{benton-hofmann-nigam}
N. Benton, M. Hofmann, and V. Nigam.
\emph{Proof-Relevant Logical Relations for Name Generation}.
Logical Methods in Computer Science 14(1:25), 2018.
\href{https://doi.org/10.23638/LMCS-14(1:25)2018}{doi:10.23638/LMCS-14(1:25)2018}.
Author version: \url{https://arxiv.org/abs/1708.05193}.



\end{thebibliography}
\end{document}
